% Source-faithful theorem spine for the current reductive--compositional route.
% This section is deliberately concise and status-local.  The much longer
% chronological audit remains below it, but does not control the proof status.

\section{The current proof spine: what is a theorem and what remains}
\label{sec:compositional-proof-spine}

This section is the controlling mathematical account of the route.  It is
grounded in the working manuscript's Addenda XXV--XXVIII, the accompanying
formalization playbook, and the compositional proof note.  The chronological
material later in this companion records how the present formulation was
reached; it must not be read as overriding the statuses stated here.

The intended proof is not a finite catalogue of unavoidable configurations.
Its central object is an open planar cubic instance with a finite boundary
interface.  A corridor is cut into open cells, the behavior visible at a cut
is compressed to a finite profile, and equal profiles permit a middle segment
to be removed.  Thus the desired contradiction has the following form:
\[
  \text{large minimal obstruction}
  \Longrightarrow \text{long composable corridor}
  \Longrightarrow \text{repeated interface profile}
  \Longrightarrow \text{smaller obstruction}.
\]
The first and last arrows are geometry.  The middle arrow is finite-state
pigeonhole.  Keeping those three arrows separate is essential: finiteness of
the profile carrier does not itself construct a corridor, and equality of two
codes does not itself construct the spliced planar map.

\subsection{A status-local dependency chain}

The following table uses ``proved'' only for an argument actually supplied in
the cited layer.  A source assertion, a finite audit still to be run, and a
conditional theorem are different statuses.

There are now two logically distinct spines.  The direct Count spine is the
short route advertised by the compositional proof note and controls the
headline.  The older Kauffman--$\mathrm{WC}^{*}$ development is an alternative
source branch; progress there does not discharge a direct-Count premise.

\begin{center}
\small
\begin{tabular}{@{}p{0.08\textwidth}p{0.45\textwidth}p{0.34\textwidth}@{}}
\toprule
Node & Direct Count spine & Present status \\
\midrule
$D_0$ & Tait reduction and least-counterexample normal form & the genuine
spherical headline, least-counterexample selection, semantic two-sidedness,
simple-graph backing, reverse Tait construction, cyclic five-edge-connectivity,
minimum face length five, and the exact passage from a non-four-colourable
spherical presentation to a zero-Count stellar dual are machine-checked.  Only
the standard external
bridge from abstract planarity to a spherical rotation presentation remains
named conditional \\
$D_1$ & Every sufficiently large least obstruction supplies a target-forced
composable decomposition & open and the highest structural risk.  The
manuscript's linear-corridor assertion remains unsupported; in the exact
tree reformulation the missing global statement is a uniform branchwidth
bound, equivalently exclusion of one fixed wall from least obstructions.
The generic finite-tree descent, conversion of connected edge shores to
connected vertex shores without width loss, the structural physical splice,
and the mesh obstruction are machine-checked.  A boundary-separated
three-interface source rail segment now supplies a canonically oriented
saturated cut, connected edge shores, and a nonempty proper majority core.
When that source trail is obtained by deleting its distinguished edge from a
closed cubic map, the same vertex partition now supplies connected
complementary edge shores and a nonempty proper majority core in the closed
ambient map as well.
The standard rooted branch-decomposition interface and its conversion to the
literal checked consumer are machine-checked.  The Fomin--Fraigniaud--Thilikos
width-preserving connectedization theorem is now machine-checked in that
interface.  Wall exclusion and the arbitrary-length nested exact finite-state
supply are not proved \\
$D_2$ & Count across every actual ordered noose cut factors through a finite
boundary-word relation, and cumulative support is deletable on repetition &
the physical port-tangle gluing bijection, matrix law, support
determinization, exact $2^{3^k}$ support bound, repeated-block deletion, and
relabelling invariance are machine-checked.  The semantic closed composite is
also proved equivalent to Tait-colourability of the literal sewn rotation
system.  The
specialized Cell--3 closure is now an optional state-minimization and
finite-base instrument, not a correctness premise of the direct descent \\
$D_3$ & Equal cumulative rooted states admit a strictly smaller physical
deletion splice in the same bridgeless cubic spherical counterexample class,
preserving boundary order and zero accepting Count & the ordinary noose
construction, port-tangle gluing, bridgelessness proof, rotation-system to
multigraph adapter, strict vertex accounting, exact rotation-system
composition, complementary boundary order, orientation-reversed seam match,
and sphericity are machine-checked on the literal shore carriers.  Support
replacement, the structural class, strict decrease and the minimality
contradiction are now machine-checked together for equal normalized nested
majority-shore states.  The standard decomposition-to-literal-shore adapter
and the terminating digon/parallel-seam normalization are machine-checked;
the remaining obligation is global production of the supplied nested states,
which belongs to $D_1$, not to the local splice \\
$D_4$ & The route-native finite floor contains no zero-Count obstruction &
open; the refuted GWCO and $A\ne\varnothing$ systems cannot be inherited.
The exact finite-state reflection joint is now machine-checked: a trace,
reachable-set coverage, counterexample adequacy, and an empty filtered
reachable set imply the literal spherical finite-base premise.  Constructing
those three witnesses and replaying the audit remain open \\
$D_5$ & Well-founded assembly from $D_0$--$D_4$ & proved as the explicit
conditional Theorem~\ref{thm:direct-count-reductive-assembly} \\
\bottomrule
\end{tabular}
\end{center}

\begin{center}
\small
\begin{tabular}{@{}p{0.08\textwidth}p{0.45\textwidth}p{0.34\textwidth}@{}}
\toprule
Node & Alternative $\mathrm{WC}^{*}$/Lift spine & Present status \\
\midrule
$W_1$ & Kauffman closed-cubic parity and the proposed deleted-edge
prime/factored trail start & parity is proved, but the trail start is refuted
by the two-defect boundary-flux identity; this branch cannot presently start \\
$W_2$ & Meyniel rotor, digon, two-cut, and square transport & substantial
ordinary and Lean components are proved; integration is incomplete \\
$W_3$ & Pentagon Lift for the realized good-state move graph & the
four-block-edge criterion is proved below; one seed in every fibre is a
strictly stronger condition, false for arbitrary annuli and open in normal
form \\
$W_4$ & Route-specific square transport for $\mathrm{WC}^{*}$ & the
route-specific proof is written; arbitrary-target transport,
$A\ne\varnothing$, and universal SQ3 connectivity have refutation fixtures \\
\bottomrule
\end{tabular}
\end{center}

\subsection{Obligation-weighted recount after retirement of the closed-web branch}

The percentage attached to this route is now computed from the following
fixed one-hundred-point ledger.  It counts logical obligations in the direct
Count dependency chain, not files, declarations, elapsed work, or historical
effort.  A point is earned only by a sealed statement usable by that chain.
An open or refuted premise has score zero.  A conditional result with
conjunctive premises receives its assigned weight times the minimum score of
its premises; consequently one refuted premise makes its contribution zero.
This deliberately prevents a large proved implication from receiving credit
when its antecedent is unavailable.

\begin{center}
\small
\begin{tabular}{@{}p{0.12\textwidth}p{0.50\textwidth}r@{/}l@{}}
\toprule
Node & Obligation & Weight & earned \\
\midrule
$D_0$ & spherical Tait bridge and least-counterexample normal form & 12 & 12 \\
$D_{1a}$ & generic mesh, connected-decomposition, cut, and state-supply infrastructure & 14 & 14 \\
$D_{1b}$ & target-forced compositional cut chain with computable constants & 10 & 0 \\
$D_{1c}$ & bounded-face geodesic and long-face perimeter transfers & 6 & 0 \\
$D_{1d}$ & planar wall reduction and fixed-wall exclusion & 8 & 0 \\
$D_{2a}$ & raw exact Count semantics, finite profiles, composition, and pumping & 10 & 10 \\
$D_{2b}$ & executable optimized closure certificate & 4 & 0 \\
$D_3$ & physical splice, class closure, strict decrease, and normalization & 14 & 14 \\
$D_{4a}$ & route-native finite-base reflection interface & 7 & 7 \\
$D_{4b}$ & concrete trace, coverage, adequacy, and empty finite-base audit & 9 & 0 \\
$D_5$ & headline assemblies with all premises discharged & 6 & 0 \\
\midrule
& Total & 100 & 57 \\
\bottomrule
\end{tabular}
\end{center}

Thus the honest baseline is
\[
                 \boxed{57\%}.
\]
The former closed-web/GWCO/Cell--3 development earns zero points.  Its former
structural allocation is represented by the still-open target-forced cut
chain $D_{1b}$ rather than being deleted from the denominator or credited to
a dead route.  In particular, the hundreds of checked local closed-web lemmas
do not raise this percentage.

For completeness, every conditional marker in this spine is accounted for as
follows.  The relative wall-exclusion reduction ($D_{1d}$) scores zero because
fixed-wall exclusion and the finite base are open.  The global-corridor
alternative ($D_{1c}$) scores zero because both geometric transfers are open.
The Cell--3 rooted factorization has ledger weight zero and score zero because
its global premise (U) is refuted.  The direct Count assembly and the
wall-exclusion assembly split $D_5$ equally and both score zero because their
conjunctive state-supply, wall, representation, decomposition, or finite-base
premises remain open.  The two explicit open markers are $D_{1b}$ and
$D_{2b}$, with weights ten and four respectively.  Future recounts must keep
these weights fixed; only the earned column may change when a cited premise is
sealed or refuted.

For scheduling only, the remaining obligations are ranked by fixed ledger
weight times a pre-audit tractability estimate.  These estimates never change
the ledger.  The first D1c falsification gate lowered its estimate because the
checked local carrier does not supply the complete cumulative interface on the
Goldberg test family.  The current order is therefore

\begin{center}
\small
\begin{tabular}{@{}p{0.12\textwidth}p{0.49\textwidth}rrr@{}}
\toprule
Node & Obligation & weight & estimate & product \\
\midrule
$D_{2b}$ & executable optimized closure certificate & 4 & $.88$ & $3.52$ \\
$D_{1b}$ & target-forced cut chain and constants & 10 & $.32$ & $3.20$ \\
$D_{1c}$ & geodesic and perimeter transfers, now after $D_{1b}$ & 6 & $.42$ & $2.52$ \\
$D_{4b}$ & concrete finite-base trace and empty audit & 9 & $.25$ & $2.25$ \\
$D_{1d}$ & fixed-wall exclusion & 8 & $.18$ & $1.44$ \\
$D_5$ & assemblies, gated by conjunctive premises & 6 & $0$ & $0$ \\
\bottomrule
\end{tabular}
\end{center}

The estimates are explicit search policy, not confidence annotations on
theorems.  In particular the D1c gate changes the attack order, not the fixed
six-point weight or its current earned score of zero.
They are the initial scheduling estimates, not a current probability model;
Section~\ref{sec:absolute-topology-gate} revisits the schedule against the
actual direct-headline dependencies without changing any ledger weight.

Consequently the route is a conditional proof, not yet a proof of the
Four-Color Theorem.  The point of the remainder of this section is that the
conditions are now explicit enough to attack independently and cannot be
blurred by prose.

The formal M0 boundary now has one exact input: the compositional core must
exclude graph-backed vertex-minimal non-Tait-colourable spherical maps.
From that proposition, least-counterexample selection, two-sidedness, graph
backing, Tait colourability of every bridgeless spherical cubic map, and the
spherical Four-Colour headline are all checked.
\Sverified{Mettapedia.GraphTheory.FourColor.GoertzelV24SphericalMinimalCounterexampleSelection.everyBridgelessSphericalCubicTaitColorable_of_no_minimal}
\Sverified{Mettapedia.GraphTheory.FourColor.GoertzelV24SphericalMinimalCounterexampleSelection.sphericalFourColorStatement_of_no_minimal}
\Sverified{Mettapedia.GraphTheory.FourColor.GoertzelV24RotationTaitCount.rotationSystemTaitCount_eq_zero_iff}
\Sverified{Mettapedia.GraphTheory.FourColor.GoertzelV24StellarTaitReduction.stellarDual_taitCount_eq_zero_of_not_colorable}
The refuted legacy \texttt{IsPlanar} predicate and the refuted closed-web
headline are not imported anywhere in this chain.

\subsection{Entering the normal form without the v23 degree-four slip}

The original v23 proof of minimum dual degree contains a small but real error:
from ``a deleted vertex has at most four neighbours'' it concludes that four
colours leave an unused colour.  This is false when four neighbours use all
four colours.  The current route does not consume that sentence.

Start instead with a vertex-minimal uncolourable connected spherical cubic
map, as supplied by the Tait bridge below.  A cyclic cut of size two, three, or
four splits the map into two smaller open cubic pieces.  Cap both pieces,
colour the smaller closed maps by minimality, and compare the colours on the
cut.  For two and three ports the boundary pattern is forced up to a colour
permutation.  At four ports the Klein-group boundary sum is zero, leaving the
all-equal class and the two non-crossing adjacent-pair classes; if one side is
all-equal, a planar Kempe escape moves it into an adjacent class.  Thus the
two sides always admit compatible boundary classes and glue to a Tait
colouring, a contradiction.  The least counterexample is therefore
cyclically five-edge-connected.

This complete cut-and-glue argument is machine-checked.
\Sverified{Mettapedia.GraphTheory.FourColor.GoertzelV24FourEdgeCutGluing.cyclicallyFiveEdgeConnected_of_vertexMinimalTaitCounterexample}

The finite logical heart of the corresponding five-boundary argument can be
stated without enumerating boundary colourings.

\begin{lemma}[abstract Birkhoff propagation on a five-cycle]
\label{lem:abstract-birkhoff-five-cycle-propagation}
\Sverified{Mettapedia.GraphTheory.FourColor.Compositional.FiveCycleBoundaryClosure.exists_drop}
\Sverified{Mettapedia.GraphTheory.FourColor.Compositional.FiveCycleBoundaryClosure.transfer_previous}
\Sverified{Mettapedia.GraphTheory.FourColor.Compositional.FiveCycleBoundaryClosure.not_pointwise_disjoint}
For each of two boundary languages, let $P_i,Q_i$ be two cyclic families of
behaviours, indexed modulo five.  Suppose each language realizes some $P_i$
and, for each language separately,
\[
 P_i\wedge\neg P_{i+1}\Longrightarrow Q_i,
 \qquad
 \neg P_{i+1}\wedge\neg Q_{i+1}\Longrightarrow P_i.
\]
Then the two languages overlap in some $P_i$ or in some $Q_i$.
\end{lemma}

\begin{proof}
Assume the languages are pointwise disjoint.  The first language realizes
some $P_i$ and misses some $P_j$, because the second realizes $P_j$.
Walking once around the five-cycle therefore finds an index $i$ at which
$P_i$ holds and $P_{i+1}$ fails.  The first implication gives $Q_i$.

Disjointness says that the second language realizes neither $P_i$ nor $Q_i$.
The second implication, followed by the first, forces it to realize both
$P_{i-1}$ and $Q_{i-1}$.  Repeating this transfer alternately between the two
languages four times forces the first language to realize $P_{i-4}=P_{i+1}$,
contradicting the chosen drop.
\end{proof}

\begin{lemma}[physical support semantics of the five-cycle classes]
\label{lem:five-cycle-class-support-semantics}
\Sverified{Mettapedia.GraphTheory.FourColor.Compositional.BoundaryKempeSwitch.switchBoundaryWord_mem_supportInCoordinates}
\Sverified{Mettapedia.GraphTheory.FourColor.Compositional.BoundaryColorOrbit.relabelBoundaryWord_mem_supportInCoordinates_iff}
\Sverified{Mettapedia.GraphTheory.FourColor.Compositional.BoundaryColorOrbit.exists_common_word_of_equivalent}
\Sverified{Mettapedia.GraphTheory.FourColor.Compositional.FiveCycleColorClass.equivalent_of_primary}
\Sverified{Mettapedia.GraphTheory.FourColor.Compositional.FiveCycleColorClass.equivalent_of_secondary}
\Sverified{Mettapedia.GraphTheory.FourColor.Compositional.FiveCycleSupportOverlap.exists_common_word_of_closure}
For a finite cubic open tangle, switching any union of complete physical
bichromatic components preserves its exact boundary support.  The same support
is invariant under one global permutation of the three nonzero Tait colours.
On a cyclic boundary of length five, let $P_i$ be the colour orbit of the
rotated $(3,1,1)$ pattern and let $Q_i$ be the colour orbit of the rotated
$(2,1,1,1)$ pattern.  Two boundary words in the same $P_i$ class, or in the
same $Q_i$ class, differ by such a global colour permutation.  Consequently,
if the two shore supports realize the same $P_i$ or $Q_i$ class, then after
globally relabelling one realizing colouring the supports contain a literal
common boundary word.  In particular, if both shore languages satisfy the two
closure laws in
Lemma~\ref{lem:abstract-birkhoff-five-cycle-propagation} and each realizes a
primary class, their exact supports contain a literal common boundary word.
\end{lemma}

\begin{proof}
A physical component-union switch is the restriction to the ports of the
usual Kempe switch on the selected components.  Reindexing the ports does not
change this statement.  Global colour relabelling transports a proper
colouring to a proper colouring and hence preserves exact support.  Finally,
an ordered Tait triple together with zero enumerates all four elements of the
Klein colour group; the induced zero-fixing equivalence carries either
five-boundary pattern to the corresponding pattern with the new colour names,
and commutes with cyclic rotation.
\end{proof}

\begin{lemma}[cyclic order on an exact five-edge bond]
\label{lem:exact-five-cut-cyclic-boundary}
\Sverified{Mettapedia.GraphTheory.FourColor.Compositional.CutPairDeletionConnectivity.deleteEdges_pair_connected_of_connected_shores}
\Sverified{Mettapedia.GraphTheory.FourColor.Compositional.FiveEdgeCutBoundaryPrerequisites.fixedPointFree_perm_card_five_transitive_or_twoCycle}
\Sverified{Mettapedia.GraphTheory.FourColor.Compositional.FiveEdgeCutBoundaryPrerequisites.deleteEdges_pair_connected_of_exactCyclicFiveCut}
\Sverified{Mettapedia.GraphTheory.FourColor.Compositional.FiveEdgeCutBoundaryOrder.deletedBoundarySuccessor_sameCycle_of_exactCyclicFiveCut}
\Sverified{Mettapedia.GraphTheory.FourColor.Compositional.FiveEdgeCutBoundaryOrder.retainedBoundarySuccessor_eq_deleted_inverse_of_exactCyclicFiveCut}
\Sverified{Mettapedia.GraphTheory.FourColor.Compositional.FiveEdgeCutBoundaryOrder.exists_boundaryOrder_of_exactCyclicFiveCut}
\Sverified{Mettapedia.GraphTheory.FourColor.Compositional.FiveEdgeCutBoundaryOrder.exists_retainedBoundaryOrder_of_exactCyclicFiveCut}
\Sverified{Mettapedia.GraphTheory.FourColor.Compositional.FiveEdgeCutBoundaryOrder.exists_cyclicBondBoundaryData_of_exactCyclicFiveCut}
Let an exact cyclic five-edge cut lie in a connected spherical graph-backed
map with cyclic vertex rotations and two-sided faces.  If both induced shores
are connected, then the five exposed darts occur in one cyclic facial-return
order on either shore, and the two orders have opposite orientation.  In
particular, either retained shore supplies the generic cyclic-bond boundary
data consumed by the physical Kempe-chain theory.
\end{lemma}

\begin{proof}
The deleted-side first-return permutation has no fixed point: a fixed exposed
dart would put both orientations of its edge on one face.  A fixed-point-free
permutation of five points is either one five-cycle or contains a two-cycle.
In the second case the two corresponding cut edges occur on the same pair of
ambient faces, so deleting them disconnects the primal graph.  On the other
hand, deleting any two crossing edges leaves both induced shores connected,
and any one of the three remaining cut edges still joins them.  This is a
contradiction.  Thus the first-return permutation is one five-cycle.  The
planar bond theorem makes the complementary first-return order its inverse,
and finite cyclic coordinates package the retained shore as a cyclic bond.
\end{proof}

\begin{lemma}[physical noncrossing on an exact five-edge shore]
\label{lem:exact-five-cut-physical-noncrossing}
\Sverified{Mettapedia.GraphTheory.FourColor.Compositional.CyclicBoundaryMatching.physicalMateInCoordinates_involutive}
\Sverified{Mettapedia.GraphTheory.FourColor.Compositional.CyclicBoundaryMatching.physicalMateInCoordinates_isNoncrossing_five}
\Sverified{Mettapedia.GraphTheory.FourColor.Compositional.FiveEdgeCutPhysicalNoncrossing.exists_physicalMateInCoordinates_noncrossing_of_exactCyclicFiveCut}
Fix one shore of an exact cyclic five-edge cut and a proper Tait colouring of
that shore.  For each unordered pair of nonzero colours, following the
physical bichromatic component from a boundary port defines an involutive
mate of the five ports.  In the facial cyclic order supplied by
Lemma~\ref{lem:exact-five-cut-cyclic-boundary}, no two mate chords cross.
\end{lemma}

\begin{proof}
The physical mate is involutive because a bichromatic component has exactly
the same two boundary endpoints when traversed from either end.  Suppose that
two mate chords crossed.  Their endpoints then alternate in the five-port
cyclic order.  One component together with either boundary interval between
its endpoints forms a separator.  The other component has one endpoint on
each interval, while its interior is disjoint from the separator.  The
two-sided-face and connected-dual hypotheses therefore forbid the second
component from joining those endpoints.  This contradiction proves
noncrossing.  For an exact five-cut, connectedness of both shores, connected
facial duality, and the cyclic coordinates are derived from the ambient
spherical map, so no extra boundary-order hypothesis remains.
\end{proof}

Lemma~\ref{lem:abstract-birkhoff-five-cycle-propagation} is a checked abstract
sufficient condition, but it is not the physical heart of the five-cut proof:
the support of a single $Y$-cap already refutes its proposed per-shore closure
laws.  The completed argument therefore does not derive or consume those laws.
Instead it retains the physical content proved above---global colour closure
and noncrossing bichromatic mates---and performs Birkhoff's short case argument
directly on the five boundary positions.

\begin{theorem}[compositional exact-five-cut theorem]
\label{thm:compositional-exact-five-cut}
\Sverified{Mettapedia.GraphTheory.FourColor.Compositional.FiveCutWordHeart.exists_common_word}
\Sverified{Mettapedia.GraphTheory.FourColor.Compositional.FiveCutShorePhysical.boundaryLanguage_retainedShore}
\Sverified{Mettapedia.GraphTheory.FourColor.FiveCutShoreAssembly.exists_common_word_on_shores}
\Sverified{Mettapedia.GraphTheory.FourColor.FiveCutSplice.false_of_nontrivial_exactCyclicFiveEdgeCut}
\Sverified{Mettapedia.GraphTheory.FourColor.FiveCutSplice.no_nontrivial_exactCyclicFiveEdgeCut_of_vertexMinimalTaitCounterexample}
\Sverified{Mettapedia.GraphTheory.FourColor.FiveCutSplice.exactCyclicFiveEdgeCut_has_small_shore_of_vertexMinimalTaitCounterexample}
A graph-backed vertex-minimal Tait counterexample contains no exact cyclic
five-edge cut whose two vertex shores both have at least six vertices.
Equivalently, every exact cyclic five-edge cut in such a counterexample has a
shore of cardinality at most five.
\end{theorem}

\begin{proof}
Use an exact-size-five cut carrier; this is distinct from the earlier
\texttt{SmallCyclicEdgeCut}, whose defining cardinality bound is four.  Pin the
facial cyclic order on one shore and transport the opposite orientation to the
other.  Capping either shore by the pentagon and by three consecutive $Y$-caps
produces smaller bridgeless spherical cubic maps.  Vertex minimality colours
those maps, so each shore language meets the four corresponding cap supports.
The physical Kempe theorem makes each shore language closed under every
noncrossing bichromatic mate, while global colour relabelling gives colour
closure.  The twelve-leaf Birkhoff word argument then forces the two shore
languages to contain one literal common boundary word.  Realizing colourings
with that word agree on all five cut edges; gluing them dart by dart gives a
Tait colouring of the ambient map, contradicting minimality.
\end{proof}

This is one generic bounded-interface normal-form theorem, not a covering
argument over planar configurations.  It strengthens the entry normal form but
does not settle the high-width decomposition theorem below.

Together with the already checked exclusion of cyclic cuts of size at most
four, this yields the cubic-side form of internal six-connectivity.  Every
realized cyclic cut with at most five edges has a shore with at most five
vertices; in the exact-five case the bounded shore has exactly five vertices
and is the support of a simple cycle of length five.  The name remains
explicitly dual-facing until a formal planar-duality bridge identifies this
statement with the standard internally-six-connected triangulation language.
\Sverified{Mettapedia.GraphTheory.FourColor.cubicDualInternallySixConnected_of_vertexMinimalTaitCounterexample}
\Sverified{Mettapedia.GraphTheory.FourColor.CyclicEdgeCutRealization.hasShoreOfSizeAtMost_five}
\Sverified{Mettapedia.GraphTheory.FourColor.exactCyclicFiveEdgeCut_has_pentagonal_shore_of_vertexMinimalTaitCounterexample}

A facial cycle of length below five would expose such a small cyclic cut (the
finitely small exceptional carriers are directly colourable), so every face
has length at least five; this consequence is checked separately.
\Sverified{Mettapedia.GraphTheory.FourColor.GoertzelV24MinimalFaceSize.orbitFaceMinimumFive_of_vertexMinimalTaitCounterexample}
Hence the curvature calculation below is entered through a sound normal form,
not through the erroneous four-neighbour extension.

\subsection{The abstract reductive pumping theorem}

We isolate the exact ordinary-mathematics theorem used by either a direct
zero-Count target or a Lift/Seed target once its physical cut chain is given.

\begin{theorem}[Finite-profile reductive pumping]
\label{thm:finite-profile-reductive-pumping}
\Sverified{Mettapedia.GraphTheory.FourColor.GoertzelV24FiniteProfilePumping.forall_P_of_base}
Let $\mathcal T$ be a class of finite open instances with a size
$|\!\cdot\!|\colon\mathcal T\to\mathbb N$ and a predicate $P$.  Assume:
\begin{enumerate}
\item every $\neg P$ instance having a ladder-reducible local feature has a
strictly smaller $\neg P$ reduction in $\mathcal T$;
\item there are constants $V_0,N$ such that every irreducible
$X\in\mathcal T$ with $|X|>V_0$ has a corridor with more than $N$
\emph{pairwise suitably separated} admissible cut positions;
\item every cut has a profile in a set $Q$ with $|Q|\le N$;
\item whenever two suitably separated cuts have the same profile, deleting
the material strictly between them constructs $X'\in\mathcal T$ with
$|X'|<|X|$ and
\[
                  \neg P(X)\Longrightarrow\neg P(X').
\]
\end{enumerate}
Then every $\neg P$ instance descends to one of size at most $V_0$.  If the
finite base $\{X\in\mathcal T:|X|\le V_0\}$ contains no $\neg P$ instance,
then $P$ holds throughout $\mathcal T$.
\end{theorem}

\begin{proof}
Suppose $\neg P$ occurs and choose such an $X$ of least size.  Hypothesis~1
makes $X$ irreducible.  If $|X|>V_0$, Hypothesis~2 supplies more than $N$
pairwise suitably separated cut positions.  Their profiles lie in $Q$, so
two are equal by the pigeonhole principle.  This pair satisfies the
separation premise of Hypothesis~4, which then constructs a strictly smaller
$\neg P$
instance, contrary to minimality.  Hence $|X|\le V_0$.  The final assertion
is immediate from emptiness of the base.  \qedhere
\end{proof}

This theorem is elementary; that is a virtue.  It shows exactly where the
difficult mathematics lives.  The finite Count functor supplies
Hypothesis~3 and helps formulate Hypothesis~4.  It does not supply
Hypothesis~2, and a quotient-code equality supplies Hypothesis~4 only after a
physical splice theorem proves that the decoded pieces really glue.

\subsection{What self-loops do and do not buy}

The source's L2 flag is a self-loop property on relevant profiles.  It gives
a useful monotonicity theorem in a genuinely periodic corridor, but it does
not erase the need to locate that periodic geometry.

\begin{lemma}[One-rung padding]
\label{lem:one-rung-padding}
\Sverified{Mettapedia.GraphTheory.FourColor.GoertzelV24PumpingEngine.CountData.count_succ_pos_of_count_pos}
Let $S$ be one fixed corridor rung with nonnegative transfer matrix $M$ on a
finite profile set.  Let $A,B$ be fixed end pieces with boundary vectors
$v,u$.  Suppose every profile occurring on an accepted path through one copy
of $S$ has a positive diagonal entry in $M$.  Then
\[
                    v^{\mathsf T}Mu>0
       \quad\Longrightarrow\quad
                    v^{\mathsf T}M^2u>0.
\]
\end{lemma}

\begin{proof}
Choose $i,j$ with $v_iM_{ij}u_j>0$.  The profile $i$ is alive on this accepted
path, so $M_{ii}>0$.  The term
$v_iM_{ii}M_{ij}u_j$ is a positive summand of $v^{\mathsf T}M^2u$.
\end{proof}

Thus a minimal zero-count obstruction cannot contain two consecutive copies
of the same physically removable rung, once the exact factorization and
re-insertion theorem are available.  For a corridor word
$M_L\cdots M_1$ with varying letters, however, a self-loop for one local
matrix does not show that deleting an arbitrary letter reconstructs the
original instance.  One must either find repeated physical rung data or use
the equal-profile splice of Theorem~\ref{thm:finite-profile-reductive-pumping}.
Accordingly profile cardinality has not been removed from the general source
argument.

\subsection{The curvature calculation and the missing global implication}

For a connected cubic plane map whose face lengths are at least five, write
$p_k$ for the number of $k$-faces and
\[
             W_-:=\sum_{k\ge7}(k-6)p_k.
\]
Euler's formula and double counting give
\begin{equation}
  \sum_f(6-|f|)=12,
  \qquad\text{hence}\qquad
  p_5=12+W_-.
  \label{eq:signed-curvature-balance}
\end{equation}
This equality is machine-checked in the rotation-system development.
\Sverified{Mettapedia.GraphTheory.FourColor.GoertzelV24CurvatureScope.faceCyclePentagonCount_eq_twelve_add_negativeCurvatureWeight}

Equation~\eqref{eq:signed-curvature-balance} is a \emph{signed} balance.  It
does not bound $p_5$, $W_-$, or the total number of nonhexagonal faces.  For
every $m$ the formal face multiset
\[
                p_5=m+12,\qquad p_7=m,\qquad p_k=0\ (k\ne5,7)
\]
satisfies the same identity and has $2m+12$ nonhexagonal faces and no
hexagons.  Lean records this precise arithmetic obstruction.
\Sverified{Mettapedia.GraphTheory.FourColor.GoertzelV24CurvatureScope.exists_arbitrarily_large_curvatureBalance_without_hexagons}

The same failure is realized by actual normal-form maps, rather than only by
formal face multisets.

\begin{proposition}[dodecahedral chains realize unbounded signed-defect cancellation]
\label{prop:dodecahedral-chain-no-hexagons}
\Sprose
For every $s\geq0$ there is a connected simple cellular spherical cubic map
$D_s$ with girth five and edge-connectivity three such that
\[
 \begin{aligned}
 |V(D_s)|&=18s+20,& |E(D_s)|&=27s+30,\\
 p_5(D_s)&=6s+12,& p_8(D_s)&=3s,
 \end{aligned}
\]
and $D_s$ has no faces of any other length.  Moreover, two disjoint untouched
pentagonal faces may be designated as boundary faces.  The resulting annular
map has $6s+10$ internal pentagons, $3s$ internal octagons, no internal
hexagon, and
\[
             \sum_{f\ \mathrm{internal}}(6-|f|)=10
\]
independently of $s$.
\end{proposition}

\begin{proof}
Let $D$ be the dodecahedral map and choose opposite vertices $a,b$.  Their
incident face sets are disjoint: vertices on one pentagonal face have graph
distance at most two, whereas $d_D(a,b)=5$.  Arrange $s+1$ disjoint copies
$D_0,\ldots,D_s$ in a row.  At the $i$th join delete $b$ from $D_i$ and $a$
from $D_{i+1}$, and match the three exposed neighbours in reverse cyclic
order.  Topologically this removes one open vertex-disc from each sphere and
glues the two boundary circles orientation-reversingly, so the result is
again a cellular map on the sphere.  Every exposed neighbour loses one edge
and gains one seam edge, hence the result remains cubic and simple.

Each join removes two vertices and six old edges and inserts three seam
edges.  It also pairs the six pentagonal face wedges at the deleted vertices
into three faces.  A pentagonal wedge left after deleting its vertex is a
three-edge path; two such paths and the two intervening seam edges form an
octagon.  Since the $a$- and $b$-incident faces in an internal block are
disjoint, no face is altered at two joins.  Starting from $s+1$ copies, this
gives
\[
 \begin{aligned}
 V&=20(s+1)-2s=18s+20,\\
 E&=30(s+1)-3s=27s+30,\\
 p_5&=12(s+1)-6s=6s+12,\qquad p_8=3s.
 \end{aligned}
\]
Thus $F=9s+12$, and both $3V=2E$ and $V-E+F=2$ hold.

For completeness, the local normal-form properties survive the join.  The
dodecahedral graph is $3$-edge-connected.  Restoring a deleted connector and
placing it on the side containing a majority of its three neighbours turns
any edge cut of size at most two in a three-edge sum into a cut of size at
most two in one factor.  That factor cut is trivial; iterating along the row
shows that the original cut was trivial.  Hence the sum is
$3$-edge-connected.  A cycle lying in one block has length at least five.  A
cycle using a seam uses at least two seam edges; between two distinct
neighbours of a deleted dodecahedral vertex every path has length at least
three, since a shorter path together with that vertex would give a triangle
or quadrilateral in $D$.  Such a cycle therefore has length at least eight.
This proves the girth assertion.

The dodecahedron has a Tait colouring.  At a deleted cubic vertex its three
incident colours are distinct.  Before each join, globally permute the three
colour names on the new block so that corresponding exposed half-edges have
the same colour, and give the new seam edge that common colour.  Induction
therefore also gives a Tait colouring of every $D_s$.

At the two free ends choose a pentagon incident with the unused connector
vertex.  These two faces are disjoint and were not altered by a join.
Declaring their interiors to be the two holes removes two pentagons from the
spherical census.  The internal curvature is consequently
\[
       (6s+10)(6-5)+(3s)(6-8)=10,
\]
while all $9s+10$ internal faces are nonhexagonal.  This completes the
construction.
\end{proof}

The census arithmetic is independently checked in Lean:
\Sverified{Mettapedia.GraphTheory.FourColor.GoertzelV24DodecahedralChainCensus.sphere_cubic_incidence}
\Sverified{Mettapedia.GraphTheory.FourColor.GoertzelV24DodecahedralChainCensus.sphere_euler}
\Sverified{Mettapedia.GraphTheory.FourColor.GoertzelV24DodecahedralChainCensus.sphereFaceSizes_curvature}
\Sverified{Mettapedia.GraphTheory.FourColor.GoertzelV24DodecahedralChainCensus.sphereFaceSizes_hexagonCount}
\Sverified{Mettapedia.GraphTheory.FourColor.GoertzelV24DodecahedralChainCensus.annularInternalFaceSizes_curvature}
\Sverified{Mettapedia.GraphTheory.FourColor.GoertzelV24DodecahedralChainCensus.exists_arbitrarily_large_annular_census_without_hexagons}

Proposition~\ref{prop:dodecahedral-chain-no-hexagons} refutes precisely the
geometry-only implication used in the source L1 and L6: fixed signed excess,
girth five, and absence of $1$- or $2$-edge cuts do not force even one
hexagonal face, much less a long hexagonal ladder.  The family is
Tait-colourable, so it does \emph{not} refute a target-aware M1 theorem for a
least zero-Count obstruction.  It shows why that theorem must use zero Count
and minimality essentially rather than repeating the signed-defect argument.

The annular and framed Excess Identities have the same sign issue, with
boundary surplus terms added.  They bound a difference of positive and
negative curvature resources, not their sum.  Therefore the sentence in
Addendum~XXV saying that the Excess Identity makes the nonhexagonal defect
budget independent of $V$ is not a consequence of that identity.

The tempting repair ``add a uniform negative-curvature bound'' is not a
second proof obligation.  It is only one sufficient route to the actual
target-forced cut-chain theorem stated below.  No such bound follows from the
signed identity, and the proof should not spend effort manufacturing a
stronger hypothesis than it consumes.  The existing weighted corridor
theorems correctly retain the negative-curvature weight or a caller-supplied
pentagon bound; they do not manufacture a uniform constant.  The live global
question is therefore stated once, as
Proposition~\ref{prop:target-forced-cut-chain}.  This is the highest
mathematical risk in the route.

The strongest existing all-face extraction theorem is exact but
graph-relative.  If \(W=W_-(M)\) and a positive corridor length \(L\) is
fixed, it proves
\[
 |F(M)|\le
 (W+7)^{(((12+2W)(W+7)+1)L)-1}
 \quad\text{or}\quad
 M\text{ contains a radius-one clean hexagonal geodesic block of length }L.
\]
\Sverified{Mettapedia.GraphTheory.FourColor.GoertzelV24WeightedOrbitFaceCorridor.orbitFace_card_le_weightedThreshold_or_exists_cleanHexCorridor}
The first alternative is not a finite base: its right-hand side contains the
unbounded parameter \(W(M)\) of the very graph being bounded.  For example,
allowing \(W\ge |F|\) makes the inequality compatible with arbitrarily large
\(|F|\).  Thus this theorem is a correct conditional corridor extractor, not
a proof of a graph-independent threshold \(V_0\).

The uniform weight bound is sufficient, but it is stronger than the
compositional architecture actually needs.  The following elementary
dichotomy removes curvature from the global existence question.

\begin{theorem}[bounded-face geodesic or a long face]
\label{thm:dual-diameter-dichotomy}
\Sverified{Mettapedia.GraphTheory.FourColor.GoertzelV24OrbitFaceDualDiameterDichotomy.exists_long_face_or_geodesic}
Let $M$ be a connected cellular spherical map with two-sided simple face
boundaries and connected facial dual.  Fix integers $B\ge2$ and $L\ge2$.  If
the number of faces is greater than
\[
       1+B\sum_{i=0}^{L-2}(B-1)^i,
\]
then either $M$ has a face of length greater than $B$, or its facial dual has
a geodesic of at least $L$ edges.  In the latter case the geodesic is an
induced face path, and every central face on it has length at most $B$.
\end{theorem}

\begin{proof}
Assume every face has length at most $B$.  Assign to each neighbouring face
a shared boundary edge.  Since a primal edge is incident with at most two
quotient faces, this assignment is injective, so the facial dual has maximum
degree at most $B$.  If its diameter were at most $L-1$, breadth-first layers about any
root would contain at most $1$ vertex at level zero, at most $B$ at level
one, and at most $B(B-1)^{i-1}$ at level $i\ge1$: after the first step a
shortest path cannot immediately return along the edge just used.  Summing
through level $L-1$ gives the displayed upper bound, contrary to the face
count.  Thus two dual vertices have distance at least $L$; a shortest path
between them is the required geodesic.  Any adjacency between two
nonconsecutive vertices of that path would shorten it, so the path is
induced.
\end{proof}

For a connected cubic spherical map, $3V=2E$ and Euler's equality give
$F=V/2+2$.  Thus the face-count hypotheses in the theorem and corollary are
effective vertex-count thresholds.  In particular, the dichotomy applies
to every sufficiently large member of the cubic normal form; it is not a
statement about maps separately assumed to have many faces.

\begin{corollary}[avoiding finitely many marked faces]
\label{cor:dual-geodesic-avoids}
\Sverified{Mettapedia.GraphTheory.FourColor.GoertzelV24OrbitFaceDualDiameterDichotomy.exists_marked_face_avoiding_geodesic}
Under the bounded-face branch of
Theorem~\ref{thm:dual-diameter-dichotomy}, let $S$ be a set of $h$ forbidden
faces and put
\[
 C=1+h(B-1),\qquad
 M=1+B\sum_{i=0}^{L-2}(B-1)^i.
\]
If the facial dual has more than $h+CM$ vertices, the induced facial dual on
the complement of $S$ contains a geodesic of at least $L$ edges.  We call
this an $S$-avoiding geodesic; no claim is made that it is geodesic in the
undeleted facial dual.
\end{corollary}

\begin{proof}
Deleting one vertex of degree at most $B$ from a connected graph increases
the number of components by at most $B-1$.  Deleting all $h$ forbidden
vertices therefore leaves at most $C$ components.  More than $CM$ vertices
remain, so one component has more than $M$ vertices.  Its maximum degree is
still at most $B$; the breadth-first estimate in the preceding proof gives a
geodesic of length at least $L$ in that component, hence an $S$-avoiding
geodesic in the stated sense.
\end{proof}

Both the underlying bounded-degree graph statement and its quotient-face map
adapter---including the sharp component count $1+h(B-1)$, the Moore threshold,
and geodesicity in the deleted graph---are machine-checked.
\Sverified{Mettapedia.GraphTheory.FourColor.GoertzelV24DeletionComponents.exists_avoiding_geodesic}
The verified theorem does not silently assert that the resulting path is
geodesic before deletion.

Thus, within the bounded-face branch, the cap faces and any fixed collection
of target or collar faces can be put into $S$ at an explicit multiplicative
cost.  Bounded-degree dual connectivity gives this finite avoidance directly;
it does not require an unsigned defect bound.

Consequently the bounded-face branch is a genuine chain of cells drawn from
a finite alphabet depending only on $B$.  A dual step records its chosen
primal edge incidence; each central face has at most $B$ such choices, and
the geodesic has no nonconsecutive face adjacency.  This conclusion does
\emph{not} require the stronger, presently unused assertion that two faces
share at most one edge.  The existing hexagonal Cell is the $B=6$,
all-central-faces-hexagonal subalphabet, not the only possible finite
alphabet.

There is a useful local bound here, but it must not be confused with a bound
on the complete transfer interface.

\begin{lemma}[$S$-avoiding geodesic locality]
\label{lem:geodesic-locality}
\Sverified{Mettapedia.GraphTheory.FourColor.GoertzelV24SeparatedGeodesicLocalCarrier.orbitFace_locality_and_carrier_le}
Let $S$ contain at most $h$ faces and let
$F_0,F_1,\ldots,F_L$ be a geodesic in the facial dual with $S$ deleted.
If a face $H\notin S$ is adjacent in the original dual to both $F_i$ and
$F_j$, then $|i-j|\le2$.

More precisely, call an unmarked face a crossing side face at the cut after
$i$ if it is adjacent to some $F_a$ with $a\le i$ and to some $F_b$ with
$i+1\le b\le L$.  Retain the complete boundaries of (i) every unmarked
neighbour of the four-face window
$F_{i-1},F_i,F_{i+1},F_{i+2}$ (omitting nonexistent end terms), (ii) the
window faces themselves, and (iii) all marked faces.  Every crossing side
face is retained.  If all faces have length at most $B$, this radius-one
carrier has cost at most
\[
                         4B^2+4B+hB
\]
darts.  In particular the bound is independent of $L$.
\end{lemma}

\begin{proof}
Because $H\notin S$, the two dual edges $F_iH$ and $HF_j$ survive the
deletion and form a walk of length two there.  Geodesicity in the deleted
dual gives $|i-j|=d_{D-S}(F_i,F_j)\le2$.  Therefore any unmarked side face
crossing the cut is adjacent to one of
$F_{i-1},F_i,F_{i+1},F_{i+2}$ (omitting nonexistent end terms).  Those four
axis faces have together at most $4B$ boundary incidences, hence at most
$4B$ distinct unmarked neighbours.  Storing their complete boundaries and
the axis boundaries costs
at most $4B^2+4B$.  A marked face need not obey the locality conclusion, but
there are at most $h$ of them and retaining each whole boundary costs at most
$B$ further darts.  Adding these contributions gives $4B^2+4B+hB$.
\end{proof}
\Sverified{Mettapedia.GraphTheory.FourColor.GoertzelV24SeparatedGeodesicLocality.index_dist_le_two}
\Sverified{Mettapedia.GraphTheory.FourColor.GoertzelV24SeparatedGeodesicLocality.carrier_le_of_counts}

Lemma~\ref{lem:geodesic-locality} bounds the faces directly incident with the
axis near one cut.  It does \emph{not} yet bound the complete interface to the
unprocessed graph.  Components outside the radius-one slab may join side
faces met at many different axis positions; a transfer state that remembers
those pending connections can therefore grow with the prefix.  Proving that
this cannot happen in a least Seed obstruction, or
turning a deep family of such connections into a second pumpable direction,
is a separate theorem.  Consequently the lemma does not by itself terminate
the receipt-enlargement process.

This distinction is visible even in defect-free geometry.  For the Goldberg
family obtained by subdividing every icosahedral triangle with frequency
\(k\) and dualizing, one has twelve pentagons and only hexagons, so \(W_-=0\)
while the number of vertices tends to infinity.  A direct combinatorial audit
of the generated maps gives the following cap-distance cut sizes.  Here
\(B_d\) is the set of primal vertices at graph distance at most \(d\) from
the five vertices of one pentagonal cap, and \(\delta B_d\) is its edge
boundary:
\[
\begin{array}{c|rrrrrrrrrrrr}
k&1&2&3&4&5&6&7&8&9&10&11&12\\ \hline
|V|&20&80&180&320&500&720&980&1280&1620&2000&2420&2880\\
\max |\delta B_d|&10&20&30&40&50&60&70&80&90&100&110&120.
\end{array}
\]
This is numerical evidence, not a theorem about every possible cut.  Its
logical force is narrower and important: an arbitrarily large clean
hexagonal region is not automatically a one-dimensional bounded-interface
ladder.  The deterministic generator and all per-radius arrays are in
\texttt{tools/fourcolor/v24\_goldberg\_frontier\_width\_audit.py} and its
JSON certificate.

\begin{proposition}[small-case gate for local-to-complete geodesic transfer]
\label{prop:geodesic-transfer-small-case-gate}
\Sprose
The bounded radius-one carrier around one cut does not, by itself, give a
bounded complete interface for cumulative prefixes of a dual geodesic.  On
the deterministic diameter geodesic of \(\mathrm{GP}(k,0)\), for
\(1\le k\le12\), every four-face radius-one cut carrier has primal edge
frontier at most $30$.  Nevertheless the maximum frontier of the
radius-one band around a geodesic prefix is

\[
  10,30,42,54,66,78,90,102,114,126,138,150,
\]

respectively.  For $k\ge2$ this is $12k+6$.  The positive control is the
source's \((5,0)\) tube through length $20$, whose corresponding radial
frontiers remain at most $10$.  The generator also reproduces the earlier
Goldberg calibrations $|V|=20k^2$ and maximum cap-radial frontier $10k$.

Thus the direct inference ``bounded local carrier implies a bounded
one-dimensional cumulative Count interface'' fails on the tested family.
This does not refute Proposition~\ref{prop:global-corridor-alternative}: a
target-forced cut chain or a second-direction transfer may absorb the lateral
attachments.  It does show that D1c cannot be discharged from
Lemma~\ref{lem:geodesic-locality} alone, so D1b now precedes D1c in the attack
order.  The deterministic audit and full arrays are
\texttt{tools/fourcolor/v24\_geodesic\_transfer\_gate.py} and
\texttt{results/fourcolor/v24\_geodesic\_transfer\_gate.json}.
\end{proposition}

\subsection{The cap-composed Menu predicate}

The terminal predicate must be fixed before any finite closure is measured.
The source does not count bichromatic components in the open tangle alone.
It restores the pentagonal cap and then counts the components obtained by
composing the open strand pairing with the cap pairing.  These are different
objects.

\begin{definition}[cap-composed pair relation]
\label{def:cap-composed-pair-relation}
\Sverified{Mettapedia.GraphTheory.FourColor.GoertzelV24AnnularFrontierMenuCapComposed.CapComposedMenuB}
Let $w$ be a good inner word, let $x$ be a proper colouring of the five cap
cycle edges extending $w$, and fix a majority pair $(p,q)$.  Let
$P_{p,q}\subseteq\mathbb Z/5\mathbb Z$ be the four inner positions whose
spokes have colour $p$ or $q$.  On $P_{p,q}$ put two relations:
\begin{align*}
 i\mathrel{\pi_{\rm str}}j
   &\quad\Longleftrightarrow\quad
   \text{the two spokes lie in the same $(p,q)$-component of the open tangle},\\
 i\mathrel{\pi_{\rm cap}}j
   &\quad\Longleftrightarrow\quad
   \text{$i$ and $j$ are joined in the cap $5$-cycle by edges of colours
   $p,q$}.
\end{align*}
The \emph{cap-composed pair relation} is the equivalence relation generated
by $\pi_{\rm str}\cup\pi_{\rm cap}$.  The state has menu B for $(p,q)$
exactly when this relation has at least two classes on $P_{p,q}$.
\end{definition}

\begin{lemma}[finite pentagonal cap table]
\label{lem:finite-cap-pairing-table}
\Sverified{Mettapedia.GraphTheory.FourColor.GoertzelV24AnnularFrontierMenuCapTable.normalizedCapTable_transport}
Normalize the majority colour to $0$, the singleton colours to $1,2$,
and let $x_i$ colour the cap edge from $i$ to $i+1$, with indices
modulo five.  For the ten good words, the unique cap extension and the two
cap matchings are:
\[
\begin{array}{c|c|c|c}
w&x&\pi_{\rm cap}(0,1)&\pi_{\rm cap}(0,2)\\ \hline
00012&21201&03\mid12&01\mid24\\
00021&12102&01\mid24&03\mid12\\
00120&12012&01\mid24&04\mid13\\
00210&21021&04\mid13&01\mid24\\
01200&20121&04\mid13&02\mid34\\
02100&10212&02\mid34&04\mid13\\
10002&21210&03\mid12&14\mid23\\
12000&01212&02\mid34&14\mid23\\
20001&12120&14\mid23&03\mid12\\
21000&02121&14\mid23&02\mid34
\end{array}
\]
Here $ab\mid cd$ denotes the perfect matching with blocks
$\{a,b\}$ and $\{c,d\}$.  Global colour permutation gives every
unnormalized case.
\end{lemma}

\begin{proof}
Properness at cap vertex $i$ says
$\{w_i,x_{i-1},x_i\}=\{0,1,2\}$.  Thus any one cap-edge colour forces the
next, and cyclic consistency leaves the single extension displayed in each
row.  For a majority pair, each of its four active cap vertices has exactly
one cap edge in the pair.  Tracing those selected edges gives the two
displayed blocks.  The ten rows exhaust the cyclic placements of a
consecutive triple together with the order of the two singleton colours.
\end{proof}

\begin{lemma}[literal cap restoration computes the composed relation]
\label{lem:cap-restoration-composition}
\Sverified{Mettapedia.GraphTheory.FourColor.GoertzelV24AnnularFrontierMenuCapRestoration.restoredReachable_iff_composedRel}
The classes in Definition~\ref{def:cap-composed-pair-relation} are exactly
the $(p,q)$-components of the graph obtained by restoring the five cap edges,
restricted to the four active cap vertices.
\end{lemma}

\begin{proof}
Every $(p,q)$-walk in the restored graph alternates between maximal portions
inside the open tangle and maximal portions inside the restored cap.  Its
successive active cap vertices are therefore related alternately by
$\pi_{\rm str}$ and $\pi_{\rm cap}$.  Hence its endpoints lie in the
equivalence closure of their union.  Conversely, every generator of
$\pi_{\rm str}$ is witnessed by a walk in the open tangle and every generator
of $\pi_{\rm cap}$ by a walk in the cap.  Concatenating the witnesses for a
chain in the generated equivalence relation gives a $(p,q)$-walk in the
restored graph.
\end{proof}

\begin{lemma}[transverse-chord criterion]
\label{lem:menu-b-transverse-chord}
\Sverified{Mettapedia.GraphTheory.FourColor.GoertzelV24AnnularFrontierMenuCapComposed.capComposedMenuB_iff_no_transverseChord}
For either majority pair, let the two blocks of
$\pi_{\rm cap}$ be $A\mid B$.  Then the cap-restored state has menu B
for that pair if and only if no pair joined by $\pi_{\rm str}$ has one
endpoint in $A$ and the other in $B$.
\end{lemma}

\begin{proof}
The cap relation is a perfect matching on the four active positions, while
the open-strand relation is a partial matching: a bichromatic component of
the open tangle has at most two inner degree-one ends.  If every strand pair
stays within $A$ or within $B$, each generator of
$\langle\pi_{\rm cap}\cup\pi_{\rm str}\rangle$ preserves the partition
$A\mid B$, so the generated relation has at least two classes.  Conversely,
one transverse strand joins a vertex of $A$ to a vertex of $B$; the two
cap pairs then connect all four active positions into one class.  By
Lemma~\ref{lem:cap-restoration-composition}, these are exactly the two
menu cases.
\end{proof}

\subsubsection{The actual Lift target is weaker than one seed per fibre}

Write a normalized good word uniquely as
\[
 w(s;m,x,y)_{s,s+1,s+2,s+3,s+4}=(m,m,m,x,y),
 \qquad s\in\mathbb Z/5\mathbb Z,
\]
where $(m,x,y)$ is an ordering of the three colours.  The coordinate $s$
is its \emph{block}.  Thus the thirty good words form five blocks of six.
The singleton-pair switch and a majority-pair menu-A switch change only the
ordering $(m,x,y)$ and not $s$.  A menu-B switch for $(m,x)$ joins block
$s$ to $s+1$, while a menu-B switch for $(m,y)$ joins block $s$ to $s-1$.
Equivalently, if $v$ is the unique inactive singleton position for the
chosen majority pair, the installed edge of the five-cycle of blocks is
\[
                         e=v+1\pmod 5,
\]
where edge $e$ joins blocks $e$ and $e+1$.

\begin{theorem}[four-block-edge Lift criterion]
\label{thm:four-block-edge-lift}
\Sverified{Mettapedia.GraphTheory.FourColor.GoertzelV24FourBlockEdgeLift.oneClass_of_fourBlockEdges_word}
Let $\mathcal C$ be the good-word colouring states of one fixed annular
tangle.  Suppose:
\begin{enumerate}
\item every nonempty fixed-word fibre is connected by inner-avoiding Kempe
moves;
\item in each of the five blocks, the nonempty realized word vertices form a
connected subgraph under within-block moves; and
\item cap-composed menu-B moves realize at least four of the five adjacent
edges of the block cycle.
\end{enumerate}
Then all states in $\mathcal C$ lie in one Kempe class.  In particular the
pentagon Lift step holds for this tangle.
\end{theorem}

\begin{proof}
Four edges of a five-cycle form a spanning path, so Hypotheses~2 and~3 make
the realized word graph connected.  Let two colouring states be given and
choose a path between their words.  At each word edge, Hypothesis~1 first
moves within the current fibre to the state witnessing that edge; perform
the corresponding cap-composed move, and continue in the next fibre.
Hypothesis~1 finally reaches the prescribed target state.  Thus the full
state graph is connected.  Restoring the uniquely coloured pentagonal cap
identifies these states and moves with the colourings and Kempe moves in the
frontier before the cap was opened, which is precisely the Lift conclusion.
\end{proof}

The later ``one seed in every nonempty good fibre'' formulation is therefore
a sufficient strengthening, not the target consumed by the word-graph
argument.  A source-faithful finite audit should retain the majority-pair
label of every seed and ask which of the five block edges it installs.  The
reproducible pair-resolved audit is
\texttt{tools/fourcolor/v24\_plantri\_lift\_bridge\_audit.py}.  This weaker
criterion does not erase the normal-form issue below: the triangular witness
realizes no block edge and the square witness realizes only edge~$3$.

On the normal-form generator class, the pair-resolved audit has so far checked
all $568$ connected annuli at $V=24,26,28,30$.  Every one realizes at
least four block-cycle edges; no failure occurs.  The exact per-size counts,
solver version, and scope warning are recorded in
\texttt{results/fourcolor/v24\_plantri\_lift\_bridge\_audit\_v24\_30.json}.
This is evidence for Hypothesis~3 of
Theorem~\ref{thm:four-block-edge-lift}, not a proof of it or of the two
within-fibre connectivity hypotheses.

\begin{proposition}[the open-tangle implementation predicate is vacuous]
\label{prop:open-menu-predicate-vacuous}
\Sprose
For well-formed annular boundary data, a proper Tait colouring whose inner
word has shape $(3,1,1)$ has at least two \emph{open-tangle} bichromatic
components meeting the inner boundary, for each of its two majority pairs.
Consequently the present Lean predicates
\texttt{AnnularFrontierMenuBForPair} and
\texttt{AnnularFrontierMenuBForPairOfProfile}, which use only the open
component relation (respectively its \texttt{connectionTable} encoding), are
true on every realizable colouring/profile in every nonempty good fibre and
are not the source's menu-B predicate.
\end{proposition}

\begin{proof}
Fix one majority pair.  Exactly four inner spokes carry one of its two
colours.  Their four boundary vertices are distinct and have degree one in
the corresponding bichromatic support graph.  Every vertex of that support
graph has degree at most two: an interior cubic Tait vertex has exactly the
two selected colours, while a boundary vertex has ambient degree one.
A finite connected graph of maximum degree two has at most two degree-one
vertices.  Thus one component cannot contain all four active inner stubs;
at least two components meet the inner boundary.  The same argument applies
to the other majority pair.
\end{proof}

This is a semantic correction, not a refutation of the Seed Lemma.  It does
mean that the current Lean seed count and its profile-factorization theorems
count the whole good fibre rather than the source's cap-composed seeds, so
they cannot discharge the Seed Lemma as presently named.  The rolling engine
still computes useful data: it already retains the open strand relation
$\pi_{\rm str}$, while $\pi_{\rm cap}$ is a five-port finite function of the
fixed boundary word and its cap extension.  The repair is therefore to add
the finite equivalence closure of
Definition~\ref{def:cap-composed-pair-relation} to terminal acceptance.
Lemma~\ref{lem:cap-restoration-composition} now verifies that this finite
closure agrees pointwise, on active ports, with reachability in the literally
restored graph; the repair is not a reason to discard the one-cell
factorization.

\begin{proposition}[cap-correct terminal acceptance is profile-local]
\label{prop:cap-correct-terminal-profile}
\Sverified{Mettapedia.GraphTheory.FourColor.GoertzelV24AnnularFrontierCapComposedProfileSemantics.restoredMenuBState_iff_capComposedOfUniformProfile}
\Sverified{Mettapedia.GraphTheory.FourColor.GoertzelV24AnnularFrontierCapComposedProfileSemantics.restoredMenuBState_iff_of_uniformProfile_eq}
For an actual Tait colouring of the open annular tangle, let the finite
uniform profile record its inner boundary word and, for each selected colour
pair, the Boolean relation saying which active ports lie in the same open
component.  The source's Menu-B test after literally restoring the pentagonal
cap is a predicate of this profile alone.  More precisely, it is the
cap-composed predicate of
Definition~\ref{def:cap-composed-pair-relation}; consequently two colourings
with equal uniform profiles either both pass the restored-cap terminal test or
both fail it.
\end{proposition}

\begin{proof}
For a majority pair the active support has four of the five cap positions, so
its complement is the unique inactive position.  A true entry of the profile's
component table joins two active positions.  Apply
Lemma~\ref{lem:cap-restoration-composition} to identify restored reachability
with the equivalence closure of the open component relation and the cap
pairing.  This description uses only the two stored profile fields, and hence
is invariant under profile equality.
\end{proof}

\begin{corollary}[cap-correct Seed count factors through finite profiles]
\label{cor:cap-correct-seed-profile-count}
\Sverified{Mettapedia.GraphTheory.FourColor.GoertzelV24AnnularFrontierCapComposedProfileCount.capComposedSeedCount_pos_iff_exists_profile}
\Sverified{Mettapedia.GraphTheory.FourColor.GoertzelV24AnnularFrontierCapComposedProfileCount.capComposedSeedCount_eq_zero_iff_forall_profile_count_eq_zero}
Define the cap-correct Seed count at a fixed word to be the number of proper
Tait colourings in that fibre which pass the restored-cap predicate of
Proposition~\ref{prop:cap-correct-terminal-profile}.  This count is positive
if and only if some realized finite profile at the word has positive profile
count and satisfies the cap-composed terminal predicate.  Equivalently, it is
zero if and only if every cap-composed accepting profile at that word has
literal profile count zero.
\end{corollary}

\begin{proof}
Partition the finite Tait-colouring fibre by its uniform profile.  Terminal
acceptance is constant on each part by
Proposition~\ref{prop:cap-correct-terminal-profile}; taking cardinalities gives
the positive form, and negating it gives the zero form.
\end{proof}

\begin{remark}[cap-correct Seed audit]
\label{rem:seed-census-all-pentagon-pairs}
\Sprose
The Seed conclusion has now been tested with the inner cap restored before
Kempe components are counted, as required by the Reformulation Theorem.
For every pair of vertex-disjoint pentagonal faces in the three source
laboratories, exhaustive enumeration of the literal five-by-five boundary
data gives
\[
\begin{array}{c|r|r|r|r|r|r}
 &\text{pairs}&\text{disconnected}&\text{good fibres}&\text{good states}&
   \text{menu-B states}&\text{seedless fibres}\\ \hline
C_{24}&42&24&1260&10080&9036&0\\
C_{30}&46&15&1380&23400&17340&0\\
C_{40}&46&10&1380&118320&89580&0\\ \hline
\text{total}&134&49&4020&151800&115956&0.
\end{array}
\]
Here ``disconnected'' means that opening the two cap cycles isolates one
edge joining an inner stub to an outer stub.  These rows are retained as a
boundary-data regression test because connectivity is absent from
\texttt{AnnularBoundaryData.WellFormed}; they are not controlling source
instances, since the full \texttt{ClosedWebAnnularEmbedding} contains the
hypothesis that the graph is connected.  On the $85$ connected source
instances the audit has $2550$ good fibres, $114840$ good states, $84066$
menu-B states, and no seedless fibre.
The polar cases reproduce the manuscript's published menu-B counts
\(228/240\), \(390/600\), and \(1680/3120\), which is a fidelity check on
the operational definition.  The deterministic audit is
\texttt{tools/fourcolor/v24\_seed\_lemma\_census.py}.
\end{remark}

This is genuine evidence for the Seed Lemma, but it is not its proof: these
are three finite fullerene families and all have five outer stubs.  It shows
that a valid global corridor theorem may profitably use the seedless
hypothesis; the manuscript's present derivation still does not construct the
required bounded-interface chain from it.

\paragraph{Wide Goldberg audit.}
The same cap-restored predicate has also been encoded directly in SAT on the
icosahedral Goldberg family, with a second literal component traversal used
to replay every satisfying witness.  On $\mathrm{GP}(1,0)$ ($20$ vertices),
all $72$ ordered vertex-disjoint cap annuli and all ten good-word
representatives modulo global colour permutation are nonempty and seeded:
$720/720$.  On $\mathrm{GP}(2,0)$ ($80$ vertices), the corresponding result is
$1320/1320$ over $132$ ordered annuli.  A maximally separated cap pair of
$\mathrm{GP}(3,0)$ ($180$ vertices, cap-vertex distance $15$) is seeded for
all ten representatives.  The exact constraints, graph hashes, and results
are generated by
\texttt{tools/fourcolor/v24\_goldberg\_seed\_sat.py}.

\begin{proposition}[the unrestricted Seed statement is false]
\label{prop:unrestricted-seed-refuted}
\Srefuted{connected triangle- and square-boundary witnesses}
Nonemptiness of a good fibre does not imply a cap-composed menu-B seed for
an arbitrary connected annular cubic tangle.
\end{proposition}

\begin{proof}
There is a connected $14$-vertex plane cubic map with a pentagonal face
disjoint from a triangular face such that opening those two faces gives
nonempty fibres at the normalized good words \(00012\) and \(00021\), while
both fibres contain zero cap-composed seeds.  Its rotation-system hash is
\[
  \texttt{96668268249e8043817aa0c8c1fff8adf23f9445fa9849a85c37de6654caa449}.
\]
The phenomenon is not removed merely by forbidding triangles: a connected
$18$-vertex example of minimum face length four, opened at a pentagon and
a square, has six nonempty normalized good fibres and is seedless at
\(10002\) and \(20001\).  Its hash is
\[
  \texttt{6fb5266b12bbf4e9c395e6ab6ea1a8ebf6dde039b48de5e0f97a67443865e039}.
\]
In both cases SAT supplies proper colourings and an independent literal
traversal restores the cap and recomputes the two majority-pair component
relations.  The complete embeddings, face cycles, generator hash, and
fibre lists are recorded in
\texttt{results/fourcolor/v24\_plantri\_seed\_scope\_triangle.json} and
\texttt{results/fourcolor/v24\_plantri\_seed\_scope\_square.json}.
\end{proof}

Thus the hypothesis silently supplied by Step~1 of Addendum~XXV is
load-bearing.  The route needs the Seed conclusion only for a frontier
obtained from the cyclically five-edge-connected, face-length-at-least-five
normal form of a least counterexample.  That is the theorem which
Proposition~\ref{prop:target-forced-cut-chain} now targets.  Alternatively,
one could retain the manuscript's unrestricted statement only after proving
that triangle/digon/two-cut/square reduction transports the
\emph{cap-composed} zero count.  The current rolling square identities use
the open-tangle predicate of
Proposition~\ref{prop:open-menu-predicate-vacuous}, so they do not provide
that proof.

\paragraph{Exhaustive small normal-form audit.}
An independent generator audit uses
\texttt{plantri -a -d -m5 -c5} to enumerate every generated simple cubic
spherical map of face length at least five and cyclic edge-connectivity at
least five through $34$ primal vertices.  Across $45$ maps it opens every
ordered pair consisting of a pentagonal inner face and an arbitrary
vertex-disjoint outer face of length $5$ through $8$.  The totals are
\[
\begin{array}{c|r|r|r|r}
V&\text{maps}&\text{boundary openings}&\text{connected annuli}
 &\text{seeded/nonempty good fibres}\\ \hline
20&1&72&12&720/720\\
22&0&0&0&0/0\\
24&1&96&36&960/960\\
26&1&108&48&1080/1080\\
28&3&380&190&3800/3800\\
30&4&539&294&5390/5390\\
32&12&1872&1092&18720/18720\\
34&23&3978&2448&39780/39780\\ \hline
\text{total}&45&7045&4120&70450/70450.
\end{array}
\]
On the connected source-annulus subset the corresponding total is
$41200/41200$.  Every SAT seed witness is replayed by an independent
literal cap-restored component traversal.  The generator hash, solver
version, aggregate counts, and explicit warning that this finite absence is
not a universal proof are recorded in
\texttt{results/fourcolor/v24\_plantri\_seed\_audit\_v2\_20\_32.json}
and \texttt{results/fourcolor/v24\_plantri\_seed\_audit\_v2\_34.json};
the reproducible driver is
\texttt{tools/fourcolor/v24\_plantri\_seed\_audit.py}.  This substantially
strengthens falsification coverage while remaining a test, not a replacement
for Proposition~\ref{prop:target-forced-cut-chain}.

At $V=36$, a connected-only run over all $71$ generated maps adds
$8721$ connected annuli and $87210/87210$ seeded/nonempty normalized
good fibres, with no seedless witness.  Thus the cumulative connected audit
through $36$ vertices is $128410/128410$ fibres over $12841$ annuli.
This longer run is recorded separately in
\texttt{results/fourcolor/v24\_plantri\_seed\_audit\_v36\_connected.json}
because disconnected openings were connectivity-filtered before SAT rather
than included in the displayed all-opening table.

These data are deliberately not promoted to a theorem.  They do, however,
separate the two phenomena which L1 had conflated: the same Goldberg family
has growing radial frontier width, while its tested good fibres have seeds.
Thus a successful global argument should use the assumption that the target
count is zero; a geometry-only claim that every large hexagonal patch is a
fixed-width tube is both stronger than needed and contradicted by the radial
audit.

\begin{proposition}[target-forced compositional cut chain]
\label{prop:target-forced-cut-chain}
\Sopen
For each fixed boundary and acceptance type of the chosen direct-Count or
Lift system, there are computable constants \(b,m,V_0\) with the
following property.  Every ladder-irreducible target counterexample of size
greater than \(V_0\) either has a strictly smaller target counterexample, or
contains more than \(m\) pairwise nested simple ordered transversals, each
meeting at most \(b\) graph edges, such that:
\begin{enumerate}
\item consecutive transversals bound composable slabs containing at least one
graph vertex;
\item the designated holes and frozen boundary lie outside the slabs;
\item every transversal has an exact saturated deletion side and a cycle on
each complementary side; and
\item all transversals carry one common finite ordered geometric seam type,
and cumulative target Count factors exactly through the associated boundary
state.
\end{enumerate}
Here \(m\) is at least the cardinality of the finite cumulative profile
carrier at width \(b\).
\end{proposition}

This proposition is the precise L1--splice junction.  It deliberately states
the geometric outputs which the existing Lean interfaces consume: a simple
transversal, its crossed-edge set, an exact deletion component boundary, two
cyclic sides, and the correspondence used for surgery.  A chord wall, an
induced dual path, or a bounded local collar supplies only part of this data.

\begin{lemma}[literal source-Cell prefixes are laminar]
\label{lem:source-cell-prefix-laminarity}
\Sverified{Mettapedia.GraphTheory.FourColor.GoertzelV24FramedCorridorSerialPrefixLaminarity.sourceCorridorSerialPrefixVertexSide_mono}
\Sverified{Mettapedia.GraphTheory.FourColor.GoertzelV24FramedCorridorSerialPrefixLaminarity.sourceCorridorSerialPrefixVertexSide_ssubset_succ_of_fresh}
\Sverified{Mettapedia.GraphTheory.FourColor.GoertzelV24FramedCorridorSerialPrefixLaminarity.sourceCorridorSerialPrefixIncidentShore_mono}
\Sverified{Mettapedia.GraphTheory.FourColor.GoertzelV24FramedCorridorSerialPrefixLaminarity.sourceCorridorSerialPrefixIncidentShore_ssubset_succ_of_fresh_edge}
Let \(C_0,C_1,\ldots\) be the literal source Cells of a realized corridor
and let
\[
                     U_r=\bigcup_{i<r}V(C_i).
\]
Then \(r\le s\) implies \(U_r\subseteq U_s\), and the canonical incident-edge
shores satisfy
\[
       \{e:e\cap U_r\ne\varnothing\}
       \subseteq
       \{e:e\cap U_s\ne\varnothing\}.
\]
If \(C_r\) contains a vertex not in \(U_r\), then
\(U_r\subsetneq U_{r+1}\).  If that Cell also contains an endpoint of an
edge not assigned to the old incident-edge shore, then the two consecutive
edge shores are strictly ordered.  Thus laminarity is automatic, while strict
shore growth has been reduced to this concrete fresh-edge condition.
\end{lemma}

\begin{lemma}[connected cyclic cuts are literal exact-state interfaces]
\label{lem:connected-cyclic-cut-exact-state-interface}
\Sverified{Mettapedia.GraphTheory.FourColor.CyclicEdgeCutRealization.card_middle_incidentEdgeShore_le_edgeCut}
\Sverified{Mettapedia.GraphTheory.FourColor.CyclicEdgeCutRealization.exists_connectedSides_subcut}
\Sverified{Mettapedia.GraphTheory.FourColor.CyclicEdgeCutRealization.exists_connectedAtWidth_subcut}
\Sverified{Mettapedia.GraphTheory.FourColor.CyclicEdgeCutRealization.toConnectedShoreNodeOfConnectedSides}
\Sverified{Mettapedia.GraphTheory.FourColor.CyclicEdgeCutRealization.toConnectedShoreNodeOfVertexMinimal}
\Sverified{Mettapedia.GraphTheory.FourColor.CyclicEdgeCutRealization.exists_connectedShoreNode_of_hasCyclicEdgeCutOfSizeAtMost}
\Sverified{Mettapedia.GraphTheory.FourColor.CyclicEdgeCutRealization.exists_endpointSimple_replacement_of_nested_connected_cyclicCuts}
Let an exact cyclic edge cut in a connected finite cubic graph have cut size
at most (b).  If either induced vertex side is disconnected, retain the
reachable component containing its certified cycle and recompute the exact
boundary.  The new boundary is a subcut of the old one and is strictly
smaller whenever the chosen side was disconnected.  Finite strict induction
therefore produces a cyclic subcut of size at most (b) whose two induced
vertex sides are connected.  Assign to the chosen side every edge with at
least one endpoint there.  This is a connected shore
whose complementary edge shore is connected, whose two majority regions are
nonempty, and whose middle set has cardinality at most (b).  Consequently
the cut canonically carries the literal phased exact-support state used by
the physical splice.  In a least counterexample, two strictly nested such
cuts with equal phased states produce a strictly smaller endpoint-simple
counterexample.
\end{lemma}

\begin{proof}
The saturation step terminates because each disconnected-side repair strictly
shrinks the finite cut support; repairing the complementary side is the same
construction applied to the complementary realization.  Cut containment
preserves the original width bound.  Connectivity of the two resulting
induced vertex sides transports to connectivity of
the incident-edge shore and its complement.  A cycle on either side supplies
a vertex in the corresponding majority region.  Every middle vertex lies on
the outer endpoint of a chosen dart leaving the selected vertex side;
distinct middle vertices give distinct darts.  Sending each such dart to its
underlying edge injects the outgoing darts into the exact cut, so the middle
cardinality is at most the cut cardinality.  These are precisely the fields
of the connected-shore node.  Completing that node constructs its roots,
ports, boundary order and exact support internally; the existing physical
equal-state replacement theorem then gives the last assertion.
\end{proof}

\begin{proof}
The defining index set for \(U_r\) is contained in that for \(U_s\) whenever
\(r\le s\), so every vertex contributed before the first cut is also
contributed before the second.  Assigning an edge to a shore when either
endpoint lies in the chosen vertex side is monotone under inclusion of vertex
sides.  Finally a vertex of \(C_r\setminus U_r\) belongs to \(U_{r+1}\) but not
to \(U_r\), proving strictness of the vertex sides.  An edge with such an
endpoint belongs to the new incident-edge shore; if it was absent from the old
shore, monotonicity plus this witness makes the shore inclusion strict.
\end{proof}

Consequently the source-native M1 residue can be stated without a separate
laminarity hypothesis: one must show that sufficiently many literal Cells
contribute fresh shore edges while the true graph frontier of their cumulative
prefix remains uniformly bounded and both complementary sides stay connected.
The displayed moving two-edge cut is not yet known to equal that true
frontier.

The following checked lemma supplies one of those outputs for a literal
finite segment of the source corridor.  It is useful precisely because its
hypotheses expose, rather than suppress, the remaining global geometry.

\begin{lemma}[a boundary-separated source rail segment supplies connected shores]
\label{lem:source-rail-segment-exact-shore}
\Sverified{Mettapedia.GraphTheory.FourColor.GoertzelV24FramedCorridorRailChainComponent.exists_alignedBoundary_component_exactBoundary_connectedEdgeShores}
Let three consecutive transverse source interfaces be joined on their two
sides by the literal facial-dual rails of the intervening corridor cells.
Assume that each composed rail is a simple path, that the two rail supports
are disjoint, that every face on either rail is an internal annular face,
and that neither rail visits the centre face of either endpoint interface.
Close the two rails by the transverse two-edge paths at the
endpoint interfaces, and let $W$ be the resulting closed facial-dual walk.
If $X(W)$ is the set of primal edges crossed by $W$, then $W$ is a
simple dual cycle and some connected component of $G-X(W)$ has graph
boundary exactly $X(W)$.  Both this component side and its complementary
vertex side induce connected subgraphs.  Moreover, if every edge incident
with the component side is assigned to one edge shore, then that edge shore
and its complementary edge shore are connected in the edge-walk sense used
by the majority-shore descent.
\end{lemma}

\begin{proof}
Split $W$ into the path formed by the first endpoint transversal followed
by one rail, and the path formed by the reverse last transversal followed by
the reverse other rail.  Simplicity of the rails, avoidance of the two centre
faces, and mutual rail disjointness imply that the two path tails are
disjoint.  Appending them therefore gives a simple cycle.  This first step is
checked as
\Sverified{Mettapedia.GraphTheory.FourColor.GoertzelV24FramedCorridorRailChainComponent.alignedBoundaryWalk_isCycle}.

For completeness, the saturation argument is not a Jordan-curve picture.
The framed dual-cycle separator first proves that deleting $X(W)$
disconnects the primal graph and supplies a component with nonempty computed
boundary.  Every computed boundary edge lies in $X(W)$.  Conversely,
two-sided face parity around the simple dual cycle says that a nonempty cut
contained in $X(W)$ must contain every crossed edge.  Hence the inclusion is
an equality.  The generic form, independent of the corridor realization, is
checked as
\Sverified{Mettapedia.GraphTheory.FourColor.GoertzelV24FramedDualCycleExactBoundary.exists_component_exactBoundary_of_dualCycle}.

The same saturation argument applies to every component of $G-X(W)$.
Because $X(W)$ is nonempty, a component distinct from the chosen one also
has nonempty boundary; parity therefore forces its boundary to be all of
$X(W)$.  Two distinct such components cannot coexist, since one crossed
edge has only two endpoints.  Thus the complement of the chosen component
is connected.  This generic strengthening is checked as
\Sverified{Mettapedia.GraphTheory.FourColor.GoertzelV24FramedDualCycleExactBoundary.exists_component_exactBoundary_connectedSides_of_dualCycle}.

Finally put in the chosen edge shore every ambient edge with at least one
endpoint in the chosen component.  A path in the component connects the
component endpoints of any two such edges, and the two edges themselves
attach their possible outside endpoints.  Hence this edge shore is
connected.  Its complement consists exactly of the edges having both
endpoints in the complementary vertex side, so connectivity of that induced
side proves connectivity of the complementary edge shore.  This construction
is canonical and monotone under inclusion of vertex sides.  Its generic form
is checked as
\Sverified{Mettapedia.GraphTheory.FourColor.GoertzelV24ConnectedVertexSideEdgeShore.connected_edgeShores_of_connected_vertexSides}.
\end{proof}

\begin{corollary}[the source rail segment has an oriented majority-shore certificate]
\label{cor:source-rail-segment-majority-shore}
\Sverified{Mettapedia.GraphTheory.FourColor.GoertzelV24FramedCorridorRailChainComponent.alignedBoundary_innerIncidentShore_certificate}
Assume in addition that the source trail is well formed.  Let
\(C_{\mathrm{out}}\) be the component of \(G-X(W)\) containing the
named outer-container root, let \(U_{\mathrm{in}}\) be its vertex
complement, and put in \(S_{\mathrm{in}}\) every edge incident with
\(U_{\mathrm{in}}\).  Then
\[
 \partial C_{\mathrm{out}}=X(W),\qquad
 S_{\mathrm{in}}\text{ and }E(G)\setminus S_{\mathrm{in}}
 \text{ are connected edge shores},
\]
and the majority vertex side of \(S_{\mathrm{in}}\) is nonempty and
proper.
\end{corollary}

\begin{proof}
The completed boundary visits only internal faces: this is an explicit
hypothesis on the two rails, while each endpoint transversal is internal by
the boundary-clean source realization.  Consequently \(X(W)\) is disjoint
from the outer-hole boundary.  The named outer-container cycle therefore
survives the deletion and lies in the component rooted at its own basepoint.
This fixes the annular orientation without choosing a component after the
fact.  These two steps are checked as
\Sverified{Mettapedia.GraphTheory.FourColor.GoertzelV24FramedCorridorRailChainComponent.alignedBoundaryCutEdges_disjoint_outerHoleBoundary}
and
\Sverified{Mettapedia.GraphTheory.FourColor.GoertzelV24FramedCorridorRailChainComponent.outerContainer_hasCycleOn_alignedBoundaryOuterComponent}.

The prescribed-component form of the dual-cycle separator now gives the
exact boundary and connectedness of both oriented vertex sides:
\Sverified{Mettapedia.GraphTheory.FourColor.GoertzelV24FramedDualCycleExactBoundary.component_exactBoundary_connectedSides_of_dualCycle}
and
\Sverified{Mettapedia.GraphTheory.FourColor.GoertzelV24FramedCorridorRailChainComponent.alignedBoundaryOuterComponent_exactBoundary_connectedSides}.
The incident-edge construction from the preceding lemma then gives the two
connected edge shores.

It remains to check that ``majority'' is not vacuous.  A simple dual cycle has
length at least three and, by unique shared-edge incidence, crosses that many
distinct primal edges.  Choose two.  If their endpoints on
\(U_{\mathrm{in}}\) coincide, that vertex already has two incident edges;
otherwise a path in the connected inner side starts with an internal edge
which, together with one crossing edge, gives the same conclusion.  Hence
some inner vertex has at least two incident edges, all in
\(S_{\mathrm{in}}\), and is a majority vertex.

For the other direction, every well-formed source-trail vertex has degree at
most three: ordinary internal vertices have degree three, the two defects
degree two, and frozen stubs degree one.  At a vertex of the surviving outer
cycle, two distinct incident cycle edges lie outside
\(S_{\mathrm{in}}\); at most one incident edge can lie inside it, so that
vertex is not a majority vertex.  The local degree and subcubic steps are
checked as
\Sverified{Mettapedia.GraphTheory.FourColor.GoertzelV24ConnectedVertexSideEdgeShore.exists_two_le_degree_of_connected_side_of_two_crossing_edges},
\Sverified{Mettapedia.GraphTheory.FourColor.GoertzelV24SourceTrailAlignment.SourceTrail.incidentEdgeFinset_card_le_three},
and
\Sverified{Mettapedia.GraphTheory.FourColor.GoertzelV24ConnectedVertexSideEdgeShore.exists_not_majorityVertexSide_incidentEdgeShore_of_complement_cycle_of_subcubic}.
\end{proof}

\begin{corollary}[the oriented source shore survives closure of the missing edge]
\label{cor:source-rail-segment-ambient-majority-shore}
\Sverified{Mettapedia.GraphTheory.FourColor.GoertzelV24FramedCorridorRailChainAmbient.alignedBoundary_ambientInnerIncidentShore_certificate}
Let the source graph in Corollary~\ref{cor:source-rail-segment-majority-shore}
be $H=G-e$, where $e=uv$ is an edge of a finite cubic ambient graph
$G$.  Use the same outer-root component $C_{\mathrm{out}}$, the same
vertex partition
\[
 U_{\mathrm{in}}=V(G)\setminus C_{\mathrm{out}},
 \qquad U_{\mathrm{out}}=C_{\mathrm{out}},
\]
but recompute the edge shore in the ambient graph by
\[
 S^G_{\mathrm{in}}
   =\{f\in E(G): f\text{ has an endpoint in }U_{\mathrm{in}}\}.
\]
Then $G[U_{\mathrm{in}}]$ and $G[U_{\mathrm{out}}]$ are connected,
$S^G_{\mathrm{in}}$ and its edge complement are connected, and the
majority vertex side of $S^G_{\mathrm{in}}$ is nonempty and proper.
The exact source equality
$\partial_H C_{\mathrm{out}}=X(W)$ remains available as the provenance of
the partition.
\end{corollary}

\begin{proof}
Restoring $e$ only adds an adjacency, so connectedness of both induced
vertex sides in $H$ implies their connectedness in $G$.  The generic
incident-edge-shore construction therefore gives connectedness of
$S^G_{\mathrm{in}}$ and its complement.  The inner witness from the source
certificate remains an inner vertex, and cubicity of $G$ makes it a
majority vertex for $S^G_{\mathrm{in}}$.  Likewise the surviving
outer-container cycle in $H$ is still a cycle in $G$; two cycle edges at
one of its vertices witness failure of majority on the outer side.

It is important that the edge shore is recomputed in $G$.  The restored
edge $e$ is allowed to cross the vertex partition, so the ambient cut-edge
set need not equal the source set $X(W)$.  No noncrossing premise is being
silently added.
\end{proof}

These results are not yet Proposition~\ref{prop:target-forced-cut-chain}.
They now construct, for one three-interface segment, a saturated boundary with a
canonical outer-root orientation, complementary connected edge shores, and a
nonempty proper majority core, and transport that oriented partition to the
closed ambient cubic map.  The internal-support premise is load-bearing:
rail disjointness alone would not prevent the completed boundary from running
through a cap.

It would be incorrect to obtain the required chain by fixing the left
interface and composing these exterior rails to successively later right
interfaces.  The resulting closed dual walk crosses the two endpoint
transversals and both longitudinal rails, so its number of crossed primal
edges grows linearly with the distance between the endpoints.  Such walks do
not have the uniform width (b) required in
Proposition~\ref{prop:target-forced-cut-chain}.

The already constructed sliding local windows do not repair this defect by
themselves.  Their exact failure mode is checked rather than inferred from a
picture.

\begin{lemma}[the local-window seriality dichotomy]
\label{lem:source-local-window-seriality}
For two consecutive aligned two-tile source windows, the output transverse
edge at each of the two positions is the input transverse edge of the next
window.  At each such edge, either the two enclosed-side boundary darts are
$\alpha$-opposite, as required for a literal serial seam, or the enclosed
vertex carriers of the two windows share the base vertex of that dart.  In
the $\alpha$-opposite branch, the common edge is internal after the two
enclosed vertex sides are united, hence it is absent from the true crossing
frontier of their union.  Moreover, the two closed facial-dual window walks
have intersecting support.
\end{lemma}

\begin{proof}
Equality of the two transverse edge coordinates and the fact that every
edge fibre consists of a dart and its $\alpha$-mate give the dichotomy.  If
the displayed darts are equal, their common base vertex belongs to both
enclosed carriers.  These statements are checked as
\Sverified{Mettapedia.GraphTheory.FourColor.GoertzelV24FramedCorridorAlignedTwoTileSuccessorBoundaryDart.sourceTwoTileAlignedEnclosed_alphaSeam_or_vertexOverlap}.
For opposite darts the two dart bases are the two endpoints of the common
edge, one in each enclosed side, so adjoining the sides absorbs that edge.
The resulting exact alternative---absorption or a concrete shared
vertex---is checked as
\Sverified{Mettapedia.GraphTheory.FourColor.GoertzelV24FramedCorridorAlignedTwoTileSuccessorAbsorption.sourceTwoTileAlignedEnclosed_sharedEdge_absorbed_or_vertexOverlap}.
The closed windows also share the centre of their common transverse
interface, so their supports are not disjoint; this is checked as
\Sverified{Mettapedia.GraphTheory.FourColor.GoertzelV24FramedCorridorAlignedTwoTileBoundaryOverlap.sourceTwoTileAlignedBoundaryWalks_consecutive_not_disjoint}.
\end{proof}

\begin{lemma}[two-window serial frontier bound]
\label{lem:source-two-window-serial-frontier}
\Sverified{Mettapedia.GraphTheory.FourColor.GoertzelV24VertexSetBoundaryUnionCancellation.vertexSetCrossingEdges_union_subset_sdiff_union_sdiff}
\Sverified{Mettapedia.GraphTheory.FourColor.GoertzelV24FramedCorridorAlignedTwoTileSerialBoundary.sourceTwoTileAlignedEnclosed_unionBoundary_subset_without_shared_of_alpha}
\Sverified{Mettapedia.GraphTheory.FourColor.GoertzelV24FramedCorridorAlignedTwoTileSerialBoundary.sourceTwoTileAlignedEnclosed_unionBoundary_card_le_eight_of_alpha}
Let \(A_i,A_{i+1}\) be the enclosed vertex sides of two consecutive aligned
two-tile source windows, and let \(T_i\) be their common two-edge transverse
interface.  If, at both positions of \(T_i\), the displayed boundary darts
are \(\alpha\)-opposite, then
\[
  \partial(A_i\cup A_{i+1})
  \subseteq
  \bigl(\partial A_i\setminus T_i\bigr)
  \cup
  \bigl(\partial A_{i+1}\setminus T_i\bigr),
  \qquad
  \bigl|\partial(A_i\cup A_{i+1})\bigr|\le 8.
\]
Here \(\partial A\) is the true primal edge frontier of the vertex side
\(A\), not merely the edge set displayed by a chosen dual walk.
\end{lemma}

\begin{proof}
For arbitrary vertex sides \(A,B\), an edge crossing \(A\cup B\) crosses at
least one of \(A\) and \(B\): choose its endpoint in the union and retain the
constituent side containing that endpoint.  If every endpoint of an edge in
a set \(T\) belongs to \(A\cup B\), no edge of \(T\) crosses the union.
Consequently
\[
  \partial(A\cup B)
  \subseteq
  (\partial A\setminus T)\cup(\partial B\setminus T).
\]

In the source specialization, the two outgoing transverse coordinates of
the first window are literally the two incoming coordinates of the second,
and they are distinct; hence \(|T_i|=2\).  Each belongs to both certified
six-edge frontiers.  The \(\alpha\)-opposite premise puts the two endpoints
of each shared edge in the united enclosed side, so the preceding generic
argument applies.  Finally
\[
  |\partial(A_i\cup A_{i+1})|
  \le |\partial A_i\setminus T_i|
     +|\partial A_{i+1}\setminus T_i|
  =(6-2)+(6-2)=8.
\]
Only the displayed inclusion is used: no converse assertion that every
unshared old frontier edge remains on the joined frontier is required.
\end{proof}

\begin{lemma}[separated local rail supports grow linearly]
\label{lem:source-local-rail-support-growth}
Let \(S\) be a finite set of aligned two-tile window starts such that any two
distinct starts differ by at least four positions.  For each \(i\in S\), let
\(R_i\) be the two primal edges displayed by the two rail coordinates of the
local six-port profile.  Then the sets \(R_i\) are pairwise disjoint and
\[
                         \left|\bigcup_{i\in S}R_i\right|=2|S|.
\]
This is a statement about the displayed local rail carriers, not yet a claim
that every such edge survives in the true crossing frontier after the
enclosed vertex sides are united.
\end{lemma}

\begin{proof}
At starts separated by at least four positions, all four possible pairs of
the two short facial-dual rail tracks have disjoint vertex support.  In a
two-sided rotation system, dual walks with disjoint support cross disjoint
primal edge sets.  Each local rail carrier has exactly two edges, so finite
additivity gives the formula.  The pairwise-disjointness and cardinality
statement are checked as
\Sverified{Mettapedia.GraphTheory.FourColor.GoertzelV24FramedCorridorSeparatedRailFrontierGrowth.card_biUnion_sourceTwoTileRailCrossingEdgesAt_eq_two_mul_card}.
\end{proof}

Thus coordinate agreement of the rolling local profiles is not yet physical
serial composition, and choosing one saturated component independently at
each window does not supply a nested shore chain.  The transverse portal is
absorbed in the oriented branch, but
Lemma~\ref{lem:source-local-rail-support-growth} shows that the displayed
lateral carriers themselves do not telescope: a separate theorem would have
to absorb all but uniformly many of them from the true frontier of the united
side.  The remaining theorem is therefore target-forced, not a geometry-only
source adapter: it must either prove precisely that additional absorption or
construct, at each selected corridor position, a different individually
bounded saturated ambient cut.
It must prove that these cuts are simple and pairwise ordered, that their
outer-root shores are strictly nested with nonempty intervening slabs, and
that both holes and the frozen interface remain outside those slabs.  It must
also decide the $\alpha$-seam branch of
Lemma~\ref{lem:source-local-window-seriality}, transport the cuts to the
ambient graph, prove connected complementary edge shores, nonempty majority
sides, and a uniform middle-set bound, and identify the recomputed ambient
shore boundaries with ordered seams.  Lemma~\ref{lem:connected-shore-literal-node}
then supplies the roots and ports, while the chain consumer computes the
ordered Count states used by
Theorem~\ref{thm:literal-majority-shore-chain-descent}.

There is a second, equally compositional way to state the same structural
risk.  It replaces a linear corridor by a bounded-width decomposition tree.
This is not asserted by the working manuscript, but it is a useful repair
normal form because its generic half is elementary and its target-specific
half has one unmistakable statement.

\begin{theorem}[finite-interface pumping on a decomposition tree]
\label{thm:finite-tree-interface-pumping}
\Sverified{Mettapedia.GraphTheory.FourColor.GoertzelV24FiniteTreeInterfacePumping.vertexCount_le_of_minimal}
Let \(X\) be an instance with a reduced rooted binary interface decomposition
\(T\).  Suppose:
\begin{enumerate}
\item each node introduces at most \(c\) vertices, every vertex is introduced
at exactly one node, and every proper ancestor--descendant segment contains
some introduced material;
\item every node carries one of finitely many ordered interface types and an
exact context state, with \(q\) typed states in total; and
\item whenever a proper ancestor \(u\) and descendant \(v\) have the same
typed state, replacing the piece below \(u\) by the piece below \(v\) is a
valid composition, strictly decreases size, and preserves the target truth
value in the exterior context.
\end{enumerate}
If \(X\) is minimal among target counterexamples, then
\[
                         |V(X)|\le c(2^q-1).
\]
\end{theorem}

\begin{proof}
No root-to-leaf path in \(T\) contains more than \(q\) nodes.  Otherwise two
nodes on that path have the same typed state.  They are comparable, and
Hypotheses~1 and~3 replace the higher piece by the lower one to produce a
strictly smaller target counterexample, contrary to minimality.

A rooted binary tree whose longest root-to-leaf path contains at most \(q\)
nodes has at most
\(1+2+\cdots+2^{q-1}=2^q-1\) nodes.  Since every graph vertex is introduced
once and a node introduces at most \(c\) vertices, the displayed bound
follows.
\end{proof}

The tree theorem deliberately leaves its target-specific input visible.  For
the direct Tait target, the exact semantic part of that input is already a
theorem.

\begin{theorem}[closed Count depends only on the inner seam support]
\label{thm:closed-count-support-replacement}
\Sverified{Mettapedia.GraphTheory.FourColor.GoertzelV24ClosedCountReplacement.not_closedColorable_of_innerSupport_eq}
Let \(O\) be an outside port tangle and let \(I,I'\) be inside port tangles
with the same ordered seam.  If \(I\) and \(I'\) realize exactly the same
seam-colour words, then
\[
  O\circ I\text{ is not Tait-colourable}
  \quad\Longrightarrow\quad
  O\circ I'\text{ is not Tait-colourable}.
\]
\end{theorem}

\begin{proof}
A colouring of \(O\circ I\) is uniquely a seam word together with a
colouring of \(O\) and a colouring of \(I\) realizing that word.  Thus the
closed composite is colourable exactly when the outer and inner supports
intersect.  Replacing the inner support by an equal set does not change this
intersection.  The Lean theorem proves the displayed implication; the same
argument gives the converse as well.
\end{proof}

\begin{theorem}[the semantic closed composite is the literal sewn map]
\label{thm:physical-closed-count-bridge}
\Sverified{Mettapedia.GraphTheory.FourColor.GoertzelV24PhysicalClosedCountBridge.closedColorable_iff_composeRotationSystem_taitColorable}
Let \(L\) and \(R\) be finite literal open rotation tangles and let
\(m:B_L\simeq B_R\) match their complete boundary-dart interfaces.  Regard
\(L\) as a port tangle with empty input and output interface \(B_L\), and
regard \(R\) as a port tangle with input interface \(B_L\), relabelled through
\(m\), and empty output.  Then the closed serial composite used by
\texttt{ClosedColorable} is colourable if and only if the literal rotation
system \texttt{composeRotationSystem} is Tait-colourable.
\end{theorem}

\begin{proof}
Write \(D_L,D_R\) for the intact dart carriers.  After deleting the two empty
external interfaces, the serial Count construction has interior dart carrier
\[
              (D_L\sqcup D_R)\sqcup(B_L\sqcup B_L),
\]
whereas the literal splice has carrier
\[
              (D_L\sqcup D_R)\sqcup(B_L\sqcup B_R).
\]
The required equivalence is the identity on intact darts and on the left seam
half, and sends the right seam half \(b\) to \(m(b)\).  It preserves the base
vertex of every dart.  It also conjugates the serial edge flip to the literal
matched-seam flip: old edges retain their old flips, while the serial pair
\((b_L,b_R)\) is sent to the literal pair \((b,m(b))\).

Consequently a proper nonzero serial dart-colouring pulls across this
equivalence to a proper nonzero colouring invariant under the literal flip.
Such a colouring descends to the two-dart edge orbits of the literal splice,
giving a Tait edge-colouring.  Conversely, a Tait edge-colouring of the
literal splice colours each serial dart by the colour of the literal edge
orbit containing its image.  Flip-invariance is automatic.  If two distinct
serial darts based at one vertex received the same colour, their image edges
would be distinct adjacent edges of the literal splice with the same colour,
contradicting properness.  This proves the biconditional.  Notice that the
vertex rotation is not used: this is exactly the representation-invariance
needed to join Count semantics to the verified structural splice.
\end{proof}

\begin{corollary}[support replacement preserves the physical zero Count]
\label{cor:physical-closed-count-replacement}
\Sverified{Mettapedia.GraphTheory.FourColor.GoertzelV24PhysicalClosedCountBridge.not_composeRotationSystem_taitColorable_of_taitInnerSupport_eq}
Fix a literal outside tangle \(L\).  Let \(R,R'\) be two literal inside
tangles whose boundaries are both matched to the same ordered seam of \(L\).
If \(R\) and \(R'\) realize exactly the same genuine nonzero Tait cut words,
then
\[
 L\mathbin{\#}R\text{ not Tait-colourable}
 \quad\Longrightarrow\quad
 L\mathbin{\#}R'\text{ not Tait-colourable},
\]
where both occurrences of \(\#\) are the literal matched-seam rotation-system
splice.
\end{corollary}

\begin{proof}
Use Theorem~\ref{thm:physical-closed-count-bridge} to move the first literal
splice to \texttt{ClosedColorable}.  Every boundary word realized by a proper
port colouring is pointwise nonzero, so equality of the finite sets of genuine
Tait cut words is already equality of the raw seam supports.  Apply
Theorem~\ref{thm:closed-count-support-replacement}, and use the same
biconditional in the reverse direction for the second literal splice.
\end{proof}

The following is the precise finite-tree statement currently proved.  It is
a theorem about a supplied bounded-width decomposition, not yet a theorem
constructing that supply from a plane map.

\begin{theorem}[typed bounded-width descent]
\label{thm:typed-sphere-cut-descent}
\Sverified{Mettapedia.GraphTheory.FourColor.GoertzelV24SphereCutDescent.vertexCount_le_of_widthSupply'}
Fix a literal seam bound \(k\).  Suppose an instance \(X\) is supplied with a rooted binary
decomposition whose nodes have states
\[
  (j,s,S),\qquad 0\le j\le k,\quad
  s\in\mathrm{SeamType}_j,\quad
  S\subseteq\{r,b,p\}^{j},
\]
and suppose:
\begin{enumerate}
\item \(|V(X)|\) is at most twice the number of labelled tree nodes; and
\item equal states at a proper ancestor and descendant produce a strictly
smaller target instance.
\end{enumerate}
If \(X\) is size-minimal among target instances, then
\[
 |V(X)|\le 2(2^q-1),\qquad
 q=\sum_{j=0}^{k}|\mathrm{SeamType}_j|\,2^{3^j}.
\]
\end{theorem}

\begin{lemma}[cubic spacing forces material between two shores]
\label{lem:cubic-sphere-cut-material}
\Sverified{Mettapedia.GraphTheory.FourColor.GoertzelV24SphereCutMaterial.exists_strict_slab_vertex_of_cubic}
Let \(A\subseteq B\subseteq E(G)\) be two edge shores of a finite cubic
graph.  For an edge set \(C\), let

\[
 \partial C=\{v:\text{\(v\) is incident with an edge of \(C\) and an edge
 outside \(C\)}\}.
\]

If \(|\partial A|,|\partial B|\le w\) and

\[
                         |B\setminus A|>6w,
\]

then there is a vertex \(v\) incident with an edge of \(B\setminus A\)
such that every edge incident with \(v\) lies in \(B\setminus A\).
Thus \(v\) lies strictly in the material between the two shores, rather
than on either middle set.
\end{lemma}

\begin{proof}
For every edge of \(B\setminus A\), choose one endpoint.  Suppose every
chosen endpoint belonged to \(\partial A\cup\partial B\).  That union has
at most \(2w\) vertices, and cubicity allows at most three chosen edges at
each vertex.  Hence \(|B\setminus A|\le6w\), a contradiction.  We may
therefore choose an edge \(e\in B\setminus A\) whose chosen endpoint \(v\)
belongs to neither middle set.

Because \(e\notin A\), if any edge of \(A\) met \(v\), then \(v\in\partial
A\).  Thus no edge at \(v\) lies in \(A\).  Because \(e\in B\), if any edge
outside \(B\) met \(v\), then \(v\in\partial B\).  Thus every edge at \(v\)
lies in \(B\).  Together these statements put every incident edge in
\(B\setminus A\), as required.
\end{proof}

\begin{theorem}[majority shores from a connected edge decomposition]
\label{thm:connected-edge-shore-majority}
\Sprose
Let \(G\) be a finite cubic graph and let \(A\subseteq E(G)\).  Write
\[
 d_A(v)=|\{e\in A:v\in e\}|,\qquad
 K_A=\{v:d_A(v)\ge2\},\qquad
 M_A=\{v:0<d_A(v)<3\}.
\]
Suppose that the edge-induced graphs on \(A\) and on
\(\bar A=E(G)\setminus A\), with isolated vertices discarded, are connected.
Then:
\begin{enumerate}
\item \(K_{\bar A}=V(G)\setminus K_A\);
\item if \(K_A\ne\varnothing\), then \(G[K_A]\) is connected, and likewise
      for \(K_{\bar A}\);
\item the edge boundary of \(K_A\) has at most \(|M_A|\) edges; and
\item if \(A\subseteq B\), then \(K_A\subseteq K_B\).  Moreover, if every
      edge incident with \(v\) lies in \(B\setminus A\), then
      \(v\in K_B\setminus K_A\).
\end{enumerate}
Consequently, a connected branch decomposition of width \(w\), restricted
to non-leaf cuts for which both shores contain at least two edges,
canonically supplies nested connected complementary vertex shores with at
most \(w\) crossing edges.  The shore is a function of the edge set alone;
no independently chosen noose or perturbation is involved.
\end{theorem}

\begin{proof}
At every vertex cubicity partitions its three incident edges between \(A\)
and \(\bar A\), so
\[
                         d_A(v)+d_{\bar A}(v)=3.
\]
Thus \(d_{\bar A}(v)\ge2\) exactly when \(d_A(v)\le1\), proving the first
claim.  For the second, take \(x,y\in K_A\).  Connectedness of the
edge-induced \(A\)-graph gives an \(A\)-edge walk from \(x\) to \(y\).
Erase repeated vertices to obtain a path.  Its endpoints already lie in
\(K_A\), and every internal vertex is incident with the two distinct path
edges adjacent to it.  Both edges lie in \(A\), so the whole path lies in
\(K_A\).  This proves that \(G[K_A]\) is connected.  Apply the same argument
to \(\bar A\).

For the boundary bound, orient every edge of
\(\delta(K_A)\) from \(K_A\) to its complement.  If the edge lies in \(A\),
send it to its outer endpoint.  That endpoint has one incident \(A\)-edge
but not two, hence belongs to \(M_A\); and two different such edges cannot
have the same endpoint, since that would give two incident \(A\)-edges
there.  If the edge lies in \(\bar A\), send it to its inner endpoint.
The inner endpoint has at least two incident \(A\)-edges and therefore, by
cubicity, exactly one incident \(\bar A\)-edge; so this assignment is again
injective.  The images of the two cases are disjoint because the first
case lands outside \(K_A\) and the second inside it.  We have constructed an
injection \(\delta(K_A)\hookrightarrow M_A\).

Finally, \(A\subseteq B\) implies \(d_A(v)\le d_B(v)\), proving monotonicity.
If all three edges at \(v\) lie in \(B\setminus A\), then
\(d_B(v)=3\) and \(d_A(v)=0\), which proves the strict statement.

For a connected branch decomposition, \(M_A\) is its usual middle set and
therefore has size at most \(w\).  An internal shore containing at least two
edges has a vertex incident with two of them: otherwise its edge-induced
graph would be disconnected.  Hence both majority shores are nonempty on
the stated cuts, and the preceding claims apply.  Descendant edge shores
are nested, so their majority shores are nested as well.
\end{proof}

The graph-local content of this theorem is machine-checked:
\Sverified{Mettapedia.GraphTheory.FourColor.GoertzelV24ConnectedEdgeShoreMajority.majorityVertexSide_complement_iff}
\Sverified{Mettapedia.GraphTheory.FourColor.GoertzelV24ConnectedEdgeShoreMajority.majorityVertexSide_mono}
\Sverified{Mettapedia.GraphTheory.FourColor.GoertzelV24ConnectedEdgeShoreMajority.strict_majority_material_of_incident_sdiff}
\Sverified{Mettapedia.GraphTheory.FourColor.GoertzelV24ConnectedEdgeShoreMajority.card_crossingSideDart_majority_le_middle}
\Sverified{Mettapedia.GraphTheory.FourColor.GoertzelV24ConnectedEdgeShoreMajority.connected_induce_majorityVertexSide}
\Sverified{Mettapedia.GraphTheory.FourColor.GoertzelV24ConnectedEdgeShoreMajority.connected_majority_shores_and_width}
The remaining conditional in the final sentence is exactly the existence of
the connected branch decomposition itself and its finite-tree packaging; it
is not part of any of the six declarations just cited.

The conversion also feeds the physical splice without another geometric
choice.  Put $D_A=V(G)\setminus K_A$, open the ambient rotation system on
$K_A$ and $D_A$, and match the two boundary copies by the ambient edge
flip.  Suppose the ambient map is connected, bridgeless, spherical, cubic,
has cyclic vertex rotations and two-sided faces; suppose both edge shores are
connected and both majority shores are nonempty; and suppose the cut has two
distinct boundary darts.  Then the canonical first-return hub rotations on
the two shores have opposite orientation, both capped shores satisfy the
spherical Euler equality, and their literal composite is a connected,
bridgeless, spherical cubic map with cyclic vertex rotations.

Indeed, Theorem~\ref{thm:connected-edge-shore-majority} supplies the two
connected vertex shores.  The planar-bond rank theorem closes each shore by
the inverse first-return hub rotation, the complementary-shore theorem proves
that the two hub rotations are oppositely oriented under the ambient edge
flip, and the seam face-count theorem proves sphericity of the composite.
All other fields are transported directly through the literal splice.  This
entire implication is machine-checked as
\Sverified{Mettapedia.GraphTheory.FourColor.GoertzelV24ConnectedEdgeShoreStructuralData.bridgelessSphericalCubicMapData_of_majorityEdgeShore}

The roots and the two distinct boundary darts in this construction are not
additional decomposition data.  They follow from connectedness,
bridge-freeness, and cubicity.

\begin{lemma}[connected shores supply a literal descent node]
\label{lem:connected-shore-literal-node}
\Sverified{Mettapedia.GraphTheory.FourColor.GoertzelV24ConnectedShoreLiteralNode.LiteralShoreNode.ofConnectedShore}
Let \(M\) be a graph-backed vertex-minimal non-Tait-colourable connected,
bridgeless, spherical cubic rotation map.  Let \(A\subseteq E(M)\) have
connected edge shore and connected edge complement, with both majority
vertex sides nonempty.  If the middle set of \(A\) has cardinality at most
\(k\) and at most \(w\), then \(A\) supplies a literal shore
node of width at most \(k\) and middle bound \(w\), including retained-dart
roots on both sides and two distinct exterior boundary darts.
\end{lemma}

\begin{proof}
Write \(K\) for either nonempty majority vertex side.  Since its complement
is also nonempty and the primal graph is connected, some edge crosses
\(K\).  There cannot be exactly one such edge \(e\): after deleting \(e\),
any path starting in \(K\) remains in \(K\), because the first edge by which
it left would be another boundary edge.  This contradicts bridge-freeness,
which says that the endpoints of \(e\) remain connected after \(e\) is
deleted.  Hence \(K\) has at least two boundary edges, equivalently at least
two literal boundary darts; choose two distinct ones.  This generic bridge
argument is checked as
\Sverified{Mettapedia.GraphTheory.FourColor.GoertzelV24RotationBoundaryBridge.not_boundary_singleton_of_edgeBridgeFree}
and its boundary-dart consequence as
\Sverified{Mettapedia.GraphTheory.FourColor.GoertzelV24ConnectedShoreLiteralNode.two_le_card_boundaryDart_of_connected_edgeBridgeFree}.

Cubicity gives three darts at every vertex, so a chosen vertex on each
nonempty side supplies its retained-dart root.  For graph-backed rotation
data, literal boundary darts are exactly the graph darts crossing the same
vertex predicate.  The injective charge from crossing darts to the edge-shore
middle set in Theorem~\ref{thm:connected-edge-shore-majority} therefore
transfers the middle-set bound to the literal width.  Packaging these choices
gives the asserted node.
\end{proof}

Thus the remaining supplier problem is no longer ``turn a branch cut into a
spherical splice'', choose roots, or prove that a proper shore has two ports.
It is to construct and package the connected branch-decomposition cuts with
strict laminar nesting and vertex accounting; the exact target-specific
Count state is then computed by the literal node consumer.

\begin{lemma}[equal phases are separated]
\label{lem:equal-depth-phases-separated}
\Sverified{Mettapedia.GraphTheory.FourColor.GoertzelV24SphereCutMaterial.period_le_sub_of_mod_eq}
If \(i<j\) and \(i\equiv j\pmod s\), then \(s\le j-i\).
\end{lemma}

\begin{proof}
The congruence says \(s\mid(j-i)\).  Since \(j-i>0\), a positive multiple
of \(s\) is at least \(s\).  (When \(s=0\), the conclusion is immediate.)
\end{proof}

The strict-material clause in a sphere-cut supply can consequently be made
finite-state rather than pictorial.  Root the branch-decomposition tree and
augment every typed state by the node depth modulo \(6w+1\).  Equal augmented
states at distinct comparable nodes are separated by at least \(6w+1\) tree
edges by Lemma~\ref{lem:equal-depth-phases-separated}.  At each step the
discarded sibling subtree contains a leaf, and the
sibling leaf sets at different steps are disjoint; hence the two nested edge
shores differ by at least \(6w+1\) graph edges.  Lemma~\ref{lem:cubic-sphere-cut-material}
then supplies a strict slab vertex.  The price is only the explicit finite
factor

\[
                         q_{\rm spaced}=(6w+1)q
\]
\Sverified{Mettapedia.GraphTheory.FourColor.GoertzelV24SphereCutDescent.card_spacedTypedState}

in the typed-state count.  This settles material accounting once the two
shore interfaces have been constructed; it does not itself construct their
ordered open-tangle carriers or prove that the splice retains exactly those
vertex shores.

There is an equivalent phase choice which is better adapted to the literal
edge-shore carrier.  If $A\subsetneq B$ are the descendant and ancestor
edge shores, attach to the interface state of a shore $S$ the residue

\[
                         |S|\pmod{6w+1}.
\]

\begin{theorem}[cardinality-phased shore descent]
\label{thm:cardinality-phased-shore-descent}
\Sverified{Mettapedia.GraphTheory.FourColor.GoertzelV24ShoreStateDescent.vertexCount_le_of_cardPhasedShore}
Let a reduced binary decomposition tree be labelled by finite edge shores,
strictly decreasing along every descent, and suppose every middle set has
size at most $w$.  Let $Q$ be any finite exact interface-state type and
let $q(S)\in Q$ be the state of shore $S$.  If equal $Q$-states admit a
target-preserving local splice whenever the slab contains a vertex all of
whose incident edges lie in the slab, then a minimal target instance has

\[
  |V|\le 2\left(2^{(6w+1)|Q|}-1\right).
\]
\end{theorem}

\begin{proof}
Use the augmented state

\[
       \widetilde q(S)=\bigl(|S|\bmod(6w+1),q(S)\bigr).
\]

Suppose comparable shores $A\subsetneq B$ have equal augmented states.
Strict inclusion gives $|A|<|B|$, while equality of residues says that
$6w+1$ divides $|B|-|A|$.  Hence

\[
                 6w+1\le |B|-|A|=|B\setminus A|.
\]

Lemma~\ref{lem:cubic-sphere-cut-material} supplies a vertex lying strictly
inside the slab.  Equality of the second components gives $q(A)=q(B)$, so
the assumed local splice produces a smaller target instance, contradicting
minimality.  Thus no descent repeats an augmented state.  There are exactly
$(6w+1)|Q|$ augmented states, and finite-tree pumping gives the displayed
bound.
\end{proof}

This formulation has the same finite cost as depth phasing, but the phase is
a function of the edge shore itself.  Consequently the connected branch
decomposition supplier need not carry or synchronize a separate node-depth
annotation.

The colouring part of the exact state can now be stated without identifying
the literal boundary-dart types of two different shores.  Let \(T\) be an
inner open tangle with boundary carrier \(B\), let \(|B|=k\), and choose a
coordinate equivalence \(c:B\simeq [k]\).  Define

\[
 \operatorname{Supp}_{k}(T,c)
   =\bigl\{x\in\{1,2,3\}^{[k]}:
       x\text{ is realized by a proper Tait colouring of }T\bigr\},
\]

where the port at coordinate \(i\) is \(c^{-1}(i)\).  This is an actual finite
set, not a predicate retained in an ambient carrier.

\begin{theorem}[normalized exact Count replacement]
\label{thm:normalized-exact-count-replacement}
\Sverified{Mettapedia.GraphTheory.FourColor.GoertzelV24NormalizedTaitSupport.not_composeRotationSystem_taitColorable_of_normalized_eq}
Let \(L\) be an outside open tangle whose seam is indexed by \([k]\), and let
\((T,c)\), \((T',c')\) be inner open tangles with \(k\)-port coordinates.  If

\[
             \operatorname{Supp}_{k}(T,c)
             =\operatorname{Supp}_{k}(T',c'),
\]

then replacing \(T\) by \(T'\) preserves zero Count, equivalently preserves
non-Tait-colourability of the literal sewn rotation system.
\end{theorem}

\begin{proof}
After transport through \(c\) and \(c'\), both inner pieces have the same
literal seam carrier \([k]\).  Equality of the displayed finite sets is
therefore equality of the genuine-word supports consumed by the closed Count
replacement theorem.  That equality is also equality of the raw colour-word
supports, because every boundary letter of a proper colouring is nonzero.
The gluing equivalence says that a closed composite is colourable exactly when
the outside and inside supports intersect.  Hence the two composites are
colourable simultaneously.  Finally, the physical/semantic bridge identifies
this closed Count statement with Tait colourability of the literal sewn
rotation systems.
\end{proof}

There are exactly

\[
              2^{3^k}
\]

possible normalized supports
\Sverified{Mettapedia.GraphTheory.FourColor.GoertzelV24NormalizedTaitSupport.card_normalizedSupportType}.
This is the raw logical state count; reachable-state minimization is only an
optimization.  The cyclic order is intentionally not forgotten.  If
\(\rho_T\) is the canonical hub rotation of \(T\), put

\[
 \Sigma_k(T,c)=
 \bigl(c\rho_Tc^{-1},\operatorname{Supp}_k(T,c)\bigr).
\]

This state is finite, of raw size

\[
                  k!\,2^{3^k},
\]

as machine-checked by
\Sverified{Mettapedia.GraphTheory.FourColor.GoertzelV24NormalizedSeamState.card_state}.

\begin{theorem}[the normalized state is a complete replacement receipt]
\label{thm:normalized-seam-state-replacement}
\Sverified{Mettapedia.GraphTheory.FourColor.GoertzelV24NormalizedSeamState.replacement_receipt}
Suppose an outside boundary with hub rotation \(\rho_L\) is matched
orientation-reversingly to \((T,c)\), and
\(\Sigma_k(T,c)=\Sigma_k(T',c')\).  Compose the old matching with
\(c'^{-1}\circ c\), from the old inner boundary to the new one.  Then
the resulting matching to \(T'\) is still orientation-reversing, and the two
normalized Tait supports are equal.
\end{theorem}

\begin{proof}
Equality of the first fields says exactly that the coordinate change conjugates
\(\rho_T\) to \(\rho_{T'}\).  Substituting this identity into the old
orientation-reversal equation gives the new equation pointwise.  Equality of
the second fields is the support equality of
Theorem~\ref{thm:normalized-exact-count-replacement}.
\end{proof}

Thus equality of the full finite receipt supplies both semantic facts required
by replacement: zero Count and seam orientation.  For a majority shore whose
literal boundary has cardinality \(k\), choose the canonical finite-cardinality
equivalence with \(\operatorname{ULift}(\operatorname{Fin} k)\).  Applying this
coordinate to the literal vertex-side open tangle and its canonical
first-return hub rotation gives an actual state \(\Sigma_k\).  Equality of two
such states produces, without any extra receipt field, both the transported
orientation-reversing matching and equality of the exact Tait supports:
\Sverified{Mettapedia.GraphTheory.FourColor.GoertzelV24MajorityShoreNormalizedState.replacement_receipt_of_normalizedState_eq}.
Thus the coordinate-equivalence packaging is closed.  The supplier must still
supply literal connected shores and a strict material witness; the two
disc-side hypotheses are then consequences of the planar-bond theorems, not
extra receipt fields.

\begin{theorem}[constructive physical replacement at equal normalized states]
\label{thm:constructive-majority-shore-physical-replacement}
\Sverified{Mettapedia.GraphTheory.FourColor.GoertzelV24MajorityShorePhysicalReplacement.strictPhysicalReplacement_of_normalizedState_eq}
Let \(M\) be a graph-backed non-Tait-colourable connected, bridgeless,
spherical cubic rotation map with cyclic vertex rotations and two-sided face
orbits.  Let \(S_{\rm new}\subseteq S_{\rm old}\) be finite edge shores such
that each shore and its edge complement are connected and each corresponding
majority vertex shore and complement are nonempty.  Choose the retained-dart
roots and two distinct darts on the old exterior boundary.  If the two
interior boundaries have common cardinality \(k\), their normalized states
are equal, and some vertex has its complete cubic star in
\(S_{\rm old}\setminus S_{\rm new}\), then these data construct a literal
rotation system \(M'\) such that
\begin{enumerate}
\item \(M'\) is connected, bridgeless, spherical and cubic, with cyclic
      vertex rotations;
\item \(M'\) is not Tait-colourable; and
\item \(\lvert V(M')\rvert<\lvert V(M)\rvert\).
\end{enumerate}
\end{theorem}

\begin{proof}
Take \(M'\) to be the literal matched-seam composite of the exterior of the
old majority shore and the interior of the new majority shore.  The matching
first crosses the old ambient cut by inverse edge flip and then transports
through the two common finite coordinates.  Equality of normalized states
contains both facts needed by the splice: equality of the exact Tait-word
supports and conjugacy of the canonical hub rotations.  The first, together
with exact boundary reindexing and complementary-shore reassembly, proves
that \(M'\) is not Tait-colourable.  The second transports the ambient
orientation-reversal equation to the new interior.

Connected complementary edge shores give connected majority vertex shores.
The planar-bond rank theorem therefore gives spherical face-simple canonical
closures on the retained sides; the transported opposite hub orders and the
seam face-count theorem make the composite spherical.  The literal gluing
theorems give connectedness, bridge-freeness, cubicity and cyclic vertex
rotations.  Finally majority is monotone under
\(S_{\rm new}\subseteq S_{\rm old}\), while the displayed cubic star belongs
to the old majority shore but not the new one.  Thus the retained new
interior is strictly smaller than the removed old interior, proving the
vertex inequality.  These three conclusions are the fields of the returned
replacement certificate; minimality is not used.
\end{proof}

\begin{theorem}[equal normalized majority-shore states contradict minimality]
\label{thm:majority-shore-physical-replacement}
\Sverified{Mettapedia.GraphTheory.FourColor.GoertzelV24MajorityShorePhysicalReplacement.no_nested_equal_normalizedState_of_strict_material}
Let \(M\) be a graph-backed vertex-minimal non-Tait-colourable member of the
connected bridgeless spherical cubic rotation-map class, with cyclic vertex
rotations and two-sided face orbits.  Let
\(S_{\rm new}\subseteq S_{\rm old}\)
be finite edge shores such that each shore and its edge complement are
connected and each corresponding majority vertex shore and complement are
nonempty.  Root the old exterior, old interior and new interior literal
vertex-side tangles by retained darts.  Suppose the two interior boundaries
have the same cardinality \(k\), their normalized states are equal, and the
old exterior boundary contains two distinct darts.  Finally suppose some
vertex \(x\) has every incident edge in
\(S_{\rm old}\setminus S_{\rm new}\).  These data are impossible.
\end{theorem}

\begin{proof}
Apply Theorem~\ref{thm:constructive-majority-shore-physical-replacement}.
It returns a strictly smaller map in the complete cap-stable structural
class and proves that map non-Tait-colourable.  Vertex minimality says the
same returned map is Tait-colourable, a contradiction.
\end{proof}

\begin{theorem}[direct literal corridor bound]
\label{thm:literal-majority-shore-chain-descent}
\Sverified{Mettapedia.GraphTheory.FourColor.GoertzelV24MajorityShoreStateDescent.length_le_of_literalShoreChain}
Let \(M\) be a graph-backed vertex-minimal non-Tait-colourable member of the
connected bridgeless spherical cubic rotation-map class.  Fix \(k,w\ge0\).
Suppose
\[
                 S_0\supsetneq S_1\supsetneq\cdots\supsetneq S_{m-1}
\]
is a chain of actual edge shores.  At every \(S_i\), suppose the shore and
its edge complement are connected, both majority vertex sides are nonempty,
the literal interior and exterior have retained-dart roots, the literal
boundary width is at most \(k\), the edge-shore middle set has at most \(w\)
vertices, and the exterior boundary contains two distinct darts.  Then
\[
       m\le (6w+1)\sum_{0\le j\le k}j!\,2^{3^j}.
\]
\end{theorem}

\begin{proof}
Label \(S_i\) by its exact normalized seam state in the dependent sum
\(\coprod_{0\le j\le k}\Sigma_j\), together with
\(|S_i|\bmod(6w+1)\).  There are exactly
\[
        (6w+1)\sum_{0\le j\le k}j!\,2^{3^j}
\]
such labels.  If two positions \(i<j\) had the same label, the strict
inclusion \(S_j\subsetneq S_i\) and equality of phases would put a cubic
vertex wholly inside \(S_i\setminus S_j\), by
Lemma~\ref{lem:cubic-sphere-cut-material}.  Equality of the dependent seam
states gives equality of the two normalized literal states at their common
width.  Theorem~\ref{thm:majority-shore-physical-replacement} would then
splice out the nonempty slab and produce a smaller non-Tait-colourable map
in the same structural class, contradicting vertex minimality.  Hence the
label map is injective, and the displayed cardinality is the claimed bound.
\end{proof}

This is the direct consumer for an L1 corridor.  It uses cumulative states of
the actual prefixes, so it needs neither the older L2 self-loop premise nor a
caller-supplied replacement operation.  The still-open geometric adapter is
not hidden by the statement: from the source corridor positions one must
construct the full ambient edge shores and verify connectedness of both edge
shores, nonemptiness of both majority sides, roots, the width and middle-set
bounds, two distinct exterior boundary darts, and strict nesting.  The
theorem above proves only that any chain carrying those literal certificates
cannot be longer than the finite state carrier.

\begin{theorem}[bounded-width connected-shore descent]
\label{thm:literal-majority-shore-tree-descent}
\Sverified{Mettapedia.GraphTheory.FourColor.GoertzelV24ConnectedShoreLiteralNode.vertexCount_le_of_connectedShoreTree}
Let \(M\) be a graph-backed vertex-minimal non-Tait-colourable member of the
connected bridgeless spherical cubic rotation-map class.  Fix \(k,w\ge0\),
and let \(\mathcal T\) be a finite rooted binary tree.  At every node \(t\),
suppose we are given an actual finite edge shore \(S_t\) such that:
\begin{enumerate}
\item \(S_t\) and its edge complement are connected, and the corresponding
majority vertex shore and its complement are nonempty;
\item the edge-shore middle set has cardinality at most \(w\);
\item the same middle set has cardinality at most \(k\);
\item along every root-to-leaf path, a later shore is a proper subset of every
earlier shore.
\end{enumerate}
If \(|V(M)|\le2|V(\mathcal T)|\), then, with
\[
             q_k=\sum_{0\le j\le k} j!\,2^{3^j},
\]
one has
\[
             |V(M)|\le
             2\left(2^{(6w+1)q_k}-1\right).
\]
\end{theorem}

\begin{proof}
Complete every supplied node by
Lemma~\ref{lem:connected-shore-literal-node}.  Thus its roots and two distinct
boundary darts are derived from the minimal map, and its literal boundary
width is at most \(k\).
For a node of literal boundary width \(j\le k\), compute its normalized seam
state directly from its displayed shore, interior root, canonical boundary
coordinate and first-return hub rotation.  Embed it in the single dependent
sum
\[
       Q_k=\coprod_{0\le j\le k}\Sigma_j,
       \qquad |Q_k|=q_k,
\]
and augment it by \(|S_t|\bmod(6w+1)\).  Thus the tree is labelled by exactly
\((6w+1)q_k\) possible states.

Suppose two comparable nodes repeated an augmented state.  Proper nesting and
equality of the cardinality phases give, by
Lemma~\ref{lem:cubic-sphere-cut-material}, a vertex whose three incident edges
all lie in the strict edge slab between the old and new shores.  Equality in
the dependent sum first identifies the two literal widths and then transports
the two normalized states to an ordinary equality at that common width.  This
dependent transport is checked explicitly in
\Sverified{Mettapedia.GraphTheory.FourColor.GoertzelV24MajorityShoreStateDescent.boundedNormalizedState_eq_elim};
it is not a caller-supplied cast.  All remaining hypotheses of
Theorem~\ref{thm:majority-shore-physical-replacement} are fields of the two
literal tree nodes.  That theorem therefore makes the repeated state
impossible.  No root-to-leaf descent repeats a state, so finite binary-tree
pumping and \(|V(M)|\le2|V(\mathcal T)|\) give the displayed bound.
\end{proof}

This theorem removes both the abstract seam-type parameter and the former
application-specific replacement premise from the bounded-width consumer.
The finite state is computed from literal shores, and repeated states are fed
directly to the physical splice.  Its vertex-accounting premise is discharged
for a genuine edge-leaf branch decomposition by the following theorem.

\begin{theorem}[bounded-width rooted connected branch decomposition]
\label{thm:rooted-connected-branch-decomposition}
\Sverified{Mettapedia.GraphTheory.FourColor.GoertzelV24ConnectedBranchDecompositionForest.vertexCount_le_of_rootedConnectedBranchDecomposition}
Let \(M\) be as in Theorem~\ref{thm:literal-majority-shore-tree-descent}.
Root a full binary edge-leaf decomposition at one distinguished edge \(e_0\).
After removing that root leaf, write its first fork as two edge-leaf trees
\(T_L,T_R\).  Suppose
\begin{enumerate}
\item \(e_0\), followed by the leaves of \(T_L\) and \(T_R\), is a
duplicate-free enumeration of \(E(M)\);
\item at every internal fork of \(T_L\) or \(T_R\), the set \(S\) of leaves
below that fork and \(E(M)\setminus S\) are connected edge shores, their
middle set has size at most both \(k\) and \(w\).
\end{enumerate}
Set
\[
             q_k=\sum_{0\le j\le k}j!\,2^{3^j},
             \qquad B=2^{(6w+1)q_k}-1.
\]
Then
\[
                         |V(M)|\le 4B+6.
\]
\end{theorem}

\begin{proof}
An internal fork is labelled by the finite set of edge leaves below it.  Distinct
leaves make the shores of its two nontrivial children proper subsets of the
parent shore.  Transitivity therefore gives proper nesting between every two
comparable internal forks.  Erasing the edge leaves produces two trees of the
kind consumed by Theorem~\ref{thm:literal-majority-shore-tree-descent}.

It remains to check that each displayed edge cut has vertices on both
majority sides; this is a consequence, not an additional hypothesis.  The
fork shore contains at least two edges, one below each child.  Its complement
also contains at least two distinct edges: the distinguished root edge
\(e_0\), and one edge in the opposite top-level tree.  These same two outside
edges remain outside every descendant shore.  If a connected edge shore
contains two distinct edges, take a shore-edge walk between them.  Either the
first step is a different edge from the first chosen edge, or delete that
step and continue; induction produces a vertex incident with two distinct
shore edges.  Thus the shore has a majority vertex.  Applying the same
argument to the complementary edge shore gives a vertex incident with two
complement edges.  Since every vertex of \(M\) has exactly three incident
edges, that vertex is not a majority vertex for the original shore.  Hence
both majority sides are nonempty at every internal fork.  The recursive
construction of these witnesses is checked by
\Sverified{Mettapedia.GraphTheory.FourColor.GoertzelV24ConnectedBranchDecompositionForest.everyForkConnected_of_geometry}.

The node-count half of that theorem, checked separately in
\Sverified{Mettapedia.GraphTheory.FourColor.GoertzelV24ConnectedShoreLiteralNode.nodeCount_le_of_connectedShoreTree},
bounds each tree by \(B\) internal nodes.

A full binary tree has one more leaf than internal forks.  Consequently the
complete leaf enumeration gives
\[
       |E(M)|=i(T_L)+i(T_R)+3,
\]
where \(i(T)\) denotes the number of internal forks of \(T\).
Cubicity gives \(3|V(M)|=2|E(M)|\).  Hence
\[
 |V(M)|\le2\bigl(i(T_L)+i(T_R)+3\bigr)\le4B+6.
\]
The one-edge cut above the first fork is deliberately absent: its
complementary majority shore is empty, so inserting it as a root node would
make the prettier one-tree constant depend on a false interface.
\end{proof}

Thus proper nesting and vertex accounting are no longer supply fields.  What
remains on the supply side is precise: construct the complete edge-leaf tree
from the source corridor or a source-faithful connected decomposition and
provide its nodewise connected edge shores and middle-set bounds.  For every
retained internal fork, both edge sides contain at least two edges: the fork
side contains two leaf branches, while its complement contains the
distinguished root edge and at least one leaf outside the fork.  In a
connected edge shore containing two edges, some vertex is incident with two
shore edges (otherwise distinct edges would lie in different components), so
the two majority sides are nonempty.  The graph-theoretic implication and its
complementary cubic form are checked in
\Sverified{Mettapedia.GraphTheory.FourColor.GoertzelV24ConnectedBranchDecompositionForest.ConnectedCutCondition.ofConnectedTwo};
the two-edge propagation itself is checked in
\Sverified{Mettapedia.GraphTheory.FourColor.GoertzelV24ConnectedBranchDecompositionForest.EdgeLeafTree.two_le_card_shore_fork}
and
\Sverified{Mettapedia.GraphTheory.FourColor.GoertzelV24ConnectedBranchDecompositionForest.two_le_card_complement_of_two_outside}.
Accordingly, the theorem above receives no majority-nonemptiness fields.
Lemma~\ref{lem:connected-shore-literal-node} then supplies the roots, literal
width, and two-port fields.

For clarity, the standard branch-decomposition input is now separated from
this consumer.  Choose one leaf of the usual unrooted ternary tree, delete
that leaf and its incident internal vertex, and call the two remaining
binary edge-leaf trees $T_L,T_R$.  The resulting rooted encoding
$D=(e_0,T_L,T_R)$ remembers only the duplicate-free complete leaf labelling.
Write $\operatorname{Conn}(D)$ when both edge-induced shores are connected at
every tree cut, including the distinguished root-leaf cut and every remaining
leaf cut, and $\operatorname{Width}_w(D)$ when every such middle set has size
at most $w$.  These are independent predicates.  The checked adapter first
forgets only the root and leaf cuts that the descent does not consume, then
converts the remaining internal cuts to the exact literal consumer via
\Sverified{Mettapedia.GraphTheory.FourColor.GoertzelV24ConnectedBranchDecompositionAdapter.everyForkGeometry_of_connected_width}
and
\Sverified{Mettapedia.GraphTheory.FourColor.GoertzelV24ConnectedBranchDecompositionAdapter.rootedBranchDecompositionToConsumerAtWidth}.
Thus ``a bounded decomposition exists'' and ``it can be connected without
increasing its width'' can no longer be conflated in one supply field.

For an ordered boundary of size at most \(k\), the manuscript's data---edge
colours, the two noncrossing connectivity relations, capped face progress,
and the target-specific accepting bit---form a finite typed context state.
Theorem~\ref{thm:finite-tree-interface-pumping} therefore shows that
\emph{bounded interface width is already enough}; a long literal tube is one
convenient decomposition, not the logical essence of the descent.

The planar grid theorem of Robertson and Seymour
(``Graph minors III: planar tree-width'', 1984,
doi:10.1016/0095-8956(84)90013-3) says that excluding one fixed planar
wall as a minor bounds treewidth.  The Seymour--Thomas sphere-cut theory
supplies noose-respecting branch decompositions of plane graphs
(Call routing and the ratcatcher, 1994,
doi:10.1007/BF01215352).  A more direct supplier is available here, and it
is now checked.
Fomin, Fraigniaud and Thilikos (The price of connectedness in expansions,
Technical Report LSI--04--28--R, 2004,
\url{https://hdl.handle.net/2117/98590}, Theorem~1; full proof in
Barri\`ere, Flocchini, Fomin, Fraigniaud, Nisse, Santoro and Thilikos,
Connected graph searching, Information and Computation 219 (2012),
doi:10.1016/j.ic.2012.08.004, Lemma~1) prove that, given a width-$w$
branch decomposition of a \(2\)-edge-connected graph, their algorithm
\emph{Make-it-Connected} returns a connected branch decomposition of width
at most $w$.  Here ``connected'' means that both sets of edge leaves at
every tree cut induce connected subgraphs; there is no hidden factor in the
width.

\begin{theorem}[width-preserving connectedization]
\label{thm:width-preserving-connectedization}
\Sverified{Mettapedia.GraphTheory.FourColor.GoertzelV24WidthPreservingConnectedization.widthPreservingConnectedization}
Let $G$ be a finite simple \(2\)-edge-connected graph and let $T$ be a
complete rooted branch decomposition of $G$ of width at most $w$.  Then $G$
has a complete rooted branch decomposition of width at most $w$ in which
both edge shores of every tree cut, including the leaf cuts, induce
connected subgraphs.
\end{theorem}

\begin{proof}
Root the decomposition at its root leaf and carry, with every subtree, the
set of edges outside it (its \emph{up-set}).  Every internal vertex of the
ternary tree is then a fork whose three directions are the two child shores
and the up-set.  Every tree edge between two internal vertices is a fork
child of a fork.  Such a tree edge is a \emph{quartet} when one pair of
directions is vertex-disjoint while the shores can be re-paired across the
edge with nonempty boundaries.

Re-pairing is a rotation of the fork-of-fork.  Two finite boundary facts
drive the argument.  First, the two new cut sizes sum to at most the sum of
the two old cut sizes.  If only one crossed matching is a quartet, its new
boundary is contained in one old boundary.  Thus some quartet replacement
never increases the width.  Second, at a fork with directions $S_1,S_2,S_3$
put
\[
 \psi(S_1,S_2,S_3)=
 \sum_{\{i,j,\ell\}=\{1,2,3\}}
 [\partial(S_i,S_j)=\varnothing]\,(|S_\ell|-1),
\]
and sum $\psi$ over all forks.  Direct boundary calculation shows that every
allowed quartet replacement strictly decreases this natural-number
potential.  Strong induction therefore produces a quartet-free tree on the
same edge leaves without increasing its width.

It remains to prove this tree connected.  In a quartet-free tree of a
$2$-edge-connected graph, every fork has pairwise-touching directions.  A
vertex-disjoint pair can be pushed across the third direction: quartet
freeness makes the merged pair vertex-disjoint from one of the two directions
beyond it.  Pairs involving the up-set are pushed down the rooted tree and
sibling pairs are pushed up; structural induction reduces either case to a
single-edge shore.  Deleting that edge cannot disconnect a
$2$-edge-connected graph, a contradiction.  Connected directions that touch
at a vertex have connected union, so both shores of every cut, including
leaf cuts, are connected.
\end{proof}

The graph-side fact that a connected graph without bridges is
\(2\)-edge-connected is checked as
\Sverified{Mettapedia.GraphTheory.FourColor.GoertzelV24ConnectedBranchDecompositionAdapter.isEdgeConnected_two_of_connected_no_bridges}.
Theorem~\ref{thm:connected-edge-shore-majority} converts its edge shores
canonically into nested connected complementary \emph{vertex} shores.  In
a cubic graph their literal edge boundary has size at most the original
middle-set width, not twice that width.  The planar-bond cap and orientation
theorems then construct the ordered open tangles directly from the ambient
rotation system.  Thus a sphere-cut noose and its local perturbation are an
optional geometric realization, not part of the load-bearing D3 supply.
Lemma~\ref{lem:cubic-sphere-cut-material} and the depth phase above settle
strict-material accounting.  Theorem~\ref{thm:rooted-connected-branch-decomposition}
now roots the complete edge-leaf decomposition, discards its invalid
one-edge top cut, derives proper nesting in the two surviving subtrees, and
performs the full edge/vertex accounting.  The majority shores, literal
nodes, and exact ordered Count states are then computed by
Theorem~\ref{thm:connected-edge-shore-majority},
Lemma~\ref{lem:connected-shore-literal-node}, and
Theorem~\ref{thm:rooted-connected-branch-decomposition}.

\begin{proposition}[wall-exclusion reduction, relative form]
\label{prop:target-forced-wall}
\Sconditional{planar wall theorem; fixed-wall exclusion; route-native finite base}
Fix the normal-form class and one direct-Count or Lift target.  Suppose there
are computable \(r,w\) such that:
\begin{enumerate}
\item no vertex-minimal zero-target instance contains a subdivision of the
\(r\)-wall;
\item every plane instance in the class with no such wall has a complete
rooted branch decomposition of width at most \(w\);
\end{enumerate}
Then every vertex-minimal zero-target instance has an effective vertex
bound.  If the route-native finite base below that bound is empty, there is
no zero-target instance.
\end{proposition}

\begin{proof}
Choose a vertex-minimal zero-target instance \(X\).  Hypothesis~1 excludes
the fixed wall.  Since a wall is subcubic, a minor model can be pruned inside
each branch set to a subdivision: the at most three incident model edges are
joined by a minimal tree, whose three terminal paths have one common branch
point and are otherwise internally disjoint.  Thus the minor and topological
forms agree for this wall in a cubic host.  Hypothesis~2 supplies the raw
width-$w$ tree, and Theorem~\ref{thm:width-preserving-connectedization}
connectedizes it without increasing the width.  The checked
standard-to-consumer adapter supplies the complete connected edge-leaf
decomposition.  Theorem~\ref{thm:rooted-connected-branch-decomposition}
derives its two strictly nested shore trees and gives the effective vertex
bound; internally, a repeated normalized state is ruled out by the exact
strict replacement of Theorem~\ref{thm:majority-shore-physical-replacement}.
An empty finite base rules out \(X\).
\end{proof}

The high-width receipt can be split one step more finely, so that the
route-specific assertion is not conflated with the conventional planar wall
theorem.

\begin{proposition}[exact fixed-mesh factorization of the width input]
\label{prop:fixed-mesh-width-factorization}
\Sverified{Mettapedia.GraphTheory.FourColor.GoertzelV24FixedMeshWidthFactorization.rawBranchDecompositionSupply_of_fixedInjectiveMeshExclusion}
Fix positive integers $a,b,w$.  Suppose:
\begin{enumerate}
\item no graph-backed vertex-minimal Tait counterexample contains an
injective $a$-by-$b$ mesh; and
\item every such counterexample which contains no injective
$a$-by-$b$ mesh has a complete rooted branch decomposition of width at
most $w$.
\end{enumerate}
Then every graph-backed vertex-minimal Tait counterexample has a complete
rooted branch decomposition of width at most $w$.  Conversely, if that raw
width-$w$ conclusion holds and
\[
  B(w,w)<ab,
\]
then the fixed injective $a$-by-$b$ mesh is absent from every such
counterexample.
\Sverified{Mettapedia.GraphTheory.FourColor.GoertzelV24FixedMeshWidthFactorization.fixedInjectiveMeshExclusion_of_rawBranchDecompositionSupply}
With the route-native finite base through $B(w,w)$, the two displayed
hypotheses therefore imply the combinatorial spherical Four-Colour
statement.
\Sverified{Mettapedia.GraphTheory.FourColor.GoertzelV24FixedMeshWidthFactorization.combinatorialFourColorStatement_of_fixedMeshExclusion_and_base}
\end{proposition}

The second premise now has an explicit target-class supplier for
$a\geq2$, proved in Section~\ref{sec:cotree-path-size}.  Its width is
coarse but computable.  The first premise, fixed permissive-mesh exclusion,
and the finite base remain separate open inputs.

\begin{proof}
For a fixed least counterexample, the first hypothesis supplies absence of
the stated injective mesh.  Feed that fact to the second hypothesis; the
result is exactly the raw branch-decomposition supplier consumed by the
checked width-preserving connectedization and reductive assembly.

For the converse, a raw width-$w$ supplier gives the checked vertex bound

\[
                         |V(G)|\le B(w,w).
\]

An injective $a$-by-$b$ mesh embeds its $ab$ branch positions into
$V(G)$, so it gives $ab\le |V(G)|$, contrary to $B(w,w)<ab$.
Finally, compose the raw supplier from the first paragraph with the checked
finite-base assembly.
\end{proof}

The cut-counting mesh in the preceding proposition is intentionally more
permissive than a topological wall: it does not remember the order in which
its designated branch vertices occur on a path, and a path need not be
simple.  Those omissions do not affect mesh isoperimetry, but they do affect
the two-dimensional overlap transport below.  The route-specific target can
therefore be weakened to the following wall-like carrier.

\begin{proposition}[ordered fixed-mesh factorization]
\label{prop:ordered-fixed-mesh-width-factorization}
An \emph{ordered injective \(a\)-by-\(b\) mesh} is an injective mesh together
with simple row and column paths and a selected position of every branch
vertex on its row and column, strictly increasing in the other coordinate.
Its underlying permissive mesh is vertex-injective.
\Sverified{Mettapedia.GraphTheory.FourColor.GoertzelV24OrderedInjectiveMeshWidthFactorization.toMesh_isVertexInjective}

Suppose that no graph-backed vertex-minimal Tait counterexample contains an
ordered injective \(a\)-by-\(b\) mesh, and that absence of this carrier
supplies a complete rooted branch decomposition of width at most \(w\).
Then the raw width-\(w\) supplier follows.
\Sverified{Mettapedia.GraphTheory.FourColor.GoertzelV24OrderedInjectiveMeshWidthFactorization.rawBranchDecompositionSupply_of_fixedOrderedInjectiveMeshExclusion}
Conversely, a raw width-\(w\) supplier excludes the ordered mesh whenever
\(B(w,w)<ab\).
\Sverified{Mettapedia.GraphTheory.FourColor.GoertzelV24OrderedInjectiveMeshWidthFactorization.fixedOrderedInjectiveMeshExclusion_of_rawBranchDecompositionSupply}
Together with the route-native finite base through \(B(w,w)\), the two
hypotheses imply the combinatorial spherical Four-Colour statement.
\Sverified{Mettapedia.GraphTheory.FourColor.GoertzelV24OrderedInjectiveMeshWidthFactorization.combinatorialFourColorStatement_of_fixedOrderedMeshExclusion_and_base}
The stronger exclusion from
Proposition~\ref{prop:fixed-mesh-width-factorization} implies the ordered
exclusion.
\Sverified{Mettapedia.GraphTheory.FourColor.GoertzelV24OrderedInjectiveMeshWidthFactorization.fixedOrderedInjectiveMeshExclusion_of_fixedInjectiveMeshExclusion}
\end{proposition}

\begin{proof}
Forgetting the positions and path-simplicity witnesses gives the old mesh;
the branch-injectivity field is exactly the existing vertex-injectivity
predicate.  The first implication feeds ordered-mesh absence to the supplied
wall-free decomposition theorem.  The converse forgets the ordering and
uses the same injection of the \(ab\) branch positions into \(V(G)\) as in
Proposition~\ref{prop:fixed-mesh-width-factorization}.  The finite-base
conclusion then uses the unchanged checked reductive assembly.
\end{proof}

The first premise of
Proposition~\ref{prop:ordered-fixed-mesh-width-factorization} is the sole
route-specific high-width assertion used below.  Its second premise is the
conventional wall/grid-to-branchwidth implication, now expressed on a
carrier which records the rectangular order needed by the overlap
connection.  Neither premise entails the other without the quantitative
size condition in the converse, and neither is concealed in a Lean
structure field.  The permissive fixed-mesh proposition remains a valid
stronger alternative, but it is not imposed on the wall-facing proof.
The supplier in Section~\ref{sec:cotree-path-size} discharges the
\emph{permissive}-mesh decomposition premise only.  Absence of ordered
meshes does not imply absence of permissive meshes.  Section~\ref{sec:ordered-mesh-supply}
closes the ordered supplier separately: it constructs ordered meshes by
first-hit ordering and simultaneous monotone thinning, then applies the
same complete-cut descent with an explicit larger width.  Fixed ordered
exclusion and the actual base remain open.

There is an exact matching-theoretic formulation of the colour obstruction
which will be useful below.  It is global rather than a bounded-radius
configuration test.

\begin{theorem}[Tait colourings are exact three-matching decompositions]
\label{thm:tait-three-matching-decomposition}
\Sverified{Mettapedia.GraphTheory.FourColor.GoertzelV24TaitMatchingDecomposition.taitColorable_iff_nonempty_threeMatchingDecomposition}
Let \(G\) be a finite cubic simple graph.  An \emph{exact three-matching
decomposition} of \(G\) consists of three fixed-point-free involutions
\(\sigma_0,\sigma_1,\sigma_2\) of \(V(G)\) such that
\begin{enumerate}
\item \(v\sigma_i(v)\) is an edge of \(G\), for every \(v\) and \(i\);
\item the three vertices \(\sigma_0(v),\sigma_1(v),\sigma_2(v)\) are distinct,
      for every \(v\); and
\item every edge \(v w\) of \(G\) satisfies \(\sigma_i(v)=w\) for some \(i\).
\end{enumerate}
Then \(G\) is Tait colourable if and only if it has an exact three-matching
decomposition.  In particular, a cubic Tait counterexample has no such
decomposition.
\Sverified{Mettapedia.GraphTheory.FourColor.GoertzelV24TaitMatchingDecomposition.no_threeMatchingDecomposition_of_not_taitColorable}
\end{theorem}

\begin{proof}
Given a Tait colouring, take the three nonzero colour classes.  At every
cubic vertex the three incident edges have distinct nonzero colours, so each
colour occurs exactly once.  Each colour class is therefore a graph-supported
perfect matching, the three partners are distinct, and the matchings cover
all edges.

Conversely, an edge in an exact three-matching decomposition has a unique
matching index.  Colour it by the corresponding one of the three nonzero
Tait colours.  If two distinct incident edges had the same index, the
involution of that matching would give one vertex two partners, a
contradiction.  Thus the resulting edge colouring is proper and nowhere
zero.  The Lean proof constructs both translations on the literal
simple-graph and shared perfect-matching types.
\end{proof}

The three-matching formulation has a sharper two-matching form.  The third
matching is forced by cubicity rather than supplied as additional data.

\begin{corollary}[the two-matching criterion]
\label{cor:tait-two-matching-criterion}
\Sverified{Mettapedia.GraphTheory.FourColor.GoertzelV24TaitMatchingPair.taitColorable_iff_hasTwoDisjointSupportedPairings}
\Sverified{Mettapedia.GraphTheory.FourColor.GoertzelV24TaitMatchingPair.exists_shared_edge_of_not_taitColorable}
Let \(G\) be a finite cubic simple graph.  Then \(G\) is Tait colourable if
and only if it has two graph-supported perfect matchings \(\sigma\) and
\(\tau\) such that
\[
                 \sigma(v)\ne\tau(v)\qquad(v\in V(G)).
\]
Consequently, in a cubic Tait counterexample every two graph-supported
perfect matchings share an edge.
\end{corollary}

\begin{proof}
Two colour classes of a Tait colouring give the forward implication.  For
the converse, fix supported matchings \(\sigma,\tau\) with distinct partners
at every vertex.  At a vertex \(v\), the two edges
\(v\sigma(v)\) and \(v\tau(v)\) are distinct.  Since \(G\) is cubic, there
is a unique third incident edge; write it as \(v\rho(v)\).

This local definition is consistent at the other endpoint.  Indeed, if
\(v\rho(v)=vw\), then \(v\ne\sigma(w)\): otherwise involutivity of
\(\sigma\) would give \(w=\sigma(v)\), contrary to the choice of the third
edge at \(v\).  Similarly \(v\ne\tau(w)\).  Thus the same edge is the
unique third edge at \(w\), and \(\rho(w)=v\).  Adjacency also gives
\(\rho(v)\ne v\).  Hence \(\rho\) is a graph-supported fixed-point-free
involution.  The three matchings \(\sigma,\tau,\rho\) have distinct partners
at every vertex and cover all edges, so
Theorem~\ref{thm:tait-three-matching-decomposition} gives a Tait colouring.

Taking the contrapositive, two supported perfect matchings in a cubic Tait
counterexample cannot be pointwise edge-disjoint.  Equality
\(\sigma(v)=\tau(v)\) means that they share the undirected edge
\(v\sigma(v)\), proving the final assertion.
\end{proof}

The first matching required by this formulation is automatic in the normal
target class; only the disjoint second matching is difficult.

\begin{theorem}[Petersen's starting matching]
\label{thm:petersen-starting-matching}
\Sverified{Mettapedia.GraphTheory.FourColor.GoertzelV24PetersenStartingMatching.petersen_supportedPairing}
Every finite bridge-free cubic simple graph has a graph-supported perfect
matching.
\end{theorem}

\begin{proof}
We verify Tutte's condition.  Let \(S\subseteq V(G)\), and let
\(C_1,\ldots,C_q\) be the odd components of \(G-S\).  For each \(C_i\),
degree summation in the cubic graph gives
\[
  3|C_i|=2|E(C_i)|+|\delta(C_i)|.
\]
Thus \(|\delta(C_i)|\) is odd.  It is nonzero, and it cannot equal one:
the unique edge in a one-edge boundary would be a bridge.  Consequently
\(|\delta(C_i)|\geq 3\).

Every boundary edge of a component of \(G-S\) has its other endpoint in
\(S\), and boundaries of distinct components are edge-disjoint.  Hence
\[
  3q
  \leq \sum_{i=1}^{q}|\delta(C_i)|
  \leq \sum_{v\in S}\deg_G(v)
  =3|S|.
\]
Therefore \(q\leq |S|\) for every \(S\), which is Tutte's condition; Tutte's
theorem supplies a perfect matching.  Turning its unique matched neighbour at
each vertex into a fixed-point-free involution gives the graph-supported
pairing used above.
\end{proof}

\begin{corollary}[the inhabited overlap family]
\label{cor:inhabited-overlap-family}
\Sverified{Mettapedia.GraphTheory.FourColor.GoertzelV24PetersenStartingMatching.exists_startingPairing_all_supportedPairings_overlap}
Let \(G\) be a finite bridge-free cubic simple graph that is not Tait
colourable.  There is a graph-supported perfect matching \(\sigma\) such that
every graph-supported perfect matching \(\tau\) shares an edge with \(\sigma\).
\end{corollary}

\begin{proof}
Choose \(\sigma\) by Theorem~\ref{thm:petersen-starting-matching}.  If some
supported \(\tau\) shared no edge with it, the two-matching criterion would
give a Tait colouring, contrary to hypothesis.
\end{proof}

The overlap obstruction has an equivalent parity carrier which more closely
resembles the source's path--matching anatomy.

\begin{theorem}[the residual two-factor obstruction]
\label{thm:residual-two-factor-obstruction}
\Sverified{Mettapedia.GraphTheory.FourColor.GoertzelV24ResidualTwoFactor.card_neighborFinset_residualGraph_eq_two}
\Sverified{Mettapedia.GraphTheory.FourColor.GoertzelV24ResidualTwoFactor.residualGraph_isCycles}
\Sverified{Mettapedia.GraphTheory.FourColor.GoertzelV24ResidualTwoFactor.taitColorable_iff_exists_residual_supportedPairing}
\Sverified{Mettapedia.GraphTheory.FourColor.GoertzelV24ResidualTwoFactor.exists_isPerfectMatching_of_bipartite_degree_two}
\Sverified{Mettapedia.GraphTheory.FourColor.GoertzelV24ResidualTwoFactor.taitColorable_of_residualGraph_isBipartite}
\Sverified{Mettapedia.GraphTheory.FourColor.GoertzelV24ResidualTwoFactor.residualGraph_not_isBipartite_of_not_taitColorable}
\Sverified{Mettapedia.GraphTheory.FourColor.GoertzelV24ResidualTwoFactor.exists_odd_closedWalk_residual_of_not_taitColorable}
\Sverified{Mettapedia.GraphTheory.FourColor.GoertzelV24ResidualTwoFactor.exists_startingPairing_overlap_and_odd_residual}
Let \(G\) be a finite cubic simple graph and let \(\sigma\) be a
graph-supported perfect matching.  Delete the edge \(v\sigma(v)\) at every
vertex and call the resulting spanning graph \(R_\sigma\).  Then:
\begin{enumerate}
\item every vertex of \(R_\sigma\) has degree two;
\item a perfect matching \(\tau\) is supported by \(R_\sigma\) exactly when
      it is supported by \(G\) and shares no edge with \(\sigma\); and
\item if \(R_\sigma\) is bipartite, then \(G\) is Tait colourable.
\end{enumerate}
Consequently, in a cubic Tait counterexample the residual graph of every
supported perfect matching is non-bipartite and has an odd closed walk.  If
\(G\) is also bridge-free, one may choose \(\sigma\) so that every supported
perfect matching overlaps it, while \(R_\sigma\) still has such an odd closed
walk.
\end{theorem}

\begin{proof}
At each vertex of a cubic graph, deleting the one neighbour selected by
\(\sigma\) leaves exactly two neighbours.  This proves the first assertion
and identifies \(R_\sigma\) as a spanning two-factor.  The definition of its
adjacency gives the second assertion literally.  Therefore the two-matching
criterion can be rewritten as
\[
 G\text{ is Tait colourable}
 \quad\Longleftrightarrow\quad
 \text{some supported }\sigma\text{ has a perfect matching in }R_\sigma.
\]

It remains to justify the bipartite shortcut without appealing to a cycle
picture.  More generally, let \(H\) be a finite bipartite graph in which
every vertex has degree two.  For any vertex set \(S\), count incidences
between \(S\) and its neighbourhood \(N(S)\).  The vertices of \(S\)
contribute exactly \(2|S|\) incidences.  Each vertex of \(N(S)\) receives at
most two, so
\[
                         2|S|\leq 2|N(S)|.
\]
Thus \(|S|\leq |N(S)|\), and Hall's theorem gives a perfect matching of
\(H\).  Apply this to \(R_\sigma\); the resulting residual matching is
disjoint from \(\sigma\), so the two-matching criterion gives a Tait
colouring of \(G\).

The contrapositive says that every residual graph in a Tait counterexample
is non-bipartite.  The standard characterization of bipartite graphs by even
closed walks then supplies an odd closed walk.  Finally Petersen's theorem
chooses \(\sigma\) in the bridge-free target class, and
Corollary~\ref{cor:inhabited-overlap-family} supplies the simultaneous overlap
statement.
\end{proof}

Thus the high-width problem is an intersection problem for an inhabited
family of matchings, and equivalently an odd-residual-two-factor problem; it
is not an existence problem for the first matching.  This formulation also
fixes a quantifier issue in applying the source's two-paths-plus-matching
anatomy.  That anatomy is obtained inside an already coloured closed-web
carrier, so it cannot by itself manufacture the disjoint second matching in
an uncoloured zero-target counterexample.  A wall-facing use of that source
lemma must instead eliminate the odd residual obstruction or construct the
strict exact-support replacement required by the reductive descent.

The obstruction admits a well-founded potential which makes the required
wall surgery exact rather than qualitative.

\begin{theorem}[the residual-defect descent principle]
\label{thm:residual-defect-descent}
\Sverified{Mettapedia.GraphTheory.FourColor.GoertzelV24ResidualDefectDescent.residualDefect_eq_zero_iff_isBipartite}
\Sverified{Mettapedia.GraphTheory.FourColor.GoertzelV24ResidualDefectDescent.exists_supportedPairing_minimal_residualDefect}
\Sverified{Mettapedia.GraphTheory.FourColor.GoertzelV24ResidualDefectDescent.residualDefect_pos_of_not_taitColorable}
\Sverified{Mettapedia.GraphTheory.FourColor.GoertzelV24ResidualDefectDescent.taitColorable_of_residualDefect_descent}
\Sverified{Mettapedia.GraphTheory.FourColor.GoertzelV24ResidualDefectDescent.exists_positive_minimal_residualDefect_of_not_taitColorable}
\Sverified{Mettapedia.GraphTheory.FourColor.GoertzelV24ResidualDefectDescent.exists_positive_minimal_residualDefect_of_cubic_edgeBridgeFree}
Let \(G\) be a finite cubic simple graph with at least one graph-supported
perfect matching.  For a supported matching \(\sigma\), let
\[
 d_G(\sigma)=\#\{C:C\text{ is a connected component of }R_\sigma
                         \text{ and }C\text{ is not bipartite}\}.
\]
Then a supported matching minimizing \(d_G\) exists.  Every such minimizer
\(\sigma_0\) satisfies
\[
 d_G(\sigma_0)\leq d_G(\tau)
 \qquad\text{for every supported perfect matching }\tau.
\]
Moreover:
\begin{enumerate}
\item if \(G\) is not Tait colourable, then \(d_G(\sigma_0)>0\), and no
      supported perfect matching has smaller defect;
\item if every supported \(\sigma\) with \(d_G(\sigma)>0\) admits a supported
      \(\tau\) with \(d_G(\tau)<d_G(\sigma)\), then \(G\) is Tait colourable.
\end{enumerate}
\end{theorem}

\begin{proof}
There are only finitely many connected components of every residual graph,
so \(d_G(\sigma)\) is a natural number.  The well ordering of the natural
numbers chooses a minimum among the nonempty family of supported matchings.

A finite graph is bipartite exactly when each connected component is
bipartite.  Hence \(d_G(\sigma)=0\) exactly when \(R_\sigma\) is bipartite.
Theorem~\ref{thm:residual-two-factor-obstruction} then turns defect zero into
a Tait colouring of \(G\).  Thus a minimizer in a Tait counterexample has
positive defect, and its minimality forbids a supported matching of smaller
defect.

Conversely, suppose every positive-defect supported matching admits a strict
defect decrease, and choose a minimizer \(\sigma_0\).  If its defect were
positive, the assumed decrease would contradict minimality.  Therefore its
defect is zero, its residual graph is bipartite, and
Theorem~\ref{thm:residual-two-factor-obstruction} gives a Tait colouring.
\end{proof}

Thus the route-specific high-width task may be sharpened once more: put a
minimum-defect matching through the wall and either construct one supported
matching of strictly lower defect, or extract the strict exact-support
replacement used by the reductive descent.  Producing an arbitrary new
matching, without controlling this potential, is not enough.

The potential has an exact quotient interpretation which keeps the
multigraph information that a simple quotient would lose.

\begin{theorem}[residual defect is even quotient oddness]
\label{thm:residual-defect-quotient-oddness}
\Sverified{Mettapedia.GraphTheory.FourColor.GoertzelV24ResidualOddness.colorable_two_iff_even_card_of_connected_degree_two}
\Sverified{Mettapedia.GraphTheory.FourColor.GoertzelV24ResidualOddness.connectedComponent_colorable_two_iff_even_supp}
\Sverified{Mettapedia.GraphTheory.FourColor.GoertzelV24ResidualOddness.residualDefect_eq_oddResidualQuotientVertices_ncard}
\Sverified{Mettapedia.GraphTheory.FourColor.GoertzelV24ResidualOddness.residualDefect_even}
\Sverified{Mettapedia.GraphTheory.FourColor.GoertzelV24ResidualOddness.two_le_residualDefect_of_pos}
\Sverified{Mettapedia.GraphTheory.FourColor.GoertzelV24ResidualOddness.exists_exchangeRigid_residualOddness_of_cubic_edgeBridgeFree}
Let \(G\) be a finite cubic simple graph and let \(\sigma\) be a supported
perfect matching.  Contract every connected component of \(R_\sigma\) to one
quotient vertex, but retain the original vertices as incident darts and pair
those darts by the involution \(\sigma\).  Thus loops and parallel quotient
edges remain distinct.  The degree of a quotient vertex is the number of
original vertices in its residual component.  Then
\[
 d_G(\sigma)
 =\#\{q:q\text{ is a quotient vertex of odd degree}\},
 \qquad 2\mid d_G(\sigma).
\]
In particular, \(d_G(\sigma)>0\) implies \(d_G(\sigma)\geq2\).

If \(G\) is bridge-free and not Tait colourable, one can moreover choose
\(\sigma\) so that \(d_G(\sigma)\geq2\) and every exchange with a supported
perfect matching on a set closed under both partner involutions has residual
defect at least \(d_G(\sigma)\).
\end{theorem}

\begin{proof}
Each residual component is a finite connected graph of degree two, hence a
cycle.  Choose a cyclic listing \(v_0,\ldots,v_{m-1}\).  If \(m\) is even,
colour \(v_i\) by the parity of \(i\); consecutive entries, including the
wraparound pair, receive opposite colours.  Conversely, if the cycle is
bipartite, its degree-two Hall condition supplies a perfect matching.  The
partner involution partitions its vertex set into pairs, so \(m\) is even.
Therefore a residual component is non-bipartite exactly when its support has
odd cardinality.

After contraction, the darts incident with the quotient vertex represented by
a residual component \(C\) are exactly the vertices of \(C\).  Hence its
degree, counted with multiplicity, is \(|C|\), and the definition of
\(d_G(\sigma)\) gives the first displayed equality.  The component supports
partition \(V(G)\).  Pairing every original vertex with its
\(\sigma\)-partner shows that \(|V(G)|\) is even; equivalently, the usual
handshaking parity argument shows that the number of odd quotient degrees is
even.  A positive even natural number is at least two.

Finally choose a supported matching of minimum residual defect by
Theorem~\ref{thm:residual-defect-descent}.  In a Tait counterexample its defect
is positive, hence at least two by the preceding paragraph.  An exchange on a
set closed under both partner involutions is again graph-supported.  Since the
chosen matching minimizes defect among all supported matchings, none of those
exchanges can decrease it.
\end{proof}

This theorem identifies the exact high-width obstruction: an even, nonempty
set of odd residual cycles joined by the matching-dart incidence, at a global
minimum under all supported exchanges.  Parity alone does not move this
minimum.  The missing wall argument must use planar placement and arbitrary
attachments to produce a strict decrease or an exact-support replacement; it
cannot be a purely local alternating-cycle identity.

The exact effect of a local matching exchange can be stated without any
monotonicity assumption.

\begin{lemma}[residual exchange is symmetric difference]
\label{lem:residual-exchange-symmetric-difference}
\Sverified{Mettapedia.GraphTheory.FourColor.GoertzelV24ResidualExchange.residualGraph_exchange_eq_symmDiff}
\Sverified{Mettapedia.GraphTheory.FourColor.GoertzelV24ResidualExchange.exchange_supportedBy_and_residualGraph_eq_symmDiff}
\Sverified{Mettapedia.GraphTheory.FourColor.GoertzelV24ResidualExchange.residualGraph_exchange_adj_of_mem}
\Sverified{Mettapedia.GraphTheory.FourColor.GoertzelV24ResidualExchange.residualGraph_exchange_adj_of_notMem}
Let \(\sigma,\tau\) be graph-supported perfect matchings of \(G\), and let
\(S\) be a vertex set closed under both partner involutions.  Write
\(\sigma'=\sigma\triangleleft_S\tau\) for the matching obtained by using
\(\tau\) on \(S\) and \(\sigma\) off \(S\).  Let \(A_S(\sigma,\tau)\)
be the graph consisting of the edges of
\(\sigma\mathbin\triangle\tau\) whose endpoints both lie in \(S\).  Then
\[
 R_{\sigma'}=R_\sigma\mathbin\triangle A_S(\sigma,\tau).
\]
Equivalently, at a vertex in \(S\) the new residual adjacency is exactly
the residual adjacency for \(\tau\), while at a vertex outside \(S\) it is
the old residual adjacency for \(\sigma\).
\end{lemma}

\begin{proof}
The supported-exchange theorem first shows that \(\sigma'\) uses only edges
of \(G\).  It remains to identify the residual edges.
For adjacent vertices \(v,w\), membership in \(R_\sigma\) means that
\(vw\) is an edge of \(G\) but is not the \(\sigma\)-edge at \(v\).
Outside \(S\), the exchanged partner is still \(\sigma(v)\), so no
residual adjacency changes.  Inside \(S\), it is \(\tau(v)\).  If the two
partners agree, neither the matching nor the residual graph changes.  If
they differ, the old \(\sigma\)-edge enters the residual graph and the old
\(\tau\)-edge leaves it.  These are precisely the two edges of the
alternating carrier incident with \(v\), so the change is symmetric
difference with \(A_S(\sigma,\tau)\).  Closure of \(S\) rules out an edge
with exactly one endpoint in \(S\), and graph support ensures that every
toggled matching edge belongs to \(G\).
\end{proof}

This identity is deliberately weaker than a defect-decrease statement.
Symmetric difference with an alternating carrier can split, merge, or reroute
residual components, so the number of non-bipartite components need not move
monotonically.  A wall argument must therefore control the carrier's
attachment to the odd residual components; the matching-exchange algebra
alone does not supply the strict descent.

There is also an exact limitation on the matching-exchange idea.

\begin{lemma}[two-matching exchange preserves the overlap locus]
\label{lem:two-matching-exchange-preserves-overlap}
\Sverified{Mettapedia.GraphTheory.FourColor.GoertzelV24TaitMatchingPair.twoMatchingKempeExchange_ne_iff}
Let \(S\) be a vertex set closed under two perfect matchings
\(\sigma,\tau\), and let \(\sigma',\tau'\) be obtained by swapping the two
matchings on \(S\).  Then, for every vertex \(v\),
\[
  \sigma'(v)\ne\tau'(v)\quad\Longleftrightarrow\quad
  \sigma(v)\ne\tau(v).
\]
In particular, such an exchange cannot turn an intersecting pair of perfect
matchings into an edge-disjoint pair.
\end{lemma}

\begin{proof}
Inside \(S\), the two partners are interchanged; outside \(S\), both partners
are unchanged.  Thus their unordered pair, and hence whether the two entries
are equal, is unchanged at every vertex.  The Lean theorem proves the displayed
pointwise equivalence directly.
\end{proof}

Theorem~\ref{thm:tait-three-matching-decomposition} does not by itself
exclude a wall.  Its purpose is to expose a route-native high-width attack
without losing the target.  Supported perfect matchings may be exchanged along
sets closed under two matchings, and the shared matching exchange theorem
proves that such a change remains graph-supported.  However,
Lemma~\ref{lem:two-matching-exchange-preserves-overlap} shows that exchanges
between a fixed pair preserve their common-edge locus exactly.  Thus they
cannot by themselves complete a Tait colouring.  A matching-based proof of the
boxed assertion below must use the wall either to construct a genuinely new
supported matching, involving edges outside the two old matchings, or to yield
the strict support-preserving replacement required by the reductive descent.
Merely finding many perfect matchings, or exchanging a fixed pair, does not
discharge the zero-target condition.
By Corollary~\ref{cor:tait-two-matching-criterion}, the completion alternative
can equivalently be stated as the production of two supported perfect
matchings with no common edge.  Thus the matching form of the high-width crux
is now exact: an embedded wall must force either such a disjoint pair or the
strict replacement used by the reductive descent.  Counting perfect matchings
without controlling their intersections is still insufficient.

One sufficient way to establish Hypothesis~1 is the stronger statement that
every zero-target instance containing the wall has either positive target or
a strict zero-target reduction.  The route has not proved that statement.
The hard step is therefore isolated as

\[
 \boxed{\text{a fixed embedded wall in a zero-target normal-form instance
 forces a strict zero-target reduction}.}
\]

The Goldberg audits show why curvature alone cannot establish the boxed
assertion: they have only pentagons and hexagons while their widths grow.
Their tested good fibres are seeded, which is evidence that the target
predicate may exclude the wall even though the geometry does not.  It is not
a proof.  Nor is the tree reformulation permission to count generic planar
configurations or revive a classical unavoidable catalogue.

The source's formation dynamics has a strong local entrance to this boxed
problem which is worth recording with its quantifiers in the correct order.

\begin{lemma}[one adjacent-pair Kempe orbit contains all constant words]
\label{lem:adjacent-pair-common-constant-orbit}
\Sverified{Mettapedia.GraphTheory.FourColor.GoertzelV24AdjacentPairCommonConstantOrbit.exists_rotationOrdered_commonKempeOrbit_with_all_constantWords_of_minimal}
Let \(uv\) be any ordered adjacent pair in a graph-backed vertex-minimal Tait
counterexample.  Delete \(u\) and \(v\), leaving the four incident defect
ports in their cyclic order.  There is one rotation ordering and one proper
Tait colouring \(C\) of this deleted graph such that, for each nonzero Tait
colour \(a\), the Kempe orbit of \(C\) contains a proper colouring \(C_a\)
whose four missing boundary colours are all \(a\).
\end{lemma}

\begin{proof}
The adjacent-pair planar reduction first supplies, on the literal cyclic
ordering, a proper colouring \(C\) and a two-colour profile in which the
first two ports belong to one Kempe side and the last two ports belong to
the other.  If a selected two-colour component joined the two sides, then
reinserting (u), (v), and their common edge would extend the colouring
to the ambient graph, contrary to the counterexample assumption.  Thus the
two side requests may be changed independently inside the Kempe closure.

Fix a nonzero target colour \(a\).  Recolour the first side to \(a\) while
leaving the second side fixed, and then recolour the second side to \(a\)
while leaving the first fixed.  Both operations are Kempe switches in the
deleted graph and preserve properness, so their composite gives \(C_a\).
The cyclic ordering and the initial colouring \(C\) were selected before
\(a\); hence the three target colourings belong to one common orbit rather
than to three unrelated existentially chosen colourings.
\end{proof}

The same local data also has a matching interpretation which survives even
when the full two-vertex colouring extension is impossible.

\begin{lemma}[an adjacent-pair deletion colouring supplies a central matching]
\label{lem:adjacent-pair-central-matching}
\Sverified{Mettapedia.GraphTheory.FourColor.GoertzelV24AdjacentPairMatchingExtraction.colorClassSubgraph_isPerfectMatching_of_boundary_absent}
\Sverified{Mettapedia.GraphTheory.FourColor.GoertzelV24AdjacentPairMatchingExtraction.AdjacentPairData.exists_central_supportedPairing_of_not_taitColorable}
Let \(G\) be a finite cubic simple graph which is not Tait colourable, let
\(u,v\) be adjacent, and suppose the four other neighbours of \(u,v\) are
distinct.  Every proper Tait colouring of \(G-\{u,v\}\) determines a
graph-supported perfect matching \(\tau\) of \(G\) with

\[
                              \tau(u)=v.
\]
\end{lemma}

\begin{proof}
Write \(w_0,w_1\) for the missing colours at the two ports formerly incident
with \(u\), and \(w_2,w_3\) for those at the ports formerly incident with
\(v\).  If the two deleted vertices could be restored as a proper coloured
bridge, the resulting colouring would be a Tait colouring of \(G\).  The
exact four-port bridge table therefore gives

\[
                         w_0=w_1,\qquad w_2=w_3.
\]

At most two of the three nonzero Tait colours occur in this word.  Choose a
nonzero colour \(c\) which does not occur.  Consider the \(c\)-coloured edges
of the deleted graph.  At every retained cubic vertex exactly one incident
edge has colour \(c\).  At a boundary vertex there are two incident edges of
distinct nonzero colours; their sum is the missing third colour \(w_i\).
Since \(c\ne w_i\), one of those two incident edges has colour \(c\), and
properness makes it unique.  Thus the \(c\)-colour class is a perfect
matching of every retained vertex.

Add the edge \(uv\).  It matches the only two omitted vertices, while all
other selected edges belong to the induced deletion and hence to \(G\).
The result is the required graph-supported perfect matching \(\tau\), and by
construction \(\tau(u)=v\).
\end{proof}

\begin{corollary}[a coloured adjacent-pair site supplies a rigid exchange input]
\label{cor:adjacent-pair-rigid-exchange-input}
\Sverified{Mettapedia.GraphTheory.FourColor.GoertzelV24AdjacentPairMatchingExtraction.AdjacentPairData.exists_exchangeRigid_with_central_sitePairing_of_cubic_edgeBridgeFree}
Under the hypotheses of Lemma~\ref{lem:adjacent-pair-central-matching}, suppose
in addition that \(G\) is bridge-free.  There are graph-supported perfect
matchings \(\sigma,\tau\) such that
\[
 d_G(\sigma)\geq2,
 \qquad \tau(u)=v,
\]
and, for every vertex set \(S\) closed under both partner involutions,
\[
 d_G(\sigma)
 \leq d_G\bigl(\sigma\mathbin\triangleleft_S\tau\bigr).
\]
\end{corollary}

\begin{proof}
Choose \(\sigma\) of minimum residual defect.  Since \(G\) is a bridge-free
cubic Tait counterexample, its residual defect is positive and even, hence at
least two.  Lemma~\ref{lem:adjacent-pair-central-matching} supplies the
graph-supported matching \(\tau\) with \(\tau(u)=v\).  Closure of \(S\) under
both partner involutions makes
\(\sigma\mathbin\triangleleft_S\tau\) another graph-supported perfect
matching.  The displayed inequality is therefore the defining minimality of
\(\sigma\).
\end{proof}

\begin{corollary}[one rigid residual matching against every ordered-mesh site]
\label{cor:ordered-mesh-rigid-site-matchings}
\Sverified{Mettapedia.GraphTheory.FourColor.GoertzelV24OrderedMeshResidualSiteMatching.exists_exchangeRigid_with_central_pairing_at_every_globalMeshStep}
Let \(G\) be a graph-backed least Tait counterexample containing an ordered
mesh.  There is one graph-supported perfect matching \(\sigma\), with
\(d_G(\sigma)\geq2\), such that for every physical row or column step
\(e=uv\) of the mesh there is a graph-supported perfect matching
\(\tau_e\) satisfying
\[
  \tau_e(u)=v
  \quad\hbox{and}\quad
  d_G(\sigma)
  \leq d_G\bigl(\sigma\mathbin\triangleleft_S\tau_e\bigr)
\]
for every vertex set \(S\) closed under both partner involutions.
The matching \(\sigma\) is independent of \(e\); the matching \(\tau_e\)
may depend on \(e\).
\end{corollary}

\begin{proof}
Choose \(\sigma\) once, globally, to minimize residual defect among all
graph-supported perfect matchings of \(G\).  Bridge-freeness and cubicity give
such a matching, and non-Tait-colourability makes its residual defect positive;
parity then gives \(d_G(\sigma)\geq2\).

Now fix a global mesh step \(e=uv\).  The ordered-mesh construction chooses
one adjacent-pair deletion datum and one proper Tait colouring of
\(G-\{u,v\}\) for this physical step.  Apply
Lemma~\ref{lem:adjacent-pair-central-matching} to obtain a supported perfect
matching \(\tau_e\) with \(\tau_e(u)=v\).  If \(S\) is closed under both
partner involutions, exchanging \(\sigma\) with \(\tau_e\) on \(S\) again
gives a graph-supported perfect matching.  The displayed inequality follows
from the global minimality of the same \(\sigma\).  Since \(e\) was arbitrary,
this supplies the family for every row and column step.  A shared side of two
mesh cells denotes the same global step, so it receives the same selected
site datum rather than a second local choice.
\end{proof}

\begin{corollary}[proper alternating component at every residual mesh step]
\label{cor:ordered-mesh-proper-alternating-components}
\Sverified{Mettapedia.GraphTheory.FourColor.GoertzelV24OrderedMeshResidualSiteMatching.exists_exchangeRigid_with_proper_alternatingComponent_at_every_globalMeshStep}
\Sverified{Mettapedia.GraphTheory.FourColor.GoertzelV24OrderedMeshResidualSiteFacialBond.exists_exchangeRigid_with_facialBond_at_every_globalMeshStep}
\Sverified{Mettapedia.GraphTheory.FourColor.GoertzelV24ResidualFormationSwitch.exists_exchangeRigid_with_facialFormationSwitch_at_every_globalMeshStep}
\Sverified{Mettapedia.GraphTheory.FourColor.GoertzelV24ResidualReturnShore.exists_exchangeRigid_with_facialFormation_returnPairing_and_shore_at_every_globalMeshStep}
\Sverified{Mettapedia.GraphTheory.FourColor.GoertzelV24ResidualReturnCycleOrder.cyclePosition_card_eq_cycle_length}
\Sverified{Mettapedia.GraphTheory.FourColor.GoertzelV24ResidualReturnCycleOrder.cycleVertexOrder_adj_finRotate}
\Sverified{Mettapedia.GraphTheory.FourColor.GoertzelV24ResidualReturnCycleOrder.orderedReturnShore_partner}
\Sverified{Mettapedia.GraphTheory.FourColor.GoertzelV24ResidualReturnCycleOrder.orderedSiteReturnPairing_walk_support_disjoint_of_chord_ne}
Under the hypotheses of
Corollary~\ref{cor:ordered-mesh-rigid-site-matchings}, choose its one global
minimum-residual-defect matching \(\sigma\).  For every physical mesh step
\(e=uv\), either \(\sigma(u)=v\), or there are a graph-supported perfect
matching \(\tau_e\) and a vertex set \(A_e\) such that
\[
 \tau_e(u)=v,\qquad u,v\in A_e,\qquad |A_e|\geq4,
\]
\(A_e\) is closed under both partner involutions.  More precisely, there is
a simple cycle \(C_e\) in \(G\), containing the edge \(e\), whose vertex set
is exactly \(A_e\) and whose edges alternate between \(\sigma\) and
\(\tau_e\).  Furthermore,
\[
 d_G(\sigma)
 \leq d_G\bigl(\sigma\mathbin\triangleleft_{A_e}\tau_e\bigr).
\]
Moreover the component is proper in the following witnessed sense: there is
a vertex \(x\notin A_e\) for which
\[
 \sigma(x)=\tau_e(x)
 \quad\hbox{and}\quad
 \sigma(x)\notin A_e.
\]
Finally, \(C_e\) is an exact bond of the facial dual.  There is a set
\(F_e\) of orbit faces such that a primal edge belongs to \(C_e\) exactly
when its two incident faces lie on opposite sides of \(F_e\), and the facial
dual subgraphs induced by \(F_e\) and by its complement are both connected.

There is also a literal local formation switch at every vertex \(z\) of
\(C_e\).  Cubicity supplies a unique third neighbour \(q_e(z)\), distinct
from both \(\sigma(z)\) and \(\tau_e(z)\).  In the residual graph obtained by
deleting the \(\sigma\)-matching, the two incident residual edges at \(z\)
are
\[
             z\tau_e(z)\quad\hbox{and}\quad zq_e(z).
\]
After exchanging \(\sigma\) with \(\tau_e\) on \(A_e\), they are
\[
             z\sigma(z)\quad\hbox{and}\quad zq_e(z).
\]
The fixed third edge \(zq_e(z)\) is not an edge of \(C_e\), and its two
incident orbit faces therefore lie in the same shore of the facial bond.

The unchanged off-cycle residual material also defines a fixed-point-free
involution
\[
                         \mu_e:A_e\longrightarrow A_e.
\]
For each \(z\in A_e\), the vertices \(z\) and \(\mu_e(z)\) lie in the same
connected component of the intersection of the residual graphs before and
after the exchange.  Equivalently, they are joined by a residual path all of
whose edges survive the formation switch.  Thus \(\mu_e\) pairs the literal
third-edge excursions when they return to the alternating cycle.  Each such
return remains on one facial shore: for every \(z\in A_e\), the orbit face on
the \(zq_e(z)\) side of the third edge belongs to \(F_e\) if and only if the
corresponding orbit face on the
\(\mu_e(z)q_e(\mu_e(z))\) side belongs to \(F_e\).  No noncrossing cyclic
order or formula for the number of resulting components is asserted by this
local statement alone.

There is, however, a canonical cyclic coordinate for this pairing.  Traverse
the simple cycle once and write
\[
 \gamma_e:\mathbb Z/|A_e|\mathbb Z\longrightarrow A_e
\]
for the resulting bijection.  Consecutive coordinates are endpoints of an
actual edge of \(C_e\).  Conjugating \(\mu_e\) by \(\gamma_e\) gives a
fixed-point-free involution \(\bar\mu_e\) on the cyclic positions.  Ordering
the two endpoints of each orbit after cutting the cycle at the chosen origin
gives a canonical finite chord, independent of which endpoint names it.
The indicator of membership in \(F_e\) at the chord's oriented third dart is
a well-defined Boolean label \(\chi_e\) on these chords.  Moreover, witnessing
paths for two different canonical chords have disjoint vertex supports (and
hence disjoint edge sets).
\end{corollary}

\begin{proof}
Fix a mesh step \(e=uv\), and let \(\tau_e\) be the site matching supplied by
Corollary~\ref{cor:ordered-mesh-rigid-site-matchings}.  If
\(\sigma(u)=v\), the first alternative holds.  Otherwise
\(\sigma(u)\ne\tau_e(u)\).  Let \(A_e\) be the equivalence class of \(u\)
under the relation generated by the two moves
\[
                  z\longmapsto\sigma(z),
             \qquad z\longmapsto\tau_e(z).
\]
By construction it contains \(u\), is closed under both moves, and contains
\(v=\tau_e(u)\).  Disagreement propagates under either involution: if the two
partners differ at \(z\), applying either matching and then using
involutivity shows that they also differ at the resulting vertex.  Hence the
two matchings differ at every vertex of \(A_e\).  The four vertices
\[
 u,\quad \sigma(u),\quad \tau_e(u),\quad
 \sigma\bigl(\tau_e(u)\bigr)
\]
are distinct, so \(|A_e|\geq4\).

Consider the symmetric difference of the two one-factors carried by
\(\sigma\) and \(\tau_e\).  At an agreement vertex it has degree zero.  At a
disagreement vertex \(z\), its two neighbours are the distinct vertices
\(\sigma(z)\) and \(\tau_e(z)\), so it has degree two.  Thus every nontrivial
connected component of the symmetric difference is a simple alternating
cycle.  Its component containing \(u\) has vertex set exactly \(A_e\): a
path in the symmetric difference is a sequence of the two generating moves,
and every generating move inside the disagreement class is an edge of the
symmetric difference.  Call this cycle \(C_e\).  Since
\(\tau_e(u)=v\), it contains the prescribed mesh edge \(e=uv\).

The two supported perfect matchings cannot be edge-disjoint.  Indeed, in a
cubic graph their two edges at each vertex would leave a unique third edge;
those unused edges form a third perfect matching and the three matchings give
a Tait colouring, contrary to the choice of \(G\).  Therefore there is a
vertex \(x\) with \(\sigma(x)=\tau_e(x)\).  Such an agreement vertex cannot
belong to the disagreement component \(A_e\), and involutivity gives agreement
at its common partner as well; consequently neither \(x\) nor \(\sigma(x)\)
belongs to \(A_e\).  Finally, closure of \(A_e\) makes the displayed exchange
a supported perfect matching, and its residual defect cannot be smaller than
that of the globally minimal \(\sigma\).

The least counterexample carries a bridgeless spherical cubic rotation system
with two-sided orbit faces.  Apply the simple-primal-cycle facial-bond theorem
to \(C_e\).  Its binary edge vector is an orbit-face boundary, and minimality
of the support of a simple cycle forces each of the two face shores to be
connected.  This supplies \(F_e\) with the asserted exact boundary equation.

At a vertex \(z\in A_e\), the two alternating-cycle neighbours are the
distinct partners \(\sigma(z)\) and \(\tau_e(z)\).  Since \(G\) is cubic,
there is exactly one remaining neighbour \(q_e(z)\).  Before the exchange the
residual graph deletes the \(\sigma\)-edge; after the exchange its selected
matching edge at \(z\) is the \(\tau_e\)-edge.  This gives the two displayed
pairs of residual edges.  The third edge is not in the alternating cycle.
The exact facial-bond equation consequently says that it does not cross the
face-shore boundary, which is precisely the final assertion.

It remains to construct the return involution.  Intersect the residual graph
before the exchange with the residual graph after it.  At a vertex of
\(A_e\), the formation calculation shows that only the third edge survives,
so its degree in the intersection is one.  Away from \(A_e\), the exchange
does not change the selected matching edge; the two residual graphs agree
there and the degree is two.  Consequently every connected component meeting
\(A_e\) is a finite path with exactly two degree-one endpoints in \(A_e\).
Pairing those two endpoints defines \(\mu_e\).  Uniqueness of the other
endpoint in each component makes \(\mu_e\) involutive, and the two endpoints
are distinct, so it is fixed-point-free.  The component itself supplies the
asserted unchanged residual path.  Every edge of this path avoids \(C_e\).
Label an oriented dart by the indicator of whether its incident orbit face
belongs to \(F_e\).  The exact facial-bond equation says that this label
changes across an edge exactly when that edge lies in \(C_e\).  The label is
therefore constant along the residual return path.  At either endpoint the
unique path edge is the canonical third edge, so the two canonical third
darts have the same \(F_e\)-shore label, as claimed.

For the cyclic coordinate, delete the repeated first vertex from the support
list of the closed simple walk \(C_e\).  The remaining list contains every
vertex of \(A_e\) exactly once, has length \(|C_e|\), and its cyclic successor
follows an edge of the displayed walk, including the final wraparound edge.
Transport \(\mu_e\) through this bijection.  Sorting the two endpoints by the
linear order after the cut makes a canonical representative because applying
the partner involution only swaps the two endpoints.  The preceding
same-shore equality makes \(\chi_e\) independent of that swap.

Finally, suppose two witness paths shared a vertex.  Their two initial cycle
vertices would then be reachable from one another in the common residual
graph.  A component of the degree--\(1/2\) graph has exactly two boundary
vertices, so uniqueness of the other boundary endpoint says that the two
initial vertices are partners (or equal).  In either case the two canonical
chords coincide.  Contraposition proves vertex-disjointness for distinct
chords.
\end{proof}

\begin{lemma}[finite residual circuit states are physical]
\label{lem:residual-circuit-state-is-physical}
\Sverified{Mettapedia.GraphTheory.FourColor.Compositional.ResidualCircuitComponentDecomposition.oldResidual_component_supp_eq_iUnion_commonComponent}
\Sverified{Mettapedia.GraphTheory.FourColor.Compositional.ResidualCircuitComponentDecomposition.newResidual_component_supp_eq_iUnion_commonComponent}
\Sverified{Mettapedia.GraphTheory.FourColor.Compositional.ResidualCircuitComponentDecomposition.oldResidual_component_cardParity_eq_oldCircuitParity}
\Sverified{Mettapedia.GraphTheory.FourColor.Compositional.ResidualCircuitComponentDecomposition.newResidual_component_cardParity_eq_newCircuitParity}
\Sverified{Mettapedia.GraphTheory.FourColor.Compositional.ResidualCircuitOddness.oldCircuitParity_eq_one_iff_mem_oddComponents}
\Sverified{Mettapedia.GraphTheory.FourColor.Compositional.ResidualCircuitOddness.newCircuitParity_eq_one_iff_mem_oddComponents}
\Sverified{Mettapedia.GraphTheory.FourColor.Compositional.ResidualCircuitOddness.oldCircuitParity_eq_one_iff_component_not_colorable}
\Sverified{Mettapedia.GraphTheory.FourColor.Compositional.ResidualCircuitOddness.newCircuitParity_eq_one_iff_component_not_colorable}
\Sverified{Mettapedia.GraphTheory.FourColor.Compositional.ResidualExchangeComponentLocalization.oldResidual_component_mem_oddComponents_iff_newResidual_of_avoids_carrier}
\Sverified{Mettapedia.GraphTheory.FourColor.Compositional.ResidualExchangeDefectLocalization.residualDefect_add_newCarrier_eq_new_add_oldCarrier}
\Sverified{Mettapedia.GraphTheory.FourColor.Compositional.ResidualCircuitDefectFormula.quot_mk_mem_oldOddCircuitClasses_iff_parity}
\Sverified{Mettapedia.GraphTheory.FourColor.Compositional.ResidualCircuitDefectFormula.quot_mk_mem_newOddCircuitClasses_iff_parity}
\Sverified{Mettapedia.GraphTheory.FourColor.Compositional.ResidualCircuitDefectFormula.oldOddCircuitClasses_ncard_eq_oldCarrierOddComponents}
\Sverified{Mettapedia.GraphTheory.FourColor.Compositional.ResidualCircuitDefectFormula.newOddCircuitClasses_ncard_eq_newCarrierOddComponents}
\Sverified{Mettapedia.GraphTheory.FourColor.Compositional.ResidualCircuitDefectFormula.residualDefect_add_newOddCircuitClasses_eq_new_add_old}
\Sverified{Mettapedia.GraphTheory.FourColor.Compositional.ResidualCircuitDefectFormula.oldOddCircuitClasses_ncard_le_new}
\Sverified{Mettapedia.GraphTheory.FourColor.Compositional.ResidualCircuitDefectFormula.residualDefect_eq_exchange_iff_oldOddCircuitClasses_ncard_eq_new}
\Sverified{Mettapedia.GraphTheory.FourColor.Compositional.ResidualCircuitDefectFormula.residualDefect_lt_exchange_iff_oldOddCircuitClasses_ncard_lt_new}
\Sverified{MatchingParity.CircuitPartition.circuitParity_const_one_eq_zero}
At one proper alternating site, combine the return involution with the old
local matching \(\tau_e\), or with the new local matching \(\sigma\), to form
a finite circuit partition on the cyclic positions of \(A_e\).  For every
root position, its oriented circuit contains exactly one endpoint of each
common residual return path used by the corresponding physical residual
component.  Consequently the component support is the disjoint union of
those return-path supports, and its cardinality parity is exactly the stored
finite circuit parity.

In either state, circuit parity one is equivalent both to membership of the
physical component in the odd-component set and to failure of
two-colourability.  Moreover every residual component disjoint from \(A_e\)
has the same support and oddness before and after the exchange.  If
\(O_{\rm old}^{A_e}\) and \(O_{\rm new}^{A_e}\) denote the odd residual
components meeting \(A_e\), then the global defect therefore satisfies the
exact subtraction-free identity
\[
 d_G(\sigma)+|O_{\rm new}^{A_e}|=
 d_G\bigl(\sigma\mathbin\triangleleft_{A_e}\tau_e\bigr)
   +|O_{\rm old}^{A_e}|.
\]
No inequality between the two local terms is asserted.

Equivalently, quotient the cyclic positions by the equivalence relation
generated by one return move or one local matching move.  The resulting old
and new finite circuit classes are in bijection with the physical residual
components meeting \(A_e\).  Under these bijections, a class is odd exactly
when the circuit-parity bit of any representative is one.  Hence the two
local cardinalities in the displayed identity are literal counts of odd
finite circuit classes.

The proper-site witness is exchange-rigid: the global residual defect cannot
fall under this exchange.  Substituting that inequality into the exact
identity shows that the old odd-circuit count is at most the new one.
Moreover, equality of the two global defects is equivalent to equality of the
two finite counts, and strict growth of the global defect is equivalent to
strict growth of the finite count.  Thus the direction comes from global
minimality, while the local circuit calculation supplies its exact finite
form; parity alone is not being used as a monotonicity principle.

The comparison does not, by itself, constrain the unlabelled return matching.
Indeed, for any two perfect matchings, assigning the constant odd label to
every return edge makes every alternating-circuit parity zero: each return
contributes one and its corresponding local edge contributes one.  Therefore
the high-width argument must use the physical return-length labels and their
compatibility across sites; chord shape plus the local inequality is
insufficient.
\end{lemma}

\begin{proof}
The equivalence relation generated by the common return involution and one
of the two local matching involutions is the physical reachability relation
inside the corresponding residual component.  Its two oriented orbits are
disjoint and are exchanged by either involution; one oriented orbit is thus
a transversal of the return paths.  Distinct positions in it have distinct
common residual components, while the canonical return path exhausts the
support of its common component.  Taking the disjoint union gives the support
formula, and taking cardinalities modulo two gives the circuit-parity
formula.

Every residual component is a finite connected two-regular graph.  Such a
component is two-colourable exactly when its support has even cardinality,
which proves the two physical interpretations of the parity bit.  Outside
\(A_e\), old and new residual adjacency agree.  An off-carrier component
therefore retains exactly its support, and the induced component graphs are
isomorphic.  Partition the old and new odd-component sets into components
disjoint from \(A_e\) and components meeting \(A_e\).  The first parts are
in support-preserving bijection; cancel their equal cardinalities to obtain
the displayed identity.  Finally apply the exchange-rigidity inequality and
elementary arithmetic to obtain the weak, equality, and strict comparison
statements for the finite circuit counts.
\end{proof}

\begin{lemma}[two-sector noncrossing of physical residual returns]
\label{lem:residual-return-two-sector-noncrossing}
\Sverified{Mettapedia.GraphTheory.FourColor.GoertzelV24ResidualReturnArc.orderedReturnPath_isPath}
\Sverified{Mettapedia.GraphTheory.FourColor.GoertzelV24ResidualReturnArc.orderedReturnPath_edge_not_cycle}
\Sverified{Mettapedia.GraphTheory.FourColor.GoertzelV24ResidualReturnArc.eq_start_or_eq_finish_of_mem_orderedReturnPath_support_of_mem_carrier}
\Sverified{Mettapedia.GraphTheory.FourColor.GoertzelV24ResidualReturnArc.orderedReturnPath_support_disjoint_of_chord_ne}
\Sverified{Mettapedia.GraphTheory.FourColor.GoertzelV24ResidualReturnSectorNoncrossing.orderedReturnSeparator_isCycle}
\Sverified{Mettapedia.GraphTheory.FourColor.GoertzelV24ResidualReturnSectorNoncrossing.exactFaceCut_orderedReturnSeparator_of_minimal}
\Sverified{Mettapedia.GraphTheory.FourColor.GoertzelV24ResidualReturnSectorNoncrossing.orderedChordAmbientPath_support_disjoint_of_ne}
\Sverified{Mettapedia.GraphTheory.FourColor.GoertzelV24ResidualReturnSectorNoncrossing.orderedChordAmbientPath_edge_not_orderedReturnSeparator_of_ne}
\Sverified{Mettapedia.GraphTheory.FourColor.GoertzelV24ResidualReturnSectorNoncrossing.not_crosses_of_orderedReturnShore_eq_of_minimal}
In the noncentral alternative above, let \(\chi_e\in\{0,1\}\) be the
face-shore label on the canonical chords of \(\bar\mu_e\).  For each
\(s\in\{0,1\}\), the chords labelled \(s\) are noncrossing in the cyclic
order of \(C_e\).  Equivalently, two distinct chords whose four endpoints
strictly alternate around \(C_e\) have different shore labels.  Thus the full
return pairing is the union of two noncrossing matchings; it need not itself
be noncrossing.
\end{lemma}

\begin{proof}
Let the first chord have ordered endpoints \(a<c\), with literal return path
\(P\).  Join \(P\) to the reverse of the displayed \(a\)-to-\(c\) interval of
\(C_e\).  The two pieces meet only at their endpoints, so their union is a
simple closed trail.  Sphericity supplies an exact binary face cut for this
trail: the labels on the two orbit faces incident with an edge differ exactly
when that edge belongs to the trail.

Let the second chord have return path \(Q\).  Distinct residual components
give vertex-disjoint return paths, and every edge of \(Q\) lies off
\(C_e\).  Thus no edge of \(Q\) belongs to the first chord's closed trail.
Transport of an exact face-cut label along such a trail says that the labels
on the canonical third darts at the two ends of \(Q\) are equal.

Suppose now that the four endpoints alternate.  In either of the two possible
linear orders, the cycle interval between the endpoints of \(Q\) contains
exactly one endpoint of \(P\).  At every other cycle vertex in this interval,
the residual third edge is not in the closed trail, so both oriented
cycle-side labels are unchanged through the local cubic turn.  At the unique
endpoint of \(P\), the residual third edge is the first or last edge of
\(P\), hence belongs to the trail, and the corresponding side label flips.
The assumption that the two chords have the same shore says that their four
endpoint third darts choose the same of the two local rotation sectors.
Consequently the appropriate binary side label flips exactly once between
the endpoints of \(Q\).  Those endpoint labels are therefore different,
contradicting the transport equality along \(Q\).  Hence equal-shore chords
cannot alternate, and each fixed-label family is noncrossing.
\end{proof}

\begin{corollary}[global two-sector residual receipt]
\label{cor:global-two-sector-residual-receipt}
\Sverified{Mettapedia.GraphTheory.FourColor.Compositional.ResidualReturnNoncrossing.exists_twoSectorReturnReceipt_at_every_globalMeshStep}
For an ordered mesh in a graph-backed least counterexample, one common
minimum-residual-defect matching can be chosen so that every mesh step is
either central for that matching or carries the complete physical
formation/return/shore receipt together with the two-sector noncrossing
conclusion of Lemma~\ref{lem:residual-return-two-sector-noncrossing}.
\end{corollary}

\begin{proof}
Choose the common residual-defect minimizer and the local physical receipts
from the preceding global construction.  At every noncentral step, apply
Lemma~\ref{lem:residual-return-two-sector-noncrossing} to the receipt's
facial-bond witness and canonical return pairing.  The choice of matching is
made before the mesh step, so the quantifier order remains
\(\exists\sigma\,\forall e\), rather than making an unrelated choice of
\(\sigma\) at each site.
\end{proof}

\begin{corollary}[colouring provenance of the global residual pairing]
\label{cor:global-residual-pairing-provenance}
\Sverified{Mettapedia.GraphTheory.FourColor.Compositional.ResidualSiteProvenance.nonempty_globalSitePairingProvenance}
\Sverified{Mettapedia.GraphTheory.FourColor.Compositional.ResidualSiteProvenance.exists_exchangeRigid_with_provenance_at_every_globalMeshStep}
\Sverified{Mettapedia.GraphTheory.FourColor.Compositional.ResidualSiteProvenance.exists_exchangeRigid_with_provenanced_alternatingSite_at_every_globalMeshStep}
The common minimum-residual-defect matching in
Corollary~\ref{cor:global-two-sector-residual-receipt} can be chosen together
with the following data at every mesh step.  Retain the exact globally
selected Tait colouring of the graph with the step's central edge deleted,
choose a colour absent from its four boundary darts, and complete that colour
class by the deleted central edge.  The resulting ambient perfect matching is
supported on graph edges, contains the central edge, and is the \emph{same}
second matching used by the proper alternating-cycle witness in the
noncentral alternative.
\end{corollary}

\begin{proof}
At a step, minimality colours the graph with the central edge deleted.  The
two pairs of boundary colours on the two sides of that deleted edge agree;
the three-colour boundary algebra therefore supplies a colour absent from all
four boundary darts.  Its colour class matches every retained vertex except
the two endpoints of the deleted edge.  Adding that edge produces a supported
perfect matching through the step.

Choose the ambient matching with minimum residual defect before choosing a
mesh step.  Its exchange-rigidity applies to each of the supported matchings
just constructed.  If the ambient matching already contains the central edge
we are in the central alternative.  Otherwise the alternating component of
the two matchings through that edge supplies the proper residual cycle.  By
construction its second matching is definitionally the completed absent
colour class just recorded, so the original deletion colouring and its
absent colour survive into the residual witness rather than being hidden
behind a fresh existential choice.
\end{proof}

\begin{corollary}[finite odd-circuit monotonicity across the mesh]
\label{cor:global-residual-circuit-monotonicity}
\Sverified{Mettapedia.GraphTheory.FourColor.Compositional.ResidualMeshCircuitMonotonicity.exists_pairing_with_oddCircuitMonotonicity_at_every_globalMeshStep}
For an ordered mesh in a graph-backed least counterexample, one supported
residual-defect minimizer can be chosen before the mesh step so that every
step is either central for that matching or carries a physical proper-site
receipt whose old odd-circuit count is at most its new odd-circuit count.
The circuit counts are taken in the finite quotient of the carrier positions
from Lemma~\ref{lem:residual-circuit-state-is-physical}.
\end{corollary}

\begin{proof}
Use the common minimizer and provenanced alternating-site receipts from
Corollary~\ref{cor:global-residual-pairing-provenance}.  At a noncentral step,
Lemma~\ref{lem:residual-circuit-state-is-physical} identifies the finite odd
circuit classes with the physical odd components meeting the alternating
carrier.  Its exact balance law and the receipt's exchange-rigidity give the
claimed inequality.  No compatibility between different sites and no strict
inequality are asserted.
\end{proof}

\begin{corollary}[proper residual geometry at every internal mesh junction]
\label{cor:internal-mesh-junction-proper-residual-site}
\Sverified{Mettapedia.GraphTheory.FourColor.Compositional.MeshJunctionAlternatingGeometry.incoming_first_ne_outgoing_second}
\Sverified{Mettapedia.GraphTheory.FourColor.Compositional.MeshJunctionAlternatingGeometry.not_both_arms_central}
\Sverified{Mettapedia.GraphTheory.FourColor.Compositional.MeshJunctionAlternatingGeometry.exists_minimizer_with_proper_site_at_every_internal_row_junction}
\Sverified{Mettapedia.GraphTheory.FourColor.Compositional.MeshJunctionAlternatingGeometry.exists_minimizer_with_twoSector_site_at_every_internal_row_junction}
Let an ordered injective mesh have at least two boundary columns in addition
to its interior columns.  One common minimum-residual-defect matching
\(\sigma\) can be chosen so that, at every interior branch vertex of every
row, at least one of the incoming and outgoing row steps carries a proper
alternating-component receipt.  Its second matching is the absent-colour
completion of the globally selected deletion colouring at that physical
step.  The same arm can be chosen with the complete facial-bond, physical
formation, canonical return, face-shore, and two-sector noncrossing receipt.
\end{corollary}

\begin{proof}
Fix an interior branch vertex on a row.  The row is vertex-injective, so the
outer endpoints of the immediately incoming and outgoing steps are distinct.
Both steps meet at the branch vertex.  If \(\sigma\) contained both steps,
the uniqueness of the \(\sigma\)-partner of that branch vertex would identify
the two distinct outer endpoints, a contradiction.  Hence at least one step
is noncentral for \(\sigma\).  Apply
Corollary~\ref{cor:global-residual-pairing-provenance} to that step: its
globally selected deletion colouring supplies the supported completed
matching, and the noncentral alternative supplies the provenanced proper
alternating component.
The provenance-preserving local geometry construction then carries this same
arm, without another choice of matching or deletion colouring, to the full
two-sector return receipt.
\end{proof}

\begin{corollary}[colouring provenance survives the residual geometry]
\label{cor:global-residual-geometry-provenance}
\Sverified{Mettapedia.GraphTheory.FourColor.Compositional.ResidualSiteGeometry.exists_provenancedTwoSectorReturnReceipt_at_every_globalMeshStep}
\Sverified{Mettapedia.GraphTheory.FourColor.Compositional.ResidualSiteGeometry.provenancedTwoSectorReturn_site_tau_eq_pairing}
\Sverified{Mettapedia.GraphTheory.FourColor.Compositional.ResidualSiteGeometry.provenancedTwoSectorReturn_site_eq_selected}
The choice in Corollary~\ref{cor:global-residual-pairing-provenance} can be
retained through the complete local geometry.  Thus one common
minimum-residual-defect matching can be chosen so that every mesh step is
central or carries a facial bond, physical formation, canonical return
pairing, and two-sector noncrossing receipt whose second matching is exactly
the absent-colour completion of the globally selected deletion colouring at
that step.
\end{corollary}

\begin{proof}
At a noncentral step, use the proper alternating cycle recorded by
Corollary~\ref{cor:global-residual-pairing-provenance}.  Its simple primal
cycle has two connected facial shores.  The formation switch and canonical
return construction depend only on this supplied site and its two supported
matchings, so they introduce no new choice of colouring or matching.  Apply
the two-sector theorem to the resulting return-shore receipt.  The dependent
equalities carried by the constructors identify both the final site and its
second matching with the originally selected site and completed absent colour
class.
\end{proof}

\begin{lemma}[colour classes and completed deletion matchings]
\label{lem:deletion-colour-matching-overlap}
\Sverified{Mettapedia.GraphTheory.FourColor.Compositional.DeletionColorMatching.colorClassPairing_mem_edge_iff}
\Sverified{Mettapedia.GraphTheory.FourColor.Compositional.DeletionColorMatching.deletedColorClassPairing_mem_edge_iff}
\Sverified{Mettapedia.GraphTheory.FourColor.Compositional.DeletionColorMatching.centralCompletionPairing_mem_retainedEdge_iff}
\Sverified{Mettapedia.GraphTheory.FourColor.Compositional.DeletionColorMatching.centralCompletionPairings_agree_on_common_retained_edge_of_bit}
A Tait colouring of an adjacent-pair deletion, together with a colour absent
from all four boundary requests, determines an ambient perfect matching by
completing that colour class with the deleted central edge.  On every
retained ambient edge, membership in the completed matching is equivalent to
having the selected colour.  Consequently two such completed matchings agree
on every edge retained by both deletions whenever they select the same absent
colour and the Boolean common-restriction coordinate of their deletion
colourings is true.
\end{lemma}

\begin{proof}
The absent-colour boundary condition makes the selected colour class a
perfect matching even at the four degree-two boundary vertices.  Completing
it changes only the two deleted vertices: they are paired by their central
edge.  Thus a retained edge belongs to the completed matching exactly when
its deletion colouring has the selected colour.  If two deletion colourings
agree on their common retained carrier and select the same colour, these two
edgewise characterizations are identical.
\end{proof}

\begin{corollary}[finite matching-overlap state and its nine-site obstruction]
\label{cor:nine-site-matching-overlap-obstruction}
\Sverified{Mettapedia.GraphTheory.FourColor.Compositional.ResidualSiteMatchingOverlap.card_matchingOverlapState}
\Sverified{Mettapedia.GraphTheory.FourColor.Compositional.ResidualSiteMatchingOverlap.deletionMatchingStates_agree_on_common_retained_edge}
\Sverified{Mettapedia.GraphTheory.FourColor.Compositional.ResidualSiteMatchingOverlap.pairings_agree_on_common_retained_edge}
\Sverified{Mettapedia.GraphTheory.FourColor.Compositional.ResidualSiteMatchingOverlap.exists_incompatible_pair_in_any_nine_row_intervals}
For a pair of provenanced residual sites, record the absent-colour index at
each site and the Boolean common-restriction coordinate of their selected
deletion colourings.  There are exactly

\[
                    3\cdot 3\cdot 2=18
\]

such states.  A state is \emph{compatible} when the two absent-colour indices
are equal and the common-restriction bit is true; compatibility forces the
two completed site matchings to agree on every common retained edge.
Nevertheless, among any nine distinct intervals of one ordered mesh row,
every provenanced choice contains two sites whose matching-overlap state is
incompatible.
\end{corollary}

\begin{proof}
The state count is the displayed product.  Compatibility first identifies
the selected colour and then invokes
Lemma~\ref{lem:deletion-colour-matching-overlap} on the common retained
carrier.  For the final assertion, forget the absent-colour proofs and retain
the exact globally selected deletion colourings carried by the provenance
receipts.  The nine-site colouring-atlas obstruction supplies two whose
common-restriction bit is false.  Their matching-overlap state is therefore
incompatible regardless of which absent colours its two receipts record.
\end{proof}

\begin{lemma}[compatible overlap transports alternating geometry]
\label{lem:compatible-overlap-alternating-geometry}
\Sverified{Mettapedia.GraphTheory.FourColor.Compositional.AlternatingOverlapGeometry.alternatingGraph_adj_congr_right}
\Sverified{Mettapedia.GraphTheory.FourColor.Compositional.AlternatingOverlapGeometry.alternatingGraphs_agree_on_common_retained_edge}
\Sverified{Mettapedia.GraphTheory.FourColor.Compositional.AlternatingOverlapGeometry.provenanced_alternatingGraphs_agree_on_common_retained_edge}
\Sverified{Mettapedia.GraphTheory.FourColor.Compositional.AlternatingOverlapGeometry.source_cycle_edge_mem_target_alternatingGraph}
Fix the global residual-defect matching $\sigma$.  Let two provenanced
residual sites have a compatible matching-overlap state.  On every ambient
edge retained by both adjacent-pair deletions, the two symmetric-difference
graphs
\[
        \sigma\mathbin{\triangle}\tau_A
        \quad\hbox{and}\quad
        \sigma\mathbin{\triangle}\tau_B
\]
have identical edge membership.  In particular, every such edge of the
distinguished alternating cycle at site $A$ belongs to the alternating graph
at site $B$.
\end{lemma}

\begin{proof}
For a fixed edge, membership in a symmetric difference is the exclusive-or
of membership in its two matchings.  The first matching $\sigma$ is literally
the same at both sites.  Compatibility of the overlap state, together with
Lemma~\ref{lem:deletion-colour-matching-overlap}, says that the second
matchings $\tau_A$ and $\tau_B$ make the same decision on every edge retained
by both deletions.  The two exclusive-or values are therefore equal.

The provenance receipts identify $\tau_A$ and $\tau_B$ with the second
matchings used by the final residual-cycle constructions.  Every edge of the
distinguished cycle at $A$ already lies in
$\sigma\mathbin{\triangle}\tau_A$; transporting its membership across the
preceding equivalence places it in
$\sigma\mathbin{\triangle}\tau_B$.  This does not assert that it lies in
site $B$'s distinguished component: the common carrier may contain several
alternating components.
\end{proof}

\begin{lemma}[alternating-component changes are confined to two deletion footprints]
\label{lem:alternating-component-deletion-footprint-localization}
\Sverified{Mettapedia.GraphTheory.FourColor.Compositional.AlternatingComponentLocalization.deletionFootprint}
\Sverified{Mettapedia.GraphTheory.FourColor.Compositional.AlternatingComponentLocalization.mem_deletionFootprint_iff}
\Sverified{Mettapedia.GraphTheory.FourColor.Compositional.AlternatingComponentLocalization.card_deletionFootprint}
\Sverified{Mettapedia.GraphTheory.FourColor.Compositional.AlternatingComponentLocalization.card_overlapFootprintValues_le}
\Sverified{Mettapedia.GraphTheory.FourColor.Compositional.AlternatingComponentLocalization.provenanced_alternatingGraphs_agree_outside_deletionFootprints}
\Sverified{Mettapedia.GraphTheory.FourColor.Compositional.AlternatingComponentLocalization.provenanced_alternatingGraphs_delete_overlapFootprints_eq}
\Sverified{Mettapedia.GraphTheory.FourColor.Compositional.AlternatingComponentLocalization.card_provenanced_alternatingGraph_edgeDisagreement_le}
\Sverified{Mettapedia.GraphTheory.FourColor.Compositional.AlternatingComponentLocalization.card_provenanced_alternatingGraph_edgeDisagreement_le_twenty}
\Sverified{Mettapedia.GraphTheory.FourColor.Compositional.AlternatingComponentLocalization.BundledResidualReturnReceipt.edgeDisagreementPathCost_le}
\Sverified{Mettapedia.GraphTheory.FourColor.Compositional.AlternatingComponentLocalization.BundledResidualReturnReceipt.card_endpoint_edgeDisagreement_le}
\Sverified{Mettapedia.GraphTheory.FourColor.Compositional.AlternatingComponentLocalization.BundledResidualReturnReceipt.exists_endpoint_component_exceptionalEdges}
\Sverified{Mettapedia.GraphTheory.FourColor.Compositional.AlternatingComponentLocalization.source_cycle_edge_mem_target_or_deletionFootprints}
\Sverified{Mettapedia.GraphTheory.FourColor.Compositional.AlternatingComponentLocalization.source_cycle_common_arc_component_eq_in_target}
For an adjacent-pair deletion let its \emph{deletion footprint} be the
central edge together with the four boundary edges.  It has exactly five
edges.  Hence the union of the footprints of two residual sites has at most
ten edges.

If the two sites have compatible matching-overlap state, delete those at
most ten edge values from both matching symmetric-difference graphs.  The
two remaining graphs are equal.  Consequently their actual edge-set
disagreement has cardinality at most ten.  Equivalently, every edge of the
distinguished alternating cycle at the first site either belongs to the
alternating graph at the second site or lies in one of the two deletion
footprints.  Every arc of the first cycle avoiding both footprints therefore
has its endpoints in one connected component of the second alternating
graph.  More generally, along a chain of $r$ consecutively compatible
receipts, the first and last alternating graphs disagree on at most $10r$
edge values.  Taking those disagreements as an exceptional carrier, every
walk in the first graph that avoids the carrier has its endpoints in one
connected component of the last graph.  Thus component splitting and merging
across the whole chain is confined to at most $10r$ edge locations.
\end{lemma}

\begin{proof}
An ambient edge fails to survive an adjacent-pair deletion precisely when it
is the central edge or one of the four boundary edges.  The four boundary
edges are distinct, and none is central, so the footprint has cardinality
five.  The union bound gives ten for two sites.

Away from this union, the edge is retained by both deletions.  Lemma
\ref{lem:compatible-overlap-alternating-geometry} then identifies membership
in the two symmetric-difference graphs.  Extensionality after deleting the
footprint gives equality of the remaining graphs.  The generic finite-graph
localization lemma says that if deleting a finite edge set makes two graphs
equal, their edge-set disagreement is contained in that set; the cardinal
bound is therefore ten.  Edge disagreement is a symmetric difference, so its
cardinality satisfies the triangle inequality.  Iterating this inequality
along a dependently bundled chain of residual receipts gives the (10r)
endpoint bound.  A walk whose edges avoid the endpoint symmetric difference
transfers verbatim to the last graph, proving the component-localization
form.  Applying the same
edgewise statement along a walk contained in the first alternating cycle
transfers that walk to the second alternating graph; its endpoints are thus
reachable in one connected component.  The statement deliberately allows
components to split or merge at footprint edges.
\end{proof}

\begin{lemma}[rooted gauge for residual face sides]
\label{lem:rooted-residual-face-side-gauge}
\Sverified{Mettapedia.GraphTheory.Embedding.FaceCutGauge.basedLabel_eq}
\Sverified{Mettapedia.GraphTheory.FourColor.Compositional.ResidualReturnSide.facialBondCut_basedLabel_eq}
\Sverified{Mettapedia.GraphTheory.FourColor.Compositional.ResidualReturnSide.basedReturnSide_partner}
\Sverified{Mettapedia.GraphTheory.FourColor.Compositional.ResidualReturnSide.not_crosses_of_basedReturnSide_eq}
Fix one ambient orbit face \(f_0\).  If two exact binary face cuts have the
same selected primal edges, then their normalized labels

\[
       \widehat\lambda(f)=\lambda(f)-\lambda(f_0)
\]

agree on every orbit face.  In particular, the facial bond at a noncentral
residual site has a canonical \(f_0\)-rooted side coordinate, independent of
which of its two shores was named first.  The two endpoints of every physical
return have the same rooted coordinate, and two return chords with the same
rooted coordinate do not cross.
\end{lemma}

\begin{proof}
The facial-boundary vector of a binary face labeling records, on each primal
edge, whether the two incident labels differ.  Exact cuts with the same
selected edges therefore have equal boundary vectors.  Their difference lies
in the kernel of the facial-boundary map.  The full facial dual is connected,
so the kernel theorem says that this difference is constant on all orbit
faces.  Subtracting its value at \(f_0\) removes that constant.

Apply this to the indicator of either connected face shore supplied by the
residual facial bond.  Complementing the chosen shore merely adds the constant
(1), hence does not change the rooted label.  The previously proved
same-shore statement for the endpoints of a return gives equality of their
rooted coordinates.  Finally, equality of rooted coordinates is equivalent
to equality of the local Boolean shore labels, so
Lemma~\ref{lem:residual-return-two-sector-noncrossing} excludes crossing.
\end{proof}

This is not merely another proof that a bridgeless cubic graph has a perfect
matching.  Every physical mesh step now supplies an ambient matching through
that step, and all those site matchings are tested against one common
minimum-residual-defect matching.  What is not yet proved is the wall-facing
statement: that the matchings obtained at neighbouring sites can be
synchronized so that symmetric difference strictly lowers residual quotient
oddness, or instead yields equal exact shore supports for a physical
replacement.  In particular, the corollary has quantifier order
\(\exists\sigma\,\forall e\,\exists\tau_e\); it does not assert a coherent
simultaneous choice of the \(\tau_e\).  Corollary~\ref{cor:global-residual-pairing-provenance}
now identifies each \(\tau_e\) with the absent colour class of the exact
globally selected deletion colouring at \(e\), but it does not yet transport
that colouring between two distinct deletions.  Corollary~\ref{cor:nine-site-matching-overlap-obstruction}
shows that this is a real coherence obstruction rather than missing
bookkeeping: the evident eighteen-state pairwise interface is incompatible
at some pair in every nine-site row sample.  Equivalently, repetition of a
pairwise state cannot by itself manufacture a simultaneous section.  The
local formation receipts identify
which residual edge changes at each cycle vertex, and the common-residual
components now pair the third-edge excursions when they return to the cycle.
They also prove that each individual return remains on one facial shore.
Lemma~\ref{lem:residual-return-two-sector-noncrossing} puts the returns into
two shore-labelled noncrossing families.  What remains is to synchronize
those two-sector pairings at neighbouring sites and compute a strict change
in residual quotient oddness (or an equal exact support state yielding a
physical replacement).

This lemma is nonlocal in precisely the way the original formation argument
needs: the Kempe paths witnessing the two independent side moves may run
through the whole retained graph.  It is therefore not a bounded-radius
reducibility certificate.  Its remaining wall-facing gap is also exact.
For overlapping adjacent pairs the lemma supplies a common orbit separately
at each deletion, but it does not transport a colouring, a Kempe word, or an
exact Count state from one deleted graph to the next.  A source-faithful
Programme--G proof of fixed-wall exclusion must construct that overlap
transport (or a simultaneous multi-deletion orbit) and then extract a strict
physical replacement.  Merely pigeonholing the three constant words would
discard the graph-dependent Kempe paths and would not prove the boxed
assertion.

The first overlap object can now be stated without choosing a privileged
representative of either Kempe orbit.

\begin{lemma}[canonical finite state of two adjacent-pair orbits]
\label{lem:adjacent-pair-overlap-kempe-state}
\Sverified{Mettapedia.GraphTheory.FourColor.GoertzelV24AdjacentPairOverlap.isTaitEdgeColoring_firstDeletionCommonCoreColoring}
\Sverified{Mettapedia.GraphTheory.FourColor.GoertzelV24AdjacentPairOverlap.isTaitEdgeColoring_secondDeletionCommonCoreColoring}
\Sverified{Mettapedia.GraphTheory.FourColor.GoertzelV24AdjacentPairOverlapKempeState.overlapKempeStateSupport_eq_of_mem}
\Sverified{Mettapedia.GraphTheory.FourColor.GoertzelV24AdjacentPairOverlapKempeState.overlapKempeStateSupport_nonempty}
Let \(p=(u,v)\) and \(q=(x,y)\) be two ordered adjacent-pair data in the
same finite graph.  Their exact common carrier is the induced graph obtained
by deleting \(u,v,x,y\).  It embeds canonically in each of the two
adjacent-pair deletions, and restriction along either embedding sends a
proper nonzero Tait colouring to a proper nonzero Tait colouring of the
common carrier.

For proper base colourings \((C_p,C_q)\) of the two deletions, attach to a pair
of proper colourings in their respective Kempe closures the following finite
state:

\begin{enumerate}
\item the \(4\)-by-\(4\) Boolean incidence table recording which labelled
  port vertices of \(p\) and \(q\) coincide;
\item the exact four-port colour word at each deletion;
\item for every ordered pair of colours and every ordered pair of ports, the
  Boolean value saying whether those ports meet the same selected
  two-colour component; and
\item the Boolean value saying whether the two restrictions to the common
  carrier agree literally.
\end{enumerate}

Let \(\mathcal R(C_p,C_q)\) be the set of all such states realised by proper
colourings in the two closures.  Then \(\mathcal R(C_p,C_q)\) is nonempty.
Moreover, if \(C'_p\) and \(C'_q\) are reachable from \(C_p\) and \(C_q\),
respectively, then

\[
        \mathcal R(C_p,C_q)=\mathcal R(C'_p,C'_q).
\]
Thus this overlap relation belongs to the two Kempe orbits, not to the
chosen representatives.
\end{lemma}

\begin{proof}
Deleting all four named vertices factors through either two-vertex deletion
by the inclusion of retained vertices.  Pulling an edge colouring back along
either induced-graph embedding merely evaluates it on the corresponding
ambient edge, so nonzeroness is preserved edge by edge.

Every coordinate listed above has finite domain and finite codomain; their
product is therefore finite.  The two base colourings themselves witness
nonemptiness.  Finally, a member of a Kempe closure has the same Kempe
closure as the base: reachability is reflexive, symmetric, and transitive.
Replacing either base by a reachable representative consequently changes
neither quantified colouring set and hence changes no realised overlap
state.
\end{proof}

Applied to Lemma~\ref{lem:adjacent-pair-common-constant-orbit}, the two
marginals of this relation each contain the three nonzero constant boundary
words.  The lemma deliberately does \emph{not} assert that a state with
literal common-carrier agreement is realised, nor that two neighbouring
relations have equal Count support.  Exact agreement is false in some of the
checked intrinsic-short reentry geometries.  The remaining wall-facing task
is therefore a genuinely two-dimensional statement about composing these
finite orbit relations---either a coherent flat region yields the support
equality needed by the physical splice, or a non-coboundary holonomy yields
a structural target.  A chosen-word transport, and any additive discrepancy
of endpoint potentials, is too weak for that task.

The first exact composition object and its connection to the wall carrier
are now available.  They remove a quantifier ambiguity which is easy to miss:
pairwise relations may be witnessed by different colourings at their common
corner, whereas a compositional cell must reuse one actual colouring there.

\begin{lemma}[coherent orbit cell on an ordered mesh]
\label{lem:ordered-mesh-coherent-orbit-cell}
\Sverified{Mettapedia.GraphTheory.FourColor.GoertzelV24AdjacentPairOverlapKempeCell.overlapKempeCellStateSupport_eq_of_mem}
\Sverified{Mettapedia.GraphTheory.FourColor.GoertzelV24AdjacentPairOverlapKempeCell.overlapKempeCellStateSupport_nonempty}
\Sverified{Mettapedia.GraphTheory.FourColor.GoertzelV24AdjacentPairOverlapKempeCell.overlapKempeCellStateSupport_subset_compatibleCycleSupport}
\Sverified{Mettapedia.GraphTheory.FourColor.GoertzelV24AdjacentPairOverlapKempeCell.sharedWitnessLifting_or_obstruction}
\Sverified{Mettapedia.GraphTheory.FourColor.GoertzelV24AdjacentPairOverlapKempeCell.overlapKempeCellStateSupport_eq_compatibleCycleSupport_iff}
\Sverified{Mettapedia.GraphTheory.FourColor.GoertzelV24AdjacentPairOverlapKempeLifting.hasBoundaryKempeStateCollision_iff_not_separatesOrbit}
\Sverified{Mettapedia.GraphTheory.FourColor.GoertzelV24AdjacentPairOverlapKempeLifting.hasSharedWitnessLifting_of_boundaryKempeStateSeparatesOrbits}
\Sverified{Mettapedia.GraphTheory.FourColor.GoertzelV24AdjacentPairOverlapKempeLifting.exists_corner_boundaryKempeStateCollision_of_not_sharedWitnessLifting}
\Sverified{Mettapedia.GraphTheory.FourColor.GoertzelV24AdjacentPairOverlapKempeHolonomy.mem_overlapKempeCellStateSupport_iff_hasRelationalSection}
\Sverified{Mettapedia.GraphTheory.FourColor.GoertzelV24AdjacentPairOverlapKempeHolonomy.mem_pairwise_not_mem_realized_iff_fixedPointFreeHolonomy}
\Sverified{Mettapedia.GraphTheory.FourColor.GoertzelV24AdjacentPairOverlapKempeHolonomy.hasSharedWitnessObstruction_iff_exists_fixedPointFreeHolonomy}
\Sverified{Mettapedia.GraphTheory.FourColor.GoertzelV24AdjacentPairOverlapKempeHolonomyResidue.hasFixedPointFreeRelationalHolonomy_iff_break_or_nonempty}
\Sverified{Mettapedia.GraphTheory.FourColor.GoertzelV24AdjacentPairOverlapKempeHolonomyResidue.hasSharedWitnessObstruction_iff_exists_break_or_nonemptyHolonomy}
\Sverified{Mettapedia.GraphTheory.FourColor.GoertzelV24AdjacentPairOverlapKempeHolonomyResidue.nonemptyFixedPointFreeRelationalHolonomy_residue}
\Sverified{Mettapedia.GraphTheory.FourColor.GoertzelV24AdjacentPairOverlapKempeHolonomyResidue.equivalence_sign_eq_neg_one_of_card_eq_two}
\Sverified{Mettapedia.GraphTheory.FourColor.GoertzelV24OrderedMeshAdjacentPairSites.meshPathStep_adjacent}
\Sverified{Mettapedia.GraphTheory.FourColor.GoertzelV24OrderedMeshAdjacentPairSites.exists_nonempty_overlapKempeCellStateSupport}
Let four adjacent-pair deletions be arranged clockwise around a cell, with
proper base colourings at the north-west, north-east, south-east and
south-west corners.  Evaluate the overlap state on each side using the same
concrete corner colouring for the two sides incident with that corner.  The
set of all cell states obtained by independently ranging the four corner
colourings through their Kempe orbits is finite and nonempty, is independent
of the four chosen orbit representatives, and projects into the cycle of the
four pairwise overlap supports.  At every corner its two projected boundary
observations agree literally.

Moreover, every consecutive coordinate rectangle of an ordered injective mesh in a
graph-backed vertex-minimal Tait counterexample canonically supplies such
four deletion sites.  They are anchored at the four mesh branch vertices,
and the common-constant-orbit lemma supplies proper base colourings, hence a
nonempty coherent cell support.
\end{lemma}

\begin{proof}
For the abstract statement, choose one proper colouring at each corner and
use it in both incident pairwise overlap evaluations.  The two boundary
observations at that corner are then definitionally the same.  Taking all
four colourings in their respective Kempe closures gives a finite relation;
the base colourings witness nonemptiness.  Replacing a base colouring by a
reachable representative leaves its whole Kempe closure unchanged, so it
leaves the cell relation unchanged.  Forgetting the shared witnesses gives
the asserted inclusion in the naive pairwise-compatible cycle.

For the mesh statement, every step of a row or column path is an edge of the
rotation-system multigraph and therefore an edge of its graph backing.  At
each cell corner select the first outgoing path step in clockwise order
(equivalently the reverse of the last incoming step on the south and west
sides).  Strict monotonicity of the branch positions guarantees that these
steps exist.  Cubicity and the absence of a common neighbour at an adjacent
pair reconstruct the canonical four-port deletion data, and the selected
step begins at the corresponding branch vertex.  Finally apply
Lemma~\ref{lem:adjacent-pair-common-constant-orbit} independently at the four
sites and use the nonemptiness just proved.  The ordered-mesh carrier permits
additional row--column intersections and shared row--column path segments,
so this construction does not assert that the four intervening path segments
form the boundary of an embedded disc.
\end{proof}

\begin{lemma}[complete boundary relation for an ordered mesh cell]
\label{lem:ordered-mesh-complete-boundary-relation}
\Sverified{Mettapedia.GraphTheory.FourColor.GoertzelV24OrderedMeshBoundaryWalk.boundaryWalk_toFinset_eq_univ}
\Sverified{Mettapedia.GraphTheory.FourColor.GoertzelV24OrderedMeshBoundaryWalk.boundarySuccessor_injective}
\Sverified{Mettapedia.GraphTheory.FourColor.GoertzelV24OrderedMeshBoundaryWalk.boundaryData_second_eq_successor_first}
\Sverified{Mettapedia.GraphTheory.FourColor.GoertzelV24OrderedMeshBoundaryWalk.boundaryCommonCarrier_property_iff_three_endpoints}
\Sverified{Mettapedia.GraphTheory.FourColor.GoertzelV24OrderedMeshBoundaryWalk.successorOverlapKempeStateSupport_nonempty}
\Sverified{Mettapedia.GraphTheory.FourColor.GoertzelV24OrderedMeshBoundaryHolonomy.successorWitnessRelation_nonempty}
\Sverified{Mettapedia.GraphTheory.FourColor.GoertzelV24OrderedMeshBoundaryHolonomy.boundaryWitnessSection_or_obstruction}
\Sverified{Mettapedia.GraphTheory.FourColor.GoertzelV24OrderedMeshBoundaryHolonomy.BoundaryWitnessSection.state_eq_overlapKempeState}
For one coordinate rectangle of an ordered injective mesh, orient every
subdivision step of its top, right, bottom and left paths clockwise.  These
steps form a finite cyclic carrier: the displayed boundary list contains
every step exactly once, and clockwise successor is a permutation.  This
includes all path steps between branch vertices, not only the four steps
selected at the corners in
Lemma~\ref{lem:ordered-mesh-coherent-orbit-cell}.

Each step canonically supplies an adjacent-pair deletion.  The terminal
vertex of a step is literally the initial vertex of its successor, including
at all four turns.  Therefore the two consecutive deletions compare on the
canonical induced common carrier obtained by deleting their three named
endpoints (with repetitions allowed).  Minimality supplies one proper base
colouring, and hence one selected Kempe orbit, at every site.  The finite
overlap-state relation from each selected site to its successor is nonempty.

Choose one realised finite overlap state on every successor edge.  Exactly
one of the following holds:
\begin{enumerate}
\item there is one concrete deletion colouring at every boundary site which
  simultaneously witnesses the chosen incoming and outgoing states; or
\item no such boundary-wide family exists.
\end{enumerate}
In the first case the state on every edge is the overlap state of the two
displayed concrete colourings.  Thus a site colouring is reused, rather than
chosen afresh, in its two incident relations.
\end{lemma}

\begin{proof}
The four side lengths are the positive differences of successive branch
positions.  List the corresponding path steps in clockwise order, reversing
the bottom and left sides.  Successor increments within a side and moves to
the first step of the next side at a turn.  A case split over the four sides,
using strict monotonicity of the branch positions at the turns, proves both
injectivity of successor and equality of consecutive endpoints.  Finiteness
then makes successor a permutation.  The four lists are disjoint by their
side tags and each finite index occurs once, proving completeness of the
boundary list.

The common-carrier statement follows by substituting the shared endpoint in
the four endpoint inequalities defining the generic two-pair deletion.
Choose the adjacent-pair data and a proper base colouring once at every
site.  The two bases at a step and its successor realise an overlap state, so
each one-step support is inhabited.  Choice over the finite carrier gives a
state assignment.  A boundary section is, by definition, a dependent family
of site colourings satisfying every selected successor relation.  Excluded
middle applied to the existence of that single family gives the stated
dichotomy, and unfolding the successor relation gives the final
state-realisation equality.
\end{proof}

\begin{lemma}[one physical Kempe site per ordered-mesh step]
\label{lem:ordered-mesh-global-site-coherence}
\Sverified{Mettapedia.GraphTheory.FourColor.GoertzelV24OrderedMeshGlobalSites.boundaryToGlobalStep_injective}
\Sverified{Mettapedia.GraphTheory.FourColor.GoertzelV24OrderedMeshGlobalSites.selectedGlobalKempeSite_boundary_endpointPair}
\Sverified{Mettapedia.GraphTheory.FourColor.GoertzelV24OrderedMeshSharedSides.boundaryToGlobalStep_east_eq_west}
\Sverified{Mettapedia.GraphTheory.FourColor.GoertzelV24OrderedMeshSharedSides.boundaryToGlobalStep_south_eq_north}
\Sverified{Mettapedia.GraphTheory.FourColor.GoertzelV24OrderedMeshSharedSides.eastGlobalStepSet_eq_westGlobalStepSet}
\Sverified{Mettapedia.GraphTheory.FourColor.GoertzelV24OrderedMeshSharedSides.southGlobalStepSet_eq_northGlobalStepSet}
\Sverified{Mettapedia.GraphTheory.FourColor.GoertzelV24AdjacentPairReversal.isTaitEdgeColoring_reverseColoring}
\Sverified{Mettapedia.GraphTheory.FourColor.GoertzelV24OrderedMeshOrientedGlobalSites.orientedGlobalBoundaryKempeSite_west_eq_reverse_east}
\Sverified{Mettapedia.GraphTheory.FourColor.GoertzelV24OrderedMeshOrientedGlobalSites.orientedGlobalBoundaryKempeSite_south_eq_reverse_north}
\Sverified{Mettapedia.GraphTheory.FourColor.GoertzelV24OrderedMeshOrientedGlobalSites.orientedGlobalBoundaryKempeSite_west_base_heq_reverse_east}
\Sverified{Mettapedia.GraphTheory.FourColor.GoertzelV24OrderedMeshOrientedGlobalSites.orientedGlobalBoundaryKempeSite_south_base_heq_reverse_north}
\Sverified{Mettapedia.GraphTheory.FourColor.GoertzelV24AdjacentPairStateReversal.vertexKirchhoffSum_reverseColoring}
\Sverified{Mettapedia.GraphTheory.FourColor.GoertzelV24AdjacentPairStateReversal.kempePortsConnected_reverseColoring_iff}
\Sverified{Mettapedia.GraphTheory.FourColor.GoertzelV24AdjacentPairStateReversal.boundaryKempeState_reverseColoring}
\Sverified{Mettapedia.GraphTheory.FourColor.GoertzelV24AdjacentPairStateReversal.reverseBoundaryKempeState_involutive}
\Sverified{Mettapedia.GraphTheory.FourColor.GoertzelV24AdjacentPairOverlapStateReversal.reverseOverlapKempeState_involutive}
\Sverified{Mettapedia.GraphTheory.FourColor.GoertzelV24AdjacentPairOverlapStateReversal.portOverlapProfile_reverse_swap}
\Sverified{Mettapedia.GraphTheory.FourColor.GoertzelV24AdjacentPairOverlapStateReversal.commonRestrictionAgreementBit_reverse_swap}
\Sverified{Mettapedia.GraphTheory.FourColor.GoertzelV24AdjacentPairOverlapStateReversal.overlapKempeState_reverse_swap}
\Sverified{Mettapedia.GraphTheory.FourColor.GoertzelV24AdjacentPairKempeClosureReversal.reverseColoring_swapOnKempeComponent}
\Sverified{Mettapedia.GraphTheory.FourColor.GoertzelV24AdjacentPairKempeClosureReversal.mem_edgeKempeClosure_reverseColoring_iff}
\Sverified{Mettapedia.GraphTheory.FourColor.GoertzelV24AdjacentPairOverlapSupportReversal.overlapKempeStateSupport_reverse_swap}
\Sverified{Mettapedia.GraphTheory.FourColor.GoertzelV24OrderedMeshSharedBoundaryStates.west_boundaryKempeState_eq_reverse_east}
\Sverified{Mettapedia.GraphTheory.FourColor.GoertzelV24OrderedMeshSharedBoundaryStates.south_boundaryKempeState_eq_reverse_north}
\Sverified{Mettapedia.GraphTheory.FourColor.GoertzelV24OrderedMeshGlobalBoundaryHolonomy.globalSuccessorWitnessRelation_nonempty}
\Sverified{Mettapedia.GraphTheory.FourColor.GoertzelV24OrderedMeshGlobalBoundaryHolonomy.globalBoundaryWitnessSection_or_obstruction}
\Sverified{Mettapedia.GraphTheory.FourColor.GoertzelV24OrderedMeshGlobalBoundaryHolonomy.not_hasGlobalBoundaryWitnessObstruction_base}
Fix one ordered injective mesh in a graph-backed vertex-minimal Tait
counterexample.  There is one globally indexed forward site for every step of
every row path and every column path, consisting of ordered adjacent-pair data
and a proper base colouring of its two-vertex deletion.  Every oriented step
on the complete boundary of a coordinate rectangle injects into this global
carrier after its cell-side tag and orientation are forgotten.  Its unordered
deleted endpoint pair is unchanged.

Use the global site directly on the north and east sides of a clockwise cell
boundary.  On the south and west sides use its canonical reversal: exchange
the two deleted endpoints, exchange the two two-port blocks, and transport the
base colouring through the canonical isomorphism between the two orders of
the same vertex deletion.  This transport preserves the nonzero Tait
condition.

If two coordinate rectangles share a side, reversal of the finite position
interval identifies their side-step carriers exactly.  Under this
identification the clockwise site of one rectangle is the canonical reversal
of the clockwise site of the other, and its base colouring is exactly the
transported base colouring.  In particular, the two rectangles do not make
independent Kempe-orbit choices on their common physical side.

This identification preserves the complete finite observation at the site.
More precisely, reversal precomposes the four-port Kirchhoff word by
\(0,1,2,3\mapsto2,3,0,1\), and it preserves every labelled assertion that two
ports lie in the same component of a specified two-colour subgraph under the
same port permutation.  Hence the full boundary Kempe state on one reading of
a shared side is exactly the involutive reversal of the state on the other
reading.  This holds for east--west and north--south neighbours.

The same transport acts on the complete two-site overlap observation.  It
transposes the source--target port-coincidence matrix after the four-port
reversal, exchanges and reverses the two boundary states, and leaves unchanged
the Boolean assertion that the two colourings agree on their common
four-vertex deletion.  This action is an involution.  Consequently, reading an
ordered pair of physical sites backwards gives exactly the reversed overlap
state, rather than merely an overlap state with the same boundary words.
The retained-vertex isomorphism also transports each selected bicoloured
component, commutes with switching that component, and therefore identifies
the two complete Kempe closures.  It follows that the entire realised overlap
support in one direction is the image, under this same involution, of the
realised support in the other direction.

For each rectangle, the successor overlap supports built from these globally
chosen sites are nonempty.  Hence every finite state assignment on its
complete boundary has the same exact alternative as in
Lemma~\ref{lem:ordered-mesh-complete-boundary-relation}: either one concrete
colouring at every globally based site simultaneously witnesses all successor
states, or no such boundary section exists.

The mesh data alone do not force the second alternative.  There is a
canonical assignment which records, on every successor edge, the overlap
state realised by the two globally chosen base colourings.  Those same base
colourings form a simultaneous boundary section, so this canonical assignment
has no witness obstruction.  Consequently a wall-facing holonomy argument
must constrain the successor states by additional ambient or replacement
data; pairwise support nonemptiness and shared-site coherence by themselves
cannot produce the obstruction.
\end{lemma}

\begin{proof}
Index a physical path step by the disjoint sum of a row index and a row-step
index, or a column index and a column-step index.  The four cell-side formulas
map into this carrier by the corresponding branch-position interval.  Within
one rectangle, equality of two images either equates two positions on the
same side or would equate its distinct top and bottom row indices, its
distinct left and right column indices, or the two summands.  This proves
injectivity.  Minimality supplies one rotation-ordered adjacent-pair Kempe
site and proper base colouring at every global index; make this choice once.

The north and east formulas traverse their path steps forward.  The south and
west formulas traverse them backward, so exchange the deleted endpoints and
the two port blocks.  Swapping the two inequalities in the retained-vertex
subtype is a graph isomorphism between the two deletion orders.  Pulling the
base colouring back along this isomorphism preserves properness and preserves
the assertion that every edge has nonzero colour.

For neighbouring rectangles, the shared row or column interval has the same
length.  The elementary identity
\[
  l+p=u-1-\bigl((u-l)-(p+1)\bigr)
\]
identifies a forward position on one boundary with the reversed position on
the other.  Substitution in the four branch-position formulas proves literal
equality of their global step indices.  The statements about adjacent-pair
data and base colourings now follow by unfolding the globally selected site
once, not by comparing two choices.

The retained-vertex swap is a graph isomorphism.  It maps the incident-edge
set at every retained vertex bijectively to the incident-edge set at its image,
so the corresponding Kirchhoff sums agree.  On each pair of colours it also
restricts to an isomorphism of the two bicoloured line subgraphs; reachability
therefore agrees for every pair of ports.  These two facts give the boundary
word and connectivity-bit formulas, and applying the port permutation twice
is the identity.  Substitution of the already proved shared-side data and
base-colouring equalities gives the east--west and north--south state
identities.

For a pair of sites, exchange source and target while applying the same
retained-vertex swap to both colourings.  The common four-vertex deletion is
carried to the reversed common deletion by the identity on ambient vertices
with the four exclusion proofs permuted.  Thus literal agreement of the two
restricted colourings is invariant.  Together with the one-site state formula
and the transposed port-coincidence matrix, this proves the complete overlap
state reversal formula and its involutivity.

For the orbit statement, map a selected bicoloured component through the
line-graph isomorphism induced by the retained-vertex swap.  Membership in the
component is preserved in both directions, so switching before transport is
the same colouring as transporting before switching.  Induction on a finite
Kempe-move sequence gives one inclusion of closures; applying the same result
after a second reversal gives the converse.  Apply the concrete overlap-state
formula pointwise to the two closure sets.  This proves equality of the
realised supports under the overlap-state involution.

Finally, two proper globally based colourings realise an overlap state on
every successor edge.  Choice of one such finite state on each edge gives an
assignment.  Excluded middle applied to a single dependent family of concrete
site colourings gives the section alternative, exactly as in the preceding
lemma.  For the canonical assignment, take at each successor edge precisely
the state realised by its two global bases.  Membership follows from
reflexivity of Kempe reachability and properness of the bases.  Reusing each
base at its site then satisfies both incident successor relations
definitionally, giving the asserted section and ruling out its obstruction.
\end{proof}

This lemma closes the carrier- and orientation-coherence problem, not the wall
theorem.  What remains is now intrinsically two-dimensional: prove that the
relational returns of neighbouring cells cancel after the displayed support
involution,
or show that their failure produces a bounded cut, a strict compositional
replacement, or another global invariant forbidden in a least
counterexample.  Such a failure cannot be obtained by selecting the canonical
base-realised successor states, since those always have a section.  No
equality of opposite-shore supports, flatness, or fixed-wall exclusion is
inferred merely from sharing the physical sites.

This lemma is deliberately an incidence and quantifier theorem.  The second
alternative is not yet a geometric obstruction, and the first does not yet
identify the supports of two physical shores.  In particular, an ordered
injective mesh may have extra row--column intersections or shared path
segments, so the cyclic carrier is not asserted to bound an embedded disc.
The next wall-facing obligation is now exact: either construct a boundary
section with the support equality required by the physical splice, or turn
failure of a section into a bounded cut, a strict replacement, or a global
invariant contradicting a least counterexample.

\begin{corollary}[complete boundary relational residue]
\label{cor:ordered-mesh-complete-boundary-residue}
\Sverified{Mettapedia.GraphTheory.FourColor.GoertzelV24OrderedMeshBoundaryHolonomyResidue.boundarySuccessor_rootPredecessor}
\Sverified{Mettapedia.GraphTheory.FourColor.GoertzelV24OrderedMeshBoundaryHolonomyResidue.nonempty_boundaryWitnessSection_iff_fixedPoint}
\Sverified{Mettapedia.GraphTheory.FourColor.GoertzelV24OrderedMeshBoundaryHolonomyResidue.hasBoundaryWitnessObstruction_iff_pairwise_and_residue}
\Sverified{Mettapedia.GraphTheory.FourColor.GoertzelV24OrderedMeshBoundaryHolonomyResidue.boundaryWitnessObstruction_residue}
Cut the complete boundary immediately before its north-west first step.  Let
(H) be the relation on proper colourings at that root obtained by imposing
every successor relation except the cut relation on one boundary-wide family,
and then using the cut relation from the last west-side colouring to a second
root colouring.  A boundary section exists if and only if (H) has a fixed
point.

Consequently, if a chosen boundary-state assignment has no section, then all
of its one-step relations are nevertheless nonempty and exactly one of the
following global residues occurs:
\begin{enumerate}
\item (H) is empty: the one-step witnesses cannot be concatenated into even
  one complete clockwise return;
\item the domain and range of (H) differ;
\item (H) branches forward or backward; or
\item (H) is the graph of a fixed-point-free permutation on its nonempty
  active support.
\end{enumerate}
\end{corollary}

\begin{proof}
The final west-side step has the last west index, so its successor is the
north-west root.  A section supplies a fixed point of (H) by cutting its
family there.  Conversely, a fixed-point witness for (H) glues the two root
colourings and satisfies every successor relation, hence is a section.

If no fixed point exists, either (H) is empty or every realised return moves
its starting colouring.  In the latter case compare the active domain and
range.  If they differ, the second residue holds.  If they agree but either
forward or backward uniqueness fails, the third holds.  Otherwise the unique
successor map is a bijection of the common active support, and the absence of
a fixed point gives the fourth residue.
\end{proof}

Only one implication is claimed:
\[
 \{\hbox{states with shared corner witnesses}\}
 \subseteq
 \{\hbox{pairwise-compatible cycles}\}.
\]
The converse is an open lifting, or sheaf, condition: equal finite boundary
observations can be represented by different intermediate colourings.  Thus
Lemma~\ref{lem:ordered-mesh-coherent-orbit-cell} does not prove flatness,
trivial holonomy, support equality across a rectangle, or fixed-wall
exclusion.  It fixes the exact input of the next wall-facing theorem: either
lift compatible local data to a coherent rectangular section and obtain the
physical support replacement, or extract a non-coboundary holonomy with a
separate structural consequence.
The finite-state dichotomy itself is checked: shared-witness lifting holds,
or one compatible cycle state explicitly fails to lift.  This dichotomy is a
specification of the two possible proof obligations, not a proof that either
alternative has the required geometric consequence.

There is an exact positive criterion.  Say that the boundary observation
\emph{separates an orbit} when two proper colourings in that Kempe orbit with
the same boundary word and all the same labelled two-colour port-connectivity
bits are literally the same colouring.  If the observation separates each of
the four corner orbits, then pairwise compatibility does lift: choose witnesses
on the four sides; corner coherence makes the two observations at a corner
equal, separation identifies the two concrete colourings, and the four
identified colourings are shared witnesses for the cell.  Contrapositively,
failure of lifting forces a nontrivial fibre at one or more corners.  Thus the
missing datum is not an unspecified topological defect: it is exactly a pair of
distinct proper orbit representatives with the same recorded boundary state.

The obstruction has an exact relational normal form.  For a fixed side state
$s_i$, let $R_i$ be the relation consisting of the pairs of concrete proper
orbit representatives which realise $s_i$.  A clockwise cell state has shared
corner witnesses precisely when
\[
   (R_{
m W}\circ R_{
m S}\circ R_{
m E}\circ R_{
m N})
      \cap \Delta \ne \varnothing;
\]
that is, the four-step relation on the north-west fibre has a fixed point.
Consequently a lifting obstruction is equivalent to a corner-coherent finite
state for which every $R_i$ is nonempty but the composite is fixed-point-free.
This is the precise sense in which the remaining datum is a relational
holonomy.  It does not yet give a geometric consequence for an embedded wall.

There are two distinct residues hidden in the phrase ``fixed-point-free.''
The fourfold composite may be empty, even though each of the four side
relations is nonempty; call this a \emph{relational break}.  Otherwise the
composite is nonempty and every completed clockwise path returns to a
different north-west witness; call this \emph{proper relational holonomy}.
These alternatives are exhaustive: if the composite is empty there is a
break, while if it is inhabited then absence of a fixed point says exactly
that every one of its pairs has distinct endpoints.  This split must precede
any signed-permutation argument, because an empty relation is not a
permutation.

The nonempty horn must be refined once more before it has a sign.  For an
endorelation (R), let its domain and range be the witnesses which occur as
first and second coordinates.  Exactly one of the following descriptions
applies: the domain and range differ (support drift); (R) is not unique in
one of its two coordinates (relational branching); or the two supports agree
and (R) is unique in both coordinates.  In the last case, choice of the
unique successor constructs a bijection of the common active support, and a
fixed-point-free (R) induces a genuine fixed-point-free permutation there.
Only this last residue can be passed to a permutation-sign invariant.  If its
active support has two elements, fixed-point-freeness makes the permutation's
support the whole two-element carrier; it is therefore a transposition and
has sign \(-1\).

An exhaustive finite diagnostic gives an important warning about the scope of
this criterion.  For the four edges of a square face in the pentagonal prism,
each selected proper deletion orbit is the proper part of its full edge-Kempe
closure, yet a pairwise-compatible cycle fails to lift; two opposite corner
observations have fibres of size two.  More strongly, a failing cycle remains
when all four literal common-restriction agreement bits are true.  In that
cycle the incident sides choose opposite representatives at exactly the two
non-singleton corners, and every non-singleton fibre in the example is joined
by one direct Kempe-component switch.  This computation has not yet been
reflected into the kernel, so it is evidence rather than a formal
counterexample.  In the displayed all-agreement obstruction each side
relation has one witness pair, but the target chosen by one side is not the
source chosen by the next at two opposite corners; its fourfold composite is
therefore empty.  Thus this example lies in the relational-break horn, not in
the proper-holonomy horn.  It nevertheless rules out treating planarity,
bridgeless cubicity, bare ordered-mesh geometry, or literal pairwise agreement
as a reason to assume a global section.  A wall argument must either control these
relational breaks using the minimal-counterexample hypothesis or derive a
target-specific consequence from a proper nonempty holonomy.

There is, however, a global consequence of literal common-core agreement
which is stronger than agreement on the four sides of one chosen cell state.

\begin{lemma}[global common-core colouring atlas]
\label{lem:global-common-core-colouring-atlas}
\Sverified{Mettapedia.GraphTheory.FourColor.GoertzelV24AdjacentPairColoringAtlas.coloring_eq_on_common_retained_edge}
\Sverified{Mettapedia.GraphTheory.FourColor.GoertzelV24AdjacentPairColoringAtlas.PairDeletionColoringFamily.not_pairwiseCommonRestrictionAgrees_of_minimal}
Let \(M\) be a graph-backed least Tait counterexample.  Let
\(\{(D_i,C_i):i\in I\}\) be a family in which \(D_i\) deletes an adjacent
vertex pair and \(C_i\) is a proper nonzero edge-colouring of the retained
induced graph.  Suppose that
\begin{enumerate}
\item every ambient edge is retained by some \(D_i\);
\item every two adjacent ambient edges are retained together by some
      \(D_i\); and
\item for every \(i,j\), the restrictions of \(C_i\) and \(C_j\) to the
      graph obtained by deleting both adjacent pairs are equal.
\end{enumerate}
Then the \(C_i\) glue to a proper nonzero Tait colouring of \(M\), a
contradiction.  Equivalently, any family satisfying the two covering
conditions has two members whose common-restriction agreement bit is false.
\end{lemma}

\begin{proof}
An ambient edge retained by both \(D_i\) and \(D_j\) has neither deleted
pair as an endpoint, so it is an edge of the exact common four-vertex
deletion.  The third hypothesis therefore says that \(C_i\) and \(C_j\)
assign it the same colour.  Choose, for each ambient edge \(e\), one index
\(i(e)\) whose deletion retains \(e\), and colour \(e\) by \(C_{i(e)}(e)\).
The preceding agreement makes this value independent of the choice.

If ambient edges \(e\) and \(f\) are adjacent, the second covering
hypothesis gives one deletion retaining both.  In that retained graph their
copies are adjacent, so its proper colouring assigns them different colours;
choice-independence transfers this inequality to the glued colouring.
Similarly, the selected local Tait condition makes every glued colour
nonzero.  Hence the glued map is Tait-colourable, contrary to the defining
property of the least counterexample.
\end{proof}

\begin{corollary}[a long ordered row forces common-core disagreement]
\label{cor:ordered-row-common-core-disagreement}
\Sverified{Mettapedia.GraphTheory.FourColor.GoertzelV24OrderedMeshColoringAtlas.exists_coloring_disagreement_in_any_nine_row_intervals}
Suppose an ordered injective mesh in a graph-backed least Tait
counterexample has a row meeting at least ten ordered branch columns.  Choose
any nine distinct intervals between successive branch columns.  At each of
the nine intervals choose arbitrarily any proper nonzero colouring of its
adjacent-pair deletion graph.  Then two of the nine chosen colourings have
false common-restriction agreement bit.
\end{corollary}

\begin{proof}
The first endpoints of the nine chosen steps are pairwise distinct because
the branch positions are strictly ordered on a simple row.  Their second
endpoints are also pairwise distinct by simplicity of the row.  More
generally, for any finite forbidden vertex set \(S\), at most \(|S|\) of
these steps can have
first endpoint in \(S\), and at most \(|S|\) can have second endpoint in
\(S\).  Thus at most \(2|S|\) selected steps meet \(S\).

For one ambient edge, its endpoint set has size two, so some selected row
step avoids it.  For two ambient edges, the union of their endpoint sets has
size at most four, so among nine selected steps one avoids both edges.  The
globally selected deletion at that step therefore retains the first edge, or
retains the two edges simultaneously.  The two covering conditions of
Lemma~\ref{lem:global-common-core-colouring-atlas} follow.  If every pair of
selected colourings had true common-restriction bit, that lemma would colour
the ambient counterexample.  Hence at least one pair among the chosen nine
has false bit.
\end{proof}

The quantifier over the nine deletion colourings is load-bearing.  In
particular, replacing the chosen bases by arbitrary Kempe-reachable
representatives cannot make all common restrictions agree simultaneously.
Thus the disagreement is an obstruction of the family, not an artefact of a
poor initial representative choice.

The first structural refinement of that disagreement is independent of
planarity.

\begin{lemma}[local transpositions generate Kempe reachability]
\label{lem:local-transpositions-generate-kempe}
\Sverified{Mettapedia.GraphTheory.FourColor.GoertzelV24LocalSwapKempeGeneration.mem_edgeKempeClosure_of_locallySwapRelated}
\Sverified{Mettapedia.GraphTheory.FourColor.GoertzelV24LocalSwapKempeGeneration.exists_not_locallySwapRelated_of_not_mem_edgeKempeClosure}
\Sverified{Mettapedia.GraphTheory.FourColor.GoertzelV24LocalSwapKempeGeneration.taitKempeReachable_of_locallySwapRelated}
\Sverified{Mettapedia.GraphTheory.FourColor.GoertzelV24LocalSwapKempeGeneration.exists_not_locallySwapRelated_of_not_taitKempeReachable}
Let \(G\) be a finite simple graph and let \(C,C'\) be two proper edge
colourings of \(G\).  Suppose that for every vertex \(v\) there are colours
\(a_v,b_v\), possibly equal, such that
\[
 C'(e)=(a_v\ b_v)C(e)
 \qquad\hbox{for every edge \(e\) incident with \(v\)}.
\]
Then \(C'\) belongs to the edge-Kempe closure of \(C\).  Consequently, if
\(C'\) is not in that closure, some vertex admits no such local
transposition.  If both colourings are Tait colourings, then the generating
sequence may moreover be chosen to switch only pairs of distinct nonzero
colours; every intermediate colouring is again a Tait colouring.  Failure of
reachability by such valid-pair switches has the same local-transposition
witness.
\end{lemma}

\begin{proof}
Choose an edge \(e\) on which \(C\) and \(C'\) disagree, and put
\(a=C(e)\), \(b=C'(e)\).  The local transposition at either endpoint of \(e\)
must be exactly \((a\ b)\): a transposition is determined by one value which
it moves.  The same argument across an adjacent \(a,b\)-coloured edge shows
that \((a\ b)\) propagates along the whole \(a,b\)-Kempe component of \(e\).
Switch that component.

Every edge of the component now agrees with \(C'\).  No agreement outside
the component is lost.  Indeed, at a vertex met by the component every
incident \(a,b\)-coloured edge lies in that same component, while an incident
edge of any other colour is fixed by \((a\ b)\); at a vertex not met by the
component the colouring is unchanged.  The switched colouring therefore
satisfies the same local-transposition hypothesis relative to \(C'\), and
the finite set of disagreeing edges has strictly smaller cardinality.  Strong
induction on that cardinality produces a finite sequence of Kempe switches
from \(C\) to \(C'\).  If \(C,C'\) are Tait colourings, the two colours
chosen at a disagreement edge are respectively \(C(e)\) and \(C'(e)\), so
they are distinct and nonzero.  Switching such a pair preserves nonzeroness,
and the same induction never leaves the Tait state space.  The two final
assertions are the corresponding contrapositives.
\end{proof}

At an interior cubic Tait vertex, both colourings use all three nonzero
colours.  Their relative local permutation is therefore the identity, a
transposition, or a three-cycle.  The first two cases are precisely those
covered by the lemma; failure of its local hypothesis is the genuine
three-cycle, or branching, sector.  At a boundary vertex of smaller degree
the conclusion is deliberately weaker and no three-cycle interpretation is
asserted.

The second horn can be compared exactly with the original deletion graph.

\begin{lemma}[common-core Kempe sequences lift or reach the exposed pair]
\label{lem:common-core-kempe-lift-or-boundary}
\Sverified{Mettapedia.GraphTheory.FourColor.GoertzelV24KempeComponentEmbeddingBoundary.touchesBoundary_or_exists_liftedComponentSwitch}
\Sverified{Mettapedia.GraphTheory.FourColor.GoertzelV24KempeComponentEmbeddingBoundary.boundaryAlongLiftOrbit_or_exists_liftedColoring}
\Sverified{Mettapedia.GraphTheory.FourColor.GoertzelV24KempeComponentEmbeddingBoundary.taitBoundaryAlongLiftOrbit_or_exists_liftedColoring}
\Sverified{Mettapedia.GraphTheory.FourColor.GoertzelV24AdjacentPairCommonCoreKempeBoundary.firstTaitOrbitReachesSecondPair_or_exists_liftedColoring}
Let \(D_{uv}\) be the graph obtained by deleting an adjacent pair
\((u,v)\), and let \(H=D_{uv}-\{x,y\}\) be its common core with the deletion
of another adjacent pair \((x,y)\).  Let \(C\) be a Tait edge colouring of
\(D_{uv}\), and suppose that a Tait edge colouring \(C^*\) of \(H\) is
reachable from \(C|_H\) by switches of distinct nonzero colour pairs.  Then
one of the following holds.
\begin{enumerate}
\item At an intermediate colouring in the Kempe orbit of \(C\), one of the
      common-core components used by the sequence is adjacent, in
      \(D_{uv}\), to a same-two-colour edge outside \(H\).  Every such
      outside edge has an endpoint \(x\) or \(y\).
\item There is a Tait colouring \(\widetilde C\), reachable from \(C\) by
      switches of distinct nonzero colour pairs, such that
      \(\widetilde C|_H=C^*\).
\end{enumerate}
\end{lemma}

\begin{proof}
Consider first one switch of an \(a,b\)-component \(K\) of \(H\).
Embed its bicoloured subgraph into the corresponding bicoloured subgraph of
\(D_{uv}\).  If an image edge of \(K\) is adjacent to an \(a,b\)-edge
outside the image, the first alternative holds.  Otherwise the image of
\(K\) is closed under bicoloured adjacency.  A connected component which is
both connected and neighbour-closed is a whole connected component of the
larger bicoloured graph.  Switching that larger component therefore
restricts literally to the prescribed switch of \(K\).

Iterate this alternative along a finite sequence witnessing
Kempe-reachability.  Until the boundary alternative occurs, each switch
lifts, so induction produces a colouring in the original deletion's Kempe
orbit with the required final restriction.  The selected colour pair at
each step is distinct and nonzero, and a component switch of such a pair
preserves the Tait condition; hence this induction stays entirely inside the
Tait state space.  Finally, an edge of \(D_{uv}\)
outside the image of \(H\) cannot have both endpoints different from
\(x,y\): if it did, attaching those endpoint inequalities would construct
its preimage edge in \(H\).  Hence every outside edge in the first
alternative is incident with \(x\) or \(y\).
\end{proof}

\begin{corollary}[exact residue of common-core disagreement]
\label{cor:common-core-disagreement-residue}
\Sverified{Mettapedia.GraphTheory.FourColor.GoertzelV24AdjacentPairCommonCoreDisagreementResidue.branching_or_taitReachesSecondPair_or_strictTaitRepair}
Let \(C_{uv}\) and \(C_{xy}\) be Tait colourings of the two adjacent-pair deletion graphs,
and suppose their restrictions to the common four-vertex deletion \(H\)
are unequal.  Then at least one of the following holds.
\begin{enumerate}
\item Some vertex of \(H\) admits no single colour transposition carrying
      the first restricted colouring to the second on all incident edges.
\item At an intermediate representative reachable from \(C_{uv}\) by
      distinct-nonzero-pair Kempe switches,
      a common-core Kempe component reaches an edge incident with \(x\) or
      \(y\).
\item There is a Tait colouring \(\widetilde C_{uv}\ne C_{uv}\), reachable
      by distinct-nonzero-pair Kempe switches from \(C_{uv}\), whose
      restriction to \(H\) equals
      \(C_{xy}|_H\).
\end{enumerate}
\end{corollary}

\begin{proof}
If \(C_{xy}|_H\) is not reachable from \(C_{uv}|_H\) by valid-pair
Tait Kempe switches, apply
Lemma~\ref{lem:local-transpositions-generate-kempe} contrapositively to get
the first alternative.  If it is reachable, apply
Lemma~\ref{lem:common-core-kempe-lift-or-boundary}.  Its boundary horn is
the second alternative and its lifting horn gives the third.  The lifted
colouring is different from \(C_{uv}\), since their common-core
restrictions were assumed unequal.
\end{proof}

The boundary alternative is an unrooted encounter, not a strand required
to start near the first deletion. Subsection~\ref{sec:separated-deletion-boundary}
below supplies ambient-separated sites, rules out common-core repair
there, and proves that this boundary alternative already has a local
witness at the original colouring. The conditional repair-cycle and
no-section results below remain valid; their no-boundary hypothesis
cannot hold when the selected family contains such a separated pair.

\begin{corollary}[nine-site Tait common-core residue]
\label{cor:nine-site-tait-common-core-residue}
\Sverified{Mettapedia.GraphTheory.FourColor.GoertzelV24OrderedMeshCommonCoreDisagreementResidue.exists_taitCommonCoreResidue_in_any_nine_row_intervals}
Fix nine injectively selected branch intervals on one row of an ordered
injective mesh, and assign an arbitrary Tait colouring to the adjacent-pair
deletion at each interval.  Then two of the nine sites satisfy at least one
of the following alternatives:
\begin{enumerate}
\item their restricted common-core colourings have a local discrepancy which
      is not a single colour transposition at some common-core vertex;
\item at an intermediate Tait representative for the first site, a
      distinct-nonzero-pair Kempe component reaches the exposed pair of the
      second site; or
\item the first site's colouring admits a strict Tait repair, reachable by
      distinct-nonzero-pair Kempe switches, whose common-core restriction
      agrees literally with the second site's colouring.
\end{enumerate}
\end{corollary}

\begin{proof}
Corollary~\ref{cor:ordered-row-common-core-disagreement} supplies two sites
whose Boolean common-restriction coordinate is false.  By the exact
characterization of that coordinate, their common-core restrictions are
unequal.  Apply Corollary~\ref{cor:common-core-disagreement-residue} to this
pair.  Its three alternatives are precisely the three alternatives above,
and its third alternative stays inside the nonzero Tait state space.
\end{proof}

For the last alternative, define a \emph{nine-site Tait state} to be one
Tait colouring at each selected deletion.  A \emph{strict repair step}
chooses two distinct coordinates whose common-core restrictions disagree,
changes the first coordinate strictly by valid-pair Tait Kempe switches,
makes its new common-core restriction agree with the second coordinate, and
leaves the other eight coordinates fixed.  This is a finite state system.

\begin{corollary}[the periodic-repair residue]
\label{cor:periodic-tait-repair-residue}
\Sverified{Mettapedia.GraphTheory.FourColor.GoertzelV24OrderedMeshCommonCoreDisagreementResidue.exists_transGen_cycle_of_finite_serial}
\Sverified{Mettapedia.GraphTheory.FourColor.GoertzelV24OrderedMeshCommonCoreDisagreementResidue.exists_strictRepairCycle_of_no_branchingOrBoundary}
Fix nine injectively selected branch intervals.  Suppose that, for every
nine-site Tait state, no pair exhibits either a local non-transposition
discrepancy or a valid-pair common-core Kempe component reaching the other
pair's exposed vertices.  Then the strict-repair relation contains a
nonempty directed cycle.
\end{corollary}

\begin{proof}
By Corollary~\ref{cor:nine-site-tait-common-core-residue}, every state has a
strict-repair successor.  Choose one successor at each state.  The state
space is finite and nonempty: each coordinate is a proper colouring of a
finite deletion graph by the three nonzero Tait colours, and the selected
deletion sites already provide one state.  Iterating the chosen successor
therefore repeats a state.  The segment between its first two occurrences is
a nonempty directed cycle of strict repair steps.  Nonemptiness is
load-bearing: each step changes its repaired coordinate strictly.
\end{proof}

\begin{corollary}[the anchored agreement-trade residue]
\label{cor:anchored-agreement-trade-residue}
\Sverified{Mettapedia.GraphTheory.FourColor.GoertzelV24OrderedMeshCommonCoreDisagreementResidue.nineSiteAgreementPairs_ssubset_of_not_collateralLoss}
\Sverified{Mettapedia.GraphTheory.FourColor.GoertzelV24OrderedMeshCommonCoreDisagreementResidue.exists_anchoredAgreementTrade_of_no_branchingOrBoundary}
Under the hypotheses of
Corollary~\ref{cor:periodic-tait-repair-residue}, some strict repair step has
the following form.  Its repaired site initially disagrees with its chosen
partner and agrees with that partner after the repair, but an ordered pair
which agreed before the repair disagrees afterwards.  The lost pair contains
the repaired site.
\end{corollary}

\begin{proof}
For a nine-site state, record the finite set of ordered site pairs whose
colourings agree on their exact common deletion.  If a strict repair loses
no old agreement, this set grows strictly: every old pair remains, while the
chosen repaired pair is new because it disagreed before the step and agrees
afterwards.  Strict growth at every step is impossible around the nonempty
directed cycle supplied by
Corollary~\ref{cor:periodic-tait-repair-residue}.  Hence one step loses an old
agreement.  Only the repaired coordinate changes in that step, so a pair
not containing it has the same two colourings before and after and cannot be
the lost pair.  Thus the lost agreement is anchored at the repaired site.
\end{proof}

\begin{corollary}[the arc-consistent no-section residue]
\label{cor:arc-consistent-no-section-residue}
\Sverified{Mettapedia.GraphTheory.FourColor.GoertzelV24OrderedMeshCommonCoreArcConsistencyResidue.arcConsistentNoSectionResidue_of_no_branchingOrBoundary}
Under the same no-branching and no-boundary hypothesis, the nine deletion
sites form a finite binary constraint system with both of the following
properties.
\begin{enumerate}
\item For any two distinct sites, any Tait colouring at the first site, and
      any Tait colouring at the second, the first colouring has a
      distinct-nonzero-pair Kempe-reachable representative agreeing with the
      second on their exact common deletion.
\item No simultaneous choice of one Tait colouring at every site agrees on
      the exact common deletion for every ordered pair of sites.
\end{enumerate}
Thus every binary projection is repair-surjective, but the constraint system
has no global section.
\end{corollary}

\begin{proof}
Extend the two arbitrary supplied colourings to a nine-site state by using
the canonically selected base colourings at the other seven sites.  If the
two supplied colourings already agree, use the empty Kempe sequence.  If
they disagree, apply
Corollary~\ref{cor:common-core-disagreement-residue}.  The standing
hypothesis excludes its branching and boundary-reaching alternatives for
this extended state, so its strict-repair alternative gives the required
representative.  This proves repair-surjectivity of every distinct binary
constraint.  A global section would contradict
Corollary~\ref{cor:nine-site-tait-common-core-residue}, which supplies a
disagreeing pair for every simultaneous nine-site assignment.
\end{proof}

The absence of a global section in the preceding corollary must not be
confused with the absence of a section on a path.  Acyclic binary constraints
have a stronger elementary property which is useful here.

\begin{corollary}[a consecutive common-core repair section]
\label{cor:consecutive-common-core-repair-section}
\Sverified{Mettapedia.Logic.PathConstraint.sectionEndingAt}
\Sverified{Mettapedia.Logic.PathConstraint.section_observe_eq_last}
\Sverified{Mettapedia.GraphTheory.FourColor.OrderedMeshCommonCorePathSection.exists_consecutive_commonCoreRepair_pathSection_of_no_branchingOrBoundary}
\Sverified{Mettapedia.GraphTheory.FourColor.OrderedMeshCommonCorePathSection.card_commonCoreRepairPathFootprint_le}
\Sverified{Mettapedia.GraphTheory.FourColor.OrderedMeshCommonCorePathSection.commonCoreColor_eq}
\Sverified{Mettapedia.GraphTheory.FourColor.OrderedMeshCommonCorePathSection.exists_pair_pathDeletionMatchingState_absentColor_eq}
\Sverified{Mettapedia.GraphTheory.FourColor.OrderedMeshCommonCorePathSection.exists_pair_pathDeletionMatchings_agree_outside_footprint}
\Sverified{Mettapedia.GraphTheory.FourColor.Compositional.AlternatingOverlapGeometry.card_alternatingGraph_edgeDisagreement_le_of_pairings_agree_outside}
\Sverified{Mettapedia.GraphTheory.FourColor.OrderedMeshCommonCorePathSection.exists_pair_pathAlternatingGraph_edgeDisagreement_le}
\Sverified{Mettapedia.GraphTheory.FourColor.Compositional.DeletionColorMatching.DeletionMatchingState.pairing_partner_first}
\Sverified{Mettapedia.GraphTheory.FourColor.Compositional.AlternatingSiteGeometry.nonempty_residualDefectMinimizer}
\Sverified{Mettapedia.GraphTheory.FourColor.Compositional.AlternatingSiteGeometry.properAlternatingComponentWitness_of_partner_ne}
\Sverified{Mettapedia.GraphTheory.FourColor.OrderedMeshCommonCorePathSection.exists_pair_pathAlternatingEndpointAlternatives}
\Sverified{Mettapedia.GraphTheory.FourColor.Compositional.DeletionSiteGeometry.DeletionAlternatingSiteReceipt.ofState}
\Sverified{Mettapedia.GraphTheory.FourColor.Compositional.DeletionSiteGeometry.DeletionTwoSectorReturnReceipt.ofState}
\Sverified{Mettapedia.GraphTheory.FourColor.OrderedMeshCommonCorePathSection.pathTwoSectorSiteAlternative}
\Sverified{Mettapedia.GraphTheory.FourColor.OrderedMeshCommonCorePathSection.exists_pair_pathTwoSectorEndpointAlternatives}
Under the no-branching and no-boundary hypothesis of
Corollary~\ref{cor:arc-consistent-no-section-residue}, one can choose a proper
nonzero edge-colouring at each of the nine ordered deletion sites so that
consecutive choices agree on their exact common deletion.  Each nonterminal
choice is reachable from the canonical colouring at its site by valid-pair
Kempe moves, and the final choice may be prescribed to be the canonical final
colouring.  Moreover, two of the nine colourings may be completed using the
same absent Tait colour.  The resulting ambient perfect matchings agree outside
the union of the nine deletion footprints, a set of at most $45$ edges.  For
any fixed graph-supported perfect matching, the corresponding two
matching-symmetric-difference graphs consequently disagree on at most $45$
edges.  At each of these two endpoints the fixed residual-defect minimizer
either uses the restored edge or the exact path-selected deletion colouring
carries a proper alternating cycle, facial bond, physical formation,
canonical return, face-shore data, and two-sector noncrossing receipt.
\end{corollary}

\begin{proof}
For the edge between consecutive sites $i$ and $i+1$, relate a colouring at
$i$ to one at $i+1$ when the former is Kempe-reachable from the canonical
colouring at $i$ and the two colourings agree on their exact common deletion.
Pairwise repair-surjectivity says that this relation is right-total: every
target colouring has a related source colouring.  Starting with the prescribed
colouring at the last site and choosing a predecessor successively constructs
the section.  Formally this is the generic dependent path-constraint theorem,
proved by reverse induction on the finite path.

There is no contradiction with the preceding no-section result.  The latter
forbids one assignment satisfying every ordered pairwise constraint; the
present assignment satisfies only the eight consecutive constraints.

For any ambient edge retained at all nine sites, consecutive common-core
agreement telescopes along the path, so all nine selected colourings give that
edge the same colour.  Each boundary word is constant on its two sides and
therefore omits at least one of the three Tait colours.  Choose one omitted
colour at each site.  Among nine choices, two coincide.  Complete these two
colour classes by the respective deleted central edges.  Their matching
decisions agree on every edge retained at every site.  The exceptional set is
therefore contained in the union of the nine deletion footprints; each
footprint has five edges, so the union has cardinality at most $45$.

Taking symmetric difference with the same first matching preserves edgewise
agreement, so it cannot enlarge this exceptional carrier.

Fix an exchange-rigid residual-defect minimizer.  At each selected deletion
site its completed absent-colour matching contains the restored central edge.
Thus either the minimizer also contains that edge, or their disagreement at
its first endpoint generates a proper alternating cycle through the edge.
In the latter case the cycle has at least four vertices, a shared matching
edge lies outside it, and exchange on its carrier cannot lower residual
defect.  The two pigeonhole-selected endpoints carry this exact alternative
simultaneously with the $45$-edge disagreement bound.

The geometric construction is independent of how the deletion colouring was
chosen.  Starting from its completed absent-colour matching and the proper
alternating cycle, take the two connected facial shores of the cycle, perform
the local formation switch at each boundary vertex, and pair boundary vertices
by their canonical returns through unchanged material.  The return preserves
the facial shore, and equal-shore returns do not cross.  Indexing this receipt
by the deletion state ensures that the resulting geometry belongs to the
path-repaired colouring itself, not to the unrelated globally selected
colouring at the same physical edge.

This removes the previously stated global colour-gauge obligation.  It does
not yet prove wall exclusion: bounded edge-disagreement is now verified.  The
sparse matching-dependent refinement below disposes of the branch in which
the minimizer uses a selected central edge at the nine sites needed by the
path argument.  It remains to compare the distinguished components across
those sites, rather than merely their ambient alternating graphs, and then use
the mesh geometry to force a strict compositional descent.
\end{proof}

\begin{corollary}[a sparse noncentral mesh deletion atlas]
\label{cor:sparse-noncentral-mesh-deletion-atlas}
\Sverified{Mettapedia.GraphTheory.FourColor.Compositional.DeletionAtlasPath.pairwiseRepairSurjectivity_of_no_branchingOrBoundary}
\Sverified{Mettapedia.GraphTheory.FourColor.Compositional.DeletionAtlasPath.exists_pathSection_of_no_branchingOrBoundary}
\Sverified{Mettapedia.GraphTheory.FourColor.Compositional.DeletionAtlasPath.hasCoherentTwoSectorPair_of_no_branchingOrBoundary}
\Sverified{Mettapedia.GraphTheory.FourColor.Compositional.SparseNoncentralMeshAtlas.sparse_branchPosition_add_two_le}
\Sverified{Mettapedia.GraphTheory.FourColor.Compositional.SparseNoncentralMeshAtlas.selectedGlobalStep_noncentral}
\Sverified{Mettapedia.GraphTheory.FourColor.Compositional.SparseNoncentralMeshAtlas.coloringFamily_coversEdges}
\Sverified{Mettapedia.GraphTheory.FourColor.Compositional.SparseNoncentralMeshAtlas.coloringFamily_coversAdjacentEdges}
\Sverified{Mettapedia.GraphTheory.FourColor.Compositional.SparseNoncentralMeshAtlas.branchingOrBoundary_or_hasCoherentTwoSectorPair}
Let an ordered injective mesh have twenty column paths, fix one row, and fix a
minimum-residual-defect matching.  Select the nine internal row junctions at
columns (1,3,\ldots,17).  At each junction select the outgoing arm if the
matching does not use it, and otherwise select the incoming arm.  These nine
arms are noncentral for the same matching.  Their adjacent-pair deletion
atlas covers every ambient edge and every adjacent pair of ambient edges.

Moreover exactly one of the following conclusions is available.
\begin{enumerate}
\item Some simultaneous choice of the nine deletion colourings exposes either
      a local non-transposition common-core discrepancy or a valid-pair Kempe
      orbit reaching the other deleted pair.
\item There is a consecutive common-core repair section through the nine
      sites, ending at the canonical final colouring, in which two sites have
      the same absent Tait colour and both carry the exact two-sector
      residual-return receipt for the fixed matching.
\end{enumerate}
\end{corollary}

\begin{proof}
At an internal row junction the incoming and outgoing arms share the branch
vertex but have distinct outer endpoints.  A perfect matching cannot contain
both arms, so the stated rule always selects a noncentral arm.  Consecutive
selected junctions have one unused branch column between them.  Consequently
the chosen row-step positions are strictly increasing even when the
incoming/outgoing choice varies, and both endpoint maps from the nine sites
are injective.

An ambient edge forbids its two endpoints; an adjacent pair of ambient edges
forbids at most four vertices.  Each forbidden vertex removes at most one
site through each of the two injective endpoint maps.  Thus at most eight of
the nine sites are removed, leaving a deletion patch which retains the edge
or the adjacent pair.  This proves the two covering assertions.

If a nine-site colouring assignment exposes one of the two horns, take the
first alternative.  Otherwise the exact common-core repair theorem makes
every consecutive binary constraint right-total.  Choose predecessors from
the prescribed final colouring to obtain a path section.  Complete each path
colouring by a Tait colour absent from its four boundary requests.  Among nine
choices from three colours, two agree.  Since every selected arm is
noncentral for the fixed matching, the state-indexed residual construction
equips both sites with their exact two-sector return receipts, proving the
second alternative.
\end{proof}

\begin{corollary}[finite component state along the sparse deletion path]
\label{cor:sparse-deletion-path-localization}
\Sverified{Mettapedia.GraphTheory.FourColor.Compositional.DeletionPathLocalization.pathFootprint}
\Sverified{Mettapedia.GraphTheory.FourColor.Compositional.DeletionPathLocalization.card_pathFootprint_le}
\Sverified{Mettapedia.GraphTheory.FourColor.Compositional.DeletionPathLocalization.pathColor_eq}
\Sverified{Mettapedia.GraphTheory.FourColor.Compositional.DeletionPathLocalization.matchingStates_agree_outside_pathFootprint}
\Sverified{Mettapedia.GraphTheory.FourColor.Compositional.DeletionPathLocalization.alternatingGraphs_agree_outside_pathFootprint}
\Sverified{Mettapedia.GraphTheory.FourColor.Compositional.DeletionPathLocalization.card_alternatingGraph_edgeDisagreement_le}
\Sverified{Mettapedia.GraphTheory.FourColor.Compositional.DeletionPathLocalization.alternatingWalk_component_eq_in_target}
\Sverified{Mettapedia.GraphTheory.FourColor.Compositional.DeletionPathLocalization.exists_endpoint_edgeDisagreement_le_of_hasCoherentTwoSectorPair}
\Sverified{Mettapedia.GraphTheory.FourColor.Compositional.DeletionPathLocalization.connectedComponentMk_eq_distinguished_of_mem_targetCarrier}
\Sverified{Mettapedia.GraphTheory.FourColor.Compositional.DeletionPathLocalization.alternatingWalk_lands_in_target_distinguished_of_anchor}
\Sverified{Mettapedia.GraphTheory.FourColor.Compositional.DeletionPathConnectivity.card_pathBoundaryVertices_le}
\Sverified{Mettapedia.GraphTheory.FourColor.Compositional.DeletionPathConnectivity.pathAlternatingGraphs_deleteFootprint_eq}
\Sverified{Mettapedia.GraphTheory.FourColor.Compositional.DeletionPathConnectivity.pathAlternatingGraph_reachable_iff_boundaryClosure}
\Sverified{Mettapedia.GraphTheory.FourColor.Compositional.DeletionPathConnectivity.selectedFirstVertices_component_eq_iff_boundaryClosure}
\Sverified{Mettapedia.GraphTheory.FourColor.Compositional.DeletionPathConnectivityState.card_pathBoundarySlot}
\Sverified{Mettapedia.GraphTheory.FourColor.Compositional.DeletionPathConnectivityState.exists_pathBoundarySlot}
\Sverified{Mettapedia.GraphTheory.FourColor.Compositional.DeletionPathConnectivityState.pathConnectivityState_reachable_iff}
\Sverified{Mettapedia.GraphTheory.FourColor.Compositional.DeletionPathConnectivityState.card_pathConnectivityState_le}
In the second alternative of
Corollary~\ref{cor:sparse-noncentral-mesh-deletion-atlas}, let $P$ be the
union of the nine deletion footprints, and let $B$ be the set of endpoints
of edges in $P$.  Then
\[
             |P|\leq45,
             \qquad |B|\leq 2|P|\leq90.
\]
The two receipt-bearing completed matchings agree on every ambient edge
outside $P$, and their symmetric-difference graphs $A_i,A_j$ with the fixed
residual-defect minimizer satisfy $A_i-P=A_j-P$.

For any path coordinate $r$ and any $x,y\in B$, membership of $x$ and $y$ in
one connected component of $A_r$ is exactly the reflexive--transitive closure
on $B$ of the following two moves:
\begin{enumerate}
  \item $x$ and $y$ lie in one component of $A_r-P$;
  \item $x$ and $y$ lie in one component of the graph consisting of the
        edges of $A_r$ which belong to $P$.
\end{enumerate}
Every selected central edge belongs to $P$, so both of its endpoints lie in
$B$.  In particular, equality of the components generated by any two
selected central roots is completely determined by this relation on at most
ninety vertices.

This finite relation has a fixed-carrier presentation.  Give each of the nine
sites two slots for its central edge and two slots for each of its four
boundary edges.  The resulting set $S$ has
\[
                       |S|=9(2+4\cdot2)=90.
\]
Every vertex of $B$ is named by at least one slot; coincident names are
allowed.  For each coordinate $r$, form a simple graph $C_r$ on $S$ by
joining two distinct slots exactly when their physical vertices belong to
one component of $A_r$.  Then reachability in $C_r$ is equivalent to
reachability of the represented vertices in $A_r$, and
\[
              \#\{\hbox{possible labelled states }C_r\}
                    \leq 2^{90\cdot90}.
\]

Moreover, every walk in $A_i$ whose edges avoid $P$ transfers edge by edge to
$A_j$.  If one endpoint of that transferred walk lies in the distinguished
alternating carrier recorded by the second site's two-sector receipt, then
the transferred component is the second site's distinguished component.
\end{corollary}

\begin{proof}
Evaluate an ambient edge retained by every deletion at each coordinate of
the coherent path.  Consecutive common-core compatibility gives equality at
successive coordinates, and hence equality at any two coordinates by
telescoping along the path.  The two selected coordinates use the same
absent Tait colour, so the corresponding completed colour-class matchings
agree on such an edge.  An edge which is not retained at some coordinate is
in that coordinate's five-edge deletion footprint.  Therefore every
possible disagreement is contained in $P$, and the union bound gives
$|P|\leq9\cdot5=45$.

Taking symmetric difference with one fixed matching is an edgewise Boolean
operation, so equality outside $P$ is preserved.  The set $B$ is the union of
the two-element endpoint sets of the edges in $P$, whence
$|B|\leq2|P|\leq90$.

For a generic graph $A$, write $A_P$ for the graph consisting of the edges of
$A$ in $P$.  Edgewise,
\[
                       A=(A-P)\cup A_P.
\]
Split a walk in this union whenever it changes between the two factors.  A
nontrivial segment in $A_P$ begins and ends at endpoints of edges of $P$, so
every switching point is in $B$.  The resulting list of switching points is
the asserted closure chain.  Conversely, concatenate witnesses for the
successive closure moves.  This proves the exact finite component
factorization.  Each selected central edge is one of the footprint edges, so
its two central vertices are in $B$, giving the central-root specialization.

The fixed slots simply enumerate the named endpoints used to define each
deletion: the two central vertices, and the deleted and retained endpoint of
each boundary edge.  Membership in the explicit five-edge footprint shows
that these ninety names cover $B$.  Recording the ambient component relation
between their images gives $C_r$.  A path in $C_r$ concatenates ambient
reachability witnesses; conversely, two distinct slots whose images are
ambient-reachable are adjacent in $C_r$ (and equal slots are reachable by the
empty path).  Thus the encoding is exact.  Finally, a simple graph on ninety
labels injects into its $90\times90$ Boolean adjacency matrix, giving the
displayed coarse state bound.

Finally transfer any supplied $P$-avoiding walk edge by edge to the second
graph.  The transferred walk is a reachability witness, which is exactly
equality of its endpoint connected components.

The target receipt identifies its carrier with the alternating component
generated by the target central vertex.  Thus membership of the transferred
walk's endpoint in that carrier identifies its target connected component
with the distinguished one.  Transitivity with the transferred-walk
component equality proves the anchor conclusion.

This is a finite-state replacement for an unjustified global-anchor claim,
not an identification of distinguished components at different coordinates.
The common exterior factor is literally the same at equal-absent-colour
coordinates; all possible splitting and merging is recorded by the finite
exceptional factor on $P$.  The remaining synchronization obligation is to
obtain repetition, or a strict descent, for these boundary closure states
across the mesh.  No cross-coordinate component equality is inferred merely
from the $45$-edge disagreement bound.
\end{proof}

\begin{corollary}[fixed-colour localization of a synchronization failure]
\label{cor:fixed-colour-connectivity-transition}
\Sverified{Mettapedia.GraphTheory.FourColor.Compositional.CentralAugmentedColorClass.centralAugmentedColorClassGraph_eq_pairingGraph}
\Sverified{Mettapedia.GraphTheory.FourColor.Compositional.CentralAugmentedColorClass.fixedColorAlternatingGraphs_delete_overlapFootprints_eq}
\Sverified{Mettapedia.GraphTheory.FourColor.Compositional.FixedColorConnectivityTransition.pathFixedColorAlternatingGraph_eq_pathAlternatingGraph_of_eq_absent}
\Sverified{Mettapedia.GraphTheory.FourColor.Compositional.FixedColorConnectivityTransition.synchronized_or_fixedColorLocalTransition}
\Sverified{Mettapedia.GraphTheory.FourColor.Compositional.FixedColorConnectivityTransition.FixedColorLocalConnectivityTransitionReceipt.deleteLocalFootprint_eq}
\Sverified{Mettapedia.GraphTheory.FourColor.Compositional.FixedColorConnectivityTransition.FixedColorLocalConnectivityTransitionReceipt.card_localFootprint_le}
\Sverified{Mettapedia.GraphTheory.FourColor.Compositional.FixedColorConnectivityTransition.FixedColorLocalConnectivityTransitionReceipt.not_reachable_deleteLocalFootprint}
\Sverified{Mettapedia.GraphTheory.FourColor.Compositional.FixedColorConnectivityTransition.FixedColorLocalConnectivityTransitionReceipt.exists_walk_hitting_localFootprint}
\Sverified{Mettapedia.GraphTheory.FourColor.Compositional.FixedColorConnectivityAlternative.card_fixedColorConnectivityTransitionCode}
\Sverified{Mettapedia.GraphTheory.FourColor.Compositional.FixedColorConnectivityAlternative.FixedColorLocalConnectivityTransitionReceipt.toCode_injective}
\Sverified{Mettapedia.GraphTheory.FourColor.Compositional.FixedColorConnectivityAlternative.FixedColorLocalConnectivityTransitionReceipt.card_edgeDisagreement_le}
\Sverified{Mettapedia.GraphTheory.FourColor.Compositional.FixedColorConnectivityAlternative.branchingOrBoundary_or_synchronized_or_fixedColorLocalTransition}
Let $i$ and $j$ be the two coordinates selected by the coherent geometric
sweep, and let $c$ be their common absent Tait colour.  Follow this one fixed
colour through every coordinate of the nine-site deletion path.  At a
coordinate $r$, restore its deleted central edge with colour $c$, retain the
colours of all other inserted edges, and call the resulting colour-class graph
$H_r$.  With the fixed residual-defect matching $\sigma$, put
\[
                         A^c_r=\sigma\mathbin\triangle H_r.
\]
At the two selected coordinates, $c$ is absent from the four boundary
requests.  Hence $H_i$ and $H_j$ are exactly the completed perfect matchings
used by the residual-return construction, so $A^c_i=A_i$ and $A^c_j=A_j$.

For consecutive coordinates $r,r+1$, the common-core repair condition makes
$H_r$ and $H_{r+1}$ agree away from the union of their two deletion
footprints.  Taking symmetric difference with the same $\sigma$ preserves
this agreement.  The union contains at most ten edges.  Consequently one of
the following alternatives is available:
\begin{enumerate}
  \item the labelled boundary-connectivity states at $i$ and $j$ agree; or
  \item for some consecutive $r,r+1$ and two named boundary slots $x,y$,
        reachability of their physical vertices differs between $A^c_r$ and
        $A^c_{r+1}$.  Deleting the two local footprints makes the graphs equal
        and disconnects $x$ from $y$; in whichever graph connects them, every
        connecting walk meets one of those at most ten footprint edges.
\end{enumerate}
There is no third alternative corresponding merely to a change in the
locally chosen absent colour.

The local transition has a fixed finite presentation: its adjacent coordinate
is one of eight choices, and each of $x,y$ is one of the ninety labelled
slots.  Forgetting only proof fields injects the transition receipts into a
set of cardinality
\[
                           8\cdot90\cdot90=64800.
\]
Combining this dichotomy with the sparse mesh atlas gives the row-level
alternative in its consumer form: a branching/boundary horn, a coherent pair
with synchronized connectivity, or one of these finite local transitions.
\end{corollary}

\begin{proof}
The total insertion labelling is defined for every fixed Tait colour, whether
or not that colour is absent from the four requests at the current deletion.
When it is absent, the restored central edge and the retained colour class
form the completed perfect matching, proving the endpoint identifications.

On an edge retained by two consecutive deletions, common-core compatibility
is literal equality of the two deletion colourings.  The total insertion
labellings, their fixed-colour graphs, and finally their symmetric differences
with $\sigma$ therefore agree on that edge.  After deleting the two explicit
five-edge footprints the consecutive graphs are equal.

If the endpoint connectivity states differ, the elementary adjacent-change
lemma for a finite path supplies consecutive states which differ.  Exactness
of the fixed ninety-slot encoding then supplies $x,y$ with differing ambient
reachability.  The generic finite-edge localization theorem says that such
vertices are disconnected in the common deleted graph and that a connecting
walk on the connected side must traverse the deleted carrier.  The ten-edge
bound is the union bound for the two footprints.

The construction follows the endpoint colour $c$ throughout; it never
reselects the locally absent colour at an intermediate coordinate.  Thus a
change of local completion choice creates no separate logical branch.  The
remaining obstruction is geometric: a real component split or merge supported
on one explicitly bounded physical patch.
\end{proof}

\begin{corollary}[finite repetition of fixed-colour row transitions]
\label{cor:fixed-colour-row-transition-repetition}
\Sverified{Mettapedia.GraphTheory.FourColor.Compositional.FixedColorConnectivityPigeonhole.card_fixedColorRowTransitionCode}
\Sverified{Mettapedia.GraphTheory.FourColor.Compositional.FixedColorConnectivityPigeonhole.FixedColorRowTransitionReceipt.toCode_eq_iff}
\Sverified{Mettapedia.GraphTheory.FourColor.Compositional.FixedColorConnectivityPigeonhole.nonempty_fixedColorRowTransitionReceipt_of_no_horn_of_no_synchronized}
\Sverified{Mettapedia.GraphTheory.FourColor.Compositional.FixedColorConnectivityPigeonhole.exists_distinct_rows_with_equal_fixedColorTransitionCode}
\Sverified{Mettapedia.GraphTheory.FourColor.Compositional.FixedColorConnectivityPigeonhole.branchingOrBoundary_or_synchronized_or_repeatedFixedColorTransition}
Let an ordered injective mesh have $a>194400$ rows and twenty columns.
Fix the residual-defect minimizer used in the preceding construction.  Then
one of the following holds:
\begin{enumerate}
  \item some row exposes the branching/boundary horn of the sparse mesh
        atlas;
  \item some row has a coherent geometric sweep pair with synchronized
        labelled boundary connectivity; or
  \item two distinct rows have local fixed-colour transition receipts with
        equal finite codes.
\end{enumerate}
The repeated code records the fixed colour, one of the eight adjacent
deletion-path coordinates, and the two labelled boundary slots whose
reachability changes.  Its state space has cardinality
\[
                   3\cdot8\cdot90\cdot90=194400.
\]
Equality of these codes is only equality of the observable transition shape.
It does not identify the two ambient row pieces and is not, by itself, a
physical replacement theorem.
\end{corollary}

\begin{proof}
If either of the first two alternatives occurs, there is nothing to prove.
Otherwise the row-level alternative of
Corollary~\ref{cor:fixed-colour-connectivity-transition} supplies a local
transition receipt on every row.  Attach to its $64800$-element local code
the one fixed Tait colour followed through that row.  This gives a function
from the $a$ rows to a $194400$-element type.  The finite pigeonhole principle
gives two distinct rows with equal codes.

The receipts are retained in the conclusion rather than replaced by their
codes.  This makes the next obligation explicit: the two physical patches
must be placed in a common geometric frame before equal observables can be
used for a splice or a strict descent.
\end{proof}

\begin{corollary}[finite repetition of geometrically framed row transitions]
\label{cor:fixed-colour-row-transition-frame-repetition}
\Sverified{Mettapedia.GraphTheory.FourColor.Compositional.FixedColorConnectivityFrame.selectedIncomingArmBit_eq_true_iff}
\Sverified{Mettapedia.GraphTheory.FourColor.Compositional.FixedColorConnectivityFrame.fixedColorRowTransitionFrameCode_eq_iff}
\Sverified{Mettapedia.GraphTheory.FourColor.Compositional.FixedColorConnectivityFrame.exists_distinct_rows_with_equal_fixedColorTransitionFrameCode}
\Sverified{Mettapedia.GraphTheory.FourColor.Compositional.FixedColorConnectivityFrame.branchingOrBoundary_or_synchronized_or_repeatedFixedColorTransitionFrame}
Let $\mathcal F$ be the finite type whose elements consist of a fixed-colour
row-transition code, the two Boolean choices of row arm at its adjacent
deletions, and the complete labelled boundary-connectivity states immediately
before and after the change.  If an ordered injective mesh has more than
$|\mathcal F|$ rows, then one of the first two alternatives of
Corollary~\ref{cor:fixed-colour-row-transition-repetition} holds, or two
distinct rows carry equal elements of $\mathcal F$.

Consequently the repeated rows agree on the fixed colour, transition
coordinate, witness slots, the two local arm choices, and every reachability
query among the ninety labelled boundary slots on both sides of the
transition.  This is still a semantic boundary frame: it neither identifies
the ambient vertices of the two rows nor proves equality of their exact
colouring-support profiles.
\end{corollary}

\begin{proof}
When neither earlier alternative occurs, choose the checked local transition
receipt supplied on every row.  Extend its $194400$-element code by the two
arm-selection bits and its two fixed boundary-connectivity graphs.  The
resulting type $\mathcal F$ is finite because its boundary carrier has ninety
slots.  The finite pigeonhole principle therefore supplies two distinct rows
with equal framed codes.  Projecting that equality gives every listed
agreement.

The cardinality is kept as $|\mathcal F|$ rather than expanded into a closed
decimal numeral.  This is both mathematically exact and proof-assistant
friendly: an adjacency-matrix estimate is available, but normalizing the
corresponding power of two contributes nothing to the argument.
\end{proof}

\begin{lemma}[semantic content of an equal fixed-colour transition frame]
\label{lem:fixed-colour-transition-frame-semantics}
\Sverified{Mettapedia.GraphTheory.FourColor.Compositional.MeshJunctionAlternatingGeometry.incomingRowStep_ne_outgoingRowStep}
\Sverified{Mettapedia.GraphTheory.FourColor.Compositional.FixedColorConnectivityTransition.fixedConnectivityStateAt_reachable_iff}
\Sverified{Mettapedia.GraphTheory.FourColor.Compositional.FixedColorConnectivityFrameSemantics.selectedIncomingArmBit_eq_true_iff_selectedRowStep_eq_incoming}
\Sverified{Mettapedia.GraphTheory.FourColor.Compositional.FixedColorConnectivityFrameSemantics.semanticAgreement_of_fixedColorRowTransitionFrameCode_eq}
\Sverified{Mettapedia.GraphTheory.FourColor.Compositional.FixedColorConnectivityFrameSemantics.branchingOrBoundary_or_synchronized_or_repeatedFixedColorSemanticAgreement}
If two row-transition receipts have equal framed codes, then, under their
common labelled ninety-slot coordinates, they choose the incoming row arm at
the same endpoint coordinates and have the same fixed-colour reachability
relation immediately before and immediately after the transition.

This conclusion compares semantic observations on a fixed finite carrier.  It
does not identify the ambient vertices named by corresponding slots, and it
does not assert equality of exact colouring-support profiles.

Consequently the consumer-facing mesh alternative can state its third branch
directly in these semantic terms: there is a branching or boundary horn, there
is already a synchronized row, or two distinct rows have the agreement above.
\end{lemma}

\begin{proof}
The incoming and outgoing row steps at an internal mesh junction are distinct.
Hence each Boolean arm bit is equivalent to choosing the incoming step.
Equality of framed codes gives equality of both bits and both labelled
connectivity graphs.  The definition of a labelled connectivity state is
exact: an edge in its transitive closure is precisely fixed-colour ambient
reachability between the physical vertices named by the corresponding slots.
Projecting the frame equality and applying these two equivalences proves the
claim.
\end{proof}

The binary constraints have one further exact locality property.  Although
common-core agreement is not transitive as a relation on whole pairwise
common deletions, its failure of transitivity cannot occur away from the
middle deletion.

\begin{lemma}[three-patch common-core localization]
\label{lem:three-patch-common-core-localization}
\Sverified{Mettapedia.GraphTheory.FourColor.GoertzelV24AdjacentPairCommonCoreLocalization.commonCoreAgrees_iff_ambient}
\Sverified{Mettapedia.GraphTheory.FourColor.GoertzelV24AdjacentPairCommonCoreLocalization.commonCoreAgrees_symm}
\Sverified{Mettapedia.GraphTheory.FourColor.GoertzelV24AdjacentPairCommonCoreLocalization.commonCoreDisagreementLocalizedAt_of_chain}
\Sverified{Mettapedia.GraphTheory.FourColor.GoertzelV24AdjacentPairInsertion.AdjacentPairData.not_isRetainedAmbientEdge_iff}
Let $D_A,D_B,D_C$ be three adjacent-pair deletions of one finite simple
graph, with edge-colourings $A,B,C$ of their retained induced graphs.
Suppose $A$ and $B$ agree on their exact common deletion, and $B$ and
$C$ agree on theirs.  If $A$ and $C$ disagree on their exact common
deletion, then some ambient edge $e$ has all of the following properties:
it is retained by $D_A$ and $D_C$, it is deleted by $D_B$, and
$A(e)\ne C(e)$.  In particular $e$ is either the central edge or one of
the four boundary edges incident with the adjacent pair deleted by $D_B$.
\end{lemma}

\begin{proof}
Equality of the restrictions of two deletion colourings is equivalent to
literal equality on every ambient edge retained by both.  The forward
direction evaluates the common restriction; conversely, every edge of the
common deletion has a canonical ambient image retained by both deletion
patches, so ambient equality gives equality of the restriction functions.
This characterization also proves symmetry, despite the two orientations
using definitionally different subtype carriers.

Since the $A$- and $C$-restrictions are unequal, choose an edge of their
common deletion on which their colours differ, and regard it as an ambient
edge $e$.  If $D_B$ retained $e$, then $A(e)=B(e)$ by the first
agreement and $B(e)=C(e)$ by the second, contradicting the choice of
$e$.  Hence $D_B$ deletes $e$.  An adjacent-pair deletion removes
exactly its central edge and its four boundary edges, which gives the final
assertion.
\end{proof}

\begin{corollary}[localized anchored agreement trade]
\label{cor:localized-anchored-agreement-trade}
\Sverified{Mettapedia.GraphTheory.FourColor.GoertzelV24OrderedMeshCommonCoreLocalizationResidue.partnerDisagreementOrLocalizedCollateralLoss_of_anchoredTrade}
\Sverified{Mettapedia.GraphTheory.FourColor.GoertzelV24OrderedMeshCommonCoreLocalizationResidue.exists_partnerDisagreementOrLocalizedCollateralLoss_of_no_branchingOrBoundary}
Under the no-branching and no-boundary hypothesis of
Corollary~\ref{cor:periodic-tait-repair-residue}, some strict repair ends in
a state with a repaired site $R$, its gained partner $G$, and a lost
partner $L$ such that one of the following holds:
\begin{enumerate}
\item the colourings at $G$ and $L$ disagree on their exact common
      deletion; or
\item the new disagreement between $R$ and $L$ has a witness edge among
      the five ambient edges deleted at $G$.
\end{enumerate}
\end{corollary}

\begin{proof}
Take the anchored trade supplied by
Corollary~\ref{cor:anchored-agreement-trade-residue}, and orient its lost
pair so that the repaired site is first; symmetry from
Lemma~\ref{lem:three-patch-common-core-localization} permits this.  In the
target state, $R$ agrees with $G$ and disagrees with $L$.  If $G$
and $L$ disagree, the first alternative holds.  Otherwise $G$ agrees
with $L$, and the chain $R\sim G\sim L$, together with
$R\not\sim L$, satisfies the three-patch localization lemma.  Its witness
is retained by the $R$- and $L$-deletions but deleted by the
$G$-deletion, hence lies in the five-edge footprint of $G$.
\end{proof}

The lift receipt contains more information than equality on the common core.
If the boundary horn does not occur, the lifted switch is supported entirely
inside that core.

\begin{lemma}[support confinement for a common-core lift]
\label{lem:common-core-lift-support-confinement}
\Sverified{Mettapedia.GraphTheory.FourColor.GoertzelV24KempeLiftConfinement.touchesBoundary_or_exists_liftedComponentSwitch_confined}
\Sverified{Mettapedia.GraphTheory.FourColor.GoertzelV24KempeLiftConfinement.taitBoundaryAlongLiftOrbit_or_exists_confinedLiftedColoring}
\Sverified{Mettapedia.GraphTheory.FourColor.GoertzelV24KempeLiftConfinement.firstTaitOrbitReachesSecondPair_or_exists_confinedLiftedColoring}
Let $D_A,D_B$ be two adjacent-pair deletions, and let a valid-pair Tait
Kempe sequence on their exact common deletion start at the restriction of a
colouring $A$ of $D_A$.  Either an intermediate component reaches the
vertices deleted by $D_B$, or the sequence lifts to a Tait colouring $A'$ of
$D_A$ which has the prescribed common-core restriction and satisfies
\[
 A'(e)=A(e)
 \quad\hbox{for every edge $e$ of $D_A$ outside the common-core image.}
\]
\end{lemma}

\begin{proof}
For one switch, suppose that the small two-colour component does not reach
the embedding boundary.  Its image is then closed under adjacency in the
large two-colour line graph.  Connectedness shows that this image is one
whole large component, rather than a proper subset of one.  Switching that
large component therefore changes exactly the image edges, gives the desired
small restriction, and fixes every edge outside the image.

Induct over the finite valid-pair Kempe sequence.  At each stage, a boundary
contact gives the first alternative.  Otherwise the preceding one-switch
argument lifts the next switch; the equality outside the common core is
preserved transitively.  Ordinary Kempe switching on a valid nonzero colour
pair preserves the Tait condition throughout.  Specializing the embedding to
the exact common deletion of $D_A$ and $D_B$ gives the stated adjacent-pair
form.
\end{proof}

\begin{lemma}[port-or-circuit classification for a confined lift]
\label{lem:confined-lift-port-or-circuit}
\Sverified{Mettapedia.GraphTheory.FourColor.GoertzelV24DegreeTwoBoundaryKempeLiftClassification.circuit_or_exists_kempeComponentMeetsPort}
\Sverified{Mettapedia.GraphTheory.FourColor.GoertzelV24DegreeTwoBoundaryKempeLiftClassification.taitBoundaryAlongLiftOrbit_or_exists_portOrCircuitConfinedLift}
\Sverified{Mettapedia.GraphTheory.FourColor.GoertzelV24DegreeTwoBoundaryKempeLiftClassification.firstTaitOrbitReachesSecondPair_or_exists_portOrCircuitConfinedLift}
Let $H$ be a finite graph with a finite ordered set of degree-two ports,
and suppose that every other vertex of $H$ has degree three.  In any Tait
colouring of $H$, every connected component selected by two distinct
nonzero colours is either a circuit or contains an edge incident with a
port.

Consequently, in the non-boundary horn of
Lemma~\ref{lem:common-core-lift-support-confinement}, every switch in the
lifted sequence on $D_A$ is either circuit-shaped or meets one of the four
ports exposed by deleting the adjacent pair $D_A$.  Thus the exact
three-way alternative is: a component reaches the second deleted pair; a
lifted component reaches the first deleted pair's boundary; or the lifted
component is a circuit.  In the latter two cases the whole sequence still
fixes every edge outside the exact common-core image.
\end{lemma}

\begin{proof}
Suppose a selected component $K$ is not a circuit.  Some selected edge
$e\in K$, viewed as a vertex of the induced two-colour line graph, then
has degree different from two.  If neither endpoint of $e$ were a port,
both endpoints would be cubic.  At either endpoint, properness and the Tait
condition give exactly one incident edge in the other selected colour.
The two resulting neighbours of $e$ are distinct because the graph is
simple.  They are all of its neighbours in the selected line graph, so
$e$ has degree two, a contradiction.  Hence $K$ meets a port.

For a finite sequence, induct over its switches.  The common-core lifting
lemma either detects contact with the second deletion, or identifies the
small component with one whole component of the first deletion and changes
no edge outside the common core.  Apply the preceding classification to
that lifted component and retain its port-or-circuit tag in the step
relation.  Valid-pair switching preserves the Tait condition, and the
outside-core equalities compose along the sequence.  Finally, the first
deletion of adjacent cubic vertices has exactly four degree-two ports and
is cubic elsewhere, giving the stated specialization.
\end{proof}

\begin{lemma}[the circuit horn carries a primal facial bond]
\label{lem:confined-lift-circuit-facial-bond}
\Sverified{Mettapedia.GraphTheory.FourColor.GoertzelV24CircuitKempeComponentPrimalCycle.isPrimalCoherentClosed_of_isCycle_of_bicolored}
\Sverified{Mettapedia.GraphTheory.FourColor.GoertzelV24CircuitKempeComponentPrimalCycle.exists_primalCycle_of_isCircuitKempeComponent}
\Sverified{Mettapedia.GraphTheory.FourColor.GoertzelV24AdjacentPairCircuitKempeFacialBond.exists_ambientPrimalCycle_faceSetBond_of_deleted_isCircuitKempeComponent}
Suppose that the circuit alternative of
Lemma~\ref{lem:confined-lift-port-or-circuit} occurs in an adjacent-pair
deletion of a graph-backed spherical cubic map whose orbit faces are
two-sided.  Then the selected line-graph component contains a simple cycle
whose primal lift is a simple cycle in the deletion graph.  Every edge of
that primal cycle belongs to the selected Kempe component.

Under the retained-vertex inclusion, this cycle remains simple in the
ambient graph.  Moreover, there is a set of orbit faces for which an ambient
edge lies on the cycle exactly when its two incident faces lie on opposite
sides of the set, and the selected and unselected face sets each induce a
connected subgraph of the facial dual.
\end{lemma}

\begin{proof}
The circuit condition says that the connected induced two-colour line graph
is finite and $2$-regular, so it contains a simple line-graph cycle.  At
three successive vertices of this cycle, the first and third primal edges
have the same one of the two selected colours: both differ from the middle
edge.  They therefore cannot meet at one primal vertex, by properness of the
edge colouring.  The same argument at the closing turn shows that the whole
line-graph cycle is primal-coherent.  Its canonical primal lift is a
nonbacktracking closed walk; extracting a simple cycle from its traced
subgraph retains only edges represented by vertices of the original
line-graph component.

The deletion graph is an induced subgraph on the retained vertices, so its
vertex inclusion into the ambient graph is injective and preserves the
simple cycle.  Finally apply the primal circuit--facial bond lemma: on a
two-sided spherical rotation system the binary edge vector of a simple
primal cycle is the boundary of a face set, and support-minimality of the
cycle forces both facial shores to be connected.
\end{proof}

\begin{lemma}[exact scope of common-core confinement]
\label{lem:common-core-confinement-scope}
\Sverified{Mettapedia.GraphTheory.FourColor.GoertzelV24KempeLiftConfinement.exists_firstCommonCore_preimage_iff_target_retained}
\Sverified{Mettapedia.GraphTheory.FourColor.GoertzelV24KempeLiftConfinement.exists_firstCommonCore_preimage_of_mem_ambientDisagreementSupport}
Let $D_R,D_G$ be two adjacent-pair deletions and let $e$ be an ambient
edge retained by $D_R$.  Then the copy of $e$ in $D_R$ lies in the
image of the exact common-deletion embedding if and only if $D_G$ also
retains $e$.  Consequently every edge in a pairwise ambient disagreement
support lies inside the corresponding common-core image.

Thus the word ``confined'' in
Lemma~\ref{lem:common-core-lift-support-confinement} is an exact support
statement but not a metric-locality statement.  The common core deletes only
the four vertices of the two adjacent pairs and can contain edges arbitrarily
far from both selected mesh intervals.  In particular, confinement alone
does not put the persistent carrier of
Corollary~\ref{cor:no-horn-persistent-disagreement-carrier} in a bounded
mesh neighbourhood.
\end{lemma}

\begin{proof}
If the retained copy of $e$ has a common-core preimage, that preimage is
retained by the second deletion, and mapping it back to the ambient graph
recovers $e$.  Conversely, if both deletions retain $e$, attach to its
two endpoints the proofs that they avoid all four deleted vertices.  This
gives an edge of the common induced graph whose first embedding is exactly
the retained copy of $e$.  A disagreement-support edge is retained by
both displayed patches by definition, so the final assertion follows.
\end{proof}

\begin{lemma}[pointwise disagreement transport under confinement]
\label{lem:confined-disagreement-support-transport}
\Sverified{Mettapedia.GraphTheory.FourColor.GoertzelV24KempeLiftConfinement.ambientDisagreementSupport_nonempty_iff}
\Sverified{Mettapedia.GraphTheory.FourColor.GoertzelV24KempeLiftConfinement.ambientDisagreementSupport_subset_of_confined_left_update}
Let $R,G,L$ be three adjacent-pair deletions.  Suppose a colouring
$R_{\rm before}$ of $R$ is changed to $R_{\rm after}$ only on the exact
common deletion of $R$ and $G$.  If
$R_{\rm after}\sim G$ and $R_{\rm before}\sim L$, then
\[
 \operatorname{Disagree}(R_{\rm after},L)
 \subseteq \operatorname{Disagree}(G,L),
\]
where each disagreement support is the set of ambient edges retained by the
two displayed deletions on which their displayed colourings differ.
\end{lemma}

\begin{proof}
Take an ambient edge $e$ in the left-hand support.  If $G$ retains $e$, then
$R_{\rm after}(e)=G(e)$, so the inequality
$R_{\rm after}(e)\ne L(e)$ is the same witnessed disagreement between $G$
and $L$.

If $G$ does not retain $e$, then $e$ lies outside the embedded common
deletion of $R$ and $G$.  Confinement gives
$R_{\rm after}(e)=R_{\rm before}(e)$, while the old agreement with $L$ gives
$R_{\rm before}(e)=L(e)$.  This contradicts the choice of $e$.  Hence the
second case is impossible, and every left-hand witness belongs to the
right-hand support.
\end{proof}

\begin{lemma}[the ambient disagreement carrier cannot grow]
\label{lem:confined-disagreement-union-antitone}
\Sverified{Mettapedia.GraphTheory.FourColor.GoertzelV24KempeLiftConfinement.ambientDisagreementSupport_subset_union_of_confined_left_update}
\Sverified{Mettapedia.GraphTheory.FourColor.GoertzelV24OrderedMeshConfinedRepairPropagation.nineSiteAmbientDisagreementUnion_antitone}
For a finite family of adjacent-pair deletion colourings, let $U(A)$ be the
union, over every ordered pair of sites, of their ambient disagreement
supports.  If $A'$ is obtained from $A$ by one confined repair of a site $R$
towards a gained partner $G$, then
\[
                         U(A')\subseteq U(A).
\]
More pointwise, for every third site $X$,
\[
 \operatorname{Disagree}(R_{\rm after},X)
 \subseteq
 \operatorname{Disagree}(R_{\rm before},X)
 \cup \operatorname{Disagree}(G,X).
\]
\end{lemma}

\begin{proof}
Let $e$ witness a disagreement between $R_{\rm after}$ and $X$.  If $G$
retains $e$, the new agreement $R_{\rm after}\sim G$ makes the same edge a
disagreement between $G$ and $X$.  If $G$ deletes $e$, confinement gives
$R_{\rm after}(e)=R_{\rm before}(e)$, so $e$ was already a disagreement
between $R$ and $X$.  This proves the pointwise inclusion.

Only pairs incident with $R$ can change in a one-coordinate repair.  Apply
the pointwise inclusion in the displayed orientation when $R$ is first and,
using symmetry of disagreement support, in the reversed orientation when
$R$ is second.  Every target witness therefore belonged to one of the source
pair supports, proving $U(A')\subseteq U(A)$.
\end{proof}

\begin{lemma}[a closed confined-repair circuit freezes the disagreement carrier]
\label{lem:closed-confined-repair-carrier-rigidity}
\Sverified{Mettapedia.GraphTheory.FourColor.GoertzelV24OrderedMeshConfinedRepairPropagation.reflTransGen_nineSiteAmbientDisagreementUnion_antitone}
\Sverified{Mettapedia.GraphTheory.FourColor.GoertzelV24OrderedMeshConfinedRepairPropagation.nineSiteAmbientDisagreementUnion_eq_of_step_in_closed_run}
\Sverified{Mettapedia.GraphTheory.FourColor.GoertzelV24OrderedMeshConfinedRepairPropagation.exists_carrierPreservingStep_of_confinedRepairCycle}
Suppose a confined-repair run starts and ends at the same nine-site state.
At every displayed step $A\longrightarrow A'$ of that closed run,

\[
                              U(A')=U(A).
\]

In particular, every nonempty confined-repair cycle has a first step which
preserves the ambient disagreement carrier exactly.
\end{lemma}

\begin{proof}
Write the closed run as a prefix from the base state $A_0$ to $A$, the
displayed step $A\longrightarrow A'$, and a suffix from $A'$ back to
$A_0$.  Iterating
Lemma~\ref{lem:confined-disagreement-union-antitone} along the prefix and
suffix gives
\[
                  U(A)\subseteq U(A_0)\subseteq U(A').
\]
The displayed step gives the reverse inclusion
$U(A')\subseteq U(A)$.  Hence equality holds.  Splitting a nonempty
transitive cycle at its first step supplies the final assertion.
\end{proof}

\begin{corollary}[the no-horn residue has a nonempty persistent carrier]
\label{cor:no-horn-persistent-disagreement-carrier}
\Sverified{Mettapedia.GraphTheory.FourColor.GoertzelV24OrderedMeshConfinedRepairPropagation.exists_nonempty_persistentDisagreementCarrier_of_no_branchingOrBoundary}
If no nine-site state exposes either the branching or the boundary-reaching
horn, then there is a confined strict repair $A\longrightarrow A'$ and an
ambient edge $e$ such that
\[
                    U(A')=U(A),
        \qquad e\in U(A)\cap U(A').
\]
Thus the recurrent alternative is not merely a repeated finite state: it
carries at least one physical edge on which the selected deletion colourings
continue to disagree.
\end{corollary}

\begin{proof}
The no-horn hypothesis makes the finite confined-repair relation serial, so
it has a nonempty directed cycle.  Choose the first carrier-preserving step
from Lemma~\ref{lem:closed-confined-repair-carrier-rigidity}.  By the
definition of a strict repair, the repaired site and its gained partner
disagree before this step.  The exact support characterization of
Lemma~\ref{lem:confined-disagreement-support-transport} therefore supplies
an ambient edge $e$ in their disagreement support, hence in $U(A)$.
Carrier equality places the same edge in $U(A')$.
\end{proof}

\begin{corollary}[forced partner-disagreement propagation]
\label{cor:forced-partner-disagreement-propagation}
\Sverified{Mettapedia.GraphTheory.FourColor.GoertzelV24KempeLiftConfinement.commonCoreAgrees_of_confined_left_update}
\Sverified{Mettapedia.GraphTheory.FourColor.GoertzelV24OrderedMeshConfinedRepairPropagation.exists_partnerDisagreement_of_no_branchingOrBoundary}
Under the no-branching and no-boundary hypothesis of
Corollary~\ref{cor:periodic-tait-repair-residue}, some confined strict repair
ends in a state with a repaired site $R$, its gained partner $G$, and a lost
partner $L$ such that the colourings at $G$ and $L$ disagree on their exact
common deletion.  More precisely, every ambient edge witnessing the new
$R$--$L$ disagreement also witnesses the $G$--$L$ disagreement:
\[
 \operatorname{Disagree}(R,L)\subseteq
 \operatorname{Disagree}(G,L).
\]
\end{corollary}

\begin{proof}
Use Lemma~\ref{lem:common-core-lift-support-confinement} in the repair
construction.  The no-boundary hypothesis selects its confined lift at every
step.  The same finite-state argument as in
Corollary~\ref{cor:periodic-tait-repair-residue} gives a nonempty cycle of
confined strict repairs, and the finite agreement set must lose an old
agreement somewhere on that cycle.  Thus there is again an anchored trade
$R,G,L$.

Suppose, for contradiction, that $G$ and $L$ agree after that trade.  Consider
any ambient edge retained by the $R$- and $L$-deletions.  If the $G$-deletion
also retains it, the new agreements $R\sim G$ and $G\sim L$ give equal
colours on the edge.  If the $G$-deletion removes it, support confinement says
that the repair did not change its colour at $R$.  The $L$ coordinate was not
changed by this one-coordinate repair, and $R\sim L$ held before the trade,
so the two colours are equal in this case as well.  Hence $R$ and $L$ still
agree on every edge of their exact common deletion, contradicting that this
was the lost agreement.  Therefore $G$ and $L$ disagree.
Equivalently, Lemma~\ref{lem:confined-disagreement-support-transport} sends
the nonempty new $R$--$L$ disagreement support into the $G$--$L$ support.
\end{proof}

\begin{proposition}[agreement-pivot dynamics alone can cycle]
\label{prop:agreement-pivot-cycle-countermodel}
\Sverified{Mettapedia.GraphTheory.FourColor.GoertzelV24AgreementPivotCycleCountermodel.confinedPivotStep_three_cycle}
\Sverified{Mettapedia.GraphTheory.FourColor.GoertzelV24AgreementPivotCycleCountermodel.not_exists_strictlyIncreasing_agreementRank}
There are three sites and three reflexive symmetric agreement states
$S_{02},S_{01},S_{12}$ with the following closed repair run:
\[
 S_{02}\xrightarrow{,0:2\mapsto 1,}S_{01}
 \xrightarrow{,1:0\mapsto 2,}S_{12}
 \xrightarrow{,2:1\mapsto 0,}S_{02}.
\]
At every arrow the repaired site first disagrees with the gained partner and
agrees with the lost partner; afterwards it agrees with the gained partner
and disagrees with the lost partner, the two partners disagree, and every
agreement not incident with the repaired site is unchanged.  Consequently
no natural-number-valued rank depending only on the agreement relation can
increase strictly at every confined pivot.
\end{proposition}

\begin{proof}
Let $S_{ij}$ consist of the diagonal pairs together with the two orientations
of the single off-diagonal pair $\{i,j\}$.  Thus $S_{02}$ has the lone
nontrivial agreement $0\sim2$.  Repair site $0$ by replacing that agreement
with $0\sim1$; the untouched pair $1,2$ disagrees before and after.  The same
description, cyclically relabelled, gives
$S_{01}\to S_{12}\to S_{02}$.  These are confined pivots with partner
disagreement exactly as stated.

If a rank $r$ increased on every such pivot, the closed run would give
$r(S_{02})<r(S_{01})<r(S_{12})<r(S_{02})$, an impossibility.
\end{proof}

Applied to the pair supplied by
Corollary~\ref{cor:ordered-row-common-core-disagreement}, this trichotomy is
the exact local residue.  It still does not give support equality or wall
exclusion: Corollary~\ref{cor:anchored-agreement-trade-residue} now proves
that repairing one pair must sometimes destroy another common-core agreement
at the repaired site.  Corollary~\ref{cor:localized-anchored-agreement-trade}
first localizes the only apparent alternative to propagation, while
Corollary~\ref{cor:forced-partner-disagreement-propagation} eliminates that
alternative using the support receipt discarded by the earlier lift theorem.
The defect must propagate to the two partner sites.  The remaining
two-dimensional task is therefore to show that repeated propagation reaches
the branching or boundary horn, or that its finite recurrent dynamics define
an exact support-preserving replacement.  Proposition~\ref{prop:agreement-pivot-cycle-countermodel}
shows that this step must use mesh geometry or information richer than the
abstract agreement graph: confinement and partner disagreement alone still
admit a periodic orbit.  The
arc-consistent no-section obstruction of
Corollary~\ref{cor:arc-consistent-no-section-residue} alone does not supply
that termination or replacement.  A coherent atlas is already impossible by
Lemma~\ref{lem:global-common-core-colouring-atlas}.  No monotonicity of the
naive number of disagreeing pairs is assumed; the verified trade explains
exactly how the periodic horn defeats that monotonicity.

This does not contradict the pentagonal-prism diagnostic above.  Four true
agreement bits attached to the four sides of one compatible cell state do
not assert that one globally selected family agrees on the common deletion
for every pair of mesh sites.  The corollary identifies the exact remaining
two-dimensional task: control the forced disagreements between globally
selected deletion colourings, rather than infer a global section from local
nonempty relations.

Exhausting all $2\,792\,976$ compatible cycles in the same prism separates
the two residues: $6\,912$ have empty composite and $6\,144$ have a
nonempty fixed-point-free composite; the remaining $2\,779\,920$ lift.
Thus both obstruction horns already occur in one planar bridgeless cubic
graph.  Moreover, each of the (6,144) proper composites is a singleton
partial transport whose one-element domain differs from its one-element
range.  Hence every proper prism example lies in the support-drift residue,
not the permutation residue.  The calculation is again external evidence,
not a kernel theorem.

There is nevertheless no colour-extension gap at a compatible checkpoint.
The following statement isolates the exact planar cap which may be attached
there; it will be the closed carrier on which Kauffman's formation parity can
legitimately be evaluated.

\begin{lemma}[compatible four-port colourings extend across the paired cap]
\label{lem:adjacent-pair-compatible-cap-extension}
\Sverified{Mettapedia.GraphTheory.FourColor.GoertzelV24AdjacentPairReductionColorExtension.reductionTaitColorable_iff_exists_deletedColoring_compatible}
Fix an adjacent pair (uv), delete (u,v), and retain the four defect ports
in their cyclic order.  Choose either of the two planar square pairings of
those ports.  The corresponding paired square reduction is Tait-colourable
if and only if the four-defect graph has a proper Tait colouring whose
missing-colour word is constant on each of the two chosen pairs.
\end{lemma}

\begin{proof}
In one direction, restrict a colouring of the paired reduction to its old
edges.  At either endpoint of a new seam, the two retained incident colours
are distinct and nonzero; their missing Klein-four colour is therefore the
colour of that seam.  The two endpoints of a seam consequently have equal
missing colours, giving the required compatible word.

Conversely, start with a compatible deletion colouring.  Restore (u,v)
temporarily by assigning the four old cut edges their four missing colours
and assigning (uv) the colour zero.  This one-zero chain is proper at every
retained vertex and has no zero on any retained or cut edge.  Compatibility
says exactly that the two old cut edges identified at each selected seam
have the same colour.  Hence their dart colouring descends through the two
identifications.  The zero edge (uv) lies in the discarded region, so the
descended colouring is nowhere zero and is a Tait colouring of the paired
reduction.
\end{proof}

The four-state square used by the source is narrower than an arbitrary open
Kempe path, and all of its intermediate words stay compatible with one fixed
cap.

\begin{lemma}[the same-side Kempe square extends through one paired cap]
\label{lem:same-side-kempe-square-cap-extension}
\Sverified{Mettapedia.GraphTheory.FourColor.GoertzelV24AdjacentPairReductionKempeSquare.sameSideKempeSquare_states_extend}
Let \(C\) be a deletion colouring compatible with the pairing
\((01\mid 23)\).  Suppose an \((a,b)\)-component \(K_{01}\) has exact boundary
support \(\{0,1\}\), while another \((a,b)\)-component \(K_{23}\) has exact
boundary support \(\{2,3\}\).  Switch neither, either, or both components.
Each of the resulting four proper deletion colourings extends to a proper
Tait colouring of the same \((01\mid 23)\) paired reduction.
\end{lemma}

\begin{proof}
The four boundary words are the original word, the result of interchanging
(a,b) at ports (0,1), the result of interchanging them at ports (2,3),
and the result of interchanging them at all four ports.  Equality within each
of the pairs \(\{0,1\}\) and \(\{2,3\}\) is preserved in all four cases.
Every corner therefore satisfies the compatibility hypothesis of
Lemma~\ref{lem:adjacent-pair-compatible-cap-extension}.  Applying its explicit
extension construction at each corner gives the four capped colourings.
\end{proof}

The source needs more than endpoint extension, but less than a universal
path-lifting theorem.  Its two moves have exact same-side support, and such a
move closes to one ordinary Kempe circuit when its matching seam is inserted.

\begin{lemma}[a same-side Kempe switch closes through the paired cap]
\label{lem:same-side-kempe-switch-cap-lift}
\Sverified{Mettapedia.GraphTheory.FourColor.GoertzelV24AdjacentPairReductionKempeSquare.exists_cappedKempeComponent_restricts_sameSideSwitch}
In the setting of
Lemma~\ref{lem:same-side-kempe-square-cap-extension}, let
\(\widehat C\) be the explicit extension of \(C\) across the
\((01\mid23)\) cap.  If \(C'\) is obtained from \(C\) by switching \(K_{01}\),
then some closed \((a,b)\)-Kempe component of \(\widehat C\) contains the
\((01)\)-seam, and switching that component restricts on every old edge exactly
to \(C'\).  The same statement holds for \(K_{23}\).
\end{lemma}

\begin{proof}
Consider \(K_{01}\).  At each of ports \(0,1\), the missing colour belongs to
\(\{a,b\}\).  Of the two old edges incident with that degree-two boundary
vertex, exactly one therefore has colour \(a\) or \(b\); exact boundary
support says that this edge lies in \(K_{01}\).  The new \((01)\)-seam has the
missing colour, so in the capped graph it joins the two open ends of
\(K_{01}\).  No other selected-colour edge is available at either endpoint.
Consequently \(K_{01}\) together with the seam is a connected 2-regular
\((a,b)\)-subgraph, hence one whole closed Kempe circuit.

Switching this circuit interchanges \(a,b\) on every old edge of \(K_{01}\)
and on the new seam, and fixes every other old edge.  Consequently its
restriction to the old graph is exactly \(C'\).  This is the required closed
realization of the open switch; no separately chosen extension of \(C'\) is
needed.  The proof for \(K_{23}\) is identical.
\end{proof}

This does not assert that every open Kempe path lifts through a fixed cap:
cross-side switches can pass through incompatible boundary words, and an
explicit planar example shows that closed formation parity can then change
after recapping.  The lemma instead covers exactly the two same-side switches
used in Lemma~\ref{lem:adjacent-pair-common-constant-orbit}.  Its remaining
checked content is the graph-level identification of the capped component:
old edges are labelled by their original open component, that label is
invariant across every selected cap adjacency, and the matching seam closes
the intended component without fusing it to another one.

\begin{corollary}[spherical parity consumer for the same-side lift]
\label{cor:same-side-cap-formation-parity}
\Sverified{Mettapedia.GraphTheory.FourColor.GoertzelV24AdjacentPairReductionSphericalKauffman.exists_cappedKempeSwitch_with_restriction_and_sphericalFormationParity}
Let the ambient map be a graph-backed least Tait counterexample.  Suppose the
retained deletion is connected, its four ports occur in their physical cyclic
order, and the canonical \(01\mid23\) paired reduction is bridge-free.  Then
the capped circuit in
Lemma~\ref{lem:same-side-kempe-switch-cap-lift} may be chosen so that its
closed switch both restricts to the required open switch and preserves
Kauffman's formation parity.  Equivalently, for the third colour
\(c=a+b\), the parity of the sum of the three bicoloured-circuit counts is
unchanged.
\end{corollary}

\begin{proof}
The two new seams join ports \(0\) to \(1\) and ports \(2\) to \(3\).
Neither joined pair can already be the endpoints of a retained edge: together
with the corresponding deleted vertex, such an edge would form a triangle,
whereas a least counterexample is triangle-free.  The two seam endpoint pairs
are distinct from one another, and the old endpoint pairs remain distinct
because the ambient graph is simple.  Thus the capped rotation system has an
endpoint-injective, canonical simple-graph presentation.

The retained connectivity and physical boundary order make the reduction
spherical, connected, and cyclic at every vertex; bridge-freeness then implies
that its two dart-sides lie on distinct face orbits.  Transport the cap
colouring, its selected bicoloured component, and the component switch through
the canonical graph presentation.  Kauffman's checked closed-spherical parity
theorem gives equality of formation parity before and after that switch.
Finally, the checked same-side component lift identifies the switched colour
of every retained edge with the prescribed open Kempe switch.
\end{proof}

No abstract cross-curve receipt is used here.  The parity theorem is applied
only after the open switch has been closed inside an actual bridge-free
spherical cubic map.  This is the precise, source-faithful role of Kauffman's
closed-curve argument; no parity assertion is being extended to an open
tangle.

There is a second, checked constraint on any such overlap transport.  It
shows why the transfer state may not remember only an affine action on the
four colours.

\begin{lemma}[singleton-collar parity is not affine colour transport]
\label{lem:singleton-collar-affine-parity}
Along a nonempty closed chain of admissible adjacent-pair reentries, choose
at each state the canonical colour frame determined by its two ordered
boundary colours.  Allow every reentry an arbitrary additive translation of
the four Tait colours.  The composite affine action around the closed chain
is even.  By contrast, the rooted support monodromy of a singleton intrinsic
collar is odd.  Consequently no change of four-state coordinates conjugates
the singleton geometric monodromy to that closed affine colour action.

\Sverified{Mettapedia.GraphTheory.FourColor.GoertzelV24AllFaceEscapeStateColorGauge.transfer_sign_eq_one_of_closed}
\Sverified{Mettapedia.GraphTheory.FourColor.GoertzelV24AllFaceIntrinsicCollarAffineColorMonodromy.targetRootedSignedMonodromy_ne_closedAffineColorWordConjugate_of_singletonPrimalCutProfile}
\end{lemma}

\begin{proof}
The linear part of a reentry is the change from the source canonical colour
frame to the target canonical colour frame.  These changes telescope along
a path.  On a closed path the initial and terminal frames agree, so the net
linear part is the identity.  Translation of the four-element colour group
is even: every nontrivial translation is the product of two disjoint
transpositions.  Hence the sign of an affine word is the sign of its net
linear part, namely (+1).

For a singleton collar the two sheets of the rooted support order are
exchanged at exactly one endpoint, so its geometric permutation has sign
(-1).  Sign is invariant under conjugacy.  An even affine colour action
therefore cannot be conjugate to this odd geometric action.
\end{proof}

Lemma~\ref{lem:singleton-collar-affine-parity} is not, by itself, a
wall-exclusion theorem.  It is a state-design theorem.  A valid simultaneous
reentry invariant must retain the geometric orientation parity in addition
to the affine colour gauge, or else separately prove that singleton collars
cannot occur.  Treating the colour action as the whole support transport
would make the overlap step false exactly in the dense-defect geometry now
under investigation.

There is a further orientation gauge that must be fixed before the odd local
transport can be regarded as a charge on a particular ambient edge.

\begin{lemma}[support reversal moves the negative singleton crossing]
\label{lem:singleton-collar-orientation-gauge}
Let a singleton intrinsic collar be represented by an oriented scalar support
cycle.  Reverse only that support-cycle orientation, leaving the remote
facial-dual cycle fixed.  The reversed walk is again a valid scalar support
state.  After both crossing orders are rooted at the distinguished target
edge, corresponding positions denote the same two physical crossing edges,
but the chirality at each position is toggled.  Consequently the unique
negative crossing for the reversed orientation is a different ambient edge
from the unique negative crossing for the original orientation.

\Sverified{Mettapedia.GraphTheory.FourColor.GoertzelV24AllFaceIntrinsicCollarSignedSupportGauge.negative_crossing_edge_depends_on_support_orientation}
\end{lemma}

\begin{proof}
Reversing a simple support cycle preserves its edge set, simplicity, target
edge, second certified crossing, and the nonvanishing support condition.  It
reverses the support crossing word and leaves the remote crossing word
unchanged.  In the singleton branch these words contain exactly two entries.
Rotating either support word to start at the target therefore gives the same
target-rooted list.

At a fixed physical crossing the facial-dual step, and hence its canonical
remote dart, is unchanged.  The dart contributed by the reversed support walk
is the reverse of the original support dart.  Thus agreement of the two darts
becomes disagreement and disagreement becomes agreement: both chiralities are
toggled.  Each orientation has exactly one negative position, so those two
negative positions cannot denote the same physical edge.
\end{proof}

The sign $(-1)$ of singleton monodromy is therefore orientation-gauge
invariant, but the ambient edge carrying its displayed transposition is not.
A global noncancellation argument must either choose support orientations
coherently across reentries or formulate its charge after quotienting this
gauge.  Merely multiplying edge-labelled local transpositions is not a
well-defined global observable.

Even after a fixed carrier has been supplied, local oddness retains only the
parity of the number of reentries.

\begin{lemma}[parity boundary for a closed intrinsic reentry word]
\label{lem:closed-intrinsic-reentry-parity-boundary}
Let a closed all-intrinsic reentry chain have $n>0$ steps.  Suppose that one
fixed finite carrier has been chosen and that the geometric transport of each
step is represented on that carrier by an odd permutation.  Then the sign of
the ordered product is
\[
  (-1)^n.
\]
If this product is identified, up to a change of coordinates, with the
canonical affine colour transport around the same closed chain, then $n$ is
even.  This conclusion is sharp at the purely algebraic level: two odd
transpositions can have identity product.

\Sverified{Mettapedia.GraphTheory.FourColor.GoertzelV24AllFaceIntrinsicCycleParityBoundary.FixedCarrierOddTransport.word_sign}
\Sverified{Mettapedia.GraphTheory.FourColor.GoertzelV24AllFaceIntrinsicCycleParityBoundary.FixedCarrierOddTransport.length_even_of_wordProduct_eq_closedGaugeConjugate}
\Sverified{Mettapedia.GraphTheory.FourColor.GoertzelV24AllFaceIntrinsicCycleParityBoundary.two_odd_transpositions_can_cancel}
\end{lemma}

\begin{proof}
Induct on the explicitly counted reentry word.  The empty product has sign
$(+1)=(-1)^0$.  Appending a step multiplies the geometric product by one
odd permutation, so the sign is multiplied by $-1$.  After $n$ steps it
is therefore $(-1)^n$.

The canonical affine colour word of a closed reentry chain has identity
linear part, and its translations are even; hence its total sign is $+1$.
Sign is invariant under a change of coordinates.  Identification of the two
products thus gives $(-1)^n=+1$, which is equivalent to $n$ being even.
For sharpness, take the transposition of the two Boolean states twice: each
factor is odd and their product is the identity.
\end{proof}

Thus a coherent fixed-carrier construction is necessary but not sufficient.
To turn the all-intrinsic residue into a contradiction one must additionally
prove odd length, or construct a stronger nonabelian or geometric charge that
cannot cancel on an even orbit.  Neither property is part of the present
closed-residue data.  This is the exact interface the high-width argument must
now supply; no parity premise is silently added.

The other terminal of the all-face machine already has the literal geometry
required by the replacement layer; it is not merely a descriptive cut
profile.

\begin{lemma}[a large all-face cut is a literal facial-dual noose]
\label{lem:all-face-large-cut-literal-noose}
\Sverified{Mettapedia.GraphTheory.FourColor.GoertzelV24AllFaceLargeRemoteNoose.exists_dualNooseSideAt_of_largeRemotePrimalCutAtDart}
Suppose the all-face machine reaches its large-remote-primal-cut outcome at a
directed edge.  For every anchor vertex \(a\), there is a simple closed walk
\(W\) in the actual facial dual, of length at least five, with crossed primal
edge set \(X\), such that the component \(C_a\) of \(G-X\) containing \(a\)
satisfies

\[
  \partial C_a=X.
\]

Both \(G[C_a]\) and \(G[V(G)\setminus C_a]\) are connected.  If

\[
  S_a:=\{e\in E(G):e\text{ has an endpoint in }C_a\},
\]

then its middle-set width obeys

\[
  |\operatorname{mid}(S_a)|\le |W|.
\]
Thus this outcome supplies the literal noose-side geometry consumed by the
physical replacement theorem.  It does not, by itself, bound \(|W|\), produce
a second nested noose, or prove equality of their phased states.
\end{lemma}

\begin{proof}
Unpack the large-cut certificate.  Its witness is already a simple closed
walk \(W\) in the literal facial dual, and its long-profile field gives
\(|W|\ge5\).  Let \(X\) be the primal edges crossed by successive dual edges
of \(W\), and choose the connected component of \(G-X\) containing the
specified anchor \(a\).  These data are exactly an oriented facial-dual
noose side; no new curve or separator is chosen.

Apply spherical facial-parity saturation to this simple dual cycle.  Every
edge of \(X\) crosses the cycle once and therefore has endpoints on opposite
sides.  Conversely, a primal edge not in \(X\) has zero intersection with
the cycle and cannot leave the chosen component.  Hence \(X\) is the exact
edge boundary of \(C_a\).  The same separation argument, applied to the two
regions of the sphere cut by \(W\), shows that the chosen component and its
vertex complement are both connected.

For the width estimate, a middle vertex of \(S_a\) lies outside \(C_a\): all
edges incident with a vertex of \(C_a\) belong to \(S_a\).  Choose an
\(S_a\)-edge incident with that middle vertex.  Exactness of the boundary
orients it uniquely from \(C_a\) to the middle vertex and therefore assigns a
crossing dart of \(W\).  Distinct middle vertices have distinct terminal
endpoints, so this assignment is injective.  A simple dual cycle crosses one
distinct primal edge at each step, giving
\(|\operatorname{mid}(S_a)|\le |X|=|W|\).
\end{proof}

The preceding local facts can be assembled without adding a geometric
supply hypothesis.  The result is useful chiefly because it states the
remaining global obligation without hiding either of its two branches.

\begin{theorem}[exact residue of the all-face formation machine]
\label{thm:all-face-cycle-residue}
\Sverified{Mettapedia.GraphTheory.FourColor.GoertzelV24AllFaceEscapeStateCycleResidue.exists_reachable_structural_or_largeRemotePrimalCut_or_closedResidue}
Start the exact all-face escape-state machine at any directed edge of a
graph-backed vertex-minimal Tait counterexample.  The machine reaches one of
the following:
\begin{enumerate}
\item a certified structural target;
\item a certified large remote primal cut;
\item a nonempty closed orbit containing a literal closure-recovery face
transfer; or
\item a nonempty closed orbit all of whose reentries are intrinsic singleton
collars with odd rooted signed-support monodromy.
\end{enumerate}
In the fourth case the sum of the exact ambient colour discrepancies is zero.
Moreover, for every assignment of additive colour translations to the
reentry steps, the resulting closed affine colour word has identity linear
holonomy and even sign.
\end{theorem}

\begin{proof}
The finite exact-state theorem supplies an initial state at the chosen dart
and either one of the first two stopping outcomes or a reachable nonempty
closed reentry orbit.  In the latter case, classify each reentry by the exact
intrinsic-singleton/recovery alternative.  If a recovery step occurs, split
the closed orbit immediately before and after that step; this retains its
source, target, localized loss, restoring gain, and the two ordinary reentry
paths around it, giving item~3.  Otherwise the same orbit is a chain entirely
of intrinsic odd signed steps, giving item~4.

For the two assertions in item~4, extend each escape colouring by zero over
the six omitted adjacent-pair edges.  A step discrepancy is the sum of its
source and target extensions.  In characteristic two every intermediate
extension appears twice, so a closed chain telescopes to zero.  Independently,
the linear colour transport of a step is the change between its endpoint
canonical colour frames.  These changes also telescope on a closed chain.
The net linear map is therefore the identity; arbitrary additive translations
of the four-element colour group are even, so the entire affine word is even.
\end{proof}

Theorem~\ref{thm:all-face-cycle-residue} is not wall exclusion.  The recovery
alternative already restarts the complete adjacent-pair fusion machine at a
different ambient edge, but no well-founded global quantity is yet known to
decrease under every such restart.  In the all-intrinsic alternative, local
oddness is also insufficient: an even number of odd local transports may
cancel while both telescoping laws remain true.  The source-faithful missing
statement is consequently one of the following equivalent kinds of global
control: a fixed ambient signed-support observable whose intrinsic charges do
not cancel, or a well-founded measure that strictly decreases on recovery
restarts and converts the residual intrinsic cycle into a structural target.

\begin{remark}[a local Stone--Wales neutralization fails]
\Srefuted{exhaustive radius-one 12-port support and boundary-Kempe-closure computation}
The most direct attempt to remove a pentagon--heptagon defect pair is not a
valid shortcut.  At the common twelve-port interface, the pure hexagonal
flower and the radius-one Stone--Wales flower have exact supports of sizes
\(16533\) and \(16854\), with intersection \(9222\); neither support contains
the other.  After exhaustive boundary-Kempe closure, \(4047\) hexagonal words
remain outside the Stone--Wales closure and \(3924\) Stone--Wales words remain
outside the hexagonal closure.  In particular the all-constant word is
realized by the hexagonal flower but is not in the Stone--Wales closure.
Thus neither patch is a monotone replacement for the other at this interface,
even after all noncrossing boundary Kempe switches.  This refutes the local
reverse-Stone--Wales move, not every possible larger-scale formation
transport.
\end{remark}

\begin{theorem}[mesh isoperimetry and the bounded-cut obstruction]
\label{thm:mesh-isoperimetry}
\Sverified{Mettapedia.GraphTheory.FourColor.GoertzelV24MeshIsoperimetry.Mesh.branchIn_card_le_or_branchOut_card_le}
Let \(G\) contain \(a\) edge-disjoint row paths and \(b\) edge-disjoint
column paths, with a designated branch vertex on every row--column pair;
arbitrary additional edges and vertices of \(G\) are allowed.  If a vertex
set \(S\) has at most \(k<\min(a,b)\) boundary edges, then at most \(k^2\)
designated branch positions lie on one side of the cut.

If
\[
 S_0\subset S_1\subset\cdots\subset S_m
\]
is a nested chain of such cuts and the set of designated branch positions
inside \(S_i\) strictly grows at every step, then
\[
                         m\le 2k^2+1.
\]
\Sverified{Mettapedia.GraphTheory.FourColor.GoertzelV24MeshIsoperimetry.Mesh.chain_length_le}
\end{theorem}

The hypothesis that the branch-position set strictly grows at every step is
load-bearing.  The theorem does not say that a graph containing a wall has no
long bounded-cut chain.  It says that at most \(2k^2+1\) slabs of such a
chain can acquire a new wall branch position.  A longer chain must contain
slabs which acquire none; they may run through subdivided wall edges or
through material attached to the wall.  Consequently a large wall does not
itself furnish the manuscript's long sequence of saturated cuts.  Conversely,
the mesh theorem does not prove wall exclusion.  It explains why the missing
implication must use zero Count and minimality rather than wall size alone.

For the converse direction used in the reduction, one must distinguish mesh
positions from actual vertices.  The abstract mesh above permits two
row--column positions to designate the same vertex, because injectivity is
irrelevant to the isoperimetric inequality.  A genuine subdivided wall has an
injective branch map.  On that exact strengthening the bounded-width descent
does exclude every sufficiently large wall.

\begin{proposition}[a raw width bound excludes a sufficiently large injective mesh]
\label{prop:raw-width-excludes-injective-mesh}
\Sverified{Mettapedia.GraphTheory.FourColor.GoertzelV24InjectiveMeshWidthExclusion.not_isVertexInjective_of_rawBranchDecompositionSupply}
Fix $w$, and suppose every graph-backed vertex-minimal zero-Tait-count map
has a complete rooted branch decomposition of width at most $w$.  Put
\[
 B(w)=4\left(2^{(6w+1)
       \sum_{j=0}^{w}j!\,2^{3^j}}-1\right)+6.
\]
Then no such map contains an $a$-by-$b$ mesh whose $ab$ designated
branch positions are pairwise distinct vertices whenever $B(w)<ab$.
\end{proposition}

\begin{proof}
The checked width-preserving connectedization turns the supplied raw tree
into the connected edge-leaf decomposition consumed by the literal-shore
descent.  Its two rooted subtrees and the exact Count state give

\[
                         |V(M)|\le B(w).
\]

On the other hand, injectivity of the branch map sending $(i,j)$ to
$v_{ij}$ embeds the $ab$-element set of mesh positions into $V(M)$, and
hence $ab\le |V(M)|$.  These inequalities contradict $B(w)<ab$.
\end{proof}

Thus the checked implication is

\[
 \text{uniform raw width bound}\quad\Longrightarrow\quad
 \text{exclusion of one sufficiently large injective wall mesh}.
\]

For arbitrary planar graphs, the reverse implication is the wall/grid
theorem: exclusion of one fixed wall bounds width.  For the exact
least-counterexample class here, Sections~\ref{sec:cotree-path-size}
and~\ref{sec:ordered-mesh-supply} now supply the corresponding raw
decomposition directly, for permissive and ordered meshes respectively.
They use the target's checked Count descent, so they are not a proof of the
general planar grid theorem.  Neither excludes a fixed wall directly from
zero Count.  Fixed-wall exclusion remains the route-specific high-width
task; its width-to-descent consumer is checked.

The numerical mismatch with the direct Count state is also exact.  A literal
width-\(k\) cut has \(3^k\) Tait boundary words, hence its unminimized
cumulative support carrier has \(2^{3^k}\) states.  For every \(k\geq1\),
\[
                         2k^2+1 < 2^{3^k}.
\]
Therefore any nested width-\(k\) sweep which acquires a new mesh branch
position at every step is strictly shorter than the raw support carrier and
cannot meet the cardinality premise of the raw support pigeonhole theorem.
Both the comparison and its direct negated-pigeonhole form are checked:
\Sverified{Mettapedia.GraphTheory.FourColor.GoertzelV24MeshSupportPigeonholeObstruction.chain_length_lt_card_normalizedSupportType}
\Sverified{Mettapedia.GraphTheory.FourColor.GoertzelV24MeshSupportPigeonholeObstruction.not_card_normalizedSupportType_le_chain_length}
This does not rule out an early accidental state collision or a much smaller
proved reachable quotient.  It rules out deducing the required repetition
from wall size and raw-state finiteness alone.  In particular, the source's
L1 conclusion that an internal hexagonal ladder exists is not yet the ambient
saturated serial decomposition consumed by pumping.

There is also a version which does not require a strict gain at every cut.
Suppose a nested width-\(k\) sweep begins with at most \(k^2\) mesh branch
positions below it and ends with at most \(k^2\) above it.  If every consecutive
slab acquires at most \(d\) branch positions, then the whole mesh has at most
\(2k^2+d\) branch positions.  Consequently, when \(ab>2k^2+d\), some single
slab must acquire more than \(d\) branch positions:
\Sverified{Mettapedia.GraphTheory.FourColor.GoertzelV24MeshSupportPigeonholeObstruction.exists_large_branch_slab_of_bounded_sweep}
\Sverified{Mettapedia.GraphTheory.FourColor.GoertzelV24MeshSupportPigeonholeObstruction.not_uniformly_local_bounded_sweep}
Thus a uniformly local one-dimensional cell decomposition cannot sweep from
one side of a sufficiently large two-dimensional mesh to the other.  This is
an obstruction to the literal L1-to-pumping bridge, not a wall-exclusion
theorem: a target-aware two-dimensional replacement principle could still
use zero Count and minimality in a way that the raw linear receipt does not.

The most immediate two-dimensional candidate is to iterate the positive
planar resolution which makes the square rung work.  That candidate already
fails on one hexagonal face, at the boundary-colour level.

\begin{proposition}[hexagonal planar-pairing obstruction]
\label{prop:hexagon-planar-pairing-obstruction}
\Srefuted{Mettapedia.GraphTheory.FourColor.GoertzelV24HexagonPlanarPairingObstruction.no_weighted_planarPairing_hexagon_count_identity}
Delete the six vertices of a facial hexagon, leaving six exterior ports in
cyclic order.  There is no fixed choice of nonnegative integer weights on the
five noncrossing perfect matchings of those ports for which, for every
nonzero boundary word, the number of proper colourings of the hexagon equals
the weighted number of colour-compatible matching caps.
\end{proposition}

\begin{proof}
Write the three Tait colours as (r,b,p), and take the cyclic boundary word

\[
                         (r,b,p,r,b,p).
\]

It has the proper internal hexagon colouring

\[
                         (p,r,b,p,r,b),
\]

where entry (i) colours the edge from vertex (i) to vertex (i+1).
Indeed the three colours at every vertex are (r,b,p) in some order.  This
extension is unique.  The five noncrossing matchings, written as pairs of
port indices, are

\[
\begin{split}
 &(01)(23)(45),\qquad (01)(25)(34),\qquad (03)(12)(45),\\
 &(05)(12)(34),\qquad (05)(14)(23).
\end{split}
\]

Every one contains a pair whose two entries in the displayed boundary word
have different colours.  Hence no planar matching cap is compatible: every
weighted cap count is zero, while the hexagon extension count is one.  Lean
checks the extension, its uniqueness, all five incompatibilities, and the
result for an arbitrary natural-valued weight function.  It also checks the
opposite sharpness witness ((r,r,b,r,r,b)): one planar cap accepts that word
although the hexagon itself does not.
\end{proof}

Thus the exact positive square identity does not extend by replacing each
hexagon with a sum of planar pairing caps.  This is the local shadow of the
same high-width issue: after digon and square normalization, a hexagonal web
retains boundary information which pairings alone do not express.  The
proposition does not rule out a richer planar web basis, a target-aware
Kempe operation, or a replacement involving a genuinely two-dimensional
region.  It rules out treating a wall as a collection of independent square
resolutions.

Nor can the pairing-cap proposal be rescued merely by allowing the chosen
pairing to depend on an unrestricted exterior.

\begin{proposition}[parity-restricted target-aware pairing reduction fails]
\label{prop:hexagon-target-aware-pairing-obstruction}
\Srefuted{Mettapedia.GraphTheory.FourColor.GoertzelV24HexagonPairingTargetAwareBoundary.no_parityRestricted_targetAware_planarPairing_reduction}
Let \(U\) be the set of genuine nonzero Tait words on the six cyclic ports,
let \(P\subseteq U\) be the support of the hexagonal cell, and let
\(R_0,\ldots,R_4\subseteq U\) be the supports of its five noncrossing pairing
caps.  Let
\[
                  A=\{w\in U:\textstyle\sum_{j=0}^{5}w_j=0\}
\]
be the parity-admissible language obeyed by every physical six-edge cut.  It
is false that
\[
  \forall E\subseteq A,\quad E\cap P=\varnothing
  \Longrightarrow
  \exists i\in\{0,\ldots,4\},\quad E\cap R_i=\varnothing.
\]
Thus no rule which chooses one of the five caps after seeing an arbitrary
parity-admissible zero exterior support can replace the hexagon while
preserving zero Count.
\end{proposition}

\begin{proof}
For each of the five pairings, there is a genuine nonzero boundary word which
is constant on every paired pair but has no proper extension across the
hexagon.  The five finite witnesses are checked simultaneously by
\Sverified{Mettapedia.GraphTheory.FourColor.GoertzelV24HexagonPairingTargetAwareBoundary.every_planarPairing_accepts_nonextension}.
Consequently \(R_i\nsubseteq P\) for every \(i\).  Every pairing word has
zero total Klein colour because each of its three colours occurs in an equal
pair, so \(R_i\subseteq A\).
The cut-parity assertion itself is not a planar-picture assumption: summing
the three incident colours at every vertex of any finite open cubic tangle
cancels every interior edge twice and leaves exactly the boundary sum.
\Sverified{Mettapedia.GraphTheory.FourColor.GoertzelV24PortTangleCutParity.taitSupport_subset_zeroSum}
\Sverified{Mettapedia.GraphTheory.FourColor.GoertzelV24HexagonPairingTargetAwareBoundary.no_planarPairingSupport_subset}
\Sverified{Mettapedia.GraphTheory.FourColor.GoertzelV24HexagonPairingTargetAwareBoundary.planarPairingSupport_subset_parityAdmissible}

Now take the parity-admissible zero exterior \(E=A\setminus P\).  It is
contained in \(A\) and disjoint from \(P\), while
\(R_i\subseteq A\) and \(R_i\nsubseteq P\) say exactly that \(E\) meets every
\(R_i\).  Hence no cap works against this one exterior.  Lean records both
the explicit parity adversary and the displayed quantified negation.
\Sverified{Mettapedia.GraphTheory.FourColor.GoertzelV24HexagonPairingTargetAwareBoundary.parityAdversary_meets_every_planarPairing}
\end{proof}

The obstruction now survives the first universal physical constraint, but
``parity-admissible'' is still weaker than ``realized by a planar cubic
exterior''.  We have not shown that \(A\setminus P\), or any equally strong
hitting support, is physically realized.  For comparison, the unrestricted
complement form and its physical-class specialization are also checked:
\Sverified{Mettapedia.GraphTheory.FourColor.GoertzelV24HexagonPairingTargetAwareBoundary.no_unrestricted_targetAware_planarPairing_reduction}
\Sverified{Mettapedia.GraphTheory.FourColor.GoertzelV24HexagonPairingTargetAwareBoundary.no_physical_targetAware_planarPairing_reduction_of_complement_mem}.
Therefore the remaining loophole is now exact: a route-native target-aware
reduction must prove and exploit a structural restriction on the support
languages of actual planar exteriors.  Moving the existential quantifier or
retaining an unrestricted profile does not provide that restriction.

The first natural strengthening is to remember the boundary trace of Kempe
components.  Fix a pair of colours \(\{a,b\}\).  For a boundary word \(w\),
call a port active when it has colour \(a\) or \(b\).  A \emph{boundary Kempe
witness} is a fixed-point-free matching of the active ports, noncrossing in
their cyclic order, such that interchanging \(a\) and \(b\) on any union of
matched pairs produces another word in the same support.  A support is
\emph{boundary-Kempe closed} when every one of its words has such a witness
for each of the three colour pairs.

\begin{proposition}[parity plus boundary Kempe closure is insufficient]
\label{prop:hexagon-boundary-kempe-obstruction}
Let \(E=A\setminus P\) be the parity adversary from
Proposition~\ref{prop:hexagon-target-aware-pairing-obstruction}.  Then \(E\)
is boundary-Kempe closed.  Nevertheless \(E\cap P=\varnothing\) and
\(E\cap R_i\ne\varnothing\) for every one of the five noncrossing pairing
caps.  Hence parity together with this abstract noncrossing Kempe-closure
condition still does not force a target-aware pairing-cap replacement.
\Sverified{Mettapedia.GraphTheory.FourColor.GoertzelV24HexagonPairingBoundaryKempeClosure.parityAdversary_boundaryKempeClosed}
\end{proposition}

\begin{proof}
For each genuine six-port word in \(E\) and each colour pair, the certificate
records one of the five Catalan noncrossing matchings of the active ports.  A
verified Boolean decoder checks fixed-point-free matching on the active set,
noncrossing, and closure under all \(2^6\) selected port sets; the last test
is conditional on the selected set being a union of complete matched pairs.
The flat certificate is checked independently for the red--blue,
red--purple, and blue--purple pairs.
The decoder and these three checks give one witness uniformly for every word
and colour pair.
\Sverified{Mettapedia.GraphTheory.FourColor.GoertzelV24HexagonPairingBoundaryKempeClosure.parityAdversary_certifiedPattern}
The remaining assertions about intersections are exactly the preceding
parity-adversary calculation.
\end{proof}

This proposition is deliberately one-sided.  The combinatorial part of the
physical bridge can be stated without a picture and is now verified.  Fix one
proper colouring of one literal port tangle.  For one colour pair, declare two
active darts adjacent when they are the two halves of an interior edge or lie
at the same vertex, and take the equivalence relation generated by these local
wires.  All three such relations are therefore derived from the same tangle
and the same colouring; they are not three independently selected boundary
matchings.

\begin{theorem}[common-web component switching]
\label{thm:port-tangle-common-kempe-web}
Let $T$ be a literal port tangle, let $C$ be a proper nonzero Tait colouring of
$T$, and fix a pair of colours.  Every connected class of the bichromatic dart
relation is closed across interior edges and across the two active darts at a
vertex.  More generally, if $D$ is any union of such classes, then swapping the
two colours on precisely the darts of $D$ gives another proper colouring of
the same tangle.  Its boundary word consequently lies in the same exact Tait
support as the word of $C$.
\Sverified{Mettapedia.GraphTheory.FourColor.GoertzelV24PortTangleCommonKempeWeb.component_kempeRegion}
\Sverified{Mettapedia.GraphTheory.FourColor.GoertzelV24PortTangleCommonKempeWeb.switchRegion_isProper}
\Sverified{Mettapedia.GraphTheory.FourColor.GoertzelV24PortTangleCommonKempeWeb.properBoundaryWord_switchRegion_mem_taitSupport}
\end{theorem}

\begin{proof}
The generated relation has two kinds of wire.  Across an interior-edge wire,
the original edge-colouring equation says that both darts have the same
colour; hence they are active together.  Across a vertex wire, both endpoints
are active by definition.  Equivalence closure therefore preserves activity,
and one equivalence class, or any locally saturated union of classes, is
closed under both wire types.

After the switch the two halves of every interior edge still agree.  At a
vertex, either all active darts in the selected component union are switched
or none are.  Every inactive dart is fixed, and its colour lies outside the
swapped pair.  Thus the incident nonzero colours remain distinct (in the
cubic case these are the familiar two active darts and one third-colour dart).
Restriction to the ports is exactly the corresponding boundary switch, so
the switched colouring itself witnesses membership of the new word in the
same exact support.
\end{proof}

\begin{theorem}[two-ended physical Kempe components]
\label{thm:port-tangle-kempe-two-ended}
Let $T$ be a finite cubic port tangle with a proper nonzero Tait colouring,
and fix a colour pair.  Every bichromatic component which contains a boundary
port contains exactly two boundary ports.
\Sverified{Mettapedia.GraphTheory.FourColor.GoertzelV24FinitePathEndpointCount.card_degreeOneVertices_eq_two}
\Sverified{Mettapedia.GraphTheory.FourColor.GoertzelV24PortTangleKempeEndpoints.activeWeb_degree_eq_one_iff_isPort}
\Sverified{Mettapedia.GraphTheory.FourColor.GoertzelV24PortTangleKempeEndpoints.card_componentPortDarts_eq_two}
\end{theorem}

\begin{proof}
Replace every active dart by a vertex of an auxiliary finite graph.  Join the
two active darts based at each cubic vertex, and join the two active darts
which are the halves of each interior edge.  Properness and cubicity give
exactly two active darts at every tangle vertex.  Hence an active boundary
dart has degree one in the auxiliary graph, while an active interior dart has
degree two: it has one neighbour across its tangle vertex and one across its
interior edge.  These two neighbours are distinct because an interior edge is
not a loop.

Now restrict to the connected component containing the chosen boundary dart.
Restriction does not change degrees, since every neighbour of a vertex in a
connected component lies in the same component.  The component is finite and
connected, all its degrees are one or two, and it has at least one degree-one
vertex.  The handshaking lemma makes the number of degree-one vertices even;
connectedness gives at most two of them.  Thus there are exactly two, and the
degree characterization identifies them exactly with the boundary ports.
\end{proof}

\begin{lemma}[boundary-free physical Kempe components are circuits]
\label{lem:port-tangle-kempe-boundary-free-cycle}
\Sverified{Mettapedia.GraphTheory.FourColor.GoertzelV24PortTangleKempeEndpoints.activeComponent_isCycles_of_componentPortDarts_eq_empty}
\Sverified{Mettapedia.GraphTheory.FourColor.GoertzelV24PortTangleKempeEndpoints.exists_activeComponent_cycle_of_componentPortDarts_eq_empty}
Let $T$ be a finite cubic port tangle with a proper nonzero Tait colouring,
fix a colour pair, and choose one active dart.  If the physical bichromatic
component of that dart contains no boundary port, then its active-web graph
is a cycle.  More precisely, after restricting the active web to that
connected component, every vertex has degree two, and the component contains
a simple closed walk through the chosen dart.
\end{lemma}

\begin{proof}
The active web has one vertex for every dart carrying one of the two selected
colours.  At a tangle vertex, proper cubicity gives each such dart exactly one
other selected dart.  An interior dart has one further neighbour across its
interior edge, and these two neighbours are distinct because the tangle has
no interior loop.  A boundary dart has no across-edge neighbour and therefore
has degree one.  Passing to a connected component removes no neighbour.

By hypothesis the chosen component has no boundary dart, so every one of its
vertices is an interior dart and has degree two.  Thus the restricted finite
graph is connected and $2$-regular.  Equivalently it is a graph of cycles;
the cycle-extraction lemma applied at the chosen dart supplies a simple closed
walk.  This is only an algebraic circuit in the physical dart web.  Turning
its contraction into a separating simple cycle of the ambient spherical map
is a distinct geometric step and is not asserted here.
\end{proof}

\begin{lemma}[a boundary-free graph-backed Kempe circuit is a primal cycle]
\label{lem:graph-backed-kempe-boundary-free-primal-cycle}
\Sverified{Mettapedia.GraphTheory.FourColor.GoertzelV24PortTanglePrimalKempeGraph.primalKempeComponent_isCycles_of_boundaryFree}
\Sverified{Mettapedia.GraphTheory.FourColor.GoertzelV24PortTanglePrimalKempeGraph.exists_primalKempeCycle_of_boundaryFree}
\Sverified{Mettapedia.GraphTheory.FourColor.GoertzelV24VertexSidePrimalKempeTrail.interiorEndpointInjective_vertexSide}
\Sverified{Mettapedia.GraphTheory.FourColor.GoertzelV24VertexSidePrimalKempeTrail.exists_ambientPrimalKempeCycle_of_boundaryFree}
Let a finite simple graph with a dart rotation be cut along a vertex shore,
and give the resulting literal port tangle a proper nonzero Tait colouring.
Fix a colour pair and an active dart whose physical component contains no
boundary port.  Contract every same-vertex wire in that component.  The
result is a simple cycle in the selected two-colour primal subgraph, and its
inclusion in the ambient simple graph is again a simple cycle.
\end{lemma}

\begin{proof}
Consider a primal vertex reached from the chosen active dart.  Any active dart
based there belongs to the same physical component: primal reachability lifts
to the common dart web, and two active darts at one vertex are joined by the
same-vertex wire.  Hence no active boundary port is based there.  Proper
cubicity supplies exactly two active darts at the vertex, so both are interior
darts.  Crossing their two interior edges gives two neighbours in the selected
primal graph.

These neighbours are distinct.  Indeed, in the graph-backed tangle an
interior dart remembers an oriented edge of the ambient simple graph.  Two
interior darts with the same initial and terminal vertices are equal.  Thus
equality of the two neighbours would identify the two active darts,
contradicting their distinctness.  Conversely every selected primal edge at
the vertex arises from one of the two active darts.  The selected primal
component is therefore connected and $2$-regular, so it contains a simple
cycle through the chosen vertex and in fact is exactly that cycle.

The inclusion of the selected primal graph into the retained shore, and then
into the ambient graph, is injective on vertices and preserves adjacency.
Mapping the cycle along this inclusion therefore preserves the simple-cycle
property.  This proves the claim.  Spherical separation of the resulting
ambient cycle is the next topological consumer, not part of this lemma.
\end{proof}

\begin{lemma}[exact face-side certificate for the boundary-free primal cycle]
\label{lem:graph-backed-kempe-boundary-free-face-separator}
\Sverified{Mettapedia.GraphTheory.FourColor.GoertzelV24VertexSidePrimalKempeFaceSeparator.exists_faceSet_separates_iff_mem_walk_of_isCycle}
\Sverified{Mettapedia.GraphTheory.FourColor.GoertzelV24VertexSidePrimalKempeFaceSeparator.exists_ambientPrimalKempeCycle_faceSetSeparator_of_boundaryFree}
Assume in addition that the ambient dart rotation is a connected bridgeless
spherical cubic map and that every primal edge has two distinct incident
orbit faces.  For the ambient simple cycle supplied by
Lemma~\ref{lem:graph-backed-kempe-boundary-free-primal-cycle}, there is a
finite set of orbit faces such that an ambient edge occurs in the cycle if
and only if exactly one of its two incident faces belongs to that set.
\end{lemma}

\begin{proof}
Count the traversals of each ambient edge by the simple closed walk modulo
two.  Closedness makes the incidence boundary of this binary edge vector
zero, so it lies in the primal binary cycle space.  On a connected spherical
map, the orbit-face boundary vectors span that cycle space.  Therefore the
edge vector is the boundary of a binary coefficient on the orbit faces.

Every binary coefficient is the characteristic function of the finite set
of faces on which it is one.  Two-sidedness says that its boundary has value
one on an edge exactly when the two incident faces have different
coefficients.  Finally, the walk is a trail, so its modulo-two value on an
edge is one exactly when that edge occurs in the cycle.  Combining the two
equivalences gives the stated face-side certificate.

This is the exact combinatorial boundary identity needed from Jordan
separation.  It does not yet assert that either class of faces is connected,
nor does it construct a geometric noose; those are separate consumers.
\end{proof}

\begin{lemma}[a simple primal cycle is a bond of the facial dual]
\label{lem:graph-backed-kempe-face-sides-connected}
\Sverified{Mettapedia.GraphTheory.FourColor.GoertzelV24CycleSupportMinimality.f2CycleSpace_support_eq_cycle_of_subset}
\Sverified{Mettapedia.GraphTheory.FourColor.GoertzelV24PrimalCycleFacialBond.exists_faceSet_bond_of_isCycle}
\Sverified{Mettapedia.GraphTheory.FourColor.GoertzelV24PrimalCycleFacialBond.exists_ambientPrimalKempeCycle_faceSetBond_of_boundaryFree}
Under the hypotheses of
Lemma~\ref{lem:graph-backed-kempe-boundary-free-face-separator}, the face set
may be chosen so that the subgraph of the full facial dual induced by that
set is connected, and the subgraph induced by its complement is connected.
Both face classes are nonempty.  Thus the ambient Kempe cycle is not merely
some facial-dual cut: it is a bond of the facial dual.
\end{lemma}

\begin{proof}
We use two elementary minimality facts.  First, the edge set of a simple
cycle is a minimal nonempty support in the binary cycle space.  Indeed, let
$z$ be a nonzero binary cycle whose support is contained in a simple cycle
$C$, and choose a selected edge of $z$.  At either endpoint, the incidence
boundary of $z$ is zero.  Exactly two edges of $C$ meet that vertex, so the
other one must also be selected.  Propagating this implication around $C$
shows that every edge of $C$ is selected.  Hence
\(\operatorname{supp}z=C\).

Second, let $S$ be the face set supplied by
Lemma~\ref{lem:graph-backed-kempe-boundary-free-face-separator}.  An edge of
the full facial dual crosses from $S$ to its complement exactly when its
underlying primal edge belongs to $C$.  Since $C$ is nonempty and every
primal edge has two distinct incident faces, both $S$ and its complement are
nonempty.

Suppose that the facial-dual subgraph induced by $S$ were disconnected, and
let $K$ be one of its connected components.  The full facial dual is
connected, so the edge boundary of the nonempty proper set $K$ is nonempty.
No dual edge joins $K$ to another component of the subgraph induced by $S$.
Consequently every edge leaving $K$ leaves $S$ as well, and its underlying
primal edge belongs to $C$.  The binary boundary of $K$ is a sum of facial
boundary vectors, hence is a nonzero element of the primal binary cycle
space supported inside $C$.

There is another component $K'$ of the subgraph induced by $S$.  Its edge
boundary is likewise nonempty and supported inside $C$.  Moreover the
boundaries of $K$ and $K'$ are disjoint: an edge belonging to both would join
the two components inside $S$.  Thus the boundary of $K$ omits at least one
edge of $C$.  We have obtained a nonzero cycle-space support properly
contained in the simple cycle $C$, contradicting the first paragraph.
Therefore the subgraph induced by $S$ is connected.  Applying the same
argument to the complementary face set proves that its induced subgraph is
connected as well.

This is the finite combinatorial circuit--bond form of Jordan separation.
It proves the two connected face shores without introducing a topological
embedding primitive.  Turning the cyclic dart order along their common bond
into a geometric noose remains a separate representation step.
\end{proof}

\begin{theorem}[the physical common web induces the exact boundary matching]
\label{thm:port-tangle-physical-kempe-matching}
Let $T$ be a finite cubic port tangle with a proper nonzero Tait colouring,
and fix a colour pair.  Pair each active boundary port with the unique other
active boundary port in its bichromatic component, and fix each inactive
port.  This defines an involution which preserves the active set and has no
active fixed point.  Moreover, for every union of its active pairs, switching
the selected boundary coordinates produces another word in the exact Tait
support of $T$.

For a six-port tangle, if this physical matching is noncrossing for every
proper colouring and every colour pair, then the exact Tait support of $T$ is
boundary-Kempe-closed.  Noncrossing is the sole hypothesis in this last
implication; the matching and every required switch witness are supplied by
the one physical common web.
\Sverified{Mettapedia.GraphTheory.FourColor.GoertzelV24PortTanglePhysicalKempeClosure.activeWeb_reachable_iff_kempeConnected}
\Sverified{Mettapedia.GraphTheory.FourColor.GoertzelV24PortTanglePhysicalKempeClosure.card_physicalComponentPort_eq_two}
\Sverified{Mettapedia.GraphTheory.FourColor.GoertzelV24PortTanglePhysicalKempeClosure.physicalMate_isActiveMatching}
\Sverified{Mettapedia.GraphTheory.FourColor.GoertzelV24PortTanglePhysicalKempeClosure.swapBoundaryWord_mem_taitSupport_of_isComponentUnion}
\Sverified{Mettapedia.GraphTheory.FourColor.GoertzelV24PortTanglePhysicalKempeClosure.boundaryKempeClosed_taitSupport_of_physical_noncrossing}
\end{theorem}

\begin{proof}
Reachability in the auxiliary active-dart graph is exactly the equivalence
relation generated by the two local wire types used in the common-web
construction.  The preceding theorem therefore says that a component meeting
the boundary has exactly two boundary ports.  The other endpoint is unique;
using it on active ports and the identity on inactive ports gives the claimed
involution.  Equality of components is symmetric, so applying the operation
twice returns to the original port, and uniqueness excludes an active fixed
point.

If a set of boundary ports is a union of complete pairs, take the union of
the corresponding physical components.  This is a Kempe region in the sense
of Theorem~\ref{thm:port-tangle-common-kempe-web}.  Switching that region
therefore gives a proper colouring of the same tangle.  On an active selected
port its boundary colour is swapped; on an unselected port it is unchanged;
and on a selected but inactive port the colour is outside the chosen pair, so
the formal swap is also the identity.  Hence the new boundary word is exactly
the requested coordinate switch and belongs to the exact support.

At six ports, adding noncrossing supplies the one remaining field of the
finite boundary-Kempe witness.  Quantifying this construction over every
proper colouring and colour pair proves boundary Kempe closure of the whole
exact support.
\end{proof}

The missing geometric bridge has the following precise form.  We state it at
the level of the spherical combinatorial map, rather than adding an
uninterpreted ``planar matching'' field to the port-tangle carrier.

\begin{theorem}[noncrossing of physical Kempe components]
\label{thm:physical-kempe-components-noncrossing}
Let $M$ be a finite connected spherical cubic combinatorial map with
two-sided faces and connected face dual.  Let $S$ and its complement be the
two connected shores of an edge bond, and suppose the boundary darts based
in $S$ are enumerated in the cyclic order given by retained facial first
return.  Give the literal open tangle on $S$ a proper nonzero Tait
colouring, and fix a pair of colours.  Let one bichromatic component join
boundary ports $a$ and $c$, and let a second join $b$ and $d$.  Suppose
$b$ lies on the forward cyclic arc from the successor of $a$ to $c$, while
the forward cyclic arc from $d$ returns to $a$, neither open arc meeting
$a$ or $c$.  If $b$ and $d$ are distinct boundary ports, the two components
cannot be distinct.  Thus physical components cannot join alternating pairs
of boundary ports, even when the second component consists only of the local
vertex wire between two ports based at one retained cubic vertex.

In particular, when the bond has exactly six boundary darts, the physical
mate involution in cyclic coordinates is a noncrossing matching.  Hence the
exact ordered Tait support of the literal shore tangle is boundary-Kempe
closed: its matching and all switch witnesses are supplied by the physical
bichromatic components of the realizing colouring, not by an independently
chosen boundary certificate.
\Sverified{Mettapedia.GraphTheory.FourColor.GoertzelV24PhysicalKempeNoncrossing.no_distinct_kempe_components_between_boundary_arcs}
\Sverified{Mettapedia.GraphTheory.FourColor.GoertzelV24PhysicalKempeNoncrossingClosure.orderedPhysicalMate_isNoncrossing}
\Sverified{Mettapedia.GraphTheory.FourColor.GoertzelV24PhysicalKempeNoncrossingClosure.boundaryKempeClosed_orderedVertexSideTaitSupport}
\end{theorem}

\begin{proof}
It is enough to rule out two component trails with endpoints occurring in
cyclic order $a,b,c,d$, the first joining $a$ to $c$ and the second joining
$b$ to $d$.  Because the complementary shore is connected, join the outside
end of the cut edge at $c$ to the outside end of the cut edge at $a$ by a
trail in the complementary shore.  Appending the two cut edges and the
$a$--$c$ trail in $S$ gives a closed trail $Q$.  The four pieces have
pairwise disjoint edge sets: the two shore trails use no cut edge, the shores
are disjoint, and the two named cut edges are distinct.  Thus $Q$ is a
closed trail.  Among boundary edges it uses exactly those at $a$ and $c$.

Label every face by the parity of the number of edges of $Q$ crossed by a
dual walk from a fixed face.  Sphericity, two-sidedness, and connectedness of
the face dual make this label well defined; an edge has differently labelled
incident faces exactly when it belongs to $Q$.  At the boundary edge $a$
the label therefore changes.  Consecutive boundary darts in retained facial
first-return order expose the two sides of one ambient face.  Since all
boundary edges strictly between $a$ and $b$ avoid $Q$, the label propagates
unchanged from the face immediately after $a$ to the face at $b$.  Likewise,
along the other boundary arc from $d$ back to $a$, it propagates unchanged
from the face at $d$ to the face immediately before $a$.  The faces incident
with the boundary endpoints $b$ and $d$ consequently have opposite labels.

If the retained endpoints of $b$ and $d$ are distinct, their component trail
is nonempty and edge-disjoint from $Q$.  At either endpoint, cubic cyclic
rotation identifies the boundary-side face label with the label beside the
first or last trail dart; along the trail the label cannot change, because
none of its edges lies in $Q$.  Hence its endpoint faces have equal labels.

If the two retained endpoints instead coincide, the two boundary darts are
still distinct.  Their two boundary edges avoid $Q$.  At a cubic vertex any
two distinct incident darts are consecutive in one of the two rotation
directions.  Exact face-cut parity across the two avoided edges therefore
gives the same label to the two boundary-side faces.  In either case this
contradicts the opposite labels forced by the two boundary arcs.  Distinct
bichromatic components are vertex-disjoint and therefore edge-disjoint, while
Theorem~\ref{thm:port-tangle-kempe-two-ended} gives exactly two boundary
endpoints for every component meeting the boundary.  This proves the stated
unrestricted noncrossing conclusion for each physical colour pair.

For six boundary positions, assume two chords of the physical mate cross in
the standard cyclic order.  Orient each chord from its smaller to its larger
coordinate.  The four endpoints then interleave.  Since there are only six
coordinates, the two complementary open arcs are finite successor lists
avoiding the endpoints of the first chord.  The theorem just proved applies
to those two physical components and says that they cannot be distinct.  But
if they were the same component, uniqueness of the other endpoint would make
the two chords equal, contrary to interleaving.  Thus the physical mate is
noncrossing.  The common-web switching construction of
Theorem~\ref{thm:port-tangle-physical-kempe-matching}, transported through the
cyclic enumeration, now supplies every field of boundary Kempe closure for
the exact ordered support.
\end{proof}

The original v23 CAP5 move-realizability proof tries to pass from this
per-pair fact to a joint planar picture by contracting every edge of one
colour.  The contraction is planar, but the asserted cross-family
disjointness does not follow.

\begin{proposition}[zero-contraction does not separate different colour-pair systems]
\label{prop:v23-zero-contraction-crossing}
\Srefuted{Mettapedia.GraphTheory.FourColor.GoertzelV24V23ZeroContractionCrossing.exists_spherical_tait_zero_contraction_with_crossing_pair_systems}
There is a spherical properly Tait-coloured cubic combinatorial map and a red
colour class whose complete contraction produces a degree-four vertex at
which the surviving blue and purple through-pairs have alternating endpoints
in the vertex link.
Consequently the two post-contraction bichromatic systems are edge-disjoint
but are not mutually disjoint as embedded arc systems away from their
terminals.
\end{proposition}

\begin{proof}
Take the tetrahedron on vertices $A,B,C,D$, with rotations

\[
 (AB,AC,AD),\quad (BA,BD,BC),\quad
 (CA,CB,CD),\quad (DA,DC,DB).
\]

Give $AB,CD$ colour red, $AC,BD$ colour blue, and $AD,BC$
colour purple.  Each vertex sees the three colours, and the displayed
rotation system has four facial orbits, so $4-6+4=2$.

Contract both red edges $AB$ and $CD$.  Deleting their four darts and
taking first return in the vertex rotation joins the two punctured rotations
at the image of $AB$ into

\[
                         (AC,AD,BD,BC).
\]

The corresponding colour word in the link of the contracted vertex is

\[
                    ({\rm blue},{\rm purple},{\rm blue},{\rm purple}).
\]

The red--blue component through $AB$ therefore joins link positions $0$
and $2$, whereas the red--purple component joins positions $1$ and $3$.
Those two chords alternate.  They meet at the contracted interior point even
though the surviving blue and purple edge sets are disjoint.  Lean checks the
source Tait colouring and spherical Euler count, identifies the contracted
dart carrier with the source darts minus the four red darts, verifies the
first-return rotation law at both contracted vertices and the spherical count
$2-4+4=2$, and checks the alternating chords.  Thus the
failure is caused by contraction itself, not by a nonplanar model.
\end{proof}

This proposition does not refute per-pair noncrossing, nor by itself refute
the exceptional-CAP5 conclusion.  It refutes the particular v23 step which
treats the two different post-contraction systems as mutually disjoint and
then forms a Jordan curve from arcs belonging to both systems.  A repair must
retain their common contracted vertices (equivalently, the common trivalent
web before contraction) and prove a genuinely joint statement.

The exact local information lost by the disjoint-arc picture is one chirality
bit per contracted edge.

\begin{theorem}[local zero-contraction chirality law]
\label{thm:zero-contraction-chirality}
\Sverified{Mettapedia.GraphTheory.FourColor.GoertzelV24ZeroContractionChirality.sameColor_pairing_and_crossing_iff_endpointOrdersAgree}
\Sverified{Mettapedia.GraphTheory.FourColor.GoertzelV24ZeroContractionChirality.endpointOrdersAgree_alpha_iff}
Let \(M\) be a cubic combinatorial map with cyclic vertex rotations, let
\(C\) be a proper Tait edge-colouring, and orient an edge by a dart \(d\).
After contracting that edge, write the four surviving darts in the cyclic
link order
\[
  \rho d,\quad \rho^2d,\quad \rho\alpha d,\quad
  \rho^2\alpha d.
\]
The two same-colour pairs have exactly one of the following forms:
\[
  (0,2),(1,3)\qquad\hbox{or}\qquad (0,3),(1,2).
\]
The first, alternating, pairing occurs if and only if
\[
 C(\rho d)=C(\rho\alpha d),
\]
that is, if and only if the two endpoint rotations present the surviving
colours in the same order.  The second, adjacent, pairing occurs when the
orders are reversed.  This agreement predicate is unchanged when \(d\) is
replaced by \(\alpha d\).
\end{theorem}

\begin{proof}
At the endpoint of \(d\), cubicity and cyclicity say that
\(d,\rho d,\rho^2d\) are the three distinct incident darts.  Properness and
the nonzero Tait condition therefore put the three nonzero colours on their
three underlying edges exactly once.  The same is true of
\(\alpha d,\rho\alpha d,\rho^2\alpha d\) at the other endpoint, and the first
edges in those two triples have the same colour because \(d\) and
\(\alpha d\) are the two darts of one edge.

Consequently the ordered pair of surviving colours at the second endpoint
is either the same as the ordered pair at the first endpoint or its reverse;
there is no third possibility.  In the same-order case the four-link colour
word is \((x,y,x,y)\), so the equal-colour chords are \((0,2)\) and
\((1,3)\), whose endpoints alternate.  In the reversed case it is
\((x,y,y,x)\), so the chords are \((0,3)\) and \((1,2)\), both adjacent.
Changing the orientation of the contracted edge exchanges the two endpoint
blocks, so it does not change whether their orders agree.  Lean checks the
three-dart rotation facts, proper-colouring inequalities, the exhaustive
three-colour dichotomy, the same-colour mate at all four positions, and the
four-point chord crossing test.
\end{proof}

Thus the v23 contraction can be repaired locally, but not by deleting the
contracted vertices from the topology: a joint state must retain this
chirality at every contraction site and its compatibility along the common
web.  The theorem does not yet supply that global compatibility law.

The reduct `PortTangle` deliberately forgets the rotation and disk, so this
conclusion cannot follow at that carrier level alone.  More importantly, even
the full noncrossing bridge does not rescue the reduction, because the
adversary already passes the resulting abstract test.  We also have not
realized the whole adversary simultaneously as the exact support of one
planar cubic exterior.  The surviving loophole is therefore stronger and
genuinely joint: the three colour-pair matchings of an actual exterior arise
from one and the same embedded tangle.  A successful replacement theorem
must exploit that simultaneous physical compatibility (or some still
stronger feature of the exact exterior language), rather than parity or
three separately existential boundary matchings.

The first joint law is now checked at its proper, edge-set level.

\begin{theorem}[Kempe idemposition identity]
\label{thm:kempe-idemposition-identity}
Let $C:E\to\{0,r,b,p\}$ be any edge colouring and let $D\subseteq E$
contain only red or blue edges.  Switch red and blue exactly on $D$.  If
$M_{xy}(C)$ denotes the set of edges whose colours lie in
$\{x,y\}$, then
\[
 \begin{aligned}
 M_{rb}(C')&=M_{rb}(C),\\
 M_{rp}(C')&=M_{rp}(C)\mathbin\triangle D,\\
 M_{bp}(C')&=M_{bp}(C)\mathbin\triangle D.
 \end{aligned}
\]
\Sverified{Mettapedia.GraphTheory.FourColor.GoertzelV24KempeIdemposition.redBlue_idemposition}
The same statement is checked for arbitrary three pairwise distinct Klein
colours.
\Sverified{Mettapedia.GraphTheory.FourColor.GoertzelV24KempeIdemposition.twoColorSupport_switch_left_companion}
\Sverified{Mettapedia.GraphTheory.FourColor.GoertzelV24KempeIdemposition.twoColorSupport_switch_right_companion}
\end{theorem}

\begin{proof}
Outside $D$ no colour changes.  On $D$, every red edge becomes blue and
every blue edge becomes red.  Thus membership in $M_{rb}$ is unchanged,
whereas membership in each of $M_{rp}$ and $M_{bp}$ is toggled exactly
once.  Pointwise membership is therefore precisely symmetric difference
with $D$.
\end{proof}

This theorem identifies why three separately chosen noncrossing boundary
matchings are not yet a physical formation.  A switch in one family changes
the other two through the \emph{same interior edge set} $D$.  Its endpoint
pair on the boundary does not in general determine the connected components
of the two symmetric differences: that depends on where the three manifolds
share interior segments.  Accordingly the next realizability test must carry
one common trivalent web (or equivalent edge-incidence data) and check all
three equations simultaneously.  Treating the three matchings as independent
would discard exactly the source's idemposition coupling and cannot discharge
the surviving physical-exterior loophole.

The part of this dynamics which belongs to one fixed colour pair is completely
rigid, not merely existential.

\begin{theorem}[same-family physical Kempe persistence]
\label{thm:same-family-physical-kempe-persistence}
Let $T$ be a finite cubic port tangle with a proper Tait colouring $C$, fix a
colour pair $q$, and let $R$ be a union of connected components of the
$q$-bichromatic web.  If $C'$ is obtained by switching the two colours on
$R$, then the $q$-web has exactly the same local edges and the same connected
components in $C'$ as in $C$.  Consequently its physical boundary mate is
unchanged.
\Sverified{Mettapedia.GraphTheory.FourColor.GoertzelV24PortTanglePersistentKempe.kempeStep_switchRegion_iff}
\Sverified{Mettapedia.GraphTheory.FourColor.GoertzelV24PortTanglePersistentKempe.kempeConnected_switchRegion_iff}
\Sverified{Mettapedia.GraphTheory.FourColor.GoertzelV24PortTanglePersistentKempe.physicalMate_switchRegion_eq}
\end{theorem}

\begin{proof}
At every dart, switching $q$ exchanges the two members of $q$ and fixes every
other colour.  Membership in the unordered pair $q$ is therefore unchanged.
The two kinds of local web edge---the two active darts at a cubic vertex and
the two halves of an active interior edge---are consequently unchanged
pointwise.  Their equivalence closure, hence each connected component, is the
same relation before and after the switch.  The physical mate is defined as
the unique other boundary port in that component (and fixes inactive ports),
so it too is unchanged.
\end{proof}

This persistence is still not a boundary characterization of a literal
hexagonal patch.  The failure persists even if one remembers a deterministic
mate at every state rather than choosing a fresh witness after each switch.

\begin{proposition}[persistent pairwise mates do not determine the hexagon]
\label{prop:hexagon-persistent-pairwise-adversary}
There is a sixty-word set $A$ of six-port Tait words with the following
properties.
\begin{enumerate}
\item $A$ is nonempty and is disjoint from the exact boundary support of the
      literal hexagonal patch.
\item For every $w\in A$ and every one of the three colour pairs $q$, there is
      a deterministic noncrossing active mate $m_q(w)$.
\item Whenever a union of complete $m_q(w)$-pairs is switched, the resulting
      word $w'$ belongs to $A$ and $m_q(w')=m_q(w)$.
\end{enumerate}
The set is colour-blind: it is the union of the ten equality-pattern classes
\[
\begin{gathered}
001001,\quad001010,\quad010001,\quad010010,\quad010212,\\
010221,\quad011202,\quad011220,\quad012102,\quad012120,
\end{gathered}
\]
where equal digits require equal colours and distinct digits require distinct
colours.
\Sverified{Mettapedia.GraphTheory.FourColor.GoertzelV24HexagonPersistentKempeAdversary.persistentAdversary_card_eq_sixty}
\Sverified{Mettapedia.GraphTheory.FourColor.GoertzelV24HexagonPersistentKempeAdversary.persistentAdversary_hasPersistentCandidate}
\Sverified{Mettapedia.GraphTheory.FourColor.GoertzelV24HexagonPersistentKempeAdversary.persistentAdversary_disjoint_hexagonSupport}
\Sverified{Mettapedia.GraphTheory.FourColor.GoertzelV24HexagonPersistentKempeAdversary.persistentAdversary_nonempty}
\end{proposition}

\begin{proof}
Define $A$ by the ten displayed equality patterns.  Because the definition
uses only equality and inequality, it is invariant under renaming the three
Tait colours.  Direct finite reduction over all $3^6$ boundary words gives
$|A|=60$.  For each of the three colour pairs and each supported word, the
certificate supplies one of the five noncrossing matchings on the active
ports.  The checker verifies both the ordinary active-matching conditions and
all $2^6$ selected port sets; for every selected set which is a union of
complete pairs, it verifies membership of the switched word in $A$ and
equality of the selected mate before and after the switch.  The same complete
enumeration verifies that no word in $A$ extends across the literal hexagon.
Finally the equality pattern $001001$, instantiated by two distinct Tait
colours, supplies an explicit member of $A$.  The flat tables used to suggest
the five-way matching choices are not trusted: the Boolean witness condition
is proved equivalent to the mathematical proposition, and the kernel reduces
every row of the three separate certificates.
\end{proof}

\begin{proposition}[the persistent adversary does not obstruct the cap menu]
\label{prop:persistent-adversary-cap-menu}
Let $A$ be the sixty-word language of
Proposition~\ref{prop:hexagon-persistent-pairwise-adversary}.  In the
zero-based enumeration of the five noncrossing planar pairing caps, $A$ is
disjoint from the supports of caps $0=(01)(23)(45)$ and
$2=(03)(12)(45)$.  In particular $A$ does not meet every planar cap.
\Sverified{Mettapedia.GraphTheory.FourColor.GoertzelV24HexagonPersistentAdversaryCapMenu.persistentAdversary_disjoint_planarPairingSupport_zero}
\Sverified{Mettapedia.GraphTheory.FourColor.GoertzelV24HexagonPersistentAdversaryCapMenu.persistentAdversary_disjoint_planarPairingSupport_two}
\Sverified{Mettapedia.GraphTheory.FourColor.GoertzelV24HexagonPersistentAdversaryCapMenu.persistentAdversary_not_meets_every_planarPairing}
\end{proposition}

\begin{proof}
The finite checker tests every word in $A$ against the two displayed pairing
relations and finds no compatible word.  Either disjointness statement alone
negates the assertion that $A$ meets all five cap supports.
\end{proof}

This corrects an important scope distinction.  The sixty-word language is a
witness against compressing the common physical web to three persistent
boundary mates; it is not an obstruction to replacing a hexagon by a member
of the five-cap menu.  The larger parity adversary in
Proposition~\ref{prop:hexagon-target-aware-pairing-obstruction} is the
language which meets every cap.

Nor should the universal physical cap-menu implication be treated as a
cheap local search.  In the route's intended use, the exterior is obtained
by deleting a facial hexagon from a hypothetical bridgeless planar cubic
counterexample.  Exhibiting such an exterior already exhibits the
counterexample, while proving that one cap always works is precisely a
uniform reducibility theorem for the dual degree-six vertex.  Even that
strong result would not by itself bound width: it supplies no exclusion of
large pentagon--heptagon wall patches.  Thus the central width problem remains
global.

Thus even the exact boundary shadow of
Theorem~\ref{thm:same-family-physical-kempe-persistence}, imposed separately
for all three pairs and all of their successor states, admits a language which
no literal hexagon realizes.  What is absent is precisely the cross-family
clause in Theorem~\ref{thm:kempe-idemposition-identity}: all three updates must
be induced by one common selected interior edge set.  The next state therefore
has to retain a common trivalent web, or equivalent shared edge-incidence
data; adding further independent pairwise closure fields cannot close this
gap.

The corresponding unquotiented joint carrier is exact.  This separates two
questions which must not be conflated: whether the source's simultaneous
idemposition has been represented faithfully, and whether that representation
has already been compressed to a bounded interface receipt.

\begin{theorem}[joint physical idemposition on one carrier]
\label{thm:joint-physical-idemposition}
Let \(T\) be a cubic port tangle and let \(C\) be a proper Tait colouring of
its common dart carrier \(D\).  Write

\[
 \mathcal M(C)=\bigl(M_{rb}(C),M_{rp}(C),M_{bp}(C)\bigr).
\]

If \(R\subseteq D\) is a physical Kempe region for a colour pair \(q\), and
\(C'\) is obtained by switching \(q\) on \(R\), then the \(q\)-coordinate of
\(\mathcal M(C)\) is fixed and the other two coordinates are both replaced by
their symmetric difference with the \emph{same} set \(R\).  Moreover the map
\(C\mapsto\mathcal M(C)\) is injective on proper Tait colourings.

\Sverified{Mettapedia.GraphTheory.FourColor.GoertzelV24PortTangleJointKempeState.manifoldTriple_switchRegion}
\Sverified{Mettapedia.GraphTheory.FourColor.GoertzelV24PortTangleJointKempeState.manifoldTriple_injective_of_proper}
\end{theorem}

\begin{proof}
The first assertion is pointwise.  On \(R\), the two colours of \(q\) are
interchanged.  Membership in \(M_q\) is unchanged, while membership in each
companion bichromatic support is toggled once.  Outside \(R\), nothing
changes.  Because all three supports are subsets of the one dart carrier, the
two companion updates use literally the same \(R\), rather than two boundary
sets merely having the same endpoints.

For injectivity, a nonzero colour is recovered from its membership pattern in
the three supports:
\[
\begin{array}{c|ccc}
 &M_{rb}&M_{rp}&M_{bp}\\ \hline
r&1&1&0\\
b&1&0&1\\
p&0&1&1.
\end{array}
\]
Thus equal triples give equal colours on every dart and hence equal
colourings.
\end{proof}

This theorem is the first carrier in this section which satisfies the exact
cross-family law demanded by the source.  It is intentionally not yet a
finite-profile theorem: \(\mathcal M(C)\) remembers subsets of every interior
dart of \(T\).  The remaining interface problem is therefore precise.  One
must quotient this shared carrier to bounded boundary data while proving that
the quotient still determines simultaneous physical updates and the
replacement support.  Projecting the three coordinates separately is ruled
out by Proposition~\ref{prop:hexagon-persistent-pairwise-adversary}.

Even retaining the three projected coordinates in one boundary state does
not repair that loss.

\begin{proposition}[projected joint idemposition does not determine the hexagon]
\label{prop:hexagon-projected-joint-adversary}
Let
\[
 \partial\mathcal M(w)=
 \bigl(A_{rb}(w),A_{rp}(w),A_{bp}(w)\bigr),
\]
where \(A_q(w)\) is the set of ports whose colours belong to \(q\).
There is a nonempty sixty-word language \(A\), disjoint from the literal
hexagon support, with a deterministic noncrossing mate \(m_q(w)\) for every
\(w\in A\) and every colour pair \(q\), such that the following holds.
Whenever \(S\) is a union of active \(m_q(w)\)-pairs and \(w'\) is obtained
by switching \(q\) on \(S\), then
\[
 w'\in A,\qquad m_q(w')=m_q(w),\qquad
 \partial\mathcal M(w')
   =\partial\mathcal M(w)\mathbin{\operatorname{toggle}}_q S.
\]
In the final equation the \(q\)-coordinate is fixed and both companion
coordinates are symmetrically differenced with the same set \(S\).
\Sverified{Mettapedia.GraphTheory.FourColor.GoertzelV24HexagonProjectedJointAdversary.boundaryManifoldTriple_swapBoundaryWord}
\Sverified{Mettapedia.GraphTheory.FourColor.GoertzelV24HexagonProjectedJointAdversary.persistentAdversary_survives_projectedJointBoundary}
\end{proposition}

\begin{proof}
Use the language and deterministic mates from
Proposition~\ref{prop:hexagon-persistent-pairwise-adversary}.  That
certificate already proves \(w'\in A\) and persistence of the operated mate.
For the displayed joint equation, argue at one port.  If the port is outside
\(S\), its colour and all three membership bits are unchanged.  If it lies in
\(S\), its colour is one of the two colours of \(q\); exchanging them fixes
membership in \(q\) and toggles membership in each of the other two pairs.
This is exactly the three-coordinate toggle by the one set \(S\).
\end{proof}

The proposition does not contradict
Theorem~\ref{thm:joint-physical-idemposition}.  It identifies the information
lost by restricting that theorem to the ports.  At the boundary, the three
support coordinates are already determined by \(w\), so writing them as one
triple adds no shared-interior information.  The next receipt must retain a
quotient of how the three manifolds overlap and reconnect inside the tangle
(including the contraction chirality above); a word, three mates, and their
entire boundary transition graph are still insufficient.

There is also a sharp warning against treating an arbitrary \emph{signed}
pairing identity as if it were already a planar counting argument.  Write a
perfect matching of the six ports as, for example, \(01|23|45\), and let
\(\delta_M(w)\in\{0,1\}\) be one precisely when the boundary word \(w\) is
constant on every pair of \(M\).  Enumerate the fifteen matchings by
\[
\begin{array}{c|ccccc}
0&01|23|45&01|24|35&01|25|34&02|13|45&02|14|35\\
5&02|15|34&03|12|45&03|14|25&03|15|24&04|12|35\\
10&04|13|25&04|15|23&05|12|34&05|13|24&05|14|23.
\end{array}
\]

\begin{proposition}[a signed pairing tensor has exactly the adversary support]
\label{prop:hexagon-signed-pairing-adversary}
Define
\[
\begin{aligned}
 k_0&=\delta_8-\delta_4,\\
 k_1&=\delta_{10}-\delta_4,\\
 k_2&=-\delta_0+\delta_1+\delta_3-2\delta_4+\delta_{11},\\
 k_3&=-2\delta_4+\delta_5+\delta_9-\delta_{12}+\delta_{13},
 \qquad K=k_0+2k_1+4k_2+8k_3.
\end{aligned}
\]
For every genuine nonzero six-port Tait word \(w\),
\[
                    K(w)\ne0\quad\Longleftrightarrow\quad w\in A\setminus P.
\]
In particular every \(k_i\), and hence \(K\), vanishes on the literal
hexagon support \(P\), while the support of \(K\) is exactly the whole
parity adversary.
\Sverified{Mettapedia.GraphTheory.FourColor.GoertzelV24HexagonPairingInvariantKernel.kernelCoordinates_eq_zero_of_mem_hexagonSupport}
\Sverified{Mettapedia.GraphTheory.FourColor.GoertzelV24HexagonPairingInvariantKernel.invariantAdversaryTensor_ne_zero_iff}
\Sverified{Mettapedia.GraphTheory.FourColor.GoertzelV24HexagonPairingInvariantKernel.invariantAdversaryTensor_eq_zero_of_mem_hexagonSupport}
\end{proposition}

\begin{proof}
There are three choices at each of the six ports.  For each of the \(3^6\)
words, evaluate the fifteen Kronecker-delta products and then the four
displayed integer combinations.  Independently test the two defining
conditions for \(A\setminus P\): zero total Klein colour and failure of the
six-step hexagon extension recurrence.  The resulting two Boolean columns
agree in all \(729\) rows.  The same exhaustive evaluation checks that each
\(k_i\) is zero on every row accepted by the recurrence.  Lean evaluates
both closed tables from the displayed mate involutions and proves the
biconditional and the four coordinate equalities.
\end{proof}

The coefficients above are signed.  Consequently \(K\) is a formal integer
linear combination of pairing contractions, not the nonnegative extension
count of one planar exterior, and several of its matchings are crossing in
the cyclic boundary order.  The proposition therefore does not realize the
adversary physically.  It says something more diagnostic: ordinary linear
pairing identities alone can carry the entire abstract obstruction.  To
rule it out one must use the positive cone and the common embedded web from
which all three two-colour pairings arise.  This is exactly the extra content
suggested by the source's Penrose--Kauffman evaluation and idemposition
calculus.

\begin{proposition}[pairing-linear positivity cannot separate the adversary]
\label{prop:hexagon-no-linear-pairing-separator}
Let
\[
                   L(w)=\sum_{M}c_M\delta_M(w),\qquad c_M\in\mathbb Z,
\]
be any integer linear combination of the fifteen pairing contractions.  If
\(L\) vanishes on the literal hexagon support \(P\), then \(L\) is neither
strictly positive nor strictly negative on all of \(A\setminus P\).
\Sverified{Mettapedia.GraphTheory.FourColor.GoertzelV24HexagonPairingNoLinearSeparator.no_strictlyPositive_pairingLinearSeparator}
\Sverified{Mettapedia.GraphTheory.FourColor.GoertzelV24HexagonPairingNoLinearSeparator.no_strictlyNegative_pairingLinearSeparator}
\end{proposition}

\begin{proof}
Take the following five adversary words:
\[
\begin{aligned}
q_0&=(r,r,b,r,r,b),&q_1&=(r,r,b,r,b,r),&
q_2&=(r,r,b,p,b,p),\\
q_3&=(r,b,r,r,b,r),&q_4&=(r,b,r,p,p,b).
\end{aligned}
\]
Lean checks \(q_i\in A\setminus P\) for all five.
\Sverified{Mettapedia.GraphTheory.FourColor.GoertzelV24HexagonPairingNoLinearSeparator.adversaryWitness_mem}
Now list eleven words accepted by the hexagon:
\[
\begin{array}{c@{\;}l c@{\;}l c@{\;}l}
0&(r,r,r,r,r,r)&1&(r,r,r,r,b,b)&2&(r,r,r,b,b,r)\\
3&(r,r,b,b,r,r)&4&(r,r,b,b,b,b)&5&(r,r,b,p,p,b)\\
6&(r,b,b,r,r,r)&7&(r,b,b,r,p,p)&8&(r,b,b,b,b,r)\\
9&(r,b,p,r,b,p)&10&(r,b,p,p,b,r).&&
\end{array}
\]
\Sverified{Mettapedia.GraphTheory.FourColor.GoertzelV24HexagonPairingNoLinearSeparator.hexagonRelationWitness_mem}
For each one of the fifteen perfect matchings \(M\), direct evaluation gives
\[
\begin{split}
\sum_{i=0}^{4}\delta_M(q_i)
={}&\delta_M(p_0)-\delta_M(p_1)-\delta_M(p_3)+\delta_M(p_4)
     -\delta_M(p_5)\\
 &-\delta_M(p_6)+\delta_M(p_7)+\delta_M(p_9)+\delta_M(p_{10}).
\end{split}
\]
This finite row identity is checked for every \(M\).
\Sverified{Mettapedia.GraphTheory.FourColor.GoertzelV24HexagonPairingNoLinearSeparator.pairContraction_farkasRelation}
Multiply its \(M\)-th instance by \(c_M\) and sum over \(M\).  Linearity
turns the left side into \(\sum_iL(q_i)\), while every term on the right is
zero because each \(p_j\) lies in \(P\).  Thus
\(\sum_iL(q_i)=0\).  The five summands cannot all be strictly positive, and
they cannot all be strictly negative.
\Sverified{Mettapedia.GraphTheory.FourColor.GoertzelV24HexagonPairingNoLinearSeparator.sum_adversaryWitness_eq_zero_of_vanishesOnHexagon}
\end{proof}

Thus the most direct Penrose repair is now ruled out precisely.  Positivity
of the \emph{same fifteen pairing coordinates} cannot supply the missing
physical-exterior restriction: their entire linear span has no strict
separator.  A successful argument must enlarge the observable to a common
trivalent web (retaining shared-edge incidence and the three coupled
idemposition equations), or exploit a nonlinear positivity statement about
actual extension counts.  This is a narrowing of the high-width problem,
not evidence against the source's broader operator programme.

There is one sign convention hidden in that conclusion which must be audited
before moving on.  The source's colouring operator $F_X$ is the
\emph{unsigned} nonnegative extension-count operator.  By contrast, an open
Penrose contraction written in the fifteen pairing coordinates carries a
boundary sign.  For a six-port word with even colour multiplicities, list the
occurrences of each colour in boundary order and pair consecutive
occurrences.  Let

\[
             \sigma(w)=(-1)^{\operatorname{cr}(M(w))}
\]

be the parity of the crossings among those three chords.  If a colour has
odd multiplicity, all fifteen pairing contractions vanish and the value
assigned to $\sigma(w)$ is immaterial.

\begin{proposition}[the gauge-positive pairing cone still contains the
persistent adversary]
\label{prop:hexagon-penrose-gauge-adversary}
Put
\[
 L=\delta_4-\delta_5-\delta_9+\delta_{10}
       +\delta_{12}-\delta_{13}.
\]
Then, on every genuine nonzero six-port Tait word,
\[
 \sigma(w)L(w)\geq0,
 \qquad
 \sigma(w)L(w)>0\quad\Longleftrightarrow\quad w\in E,
\]
where $E$ is the sixty-word persistent Kempe adversary.  Moreover $L$
vanishes on the literal hexagon support $P$.  For each of the three
four-vertex plane cubic tree shapes there is a word in its support on which
$\sigma L$ is positive.
\Sverified{Mettapedia.GraphTheory.FourColor.GoertzelV24HexagonPenroseGaugeAdversary.penroseKernelTensor_eq_zero_of_mem_hexagonSupport}
\Sverified{Mettapedia.GraphTheory.FourColor.GoertzelV24HexagonPenroseGaugeAdversary.gaugeCorrectedPenroseTensor_nonnegative}
\Sverified{Mettapedia.GraphTheory.FourColor.GoertzelV24HexagonPenroseGaugeAdversary.gaugeCorrectedPenroseTensor_pos_iff}
\Sverified{Mettapedia.GraphTheory.FourColor.GoertzelV24HexagonPenroseGaugeAdversary.persistentAdversary_is_gaugePositive_pairingKernel}
\end{proposition}

\begin{proof}
In the displayed enumeration of matchings, crossing parity alternates with
the index.  For an even-multiplicity word, the first compatible matching in
that enumeration is exactly the matching obtained by pairing consecutive
occurrences of each colour.  Thus its index parity computes $\sigma(w)$.
Evaluate this sign and the six displayed pairing coordinates on the
$3^6=729$ boundary words.  The corrected value is nonnegative in every
row, and its positive rows are precisely the ten colour-blind equality
patterns defining $E$, hence precisely its sixty words.

The vanishing on $P$ also has a short algebraic check.  In the notation of
Proposition~\ref{prop:hexagon-signed-pairing-adversary},
\[
                         L=k_1-k_3,
\]
and both kernel coordinates vanish on $P$.  Finally the three positive
rows
\[
 \mathtt{rrbrrb},\qquad \mathtt{rbrrrb},\qquad \mathtt{rbrpbp}
\]
(spaces only separate the six boundary positions) are respectively accepted
by the separated-path, adjacent-path, and tripod tree supports.  Lean checks
the complete table, the kernel identity, and these three explicit extension
witnesses.
\end{proof}

This proposition is stronger than the ungauged linear-separator failure, but
its scope is equally important.  It does \emph{not} say that $L$, or its
positive support $E$, is realized by one planar cubic exterior.  It says
that sign-corrected linear invariant theory plus coordinatewise positivity
still permits exactly the already known persistent obstruction.  The next
condition must therefore use the nonlinear, integral semigroup of tensors
realized by actual common trivalent webs (or an equivalent shared-interior
receipt).  That distinction keeps this negative result aligned with the
source's unsigned operator $F_X$, rather than mistaking a formal pairing
expansion for an extension count.

The remaining gap in that last paragraph can be seen in one small graph.
The positive tensor is an honest unsigned extension count, but only after the
cyclic order of the six ports is allowed to cross.

\begin{proposition}[the gauge adversary is a crossed pentagon--leaf tensor]
\label{prop:hexagon-crossed-pentagon-leaf-adversary}
Let $Q$ be the abstract trivalent network with internal vertices
$v_0,\ldots,v_5$, internal edges
\[
 v_0v_1,\ v_0v_4,\ v_1v_2,\ v_1v_5,\ v_2v_3,\ v_3v_4,
\]
and six boundary half-edges attached, in the prescribed cyclic port order,
at
\[
                         v_3,v_0,v_2,v_4,v_5,v_5.
\]
For every nonzero Tait boundary word $w$, the number $N_Q(w)$ of proper
internal edge-colourings is
\[
              N_Q(w)=\sigma(w)L(w)=
              \begin{cases}1,&w\in E,\\0,&w\notin E.\end{cases}
\]
In particular each supported word has exactly one internal colouring.

Now adjoin boundary vertices $b_0,\ldots,b_5$ at the six half-edges and the
boundary cycle $b_0b_1\cdots b_5b_0$.  The resulting graph contains
$K_{3,3}$ as a minor.  Consequently $Q$ has no embedding in a closed disc
whose boundary ports occur in the displayed cyclic order.
\Sverified{Mettapedia.GraphTheory.FourColor.GoertzelV24HexagonCrossedPentagonLeafAdversary.extension_eq_crossedPentagonLeafCandidate}
\Sverified{Mettapedia.GraphTheory.FourColor.GoertzelV24HexagonCrossedPentagonLeafAdversary.card_crossedPentagonLeafExtensions}
\Sverified{Mettapedia.GraphTheory.FourColor.GoertzelV24HexagonCrossedPentagonLeafAdversary.existsUnique_crossedPentagonLeafExtension_iff}
\Sverified{Mettapedia.GraphTheory.FourColor.GoertzelV24HexagonCrossedPentagonLeafAdversary.gaugeCorrectedPenroseTensor_eq_extensionCount}
\Sverified{Mettapedia.GraphTheory.FourColor.GoertzelV24HexagonCrossedPentagonLeafAdversary.k33_isMinor_crossedPentagonLeafBoundaryGraph}
\Sverified{Mettapedia.GraphTheory.FourColor.GoertzelV24HexagonCrossedPentagonLeafAdversary.persistentAdversary_is_crossedPentagonLeafTensor_with_k33Minor}
\end{proposition}

\begin{proof}
Name the internal edge colours, in the order displayed above, by
$x_{01},x_{04},x_{12},x_{15},x_{23},x_{34}$.  At the pendant vertex $v_5$,
properness forces
\[
                         x_{15}=w_4+w_5.
\]
At $v_1$ the other two colours are therefore the unordered pair
$\{w_4,w_5\}$: they are distinct and nonzero, and in $\mathbb F_2^2$ two
pairs of distinct nonzero colours having the same sum are the same unordered
pair.  There are consequently only two candidate orientations,
\[
 (x_{01},x_{12})=(w_4,w_5)
       \quad\hbox{or}\quad
 (x_{01},x_{12})=(w_5,w_4).
\]
Once one is chosen, the remaining vertices force, successively,
\[
 x_{23}=x_{12}+w_2,\qquad
 x_{34}=x_{23}+w_0,\qquad
 x_{04}=x_{34}+w_3.
\]
Thus every internal colouring is one of two explicitly displayed candidates;
no enumeration over arbitrary internal colourings is involved.  Evaluating
local properness for those two candidates on the $3^6$ boundary words gives
one valid candidate precisely on the ten equality-pattern classes defining
$E$, and none elsewhere.  The same finite table agrees row by row with
$\sigma L$.  Lean proves the forcing argument, the equivalence between the
two-candidate fibre and the full colouring fibre, and the $729$-row boundary
check.

For the topological assertion, use the following six disjoint connected
branch sets in the boundary-cycle graph:
\[
\begin{array}{lll}
A_0=\{v_1,v_2,v_5\},&A_1=\{b_0,b_1\},&
A_2=\{v_4,b_3,b_4\},\\
B_0=\{v_3\},&B_1=\{b_2\},&B_2=\{b_5\}.
\end{array}
\]
Each $A_i$ is adjacent to each $B_j$: the nine witnessing edges are,
row by row,
\[
\begin{array}{lll}
v_2v_3,&v_2b_2,&v_5b_5,\\
b_0v_3,&b_1b_2,&b_0b_5,\\
v_4v_3,&b_3b_2,&b_4b_5.
\end{array}
\]
Contracting each branch set therefore gives $K_{3,3}$.  If $Q$ had the
claimed disc embedding, the six boundary-cycle edges could be drawn along
the rim, producing a plane embedding of this augmented graph.  This is
impossible because planarity is minor-closed and $K_{3,3}$ is nonplanar.
\end{proof}

This identifies the failure more sharply than positivity alone.  The tensor
$\sigma L$ is the unsigned, integral, uniquely-colourable tensor of a tiny
common trivalent network; what prevents it from being a physical exterior is
the \emph{joint cofacial order} of that network.  The proposition does not
claim that every possible realization of the same tensor has this particular
$K_{3,3}$ model.  It says that any successful local receipt must be strong
enough to distinguish a disc-realizable common web from this crossed one.

There is no hidden second pairing-linear realization with different
multiplicities on the sixty supported words.  The following calculation
shows that the preceding Penrose tensor spans the whole pairing-linear ray
with that support.

\begin{proposition}[rigidity of the Penrose ray on the persistent adversary]
\label{prop:hexagon-penrose-ray-rigidity}
\Sverified{Mettapedia.GraphTheory.FourColor.GoertzelV24HexagonPenroseRayRigidity.persistentAdversary_pairingRay_rigid}
Let \(E\) be the sixty-word persistent adversary.  Using the zero-based
enumeration of the fifteen pairings displayed above, let

\[
              T_c(w)=\sum_{i=0}^{14}c_i\delta_i(w),
              \qquad c_i\in\mathbb Z.
\]

If \(T_c(w)=0\) for every nonzero six-port Tait word outside \(E\), then

\[
 c_i=c_4d_i\quad(0\leq i<15),\qquad
 (d_0,\ldots,d_{14})
   =e_4-e_5-e_9+e_{10}+e_{12}-e_{13}.
\]

Consequently \(T_c=c_4L\), and for every \(w\in E\),

\[
                         \sigma(w)T_c(w)=c_4.
\]
\end{proposition}

\begin{proof}
The following fourteen words all lie outside \(E\):

\[
\begin{array}{llll}
\mathtt{rrrrrr},&\mathtt{rrrrbb},&\mathtt{rrrbrb},&\mathtt{rrrbbr},\\
\mathtt{rrbbrr},&\mathtt{rrbbbb},&\mathtt{rrbbpp},&\mathtt{rrbpbp},\\
\mathtt{rbrbrr},&\mathtt{rbrbbb},&\mathtt{rbrbpp},&\mathtt{rbbbrb},\\
\mathtt{rbprbp},&\mathtt{rbpprb}.&&
\end{array}
\]

Evaluate the fifteen indicators \((\delta_0,\ldots,\delta_{14})\) on these
rows.  The resulting fourteen
homogeneous equations are independent.  Successive integer elimination,
with \(c_4\) left free, gives

\[
c_5=-c_4,\quad c_9=-c_4,\quad c_{10}=c_4,\quad
c_{12}=c_4,\quad c_{13}=-c_4,
\]

and every other \(c_i\) is zero.  The checked proof evaluates the fourteen
rows directly from the displayed mate involutions and performs this
elimination over \(\mathbb Z\); hence no division or rank-over-a-field
assumption is hidden.  This proves \(T_c=c_4L\).  On \(E\), the already
checked gauge-corrected Penrose tensor is the extension count of the crossed
pentagon--leaf network and equals \(1\).  Therefore
\(\sigma(w)T_c(w)=c_4\) on all of \(E\).
\end{proof}

This is an algebraic rigidity result.  Within the pairing-linear ansatz it
closes the unequal-multiplicity loophole: support exactly \(E\) forces one
constant gauge-corrected multiplicity.  It does \emph{not} prove that a
positive multiple of this ray is realized by a plane cubic tangle.  That
remaining condition is topological, and the crossed realization above shows
why cyclic order cannot be discarded.

\begin{theorem}[the exact width criterion]
\label{thm:exact-width-criterion}
\Sprose
Let \(\mathcal M\) be the class of vertex-minimal zero-target normal-form
instances.  Suppose that for every fixed width \(k\), ordered width-\(k\)
tree decompositions carry a finite exact typed context semantics, and equal
comparable states admit a strict zero-target replacement as in
Theorem~\ref{thm:finite-tree-interface-pumping}.  Then the following are
equivalent:
\begin{enumerate}
\item the vertex counts of the members of \(\mathcal M\) have a uniform
finite bound;
\item the treewidths of the members of \(\mathcal M\) have a uniform finite
bound.
\end{enumerate}
If the width bound and the finite context semantics are effective, then the
vertex bound is effective.
\end{theorem}

\begin{proof}
A graph on at most \(N\) vertices has treewidth at most \(N-1\), proving
\(1\Rightarrow2\).  Conversely, assume every member of \(\mathcal M\) has
treewidth at most \(k\).  Convert a width-\(k\) decomposition to a reduced
nice rooted binary decomposition.  Its introduce nodes account for at most
one graph vertex each; join, forget, and leaf nodes introduce none.  By
hypothesis its ordered interfaces and exact exterior behaviours form one
finite typed state set, say of cardinality \(q(k)\).  Minimality and exact
replacement forbid a repeated state on an ancestor--descendant chain.
Theorem~\ref{thm:finite-tree-interface-pumping}, with \(c=1\), gives
\[
                         |V(G)|\le 2^{q(k)}-1
\]
for every \(G\in\mathcal M\).  This proves \(2\Rightarrow1\), and the
displayed formula is effective when the stated inputs are.
\end{proof}

Thus, \emph{relative to exact finite tree semantics and D3 replacement}, the
tree-form repair of D1 is precisely a width theorem for minimal zero-target
instances.  This is not a literal proof of the manuscript's linear
hexagonal-corridor L1.  It is a compositional reformulation of that failed
global step.  The claimed equivalence below always has these semantic and
physical hypotheses in scope.

Neither the planar wall theorem nor a branchwidth dichotomy supplies the
replacement by itself.  Model paths of a wall subdivision may contain
arbitrarily many degree-two vertices before they are viewed inside the cubic
host, and ambient material may attach along them.  Cutting only abstract
rails therefore need not give a saturated cut of the original map; saturating
it can make the interface grow with the model paths.  The mesh theorem above
sharpens this warning: only the number of slabs that acquire new designated
branch positions is bounded.  A wall becomes useful to the proof only after
zero Count and minimality turn it into a strict exact-state replacement or
exclude it altogether.

\begin{proposition}[two exhaustive geometric branches for the global corridor]
\label{prop:global-corridor-alternative}
\Sconditional{bounded-face geodesic transfer; long-face perimeter transfer}
Fix constants $B,L$ above the relevant profile thresholds.  Assume both of
the following source-preserving constructions, including all the transversal
data in Proposition~\ref{prop:target-forced-cut-chain}.
\begin{enumerate}
\item A bounded-face geodesic of length $L$ factors through a finite
  manuscript profile and equal profiles admit a physical deletion splice.
\item A face of length greater than $B$ admits an analogous perimeter
  factorization and a strictly size-decreasing, target-preserving splice.
\end{enumerate}
Then every sufficiently large irreducible obstruction is pumpable, with no
bound on $W_-$ and no claim that most faces are hexagons.
\end{proposition}

\begin{proof}
Apply Theorem~\ref{thm:dual-diameter-dichotomy}.  In its bounded-face branch,
the first assumed construction supplies the required composable cut chain and
splice.  In its long-face branch, the second does.  The vertex/face identity
$F=V/2+2$ makes the dichotomy's face threshold an effective vertex threshold.
Thus every obstruction above the maximum of the two construction thresholds
is pumpable.
\end{proof}

Theorem~\ref{thm:dual-diameter-dichotomy} proves that the two \emph{geometric
inputs}---a bounded-face geodesic or a long face---are exhaustive.  It does
not prove either transfer construction.  This is nevertheless a sharper
research front than a uniform-curvature hypothesis: the high-negative-
curvature branch is retained explicitly instead of being discarded by an
invalid unsigned-defect inference.

The original v23 manuscript proposes a different genuinely two-dimensional
move: its Theorem~4.9 claims that annular Kempe-closure face generators span
the whole boundary-zero cycle space.  That would be highly relevant here, so
we rechecked the actual proof rather than relying on its later summaries.  One
of its two concluding inferences is not valid as stated.

\begin{proposition}[the v23 annihilator inference needs a containment lemma]
\label{prop:v23-annihilator-gap}
\Srefuted{Mettapedia.GraphTheory.FourColor.GoertzelV24V23AnnihilatorGap.trivial_relative_annihilator_but_not_contained}
For subspaces $U,W$ of a finite-dimensional space with a nondegenerate
bilinear form, the condition
\[
                         U^\perp\cap W=\{0\}
\]
does not by itself imply $W\subseteq U$.
\end{proposition}

\begin{proof}
Over $\mathbb F_2$, take the ambient space $\mathbb F_2^2$ with its standard
dot product, let $W=\langle(1,0)\rangle$, and let
$U=\langle(1,1)\rangle$.  A vector $(x,0)\in W$ orthogonal to $(1,1)$ has
$x=0$, so $U^\perp\cap W=\{0\}$.  But $(1,0)\notin U$, hence
$W\nsubseteq U$.  The Lean theorem proves both assertions on literal
two-coordinate functions over $\mathbb Z/2\mathbb Z$.
\end{proof}

In v23, $W$ is $W_0(H)$ and $U$ is the span of the purified face generators.
The proof establishes the displayed annihilator condition and immediately
concludes $W_0(H)\subseteq U$.  It does not first establish $U\subseteq
W_0(H)$; indeed the definition of a purified generator explicitly permits
support on boundary edges, while $W_0(H)$ vanishes there.  The dual-forest
peeling also treats the unique \emph{tree} edge incident with a leaf face as
though it were the unique remaining dual-graph edge, which need not hold in a
spanning forest.  Proposition~\ref{prop:v23-annihilator-gap} by itself only
refutes that inference.  The literal spanning statement is in fact false for
an additional, concrete reason.

\begin{proposition}[triangular-prism counterexample to the literal v23 spanning theorem]
\label{prop:v23-triangular-prism-spanning-counterexample}
\Srefuted{Mettapedia.GraphTheory.FourColor.GoertzelV24V23TriangularPrismSpanningCounterexample.exists_boundaryZero_not_mem_generatorSpan}
Let $H$ be the triangular prism embedded as an annulus, with its two
triangular faces as the two boundary cycles and its three rectangular faces
as the internal faces.  Define the v23 generator $X_f^{ab}(C)$ exactly as in
the final Kempe-closure generator-family definition: it has coefficient equal
to the third colour on those edges of $\partial f$ whose colours lie in
$\{a,b\}$, and coefficient zero elsewhere.  Even if one permits generators
from \emph{every} proper Tait colouring of $H$, rather than only the Kempe
closure of a fixed colouring, their span does not contain $W_0(H)$.
\end{proposition}

\begin{proof}
Encode the three nonzero colours by $\mathbb F_2^2$.  Number the top triangle
edges $t_0,t_1,t_2$, the bottom triangle edges $b_0,b_1,b_2$, and the spokes
$s_0,s_1,s_2$, with rectangular face $f_i$ incident with
$t_i,b_i,s_i,s_{i+1}$ (indices modulo three).  Define a linear functional
$\lambda$ on $\mathbb F_2^2$-valued edge chains by summing the first coordinate
on
\[
                    t_0,t_1,t_2,s_0,s_1,s_2.
\]
For every proper colouring, the three colours on
$t_i,s_i,s_{i+1}$ are pairwise distinct.  The first two inequalities are
properness at the two endpoints of $t_i$; the third follows because the three
top edges have pairwise distinct colours and $s_i,s_{i+1}$ are the respective
missing colours at consecutive top vertices.  Hence any chosen colour pair
contains exactly two of these three edges.  On $X_{f_i}^{ab}(C)$ the two
selected edges carry the same third-colour coefficient, so their first
coordinates cancel in $\mathbb F_2$.  The bottom edge does not occur in
$\lambda$.  Thus $\lambda$ annihilates every generator arising from every
proper colouring, and therefore annihilates their span.

Now let $w$ have first coordinate one on $s_0$ and be zero everywhere else.
Every vertex of this annulus lies on a boundary component, so the v23
Kirchhoff conditions, which are imposed only at interior vertices, are
vacuous; moreover $w$ vanishes on both boundary cycles.  Thus $w\in W_0(H)$.
But $\lambda(w)=1$, so $w$ is not in the generator span.  The Lean model uses
all $3^9$ edge assignments satisfying the six literal cubic properness
conditions and proves this separating-functional argument without enumerating
a Kempe closure.
\end{proof}

The triangular prism is planar, cubic, and bridgeless, and its two triangular
cycles are disjoint boundary components, so it satisfies the hypotheses as
printed.  It is boundary-degenerate in the precise sense that it has no
interior vertices.  That observation does not repair the theorem: the same
failure occurs with two pentagonal boundaries and ten genuine interior
vertices.

\begin{proposition}[positive-depth dodecahedral counterexample]
\label{prop:v23-dodecahedral-spanning-counterexample}
\Srefuted{Mettapedia.GraphTheory.FourColor.GoertzelV24V23DodecahedralSpanningCounterexample.exists_boundaryZero_not_mem_generatorSpan}
Let $G$ be the dodecahedral graph in its standard spherical embedding, and
let $H$ be obtained by declaring the disjoint pentagonal faces numbered $0$
and $6$ to be its two boundary components.  Use the remaining ten faces as
the internal faces.  For the explicit Tait colouring $C_0$ recorded in the
certificate, the literal Definition--4.8 generators from the exact
three-colour Kempe closure of $C_0$ do not span $W_0(H)$.
\end{proposition}

\begin{proof}
The two selected face boundaries each contain five edges and are disjoint;
the ten vertices not on them are
\[
               6,7,8,9,10,11,12,13,14,18.
\]
In $C_0$, for each of the three unordered colour pairs the corresponding
bicoloured subgraph is connected.  A Kempe component for two distinct colours
is therefore the whole bicoloured subgraph, so switching it is just a global
transposition of those two colours.  Switching a colour with itself changes
nothing.  Induction on the Kempe reachability relation shows that every
colouring in the exact source closure has the form $\sigma C_0$ for a
permutation $\sigma$ of the three colours.

Write $e_{ij}$ for the dodecahedral edge joining vertices $i,j$.  Define the
linear functional $\lambda$ by summing the first Klein coordinate on
\[
 e_{01},e_{3,19},e_{9,10},e_{10,11},e_{11,18},e_{12,13},e_{12,16},
 e_{15,16},e_{16,17},e_{17,18},e_{18,19}.
\]
The finite certificate checks, for all six permutations $\sigma$, all ten
internal faces, and all three colour pairs, that
\[
                    \lambda\bigl(X_f^{ab}(\sigma C_0)\bigr)=0.
\]
Hence $\lambda$ annihilates the span of the literal source generators.

Let $w$ have first coordinate one on the four-edge path
\[
                  e_{0,10},e_{10,11},e_{11,18},e_{18,19}
\]
and vanish elsewhere.  Its endpoints lie on a boundary pentagon, its three
internal path vertices have even incidence, and all other internal vertices
have incidence zero.  It also vanishes on both boundary cycles.  Thus
$w\in W_0(H)$.  Exactly three of its four edges occur in the support of
$\lambda$, so $\lambda(w)=1$ over $\mathbb F_2$.  Therefore $w$ is not in the
generator span.  Lean checks the graph, the two five-edge boundaries, their
disjointness, all ten Kirchhoff equations, the exact Kempe-closure reduction,
and the separating calculation.
\end{proof}

Thus neither boundary degeneracy nor a mere positive-depth hypothesis can
repair the literal v23 theorem.  A route-realizability restriction would be a
new statement needing a new proof.  The v23 spanning theorem cannot serve as
the missing high-width M1 theorem.  A corrected need-based generator theorem
would have to prove that the actual route discrepancies lie in its generated
subspace and then connect those algebraic moves to a strict physical
replacement.  This refutation does not affect weak L2 or the later profile
certificates, which make different claims.

One elementary combinatorial ingredient of the long-boundary branch is the
following sweep lemma.

First isolate its finite part from the geometric assertion that a particular
embedded graph supplies such stacks.

\begin{lemma}[bounded-stack sweep pigeonhole]
\label{lem:bounded-stack-sweep-pigeonhole}
\Sverified{Mettapedia.GraphTheory.FourColor.GoertzelV24NoncrossingSweepPigeonhole.exists_distinct_eq_sweepState}
Let $A$ and $K$ be finite sets.  At each of $n$ positions record one element
of $A$ and, for each of $r$ families, a $K$-labelled stack of depth at most
$D$.  If
\[
 n>|A|\left(\sum_{d=0}^{D}|K|^d\right)^r,
\]
then two distinct positions carry exactly the same record.
\end{lemma}

\begin{proof}
A stack is encoded by its exact depth $d\in\{0,\ldots,D\}$ and a word in
$K^d$.  Thus the stack carrier has exactly
$\sum_{d=0}^{D}|K|^d$ elements, and the full record carrier has exactly
$|A|(\sum_{d=0}^{D}|K|^d)^r$ elements.  The conclusion is the ordinary
pigeonhole principle.  The Lean statement uses the literal dependent sum
$\Sigma_{d:\operatorname{Fin}(D+1)}\operatorname{Vector}(K,d)$, proves this
cardinality exactly, and returns the two unequal indices together with their
equal states.
\end{proof}

\begin{lemma}[noncrossing length--depth alternative]
\label{lem:noncrossing-length-depth}
\Sverified{Mettapedia.GraphTheory.FourColor.GoertzelV24NoncrossingSweepLifo.SweepData.exists_deep_family_or_distinct_eq_rawState}
Let $n$ ordered boundary positions carry a letter from a finite alphabet $A$
and, for each of $r$ families, a labelled noncrossing partial matching whose
arc labels lie in a finite set $K$.  Fix $D\ge0$ and put
\[
 S=|A|\left(\sum_{d=0}^{D}|K|^d\right)^r.
\]
If $n>S$, then either some matching has nesting depth greater than $D$, or two
cut positions have identical sweep states (the local letter together with the
ordered stacks of currently open labelled arcs in every family).
\end{lemma}

\begin{proof}
Sweep the boundary in its given order.  Noncrossing means that the open arcs
of one family close in last-in--first-out order, so its complete unresolved
connectivity at a cut is an ordered stack.  If every nesting depth is at most
$D$, one stack has at most
$\sum_{d=0}^{D}|K|^d$ possible values.  Together with the local letter and the
$r$ independent stacks, Lemma~\ref{lem:bounded-stack-sweep-pigeonhole} gives
two equal sweep states as soon as there are more than $S$ positions.  If the
depth bound fails, the first alternative holds.

In the formal statement, each arc has two ordered endpoints and a label;
distinct arcs are endpoint-disjoint and arcs are listed by increasing left
endpoint.  For two arcs open at the same cut, a later left endpoint followed
by a later right endpoint would be a crossing.  Equality of the right
endpoints is excluded by endpoint-disjointness.  Hence their right endpoints
strictly decrease, which is the claimed LIFO order.  The resulting raw lists
are encoded by the exact stack carrier of
Lemma~\ref{lem:bounded-stack-sweep-pigeonhole}; equality of codes is then
transported back to literal equality of the local letter and every stack.
\end{proof}

\begin{lemma}[restricted noncrossing sweep]
\label{lem:restricted-noncrossing-sweep}
\Sverified{Mettapedia.GraphTheory.FourColor.Compositional.RestrictedNoncrossingSweep.SweepData.exists_deep_family_on_or_distinct_eq_rawState}
\Sverified{Mettapedia.GraphTheory.FourColor.Compositional.RestrictedNoncrossingSweep.SweepData.exists_deep_family_on_or_spaced_eq_rawState}
In Lemma~\ref{lem:noncrossing-length-depth}, let $P$ be any selected finite
set of cut positions.  If
\[
 |P|>|A|\left(\sum_{d=0}^{D}|K|^d\right)^r,
\]
then either some selected position has a matching stack deeper than $D$, or
two distinct selected positions have the same raw sweep state.  After
adjoining the position phase modulo $s+1$, the corresponding threshold is
\[
 |P|>|A|(s+1)\left(\sum_{d=0}^{D}|K|^d\right)^r,
\]
and the equal states can be ordered at positions whose difference is at
least $s+1$.
\end{lemma}

\begin{proof}
Restrict the state-coding function to the subtype of positions belonging to
$P$.  If the deep alternative fails on $P$, every selected stack has depth
at most $D$, so the same exact bounded-stack carrier used in
Lemma~\ref{lem:noncrossing-length-depth} encodes all selected states.  The
pigeonhole principle on the finite subtype $P$ gives two unequal selected
positions with equal codes, and decoding gives literal equality of raw
states.  For the phased version, equal phases make the positive coordinate
difference a multiple of $s+1$.
\end{proof}

\begin{corollary}[physical residual-return length--depth alternative]
\label{cor:physical-residual-return-length-depth}
\Sverified{Mettapedia.GraphTheory.FourColor.Compositional.NoncrossingPairingSweep.noncrossingMatchingOfPairing}
\Sverified{Mettapedia.GraphTheory.FourColor.Compositional.ResidualReturnSweep.returnShoreMatching_pairwise_noncrossing}
\Sverified{Mettapedia.GraphTheory.FourColor.Compositional.ResidualReturnSweep.exists_deep_return_shore_or_distinct_eq_rawState}
\Sverified{Mettapedia.GraphTheory.FourColor.Compositional.PhasedNoncrossingSweep.SweepData.exists_deep_family_or_spaced_eq_rawState}
\Sverified{Mettapedia.GraphTheory.FourColor.Compositional.ResidualReturnSweep.exists_deep_return_shore_or_spaced_eq_rawState}
\Sverified{Mettapedia.GraphTheory.FourColor.Compositional.NoncrossingPairingSweep.openArcs_pairwise_strictlyNested}
\Sverified{Mettapedia.GraphTheory.FourColor.Compositional.ResidualReturnSweep.exists_long_strictlyNested_return_family_of_deep}
\Sverified{Mettapedia.GraphTheory.FourColor.Compositional.ResidualReturnSweep.exists_long_strictlyNested_return_family_or_spaced_eq_rawState}
\Sverified{Mettapedia.GraphTheory.FourColor.Compositional.NoncrossingPairingSweep.exists_canonical_position_of_mem_canonicalArcList}
\Sverified{Mettapedia.GraphTheory.FourColor.Compositional.ResidualReturnNestedFamily.strictlyNestedReturnFamily_of_deep_stack}
\Sverified{Mettapedia.GraphTheory.FourColor.Compositional.ResidualReturnNestedFamily.strictlyNestedReturnTriple_of_deep_stack}
\Sverified{Mettapedia.GraphTheory.FourColor.Compositional.ResidualReturnNestedFamily.strictlyNestedReturnTriple_of_nestedCarrierDeepReturnStack}
\Sverified{Mettapedia.GraphTheory.FourColor.Compositional.ResidualReturnNestedFamily.strictlyNestedReturnQuadruple_of_deep_stack}
Let a proper alternating site in a graph-backed least counterexample have
$n$ cyclic positions and the physical residual-return receipt of
Lemma~\ref{lem:residual-return-two-sector-noncrossing}.  Fix $D\geq0$.
If
\[
                         n>2(D+1)^2,
\]
then either one face shore has a residual-return stack of depth greater than
$D$, or two distinct cyclic positions have the same shore label and the
same two ordered stacks of currently open residual returns.

More generally, fix a requested spacing $s\geq0$.  If
\[
                    n>2(s+1)(D+1)^2,
\]
then the deep-stack alternative holds, or there are positions $i<j$ with
\[
                         j-i\geq s+1
\]
which have the same shore label, the same two return stacks, and the same
position phase modulo $s+1$.

The deep alternative is literal: it supplies more than $D$ canonical
residual-return chords on one shore, all open across one common cut.  In
their increasing-left-endpoint order, the right endpoints strictly decrease;
hence these chords are pairwise strictly nested.
\end{corollary}

\begin{proof}
For each return pair, retain its smaller cyclic coordinate and sort these
canonical representatives.  The partner involution makes distinct pairs
endpoint-disjoint.  Lemma~\ref{lem:residual-return-two-sector-noncrossing}
says that the returns on either fixed Boolean shore are noncrossing, so the
two sorted families are literal noncrossing matchings.  Sweep them
simultaneously around the carrier and apply
Lemma~\ref{lem:noncrossing-length-depth} with local alphabet
$A=\{0,1\}$, two matching families, and the one-element arc-label set
$K=\{*\}$.  A stack of depth at most $D$ then has $D+1$ possible
values, and the complete raw record has at most $2(D+1)^2$ values.

For the spaced form, augment the local letter by the cyclic coordinate
modulo $s+1$.  This multiplies the state count by $s+1$.  Equality of two
phases says $i\equiv j\pmod{s+1}$; since $i<j$, their difference is a
positive multiple of $s+1$ and hence at least $s+1$.

In the deep branch, filter the selected shore's canonical chord list by the
arcs open across the witnessing cut.  Its length is exactly the corresponding
stack depth.  Left endpoints remain strictly increasing under filtering, and
noncrossing plus endpoint-disjointness forces right endpoints to be strictly
decreasing.  Thus the filtered list is the asserted explicit nested family.

The formal adapter constructs the sorted arc lists directly from the
fixed-point-free return involution.  Its endpoint-disjointness proof uses
only involutivity, while noncrossing is supplied by the exact spherical
face-side argument.  Since the arc label is the one-element type, equality
of the two literal stacks is equivalently equality of their two depths; it
does not yet record the colouring support or the physical connectivity state
needed for a splice.
\end{proof}

\begin{lemma}[pushing an exact cycle cut to a physical vertex side]
\label{lem:exact-cycle-cut-push-off}
\Sverified{Mettapedia.GraphTheory.Embedding.ExactFaceCut.filledCycleSide_iff_label_of_not_mem_support}
\Sverified{Mettapedia.GraphTheory.Embedding.ExactFaceCut.filledCycleSide_iff_of_edge_endpoints_not_mem_support}
\Sverified{Mettapedia.GraphTheory.FourColor.Compositional.CyclePushOffCut.mem_edges_implies_off_cycle_and_incident}
\Sverified{Mettapedia.GraphTheory.FourColor.Compositional.CyclePushOffCut.card_edges_le_support_mul_degreeBound}
\Sverified{Mettapedia.GraphTheory.FourColor.Compositional.CyclePushOffCut.filledCycleSide_iff_of_walk_avoids_support}
\Sverified{Mettapedia.GraphTheory.FourColor.Compositional.CyclePushOffCut.exists_selected_with_complement_cycle_of_disjoint_cycle}
\Sverified{Mettapedia.GraphTheory.FourColor.Compositional.CyclePushOffCut.hasCycleOnSide_filledCycleSide}
\Sverified{Mettapedia.GraphTheory.FourColor.Compositional.CyclePushOffCut.cyclicEdgeCutRealization_of_complement_cycle}
Let $W$ be a closed primal walk in a graph-backed rotation system, and let
$\lambda$ be an exact binary face cut supported on the edges of $W$.  Fix one
label $\epsilon\in\mathbb F_2$.  Define the filled vertex side $U_\epsilon$
by putting every vertex of $W$ in $U_\epsilon$ and, away from $W$, putting a
vertex in $U_\epsilon$ exactly when an incident dart outside $W$ sees face
label $\epsilon$.  If the vertex rotations are cyclic, then this definition
is independent of the chosen off-walk dart.  Moreover every edge crossing
$U_\epsilon$ is not an edge of $W$ and has an endpoint on $W$.  If every
vertex of $W$ has degree at most $\Delta$, the number of crossing edges is at
most $\Delta\,|V(W)|$.
\end{lemma}

\begin{proof}
At a vertex outside $W$, no incident edge belongs to the support of $W$.
Rotating from one incident dart to the next therefore never crosses the exact
face cut, so all incident dart faces have the same label.  If an edge has
both endpoints outside $W$, the edge itself is also absent from the cut;
crossing it preserves the label, and hence preserves membership in
$U_\epsilon$.  Such an edge cannot cross the filled vertex side.  An edge of
$W$ cannot cross either, because both of its endpoints were filled into
$U_\epsilon$.  Thus every crossing edge is an off-walk edge with at least one
endpoint on $W$.  All crossing edges therefore lie in the union of the
incident-edge stars of the vertices of $W$.  Subadditivity of cardinality
gives the displayed $\Delta\,|V(W)|$ bound.

The Lean development separates the generic face-to-vertex conversion from
the four-colour consumer.  The robust local-degree estimate is sufficient
for finiteness; in a cubic simple cycle it can later be sharpened from the
factor $3$ to $1$, since two incident edges are already used by the cycle.
One side cycle is automatic: the filled side contains $W$ itself whenever
$W$ is simple.  Consequently a cycle in the complementary filled side is the
single remaining input needed to package the pushed-off edge set as a cyclic
edge-cut realization.
\end{proof}

\begin{corollary}[bounded push-offs of a residual-return separator]
\label{cor:residual-return-bounded-push-off}
\Sverified{Mettapedia.GraphTheory.FourColor.Compositional.ResidualReturnTransversal.exists_exactFaceCut_with_bounded_pushOff_edges}
For every physical residual-return chord whose selected carrier interval has
an interior edge, the resulting simple return separator has an exact binary
face cut.  Pushing either binary label to the filled vertex side gives a
literal finite edge interface of cardinality at most three times the number
of vertices in the separator support.
\end{corollary}

\begin{proof}
The return path and the reverse carrier interval form the simple cycle of
Lemma~\ref{lem:residual-return-two-sector-noncrossing}.  Its exact face cut is
the checked spherical closed-trail construction.  Apply
Lemma~\ref{lem:exact-cycle-cut-push-off}; cubicity bounds every incident-edge
star by three.  The statement applies independently to the two values of
$\mathbb F_2$.

This corollary constructs the bounded physical edge set rather than accepting
one as a receipt.  For a single return, promotion to a cyclic edge-cut
realization still asks for a cycle in the complementary filled side; the
selected side already contains the return separator itself.  The next
corollary supplies that opposite cycle from a strictly nested pair.
\end{proof}

\begin{corollary}[strictly nested residual returns give a degree-controlled cyclic cut]
\label{cor:nested-residual-return-cyclic-cut}
\Sverified{Mettapedia.GraphTheory.FourColor.Compositional.ResidualReturnComplementaryCycle.complementaryCycleInterval_support}
\Sverified{Mettapedia.GraphTheory.FourColor.Compositional.ResidualReturnComplementaryCycle.cycleVertexOrder_not_mem_complementaryCycleInterval_support}
\Sverified{Mettapedia.GraphTheory.FourColor.Compositional.ResidualReturnComplementaryCycle.complementaryReturnCycle_length_alternative}
\Sverified{Mettapedia.GraphTheory.FourColor.Compositional.ResidualReturnComplementaryCycle.orderedReturnSeparator_length_alternative}
\Sverified{Mettapedia.GraphTheory.FourColor.Compositional.ResidualReturnComplementaryCycle.orderedReturnSeparator_isCycle_automatic}
\Sverified{Mettapedia.GraphTheory.FourColor.Compositional.ResidualReturnComplementaryCycle.complementaryReturnCycle_support_disjoint_orderedReturnSeparator}
\Sverified{Mettapedia.GraphTheory.FourColor.Compositional.ResidualReturnCyclicCut.exists_bounded_cyclicEdgeCut_of_strictly_nested_returns}
Let two physical residual-return chords be strictly nested in cyclic
coordinates,
\[
                  \ell_o<\ell_i<r_i<r_o.
\]
Then the exact face cut of the inner return separator pushes to a cyclic
edge-cut realization.  Its crossed-edge set has cardinality at most three
times the number of vertices in the inner separator support.
\end{corollary}

\begin{proof}
Rotate the carrier cycle to start at $\ell_o$.  The carrier interval
complementary to the outer chord is exactly the suffix beginning at $r_o$
followed by the prefix ending at $\ell_o$.  Consequently every carrier
coordinate strictly between $\ell_o$ and $r_o$ is absent from that
complementary interval.

Both closures are automatically nondegenerate.  In each case the physical
return and the chosen carrier interval are nonempty.  If both pieces of one
closure had exactly one edge, simplicity of the ambient graph would make
them the same edge; this is impossible because a residual return avoids
every carrier edge.

Close the outer physical return along this complementary interval, and close
the inner physical return along the reverse interval from $\ell_i$ to $r_i$.
Distinct residual-return paths have disjoint supports, and each such path
meets the carrier only at its two endpoints.  Since both endpoints and the
whole carrier interval of the inner return lie strictly between
$\ell_o$ and $r_o$, the two resulting simple cycles are vertex-disjoint.

The generic closed-trail face-cut theorem gives the inner separator an exact
binary cut.  Choose its label opposite to the face label read by a first dart
of the disjoint outer cycle.  Exact-cut transport propagates this choice
around the outer cycle, so the outer cycle lies in the complementary filled
side; the inner separator lies in the filled side by construction.  Thus the
pushed edge set is a cyclic cut, and the cubic incidence estimate of
Corollary~\ref{cor:residual-return-bounded-push-off} gives the asserted
factor-three bound.
\end{proof}

\begin{corollary}[geometric residual-return sweep alternative]
\label{cor:geometric-residual-return-sweep-alternative}
\Sverified{Mettapedia.GraphTheory.FourColor.HasCyclicEdgeCutOfSizeAtMost}
\Sverified{Mettapedia.GraphTheory.FourColor.Compositional.ResidualReturnSweepCyclicCut.exists_physical_chord_of_mem_openArcs}
\Sverified{Mettapedia.GraphTheory.FourColor.Compositional.ResidualReturnSweepCyclicCut.hasDegreeControlledCyclicReturnCut_of_nested_openArcs}
\Sverified{Mettapedia.GraphTheory.FourColor.Compositional.ResidualReturnSweepCyclicCut.hasDegreeControlledCyclicReturnCut_of_deep_return_shore}
\Sverified{Mettapedia.GraphTheory.FourColor.Compositional.ResidualReturnSweepCyclicCut.hasDegreeControlledCyclicReturnCut_or_spaced_eq_rawState}
\Sverified{Mettapedia.GraphTheory.FourColor.Compositional.ResidualReturnSweepCyclicCut.HasNestedReturnSeparatorLargerThan}
\Sverified{Mettapedia.GraphTheory.FourColor.Compositional.ResidualReturnSweepCyclicCut.hasCyclicEdgeCutOfSizeAtMost_or_longReturnSeparator_or_spaced_eq_rawState}
\Sverified{Mettapedia.GraphTheory.FourColor.Compositional.ReturnSeparatorLength.support_toFinset_card_append_le}
\Sverified{Mettapedia.GraphTheory.FourColor.Compositional.ReturnSeparatorLength.orderedReturnSeparator_support_toFinset_card_le}
\Sverified{Mettapedia.GraphTheory.FourColor.Compositional.ReturnSeparatorLength.residualCycleInterval_support_toFinset_card}
\Sverified{Mettapedia.GraphTheory.FourColor.Compositional.ReturnSeparatorLength.long_ambientReturn_or_long_carrierInterval}
\Sverified{Mettapedia.GraphTheory.FourColor.Compositional.ResidualReturnSweepCyclicCut.HasNestedAmbientReturnSupportLargerThan}
\Sverified{Mettapedia.GraphTheory.FourColor.Compositional.ResidualReturnSweepCyclicCut.HasNestedCarrierIntervalSupportLargerThan}
\Sverified{Mettapedia.GraphTheory.FourColor.Compositional.ResidualReturnSweepCyclicCut.hasNestedAmbientReturnSupportLargerThan_or_carrierInterval_of_separator}
\Sverified{Mettapedia.GraphTheory.FourColor.Compositional.ResidualReturnSweepCyclicCut.hasCyclicEdgeCutOfSizeAtMost_or_longAmbientReturn_or_longCarrierInterval_or_spaced}
Fix a requested component bound $B\geq0$ and spacing $s\geq0$.  If a proper
alternating site has $n$ cyclic positions and
\[
                              n>8(s+1),
\]
then one of the following holds:
\begin{enumerate}
\item the site has a realized cyclic edge cut of width at most $6B$;
\item two strictly nested physical returns on one facial shore have inner
      ambient-return support larger than $B$;
\item two strictly nested physical returns on one facial shore have inner
      carrier-interval support larger than $B$;
\item there are positions $i<j$ with $j-i\geq s+1$ carrying the same phased
      two-shore sweep state.
\end{enumerate}
Thus the only obstruction to converting the deep sweep horn into a uniformly
bounded interface is now separated into a long path through the ambient graph
and a long interval along the already ordered carrier.
\end{corollary}

\begin{proof}
Apply the spaced physical length--depth alternative with depth $D=1$.
Its shallow horn is exactly the asserted repeated-state conclusion.  In its
deep horn one shore has at least two open arc records.  Each record in the
canonical sweep list remembers the canonical endpoint from which it was
constructed, so it lifts back to the corresponding physical return chord,
with both endpoints and the shore label unchanged.  The first two records in
the LIFO stack are strictly nested.  Corollary~\ref{cor:nested-residual-return-cyclic-cut}
therefore supplies a cyclic edge cut with at most three crossed edges for
each vertex of the inner separator.  That separator is the append of the
ambient return and the reverse carrier interval.  Its number of distinct
vertices is at most the sum of the two component supports.  If both component
supports have at most $B$ vertices, the separator has at most $2B$ and the cut
has width at most $6B$.  Otherwise the long component determines the second or
third horn.  The carrier component is a simple path, so its support cardinality
is exactly $r_i-\ell_i+1$.  Finally, $2(s+1)(1+1)^2=8(s+1)$.
\end{proof}

\begin{corollary}[carrier-local compression of a long nested return]
\label{cor:carrier-local-residual-return-compression}
\Sverified{Mettapedia.GraphTheory.FourColor.Compositional.ResidualReturnCarrierSweep.carrierIntervalPositions_card}
\Sverified{Mettapedia.GraphTheory.FourColor.Compositional.ResidualReturnCarrierSweep.exists_deep_return_shore_on_carrierInterval_or_spaced_eq_rawState}
\Sverified{Mettapedia.GraphTheory.FourColor.Compositional.ResidualReturnCarrierSweep.hasNestedCarrierDeepReturnStack_or_spacedSweepRepeat}
\Sverified{Mettapedia.GraphTheory.FourColor.Compositional.ResidualReturnCarrierSweep.hasCarrierCompressedGeometricAlternative}
Let $\ell<r$ be the endpoints of the inner member of a strictly nested
physical return pair.  Fix a stack bound $D$ and spacing $s$, and put
\[
                         S=2(s+1)(D+1)^2.
\]
If the carrier interval from $\ell$ through $r$ contains more than $S$
vertices, then, using only cut positions in that same closed interval,
either one return shore has stack depth greater than $D$, or two positions
$i<j$ carry the same phased two-shore state and satisfy $j-i\ge s+1$.

Consequently the geometric sweep alternative may take component bound $S$.
It then gives a cyclic edge cut of width at most $6S$, a nested ambient
return with more than $S$ vertices, a deeper stack inside the selected
carrier interval, a repeated phased state inside that interval, or a
repeated phased state on the full carrier.  There is no remaining
unstructured long-carrier horn.
\end{corollary}

\begin{proof}
The carrier interval is a simple path and has exactly $r-\ell+1$ distinct
vertices.  Its closed coordinate interval $[\ell,r]$ has the same
cardinality.  Apply Lemma~\ref{lem:restricted-noncrossing-sweep} to precisely
that finite interval, with the Boolean shore letter, two matching families,
and the one-element arc-label set.  The exact state count is
$2(s+1)(D+1)^2$.  Applying this local alternative to the long-carrier horn
of Corollary~\ref{cor:geometric-residual-return-sweep-alternative} gives the
last assertion.  The nested physical pair is retained throughout; only the
set of sweep positions is restricted.
\end{proof}

\begin{lemma}[cubic third-edge classification along an ambient return]
\label{lem:ambient-return-third-edge-classification}
\Sverified{Mettapedia.GraphTheory.CubicPathAttachment.ncard_attachmentNeighborSet_eq_one}
\Sverified{Mettapedia.GraphTheory.CubicPathAttachment.chordAttachment_or_externalAttachment}
\Sverified{Mettapedia.GraphTheory.FourColor.Compositional.ResidualReturnPathAttachment.ambientReturnChordAttachment_or_externalAttachment}
\Sverified{Mettapedia.GraphTheory.FourColor.Compositional.ResidualReturnPathAttachment.hasNestedAmbientReturnInternalPositionsMoreThan_of_support}
Let $P$ be the simple ambient path of a physical residual return in the cubic
graph.  Every strict internal position of $P$ has exactly one neighbour not
joined to it by either of the two path edges.  The corresponding third edge
either leaves the support of $P$, or is a chord whose other endpoint is a
nonconsecutive position of $P$.  Moreover, if $P$ has more than $B+2$
distinct support vertices, then it has more than $B$ strict internal
positions carrying these canonical third edges.

This is only the syntax of the ambient sweep.  It does not assert that the
external attachments lie on a common planar side, that its chords are
noncrossing, or that any set of these third edges is a bounded interface.
Those are geometric obligations, not consequences of cubicity.
\end{lemma}

\begin{proof}
Index the strict internal positions by
$\operatorname{Fin}(|P|-1)$, with coordinate $i$ denoting path coordinate
$i+1$.  The two path neighbours are distinct because $P$ is simple.  Their
neighbour set is contained in the ambient neighbour set; the former has
cardinality two and the latter has cardinality three.  Their set difference
therefore has cardinality one.  If its unique member belongs to the support
of $P$, simplicity gives its unique path coordinate, and exclusion from the
two path neighbours shows that coordinate is neither the predecessor nor the
successor.  Otherwise it is an external attachment.  Finally a simple path
has exactly two fewer strict internal positions than support vertices.
\end{proof}

\begin{lemma}[matching role and genuine exit of a residual-return attachment]
\label{lem:residual-return-attachment-matching-role}
\Sverified{Mettapedia.GraphTheory.FourColor.Compositional.ResidualReturnAttachmentMatching.ambientReturnAttachmentNeighbor_eq_partner}
\Sverified{Mettapedia.GraphTheory.FourColor.Compositional.ResidualReturnAttachmentMatching.ambientReturnPathEdge_partner_ne}
\Sverified{Mettapedia.GraphTheory.FourColor.Compositional.ResidualReturnAttachmentMatching.internalChord_partner_leftVertex_eq_rightVertex}
\Sverified{Mettapedia.GraphTheory.FourColor.Compositional.ResidualReturnAttachmentMatching.ambientReturnAttachmentNeighbor_not_mem_carrier}
\Sverified{Mettapedia.GraphTheory.FourColor.Compositional.ResidualReturnAttachmentMatching.ambientReturnExternalAttachment_leavesSeparator}
\Sverified{Mettapedia.GraphTheory.FourColor.Compositional.ResidualReturnAttachmentMatching.AttachmentExitSideReceipt.faceSide}
\Sverified{Mettapedia.GraphTheory.FourColor.Compositional.ResidualReturnPairingChordRole.ambientReturnPathEdge_not_pairingGraph}
\Sverified{Mettapedia.GraphTheory.FourColor.Compositional.ResidualReturnPairingChordRole.turnChord_pairingGraph_adj}
At every strict internal position of a physical residual return, the canonical
third edge is the edge of the reference matching.  Every edge traversed by the
return path lies outside that matching, while every internal chord of the path
is exactly the reference-matching edge at either endpoint.  Thus each
same-turn chord stack used below is a family of actual matching edges, rather
than merely a coordinate diagram.  The other endpoint of an attachment cannot
belong to the operated alternating carrier.  Consequently, if that edge is
external to the return path, it genuinely leaves the complete separator made
from the return and its complementary carrier interval.  The exact face-side
branch of its separator-exit receipt therefore always holds; the former
``returns to the carrier'' alternative is impossible.
\end{lemma}

\begin{proof}
The return path lies in the common residual graph before and after the
alternating exchange.  Its two path edges at an internal vertex are therefore
different from the reference-matching edge.  Cubicity leaves exactly one
third edge, so its neighbour is the reference partner.  Applying this identity
to the canonical left attachment of any internal path chord identifies the
chord's right endpoint with that partner.  Filtering the chord diagram by
local turn preserves this identity and proves the same statement for every
member of either turn stack.  If the partner lay on the carrier, closure of
the carrier under the reference matching and involutivity would put the
internal path vertex on the carrier.  But a simple physical return meets the
carrier only at its two endpoints, contradicting strict internality.  The
earlier separator alternative now has only its face-side branch.
\end{proof}

\begin{lemma}[innermost residual chords drain to the opposite turn]
\label{lem:residual-return-innermost-chord-drainage}
\Sverified{Mettapedia.GraphTheory.exists_innermost_orderedPathChord}
\Sverified{Mettapedia.GraphTheory.OrderedPathChord.crosses_of_hasEndpointInside_of_innermost}
\Sverified{Mettapedia.GraphTheory.FourColor.Compositional.ResidualReturnInnermostChordDrainage.exists_innermostChordDrainage}
\Sverified{Mettapedia.GraphTheory.FourColor.Compositional.ResidualReturnInnermostChordDrainage.external_or_oppositeTurnChord_of_eligible_inside}
\Sverified{Mettapedia.GraphTheory.FourColor.Compositional.ResidualReturnInnermostChordDrainage.exists_exitSideReceipt_or_all_eligible_inside_drain}
\Sverified{Mettapedia.GraphTheory.card_orderedPathChordEndpoints_le_two_mul}
\Sverified{Mettapedia.GraphTheory.FourColor.Compositional.ResidualReturnInnermostChordDrainage.card_eligibleCoordinatesInside_le_eight_mul_of_stackBounds}
\Sverified{Mettapedia.GraphTheory.CubicPathChordDiagram.OrderedPathChord.span_le_card_internalNonendpointCoordinates_inside_add_five}
\Sverified{Mettapedia.GraphTheory.FourColor.Compositional.ResidualReturnInnermostChordDrainage.exitReceipt_or_innermostChord_span_le_eight_mul_depth_add_five}
If the full internal-chord diagram of a physical residual return is nonempty,
it contains an innermost chord $c$.  This chord is an actual edge of the
reference matching.  Every distinct internal chord having an endpoint on the
open path interval cut off by $c$ crosses $c$, and therefore has the opposite
local attachment turn.  More precisely, every sweep-eligible coordinate in
that open interval either has a third edge leaving the ambient return path or
is incident to such an opposite-turn chord.  The eligibility condition removes
exactly the bounded endpoint-return exceptions already accounted for by the
sweep budget.  Uniformly over the interval, either one such external edge
supplies an exact face-side separator-exit receipt, or every eligible
coordinate drains through an opposite-turn chord.

This alternative is quantitative.  Suppose that, at each endpoint of the
selected chord, each of the two turn stacks has depth at most $d$.  In the
no-exit branch every eligible coordinate strictly inside the chord is an
endpoint of a chord open at one of those two endpoints.  A finite chord family
uses at most two endpoint coordinates per chord, while the four open stacks
contain at most $4d$ chords in total.  Hence the eligible interior contains at
most $8d$ coordinates.  Cubicity supplies a separate sharp bound of four on
the discarded endpoint-return attachments in any path interval.  Consequently
the selected chord itself has span at most $8d+5$, unless the exact separator
exit receipt exists.

This is the exact residual-return analogue of the combinatorial drainage core
of the source's Sector-Alternation Lemma.  It does not by itself construct the
source's named annular holes or prove that one of the two regions bounded by
$c$ is hole-free; those are global realization fields, not consequences of
the residual pairings alone.
\end{lemma}

\begin{proof}
Choose $c$ with minimum endpoint span.  A chord strictly inside its interval
would have smaller span.  Cubicity and simplicity make the internal chords an
endpoint-disjoint partial matching.  Hence any other chord entering the open
interval can neither be nested there nor share an endpoint with $c$; the
linear order of the four endpoints forces it to cross $c$.  The exact
face-cut crossing theorem for residual-return chords then gives unequal local
turns.  Finally
Lemma~\ref{lem:residual-return-attachment-matching-role} identifies $c$ with
the reference-matching edge at its left endpoint.
\end{proof}

\begin{lemma}[one-path laminar noncrossing has no uniform depth bound]
\label{lem:one-path-laminar-noncrossing-unbounded}
\Sverified{Mettapedia.GraphTheory.card_rainbowChordFamily}
\Sverified{Mettapedia.GraphTheory.rainbowChordFamily_pairwiseEndpointDisjoint}
\Sverified{Mettapedia.GraphTheory.rainbowChordFamily_pairwiseLaminar}
\Sverified{Mettapedia.GraphTheory.exists_constant_sector_noncrossing_family}
\Sverified{Mettapedia.GraphTheory.no_uniform_cardinality_bound_of_laminar_fiberwiseNoncrossing}
For every $d$ there is an endpoint-disjoint family of exactly $d$ ordered
chords on one path whose intervals are strictly nested, whose chords are
pairwise noncrossing, and whose Boolean sector label is constant.  Therefore
endpoint disjointness, laminarity, and same-sector noncrossing do not imply a
uniform bound on either the family cardinality or its nesting depth.
\end{lemma}

\begin{proof}
On a path with $2d+1$ positions, join position $i$ to position $2d-i$ for
$0\le i<d$.  Increasing $i$ moves the left endpoint strictly rightward and
the right endpoint strictly leftward, so every pair is nested.  The two
halves of the coordinate interval are disjoint, so distinct chords share no
endpoint.  Nested chords do not cross.  Giving every chord the same Boolean
label supplies the final claim, and choosing $d$ one larger than any proposed
uniform bound gives a contradiction.
\end{proof}

The abstract sharpness result alone does not decide the source's stronger
two-rail drainage premise.  That premise has therefore been tested directly
on the annular specimens cited by the source.

\begin{proposition}[small-case falsification gate for universal chord drainage]
\label{prop:universal-chord-drainage-refuted-on-census}
\Srefuted{Mettapedia.GraphTheory.ClosedWebSectorNestingWitness.universalDrainage_false}
\Srefuted{Mettapedia.GraphTheory.GoodWordClosedWebNestingWitness.universalDrainage_false}
\Srefuted{Mettapedia.GraphTheory.GoodWordRealGraphNestingWitness.universalDrainage_false}
The statement that \emph{every} same-rail chord drains all of its strict
interior third-colour attachments to the other rail or to a hole is false on
the cited closed-web census annuli.  In the $C_{30}$ census there are $5412$
proper colourings, $360$ totally closed webs, and $240$ two-rail webs.  Those
two-rail webs contain $360$ same-rail chords and $240$ same-sector nested
pairs, with nesting depth three.  In the $C_{40}$ census the corresponding
figures are $31704$, $720$, $480$, $1440$, and $960$, again with depth three.

On the $C_{30}$ two-rail webs, drainage holds for all $180$ innermost chords
and fails for all $180$ chords which contain another chord.  A transparent
certificate has the strict same-sector nest
\[
       (19,21)\supset(12,20)\supset(14,23).
\]
For the outer chord $(19,21)$, the third-colour edge $(12,20)$ begins and ends
strictly inside the outer interval in the same face sector.  The standalone
checker
\texttt{lean/mettapedia/scripts/validation\_lab/verify\_sector\_alternation\_depth3.py}
reconstructs proper Tait colouring, total closure, the two rails, the two
face sectors, the depth-three nest, and this explicit drainage failure from
the accompanying JSON certificate.

The polar census witness has boundary colour multiplicities $(3,2,0)$, not
the good-word type $(3,1,1)$, and the polar $C_{30}$ and $C_{40}$ censuses in
fact contain no good-word closed web.  The target population is nevertheless
also refuted by a second kernel witness: a spherical width-two annular tangle
with good inner word $(3,1,1)$ has same-side nested chords
$(5,13)\supset(7,11)$ and fails universal drainage.  This is not merely an
abstract-tangle pathology.  In the non-polar $C_{30}$ cap-pair sweep, all
$120$ closed webs with the bare $(3,1,1)$ multiplicity also satisfy the
cyclic consecutivity required of a good word; $60$ carry chords, and all $60$
carry a same-sector nested pair of depth two.  Thus universal drainage fails
on real graph specimens satisfying the source's full good-word predicate.

The source's separate raw count of $460$ chords is not reproduced by this
audit.  Any raw statistic invariant under global permutation of the three
nonzero colours has orbit sizes divisible by six on these proper colourings,
whereas $460$ is not divisible by six.  Thus that number uses an unstated
filter or quotient, or is a transcription error; it cannot be used as an
unqualified raw count.
\end{proposition}

The failure is localized.  The checked arbitrary-chord construction still
gives an exact chord cycle with a hole-free side,
\Sverified{Mettapedia.GraphTheory.FourColor.GoertzelV24ClosedWebChordHoleSideCutWitness.exists_exact_chordCycle_faceCut_with_holeFreeSide}
and every strict internal path position still has its canonical third-colour
edge,
\Sverified{Mettapedia.GraphTheory.FourColor.GoertzelV24ClosedWebRadialPathChords.exists_thirdColor_external_port_of_interior}.
What fails is the extra inference that an edge not lying on the chord-cycle
boundary must have its far endpoint outside the selected disk: a nested
chord is not a boundary edge, does not cross the outer chord, and remains
inside the disk.

For a minimum-span chord the prohibited nested case is absent, and the
checked theorem still supplies uniform drainage,
\Sverified{Mettapedia.GraphTheory.FourColor.GoertzelV24ClosedWebInnermostChordDrainage.exists_innermost_exact_holeFreeSide_with_uniform_drainage}.
It is no longer a route keystone.  Iterating it would again require a bound on
same-side nesting or a uniform interface between peeled disks, precisely the
content refuted by the good-word and real-graph witnesses.  The sole
closed-web factorization consumer is conditional on the independently
refuted global premise (U), while the surviving five-outer-stub
closed-to-menu-B theorem uses no chord layer.  The closed-web deep branch is
therefore retired: no innermost-peeling, drainage-layer, or nesting-bound
obligation is inherited by the direct Count route.

\begin{lemma}[rotation side of a cubic path attachment]
\label{lem:cubic-path-attachment-rotation-side}
\Sverified{Mettapedia.GraphTheory.FourColor.SimpleGraphDartRotation.Data.vertexRotation_apply_eq_second_or_third_of_regularThree}
\Sverified{Mettapedia.GraphTheory.FourColor.CubicPathRotation.rotationCycle_of_attachmentTurn_eq_backwardToAttachment}
\Sverified{Mettapedia.GraphTheory.FourColor.CubicPathRotation.rotationCycle_of_attachmentTurn_eq_forwardToAttachment}
\Sverified{Mettapedia.GraphTheory.FourColor.Compositional.ResidualReturnPathAttachment.ambientReturnAttachmentTurn_rotationCycle}
Orient a simple path and consider one strict internal position.  Write
$d_-$ and $d_+$ for the darts pointing to the preceding and following path
vertices, and write $d_a$ for the dart on the canonical third edge from
Lemma~\ref{lem:ambient-return-third-edge-classification}.  If the vertex
rotation is cyclic, then exactly one of the following local rotation cycles
holds:
\[
       d_-\longmapsto d_a\longmapsto d_+\longmapsto d_-,
       \qquad\hbox{or}\qquad
       d_-\longmapsto d_+\longmapsto d_a\longmapsto d_-.
\]
The corresponding two-valued attachment turn is therefore defined at every
strict internal position of every ambient residual-return path.

This turn is local rotation syntax.  In particular, the lemma does not yet
say that two attachments with the same turn lie on a common planar side, or
that chords with equal turns cannot interleave.  Those conclusions require a
separate face-side separation argument.
\end{lemma}

\begin{proof}
The predecessor, successor, and attachment darts are pairwise distinct and
have the same base vertex.  Cubicity says that they exhaust the three darts
at that vertex.  A cyclic permutation of a three-element fibre has no fixed
point, so the image of $d_-$ is either $d_a$ or $d_+$; after this first image
is chosen, the remaining two images are forced.  Apply this generic statement
to the simple ambient path witnessing the residual return.
\end{proof}

\begin{lemma}[exact face separator from an ambient-return chord]
\label{lem:ambient-return-chord-face-separator}
\Sverified{Mettapedia.GraphTheory.SamePathChordBoundary.cycleWalk_isCycle}
\Sverified{Mettapedia.GraphTheory.CubicPathAttachment.ChordAttachment.boundary}
\Sverified{Mettapedia.GraphTheory.Embedding.exists_exactFaceCut_of_samePathChordBoundary}
\Sverified{Mettapedia.GraphTheory.FourColor.Compositional.ResidualReturnPathAttachment.ambientReturnChordBoundary_cycleWalk_isCycle}
\Sverified{Mettapedia.GraphTheory.FourColor.Compositional.ResidualReturnChordSeparator.exists_exactFaceCut_of_ambientReturnChord}
Suppose the canonical third edge at an internal position of an ambient
residual-return path comes back to a nonconsecutive position of that path.
Order its two path coordinates and join the third edge to the intervening path
interval.  Their union is a simple cycle.  In a connected cellular sphere map
there is consequently an exact binary face labeling that changes value across
precisely the edges of this chord and path interval.

This separator is pointwise.  The lemma neither identifies which local turn
is the inside of the cycle nor asserts that separators produced by distinct
third edges are nested, disjoint, or a bounded interface.
\end{lemma}

\begin{proof}
The third edge cannot join consecutive path positions, since it is the unique
edge outside the two path edges at its initial position.  The interval of a
simple path between the ordered endpoints is again simple, has length greater
than one, and does not contain the third edge.  Traversing the chord and
returning along the reverse interval therefore gives a simple cycle.  Apply
the exact face-cut theorem for closed trails to this cycle and use that its
edge set is exactly the chord together with the interval edges.
\end{proof}

\begin{corollary}[pointwise geometric sweep alternatives across the mesh]
\label{cor:pointwise-mesh-residual-sweep-alternative}
\Sverified{Mettapedia.GraphTheory.FourColor.Compositional.MeshResidualReturnSweep.exists_geometricSweepReceipt_at_every_globalMeshStep}
\Sverified{Mettapedia.GraphTheory.FourColor.Compositional.MeshResidualReturnSweep.exists_minimizer_with_geometricSweep_at_every_internal_row_junction}
\Sverified{Mettapedia.GraphTheory.FourColor.Compositional.MeshResidualReturnSweep.carrierCompressedAlternative_of_geometricSweepReceipt}
\Sverified{Mettapedia.GraphTheory.FourColor.Compositional.AmbientReturnAttachmentCompression.hasNestedAmbientAttachmentSweepAlternative}
\Sverified{Mettapedia.GraphTheory.FourColor.Compositional.AmbientReturnAttachmentCompression.hasFullyCompressedGeometricAlternative}
\Sverified{Mettapedia.GraphTheory.FourColor.Compositional.MeshResidualReturnSweep.fullyCompressedAlternative_of_geometricSweepReceipt}
\Sverified{Mettapedia.GraphTheory.FourColor.Compositional.ExactStateGeometricAlternative.hasFullyCompressedExactStateAlternative}
\Sverified{Mettapedia.GraphTheory.FourColor.Compositional.MeshResidualReturnSweep.fullyCompressedExactStateAlternative_of_geometricSweepReceipt}
\Sverified{Mettapedia.GraphTheory.FourColor.Compositional.DeepReturnExactState.saturated_of_nestedCarrierDeepReturnStack}
\Sverified{Mettapedia.GraphTheory.FourColor.Compositional.ResidualReturnSeparatorNesting.orderedReturnSeparator_support_subset_filledCycleSide_of_triple_nested}
\Sverified{Mettapedia.GraphTheory.FourColor.Compositional.ResidualReturnNestedCuts.hasNestedConnectedReturnCuts_of_nestedCarrierDeepReturnStack}
\Sverified{Mettapedia.GraphTheory.FourColor.Compositional.ResidualReturnStrictNesting.shallow_ambientPath_outside_deep_closureSide}
\Sverified{Mettapedia.GraphTheory.FourColor.Compositional.ResidualReturnStrictNesting.incidentEdgeShore_deepClosure_ssubset_shallowClosure}
\Sverified{Mettapedia.GraphTheory.FourColor.Compositional.ResidualReturnNestedCuts.nestedConnectedReturnCutConstruction}
\Sverified{Mettapedia.GraphTheory.FourColor.Compositional.ResidualReturnStrictNesting.hasStrictlyNestedConnectedReturnCuts_of_strictlyNestedReturnQuadruple}
\Sverified{Mettapedia.GraphTheory.FourColor.Compositional.ResidualReturnStrictNesting.hasStrictlyNestedConnectedReturnCuts_of_nestedCarrierDeepReturnStack}
\Sverified{Mettapedia.GraphTheory.FourColor.Compositional.ResidualReturnNestedCuts.closureSide_subset_of_nested_return_separators}
\Sverified{Mettapedia.GraphTheory.FourColor.Compositional.ResidualReturnStrictFamily.hasStrictlyNestedConnectedReturnCutFamily_of_deep_stack}
\Sverified{Mettapedia.GraphTheory.FourColor.Compositional.ResidualReturnStrictFamily.exists_width_with_length_le_stateCount_of_deep_stack}
\Sverified{Mettapedia.GraphTheory.FourColor.Compositional.ResidualReturnStrictFamily.hasStrictlyNestedConnectedReturnCutFamily_of_nestedCarrierDeepReturnStack}
\Sverified{Mettapedia.GraphTheory.FourColor.Compositional.StrictFamilyGeometricAlternative.hasFullyCompressedStrictFamilyAlternative}
\Sverified{Mettapedia.GraphTheory.FourColor.Compositional.PathChordStrictFamily.hasStrictlyNestedChordFamily_of_stackAt_length_le}
\Sverified{Mettapedia.GraphTheory.FourColor.Compositional.AmbientReturnStrictChordFamily.realized_of_nestedAmbientAttachmentSweepAlternative}
\Sverified{Mettapedia.GraphTheory.FourColor.Compositional.StrictFamilyGeometricAlternative.hasFullyRealizedStrictFamilyAlternative}
\Sverified{Mettapedia.GraphTheory.FourColor.Compositional.NoncrossingSweepInterface.openArcInterfaceEquivOfPhasedRawStateEq}
\Sverified{Mettapedia.GraphTheory.FourColor.Compositional.ResidualReturnSweepInterface.SpacedResidualReturnInterfaceReceipt.openWireEquiv}
\Sverified{Mettapedia.GraphTheory.FourColor.Compositional.StrictFamilyGeometricAlternative.hasFullyWiredStrictFamilyAlternative}
\Sverified{Mettapedia.GraphTheory.FourColor.Compositional.DeepReturnExactState.hasFullyRealizedExactStateAlternative}
\Sverified{Mettapedia.GraphTheory.FourColor.Compositional.MeshResidualReturnSweep.fullyRealizedExactStateAlternative_of_geometricSweepReceipt}
For one ordered injective mesh in a graph-backed least counterexample, choose
the common minimum-residual-defect matching of
Corollary~\ref{cor:global-residual-geometry-provenance}.  Every global mesh
step is either central for this matching or carries its provenanced geometric
sweep alternative from
Corollary~\ref{cor:geometric-residual-return-sweep-alternative}, at every
requested component bound and spacing.  Moreover, at each internal row
junction at least one of the immediately incoming and outgoing arms carries
this finite alternative.  Its bounded-cut horn already returns a connected
literal exact-state shore.  If the carrier stack has depth at least two, its
deep horn likewise returns a connected cyclic subcut, retaining the fact that
it lies inside the physical return separator.
\end{corollary}

\begin{proof}
At a noncentral global step, retain the existing provenanced two-sector
receipt and apply the single-site geometric alternative to its facial bond.
This changes neither the site nor its supported matching, selected deletion
colouring, or return pairing.  At an internal row junction the common matching
cannot contain both incident row edges, because their outer endpoints are
distinct while their shared branch vertex has a unique partner.  Therefore
at least one arm is noncentral and carries the upgraded receipt.

The conclusion is pointwise.  It does not choose among the resulting
short-cut, attachment-exit, deep-stack, and repeated-state horns uniformly,
nor does it assert that choices made at adjacent junctions define nested or
disjoint shores.  Both potentially long coordinates are nevertheless
compressed on every receipt.  The carrier interval has the restricted LIFO
sweep of Corollary~\ref{cor:carrier-local-residual-return-compression}; the
ambient return path has a second sweep of its cubic third-edge attachments.
Thus neither survives as an unstructured long-path alternative.

The depth-two form of the carrier-stack horn is now physical as well as
combinatorial.  Its first three strictly nested open arcs lift through the
canonical sweep-to-chord bridge to three original residual-return chords on
one shore.  The middle and inner chords give exact face-cut separators; using
the outer complementary return cycle as a common exterior root and saturating
both cuts at one explicit common width produces connected shores with the
inner shore pointwise contained in the middle shore.  This removes the former
face-side comparison gap for a nested pair.

At stack depth four the next local obligation is also checked.  The outermost and
innermost returns orient the two middle separators.  Exact face-label
transport places the shallower middle return path inside the shallower rooted
shore but in the exterior component of the deeper rooted shore, even after
saturation.  Its first edge therefore witnesses strict inclusion of the two
incident-edge shores.  A provenance-retaining constructor now packages the
two saturated separators themselves, so the depth-three carrier horn returns
two actual \texttt{ConnectedAtWidth} receipts at one common bound with strict
inclusion of their incident-edge shores.

The same construction is now uniform over every finite prefix of a deep
stack.  Keep the first return as a common exterior root and the last as a
common orientation guard; each intervening return gives one connected cut.
A generic face-side comparison proves every later shore strictly contained
in every earlier shore.  Widening the finitely many receipts to the sum of
their local bounds gives one common width, and the resulting family feeds the
cumulative exact-state repetition theorem directly.  This closes family
coherence.  The remaining high-width obligation is quantitative: replace the
finite but graph-dependent common width by a bound controlled independently
of the ambient counterexample.

Composing this family construction with the both-coordinate compression
theorem removes the deep-stack predicate from the consumer-facing
alternative.  For every requested family length, a sufficiently long
residual site now yields either a bounded connected exact-state shore, a
strict family of that length, a structured ambient-attachment sweep, or one
of the two materially spaced repeated-state horns.  This is a qualitative
decomposition theorem: it leaves all three unresolved horns visible and does
not treat the finite, graph-dependent family width as a uniform bound.

The structured ambient horn has now also crossed its representation seam.
A generic graph-free theorem reconstructs the original ordered path chords
from the open arcs counted by a deep LIFO stack.  Applied to the two
same-turn attachment families, it returns a strictly nested family of actual
internal chords of the ambient residual path.  Consequently both deep-stack
outcomes of the fully compressed alternative now carry original graph
objects.  The external-attachment outcome is now a genuine separator exit
with an exact face-side receipt; the impossible return-to-carrier branch has
been removed by
Lemma~\ref{lem:residual-return-attachment-matching-role}.  That exit and the
two repeated-state outcomes remain explicit; in particular, equality of the
present unlabelled stack shapes is not asserted to be a colour-preserving
splice interface.

The repeated residual-return outcomes now expose their exact topological
content as well.  Equality of raw states identifies, family by family, the
cardinality of the ordered lists of open arcs, and therefore gives a
canonical equivalence of their ordered open-wire slots.  The carrier-local
receipt retains membership of both cuts in the same inner carrier interval.

The repeated ambient-attachment outcome now has the same proof-carrying
interface.  Its receipt retains the two eligible internal cut positions,
their prescribed material spacing, and equality of the phased raw sweep
states.  That equality induces a canonical equivalence between the two
ordered open-attachment interfaces.
\Sverified{Mettapedia.GraphTheory.FourColor.Compositional.AmbientReturnSweepInterface.SpacedAmbientReturnInterfaceReceipt.openWireEquiv}
Consequently all three repeated-sweep branches now expose ordered physical
wire slots rather than bare equalities of implementation states.

The ambient sweep now also controls the size of the corresponding physical
cuts.  At an internal cut position, every edge crossing the strict path
prefix belongs to one of four explicit classes: an edge incident to the
initial endpoint, an edge incident to the terminal endpoint, the forward
path edge, or an internal chord that is open at the cut.  The LIFO sweep
stack counts those open chords exactly.  Hence, if both same-turn stacks have
depth at most (d), the physical boundary has size at most (7+2d) at each
of the two repeated positions.
\Sverified{Mettapedia.GraphTheory.FourColor.Compositional.PathPrefixBoundary.card_crossingEdgeFinset_pathPrefixSide_le}
\Sverified{Mettapedia.GraphTheory.FourColor.Compositional.AmbientReturnSweepWidth.card_ambientReturnPrefixBoundary_le_stack_lengths}
\Sverified{Mettapedia.GraphTheory.FourColor.Compositional.AmbientReturnSweepBoundedRepeat.BoundedSpacedAmbientReturnInterfaceReceipt.firstBoundaryWidth}
\Sverified{Mettapedia.GraphTheory.FourColor.Compositional.AmbientReturnSweepBoundedRepeat.BoundedSpacedAmbientReturnInterfaceReceipt.secondBoundaryWidth}
Thus a sufficiently long ambient return yields an exact face-sided separator
exit, a strictly nested physical chord family, or a materially spaced
raw-state repeat whose two actual cut boundaries share one graph-independent
bound.
\Sverified{Mettapedia.GraphTheory.FourColor.Compositional.AmbientReturnBoundedAlternative.hasBoundedRealizedNestedAmbientAttachmentAlternative}

The second, semantic compression is now checked without identifying a raw
sweep state with a colouring interface.  First discard a margin of
\(k=7+2d\) positions at both ends.  Cubic tree counting promotes every
remaining shallow prefix to an actual cyclic cut.  Choose the last such cut,
take a cycle in its suffix component, and use that cycle and one exterior root
to saturate every earlier prefix.  The resulting connected shores are
strictly nested: the first forward path edge after an earlier cut belongs to
the later incident-edge shore but not the earlier one.
\Sverified{Mettapedia.GraphTheory.FourColor.Compositional.PathPrefixCyclicCut.hasCycleOnSide_prefixExteriorComponent}
\Sverified{Mettapedia.GraphTheory.FourColor.Compositional.PathPrefixRootedNesting.incidentEdgeShore_closureSide_ssubset}

In a vertex-minimal counterexample, two strictly nested literal shores cannot
have the same complete phased normalized Count profile.  Therefore a shallow
central family at width \(k\) has cardinality at most
\[
  N(k)=(6k+1)\sum_{j=0}^{k} j!\,2^{3^j},
\]
the exact cardinality of the already checked profile carrier.
\Sverified{Mettapedia.GraphTheory.FourColor.Compositional.AmbientReturnConnectedCutProfileBound.ShallowCentralCutFamilyReceipt.card_le_profileStateCount}
Consequently an ambient-return support larger than
\(6+N(7+2d)+2(7+2d)\) has an exact face-sided separator exit or contains
\(d+1\) strictly nested original internal chords; the shallow repeat branch
has disappeared rather than been postulated colour-compatible.
\Sverified{Mettapedia.GraphTheory.FourColor.Compositional.AmbientReturnFullProfileAlternative.exists_externalAttachment_or_strictlyNestedChordFamily}
The cardinality argument is path-generic: any finite family of bounded cubic
path prefixes with the same two margins admits the same common-root
saturation and complete-profile injection.
\Sverified{Mettapedia.GraphTheory.FourColor.Compositional.PathPrefixConnectedCutProfileBound.card_prefixFamily_le_profileStateCount}
The ambient-return theorem is now only the adapter which supplies those
hypotheses from its attachment sweep.

This strengthened single-site alternative has also been lifted through the
existing global-mesh and coherent-deletion provenance without reselecting a
colouring, matching, site, or path.  Its carrier deep-stack branch is
realized as a strict family of connected literal cut shores; its ambient
branch already contains either a proof-carrying separator exit or a strict
family of original same-turn path chords.
\Sverified{Mettapedia.GraphTheory.FourColor.Compositional.FullProfileGeometricAlternative.hasFullProfileStrictFamilyAlternative}
\Sverified{Mettapedia.GraphTheory.FourColor.Compositional.MeshFullProfileAlternative.exists_minimizer_with_fullProfileSweep_at_every_internal_row_junction}
\Sverified{Mettapedia.GraphTheory.FourColor.Compositional.CoherentFullProfileSweep.branchingOrBoundary_or_hasCoherentFullProfileSweepPair}
The same unchanged coherent pair can simultaneously retain this complete-profile
geometry and the fixed-colour connectivity alternative already obtained from the
deletion path: either its endpoint connectivity states agree, or a certified
local transition records their first difference.
\Sverified{Mettapedia.GraphTheory.FourColor.Compositional.FullProfileConnectivityAlternative.branchingOrBoundary_or_fullProfilePair_with_fixedColorAlternative}
The same pair now also carries the conservative three-colour refinement.  Either
all three labelled connectivity relations agree at its endpoints, or one
explicit colour and one adjacent path step exhibit a reachability transition;
the two graphs involved are literally equal away from at most ten ambient
edges.
\Sverified{Mettapedia.GraphTheory.FourColor.Compositional.FullProfileConnectivityAlternative.branchingOrBoundary_or_fullProfilePair_with_colorIndexedAlternative}
\Sverified{Mettapedia.GraphTheory.FourColor.Compositional.ColorIndexedConnectivityTransition.ColorLocalConnectivityTransitionReceipt.deleteLocalFootprint_eq}
\Sverified{Mettapedia.GraphTheory.FourColor.Compositional.ColorIndexedConnectivityTransition.ColorLocalConnectivityTransitionReceipt.card_edgeDisagreement_le}
This three-colour object is an internal refinement, not an identification with
the source's two complementary closed-web roles.  Such an identification still
requires the closed-web/no-hole sector hypotheses, so no sector or boundary
conclusion is inferred from this corollary alone.
There is nevertheless an exact local representation bridge.  On the carrier of
every proper alternating-component witness, the two residual matching roles
are pointwise distinct.  Their symmetric-difference carrier is therefore
literally the union of the two restricted one-factor graphs, just as a selected
two-colour graph is the union (equivalently, the symmetric difference) of its
two disjoint colour classes.
\Sverified{Mettapedia.GraphTheory.FourColor.Compositional.AlternatingCarrierUnion.ProperAlternatingComponentWitness.exchangeCarrier_eq_pairingUnion}
\Sverified{Mettapedia.GraphTheory.FourColor.Compositional.AlternatingCarrierUnion.colorClassGraph_symmDiff_eq_colorPairGraph}
Every edge of the distinguished residual cycle belongs to this literal union.
\Sverified{Mettapedia.GraphTheory.FourColor.Compositional.AlternatingCarrierUnion.ProperAlternatingComponentWitness.cycle_edges_pairingUnion}
This bridge is deliberately carrier-local: it justifies importing
two-matching arguments on an actual residual component, but it does not assert
that the full ambient sweep graph has the source closed-web/no-hole structure.
The local transition has now been put into a row-independent finite frame.  Its
code retains the explicit colour, adjacent path index, two boundary slots, both
row-arm choices, and the complete labelled connectivity relations in all three
colours immediately before and after the change.  Therefore, if the mesh has
more rows than such frames, there is a branching/boundary horn, a coherent pair
whose three connectivity states already agree, or two distinct physical rows
with the same framed transition.  In the last case the transition code and arm
directions agree, and reachability between every pair of the ninety labelled
boundary slots agrees in every colour on both sides of the transition.
\Sverified{Mettapedia.GraphTheory.FourColor.Compositional.FullProfileColorConnectivityFrame.branchingOrBoundary_or_colorSynchronized_or_repeatedFullProfileColorFrame}
\Sverified{Mettapedia.GraphTheory.FourColor.Compositional.FullProfileColorConnectivityFrame.semanticAgreement_of_fullProfileColorTransitionFrame_eq}
The threshold is left as the symbolic cardinality of the finite frame.  This
avoids elaborating its enormous adjacency-matrix cardinal as a numeral, and the
conclusion is still a repeated residual-matching observable rather than a
splice or a source closed-web profile.
The remaining work is geometric and mesh-level: consume the external
attachment horn, add the source's closed-web sector hypotheses before using
sector alternation on a same-turn chord family, and construct the
component-saturated carrier shore whose boundary is measured by the residual
return stacks.  A raw carrier path prefix is not such a shore: its off-carrier
edges leave the path one vertex at a time, even when their residual components
return later.
The first saturation layer is now explicit.  Given a carrier cut, retain the
earlier carrier vertices and fill every common-residual component whose paired
return endpoints are both earlier than the cut.  Any common-residual edge
leaving this saturated side must then start at an endpoint of a return open
across the cut.  The two physical shore stacks contain those endpoints, and
their combined depth bounds the number of outgoing common-residual darts.
\Sverified{Mettapedia.GraphTheory.PairingBoundarySaturation.card_pairingBoundaryExitDart_le_openEndpoint}
\Sverified{Mettapedia.GraphTheory.FourColor.Compositional.ResidualReturnComponentSaturation.exists_open_shore_arc_of_commonResidual_exit}
\Sverified{Mettapedia.GraphTheory.FourColor.Compositional.ResidualReturnComponentSaturation.card_commonResidual_exit_le_stackLengths}
This does not yet bound the full ambient boundary: reference-matching
attachments form a second layer, represented by the external-exit and
same-turn-chord alternatives above.  Keeping that layer separate prevents the
common-residual return pairing from being mistaken for an ambient cut.
The full ambient boundary has nevertheless been decomposed without an annular
assumption.  Every outgoing dart is either a common-residual return dart, a
reference-matching dart, or a local incoming-matching dart on the alternating
carrier.  Moreover saturation acquires no additional carrier vertices: on the
carrier it is exactly the displayed prefix.  Consequently the full boundary
cardinality is bounded by the two return-stack depths plus the two explicitly
named matching frontiers, and the local incoming-matching frontier injects
into the open endpoints of its transported carrier pairing.
\Sverified{Mettapedia.GraphTheory.FourColor.Compositional.CarrierSaturationBoundaryDecomposition.ambientExit_classification}
\Sverified{Mettapedia.GraphTheory.FourColor.Compositional.CarrierSaturationBoundaryDecomposition.cycleVertex_mem_carrierPrefixReturnSaturation_iff}
\Sverified{Mettapedia.GraphTheory.FourColor.Compositional.CarrierSaturationBoundaryDecomposition.card_localTauExitDart_le_openEndpoint}
\Sverified{Mettapedia.GraphTheory.FourColor.Compositional.CarrierSaturationBoundaryDecomposition.card_ambientExitDart_le_stackLengths_add_reference_add_localTau}
Thus the remaining unbounded object is no longer ``the ambient boundary'' in
general.  The reference-matching frontier itself splits into a carrier-local
part and a genuinely off-carrier part.  The carrier-local part injects into
the open endpoints of the transported reference pairing, just as the incoming
matching part injects into its own open endpoints.  Hence the only unbounded
term in the refined estimate is the off-carrier reference-matching frontier.
This is the exact object the external-exit/same-turn-chord horn must consume.
\Sverified{Mettapedia.GraphTheory.FourColor.Compositional.CarrierSaturationBoundaryDecomposition.card_localSigmaExitDart_le_openEndpoint}
\Sverified{Mettapedia.GraphTheory.FourColor.Compositional.CarrierSaturationBoundaryDecomposition.card_referenceMatchingExitDart_le_local_add_offCarrier}
\Sverified{Mettapedia.GraphTheory.FourColor.Compositional.CarrierSaturationBoundaryDecomposition.card_ambientExitDart_le_stackLengths_add_offCarrierReference_add_local}
The off-carrier term now has the intended physical meaning, rather than merely
the right name.  Saturation puts its inside endpoint in a completed
common-residual return component; the canonical return path exhausts that
component, so the endpoint occurs strictly inside the path.  Cubicity then
identifies the outgoing reference edge with the path's unique third edge, and
the fact that its outside endpoint is not saturated proves that this third edge
is an external attachment.  It therefore produces the exact separator
face-side receipt already consumed by the ambient-return alternative.
\Sverified{Mettapedia.GraphTheory.FourColor.Compositional.OffCarrierReferenceAttachment.exists_completedReturnExitWitness}
\Sverified{Mettapedia.GraphTheory.FourColor.Compositional.OffCarrierReferenceAttachment.exists_offCarrierAttachmentWitness}
\Sverified{Mettapedia.GraphTheory.FourColor.Compositional.OffCarrierReferenceAttachment.exists_attachmentExitSideReceipt_of_offCarrierReferenceExit}
Combining the two cases gives the route-level boundary dichotomy: either some
completed return supplies such a physical receipt, or the full ambient
boundary is bounded by the two return stacks and the open endpoints of the two
carrier-local pairings.  Thus no unnamed ambient-error term remains.
\Sverified{Mettapedia.GraphTheory.FourColor.Compositional.OffCarrierReferenceAttachment.attachmentReceipt_or_card_ambientExitDart_le_stacks_add_local}
\end{proof}

For the route, the two Kempe families suggest the matching stacks and capped
face progress supplies additional finite labels.  However those matchings
belong to a colouring state, whereas the physical graph surgery must be
defined before a global colouring exists.  In the shallow case an equal
sweep state is therefore only a candidate length interface.  A short
separator now produces a connected exact-state shore directly.  A deep
carrier stack produces a physical cyclic cut from two nested return chords
and saturates it to connected complementary shores, with an explicit bound
equal to three times the chosen separator support.  Thus realization of that
horn is closed, while uniform control of its separator-dependent width is
not.  A long carrier interval is compressed by the restricted LIFO sweep;
its remaining shallow outcome is a materially spaced repeated state inside
that same interval.  A long ambient return is likewise compressed by the
attachment sweep.  Complete-profile compression now removes its shallow
repeated-state outcome: beyond the explicit profile threshold it yields an
external separator exit or a deep same-turn chord nest of prescribed length.
What remains is to consume these two geometric horns, obtain a uniform bound
for the realized deep-return cuts, and choose compatible interfaces coherently
across neighbouring mesh sites.  The corollary is therefore a graph-level
local alternative at one alternating site; it proves neither a coherent
mesh-level family nor the final splice.

\subsection{The exact finite interface}

The colouring part of the finite-state claim does not depend on the current
Cell--3 carrier.  It is the following elementary fibre-product theorem.

\begin{theorem}[bounded-interface Count is a strict composition law]
\label{thm:bounded-interface-count-composition}
\Sprose
Let \(\mathcal C=\{r,b,p\}\).  A finite tangle \(T\) has disjoint ordered
input and output port sets \(I,O\), and
\(\operatorname{Col}_T(x,y)\) denotes its proper nonzero Tait colourings with
input word \(x\in\mathcal C^I\) and output word
\(y\in\mathcal C^O\).  Put
\[
        M_T(x,y)=|\operatorname{Col}_T(x,y)|.
\]
Suppose \(T_2\circ T_1\) is obtained by identifying the output ports \(J\)
of \(T_1\) with the equally ordered input ports of \(T_2\), deleting the
temporary stub vertices, and making no other identification.  Then
\[
 \operatorname{Col}_{T_2\circ T_1}(x,z)
 \;\cong\;
 \coprod_{y\in\mathcal C^J}
       \operatorname{Col}_{T_1}(x,y)\times
       \operatorname{Col}_{T_2}(y,z),                 \tag{C1}
\]
naturally in \(x,z\).  Consequently
\[
                   M_{T_2\circ T_1}=M_{T_1}M_{T_2}.   \tag{C2}
\]
After Booleanizing positive entries, the support relation
\[
 R_T=\{(x,y):M_T(x,y)>0\}
\]
satisfies \(R_{T_2\circ T_1}=R_{T_1};R_{T_2}\).  For port width at most
\(b\), there are at most \(2^{3^{2b}}\) such relations.
\end{theorem}

\begin{proof}
Restrict a colouring of the composite to the two tangles.  The restrictions
have one common colour on every identified port, hence determine a unique
middle word \(y\) and an element of the corresponding summand on the
right-hand side of (C1).

Conversely, take two proper colourings in one summand.  Their common middle
word says exactly that the two colours on every pair of identified stubs
agree.  Delete the temporary stubs and give the resulting seam edge that
common colour.  Every old internal vertex retains its three distinct
nonzero incident colours, and no new vertex is created, so the glued
colouring is proper.  Restriction and gluing are inverse operations.  Taking
cardinalities of the disjoint union gives (C2), and replacing positive
natural numbers by truth values gives relational composition.  There are
\(3^b\) possible words on each side, hence \(3^{2b}\) possible ordered word
pairs and at most \(2^{3^{2b}}\) relations.
\end{proof}

\begin{corollary}[exact cumulative support update]
\label{cor:bounded-interface-support-update}
\Sprose
For a serial word \(T_1,\ldots,T_n\), let \(S_i\subseteq\mathcal C^J\)
be the boundary words extendable across the first \(i\) pieces.  Then
\[
 S_{i+1}=\{y:\exists x\in S_i,\ (x,y)\in R_{T_{i+1}}\}.              \tag{C3}
\]
If \(S_i=S_j\) for \(i<j\), deleting the block
\(T_{i+1},\ldots,T_j\) leaves the final accepting support unchanged.
Port reordering merely conjugates every relation by the induced permutation,
so canonical relabelling is representation-invariant.
\end{corollary}

\begin{proof}
Equation (C3) is the Boolean form of (C2).  Equal supports followed by the
same suffix remain equal by induction over that suffix.  For a port
permutation \(\sigma\), restriction reads the relabelled word as
\(x\circ\sigma^{-1}\); applying this bijection at both ends conjugates the
relation and preserves composition and terminal nonemptiness.
\end{proof}

\begin{remark}[machine-checked algebra after the physical gluing bijection]
\label{rem:bounded-interface-count-lean}
The graph-level restriction/gluing bijection (C1) is now machine-checked on
the in-tree port-tangle model.
\Sverified{Mettapedia.GraphTheory.FourColor.GoertzelV24PortTangleGluing.gluingEquiv}
Its cardinal convolution is machine-checked as well.
\Sverified{Mettapedia.GraphTheory.FourColor.GoertzelV24BoundedInterfaceCountLaw.countMatrix_eq_matMul}
Boolean support of the convolution is relational composition.
\Sverified{Mettapedia.GraphTheory.FourColor.GoertzelV24BoundedInterfaceCountLaw.support_matMul_iff}
Equal cumulative supports remain equal after deleting the intervening block
and appending a common suffix.
\Sverified{Mettapedia.GraphTheory.FourColor.GoertzelV24BoundedInterfaceCountLaw.runSupport_delete_block}
Port relabelling commutes with matrix composition.
\Sverified{Mettapedia.GraphTheory.FourColor.GoertzelV24BoundedInterfaceCountLaw.relabel_matMul}
Finally, the Boolean support carrier has cardinality
\(2^{|X||Z|}\).
\Sverified{Mettapedia.GraphTheory.FourColor.GoertzelV24BoundedInterfaceCountLaw.card_boolSupports}
\end{remark}

Thus the direct zero-Count automaton never needs a global colouring of the
unprocessed graph.  Its state is the \emph{set of boundary words currently
extendable through the processed prefix}; its letters are literal tangle
relations.  Connectivity, bridgelessness, cyclic side data, and capped face
progress are additional finite \emph{geometric} interface fields needed by
the physical replacement theorem.  They are not licences to ask for a
whole-annulus Tait colouring.  At fixed width each field has finitely many
values (partitions of a finite port set, bounded incidence bits, and capped
lengths), and its update must be proved from the literal composition just as
(C1) was.

The manuscript's profile has three kinds of observable data:
\begin{enumerate}
\item the colours of edges crossing the cut;
\item connectivity for the $\alpha\beta$ and $\alpha\gamma$ Kempe families;
\item capped progress of faces crossing the cut.
\end{enumerate}

\begin{proposition}[global opened-annulus face uniqueness is false]
\label{prop:opened-global-face-uniqueness-refuted}
\Srefuted{dodecahedral two-cap opening; ten shared interior edges}
Pairwise uniqueness of shared interior face edges does not hold on the
literal annular carrier obtained by opening two faces of a normal-form map.
In particular it cannot be supplied globally to the current rooted
factorization interface.
\end{proposition}

\begin{proof}
Open two opposite pentagonal faces of the dodecahedral map, deleting the ten
cap-cycle edges but retaining their vertices as degree-one stubs.  The opened
map is connected, with $V=20$, $E=20$, and hence $F=2$.  Its ten stub
edges are bridges, while each of the other ten edges is incident to both face
orbits.  Thus the two distinct opened faces share ten interior edges, not at
most one.

For exact replay, the two deleted facial cycles are

\[
 (0,1,5,6,2),\qquad (14,18,19,16,15),
\]
and the ten shared edges are
\[
\begin{split}
 &(3,8),(3,9),(4,9),(4,10),(7,8),\\
 &(7,13),(10,11),(11,17),(12,13),(12,17).
\end{split}
\]
The opened rotation has SHA-256 fingerprint
\texttt{f109f292c33df0ae4eff0403a23efb424979ddba8ecc8ef2d14d16c54360c862}.
The independent rotation traversal also checked every connected normal-form
opening generated by \texttt{plantri -a -d -m5 -c5} through $34$ vertices:
all $4120$ violate the global premise.  The reproducible falsification
audit and its compact certificate are
\texttt{tools/fourcolor/v24\_open\_annulus\_face\_uniqueness\_audit.py} and
\texttt{results/fourcolor/v24\_open\_annulus\_face\_uniqueness\_20\_34.json}.
The single dodecahedral witness proves the proposition; the larger sweep only
shows that the problem is structural rather than exceptional.
\end{proof}

\begin{theorem}[conditional Cell--3 rooted factorization of one positive Cell]
\label{thm:exact-rooted-cell-factorization}
\Sconditional{ClosedWebAtGoodWord carrier; global premise (U), refuted on its opened annulus}
Assume the positioned Cell geometry satisfies
\[
  \tag{U}
  F\ne F'\quad\Longrightarrow\quad
  |\operatorname{sharedInteriorEdges}(F,F')|\le1.
\]
Fix a Cell--3 closed-web-at-good-word instance with global proper Tait
colouring $p$, and an executable corridor position together with the next two
positions needed by the rolling root.  Let $c$ be a positive literal Cell
colouring which agrees with the restriction of $p$ on their actual overlap.
Write $s(p)$ for the canonical rooted state at
the incoming cut, $a(p,c)$ for the finite Cell--rebase letter extracted from
the two colourings, and $p\diamond c$ for their right-biased splice.  Then:
\begin{enumerate}
\item $a(p,c)$ accepts $s(p)$ and its successor is exactly the canonical
  rooted state $s(p\diamond c)$ at the following cut;
\item every positive entry of the literal five-field Cell count supplies such
  a Cell colouring $c$, so every source-positive Cell behaviour occurs in the
  finite transition;
\item the induced transition is independent of the ambient representative:
  two witnesses with the same rooted source, literal Cell code, successor
  source-profile, and four rebase-role colours induce the same source and
  target rooted states, hence the same transition relation; and
\item the realized initial set and the realized directed-edge set have exact
  witness semantics.  Membership is equivalent, respectively, to the
  existence of a compatible first-position witness and to the existence of a
  compatible positioned prefix/Cell witness.
\end{enumerate}
After projection to the manuscript's two Kempe families, terminal-aware
compression preserves precisely the equality fibres of both endpoints.
Consequently it changes neither the realized edge count nor the realizable
initial count.
\end{theorem}

\Sverified{Mettapedia.GraphTheory.FourColor.GoertzelV24ClosedWebLocalLayerSerialRootedInteractionRollingTransitionColorParametric.sourceLocalLayerSerialRootedInteractionRollingCellFactorForColorAt_successor_exact}
\Sverified{Mettapedia.GraphTheory.FourColor.GoertzelV24ClosedWebLocalLayerSerialCellCountRootedRollingFactorization.exists_rootedRollingFactorizationForColor_of_count_pos}
\Sverified{Mettapedia.GraphTheory.FourColor.GoertzelV24ClosedWebLocalLayerSerialRootedInteractionRealizableTransition.sourceLocalLayerSerialRootedInteraction_targetCongruentAt}
\Sverified{Mettapedia.GraphTheory.FourColor.GoertzelV24ClosedWebLocalLayerSerialRootedInteractionRealizableTransition.sourceLocalLayerSerialRootedInteraction_transition_invariant_of_targetCongruent}
\Sverified{Mettapedia.GraphTheory.FourColor.GoertzelV24ClosedWebLocalLayerSerialRootedInteractionPositionedSourceClosure.mem_sourceLocalLayerSerialRealizedInitialSourceCompressedProfileSet_iff}
\Sverified{Mettapedia.GraphTheory.FourColor.GoertzelV24ClosedWebLocalLayerSerialRootedInteractionPositionedSourceClosure.mem_sourceLocalLayerSerialPositionedRealizedSourceCompressedProfileEdgeSet_iff}

\begin{proof}
For the first assertion, form the literal Cell colour on the ambient edge
carrier, splice it with $p$ on the Cell region, and then perform the boundary
rebase.  Positivity makes every one of the four rebase-role colours nonzero.
The finite letter first checks the physical Cell profile and the equality of
the shared intermediate profile.  Its tracked and facial rebase maps then
compute the two successor exterior relations, while the rolling maps compute
the next interaction roots.  Projection along the fixed carrier coordinates
recovers the successor cut colours, the two connectivity relations, and the
capped face-progress field.  Equality in each field therefore identifies the
computed state with $s(p\diamond c)$.

For the second assertion, a positive count is a nonempty finite set of literal
open Tait colourings with the prescribed left and right five-field profiles.
Choose one member.  The overlap hypothesis is exactly the condition required
to splice it with $p$, so the first assertion applies.  Conversely, letters in
the realized alphabet are defined as images of these compatible literal
witnesses; hence no accepted behaviour is introduced by the extraction.

For representation invariance, retain only the rooted source, literal Cell
code, output profile, and rebase-role colours.  Equality of these data gives
equality on the active part of the fixed rebase collar.  The tracked local
graphs therefore agree.  The facial output is already determined by the
common output profile, and the four role colours determine the transported
cut colouring.  Extensionality of the rooted state gives equal targets, so
the two singleton transition relations coincide.  Finally, both realized
sets are finite images.  Membership in an image is equivalent to existence
of a preimage, giving the two witness characterizations.  On terminal graph
profiles the compression map is injective, endpoint by endpoint, which proves
the final cardinality statement.
\end{proof}

\begin{proposition}[the Cell--3 carrier is not supplied by a direct Count obstruction]
\label{prop:closed-web-carrier-not-direct-count}
\Sprose
An arbitrary zero-accepting-Count instance does not supply an object of type
\texttt{ClosedWebAtGoodWord.Instance} on its opened annulus.  If the annulus
itself has zero Tait count then such an object is impossible; even when the
annulus is colourable, the direct hypothesis supplies neither a totally
closed colouring nor a good inner word.
\end{proposition}

\begin{proof}
The type contains a proper edge colouring of the whole opened graph;
properness is bundled into \texttt{EdgeColoring}, while its \texttt{tait}
field supplies nonzeroness.  Hence an inhabitant already proves that the open
graph has positive Tait count.  In the colourable case the additional
\texttt{totallyClosed} and \texttt{goodWord} fields are still extra
existential witnesses, absent from the direct obstruction premise.  Neither
follows from zero terminal acceptance of the corresponding closed map.
\end{proof}

The six cited Lean declarations are machine-checked with (U) as an explicit
hypothesis; they do not prove (U).  Proposition~\ref{prop:opened-global-face-uniqueness-refuted}
shows that no cap-deletion adapter can establish it: opening the caps splices
ambient faces and destroys global uniqueness on the intended carrier.
Therefore this conditional theorem is not presently an on-route
factorization result, and no closure number derived from it is authoritative.

There are therefore two independent repairs.  First, factorization must be
generalized from the special globally coloured
\texttt{ClosedWebAtGoodWord.Instance} to a colouring-independent annular Cell
context together with the partial positive prefix and literal Cell witnesses
used by cumulative Count.  A total colouring of the opened annulus may be
used to prove that one particular run is sound, but it cannot define the
autonomous letter relation: doing so would make the available transitions
depend on extendability through the unprocessed suffix.  In particular no
theorem on the direct path may obtain prefix properness from the type of the
ambient \texttt{EdgeColoring} or prefix nonzeroness from \texttt{web.tait}.
Second, face lookup must be local, not a stronger global hypothesis.  The
closed minimal map does satisfy ambient pairwise uniqueness.
\Sverified{Mettapedia.GraphTheory.FourColor.GoertzelV24MinimalFaceIntersections.pairwiseUniqueSharedInteriorEdges}
If two corridor faces are fully retained by the opening, their open images
still share at most one edge; the tree proves this both for arbitrary retained
ambient faces and for every pair of retained corridor faces.
\Sverified{Mettapedia.GraphTheory.FourColor.GoertzelV24OpenRegionFaceIntersectionTransport.openMinimalFaceOrbits_sharedInteriorEdges_card_le_one}
\Sverified{Mettapedia.GraphTheory.FourColor.GoertzelV24OpenRegionFaceIntersectionTransport.openMinimalCorridorFaces_sharedInteriorEdges_card_le_one}
The selected common-neighbour incidence calculation likewise identifies the
two local flanks without reinstating global uniqueness.
\Sverified{Mettapedia.GraphTheory.FourColor.GoertzelV24ClosedWebLocalCommonNeighborIncidence.face_eq_before_or_after_of_commonNeighborVertexIncidenceAt}
The remaining Lean task is consequently to inventory both dependencies:
replace every use of the global colouring by its rooted prefix/Cell witness,
replace every globally chosen shared edge by the appropriate retained-face or
pointwise-incidence theorem, and keep boundary-spliced faces out of the simple
interior-dual lookup.  This is a carrier-generalizing,
hypothesis-minimizing refactor of the factorization, not a new geometric
assumption.

The tracked half of this refactor is now complete on the corrected
pointwise-selected source geometry.  At each corridor position the selected
literal Cell has two vertices and a unique internal bond.  A crossing of a
strictly earlier Cell can neither equal nor meet an outgoing crossing of the
successor Cell: the corresponding corridor faces are separated along the
clean dual geodesic.  If an edge were internal to both Cells, equality of the
two endpoint sets would reduce to the same contradiction.  Consequently any
historical Cell edge which reappears in the successor rolling carrier is one
of the four rebase-role edges, and every active successor edge is either a
rebase-role edge or belongs to the current literal Cell.
\Sverified{Mettapedia.GraphTheory.FourColor.GoertzelV24ClosedWebPointwiseSelectedCellPastOverlap.pointwiseSelectedSourceLocalLayerSerialCellRebase_pastCellOverlap_mem_switch}
\Sverified{Mettapedia.GraphTheory.FourColor.GoertzelV24ClosedWebPointwiseSelectedCellPastOverlap.pointwiseSelectedSourceLocalLayerSerialCellRebase_activeTarget_mem_switch_or_currentCell}

The resulting finite tracked factor uses the actual interaction carrier
(at most forty-nine edges), a four-role deletion mask, the literal local
adjacency table, and a partial coordinate map into the next rolling carrier
(at most twenty-one edges).  Its update is rowwise equal to the literal
successor graph on the interaction carrier.  An unmapped target coordinate is
isolated by the historical-overlap theorem above, so generic partial
contraction is exact for vertex equality, direct adjacency, and exterior
connectivity.  Thus the factor's executable target state is exactly the
canonical tracked exterior code on the next literal rolling carrier.
\Sverified{Mettapedia.GraphTheory.FourColor.GoertzelV24ClosedWebPointwiseSelectedTrackedRebaseFactor.pointwiseSelectedSourceLocalLayerSerialTrackedRebaseFactorAt_uniformSuccessorCode_exact}
\Sverified{Mettapedia.GraphTheory.FourColor.GoertzelV24ClosedWebPointwiseSelectedTrackedTargetCode.pointwiseSelectedSourceLocalLayerSerialTrackedRebaseFactorAt_targetState_code_exact}

The combined finite projection now closes both rolling roots on this
pointwise-selected carrier.  For every compatible positive prefix witness and
literal Cell colouring, the tracked target root computed by the finite rebase
factor is exactly the canonical tracked prefix state at the next cut.

\Sverified{Mettapedia.GraphTheory.FourColor.GoertzelV24ClosedWebPointwiseSelectedRollingProjectionExact.pointwiseSelectedSourceLocalLayerSerialRollingProjectionFactorAt_tracked_exact}

The facial target root is exact as well.  Its fixed twenty-four interface
slots and four persistent-port slots recover the next-cut facial connectivity,
presence bits, and cap-at-five component sizes from the cap-at-six target
code.  The proof uses only the pointwise-selected face carrier and the local
prefix/Cell witnesses; it does not reinstate global opened-annulus face
uniqueness.

\Sverified{Mettapedia.GraphTheory.FourColor.GoertzelV24ClosedWebPointwiseSelectedRollingProjectionExact.pointwiseSelectedSourceLocalLayerSerialRollingProjectionFactorAt_face_exact}

\begin{proposition}[the selected finite Cell--rebase machine is formation-local]
\label{prop:selected-cell-rebase-formation-local}
\Sverified{Mettapedia.GraphTheory.FourColor.GoertzelV24ClosedWebPointwiseSelectedCellRebaseFiniteTransition.pointwiseSelectedSourceLocalLayerSerialCellRebaseFiniteLetterAt_transition}
Fix only the colouring-independent annular formation, a pointwise-selected
literal Cell, and the following boundary rebase.  Every positive prefix
colour function and compatible positive literal Cell colouring extracts a
supported finite letter whose source and target are its exact terminal-aware
input and rebased output profiles.  In particular the finite carrier, colour
tables, rooted interaction receipt, supported-letter predicate, and complete
Cell--rebase transition consume no global Tait colouring, total-closure
witness, or good inner word.  Specializing a former
\texttt{ClosedWebAtGoodWord.Instance} factors through its
colouring-independent \texttt{toFormation} field and gives the same
transition.
\Sverified{Mettapedia.GraphTheory.FourColor.GoertzelV24ClosedWebPointwiseSelectedCellRebaseFiniteTransition.pointwiseSelectedSourceLocalLayerSerialCellRebaseFiniteLetterAt_transition_of_instance}
\end{proposition}

\begin{proof}
The finite Cell letter is computed from the selected formation, the positive
prefix values, and the literal open Tait colouring of the Cell.  Compatibility
on the terminal overlap proves physical Cell support.  The right-biased splice
is positive on each of the four rebase roles, so the normalized rebase letter
is supported; its input is definitionally the Cell letter's output.  Their
pair therefore witnesses the finite transition between the displayed source
and target.  The compatibility theorem applies this construction to
\texttt{web.toFormation}; none of the three additional fields of
\texttt{web} occurs in the resulting statement or proof.
\end{proof}

\begin{proposition}[selected Cell--rebase realizations are one-step local]
\label{prop:selected-cell-rebase-witness-local}
\Sverified{Mettapedia.GraphTheory.FourColor.GoertzelV24ClosedWebPointwiseSelectedCellRebaseLocalWitness.pointwiseSelectedLocalRealizedTransitionAt_iff_ambient}
Fix a selected Cell and its following boundary rebase.  Let the observed
prefix support be the union of the current terminal-aware cumulative region
and the rebase switch.  The latter is the image of the four fixed rebase
roles, and therefore has cardinality at most four.
\Sverified{Mettapedia.GraphTheory.FourColor.GoertzelV24ClosedWebPointwiseSelectedCellRebaseLocalWitness.card_observedPrefixRegion_sdiff_terminal_le_four}
Replace the former positive prefix function on the whole ambient edge set by
a positive function on this observed support, retaining the same compatible
literal Cell colouring.  Then the realized source--target transition image is
unchanged.  Moreover every transition in this local realized image is a
transition of the canonical finite Cell--rebase machine.
\Sverified{Mettapedia.GraphTheory.FourColor.GoertzelV24ClosedWebPointwiseSelectedCellRebaseLocalWitness.pointwiseSelectedLocalRealizedTransitionAt_finiteTransition}
\end{proposition}

\begin{proof}
Restriction is immediate.  In the other direction extend a local function by
one fixed nonzero colour outside the observed support; this extension is only
an adapter to the former API.  The source profile reads only the current
terminal region.  For the target, adjoin the common literal Cell.  On the
next terminal region, an edge in the rebase switch is already observed.  An
edge outside the switch belongs to the old pre-rebase region if and only if it
belongs to the next terminal region.  Since the pre-rebase region is exactly
the union of the current terminal region and the Cell, its colour is therefore
either an observed prefix colour or the common Cell colour.  Hence the two
spliced colourings agree throughout the next terminal region.  Regional
profile congruence gives equality of both exact endpoints, so restriction and
extension prove the two inclusions between realized transition images.  The
finite-letter soundness theorem applied to the extended local witness gives
the final assertion.  No converse from an arbitrary abstract supported letter
to a physical local witness is asserted here.
\end{proof}

Thus the finite one-Cell transition has now been proved exact on both of the
strong rolling roots from which the manuscript's cumulative profile is
projected, its finite letter semantics has been separated from the global
closed-web colouring, and its physical realization witness carries no colour
coordinate on the unprocessed suffix.  This does not yet prove the full
theorem above.  The realizable initial image still sits under the surrounding
closed-web parameter.  The source collar or three-Cell word must therefore be
made into a formation directly, and the resulting local realized transitions
and initials must be connected soundly and completely to the cumulative Count
and Seed semantics.  In particular, completeness of the abstract supported
finite-letter relation with respect to physical local witnesses remains a
separate obligation.  Global opened-annulus face uniqueness is no longer an
obligation of either rolling transition.

The Lean transition uses a conservative three-pair refinement internally and
then projects to exactly the two families named by the manuscript.  It proves
that every compatible positive literal Cell witness followed by the boundary
rebase produces one exact source and target profile edge, and that terminal-
aware compression is lossless relative to that two-pair source profile:
compressed edges agree exactly when their projected two-pair edges agree.
\Sverified{Mettapedia.GraphTheory.FourColor.GoertzelV24ClosedWebLocalLayerSerialRootedInteractionPositionedSourceClosure.sourceLocalLayerSerialPositionedSourceCompressedProfileEdge_eq_iff}
The realized edge image and realizable initial image both preserve cardinality
under this compression.
\Sverified{Mettapedia.GraphTheory.FourColor.GoertzelV24ClosedWebLocalLayerSerialRootedInteractionPositionedSourceClosure.card_sourceLocalLayerSerialPositionedRealizedSourceCompressedProfileEdgeSet_eq}
\Sverified{Mettapedia.GraphTheory.FourColor.GoertzelV24ClosedWebLocalLayerSerialRootedInteractionPositionedSourceClosure.card_sourceLocalLayerSerialRealizedInitialSourceCompressedProfileSet_eq}

The conservative third family has not been smuggled into the source count.  At
the live $(2,1,4)$ shape its carrier is exactly sixteen times larger:
\[
  12{,}556{,}911{,}218{,}688
  =16\cdot784{,}806{,}951{,}168.
\]
This equality is machine-checked.
\Sverified{Mettapedia.GraphTheory.FourColor.GoertzelV24SourceTerminalAwareProfileCompression.boundedTerminalAwareCompressedProfileCount_two_one_four_eq}

The local letter and the cumulative terminal receipt have different shapes,
and this distinction is load-bearing.  The rolling Cell letter has shape
$(2,1,4)$ because it remembers one local seam terminal.  Terminal acceptance
is read only after passing to the already-constructed cumulative source
crosscut carrier of shape $(2,5,8)$, whose five terminals are the five fixed
cap-foot edges.  The following proposition verifies that no additional
terminal datum is needed there.

\begin{proposition}[cap acceptance is a projection of the five-terminal source receipt]
\label{prop:five-terminal-source-cap-view}
\Sverified{Mettapedia.GraphTheory.FourColor.GoertzelV24SourceTerminalCapView.standardGaugeCapComposedMenuBState_sourceCompressedCodeOfCell3_iff}
In the standard colour-name gauge, project a graph-semantic state of the
existing \texttt{BoundedCorridorCutProfile 2 5 8} to the two Kempe families named in
the manuscript and apply the lossless terminal-aware compression.  The
compressed state determines:
\begin{enumerate}
\item the colours of all five cap feet; and
\item the complete terminal connectivity relations for the
      $\alpha\beta$ and $\alpha\gamma$ families.
\end{enumerate}
Consequently it determines exactly whether either source pair satisfies the
cap-composed Menu-B predicate.  No graph lookup and no enlargement of the
one-Cell letter is required.
\end{proposition}

\begin{proof}
For one terminal, the two diagonal bits record membership in
$\{\mathrm{red},\mathrm{blue}\}$ and
$\{\mathrm{red},\mathrm{purple}\}$.  On a nonzero Tait colour the possible
bit pairs are therefore
\[
 (1,1)\mapsto\mathrm{red},\qquad
 (1,0)\mapsto\mathrm{blue},\qquad
 (0,1)\mapsto\mathrm{purple};
\]
the remaining pair $(0,0)$ is impossible.  This finite decoding is checked
as
\Sverified{Mettapedia.GraphTheory.FourColor.GoertzelV24SourceTerminalCapView.terminalColorOfSourceBits_eq},
and on a literal graph-derived regional profile it gives exactly the colours
of the five named terminal edges:
\Sverified{Mettapedia.GraphTheory.FourColor.GoertzelV24SourceTerminalCapView.terminalCapView_word_toSourceProfile_regionalProfile}.

Terminal diagonal entries are stored explicitly by the compressor.  Every
off-diagonal entry is stored once on the corresponding unordered pair, and
symmetry reconstructs its opposite orientation.  Hence compression preserves
the whole cap-facing view:
\Sverified{Mettapedia.GraphTheory.FourColor.GoertzelV24SourceTerminalCapView.terminalCapView_compressBounded}.
The cap-composed predicate is a function only of this five-word and these two
relations, so rewriting by that equality proves the proposition.
\end{proof}

The preceding statement used literal colour names only to choose coordinates.
The manuscript instead names the two connectivity roles relative to the
ordered majority/singleton triple of the current good word and then says
``by global colour relabeling'' to use the standard names.  That normalization
is also exact.

\begin{proposition}[global colour normalization preserves cap acceptance]
\label{prop:source-cap-colour-gauge}
\Sverified{Mettapedia.GraphTheory.FourColor.GoertzelV24SourceTerminalCapGauge.standardGaugeCapComposedMenuBState_normalized_iff}
Let $(\alpha,\beta,\gamma)$ be the ordered majority, first-singleton and
second-singleton colours of a five-terminal cap view.  There is a canonical
zero-fixing permutation of the four Klein-group colours sending
$(\alpha,\beta,\gamma)$ to
$(\mathrm{red},\mathrm{blue},\mathrm{purple})$.  Relabel the boundary word by
this permutation while retaining the two abstract strand roles
$\alpha\beta$ and $\alpha\gamma$.  The arbitrary-gauge cap-composed Menu-B
predicate holds before relabeling if and only if the standard-gauge predicate
holds afterward.

Thus the two-pair source receipt is sufficient in every colour gauge.  This
uses the fact that its pair labels are roles relative to the ordered triple;
it does \emph{not} assert that two fixed literal colour pairs suffice without
normalization.
\end{proposition}

\begin{proof}
The inverse of the canonical Tait-triple equivalence sends
$\alpha,\beta,\gamma$ to the three standard nonzero colours and fixes zero.
Relabeling preserves the three boundary multiplicities, so the normalized
word has the standard majority/singleton triple:
\Sverified{Mettapedia.GraphTheory.FourColor.GoertzelV24SourceTerminalCapGauge.normalizedCapView_standardTriple}.
For either abstract source role, relabeling carries its unique inactive
position to the same cap index and leaves its Boolean strand relation
unchanged.  Therefore its cap-composed predicate is equivalent before and
after normalization:
\Sverified{Mettapedia.GraphTheory.FourColor.GoertzelV24SourceTerminalCapGauge.capComposedMenuBForPair_normalized_iff}.
Taking the disjunction over the two roles proves the claim.  The independent
cap-extension table is equivariant under the same zero-fixing permutation:
\Sverified{Mettapedia.GraphTheory.FourColor.GoertzelV24AnnularFrontierMenuCapTable.normalizedCapTable_transport}.
\end{proof}

The corrected all-colour Seed count and the compressed cumulative receipt now
meet at exactly this normalized terminal view.  It is useful to state the two
directions separately, because the remaining realization problem lies between
them rather than inside either statement.

\begin{proposition}[corrected Seed acceptance and the compressed accepting set]
\label{prop:source-seed-compressed-acceptance}
\Sverified{Mettapedia.GraphTheory.FourColor.GoertzelV24SourceTerminalCapCount.capComposedSeedCount_pos_iff_exists_normalized_sourceView}
For a fixed cap word, the corrected cap-composed Seed count is positive if and
only if there is a realized all-colour uniform profile of positive count and an
ordered majority/singleton triple whose normalized two-role terminal view is
accepted.

For every graph-semantic bounded cumulative source profile, define its
compressed code to be accepting precisely when the standard-gauge terminal
view decoded from the code satisfies cap-composed Menu~B.  Lossless compression
preserves this predicate exactly:
\Sverified{Mettapedia.GraphTheory.FourColor.GoertzelV24SourceTerminalCapCount.compressBounded_mem_acceptedSourceCompressedProfileSet_iff}.
In particular this holds on the existing Cell--3 cumulative carrier of shape
$(2,5,8)$:
\Sverified{Mettapedia.GraphTheory.FourColor.GoertzelV24SourceTerminalCapCount.sourceCompressedCodeOfCell3_mem_accepted_iff}.
\end{proposition}

\begin{proof}
Project an all-colour profile to the two source roles
$\alpha\beta$ and $\alpha\gamma$.  For either role, the resulting cap relation
is definitionally the corresponding component of the all-colour connection
table.  The corrected Seed predicate therefore says exactly that one of these
two projected roles is cap-accepted.  Proposition~\ref{prop:source-cap-colour-gauge}
then transports this disjunction to the standard colour gauge.  Rewriting the
already verified positive-count factorization by this equivalence proves the
first assertion.

For the second assertion, the accepting set is the finite subset obtained by
filtering the compressed carrier by the decoded standard-gauge cap predicate.
Proposition~\ref{prop:five-terminal-source-cap-view} says that decoding after
compression gives the same terminal word and the same two strand relations as
decoding before compression, so membership rewrites to the uncompressed
predicate.  Specializing the bounds to $(2,5,8)$ gives the last statement.
\end{proof}

The two profile presentations in the preceding proposition are not merely
abstract finite carriers.  On a literal annulus their terminal coordinates
coincide.

\begin{proposition}[the counted terminal profile is cumulative]
\label{prop:annular-terminal-profile-realization}
\Sverified{Mettapedia.GraphTheory.FourColor.GoertzelV24AnnularTerminalProfileRealization.capComposedSeedCount_pos_iff_exists_realizedSourceProfile}
Let an annular boundary datum be backed by a rotation system on its underlying
simple graph.  Give the whole annulus the regional edge set $E(G)$, no crossing
ports, its five cap-foot edges as terminals, and no face-fragment coordinates.
Then every Tait colouring counted by the uniform Menu--B profile determines a
graph-derived cumulative source profile whose five-terminal cap view is
exactly the counted view.  After the manuscript's global colour normalization,
the equality remains literal.  Consequently the corrected Seed count is
positive if and only if an accepted normalized cap view is realized by such a
cumulative source profile.
\end{proposition}

\begin{proof}
Fix two selected colours.  The counted profile puts two cap stubs in the same
component precisely when they are connected in the subgraph consisting of
edges of those colours.  Each cap stub is incident with its named cap-foot
edge.  In any graph, two edges with chosen incident endpoints are connected
in the line graph if and only if those endpoints are connected in the
original graph:
\Sverified{Mettapedia.GraphTheory.FourColor.GoertzelV24AnnularTerminalProfileRealization.boundaryLineGraph_reachable_iff_vertexReachable}.
Applying this to the selected-colour subgraph identifies the source component
label with tracked-edge reachability:
\Sverified{Mettapedia.GraphTheory.FourColor.GoertzelV24AnnularTerminalProfileRealization.activePairComponents_eq_iff_bicoloredReachable}
and
\Sverified{Mettapedia.GraphTheory.FourColor.GoertzelV24AnnularTerminalProfileRealization.sourcePair_connectionTable_eq_regionalTrackedConnectivity}.

Because the regional edge set is all of $E(G)$, the regional guards add no
restriction.  The diagonal entries reconstruct the five nonzero terminal
colours, while the off-diagonal entries are the two source connectivity
relations.  Thus the entire terminal views agree.  Relabel the actual edge
colouring by the canonical zero-fixing colour equivalence, rather than merely
renaming the view.  The selected two-colour subgraphs before and after
relabeling are isomorphic, so the complete connection table is preserved.
This proves the normalized equality
\Sverified{Mettapedia.GraphTheory.FourColor.GoertzelV24AnnularTerminalProfileRealization.normalizedUniformCapView_eq_wholeAnnulusSourceProfile}.
Finally, positivity of a uniform-profile count supplies an actual Tait
colouring, and conversely that colouring lies in its own profile fibre;
substitution gives the displayed equivalence.
\end{proof}

The initial object of the literal serial calculation can likewise be made
genuinely local.  This matters because the source does not claim that a wide
annulus is one bounded-width serial corridor: its finite calculation is made
on a selected depth-one collar (and, at the local audit, on three consecutive
cells).  A witness for the first serial position must therefore not carry a
second arbitrary positive-colour function on the unprocessed part of the
annulus.

\begin{proposition}[realizable initials are prefix-local]
\label{prop:serial-initial-prefix-local}
\Sverified{Mettapedia.GraphTheory.FourColor.GoertzelV24ClosedWebLocalLayerSerialLocalInitial.sourceLocalLayerSerialLocalInitialSourceProfileSet_eq}
Fix a literal source corridor and its first terminal-aware serial position.
Replace the positive colour function on all edges of the opened graph by a
positive colour function whose domain is exactly the active prefix region,
retaining the same compatible literal Cell colouring.  The resulting finite
image of initial cumulative source profiles is exactly the former realizable
initial image.
\end{proposition}

\begin{proof}
Restrict an old witness to the active prefix.  Conversely, extend a local
witness by one fixed nonzero colour outside that prefix; this is only an API
adapter and consults no suffix datum.  Restriction after extension is literal.
For extension after restriction, every field of the rooted source receipt
observes the colour function only on its stated terminal input region:
\Sverified{Mettapedia.GraphTheory.FourColor.GoertzelV24ClosedWebLocalLayerSerialRootedInteractionRealizableTransition.sourceLocalLayerSerialRootedInteractionStateForColorAt_eq_of_eq_on_terminal}.
Thus the complete rooted source state, and hence its bounded cumulative
projection, is unchanged.  The two inclusions between the finite images follow
by extension and restriction.
\end{proof}

Together with Proposition~\ref{prop:annular-terminal-profile-realization}, this
closes the narrower representation-level excess in the initial witness: its
additional positive-colour coordinate is now prefix-local.  It does \emph{not}
remove the global Tait colouring, total closure, or good word already contained
in the surrounding \texttt{ClosedWebAtGoodWord.Instance}.  Proposition~
\ref{prop:selected-cell-rebase-formation-local} completes the corresponding
refactor for the finite carrier, state, supported letter, and one-step
transition: those objects now require only the colouring-independent
formation.  Proposition~\ref{prop:selected-cell-rebase-witness-local} also
removes the transition witness's ambient suffix-colour coordinate.  What
remains is narrower than whole-annulus coverage but still substantive:
construct that formation directly from the radius-$1$ collar or three-Cell
word actually selected by the source, identify both local realized images with
the cumulative Count and Seed semantics consumed at the local target, and
prove physical completeness for the abstract finite letters retained by the
closure.  No further equivalence between vertex-component and edge-component
semantics is owed.  The filtered finite accepting set has also not yet been
emitted as an optimized executable certificate.

Theorem~\ref{thm:exact-rooted-cell-factorization} checks the Cell--3 algebra
under (U), but Propositions~\ref{prop:opened-global-face-uniqueness-refuted}
and~\ref{prop:closed-web-carrier-not-direct-count} prevent applying it to the
direct source target.  Even after both refactors it will not by itself justify
a numerical closure count: the exact realized letter image is still specified
as a finite image of literal witnesses, and the initial image, although local
in its additional positive-colour coordinate, is still parameterized by the
closed-web instance.  Neither is yet a practical enumerator.

\begin{proposition}[Executable optimized closure certificate]
\label{prop:exact-closure-gate}
\Sopen
For each relevant corridor type and target semantics, produce a finite
raw-profile array $Q$, a finite array $A$ of \emph{physical} Cell types, an
initial array $I\subseteq Q$, an accepting array $F\subseteq Q$, and a
letter-labelled edge array
$E\subseteq Q\times A\times Q$ such that
\[
 q\in I\quad\Longleftrightarrow\quad
 q\text{ is produced by a compatible literal first-position witness},
\]
and
\[
 (q,a,q')\in E\quad\Longleftrightarrow\quad
 (q,q')\text{ is produced by a compatible positive prefix/Cell witness
 of physical type }a.
\]
Also prove
\[
 q\in F\quad\Longleftrightarrow\quad
 q\text{ has a literal terminal realization satisfying the chosen target.}
\]
The letter index must forget the colouring witness while retaining every
geometric choice on which the literal Cell count depends; an unlabelled union
of edges is insufficient for cumulative Count.

For $a\in A$ let
$R_a(S)=\{q'\mid\exists q\in S,\ (q,a,q')\in E\}$.
Prove the displayed biconditionals in the kernel, replay raw reachability,
then replay the determinized closure from $I$ under the maps $R_a$, and
finally minimize the reachable deterministic system by equality of accepted
right languages.  Compare all three fingerprints and counts with an
independently implemented generator:
\[
 |Q_{\rm raw}^{\rm reach}|,\qquad
 |\operatorname{Reach}(\mathcal P(Q))|,\qquad
 |\operatorname{Reach}(\mathcal P(Q))/\mathord{\sim}_{\rm Nerode}|.
\]
Then, and only then, these array lengths are the exact raw, cumulative-support,
and observational state counts for that corridor type.  This certificate is
needed when its minimized count is used to sharpen an effective threshold or
to replay a finite base.  It is not needed for the existence of a finite
direct-Count state space once actual ordered noose cuts have been supplied:
Theorem~\ref{thm:raw-noose-count-pumping} below gives the conservative state
space directly.
\end{proposition}

There are two different accepting arrays.  For the Lift/Seed system, $F$ is
the cap-composed predicate of
Definition~\ref{def:cap-composed-pair-relation}.  Replaying that closure
against the open-tangle predicate of
Proposition~\ref{prop:open-menu-predicate-vacuous} would certify only that a
nonempty good fibre is nonempty; it would not test the Seed conclusion.  For
the direct Count system, $F$ instead consists of terminal profiles whose
literal boundary closure is a proper Tait colouring of the closed instance.
The direct descent consumes
\[
  S_n\cap F=\varnothing
  \quad\Longleftrightarrow\quad
  \text{the restored closed map has zero Tait count},
\]
not menu B.  The transition carrier may be shared, but these accepting
predicates and their Nerode quotients must not be conflated.

The measured $(5,0)$ certificates remain valuable regression tests and
base-bound optimizations, but they are not substitutes for this equality
theorem when their numerical sizes are quoted: their own headers correctly
say that identifying the decoded transition with every source-legal corridor
is a separate fidelity obligation.

\subsection{Direct cumulative-Count descent}

The finite-state object needed to pump a zero count is not one raw profile.
An uncolourable closed instance has no accepting colouring from which one
could select a profile path.  Its prefix invariant is instead the \emph{set}
of all raw profiles realized by positive prefix colourings.  This
determinization costs a powerset in the crude bound and is why the three
counts in Proposition~\ref{prop:exact-closure-gate} must not be conflated.

\begin{definition}[support action of a positive Count letter]
\label{def:count-support-action}
\Sprose
Let $Q$ be finite and let $M_a\in\operatorname{Mat}_Q(\mathbb N)$ be the
Count matrix of a physical Cell $a$.  Its Boolean support relation is
\[
       q\,R_a\,q' \quad\Longleftrightarrow\quad M_a(q,q')>0.
\]
For $S\subseteq Q$ put
\[
       R_a[S]=\{q'\in Q:\exists q\in S,\ q\,R_a\,q'\}.
\]
If $x\in\mathbb N^Q$ is a cumulative Count row vector and
$S=\{q:x_q>0\}$, then the support of $xM_a$ is exactly $R_a[S]$.
\end{definition}

\begin{proof}
The $q'$ entry of $xM_a$ is
$\sum_q x_qM_a(q,q')$.  Every summand is a natural number, so the sum is
positive exactly when at least one summand is positive.  A product of natural
numbers is positive exactly when both factors are positive.  These two
equivalences give precisely $q'\in R_a[S]$.
\end{proof}

\begin{theorem}[determinized support-state deletion]
\label{thm:determinized-support-deletion}
\Sprose
Let $a_1,\ldots,a_n$ be physical Cell letters, let $I\subseteq Q$ be the
initial support, and define
\[
  S_0=I,\qquad S_{k+1}=R_{a_{k+1}}[S_k].
\]
If $S_i=S_j$ for some $0\le i<j\le n$, delete the block
$a_{i+1},\ldots,a_j$.  The shortened word has exactly the same final support
$S_n$.  In particular, for any accepting set $F\subseteq Q$, it has an
accepting colouring if and only if the original word does:
\[
                         S_n\cap F\ne\varnothing .
\]
Consequently, if more prefix positions occur than there are reachable
determinized support states, some nonempty block is deletable without
changing zero versus positive total Count.
\end{theorem}

\begin{proof}
After the deletion, the support at the new position $i$ is $S_i=S_j$.
Apply the same remaining functions
$R_{a_{j+1}},\ldots,R_{a_n}$ to equal sets.  Induction over this common
suffix gives equality of every later support, hence equality of the final
support and of its intersection with $F$.  The last assertion is the
pigeonhole principle applied to the finite reachable subset of
$\mathcal P(Q)$.
\end{proof}

The repeated-block equality and terminal nonemptiness equivalence are
machine-checked for arbitrary finite cut-word relations.
\Sverified{Mettapedia.GraphTheory.FourColor.GoertzelV24RawNooseCountPumping.terminal_nonempty_congr}
The exact support cardinality and its pigeonhole consequence are checked
separately.
\Sverified{Mettapedia.GraphTheory.FourColor.GoertzelV24RawNooseCountPumping.card_supports_set}
\Sverified{Mettapedia.GraphTheory.FourColor.GoertzelV24RawNooseCountPumping.exists_eq_prefixSupport}

This theorem uses no self-loop and no accepting witness.  It is the direct
zero-count pumping law advertised by the source's monoidal Count semantics.
It is also exact about its price: raw-profile reachability alone does not
justify it, and the worst-case state count is $2^{|Q|}$ before reachable-state
minimization.

\begin{theorem}[unrestricted support right languages separate states]
\label{thm:unrestricted-support-nerode-separation}
\Sverified{Mettapedia.GraphTheory.FourColor.GoertzelV24SupportNerodeLowerBound.code_injective_of_exact_suffix_experiments}
Let $W$ be a finite set of boundary words.  Suppose a deterministic encoding
$e\colon\mathcal P(W)\to X$ and an observer $A\colon X\times\mathcal P(W)\to
\{\mathsf{false},\mathsf{true}\}$ answer every abstract suffix experiment
exactly:
\[
 A(e(S),T)\quad\Longleftrightarrow\quad S\cap T\ne\varnothing
 \qquad(S,T\subseteq W).
\]
Then $e$ is injective.  Hence, for $W=\mathcal C^k$ with
$|\mathcal C|=3$, every such encoding has at least
$2^{3^k}$ states.
\end{theorem}

\begin{proof}
If $S\ne S'$, choose a word $x$ in their symmetric difference.  The singleton
suffix $T=\{x\}$ accepts exactly one of $S,S'$, because
$S\cap\{x\}\ne\varnothing$ is equivalent to $x\in S$.  Thus equal codes
would give contradictory answers to this suffix experiment, so $e$ is
injective.  The powerset of $\mathcal C^k$ has cardinality
$2^{|\mathcal C^k|}=2^{3^k}$.
\end{proof}

The cardinal specialization is checked separately.
\Sverified{Mettapedia.GraphTheory.FourColor.GoertzelV24SupportNerodeLowerBound.two_pow_three_pow_le_card_code}
This lower bound says that the raw powerset is not merely a wasteful encoding
of the \emph{unrestricted} support algebra.  It does not say that every
singleton suffix, or every raw support, is realized by a planar cubic tangle.
A smaller automaton remains possible after restricting both prefix and suffix
experiments to the physically reachable language.  Consequently this theorem
does not discharge the high-width supply: that would require a proved bound on
the reachable Myhill--Nerode quotient, or a genuinely two-dimensional
replacement theorem.

There is a second boundary on that proposed minimization.  Let
$\{L_a:a\in A\}$ and $\{R_b:b\in B\}$ be indexed families of literal open
tangles with one common ordered seam and fixed orientation-reversing matching,
and define the physical right language
\[
 \mathcal L(R_b)=\{a\in A:L_a\mathbin{\#}R_b
                       \text{ is Tait-colourable}\}.
\]
The physical/semantic bridge gives, exactly,
\[
 a\in\mathcal L(R_b)\quad\Longleftrightarrow\quad
 \operatorname{Supp}_{\rm out}(L_a)\cap
 \operatorname{Supp}_{\rm in}(R_b)\ne\varnothing.
\]
\Sverified{Mettapedia.GraphTheory.FourColor.GoertzelV24PhysicalRightLanguageBoundary.mem_physicalRightLanguage_iff_supports_meet}
Consequently
\[
 \bigl(\forall b,\ \mathcal L(R_b)=A\bigr)
 \quad\Longleftrightarrow\quad
 \bigl(\forall a,b,\ L_a\mathbin{\#}R_b
                    \text{ is Tait-colourable}\bigr).
\]
\Sverified{Mettapedia.GraphTheory.FourColor.GoertzelV24PhysicalRightLanguageBoundary.all_physicalRightLanguages_universal_iff_all_splices_taitColorable}
Thus physical suffix restriction can genuinely identify raw support states,
but universality of the right languages on a target-complete family is the
colourability target itself, not an independent compression lemma.  Any
reachable quotient used in the descent must therefore be certified from its
literal transition table and terminal predicate; it cannot be justified just
by saying that all testers are physical.

The smallest genuinely two-dimensional repair one might try is to replace
one hexagonal cell at a time.  Its exact boundary behaviour is already large
enough to rule out that local move.

Let \(H_6\) be the open tangle consisting of a six-cycle with one dangling
port at each cycle vertex, the ports ordered cyclically.  If \(b_i\) is the
colour on port \(i\), and \(x_i\) is the colour on the cycle edge leaving
vertex \(i\), properness at that vertex is equivalent, in
\(\mathbb F_2^2\), to
\[
                         x_{i-1}+b_i+x_i=0.
\]
Thus, after choosing the nonzero colour \(x_5\), the other cycle-edge
colours are forced successively; the word extends precisely when none of
them is zero and the last equation closes the cycle.  The literal port-tangle
semantics and this recurrence are machine-checked to be equivalent.
\Sverified{Mettapedia.GraphTheory.FourColor.GoertzelV24HexCellPortTangle.nonempty_col_hexPortTangle_iff}
Exactly \(63\) of the \(3^6=729\) ordered boundary words extend across
\(H_6\).
\Sverified{Mettapedia.GraphTheory.FourColor.GoertzelV24HexCellPortTangle.card_hexPortSupport}

\begin{proposition}[a facial cycle has no smaller exact connected replacement]
\label{prop:facial-cycle-no-smaller-exact-replacement}
\Sprose
For \(n\geq4\), let \(H_n\) be the open tangle consisting of a labelled
\(n\)-cycle with one cyclically ordered port at each cycle vertex.  No
connected cubic \(n\)-port tangle with fewer than \(n\) internal vertices
has the same ordered boundary support as \(H_n\).
\end{proposition}

\begin{proof}
First count the support of \(H_n\).  A proper colouring of its cycle edges
is a cyclic word \(x=(x_0,\ldots,x_{n-1})\) in the three nonzero Klein
colours with \(x_{i-1}\ne x_i\).  It induces the port word
\[
                             b_i=x_{i-1}+x_i.
\]
Conversely this equation is exactly the local Tait equation at every cycle
vertex, so the image of this derivative map is
\(\operatorname{Supp}(H_n)\).

The number of proper colourings of a labelled \(n\)-cycle with three colours
is
\[
                         2^n+2(-1)^n.
\]
For completeness, fix the colour at one vertex.  If \(a_j\) and \(b_j\)
count length-\(j\) proper paths ending respectively at that colour and at
either one specified other colour, then
\[
 (a_{j+1},b_{j+1})=(2b_j,a_j+b_j),\qquad (a_0,b_0)=(1,0).
\]
Induction gives \(a_j=(2^j+2(-1)^j)/3\); multiplying by the three choices
of the fixed colour gives the displayed cycle count.

If two cycle colourings \(x,y\) induce the same boundary word, then
\[
 x_{i-1}+x_i=y_{i-1}+y_i,
 \qquad\text{hence}\qquad
 x_i+y_i=x_{i-1}+y_{i-1}.
\]
Thus \(y_i=x_i+c\) for one constant Klein colour \(c\).  When \(x\) uses
all three nonzero colours, nonzeroness of \(y\) forces \(c=0\), so the
derivative fibre is a singleton.  A proper cycle colouring using only two
colours exists exactly when \(n\) is even.  There are then six such
alternating colourings; translation by the omitted third colour pairs them
into three derivative fibres.  Therefore
\Sverified{Mettapedia.GraphTheory.FourColor.GoertzelV24FacialCycleBoundaryDerivative.boundaryDerivative_eq_iff_exists_translation}
\Sverified{Mettapedia.GraphTheory.FourColor.GoertzelV24FacialCycleBoundaryDerivative.eq_of_boundaryDerivative_eq_of_usesAll}
\[
|\operatorname{Supp}(H_n)|=
 \begin{cases}
   2^n-1,&n\ \text{even},\\
   2^n-2,&n\ \text{odd}.
 \end{cases}
\]
For \(n=6\) this recovers the checked value \(63\).

Now let a proposed connected cubic replacement have \(r<n\) internal
vertices and \(m\) internal edges.  Counting internal incidences and the
\(n\) ports gives
\[
                             3r=2m+n.
\]
Connectivity gives \(m\geq r-1\), hence \(r\geq n-2\); the incidence
equation gives \(r\equiv n\pmod2\).  Since \(r<n\), necessarily
\(r=n-2\) and \(m=r-1\).  Its internal multigraph is therefore a tree.

Root that tree.  A proper colouring has at most \(3!=6\) choices for the
three labelled incidences at the root.  At each of the other \(r-1\)
vertices the parent-edge colour is fixed, and the two remaining labelled
incidences have at most two possible orders.  Hence the replacement has at
most
\[
                    6\cdot2^{r-1}=6\cdot2^{n-3}
                    =3\cdot2^{n-2}
\]
proper colourings, and no more boundary words than colourings.  For
\(n\geq4\),
\[
                3\cdot2^{n-2}<2^n-2
                \leq |\operatorname{Supp}(H_n)|.
\]
\Sverified{Mettapedia.GraphTheory.FourColor.GoertzelV24FacialCycleBoundaryDerivative.treeColouringUpper_lt_cycleSupportLower}
The two ordered supports cannot be equal.
\end{proof}

This is a generic perimeter obstruction, not wall exclusion.  It rules out
an exact support-preserving reduction which replaces a single facial cycle
by a smaller connected cubic tangle on the same ports; increasing the face
length does not repair the deficit.  It does not rule out replacing a larger
two-dimensional patch, using a target-aware quotient weaker than full
support equality, or performing a move whose size decrease occurs elsewhere.
Those are the remaining source-faithful possibilities.

For a hexagonal cell, the stronger monotone replacement question can also be
settled completely.

\begin{theorem}[no smaller connected cubic monotone replacement for a hexagon]
\label{thm:hexagon-no-smaller-tree-support}
\Sprose
Let \(R\) be a connected plane cubic tangle with six boundary ports in a
specified cyclic order, all on the outer face.  If \(R\) has fewer than six
internal vertices, then
\[
       \operatorname{Supp}(R)\not\subseteq\operatorname{Supp}(H_6).
\]
Thus no smaller connected cubic six-port tangle can replace one hexagonal
cell by the monotone support-inclusion criterion.
\end{theorem}

\begin{proof}
Let \(r\) and \(m\) be the numbers of internal vertices and internal edges of
\(R\).  Counting incidences at the cubic internal vertices gives
\[
                             3r=2m+6.
\]
Connectivity gives \(m\ge r-1\), hence \(r\ge4\), while the incidence
equation gives \(r\equiv6\pmod2\).  Since \(r<6\), necessarily \(r=4\) and
\(m=3\).  The internal graph is therefore a tree.

There are only two abstract trees on four vertices: the path and the tripod.
At an endpoint of the path two ports occur consecutively on the outer face;
at either internal path vertex one port occurs.  Up to a dihedral change of
the six cyclic port positions, the two singleton ports are either separated
by the two endpoint blocks or adjacent to one another.  These are the two
path types.  For the tripod, the central vertex has no port and each leaf has
two consecutive ports, giving one further type.  Thus every \(R\), after
renaming internal vertices and applying a dihedral boundary relabelling, has
one of the following port-block orders:
\[
    S,P,S,P,\qquad P,S,S,P,\qquad P,P,P,
\]
where \(S\) is a singleton and \(P\) is the pair at a tree leaf.

For these three types respectively, the following boundary words are
properly realised by the tree but do not extend across the hexagon:
\[
 (r,r,b,r,r,b),\qquad
 (r,b,r,r,r,b),\qquad
 (r,b,r,p,b,p).
\]
For the two path types, order the internal edges as \((01,02,13)\) and colour
them \((b,p,p)\).  For the tripod, colour the three centre--leaf edges
\((p,b,r)\).  Inspection at the four internal vertices gives the three
distinct colours in every incident triple.  The six-step hexagon recurrence
rejects each displayed boundary word.  These three finite checks, stated as
exact support relations rather than drawings, are machine-checked by
\Sverified{Mettapedia.GraphTheory.FourColor.GoertzelV24HexagonTreeReplacementObstruction.every_planeTreeShape_has_nonextension}
and
\Sverified{Mettapedia.GraphTheory.FourColor.GoertzelV24HexagonTreeReplacementObstruction.no_treeSupport_subset_hexagonRawSupport}.
Finally, hexagon extension is invariant under cyclic rotation and reflection
of the boundary, so the three canonical failures cover every plane embedding
of the tree.
\end{proof}

This closes the one-cell monotone loophole left by
Proposition~\ref{prop:facial-cycle-no-smaller-exact-replacement}: equality was
not the issue.  The source-faithful high-width repair must replace a larger
two-dimensional patch, exploit the restricted support language of the
\emph{actual} exterior, or make its strict size decrease elsewhere.

\begin{theorem}[monotone and target-aware support replacement]
\label{thm:target-aware-support-replacement}
Let \(E\) be a fixed exterior tangle, \(P\) the piece removed from it, and
\(R\) a replacement with the same ordered interface.  If
\[
                  \operatorname{Supp}(R)\subseteq
                  \operatorname{Supp}(P),
\]
then
\[
       E\mathbin{\#}P\text{ not Tait-colourable}
       \quad\Longrightarrow\quad
       E\mathbin{\#}R\text{ not Tait-colourable}.
\]
\Sverified{Mettapedia.GraphTheory.FourColor.GoertzelV24TargetAwareSupportReplacement.not_closedColorable_of_innerSupport_subset}
For this one fixed exterior, the sharp condition is
\[
 E\mathbin{\#}R\text{ not Tait-colourable}
 \quad\Longleftrightarrow\quad
 \operatorname{Supp}(E)\cap\operatorname{Supp}(R)=\varnothing.
\]
\Sverified{Mettapedia.GraphTheory.FourColor.GoertzelV24TargetAwareSupportReplacement.not_closedColorable_iff_disjoint}
Both statements hold for the literal sewn rotation systems after restricting
to genuine nonzero Tait words.
\Sverified{Mettapedia.GraphTheory.FourColor.GoertzelV24TargetAwareSupportReplacement.not_composeRotationSystem_taitColorable_of_taitInnerSupport_subset}
\end{theorem}

\begin{proof}
By the physical gluing bijection, a colouring of a sewn tangle is exactly a
boundary word realized on both sides.  If a colouring of
\(E\mathbin{\#}R\) existed under the displayed inclusion, its seam word
would also lie in \(\operatorname{Supp}(P)\), and hence would colour
\(E\mathbin{\#}P\), a contradiction.  Without requiring inclusion, the
same gluing statement says directly that a colouring exists exactly when the
two fixed support sets meet.
\end{proof}

The first nontrivial physical restriction on the exterior support can now be
made exact.  Say that a support language is \emph{single-orbit} if all of its
words differ only by one global permutation of the three nonzero Tait colour
names (equivalently, by a colour equivalence fixing \(0\)).

\begin{theorem}[a single-orbit exterior admits a strict hexagon replacement]
\label{thm:hexagon-single-orbit-replacement}
Let \(E\subseteq\mathcal C^6\) be the ordered boundary support of the fixed
exterior of one facial hexagon, and suppose that \(E\) is single-orbit.  If

\[
                     E\cap\operatorname{Supp}(H_6)=\varnothing,
\]

then one of the three canonical plane cubic trees on four internal vertices
has support disjoint from \(E\).  Replacing the six-vertex hexagonal cell by
that tree therefore produces a strictly smaller zero-Count splice.
\Sverified{Mettapedia.GraphTheory.FourColor.GoertzelV24HexagonSingleOrbitReplacement.exists_strict_tree_replacement_of_singleGlobalColorOrbit}
\end{theorem}

\begin{proof}
Write \(T_s,T_a,T_t\) for the supports of the separated-path,
adjacent-path, and tripod replacements classified in
Theorem~\ref{thm:hexagon-no-smaller-tree-support}.  Their common intersection
is empty.  There is a short structural proof, avoiding a boundary-word
census.

Let \(w=(w_0,\ldots,w_5)\) extend through both path trees.  Write
\(s_0,s_1,s_2\) for the three internal colours of the separated path and
\(a_0,a_1,a_2\) for those of the adjacent path, in the edge orders used in
the preceding theorem.  The shared leaf with port colours \(w_4,w_5\) gives
\(s_1=a_2\), because two entries of a proper three-colour triple determine
the third.  Comparing the two proper triples at \(w_3\) gives equality of
their remaining unordered colour pairs.  The exchanged alternative would
give \(s_0=a_2=s_1\), impossible at the vertex carrying
\((w_0,s_0,s_1)\).  Hence

\[
                         s_0=a_0,\qquad s_2=a_2=s_1.
\]

Next compare the two triples at \(w_0\).  Again the remaining pairs agree up
to exchange.  The exchanged alternative would give \(s_1=w_1\), impossible
at the separated-path vertex carrying \((w_1,w_2,s_2)\), since
\(s_2=s_1\).  Thus \(s_1=a_1\).  The adjacent-path triples at \(w_2,w_3\)
are consequently \((w_2,s_0,s_1)\) and \((w_3,s_0,s_1)\), whereas the
separated path has \((w_0,s_0,s_1)\).  Uniqueness of the remaining colour
gives

\[
                              w_2=w_0=w_3.
\]

But any tripod extension has a proper triple
\((w_2,w_3,t_1)\) at one leaf, and therefore requires \(w_2\ne w_3\).  No
word lies in all three supports.  The two path implication, tripod
implication, and empty intersection are checked respectively by
\Sverified{Mettapedia.GraphTheory.FourColor.GoertzelV24HexagonSingleOrbitReplacement.path_extensions_force_port_two_eq_three},
\Sverified{Mettapedia.GraphTheory.FourColor.GoertzelV24HexagonSingleOrbitReplacement.tripod_extension_forces_port_two_ne_three},
and
\Sverified{Mettapedia.GraphTheory.FourColor.GoertzelV24HexagonSingleOrbitReplacement.no_word_mem_all_treeSupports}.

Each \(T_i\) is invariant under global colour renaming.  If \(E\) met all
three \(T_i\), choose one word from each intersection and rename each back to
the common orbit representative of \(E\).  The representative would then
belong to \(T_s\cap T_a\cap T_t\), a contradiction.  Hence some \(T_i\) is
disjoint from \(E\).  The physical support-gluing equivalence in
Theorem~\ref{thm:target-aware-support-replacement} says exactly that the
corresponding replacement remains zero-Count, and it removes two internal
vertices.
\end{proof}

Thus a locally irreducible zero-target wall cannot present a single global
colour orbit to every hexagonal cell.  At each cell where all three strict
tree reductions fail, its actual exterior must realize at least two distinct
global colour orbits.  This is a genuine restriction supplied by the target,
not by wall geometry.  It does not yet exclude a wall: the remaining M1 task
is to show that a large common web cannot maintain this multi-orbit boundary
diversity everywhere, or else to use that diversity to construct a larger
two-dimensional replacement.

\begin{theorem}[unrestricted target-awareness collapses to support inclusion]
\label{thm:target-aware-quantifier-collapse}
Let \(U\) be the boundary-word universe, let \(P\subseteq U\) be the support
of the removed piece, and let \((R_i)_{i\in I}\) be the supports of all
candidate replacements.  Then
\[
 \bigl[
   \forall E\subseteq U,\ E\cap P=\varnothing
   \Longrightarrow
   \exists i\in I,\ E\cap R_i=\varnothing
 \bigr]
 \quad\Longleftrightarrow\quad
 \bigl[\exists i\in I,\ R_i\subseteq P\bigr].
\]
\Sverified{Mettapedia.GraphTheory.FourColor.GoertzelV24TargetAwareReplacementQuantifiers.unrestricted_targetAware_iff_exists_subset}
\end{theorem}

\begin{proof}
For the forward implication, test the maximal abstract zero exterior
\(E=U\setminus P\).  A candidate disjoint from \(E\) is contained in \(P\).
Conversely, if \(R_i\subseteq P\), then every set disjoint from \(P\) is
also disjoint from \(R_i\).  The same equivalence holds for a restricted
class of physical exterior supports whenever that class itself contains
\(U\setminus P\).
\Sverified{Mettapedia.GraphTheory.FourColor.GoertzelV24TargetAwareReplacementQuantifiers.physical_targetAware_iff_exists_subset_of_compl_mem}
\end{proof}

Retaining the full corridor profile does not by itself evade this quantifier
boundary.  This point is easy to blur because the richer state really is
essential for advancing through the next Cell.  It is nevertheless invisible
to the final colouring splice.

\begin{theorem}[rich-profile gluing erases to the boundary word]
\label{thm:target-aware-profile-erasure}
Let \(X\) and \(Y\) be arbitrary profile types with boundary-word projections
\(\pi_X:X\to U\) and \(\pi_Y:Y\to U\).  For profile languages
\(A\subseteq X\) and \(B\subseteq Y\), define them to be compatible when some
\(a\in A\) and \(b\in B\) have \(\pi_X(a)=\pi_Y(b)\).  Then
\[
 A\text{ and }B\text{ are compatible}
 \quad\Longleftrightarrow\quad
 \pi_X[A]\cap\pi_Y[B]\ne\varnothing.
\]
\Sverified{Mettapedia.GraphTheory.FourColor.GoertzelV24TargetAwareProfileErasure.compatible_iff_projectedSupports_meet}
In particular, profile languages with equal projected word support are
indistinguishable by every exterior profile language.
\Sverified{Mettapedia.GraphTheory.FourColor.GoertzelV24TargetAwareProfileErasure.compatible_iff_of_projectedSupport_eq}

If \((B_i)_{i\in I}\) is a family of rich replacement-profile languages and
the exterior ranges over arbitrary word supports, target-aware safety again
holds exactly when
\[
                 \pi_Y[B_i]\subseteq\pi_X[A]
\]
for some \(i\).
\Sverified{Mettapedia.GraphTheory.FourColor.GoertzelV24TargetAwareProfileErasure.unrestricted_profileTargetAware_iff_exists_projected_subset}
\end{theorem}

\begin{proof}
A compatible pair of profiles supplies their common projected word, and a
word in both projected images supplies a compatible pair of profiles.  This
proves the first equivalence and the indistinguishability statement.  Apply
Theorem~\ref{thm:target-aware-quantifier-collapse} to the two projected word
supports for the last assertion.  Equivalently, the complement of
\(\pi_X[A]\) is still the maximal abstract exterior adversary.
\Sverified{Mettapedia.GraphTheory.FourColor.GoertzelV24TargetAwareProfileErasure.projected_complement_adversary_of_no_subset}
\end{proof}

This changes the correct formulation of the two-dimensional crux.  Exact
support equality is useful for pumping because it works against every suffix,
but a minimal-counterexample descent only needs a replacement whose support
misses the \emph{actual} exterior.  Conversely, if the original splice is
uncolourable while a proposed smaller admissible splice is colourable, then
the replacement support cannot be contained in the original support:
\Sverified{Mettapedia.GraphTheory.FourColor.GoertzelV24TargetAwareSupportReplacement.not_innerSupport_subset_of_original_zero_of_replacement_colorable}
it must introduce some boundary word that the removed piece did not realize.
Thus Proposition~\ref{prop:facial-cycle-no-smaller-exact-replacement} does not
kill target-aware one-cell reduction; it shows precisely why such a reduction
cannot be justified by exact support preservation.  Theorem~
\ref{thm:target-aware-quantifier-collapse} adds an essential quantifier
warning: allowing the replacement to depend on an \emph{arbitrary} exterior
does not help either.  The gain can come only from a structural theorem
restricting which support sets are realized by actual planar exteriors.
Theorem~\ref{thm:target-aware-profile-erasure} adds that merely enriching each
realization by its Kempe-component pairing or face-progress coordinates does
not supply such a restriction: those coordinates matter only insofar as their
\emph{globally realizable family} constrains the projected word support.
The remaining constructive problem is to exploit that physical restriction
while preserving the topological class and decreasing size.

\begin{theorem}[raw noose Count pumping]
\label{thm:raw-noose-count-pumping}
\Sprose
Let a closed cubic map be cut by a chain of pairwise nested ordered nooses,
each crossing exactly \(k\) edges, into a left end, serial slabs, and a right
end.  Identify every ordered cut word with an element of
\(\mathcal C^k\), where \(|\mathcal C|=3\).  For the prefix ending at cut
\(i\), let
\[
 S_i=\{x\in\mathcal C^k:
       \text{the prefix has a proper nonzero Tait colouring with word }x\}.
\]
Then:
\begin{enumerate}
\item each slab acts on supports by the Boolean relation of its physical
port-tangle Count;
\item there are at most \(2^{3^k}\) possible cumulative supports \(S_i\);
and
\item if \(S_i=S_j\) for \(i<j\), deleting the intervening slabs and gluing
cut \(i\) to cut \(j\) preserves whether the complete closed map has a Tait
colouring.
\end{enumerate}
Thus more than \(2^{3^k}\) cut positions force a Count-preserving repeated
support.  If the original map has zero accepting Count and the cuts also
satisfy the geometric hypotheses of
Theorem~\ref{thm:physical-noose-deletion}, that repetition produces a
strictly smaller connected bridgeless cubic spherical zero-Count instance.
\end{theorem}

\begin{proof}
Cutting a slab open gives a port tangle with \(k\) ordered input and \(k\)
ordered output incidences.  The physical restriction/gluing bijection (C1)
identifies its Count support with a relation
\[
                         R_i\subseteq\mathcal C^k\times\mathcal C^k.
\]
Consequently \(S_{i+1}=R_i[S_i]\).  Since
\(|\mathcal C^k|=3^k\), its powerset has cardinality \(2^{3^k}\).
Pigeonhole gives equal supports when there are more cut positions.
Theorem~\ref{thm:determinized-support-deletion} then applies the unchanged
right-end suffix to those equal supports and proves equality of terminal
acceptance.  Finally Theorem~\ref{thm:physical-noose-deletion} supplies all
the asserted graph-class and strict-size conclusions.
\end{proof}

All algebraic assertions in this theorem are now machine-checked: physical
serial gluing becomes relational composition, its matrix support is the same
relation, and the cut alphabet has exactly three letters per crossed edge.
\Sverified{Mettapedia.GraphTheory.FourColor.GoertzelV24RawNooseCountPumping.supportRel_series}
\Sverified{Mettapedia.GraphTheory.FourColor.GoertzelV24RawNooseCountPumping.support_countMatrix_eq_supportRel}
\Sverified{Mettapedia.GraphTheory.FourColor.GoertzelV24RawNooseCountPumping.card_cutWord}
The final spherical-map sentence remains an application of the physical
noose theorem, whose representation-specific Lean constructor is D3 rather
than part of this algebra module.

This theorem is deliberately insensitive to the number or internal size of
the slabs.  It needs no finite alphabet of physical Cell graphs and no
globally chosen colouring.  The only geometric information used by the
physical replacement is carried by D1: the actual ordered nooses, connected
saturated retained sides, a nonempty deleted region, and the common finite
seam type.  The richer rooted Cell carrier remains useful only to reduce the
astronomical conservative bound \(2^{3^k}\), thereby making D4's finite
floor computationally realistic.

\begin{lemma}[two boundary-essential sides glue bridgelessly]
\label{lem:boundary-essential-gluing}
\Sverified{Mettapedia.GraphTheory.FourColor.GoertzelV24BoundaryEssentialGluing.glue_bridgeless}
Let \(H\) and \(K\) be finite connected graphs, each with \(k\ge2\) ordered
boundary half-edges.  Call a side \emph{boundary-essential} if, for every
bridge of the side, both components left after deleting that bridge contain
a boundary half-edge.  Glue the \(i\)-th half-edge of \(H\) to the \(i\)-th
half-edge of \(K\) for every \(i\).  Then the glued graph is connected and
bridgeless.
\end{lemma}

\begin{proof}
Connectivity follows because each side is connected and at least one seam
edge joins them.  Consider an old edge \(e\) of \(H\).  If \(e\) is not a
bridge of \(H\), an \(H\)-path already bypasses it.  If it is a bridge, its
two components contain boundary ports \(i\) and \(j\).  Travel from one end
of \(e\) inside its component to port \(i\), cross to \(K\), use
connectedness of \(K\) to reach port \(j\), and cross back to the other
component.  This is a path between the ends of \(e\) avoiding \(e\).
The same argument handles old edges of \(K\).

Finally take a seam edge \(i\).  Choose \(j\ne i\).  Connectedness of \(H\)
and \(K\), together with seam edge \(j\), gives a path between the ends of
seam edge \(i\) which avoids it.  Thus no edge of the glued graph is a
bridge.
\end{proof}

\begin{lemma}[a saturated side of a bridgeless graph is boundary-essential]
\label{lem:saturated-side-boundary-essential}
\Sverified{Mettapedia.GraphTheory.FourColor.GoertzelV24BoundaryEssentialGluing.inducedSide_boundaryEssential_of_saturated}
Let \(G\) be bridgeless and let \(H\) be a connected vertex side of an edge
cut, with every edge from \(H\) to \(G-H\) represented by a boundary
half-edge.  Then \(H\) is boundary-essential.
\end{lemma}

\begin{proof}
If a bridge \(e\) of \(H\) left a component \(C\) containing no boundary
half-edge, saturation says that no edge of \(G\) leaves \(C\) except \(e\).
Then \(e\) would be a bridge of \(G\), a contradiction.
\end{proof}

\paragraph{Hub closures and the exact face count.}
Let \(T\) be an open rotation tangle with boundary set \(B\).  Its
\emph{hub closure} \(\widehat T_\rho\) has one new hub dart
\(\bar b\) for every \(b\in B\), pairs \(b\) with \(\bar b\), and
rotates the hub darts by a permutation \(\rho\) giving their boundary
order.  The old vertex rotations and old internal pairings are unchanged.
Write
\[
 \widehat\alpha_T(b)=\bar b,
 \qquad
 \widehat\alpha_T(\bar b)=b,
 \qquad
 \widehat\phi_{T,\rho}=\widehat\rho_T\widehat\alpha_T.
\]
We call the cap \emph{face-simple} when distinct hub darts lie in distinct
\(\widehat\phi_{T,\rho}\)-orbits.  This is the exact combinatorial
condition called \texttt{HubFacesDistinct} in the implementation; it must
not be inferred merely from the hub being one vertex.

\begin{lemma}[orbit exchange]
\label{lem:orbit-exchange}
\Sverified{Mettapedia.GraphTheory.FourColor.GoertzelV24OrbitExchange.orbitCount_exchange_univ}
Let a finite family of permutations be obtained by processing ports one at
a time.  At each port first transpose two points in different current
orbits, and then transpose two distinct points in the resulting common
orbit.  Suppose the unprocessed merge orbits remain pairwise distinct.
Then processing all ports does not change the number of orbits.
\end{lemma}

\begin{proof}
A transposition of points in different cycles merges precisely those two
cycles, decreasing the orbit count by one.  A transposition of distinct
points in the same cycle cuts that cycle into two, increasing the count by
one.  Thus one exchange has net change zero.  The stated invariant ensures
that the same hypotheses hold after each processed port, so induction on the
finite port set proves the result.
\end{proof}

\begin{theorem}[the hub-capped seam face formula]
\label{thm:hub-seam-face-formula}
\Sverified{Mettapedia.GraphTheory.FourColor.GoertzelV24CompositeSphericity.card_orbitFace_composite_add}
Let \(L,R\) be open rotation tangles with equally large boundary sets,
let \(m:B_L\simeq B_R\) match their ports, and give their hubs rotations
\(\rho_L,\rho_R\).  Assume both hub closures are face-simple and
\[
             \rho_R(mb)=m(\rho_L^{-1}b)
             \qquad(b\in B_L).                 \tag{25}
\]
If \(k=|B_L|\), and \(L\#_mR\) is the closed rotation system obtained by
pairing matched real boundary darts, then
\[
 F(L\#_mR)+k=F(\widehat L_{\rho_L})+
                    F(\widehat R_{\rho_R}).    \tag{26}
\]
Here every \(F\) is the number of orbits of the face permutation
\(\phi=\rho\alpha\).
\end{theorem}

\begin{proof}
Begin with the disjoint union of the two hub closures.  For each left port
\(b\), replace the two cap pairings
\[
  (b,\bar b),\qquad (m b,\overline{m b})
\]
by the real seam pairing \((b,m b)\) and the hub pairing
\((\bar b,\overline{m b})\).  On the face permutation this replacement is
one merging transposition followed by one splitting transposition.  The
merge joins the left hub face at \(b\) to the right hub face at \(m b\).
Face-simplicity says that merge faces belonging to distinct unprocessed
ports are distinct.  The split belonging to \(b\) lies on the face created
by the merge belonging to \(\rho_L^{-1}b\); equation~(25) is exactly the
identity needed on the right hub to identify those two points.  Therefore
Lemma~\ref{lem:orbit-exchange} applies and the fully exchanged face
permutation has
\[
 F(\widehat L_{\rho_L})+F(\widehat R_{\rho_R})
\]
orbits.  This intermediate equality is also checked directly by
\texttt{GoertzelV24SeamExchange.orbitCount\_phiS\_univ}.

After all ports have been exchanged, real darts never map to hub darts and
hub darts never map to real darts.  On the real summand the permutation is,
up to the carrier reassociation, the face permutation of \(L\#_mR\).  On
the \(2k\) hub darts it is the face permutation of the two-hub \(k\)-edge
dipole:
\[
 \bar b_L\longmapsto \rho_R(\overline{m b})_R,
 \qquad
 \bar c_R\longmapsto
       \rho_L(\overline{m^{-1}c})_L.
\]
Equation~(25) makes this permutation an involution.  It is fixed-point-free
because it exchanges the left and right summands, so its \(2k\) elements
form exactly \(k\) two-cycles.  Orbit counts add over disjoint sums and are
invariant under the carrier reassociation.  Subtracting these \(k\) dipole
orbits gives~(26).  The dipole count and the reassociation are machine-checked
in \texttt{GoertzelV24CompositeSphericity.orbitCount\_dipolePhi} and in the
theorem named by the verification badge.
\end{proof}

\begin{corollary}[hub-capped seam sphericity]
\label{cor:hub-seam-sphericity}
\Sverified{Mettapedia.GraphTheory.FourColor.GoertzelV24CompositeSphericity.composite_euler}
Under the hypotheses of Theorem~\ref{thm:hub-seam-face-formula}, suppose
both hub closures are spherical.  Then the composite satisfies
\[
       V(L\#_mR)+F(L\#_mR)=E(L\#_mR)+2.
\]
\end{corollary}

\begin{proof}
Let \(e_L,e_R\) be the numbers of intact edges of the two open sides and
write \(F_L,F_R\) for the face counts of their hub closures.  The two cap
equalities are
\[
 (V_L+1)+F_L=(e_L+k)+2,
 \qquad
 (V_R+1)+F_R=(e_R+k)+2.
\]
The composite has \(V_L+V_R\) vertices and \(e_L+e_R+k\) edges, while
Theorem~\ref{thm:hub-seam-face-formula} gives
\(F=F_L+F_R-k\).  Adding the two cap equalities and substituting these three
counts yields
\[
 (V_L+V_R)+F=(e_L+e_R+k)+2,
\]
as required.
\end{proof}

\begin{lemma}[opposite first-return orders give the seam orientation]
\label{lem:canonical-hub-orientation}
\Sverified{Mettapedia.GraphTheory.FourColor.GoertzelV24VertexSideOpenTangle.orientationReversing_canonicalHubRotation_of_opposite}
For an open shore \(T\), let \(s_T\) be the first-return permutation on its
boundary darts and choose the canonical hub rotation
\(\rho_T=s_T^{-1}\).  Let \(m:B_L\to B_R\) match two boundary interfaces.
If
\[
        s_R\bigl(m(s_L(b))\bigr)=m(b)
        \qquad(b\in B_L),
\]
then \(m\) is orientation-reversing for the two canonical hub rotations:
\[
        \rho_R(m(b))=m(\rho_L^{-1}(b)).
\]
\end{lemma}

\begin{proof}
Apply \(s_R^{-1}\) to the displayed opposite-return equation.  Its
left-hand side becomes \(m(s_L(b))\), while
\(\rho_L^{-1}=s_L\) and \(\rho_R=s_R^{-1}\).  This is exactly the required
identity.
\end{proof}

\begin{lemma}[inverse first return makes the cap face-simple]
\label{lem:canonical-hub-face-simple}
\Sverified{Mettapedia.GraphTheory.FourColor.GoertzelV24CanonicalHubClosure.hubFacesDistinct_canonical}
Let \(T\) be a literal vertex shore, let \(s_T\) be first return of its
retained facial walk to the exposed boundary darts, and put
\(\rho_T=s_T^{-1}\).  In the hub closure
\(\widehat T_{\rho_T}\), distinct hub darts lie on distinct faces.
\end{lemma}

\begin{proof}
Transport the hub closure to the carrier consisting of the old retained
darts and one copy of the boundary.  Eliminate the inserted hub darts by
taking first return to the old retained carrier.  An internal dart makes the
same one-step facial move as before.  At a boundary dart \(b\), the walk
visits the hub dart \(\rho_T(b)\) and returns to the old carrier after two
steps.  Consequently the transported return permutation is
\[
  \phi_T^{\mathrm{cap}}\,\widetilde{\rho_T},
\]
where \(\widetilde{\rho_T}\) is \(\rho_T\) on the boundary darts and the
identity elsewhere.  Taking first return once more, now to the boundary,
gives
\[
       s_T\rho_T=s_Ts_T^{-1}=1.                 \tag{27}
\]
Every hub dart is facially adjacent to an old boundary dart.  Thus two hub
darts on the same capped face would give two boundary darts in the same
cycle of the return permutation in~(27).  The identity permutation has only
singleton cycles, so those boundary darts, and hence the original hub darts,
are equal.
\end{proof}

This lemma is deliberately narrower than saying that the cap is one new
vertex.  Equation~(27) proves the facial statement
\texttt{HubFacesDistinct} for every finite boundary set.  Inverse first
return alone does not make \(\rho_T\) one cycle.  For the connected shores
used by the splice, however, the exact side rank forces that conclusion by
Euler's inequality; this is Theorem~\ref{thm:exact-rank-forces-one-hub} below.

\begin{theorem}[exact canonical-hub face count]
\label{thm:canonical-hub-face-count}
\Sverified{Mettapedia.GraphTheory.FourColor.GoertzelV24CanonicalHubClosure.card_canonicalHub_face_eq_boundary_add_ambient}
Let \(T\) be a literal vertex shore with boundary darts \(B\), and let
\(F_{\rm pure}(T)\) be the number of ambient faces all of whose darts lie in
that shore.  Closing \(T\) with the canonical hub rotation
\(\rho_T=s_T^{-1}\) gives
\[
       F(\widehat T_{\rho_T})=|B|+F_{\rm pure}(T).       \tag{28}
\]
\end{theorem}

\begin{proof}
First return of a capped face to the old retained darts partitions the faces
without changing their number: every inserted hub dart returns to the old
carrier after one step.  Split those return cycles according to whether they
meet a boundary dart.  By equation~(27), the boundary-meeting part has one
cycle for each \(b\in B\).  On a cycle avoiding the boundary, the inserted
hub order is never consulted; every step is literally the old ambient facial
step.  Forgetting the retainedness tag therefore gives a bijection from
boundary-avoiding capped faces to ambient faces wholly contained in the
shore.  Adding the two disjoint classes proves~(28).
\end{proof}

\begin{corollary}[canonical hub closure from the exact side rank]
\label{cor:canonical-hub-spherical}
\Sverified{Mettapedia.GraphTheory.FourColor.GoertzelV24CanonicalHubClosure.closedSideSpherical_canonical_of_faceRank}
Let \(V_T\) and \(E_T\) be the retained vertices and intact edges of the
shore.  If
\[
       F_{\rm pure}(T)+|V_T|=|E_T|+1,                 \tag{29}
\]
then its canonical hub closure satisfies the spherical Euler equality.
For a connected vertex shore and connected complement in a connected
graph-backed spherical map, equation~(29), and hence the conclusion, follows
from the planar-bond rank theorem.
\Sverified{Mettapedia.GraphTheory.FourColor.GoertzelV24CanonicalHubClosure.closedSideSpherical_canonical_of_planarBond}
\end{corollary}

\begin{proof}
The hub closure has \(|V_T|+1\) displayed vertices.  Its pairing orbits are
the \(E_T\) intact edges and one hub spoke for every boundary dart, so it has
\(|E_T|+|B|\) edges.  By Theorem~\ref{thm:canonical-hub-face-count} it has
\(|B|+F_{\rm pure}(T)\) faces.  Substitution of~(29) gives
\[
 (|V_T|+1)+(|B|+F_{\rm pure}(T))
      =(|E_T|+|B|)+2.
\]
The final assertion is exactly the connected-side equality already obtained
by splitting spherical Euler across a planar bond.
\end{proof}

\begin{theorem}[connected exact rank forces one hub]
\label{thm:exact-rank-forces-one-hub}
\Sverified{Mettapedia.GraphTheory.FourColor.GoertzelV24CanonicalHubCyclicity.hubRotation_isCycle_of_closedSideSpherical}
Let \(T\) be a connected cubic open rotation tangle whose old darts at each
vertex form one rotation cycle.  Let \(\rho\) be any permutation of its
boundary darts, and suppose the boundary contains two distinct darts.  If
the cap is assigned the exact one-hub Euler rank
\[
       |V_T|+1+F(\widehat T_\rho)=E(\widehat T_\rho)+2,
\]
then \(\rho\) is one nontrivial cycle on the entire boundary carrier.
\end{theorem}

The inequality used here is itself available directly at map level.  If
\(R\) is a connected finite rotation system and the rotation darts at every
stored vertex form one cycle, then
\[
       \operatorname{orb}(R.\rho)+\operatorname{orb}(R.\phi)
          \le |E(R)|+2.                                  \tag{29a}
\]
\Sverified{Mettapedia.GraphTheory.FourColor.GoertzelV24ConnectedMapEulerBound.orbitCount_rho_add_orbitCount_phi_le_edge_add_two}
Indeed, the edge involution has an explicit transposition-list presentation;
connectedness makes the associated word action transitive, so the general
permutation Euler inequality has exactly one component.  Thus~(29a) carries
no cellular-embedding or genus hypothesis.

\begin{proof}
Present the fixed-point-free hub edge involution as one transposition per
edge.  The old rotation, these transpositions, and the hub spokes act
transitively on the capped dart carrier: cyclic old rotations move within
one old vertex, connectedness moves across the intact side edges, and each
hub dart crosses its spoke to its old boundary dart.  Hence the generic
permutation Euler inequality has one component and gives
\[
 \operatorname{orb}(\rho_T)+\operatorname{orb}(\rho)
       +F(\widehat T_\rho)
       \le E(\widehat T_\rho)+2.
\]
The old rotation has exactly one orbit per old vertex, so
\(\operatorname{orb}(\rho_T)=|V_T|\).  Comparing with the displayed exact
rank yields
\(\operatorname{orb}(\rho)\le1\).  The nonempty boundary gives the reverse
inequality, so \(\operatorname{orb}(\rho)=1\).  With two distinct boundary
darts this unique orbit is nontrivial; therefore \(\rho\) is a single cycle
with full support.
\end{proof}

\begin{corollary}[a connected planar bond supplies the canonical cap]
\label{cor:planar-bond-canonical-cap}
\Sverified{Mettapedia.GraphTheory.FourColor.GoertzelV24CanonicalHubCyclicity.canonicalHubDiscData_of_planarBond}
Let a connected spherical graph-backed cubic map have cyclic vertex
rotations and two-sided faces.  If a vertex bipartition has connected induced
graphs on both sides and at least two boundary darts, then closing either
side by the inverse of its facial boundary first-return permutation gives a
face-simple spherical cap with one cyclic hub vertex.
\end{corollary}

\begin{proof}
The planar-bond rank theorem supplies equation~(29), so
Corollary~\ref{cor:canonical-hub-spherical} supplies the exact one-hub rank.
Lemma~\ref{lem:canonical-hub-face-simple} supplies face-simplicity without
any topological hypothesis.  Theorem~\ref{thm:exact-rank-forces-one-hub}
then proves that the canonical hub rotation is a single cycle.
\end{proof}

\begin{theorem}[complementary planar-bond orders are opposite]
\label{thm:planar-bond-orders-opposite}
\Sverified{Mettapedia.GraphTheory.FourColor.GoertzelV24ComplementaryShoreBoundaryOrder.boundarySuccessors_opposite_of_planarBond}
In the setting of Corollary~\ref{cor:planar-bond-canonical-cap}, let
\(s_S\) and \(s_{\bar S}\) be the actual facial first-return permutations
on the two literal shores.  Match their boundary darts by the ambient edge
flip
\[
                         m(d)=\alpha d.
\]
Then
\[
              s_{\bar S}\bigl(m(s_S(d))\bigr)=m(d).       \tag{29b}
\]
Consequently the two canonical inverse-return hub rotations satisfy the
orientation-reversing seam equation.
\Sverified{Mettapedia.GraphTheory.FourColor.GoertzelV24ComplementaryShoreBoundaryOrder.orientationReversing_canonicalHubRotation_of_planarBond}
Together with the two canonical cap equalities, the natural seam composite
therefore has the complete connected, bridge-free, cubic, rotation-cyclic
spherical-map package.
\Sverified{Mettapedia.GraphTheory.FourColor.GoertzelV24ComplementaryShoreBoundaryOrder.bridgelessSphericalCubicMapData_of_planarBond_canonical}
\end{theorem}

\begin{proof}
Transport the deleted shore's capped facial first return to the retained
shore's boundary carrier by reversing every cut edge.  On that common
carrier the planar-bond boundary theorem says that the retained return
\(s_S\) is the inverse of the transported deleted return.  Transporting
this identity back through \(m\) gives exactly~(29b).

The canonical hub rotations are \(s_S^{-1}\) and
\(s_{\bar S}^{-1}\), so~(29b) is precisely the orientation equation used by
the seam face-count theorem.  Corollary~\ref{cor:planar-bond-canonical-cap}
supplies the two face-simple spherical caps, while the literal-shore
structure lemmas supply connectedness, boundary-essentiality, cubicity, and
cyclic old-vertex rotations.  The structural seam theorem packages these
facts on the literal composite.
\end{proof}

\begin{corollary}[canonical literal-shore splice]
\label{cor:canonical-literal-shore-splice}
\Sverified{Mettapedia.GraphTheory.FourColor.GoertzelV24CanonicalHubClosure.bridgelessSphericalCubicMapData_ofVertexSides_canonical}
Let two connected vertex shores of one bridgeless cubic rotation system have
at least two matched boundary darts.  Suppose their matching reverses the
two boundary first-return orders, both shores satisfy the exact side-rank
equation~(29), and the ambient vertex rotations are cyclic.  Then gluing the
literal shores along the matching produces a connected, bridgeless, cubic,
rotation-cyclic spherical map.
\end{corollary}

\begin{proof}
Lemma~\ref{lem:saturated-side-boundary-essential} supplies boundary
essentiality.  Lemma~\ref{lem:canonical-hub-face-simple} supplies the two
face-simple caps, Lemma~\ref{lem:canonical-hub-orientation} supplies their
orientation reversal, and Corollary~\ref{cor:canonical-hub-spherical}
supplies their spherical Euler equalities.  The structural seam theorem then
gives connectedness, bridge-freeness, cubicity, cyclic old-vertex rotations,
and sphericity simultaneously.
\end{proof}

\begin{definition}[cyclic bond-boundary datum]
\label{def:cyclic-bond-boundary}
\Sverified{Mettapedia.GraphTheory.FourColor.GoertzelV24CyclicBondBoundary.CyclicBondBoundaryData}
Let \(T\) be a literal vertex shore, let \(B_T\) be its boundary darts, and
let \(s_T\) be first return of the retained facial walk to \(B_T\).  A
\emph{cyclic bond-boundary datum} for \(T\) consists of an integer \(k\ge2\)
and a bijection
\[
             \eta:\mathbb Z/k\mathbb Z\longrightarrow B_T
\]
such that
\[
             s_T(\eta(i))=\eta(i+1)                 \tag{30}
\]
for every \(i\in\mathbb Z/k\mathbb Z\).  Thus the datum records both that
the displayed ports exhaust the boundary and that their facial first-return
order is one cyclic order, rather than a disjoint union of several orders.
\end{definition}

\begin{lemma}[a cyclic bond boundary is one canonical hub]
\label{lem:cyclic-bond-one-hub}
\Sverified{Mettapedia.GraphTheory.FourColor.GoertzelV24CyclicBondBoundary.CyclicBondBoundaryData.canonicalHubRotations_singleCycle_of_opposite}
If \(T\) has a cyclic bond-boundary datum, then both \(s_T\) and its inverse
\(\rho_T=s_T^{-1}\) are single cycles on \(B_T\).  Consequently the
canonical closure used above really adds one vertex whose rotation is
\(\rho_T\), not one formal ``hub'' split into several vertices.  This gives
an explicit coordinate certificate for the same conclusion as
Theorem~\ref{thm:exact-rank-forces-one-hub}; it is no longer a hypothesis of
the physical splice.

Suppose moreover that \(m:B_L\to B_R\) obeys the opposite-return equation
\[
       s_R\bigl(m(s_L(b))\bigr)=m(b).                \tag{31}
\]
If the left shore has a cyclic bond-boundary datum, then the right
first-return order is the transported inverse of the left order.  In
particular both canonical hub rotations are single cycles on their entire
boundary carriers and the matching is orientation-reversing.
\end{lemma}

\begin{proof}
Equation~(30) says exactly that
\[
       s_T=\eta\circ(i\mapsto i+1)\circ\eta^{-1}.
\]
The shift on \(\mathbb Z/k\mathbb Z\) is one cycle when \(k\ge2\); conjugacy
preserves being a cycle, and the inverse of a cycle is a cycle.  This proves
the first assertion.

For the second assertion, equation~(31) is equivalent to
\[
       s_R=m\circ s_L^{-1}\circ m^{-1}.
\]
Hence \(s_R\) is conjugate to the inverse of \(s_L\).  Transporting the
enumeration by \(m\) and reversing its indices gives equation~(30) on the
right.  The final orientation statement is
Lemma~\ref{lem:canonical-hub-orientation}.
\end{proof}

\begin{lemma}[a tight transverse noose supplies the cap data]
\label{lem:tight-noose-cap-data}
\Sprose
Let \(M\) be a cellular spherical map and let \(\gamma\) be a simple closed
curve which avoids vertices, crosses \(k\) edges transversely, and is
\emph{tight}: its intersection with the interior of every face is empty or
one connected arc.  Cut the crossed edges at \(\gamma\), retain either
closed-disc side, and close its \(k\) boundary darts by a hub in their cyclic
order.  Each hub closure is spherical and face-simple.  Under the natural
matching of the two halves of every crossed edge, the two hub rotations
satisfy~(25).
\end{lemma}

\begin{proof}
By the Jordan curve theorem, \(\gamma\) cuts the sphere into two open discs
with common boundary \(\gamma\).  Collapse the boundary circle of either
closed disc to a point.  The quotient is a sphere; the cut endpoints become
the darts at the new hub, in their cyclic boundary order.  This is precisely
the hub closure, so its cellular Euler equality is the spherical one.

Enumerate the crossings once around \(\gamma\).  Tightness makes the facial
arc inside the retained disc carry each boundary dart to the next crossing,
so this enumeration is a cyclic bond-boundary datum in the sense of
Definition~\ref{def:cyclic-bond-boundary}.  In particular
Lemma~\ref{lem:cyclic-bond-one-hub} proves at permutation level that the cap
adds one cyclic hub vertex.

The sector at the hub between two consecutive cut darts is the image of the
unique boundary arc of \(\gamma\) between the corresponding crossings.  It
lies in one old face.  Tightness says that no old face contains two such
arcs.  Hence no face of the capped rotation system contains two hub darts,
which is face-simplicity.  Equivalently, the boundary first-return order is
the cyclic order of crossings around the one boundary circle, so its inverse
is a single cycle and Lemma~\ref{lem:canonical-hub-face-simple} supplies the
same face-simplicity directly at permutation level.  Finally, the two boundary orientations induced on
the common boundary of the two discs are opposite.  Thus their computed
first-return maps obey \(s_R\circ m\circ s_L=m\).
Lemma~\ref{lem:canonical-hub-orientation} converts this equation into
equation~(25).
\end{proof}

This tight-noose lemma is now an optional geometric presentation, not the
load-bearing D3 adapter.  For a supplied vertex bipartition with both induced
sides connected, Corollary~\ref{cor:planar-bond-canonical-cap} obtains the
spherical face-simple one-hub cap directly from the planar-bond rank and
Euler's inequality; no Jordan theorem and no per-face tightness are needed.
What the sphere-cut supplier must still provide is the combinatorial data
actually used by the splice: nested connected edge shores and an equality of
ordered seam states which transports the retained first-return order.
Lemma~\ref{lem:connected-shore-literal-node} derives the roots and at least two
boundary darts.  Theorem~\ref{thm:planar-bond-orders-opposite} then makes
the natural edge-flip to the complementary shore orientation-reversing
automatically.  A tight noose is one sufficient way to manufacture the
ordered shore data, but it is not part of the algebraic seam theorem.

\begin{lemma}[one cubic sphere cut can be moved to edges]
\label{lem:cubic-noose-edge-perturbation}
\Sprose
Let \(N\) be a tight sphere-cut noose of a cubic plane graph.  Thus \(N\)
meets the graph only in a middle set \(Z\) of vertices, and at every
\(z\in Z\) each of the two edge sides has at least one incident edge.
After choosing one of the two disc sides, \(N\) can be moved off the
vertices to a tight transverse noose \(\gamma\) crossing at most
\(2|Z|\) edges.  The vertices strictly in the chosen side remain there.
If a laminar family of such nooses has pairwise ordered collars, these moves
can be made inside the collars and preserve laminarity.
\end{lemma}

\begin{proof}
Choose disjoint small vertex discs around the points of \(Z\).  At a cubic
middle vertex, the chosen side contains either one or two consecutive
incident edge germs.  Replace the portion of \(N\) through that vertex by
an arc in its vertex disc pushed toward the chosen side.  It crosses exactly
those one or two germs and no others.  The replacements are disjoint, so the
result is still a simple closed curve, avoids every vertex, preserves every
strictly separated vertex, and has at most \(2|Z|\) transverse crossings.

Two harmless simplifications may now be needed.  If the curve crosses the
same edge twice with no vertex on the intervening edge segment, a
closed regular neighbourhood of that edge segment is a rectangle disjoint
from every vertex and from every other edge.  Replace the curve arc on one
side of the rectangle by the parallel arc on the other side.  This is a
local isotopy followed by cancellation of the two transverse crossings; it
does not move the curve across a graph vertex.  If the curve visits
one face in more than one arc, choose two consecutive arcs in that open
disc.  The subdisc between them contains no graph; replacing the appropriate
curve subarc along that subdisc removes a repeated visit without changing
which graph vertices lie on either side and without introducing an edge
crossing.  Each operation strictly lowers the finite number of crossings or
face visits, so they terminate.  The resulting curve is transverse and
tight, with the same bound.

For an already laminar finite family, choose the vertex discs and the
subsequent face surgeries inside mutually ordered collar neighbourhoods,
starting with the innermost curve.  An isotopy supported in one such collar
cannot cross a curve outside that collar, so the nesting order is preserved.
\end{proof}

This lemma settles the local factor \(2\) in the state count.  It does
\emph{not} assert that the nooses attached independently to all edges of an
arbitrary sphere-cut decomposition have already been chosen as one laminar
family, nor that every proper tree segment contains a retained vertex.
Those two global compatibility clauses remain explicit fields of the
sphere-cut supply consumed by Theorem~\ref{thm:typed-sphere-cut-descent}.

\begin{theorem}[physical deletion between equal nested noose states]
\label{thm:physical-noose-deletion}
\Sprose
Let \(G\) be a connected bridgeless cubic spherical map.  Let
\(\gamma_i,\gamma_j\), with \(i<j\), be disjoint nested simple nooses which
avoid vertices and each cross exactly \(k\ge2\) edges.  Suppose the two
vertex shores of each cut induce connected graphs, the two shores retained
for the splice are saturated, their ordered geometric seam types identify
the retained-side boundary first-return orders, and the annular slab between
them contains a graph vertex.  Delete that slab and glue the paired boundary
half-edges, using the ambient edge flip on the outer cut to pass to its
complementary shore.  Then the result \(G'\) is a strictly smaller connected
bridgeless cubic spherical map.

If, in addition, the two cuts have the same exact cumulative Count-support
state and the same finite geometric seam type, then \(G'\) has an accepting
Tait colouring exactly when \(G\) does.  Hence zero accepting Count is
preserved.
\end{theorem}

\begin{proof}
Retain the vertex shore inside the inner cut and the complementary vertex
shore outside the outer cut.  For each cut, the retained shore and its
complement are connected; Corollary~\ref{cor:planar-bond-canonical-cap}
therefore supplies a spherical, face-simple canonical closure with one hub.
The equality of ordered seam states transports the retained-shore return at
the inner cut to the retained-shore return at the outer cut.
Theorem~\ref{thm:planar-bond-orders-opposite} says that ambient edge flip at
the outer cut transports that order to the inverse order on the complementary
shore.  Their composite is therefore the orientation-reversing seam matching.
Hence
Corollary~\ref{cor:hub-seam-sphericity} proves, by face-orbit counting rather
than by an implicit drawing, that the glued rotation system is spherical.
Every retained cubic
vertex loses exactly the cut incidences recorded at its boundary and gains
the corresponding seam incidences, so its degree remains three.  The two
retained sides are connected.  By
Lemma~\ref{lem:saturated-side-boundary-essential} they are
boundary-essential, and Lemma~\ref{lem:boundary-essential-gluing} makes the
composite bridgeless.  The deleted slab contains a vertex and none is added,
so \(G'\) is strictly smaller.

The geometric seam type records the ordered endpoint identifications and any
local equality data needed by the chosen graph representation.  The literal
composite is loopless.  A simple-graph implementation may reject a duplicate
edge in its seam type, or the loopless-multigraph output may instead be
normalized by repeated digon suppression as explained below.  Equality of
cumulative support states is exactly the hypothesis of
Theorem~\ref{thm:determinized-support-deletion}.
Corollary~\ref{cor:physical-closed-count-replacement} identifies that support
replacement with the literal matched-seam rotation system used in the
preceding structural argument.  Thus the unchanged exterior and the
replacement interior have the same accepting support in \(G'\) and \(G\),
proving the last assertion without a second, merely semantic composite.
\end{proof}

The complete replacement kernel of this argument is checked on the literal
majority-shore data in
Theorem~\ref{thm:majority-shore-physical-replacement}.  The present theorem
remains marked \textsc{prose} only because it is the optional continuous
Jordan-curve presentation.  The load-bearing discrete version is now fully
checked in Corollary~\ref{cor:dual-noose-literal-replacement}: it starts with
an actual simple closed walk in the facial dual and a chosen deletion
component, derives the exact planar bond and connected complementary shores,
computes the canonical phased state, performs the literal splice, and
terminates parallel-seam normalization.  Thus the remaining continuous
noose-to-facial-dual conversion belongs to the bounded-width M1 supplier; no
colouring, cap, Euler, normalization or minimality step remains hidden in M2.

\begin{theorem}[the implemented seam composite has the non-topological map fields]
\label{thm:implemented-seam-composite-fields}
Let two finite open rotation tangles be glued along a nonempty matched seam.
If both side multigraphs are connected, the composed rotation system is
primal-connected.
\Sverified{Mettapedia.GraphTheory.FourColor.GoertzelV24CompositeSeamMultigraph.primalConnected_composeRotationSystem}
If, in addition, both sides are boundary-essential and the seam has at least
two ports, it is edge-bridge-free.
\Sverified{Mettapedia.GraphTheory.FourColor.GoertzelV24CompositeSeamMultigraph.edgeBridgeFree_composeRotationSystem}
If each open side has exactly three incident darts at every retained vertex,
counting its boundary half-edges, then the composite is cubic.
\Sverified{Mettapedia.GraphTheory.FourColor.GoertzelV24CompositeSeamCubic.isCubic_composeRotationSystem}
If the darts at each side vertex form one rotation cycle, the same is true in
the composite.
\Sverified{Mettapedia.GraphTheory.FourColor.GoertzelV24CompositeSeamCubic.vertexRotationCyclic_composeRotationSystem}
\end{theorem}

\begin{lemma}[literal vertex shores inherit the local map structure]
\label{lem:vertex-shore-local-structure}
Let \(M=(D,\alpha,\rho)\) be a finite cubic rotation system whose darts at
each vertex form one \(\rho\)-cycle, and let \(S\) be any set of vertices.
Cut every edge with exactly one endpoint in \(S\), retaining its dart based
in \(S\) as an unpaired boundary dart.  The resulting literal open tangle is
cubic when its boundary half-edges are counted.
\Sverified{Mettapedia.GraphTheory.FourColor.GoertzelV24VertexSideOpenTangle.openIsCubic_ofVertexSide}
Its darts at each retained vertex still form one rotation cycle.
\Sverified{Mettapedia.GraphTheory.FourColor.GoertzelV24VertexSideOpenTangle.openRotationCyclic_ofVertexSide}
If the graph induced by \(S\) is connected, then the interior-edge
multigraph of the open tangle is connected.
\Sverified{Mettapedia.GraphTheory.FourColor.GoertzelV24VertexSideOpenTangle.sideMultigraph_connected_of_induce_connected}
Moreover its full boundary-dart carrier is saturated: every ambient edge
leaving \(S\) contributes its unique retained orientation as a boundary
port.
\Sverified{Mettapedia.GraphTheory.FourColor.GoertzelV24VertexSideOpenTangle.boundaryDarts_saturate_toMultigraph}
If the ambient edge multigraph is bridgeless, then this literal interior-edge
multigraph is boundary-essential with respect to those boundary darts.
\Sverified{Mettapedia.GraphTheory.FourColor.GoertzelV24VertexSideOpenTangle.boundaryEssential_ofVertexSide}
\end{lemma}

\begin{proof}
The retained-dart partition
\[
 \{d\in D: \operatorname{vert}(d)\in S\}
   \;\cong\; D_{\rm int}(S)\sqcup D_{\partial}(S)
\]
only records whether the opposite endpoint is also in \(S\); it neither
adds nor removes a dart based at a retained vertex.  Hence the open darts
based at \(v\in S\) are in bijection with the ambient darts based at \(v\),
and there are exactly three of them.

The ambient rotation preserves the base vertex, so it restricts to the
retained darts.  The open rotation is exactly this restriction conjugated
by the displayed partition equivalence.  Two open darts based at the same
retained vertex therefore correspond to two ambient darts based at the same
vertex; ambient rotation-cyclicity puts them in one \(\rho\)-orbit, and
restriction followed by conjugation puts the original open darts in one
open-rotation orbit.  Finally, every adjacency of the graph induced by \(S\)
is represented by an ambient dart whose two endpoints lie in \(S\), hence by
one internal dart and one interior edge of the open tangle.  An induced path
therefore maps edge by edge to a path in its interior-edge multigraph.  If an
ambient edge has exactly one endpoint in \(S\), its dart based at that endpoint
is, by definition, a boundary dart; conversely every boundary dart arises in
this way.  This is precisely saturation.  Finally, identify an internal
ambient edge with the involution-orbit of either of its two internal darts.
This is a bijection preserving the unordered endpoint pair.  The preceding
saturated-side lemma makes the induced ambient side boundary-essential, and
transport across this multigraph isomorphism gives boundary-essentiality on
the literal open-tangle carrier used by the splice.
\end{proof}

\begin{lemma}[complementary vertex shores reassemble the ambient Tait problem]
\label{lem:complementary-shore-tait-reassembly}
\Sverified{Mettapedia.GraphTheory.FourColor.GoertzelV24VertexSideReassembly.rotationSystemTaitColorable_complementaryComposite_iff}
Let \(M=(D,\alpha,\rho)\) be a finite loopless rotation system and let
\(S\subseteq V(M)\).  Form the two literal vertex-side open tangles on \(S\)
and \(V(M)\setminus S\), and match their boundary darts by the ambient edge
flip \(d\mapsto\alpha d\).  Their closed seam composite is Tait-colourable if
and only if \(M\) is Tait-colourable.
\end{lemma}

\begin{proof}
Partition each shore's retained darts into internal and boundary darts.  The
composite dart carrier is
\[
  (D_{\rm int}(S)\sqcup D_{\rm int}(S^c))
    \sqcup(D_{\partial}(S)\sqcup D_{\partial}(S^c)).
\]
Undoing these two partitions and then the predicate partition \(S\sqcup S^c\)
gives a canonical bijection from this carrier to \(D\).  On an internal dart
the composite edge flip is the restricted old flip.  On a boundary dart it
is the seam matching, which was defined to be the old flip.  Hence the
bijection conjugates the two edge involutions.  The analogous predicate
partition of the vertex carrier identifies the composite base vertex of a
dart with its old base vertex.

Read any Tait edge-colouring on either rotation system at its darts and
transport that dart-colouring across the displayed bijection.  Conjugacy of
the flips makes it constant on the opposite pair of every edge.  Preservation
of base vertices makes it locally proper, and nonzeroness is unchanged.
The generic descent from an invariant locally proper dart-colouring to edge
orbits therefore gives a Tait edge-colouring on the other rotation system.
Applying the construction in both directions proves the equivalence.  Notice
that this argument uses no face or embedding claim: those are supplied
separately by the planar-bond structural theorem.
\end{proof}

\begin{corollary}[two literal shores feed the structural splice]
\label{cor:literal-shores-structural-splice}
\Sverified{Mettapedia.GraphTheory.FourColor.GoertzelV24VertexSideOpenTangle.bridgelessSphericalCubicMapData_ofVertexSides}
Let \(S_L,S_R\) be connected induced vertex shores of one bridgeless cubic
rotation system with cyclic vertex rotations.  Give their complete cut-dart
interfaces a matching with at least two ports.  If their hub closures are
spherical and face-simple and the matching reverses the two hub orders, then
the literal seam composite is a connected, bridge-free, cubic spherical map
with cyclic vertex rotations.
\end{corollary}

\begin{proof}
Lemma~\ref{lem:vertex-shore-local-structure} supplies connectedness,
boundary-essentiality, open cubicity, and open rotation-cyclicity on the
exact two side carriers.  The hub hypotheses supply the face-count input.
Applying the implemented structural seam theorem therefore gives all four
fields on the literal composite rotation system.
\end{proof}

\begin{corollary}[the implemented seam composite packages the full structural class]
\label{cor:implemented-seam-composite-structure}
\Sverified{Mettapedia.GraphTheory.FourColor.GoertzelV24CompositeStructuralData.bridgelessSphericalCubicMapData_composeRotationSystem}
Let the two finite open rotation tangles satisfy the connectedness,
boundary-essentiality, open-cubicity, and vertex-rotation hypotheses of the
preceding theorem.  Suppose the seam contains two distinct ports.  If the
two hub closures are spherical and face-simple and their hub orders obey
the orientation-reversing equation~(25), then the literal composed rotation
system is a connected, bridge-free, cubic spherical map with one cyclic
rotation at every vertex.
\end{corollary}

\begin{proof}
Theorem~\ref{thm:implemented-seam-composite-fields} supplies connectedness,
bridge-freeness, cubicity, and cyclic vertex rotations.  Corollary~\ref{cor:hub-seam-sphericity}
supplies the spherical Euler equality on the same composite carrier.  These
are exactly the four fields of the structural record, so no transport to a
second representation is required.
\end{proof}

The implementation obtains the first two conclusions by an explicit edge
equivalence between the multigraph glue and
\texttt{composeRotationSystem}; the latter two follow by transporting the
side rotations through the composite carrier.  Thus connectedness,
bridgelessness, cubicity, rotation-cyclicity, and strict vertex accounting
are no longer D3 gaps.  Theorem~\ref{thm:hub-seam-face-formula} and
Corollary~\ref{cor:hub-seam-sphericity} also close the face-orbit calculation
once face-simple spherical caps and the orientation equation are supplied.
The bare matching type cannot manufacture those witnesses.  On literal
graph-backed shores they now come from connected planar-bond data by
Corollary~\ref{cor:planar-bond-canonical-cap};
Theorem~\ref{thm:planar-bond-orders-opposite} derives the opposite
first-return equation from the natural ambient edge-flip matching.

\begin{lemma}[the colouring half of digon suppression]
\label{lem:digon-suppression-colour}
\Sverified{Mettapedia.GraphTheory.FourColor.GoertzelV24DigonSuppressionColor.RotationSystem.TwoEdgeCutPairData.not_complementCap_taitColorable_of_not_ambient}
Let a cubic rotation system be separated by an exact non-colliding two-edge
cut, and suppose that the chosen side contains exactly two vertices.  If the
ambient rotation system is not Tait-colourable, then the complementary cap
is not Tait-colourable.
\end{lemma}

\begin{proof}
The cap of the two-vertex side is cubic and still has two vertices.  Counting
darts in two ways gives
\[
  2|E|=|D|=3|V|=6,
\]
so it has exactly three edges.  Choose a bijection from these edges to the
three nonzero elements of $\mathbb F_2^2$.  Distinct adjacent edges then have
distinct colours, and every edge has nonzero colour, so this cap is
Tait-colourable.

If the complementary cap also had a Tait colouring, globally permute its
three nonzero colours so that the two cap edges receive the same colour.
The exact two-edge-cut gluing theorem then joins the two cap colourings to a
Tait colouring of the ambient rotation system, a contradiction.  Thus this
is the formal version of ``colour the replacement edge, copy its colour to
the two exterior edges, and use the other two colours on the digon.''
\end{proof}

\begin{corollary}[parallel seams can be normalized after gluing]
\label{cor:parallel-seam-normalization}
\Sverified{Mettapedia.GraphTheory.FourColor.GoertzelV24IteratedDigonNormalization.exists_endpointSimpleCounterexampleBelow_of_areParallelEdges}
Let \(M\) be a finite connected bridgeless cubic spherical loopless
multigraph which is not Tait-colourable, and suppose that \(M\) contains a
parallel pair.  Repeatedly suppress a parallel pair by the digon splice of
Lemma~\ref{lem:digon-splice}.  The process terminates in a strictly smaller
connected bridgeless cubic spherical rotation map with injective endpoint
map which is still not Tait-colourable.  Equivalently, the terminal map has
the canonical simple-graph presentation used in the graph-backed layer.
\end{corollary}

\begin{proof}
If three parallel edges join the same two vertices, connectedness and
cubicity make \(M\) the theta graph, which is Tait-colourable.  Hence at every
non-simple, non-theta stage some repeated endpoint pair has multiplicity
exactly two.
Lemma~\ref{lem:digon-splice} deletes its two cubic endpoints, preserves
connectedness, bridgelessness, cubicity and sphericity, and introduces no
loop.  The two exterior edges are an exact two-edge cut whose digon side has
two vertices, and its complementary cap is precisely the suppressed map.
Lemma~\ref{lem:digon-suppression-colour} therefore proves, without an
additional colour calculation, that uncolourability is preserved.  Each
suppression removes two vertices, so
well-founded induction on the vertex count terminates.  The terminal graph
has no loop and no parallel pair, hence has injective endpoint map, and the
first suppression already made it strictly smaller than \(M\).  The canonical
simple-graph backing then preserves its structural class and Tait semantics.
\end{proof}

Lean carries out this well-founded normalization uniformly, returning an
endpoint-injective rotation system in the same cap-stable class with strictly
fewer vertices.  Endpoint injectivity is exactly the condition consumed by
the canonical simple-graph backing.  The checked minimal form of the same
argument is also available:
\Sverified{Mettapedia.GraphTheory.FourColor.GoertzelV24DigonSpliceNormalization.no_wellFormed_digonPatchData_of_minimal}
a vertex-minimal zero-Tait-count member of the bridgeless spherical cubic
class, with two-sided face orbits, admits no well-formed digon patch.  The
arbitrary-pair extraction and the terminating iteration are therefore no
longer representational gaps.

Therefore a duplicate seam edge need not enlarge the finite receipt.  The
replacement is performed in the loopless rotation-multigraph category, which
is exactly the class quantified over by vertex minimality.  In that class no
post-splice simplification is needed: the sewn map itself is the strictly
smaller zero-Count competitor.  If a later presentation insists on a simple
graph, Corollary~\ref{cor:parallel-seam-normalization} removes a parallel seam,
but this is an optional representation change rather than a premise of the
minimality contradiction.

Thus M2 is closed in the literal-shore representation.  The full finite state
is the canonical hub order together with the set of genuine Tait boundary
words, and Theorem~\ref{thm:majority-shore-physical-replacement} consumes its
equality directly.  What remains belongs to M1: a connected corridor or
branch-decomposition supplier must construct the complete rooted edge-leaf
decomposition and its nodewise connected cuts and middle bounds.  The two
nested shore trees, edge/vertex accounting, roots, literal widths, two
boundary ports, exact states, and strict cubic star between repeated states
are now derived.  No ambient graph data are reread after the state has been
formed.
Moreover every constructive equal-state splice has a strictly smaller
endpoint-simple zero-Count descendant:
\Sverified{Mettapedia.GraphTheory.FourColor.GoertzelV24IteratedDigonNormalization.exists_strict_endpointSimple_replacement_of_normalizedState_eq}.
The cap face count, cap sphericity,
one-hub property, orientation reversal, connectedness, bridge-freeness,
cubicity, and cyclic vertex rotations are no longer open.

\begin{corollary}[phased literal-shore replacement]
\label{cor:phased-literal-shore-replacement}
\Sverified{Mettapedia.GraphTheory.FourColor.GoertzelV24LiteralShoreReplacement.exists_endpointSimple_replacement_of_cardPhasedState_eq}
Let \(A\subsetneq B\) be two certified literal connected shores in a finite
bridgeless spherical cubic map of zero Tait Count.  Suppose their middle sets
have size at most \(w\), their literal boundary widths have size at most
\(k\), and their phased canonical states agree:
\[
 \bigl(|B|\bmod(6w+1),\,q(B)\bigr)
   =\bigl(|A|\bmod(6w+1),\,q(A)\bigr).
\]
Then there is a strictly smaller endpoint-simple bridgeless spherical cubic
map of zero Tait Count.
\end{corollary}

\begin{proof}
Proper containment and equality of the cardinality phases imply
\(|B\setminus A|\ge 6w+1\).  The cubic spacing lemma therefore finds a
vertex whose full three-edge star lies in \(B\setminus A\); this is the
strict material deleted by the splice.  Equality of the second coordinates
of the phased states unpacks to one common literal width and equality of the
ordinary normalized states.  The constructive majority-shore replacement
then gives a strictly smaller zero-Count rotation map in the complete
bridgeless spherical cubic class.  Finally the terminating digon
normalization removes any parallel seam while preserving that class,
zero Count, and strict decrease.
\end{proof}

\begin{corollary}[nested facial-dual nooses physically shorten the map]
\label{cor:dual-noose-literal-replacement}
\Sverified{Mettapedia.GraphTheory.FourColor.GoertzelV24DualNooseLiteralReplacement.exists_endpointSimple_replacement_of_nested_dualNooses}
Let \(M\) be a graph-backed vertex-minimal zero-Tait-count member of the
bridgeless spherical cubic class.  For \(r\in\{0,1\}\), let \(W_r\) be a
simple closed walk of length at most \(w\) in the actual facial dual of
\(M\), let \(X_r\) be the primal edges crossed by \(W_r\), and choose a
connected component \(C_r\) of \(M-X_r\).  Put

\[
 S_r:=\{e\in E(M): e\text{ has an endpoint in }C_r\}.
\]

Assume that \(E(M)\setminus S_r\) contains at least two edges for both
values of \(r\), that \(S_1\subsetneq S_0\), and that the two computed
phased canonical states agree.  Then there is a strictly smaller
endpoint-simple bridgeless spherical cubic map of zero Tait Count.
\end{corollary}

\begin{proof}
A simple facial-dual cycle crosses pairwise distinct primal edges.  Facial
parity saturation proves that \(X_r\) is exactly the graph boundary of
\(C_r\), not merely a set containing that boundary, and proves that both
\(C_r\) and its vertex complement induce connected graphs.  Consequently
\(S_r\) and \(E(M)\setminus S_r\) are connected edge shores.
\Sverified{Mettapedia.GraphTheory.FourColor.GoertzelV24DualNooseLiteralReplacement.DualNooseSide.exactBoundary_connectedSides}

It remains to compare the geometric length with the middle-set width used by
the finite state.  A middle vertex of \(S_r\) cannot lie in \(C_r\): every
edge incident with a vertex of \(C_r\) belongs to \(S_r\).  Choose one
\(S_r\)-edge incident with the middle vertex.  Its other endpoint lies in
\(C_r\), so orienting it outward gives a crossing dart of the noose.  Distinct
middle vertices give distinct darts because they are the distinct terminal
endpoints.  Hence

\[
 |\operatorname{mid}(S_r)|\le |X_r|=|W_r|\le w.
\]
\Sverified{Mettapedia.GraphTheory.FourColor.GoertzelV24DualNooseLiteralReplacement.DualNooseSide.middle_card_le_walk_length}

The chosen component is nonempty and cubic, so \(S_r\) contains at least
three edges.  Together with the assumed two exterior edges, connectedness
therefore supplies both nonempty majority sides.  Thus each noose constructs
the complete literal connected-shore node of width \(w\), including roots
and two distinct boundary darts.  The hypotheses \(S_1\subsetneq S_0\) and
equality of their computed phased states now meet
Corollary~\ref{cor:phased-literal-shore-replacement}, which performs the
physical splice and terminating digon normalization.
\end{proof}

\begin{corollary}[nested noose replacement returns graph-backed normal form]
\label{cor:dual-noose-graph-backed-replacement}
\Sverified{Mettapedia.GraphTheory.FourColor.GoertzelV24GraphBackedNooseReplacement.graphBackedCounterexampleBelow_of_endpointSimpleCounterexampleBelow}
\Sverified{Mettapedia.GraphTheory.FourColor.GoertzelV24GraphBackedNooseReplacement.exists_graphBackedCounterexampleBelow_of_nested_dualNooses}
Under the hypotheses of
Corollary~\ref{cor:dual-noose-literal-replacement}, there is a finite simple
graph \(H\), equipped with graph-backed rotation data, such that its rotation
map is connected, bridgeless, spherical and cubic with cyclic vertex
rotations, has zero Tait Count, and satisfies
\[
                         |V(H)|<|V(M)|.
\]
Thus physical noose replacement lands in the same graph-backed normal class
consumed by the decomposition supplier; graph backing is not a subsequent
unproved representation step.
\end{corollary}

\begin{proof}
Corollary~\ref{cor:dual-noose-literal-replacement} returns an endpoint-simple
rotation system \(R\) with the complete bridgeless spherical cubic structure,
zero Tait Count, and strictly fewer vertices.  Form the primal simple graph
whose edges are the endpoint pairs of the edges of \(R\).  Endpoint
injectivity makes the literal edge carrier bijective with this simple-graph
edge carrier and makes the literal darts bijective with the two oriented
incidences of those edges.  Transport the rotation through this dart
bijection.  It conjugates edge reversal, vertex rotation and facial
successor, so connectedness, bridge-freeness, cubicity, cyclic rotations and
the spherical Euler equation all pass to the graph-backed map.  Conversely,
any Tait colouring of the graph-backed map pulls back along the edge
bijection to one of \(R\), so zero Count passes as well.  The vertex carrier
is unchanged by this presentation, and hence the strict inequality remains
the one already proved for \(R\).
\end{proof}

\begin{theorem}[direct Count-reductive assembly]
\label{thm:direct-count-reductive-assembly}
\Sconditional{ladder transport; target-forced literal connected-shore state supply; route-native finite base}
Assume that every sufficiently large least Tait counterexample can be put,
by the source's target-preserving local reductions, into a composable corridor
whose cut roots satisfy Proposition~\ref{prop:target-forced-cut-chain}.
Use the raw finite boundary-word semantics of
Theorem~\ref{thm:raw-noose-count-pumping}; an optimized closure from
Proposition~\ref{prop:exact-closure-gate} may replace its conservative bound
but is not a logical premise.  Work in the loopless plane-multigraph category
through the splice and apply
Corollary~\ref{cor:parallel-seam-normalization} if a parallel seam appears.
Finally,
verify by the same route-native semantics that no
counterexample exists at or below the resulting effective threshold.  Then
every bridgeless cubic plane graph is Tait colourable, and hence every planar
graph in the finite combinatorial-map convention is four-colourable.
\end{theorem}

\begin{proof}
Suppose otherwise and choose a least counterexample.  The finite base puts it
above the threshold.  Apply the target-preserving ladder reductions and the
global cut-chain theorem.  Exact Count factorization and
Definition~\ref{def:count-support-action} identify the successive cumulative
supports with a path in the finite powerset system.  The cut chain is longer
than \(2^{3^k}\) (or than the separately verified optimized reachable-state
count), so two cumulative support states repeat.
Theorem~\ref{thm:determinized-support-deletion} preserves the
empty accepting support.  The supplied literal cut data meet the hypotheses
of Theorem~\ref{thm:majority-shore-physical-replacement}, which realizes the
deletion as a strictly smaller noncolourable member of the same bridgeless
cubic spherical rotation-system class.  This
contradicts minimality.  The Tait bridge gives the Four-Colour Theorem.
\end{proof}

\begin{theorem}[route-native finite-base reflection]
\label{thm:route-native-finite-base-reflection}
\Sverified{Mettapedia.GraphTheory.FourColor.GoertzelV24SphericalRouteNativeBaseReflection.TaitBaseReflection.taitBaseVerifiedAt_of_filter_eq_empty}
Fix a vertex bound \(B\), a finite state set \(Q\), a predicate
\(\operatorname{Necessary}\subseteq Q\), and a computed reachable set
\(R\subseteq Q\).  Suppose every graph-backed bridgeless spherical cubic map
\(M\) with \(|V(M)|\leq B\) has a trace \(\tau(M)\in Q\) such that:
\begin{enumerate}
\item \(\tau(M)\in R\); and
\item if \(M\) is not Tait-colourable, then
      \(\operatorname{Necessary}(\tau(M))\).
\end{enumerate}
If
\[
             \{q\in R:\operatorname{Necessary}(q)\}=\varnothing,
\]
then every such \(M\) is Tait-colourable.  Consequently the exact
finite-base premise at bound \(B\) used by the spherical reductive assembly
holds.
\end{theorem}

\begin{proof}
If a bounded map \(M\) were not Tait-colourable, adequacy would make
\(\operatorname{Necessary}(\tau(M))\) true, while coverage would put
\(\tau(M)\) in \(R\).  Thus \(\tau(M)\) would belong to the displayed
filtered set, contradicting its emptiness.  This is exactly the proposition
\(\operatorname{TaitBaseVerifiedAt}(B)\) consumed by the checked spherical
assembly.
\end{proof}

This theorem fixes the certificate boundary without pretending to have run
the certificate.  The trace must come from the compositional profile system,
its adequacy theorem must be proved from the actual bounded map, and the
reachable set must cover all such traces.  It may safely overapproximate the
realizable states.  The final obligation is then one explicit finite
emptiness check, suitable for the flat-payload verified-reflection discipline;
it is not a classical catalogue of unavoidable configurations.

\begin{theorem}[checked spherical reductive assembly]
\label{thm:spherical-reductive-assembly}
\Sverified{Mettapedia.GraphTheory.FourColor.GoertzelV24SphericalReductiveAssembly.combinatorialFourColorStatement_of_supply_and_base}
Fix interface bounds $k,w$, and put
\[
 B(k,w)=4\left(2^{(6w+1)
   \sum_{0\le j\le k}j!\,2^{3^j}}-1\right)+6.
\]
Suppose that every graph-backed vertex-minimal Tait counterexample carries
the complete rooted connected edge-leaf decomposition of bounds $k,w$
required by Theorem~\ref{thm:rooted-connected-branch-decomposition}.
Suppose also that every graph-backed bridgeless spherical cubic map with at
most $B(k,w)$ vertices is Tait colourable.  Then every finite combinatorially
planar simple graph is vertex-four-colourable, component by component.
\end{theorem}

\begin{proof}
If a Tait counterexample exists, choose one with the least number of
vertices.  Bridge-freeness and the spherical Euler equation make every edge
two-sided; minimality excludes parallel edges, so this least rotation map has
a canonical simple-graph backing.  The supplied connected decomposition and
Theorem~\ref{thm:rooted-connected-branch-decomposition} give

\[
                         |V|\le B(k,w).
\]

The bounded-base hypothesis therefore Tait-colours the same map, a
contradiction.  Thus no graph-backed vertex-minimal Tait counterexample
exists.  The checked minimal-counterexample selection and stellar Tait
reduction then give the connected spherical Four-Colour statement, and the
componentwise wrapper gives the displayed conclusion.  Lean checks this
entire implication; its only open inputs are the decomposition-producing
function and the bounded-base proposition stated above.
\end{proof}

\begin{corollary}[factored bounded-width assembly]
\label{cor:factored-bounded-width-assembly}
\Sverified{Mettapedia.GraphTheory.FourColor.GoertzelV24SphericalReductiveAssembly.combinatorialFourColorStatement_of_raw_and_base}
Fix $w$ and put $B=B(w,w)$.  Suppose every graph-backed vertex-minimal
Tait counterexample has a complete rooted branch decomposition of width at
most $w$, and suppose every graph-backed bridgeless spherical cubic map with
at most $B$ vertices is Tait colourable.  Then every finite combinatorially
planar simple graph is vertex-four-colourable.
\end{corollary}

\begin{proof}
A least counterexample is connected and bridge-free, hence
$2$-edge-connected.  Apply
Theorem~\ref{thm:width-preserving-connectedization} to the raw width-$w$
tree, then apply the checked rooted adapter to obtain the exact decomposition
supply of Theorem~\ref{thm:spherical-reductive-assembly} with $k=w$.  That
theorem and the bounded base give the conclusion.  Lean checks this
composition; the only remaining mathematical inputs at this boundary are
the raw width bound and the route-native finite base.
\end{proof}

\begin{corollary}[factored assembly from a route-native base audit]
\label{cor:factored-bounded-width-reflection-assembly}
\Sverified{Mettapedia.GraphTheory.FourColor.GoertzelV24SphericalRouteNativeBaseReflection.TaitBaseReflection.combinatorialFourColorStatement_of_raw_and_filter_audit}
Fix \(w\), put \(B=B(w,w)\), and suppose every graph-backed
vertex-minimal Tait counterexample has a complete rooted branch decomposition
of width at most \(w\).  Suppose also that a finite-state reflection at
bound \(B\) satisfies the coverage and adequacy hypotheses of
Theorem~\ref{thm:route-native-finite-base-reflection}, and that its filtered
reachable set is empty.  Then every finite combinatorially planar simple
graph is vertex-four-colourable.
\end{corollary}

\begin{proof}
Theorem~\ref{thm:route-native-finite-base-reflection} converts the finite
audit into the exact bounded-base proposition.  The checked
width-preserving connectedization converts the raw decomposition into the
connected decomposition consumed by
Theorem~\ref{thm:spherical-reductive-assembly}.  The Lean declaration cited
above checks this composition without any additional premise.
\end{proof}

\begin{corollary}[wall-exclusion assembly]
\label{cor:wall-exclusion-assembly}
\Sconditional{abstract-planar representation theorem; fixed-wall exclusion for least counterexamples; bounded-width decomposition; route-native finite base}
Assume the conventional representation theorem from finite abstract planar
graphs to spherical rotation presentations.  Assume there is a fixed wall
excluded from every least Tait counterexample, and that wall-free members
of the normal-form class carry a complete rooted branch decomposition of
width at most $w$.  Assume that the resulting route-native finite base
contains no counterexample.  Then every planar graph is four-colourable.
\end{corollary}

\begin{proof}
Proposition~\ref{prop:target-forced-wall} bounds the size of a least Tait
counterexample and constructs the raw width supply consumed by
Corollary~\ref{cor:factored-bounded-width-assembly}.
Theorem~\ref{thm:width-preserving-connectedization} converts it to the
literal decomposition supply, and the finite-base hypothesis is the
remaining premise.  That corollary gives the spherical Four-Colour statement;
the separate classical representation theorem from an abstract planar graph
to a spherical rotation presentation gives the customary abstract-planar
formulation.
\end{proof}

The word ``bridgeless'' in the splice conclusion is load-bearing.  A cubic
plane graph with a bridge has zero Tait count for a trivial parity reason, so
a splice producing a bridge would not contradict minimality inside the class
to which Tait's equivalence applies.  Planarity, cubicity, connectedness,
strict size decrease, and zero Count are proved together with bridgelessness
in the ordinary argument of
Theorem~\ref{thm:physical-noose-deletion}; the corresponding discrete
facial-dual construction is checked in
Corollary~\ref{cor:dual-noose-literal-replacement}.  Lean checks the literal
rotation-system composite's connectedness, bridge-freeness, cubicity, and
rotation-cyclicity in
Theorem~\ref{thm:implemented-seam-composite-fields}.  The
orientation-aware face calculation and spherical Euler equality are now
verified in Theorem~\ref{thm:hub-seam-face-formula} and
Corollaries~\ref{cor:hub-seam-sphericity} and
\ref{cor:implemented-seam-composite-structure} package these facts on the
literal composite.  Corollary~\ref{cor:planar-bond-canonical-cap} derives
the cap hypotheses from connected graph-backed shores and the exact
planar-bond rank, without a Jordan/noose premise.  The literal
composite is loopless, and Corollary~\ref{cor:parallel-seam-normalization}
removes any parallel seam if a later simple-graph presentation requires it.
Thus the remaining cut-side work is M1 rather than M2: produce the complete
rooted edge-leaf decomposition and its nodewise connected cuts and middle
bounds from the chosen corridor or connected branch decomposition.  The two
strictly nested trees, vertex accounting, roots, ports, strict material,
ordered first-return data, and exact supports are computed by the checked
consumer; the supplier does not assume in advance that two such fields are
equal.

Theorem~\ref{thm:direct-count-reductive-assembly} is the shortest literal
implementation of the companion note's instruction to open the instances,
count through the finite interface, compose, and pump.  It does not use Trail
Completability or Kauffman factor gluing.  Those remain a source-faithful
alternative assembly and a valuable explanation of the formation calculus,
but their synchronization seam is not on the direct reductive critical path.
No classical unavoidable set, discharging rule, or configuration catalogue
enters the replacement.

\subsection{GWCO: the refuted shortcut and the surviving lemma}

The Good-Word Closed-web Obstruction is false.
\Srefuted{dodecahedral two-cap witness; growing width-two family}
Good-word totally closed webs occur already at ten internal vertices and in a
growing family.  Consequently the old False-valued GWCO base, its numerical
threshold, and any reductive supplier targeting that base are retired.
The dodecahedral opening is the smallest witness but its two cap interfaces
share one boundary edge.  A non-polar \(C_{30}\) opening at cap faces (4)
and (12) is the clean annular witness: the cap cycles are vertex- and
edge-disjoint, and exact enumeration finds six good-word totally closed webs.
This does not contradict the manuscript's polar \(C_{30}\) and \(C_{40}\)
censuses, which are reproduced exactly and contain only bad-word closed webs;
it refutes the universal extrapolation from those laboratories.

The purpose of GWCO was narrower than its statement: it was meant to exclude
a fibre frozen in an adversarial menu-A state.  At the five-outer-stub
boundary used by the closed-web system, the following replacement is valid.

\begin{lemma}[Totally closed at five outer stubs implies menu B]
\label{lem:closed-implies-menu-b}
\Sverified{Mettapedia.GraphTheory.FourColor.GoertzelV24ClosedWebCapComposedMenu.both_restoredMenuBForPair_of_totallyClosed_at_five}
Let an annular frontier tangle have five inner and five outer stubs and a good
inner word.  If a colouring has a totally closed web, then it is menu B for
each majority pair.
\end{lemma}

\begin{proof}
The Cubic Closure Theorem makes every two-colour component inner-touching.
At the five--five boundary, the exact endpoint census then rules out two
distinct inner stubs in one component; equivalently, every such component is
radial, with one inner and one outer end.  Thus, for either majority pair, the
open strand relation on the four active inner positions is diagonal: no
strand joins two distinct positions.  Restoring the pentagonal cap contributes
exactly its two disjoint active blocks.  Since no open strand can join those
blocks, their generated equivalence relation still has at least two classes.
The transverse-chord criterion identifies this condition precisely with
cap-composed menu B.  Lean checks the entire chain on the literal graph-derived
profile relation, including both majority pairs.
\end{proof}

For six or more outer stubs the source's Cubic Closure counting already makes
total closure impossible.  Boundaries with fewer than five outer stubs are
not covered by Lemma~\ref{lem:closed-implies-menu-b}; they must be discharged
by their own finite boundary analysis if they occur in the final assembly.
This qualification is load-bearing and prevents the five-stub argument from
being silently promoted to an all-boundary theorem.

\subsection{The exact conditional conclusion}

\begin{remark}[retired Trail-Completability assembly]
\label{thm:conditional-compositional-route}
\Srefuted{two-defect boundary-flux identity and cubic/trail parity mismatch}
The earlier conditional route theorem passed from a deleted edge to a
coloured uncompletable trail and later applied the closed-cubic Kauffman
parity theorem to defect-path moves.  Both steps fail, independently, as
shown in Section~\ref{sec:kauffman-trail-spine}.  Merely listing Trail
Completability and a relative parity theorem as assumptions would conceal
rather than repair the missing mathematics, so that assembly is withdrawn.
\end{remark}

The surviving conditional conclusion is
Theorem~\ref{thm:direct-count-reductive-assembly}, with the tree-form
alternative isolated in Corollary~\ref{cor:wall-exclusion-assembly}.  The
linear dependency chain is: the Tait/minimal-counterexample bridge;
source-preserving local ladder transport;
Proposition~\ref{prop:target-forced-cut-chain}; the raw noose Count theorem
and physical support-state splice; and a route-native finite base.  The
tree-form chain replaces the global corridor proposition by fixed-wall
exclusion plus a raw bounded-width decomposition.  The checked
width-preserving connectedization theorem supplies the connected tree.  Both chains
use the same checked majority-shore D3 splice and the
same finite base.

At present the decisive open premises are therefore: either the
target-forced global cut chain or fixed-wall exclusion together with the raw
bounded-width decomposition supplier; and the route-native finite base.
Exact Count replacement on supplied normalized majority shores is now
checked.  A compatible tight vertex-noose to edge-noose
perturbation is optional rather than load-bearing.  D3's face count, Euler
arithmetic, and majority-shore structural packaging are no longer open.  The mesh
isoperimetry theorem proves that a large wall does not automatically supply
the linear chain, but it does not exclude the wall.  An executable minimized
Cell closure remains important for making the base small enough to replay,
while the raw finite semantics no longer depends on the refuted closed-web
carrier.  Neither Trail Completability nor Kauffman factor gluing lies on
this direct reductive critical path.

\subsection{Generic handoff integration and strict nested-cut transport}
\label{sec:generic-handoff-integration}

The generic part of the development at lab revision
\texttt{5ed2e5968} has been integrated selectively.  The imported consumers
are exact side-colouring semantics, arbitrary-shape slab and node
constructions, colour-orbit state compression, and a sound and complete
enumerator for a presented tangle.  Configuration catalogues, flower
certificates, and their production sweeps are not part of this integration.
Concrete width-specific zigzag tables remain on the lab branch; the generic
consumers do not require their decision computations.

\begin{proposition}[compressed bound for supplied nested sides]
\label{prop:compressed-supplied-nested-sides}
\Sverified{Mettapedia.GraphTheory.FourColor.TubeSlab.le_of_nested_nodes_orbit}
\Sverified{Mettapedia.GraphTheory.FourColor.TubeSlab.le_of_nested_sides_five}
In a graph-backed vertex-minimal Tait counterexample, a strictly nested
family of literal nodes of exact boundary width $k$ and middle bound $w$
has at most
\[
                 (6w+1)k!\,2^{q_k}
\]
members, where $q_k$ is the number of global colour-permutation orbits of
zero-sum nonzero boundary words.  In particular, $q_5=10$, and a family of
good sides of exact width five with connected internal edge shores and
connected complementary edge shores has at most
\[
                         31\cdot120\cdot2^{10}=3{,}809{,}280
\]
members, provided its internal edge shores are strictly nested.
\end{proposition}

\begin{proof}
The actual support of a cubic vertex side has zero colour sum and is
invariant under the six global permutations of nonzero colours.  Its set
of orbit keys therefore determines the support exactly.  Keep the hub
rotation and shore-size phase, and apply pigeonhole to these data and the
key set.  Equality recovers the original phased state, so the existing
strict physical replacement contradicts minimality.
\end{proof}

This bound counts supplied sides, not vertices of the ambient graph.  It is
universe-polymorphic.  Boundary coordinates are supplied bijections, with
their induced hub rotation retained in the state; cyclic order has not been
discarded.  Neither the good-side condition nor strict shore growth follows
merely from having local patches of small width.

\begin{lemma}[connected exterior forces strict shore growth]
\label{lem:connected-exterior-strict-shore}
\Sverified{Mettapedia.GraphTheory.FourColor.NestedCyclicCutStrictness.incidentEdgeShore_ssubset}
\Sverified{Mettapedia.GraphTheory.FourColor.NestedCyclicCutStrictness.cyclicCut_nodes_strict}
Let $S\subsetneq T\subsetneq V(G)$, and suppose $G[V(G)\setminus S]$
is connected.  Assign to a vertex side every edge with at least one endpoint
there.  Then its incident-edge shores satisfy $I(S)\subsetneq I(T)$.
Consequently the connected-cyclic-cut node constructor preserves strict
vertex-side nesting, without an additional fresh-edge hypothesis.
\end{lemma}

\begin{proof}
Choose $x\in T\setminus S$ and $y\notin T$.  A path from $x$ to $y$
inside the old complement crosses from $T\setminus S$ to $V(G)\setminus T$.
The crossing edge belongs to $I(T)$ but not to $I(S)$.  The existing
incident-shore map is monotone, proving strictness.  For realized cyclic
cuts the outer cycle certifies $y$; their widths, vertex sides and boundary
order are unchanged by this argument.
\end{proof}

The pre-formalization gate checks all labelled simple graphs through five
vertices and three cubic controls, including $594$ nested pairs with cycles
on both sides in the hexagonal prism.  The three-vertex path shows that
dropping connectedness of the old exterior permits equal shores at strictly
nested vertex sides.  These are generic geometric tests, not target
counterexamples.  The implementation and receipts are
\texttt{tools/fourcolor/v24\_nested\_shore\_strictness\_gate.py} and
\texttt{results/fourcolor/v24\_nested\_shore\_strictness\_gate.json}.

\begin{proposition}[presented enumeration has literal tangle semantics]
\label{prop:presented-enumeration-literal-semantics}
\Sverified{Mettapedia.GraphTheory.FourColor.TubeSlab.Presented.Pres.enumMask_testBit_iff}
\Sverified{Mettapedia.GraphTheory.FourColor.TubeSlab.Presented.Pres.accepts_iff}
For a valid finite presentation with $k>0$ ports on each side, bit $Y<3^k$
of its computed output mask at input $X$ is set exactly when there is a
proper nonzero internal colouring extending the two encoded boundary words;
equivalently the constructed two-sided tangle accepts those words.
\end{proposition}

This is an enumeration correctness theorem, not a certificate of closure of
the route's physical alphabet.  Initial states, accepting states, coverage,
and the concrete finite-base audit remain separate obligations.

\paragraph{Trail-refutation scope.}
The lab's \texttt{KauffmanTrailCounterexample.all\_ok} checks its explicitly
defined finite transition certificate: proper full colouring, a list of
$36$ defect states, component partitions, switch closure, and absence of a
component joining the two tips.  Its Boolean statement does not itself
construct a spherical rotation system or identify its transition system
with the manuscript's frozen-boundary trail predicate.  These semantic
qualifications must accompany any use of that certificate.  The trail
terminal assembly above remains retired; no stronger interpretation of this
additional finite certificate is needed by the direct Count assembly.
The latter still explicitly takes raw bounded-width decomposition and the
route-native finite base as premises.  Their absence is not cured by the
fact that its import graph is unchanged.

The fixed obligation-weighted ledger remains $57/100$.  The results here
strengthen already credited generic infrastructure; they do not discharge
D1b's long-chain supply, either D1c transfer, D2b's executable physical
closure, or D4b's base.  In particular, the strictness lemma does not bound
the complete interface of a long geodesic or control its lateral material.

\subsection{Optimizing the complete geodesic interface}
\label{sec:canonical-geodesic-cuts}

A large frontier for one chosen sweep is not a lower bound on the best
interface.  To test whether rerouting can repair the geodesic construction,
allow every off-axis vertex to move and minimize the complete cut.  The
following generic construction removes arbitrary choices from that attempt.

\begin{theorem}[canonical terminal cuts nest]
\label{thm:canonical-terminal-cuts}
\Sverified{Mettapedia.GraphTheory.FourColor.CanonicalTerminalCuts.exists_canonical}
\Sverified{Mettapedia.GraphTheory.FourColor.CanonicalTerminalCuts.canonical_mono}
\Sverified{Mettapedia.GraphTheory.FourColor.CanonicalTerminalCuts.canonical_ssubset}
In a finite multigraph, let $A,B$ be disjoint vertex sets.  Among all sides
$S$ containing $A$ and disjoint from $B$, choose one minimizing
$|\delta(S)|$, then $|S|$.  Such a canonical side exists and is unique.
If $A\subseteq A'$ and $B'\subseteq B$, the canonical sides satisfy
$S\subseteq S'$.  If also $B\cap A'\ne\varnothing$, then
$S\subsetneq S'$.
\end{theorem}

\begin{proof}
Finite minimization of $(|V|+1)|\delta(S)|+|S|$ gives existence with the
stated tie-break.  The edge-cut cardinality is submodular, as checked one
edge at a time.  For canonical sides $S,T$ at the two terminal pairs,
$S\cap T$ is feasible for the earlier pair and $S\cup T$ for the later
one.  Minimality and submodularity force equality of the respective cut
cardinalities.  The first side's cardinality tie-break then gives
$S\cap T=S$.  Applying nesting in both directions proves uniqueness.
A vertex in $B\cap A'$ belongs to the later side and not the earlier one,
proving strictness.
\end{proof}

This is a finite construction, not an assumed interface supply.  Lean's
selection uses classical choice from the proved finite minimum; the lab's
flow algorithm is a separate executable implementation, not yet identified
with that selection by a Lean theorem.  Neither the construction nor its
nesting theorem supplies a uniform width bound, connectedness, or boundary
order.

\begin{lemma}[path packing lower-bounds every feasible cut]
\label{lem:terminal-path-packing-bound}
\Sverified{Mettapedia.GraphTheory.FourColor.CanonicalTerminalCuts.paths_le_boundary}
If $r$ pairwise edge-disjoint paths each meet $A$ and $B$, every side
containing $A$ and avoiding $B$ has at least $r$ boundary edges.
\end{lemma}

\begin{proof}
Each path crosses the boundary.  Choose a crossed edge on each path;
edge-disjointness makes the choice injective into the complete cut.
\end{proof}

\paragraph{Calibrated laboratory.}
Along a facial-dual diameter geodesic $F_0,\ldots,F_L$, fix all vertices
of $F_0,\ldots,F_i$ inside and all vertices of
$F_{i+2},\ldots,F_L$ outside.  The intervening face is unconstrained.
In the tested cubic maps these terminal sets are disjoint, connected, and
contain face cycles.  A terminal changes sides after two steps, so canonical
cuts give strict progress at that spacing.  The lab checks connectivity of
both resulting sides and traces the complete cut as one simple dual cycle,
recording its edge order.  It also supplies an equally large family of
edge-disjoint terminal paths: the displayed cut is an upper bound, and the
packing is a matching lower bound on \emph{every} feasible cut.

For the source tube controls of lengths $1$ through $12$, all optimum
widths are five.  On Goldberg controls the maximum optimum widths are:
\[
\begin{array}{c|rrrrrr}
\text{frequency}&2&4&6&8&10&12\\ \hline
\text{optimum complete width}&8&14&20&26&32&38\\
\text{distinct nested sides}&5&11&17&23&29&35.
\end{array}
\]
Thus the local-carrier bound $30$ does not even bound the optimally rerouted
interfaces for these terminal requirements.  Unlike the earlier large
measured frontiers, these are lower bounds for the specified edge-separation
problem.  They count cut edges (ports), not middle-set vertices; they are
not lower bounds on optimum branchwidth.  If only the final
cap is required to stay outside while the prefix grows, all tested Goldberg
widths are five, but only two distinct canonical sides remain.  Small width
alone therefore does not give the needed long chain.

The script \texttt{tools/fourcolor/v24\_geodesic\_mincut\_gate.py} and
\texttt{results/fourcolor/v24\_geodesic\_mincut\_gate.json} record these
tests.  The small-case gate also checks every graph through four vertices
and every monotone terminal transition; its flow solver is compared against
exhaustive minimization on connected terminal controls.  Saved worst-cut
certificates can be checked without rerunning the flow solver, and duplicated
path witnesses are rejected.  These numerical graph witnesses are checked
by the executable verifier, not instantiated as kernel theorems about the
Goldberg graphs.  The two generic results above are kernel-checked.

No tested graph is a claimed zero-target counterexample.  A finite list of
growing optima is not itself a theorem of unbounded growth, and these fixed
terminal requirements do not exhaust every possible compositional
decomposition.  What fails here is the proposed repair that lets arbitrary
lateral reassignment turn this ordered geodesic sweep into a width-$30$
chain.  D1b/D1c still need a target-specific width-and-progress argument or
a different source-faithful global construction.  The fixed ledger remains
$57/100$.

\subsection{Dependency-directed scheduling and a genuinely absolute control}
\label{sec:absolute-topology-gate}

The fixed ledger is coarse: it deliberately credits complete obligations,
not each internal construction. Consequently a plateau is compatible with
useful internal work, but is not evidence that the remaining obligations
are nearly proved. Nor are the remaining $43$ points independent jobs.
In the existing direct specialization, fixed-mesh exclusion and the
route-native base at \texttt{meshFreeVertexBound} are sufficient; the
mesh-free raw-decomposition adapter is already constructed in
\texttt{SphericalMeshFreeRawDecomposition}. In particular $D_{2b}$ denotes
optimized executable physical closure, not that adapter or a missing
general grid theorem.

Proposition~\ref{prop:exact-closure-gate} explicitly makes optimized closure
optional when the conservative finite state bound is retained. Closing its
four-point row would still require the physical alphabet, labelled steps,
literal initial and accepting sets, their soundness and completeness, and
the three independently replayed closure counts. It is not just running
an optimizer on the available unlabelled relation. Likewise the final
finite-base audit cannot be credited by an audit at a smaller convenient
threshold. The immediate mathematical priority is an ordered-mesh
reduction using actual target hypotheses, with a genuine decrease or
compatible smaller-piece gluing, not a smaller state count alone.

\paragraph{A stronger falsification control.}
Colourable closed controls have zero disagreement at unrestricted absolute
minima. Positive Kempe-orbit minima test a weaker hypothesis. To test
absolute minimality nonvacuously while deliberately withholding planarity,
the new lab uses the following single cubic graph. Its vertices are
$(i,s)$ for $i\in\mathbb Z/5\mathbb Z$ and $s\in\{0,1,2,3\}$.
Join $(i,0)$ to the other three vertices in its block, and add edges
\[
 (i,1)(i+1,1),\qquad (i,2)(i+1,3),\qquad (i,3)(i+1,2).
\]
Two complete independent enumerators agree that this graph has no Tait
colouring. Direct checks confirm connectedness, absence of bridges, and
girth five. Vertex numbers below mean $4i+s$.

Fix the deleted edge $(0,1)$ and inspect all six target edges whose closed
endpoint neighbourhoods are disjoint from its own. Both deletion graphs
are colourable in each case. Across all $50544$ pairs of deletion
colourings, there are $120$ absolute minimizers. The minimum is three for
four target edges and five for the other two. Every minimizing
disagreement graph is exactly two paths, each meeting both deletion-port
sets; none has a branch or a cycle. Thus even a secondary minimization of
branch count or cycle rank leaves these positive absolute minima intact.

The omitted planar hypothesis is exhibited by a literal subdivision of
$K_{3,3}$, with branch classes $\{8,12,16\}$ and $\{13,14,15\}$.
Its nine paths are
\[
\begin{array}{lll}
8,9,13 & 8,11,14 & 8,10,15\\
12,13 & 12,14 & 12,15\\
16,17,13 & 16,19,14 & 16,18,15.
\end{array}
\]
The checker verifies actual edges, every required branch pair, and
pairwise disjoint path interiors. No planar-embedding oracle is trusted
for this receipt.

This is an independently cross-checked computational control, not a Lean
kernel certificate or a least spherical counterexample. It refutes only
a generic conclusion of zero disagreement from absolute minimality and
two spanning paths without additional hypotheses. It does not refute the
spherical, ordered-mesh, or minimal-counterexample route. Its purpose is
to reject exchange arguments which never use those missing hypotheses.
The reproducible files are
\texttt{tools/fourcolor/v24\_absolute\_topology\_gate.py} and
\texttt{results/fourcolor/v24\_absolute\_topology\_gate.json}.
No configuration catalogue or reducibility certificate is involved. No
ledger row is newly discharged: $57/100$.

\subsection{Least-size selection does not compute the structural threshold}
\label{sec:least-size-bound-audit}

One attempted use of the target's extra hypotheses is to appeal directly to
global vertex minimality for a uniform bound.  This proves an existential
statement, but does not yet give the numerical threshold required by the
reductive route.  The distinction can be checked against the live definitions.

\begin{proposition}[semantic least-size bound]
\Sverified{Mettapedia.GraphTheory.FourColor.LeastSizeBoundAudit.card_eq_semanticThreshold}
\Sverified{Mettapedia.GraphTheory.FourColor.LeastSizeBoundAudit.exists_uniform_minimal_size_bound}
Fix a universe and let $N_*$ be the infimum in $\mathbb N$ of the vertex
counts of non-Tait-colourable bridgeless spherical cubic maps in the
minimality class.  Every graph-backed vertex-minimal counterexample has
exactly $N_*$ vertices.  In particular, there exists a uniform bound on the
sizes of such minimal counterexamples, independently of corridor geometry.
\end{proposition}

\begin{proof}
Its vertex count belongs to the bad-size set.  If that count exceeded the
least element, the corresponding smaller bad map would contradict its
\texttt{smallerColorable} field.  If the set is empty there is no minimal
counterexample to bound.
\end{proof}

\begin{proposition}[the semantic base carries the universal target]
\Sverified{Mettapedia.GraphTheory.FourColor.LeastSizeBoundAudit.base_at_semanticThreshold_iff_universal}
The live predicate \texttt{TaitBaseVerifiedAt} at $N_*$ is equivalent to
\texttt{EveryBridgelessSphericalCubicTaitColorable}.
\end{proposition}

The equivalence itself is not an obstruction: a successful reductive proof
also gives a biconditional with its finite base.  What this particular
construction lacks is an explicit numerical bound derived independently of
the unknown bad-size set.  No value of $N_*$ is computed here, no finite audit
is performed, and no raw decomposition supplier is obtained.  This is not a
claim that such a value is uncomputable in principle, nor a refutation of the
compositional route.  It records why bare least-size selection does not
advance the effective D1b/D1c or D4b obligations.  The diagnostic module is
not added as a premise supplier to the headline; the ledger stays $57/100$.

\subsection{Narrow decomposition does not bound a prescribed whole-face cut}
\label{sec:thin-annulus-forced-cut-audit}

The next attempted repair was to transfer a small decomposition width to the
canonical cuts prescribed by the facial geodesic.  With long faces this
transfer fails even on thin spherical cubic graphs.  This is a diagnostic
of the routing requirements, not a new counterexample to the target theorem.

The gate builds the thickness-two annulus with two cap faces of length $q$,
using six vertices per rung.  Besides the caps it has $2q$ pentagons and
$q$ hexagons.  It checks the complete spherical rotation data, connectedness,
girth at least five, and a cap-to-cap facial-dual geodesic of length four.
Require the first cap's vertices inside and the future geodesic faces after
one free face outside.  The $q$ disjoint rungs join these terminals, so every
feasible cut has at least $q$ edges.  The first cap itself realizes equality;
its entire boundary and its cyclic order are checked.

\[
\begin{array}{c|rrrrr}
q&5&8&16&32&64\\ \hline
|V|&30&48&96&192&384\\
\text{optimum forced cut}&5&8&16&32&64\\
\text{displayed decomposition width}&8&8&8&8&8\\
\text{width-eight perimeter sides}&2&5&13&29&61.
\end{array}
\]

The decomposition certificate lists each edge exactly once, takes the first
edge as root leaf, the second as the other leaf at the root fork, and the
remaining edges as a comb.  Every tree cut is checked: internal shores are
suffixes, whose middle sets equal those of the complementary prefixes;
singleton leaf cuts are checked separately.  This proves an upper bound on
branchwidth, not optimality.  Some raw tree shores are disconnected, and
the receipt records them rather than asserting a connected decomposition.

For comparison, intervals of two through $q-2$ complete rungs give strictly
nested vertex sides and strictly nested internal edge shores.  Both sides
are connected, contain cycles, and have internal vertex degree at least two.
Their full cuts have eight edges, including both ends of the strip and both
cap-face crossings.  The entire simple dual-cycle order is retained.

The executable witnesses and checker are
\texttt{results/fourcolor/v24\_thin\_annulus\_cut\_gate.json} and
\texttt{tools/fourcolor/v24\_thin\_annulus\_cut\_gate.py}.  Independent replay
uses no cut optimizer; corrupted leaf lists, packed paths, and boundary
orders are rejected.  The graph-specific assertions are executable checks,
not newly kernel-instantiated Lean theorems.  The generic path-packing
lower-bound lemma above supplies their mathematical justification.

Thus even a displayed raw width of eight does not put these prescribed
whole-face cuts within the earlier local budget thirty.  No exclusion of the
stack's injective mesh is certified here; no zero Count is asserted.  The
long cap faces also prevent this family from refuting a uniform
\emph{bounded-face} transfer.  The successful perimeter cuts are controls,
not a substitute for supplying comparable cuts in every large target
instance.  D1b/D1c and the $57/100$ ledger are unchanged.

\subsection{Constructed distance components and contour--linkage meshes}
\label{sec:contour-linkage-mesh-extraction}

The next extraction step uses the exact mesh carrier rather than first
building a grid-minor model.  The external constructive branchwidth argument
of Gu and Tamaki, \emph{Improved Bounds on the Planar Branchwidth with
Respect to the Largest Grid Minor Size}, Algorithmica 64 (2012), 416--453,
\href{https://doi.org/10.1007/s00453-012-9627-5}{doi:10.1007/s00453-012-9627-5},
Section~4, uses nested contours and vertex-disjoint crossing paths.  Here
those data can supply the weaker, exact \texttt{Mesh} carrier directly;
this observation does not import their global decomposition theorem.

\begin{proposition}[deep face components and disjoint frontiers]
\Sverified{Mettapedia.GraphTheory.FourColor.FaceDistanceFrontiers.deepRegion_connected}
\Sverified{Mettapedia.GraphTheory.FourColor.FaceDistanceFrontiers.deepRegion_compl_connected}
\Sverified{Mettapedia.GraphTheory.FourColor.FaceDistanceFrontiers.mixed_disjoint}
In a connected graph $H$ on faces, fix a root $r$ and a far face $z$.
Define $R_i$ by walks from $z$ all of whose vertices have distance strictly
greater than $i$ from $r$.  If $d(r,z)>i$, this constructs a nonempty
connected component.  Its complement is also connected: every exterior
face has a distance-decreasing walk to $r$ which stays exterior.
The regions nest with increasing $i$; every cut edge joins distances
$i+1$ and $i$.

Suppose the faces incident at each primal vertex form a clique in $H$.
A vertex incident to both $R_i$ and its complement then has minimum
incident-face distance exactly $i$.  Thus the complete mixed-vertex
frontiers at different depths are disjoint.
\end{proposition}

\begin{proposition}[constructing the exact injective mesh]
\Sverified{Mettapedia.GraphTheory.FourColor.ContourMeshExtraction.ofPaths_isVertexInjective}
\Sverified{Mettapedia.GraphTheory.FourColor.ContourMeshExtraction.exists_injectiveMesh_of_frontiers}
\Sverified{Mettapedia.GraphTheory.FourColor.ContourMeshExtraction.exists_injectiveMesh_of_distanceFrontiers}
Let $a$ row walks and $b$ column paths meet on every row--column pair,
with each family vertex-disjoint.  Choosing one meeting vertex per pair
constructs a \texttt{Mesh} with \texttt{IsVertexInjective}; within-family
edge-disjointness follows from vertex-disjointness.  Closed contour walks
are allowed by the existing carrier.  The distance-frontier result above
discharges row disjointness when the rows lie on distinct such frontiers.

Alternatively, intersections are derived from regions whose complete
outer edge frontiers are covered by the rows: a column meeting the region
and its exterior must meet that row.  Columns already meeting a row,
including at an extreme contour endpoint, are handled directly.
\end{proposition}

The preregistered gate constructs these data in complete cubic spherical
$GP(k,0)$ graphs for $k=1,\ldots,8$, with only pentagonal and hexagonal
faces.  The root is face zero and the far face is the first at maximum
distance.  Boundaries of its deep components are traced in full, and the
middle half of the contours is selected.  A vertex-capacitated flow on
the entire ambient graph constructs the columns; both terminal sets have
unit vertex capacities too.  The resulting row and column counts are
\[
\begin{array}{c|rrrrrrrr}
k&1&2&3&4&5&6&7&8\\ \hline
\text{rows}&3&4&5&6&9&10&11&12\\
\text{columns}&5&10&15&20&20&25&30&35.
\end{array}
\]
At $k=8$ the mesh has $420$ distinct branch vertices in a graph of $1280$
vertices.  Replay checks all edges and faces, both connected dual regions,
strict region nesting, full cyclic boundaries, the incident-face clique,
the disjoint column paths and all branch positions.  A matching-size
vertex separator certifies optimality for these selected terminals, not
optimal branchwidth.  Twenty-four corruptions are rejected without running
a flow optimizer.  The executable receipts are in
\texttt{results/fourcolor/v24\_contour\_mesh\_gate.json}, with verifier
\texttt{tools/fourcolor/v24\_contour\_mesh\_gate.py}.

\paragraph{Remaining extraction boundary.}
The generic metric and mesh lemmas are kernel-checked; the eight concrete
graph receipts are independently executable checks, not Lean instances.
No zero Count is asserted.  To complete the global extraction one must
prove the small-noose versus large-linkage alternative, and assemble the recursive decompositions while
retaining their entire seams.  Connected \emph{dual face regions} above
are not automatically the connected primal shores needed for a Count
splice.  The new constructors produce the exact mesh existential used by
\texttt{RawDecompositionOfNoInjectiveMesh}; they do not prove that theorem,
fixed-mesh exclusion, or D1b/D1c.  The fixed ledger remains $57/100$.
The incidence and single-cycle frontier steps are now supplied below.

\subsection{Complete spherical distance contours}
\label{sec:complete-spherical-distance-contours}

\begin{theorem}[a facial bond has one complete primal boundary cycle]
\Sverified{Mettapedia.GraphTheory.FourColor.FacialBondPrimalCycle.boundary_support_eq_of_bond}
\Sverified{Mettapedia.GraphTheory.FourColor.FacialBondPrimalCycle.exists_cycle_of_facial_bond}
Let a finite simple graph carry the stack's connected bridgeless spherical
cubic rotation structure, with two-sided quotient faces.  If a set $S$ of
faces and its complement both induce connected facial-dual subgraphs,
then there is a simple primal cycle containing exactly the edges with one
incident face in $S$ and the other outside $S$.  Neither the cycle nor a
chosen boundary component is an input.
\end{theorem}

\begin{proof}
The binary facial boundary of $S$ is a nonzero circulation.  Cubicity
provides a simple cycle supported in this circulation.  Spherical
face-boundary/cycle-space equality represents that cycle by another face
set $T$.  Along every dual edge internal to $S$ or its complement, the
boundary of $T$ is zero, since its support is contained in the original
cut.  Its binary face labels are therefore constant on each of the two
connected sides.  The constants differ because the cycle is nonempty.
Consequently every edge of the original cut belongs to that cycle.
\end{proof}

\begin{theorem}[constructed, strictly nested distance regions and contours]
\Sverified{Mettapedia.GraphTheory.FourColor.FaceDistanceFrontiers.deepRegion_ssubset}
\Sverified{Mettapedia.GraphTheory.FourColor.SphericalDistanceContours.vertexFaces_clique}
\Sverified{Mettapedia.GraphTheory.FourColor.SphericalDistanceContours.mixed_iff_boundary_dart}
\Sverified{Mettapedia.GraphTheory.FourColor.SphericalDistanceContours.exists_distance_contour}
\Sverified{Mettapedia.GraphTheory.FourColor.SphericalDistanceContours.distance_frontiers_disjoint}
For the actual quotient facial dual $H$, fix faces $r,z$ and let $R_i$ be
the constructed deep component above.  At every $i<d_H(r,z)$ its complete
boundary is a simple primal cycle, whose vertex support is exactly the set
of vertices incident to both $R_i$ and its complement.  These cycles are
vertex-disjoint at distinct depths.  Moreover, $R_j\subsetneq R_i$ whenever
$i<j$ and $i<d_H(r,z)$.
\end{theorem}

The local incidence hypothesis is discharged, not retained: two outgoing
darts at a cyclic cubic vertex are consecutive in one orientation, so
their distinct face sectors share a dual edge.  A mixed vertex supplies a
literal boundary dart based there; the opposite edge face is also incident
there through the rotated dart.  Thus the complete cyclic frontier, not
just a subset of its vertices, is identified.  Strict nesting follows from
a walk from $z$ to $r$: its exit from $R_i$ has an inside endpoint at
distance $i+1$, which is absent from every later $R_j$.

Before formalization, the gate exhausted all $4094$ nonempty proper face
sets of the dodecahedral sphere.  Exactly $2336$ have connected selected
and complementary dual sides, and all have one complete boundary cycle.
For the $3\times3$ periodic honeycomb torus, $486$ of $510$ face sets have
both sides connected; $180$ have two boundary cycles.  This negative control
shows that the spherical hypothesis cannot simply be dropped.  Full edge,
face and rotation data are checked; the torus has $18$ vertices, $27$ edges,
nine hexagonal faces and Euler characteristic zero.  The executable gate
and replay receipt are \texttt{tools/fourcolor/v24\_facial\_bond\_gate.py}
and \texttt{results/fourcolor/v24\_facial\_bond\_gate.json}.  The concrete
controls are executable checks, not new Lean instances or zero-Count graphs.

This closes the contour-construction and cubic-incidence steps in the mesh
extraction argument.  It does not bound contour lengths, construct the
inter-contour linkage, turn the primal contours into Count splice nooses,
or assemble the recursive bounded-width decomposition.  Those geometric
obligations, fixed-mesh exclusion and the explicit route-native base remain
open.  No catalogue or reducibility computation is used.  The unchanged
obligation ledger is still $57/100$.

\subsection{Small vertex separators construct complete connected contour cuts}
\label{sec:contour-separator-bonds}

\begin{theorem}[separator-to-bond construction]
\Sverified{Mettapedia.GraphTheory.FourColor.VertexSeparatorBond.basin_connected}
\Sverified{Mettapedia.GraphTheory.FourColor.VertexSeparatorBond.filled_connected}
\Sverified{Mettapedia.GraphTheory.FourColor.VertexSeparatorBond.exists_bond_of_separator}
\Sverified{Mettapedia.GraphTheory.FourColor.VertexSeparatorBond.exists_bounded_bond}
Let $G$ be a connected finite simple graph, and let $A,B$ be disjoint
nonempty vertex sets, each inducing a connected subgraph.  Suppose every
$A$--$B$ walk meets a vertex set $X$, possibly at a terminal vertex.
If every vertex of $X$ has degree at most $d$, there is a vertex shore
$S$ with $A\subseteq S$, $B\subseteq S^c$, both induced shores connected,
and complete edge boundary $|\delta_G(S)|\le d|X|$.
\end{theorem}

\begin{proof}
Let $C$ consist of $A$ and every vertex reachable from $A$ by a walk
avoiding $X$.  Restoring all of $A$ connects these components.  The
separator property implies $B\subseteq C^c$.  Let $R$ be the component of
$G[C^c]$ containing $B$, and put $S=R^c$: all other complementary
components are filled into the first shore.  A walk from a vertex outside
$R$ to $C$ first reaches $C$ without entering $R$; otherwise that vertex
would belong to $R$.  Since $C$ is connected, $S$ is connected too.
Every edge from $R$ to $S$ has its $S$ endpoint in $C$, and every edge
leaving $C$ has an endpoint in $X$.  Charging crossing edges to the union
of the incidence sets at $X$ gives
$|\delta_G(S)|\le\sum_{x\in X}\deg_G(x)\le d|X|$.
\end{proof}

\begin{theorem}[application to the constructed spherical contours]
\Sverified{Mettapedia.GraphTheory.FourColor.SphericalContourSeparators.frontier_connected}
\Sverified{Mettapedia.GraphTheory.FourColor.SphericalContourSeparators.exists_bounded_bond_between_frontiers}
\Sverified{Mettapedia.GraphTheory.FourColor.SphericalContourSeparators.card_boundaryDart_le_edgeBoundary}
\Sverified{Mettapedia.GraphTheory.FourColor.SphericalContourSeparators.exists_ordered_bond_between_frontiers}
In the spherical cubic setting of the preceding subsection, take any two
distinct depths below $d_H(r,z)$.  A vertex separator $X$ between their
\emph{complete} mixed frontiers constructs a bond retaining the entire
first contour and excluding the entire second.  Its full edge boundary
and its actual exposed-port carrier each have size at most $3|X|$.
The computed complementary first-return boundary orders are inverse, and
each touched ambient face contains at most one retained boundary dart.
\end{theorem}

Connectivity of the terminals comes from the constructed primal cycles;
disjointness comes from the distance-incidence theorem.  Neither is an
extra supplied geometric witness.  The boundary-order conclusions use the
existing planar-bond theorem on the newly constructed shores.  The bound
is on the full ambient cut, including all lateral attachments, not on a
selected local carrier.  It is a factor-three bound, not an assertion of
a separator-length-optimal noose.

The preregistered gate exhausted all connected labelled simple graphs on
$2$ through $5$ vertices (respectively $1,4,38,728$ graphs), all ordered
disjoint connected terminal pairs, and all vertex separators.  All
$1{,}669{,}306$ cases passed; $1{,}635{,}062$ include separator vertices
on a terminal, and $306{,}400$ require filling complementary components.
The general-degree bound is attained in $14{,}266$ cases.  These are
executable generic-graph checks, not zero-Count graphs or Lean instances.
The gate and fixed-population replay are
\texttt{tools/fourcolor/v24\_separator\_bond\_gate.py}, with receipt
\texttt{results/fourcolor/v24\_separator\_bond\_gate.json}.

This discharges the conversion needed on the small-separator branch of
the contour extraction.  It does not yet construct the separator/linkage
alternative, uncross a family of separators, or assemble branch
decompositions with their complete recursive seams.  No global
counterexample supplier, fixed-mesh exclusion, or route-native base is
claimed.  The unchanged ledger remains $57/100$.

\subsection{Integral flow constructs the contour separator/linkage alternative}
\label{sec:integral-contour-linkage}

\begin{theorem}[finite integral max-flow/min-cut]
\Sverified{Mettapedia.GraphTheory.FourColor.IntegralFlow.Flow.exists_max_flow_min_cut}
For a finite directed network with natural capacities and distinct source
and sink, there exist an integer skew flow and a source--sink cut whose
values agree.  This flow maximizes the value among all feasible integer
flows, and this cut minimizes capacity among all source--sink cuts.
\end{theorem}

The proof uses boundedness of attainable integer flow values, not a
real-flow rounding theorem.  A simple residual route would augment a
maximizer by one; hence the sink is outside its residual-reachable shore.
Every outward arc of that shore is saturated.  Conservation and internal
skew cancellation identify its entire cut capacity with the flow value.
The existence proof uses classical finite mathematics; it does not claim
that the executable gate's search algorithm has been verified in Lean.

\begin{theorem}[positive-route extraction and unit-gate exclusion]
\Sverified{Mettapedia.GraphTheory.FourColor.IntegralFlow.Flow.exists_positive_route}
\Sverified{Mettapedia.GraphTheory.FourColor.IntegralFlow.Flow.exists_routes}
\Sverified{Mettapedia.GraphTheory.FourColor.IntegralFlow.Route.unit_gate_disjoint}
An integer flow of value at least $m$ supplies $m$ simple source--sink
routes using only positive-capacity arcs.  On each arc with zero reverse
capacity, their total signed usage is bounded by the original flow on
that arc.  Consequently two such routes cannot both use a unit-capacity
arc whose reverse capacity is zero.
\end{theorem}

A positive-valued flow reaches the sink through positive-flow arcs:
otherwise its reachable shore has nonpositive outward flow, contradicting
conservation.  Subtracting a simple positive route leaves a feasible
integer flow with value smaller by one.  Iteration proves the aggregate
usage bound.  The zero-reverse-capacity condition is essential to the
nonnegativity of each route's usage on a given gate.

\begin{theorem}[full-ambient vertex separator/linkage alternative]
\Sverified{Mettapedia.GraphTheory.FourColor.VertexSplitCut.separates_of_small_cut}
\Sverified{Mettapedia.GraphTheory.FourColor.VertexSplitCut.decode_route}
\Sverified{Mettapedia.GraphTheory.FourColor.VertexSplitCut.exists_linkage_of_flow}
\Sverified{Mettapedia.GraphTheory.FourColor.VertexSplitCut.Linkage.le_separator}
\Sverified{Mettapedia.GraphTheory.FourColor.VertexSplitCut.exists_linkage_or_separator}
For any finite simple graph $G$, terminal sets $A,B$ and natural number
$k$, either there are $k$ pairwise vertex-disjoint $A$--$B$ paths, or a
vertex set $X$ with $|X|<k$ meets every $A$--$B$ walk.  Endpoints count
towards disjointness and may belong to $X$.  Every such separator has
cardinality at least the size of any such linkage.  Neither connectivity
nor planarity is assumed here.
\end{theorem}

Each original vertex is split into an entry and exit joined by a unit
gate, including every terminal.  Source, sink and original adjacency
connectors have capacity $k$.  A cut of capacity less than $k$ cannot
cross a connector, so its crossing gates meet every original walk.
Otherwise the integral flow supplies $k$ simple routes.  Decoding the
entry--exit alternation produces walks in the original graph; every
visited vertex uses its own gate.  Removing repeated vertices preserves
the gate witnesses, and the unit-gate bound proves vertex disjointness.
No original adjacency is omitted, even if it lies outside a chosen
geometric band.

\begin{theorem}[constructed alternative between the spherical contours]
\Sverified{Mettapedia.GraphTheory.FourColor.SphericalContourSeparators.exists_linkage_or_ordered_bond}
In the spherical cubic setting of Section~\ref{sec:contour-separator-bonds},
take two distinct depths below $d_H(r,z)$ and $k\ge1$.  Either their
complete mixed frontiers are joined by $k$ vertex-disjoint ambient paths,
or there is a finite deleted region retaining the entire first frontier
and excluding the second, whose two induced shores are connected and
whose full edge boundary and exposed-port carrier both have size at most
$3(k-1)$.  Its touched-face uniqueness and opposite first-return boundary
orders are proved on that constructed region.
\end{theorem}

Here neither a separator nor a linkage is a supplied premise.  The small
branch composes the new alternative with the preceding component-filling
construction.  The large branch supplies actual paths in the complete
ambient graph, not merely an optimum value of a flow problem.

The preregistered integral gate exhausted all $4164$ directed loop-free
unit-capacity networks on $2$ through $4$ vertices, comparing residual
search with exhaustive cut capacities.  A six-vertex control requires
reverse residual cancellation.  Positive-flow subtraction was separately
checked throughout that population before its formalization.  The
vertex-splitting gate exhausted all labelled simple graphs on $2$ through
$4$ vertices, all ordered disjoint nonempty terminal sets, and thresholds
$1$ through $|V|+1$: $16{,}396$ cases, with $13{,}480$ small cuts and
$2916$ large-flow cases.  It independently exhausts original simple-path
packings and vertex separators, decodes the extracted routes, and checks
their vertex disjointness.  Both receipts replay the fixed populations:
\texttt{tools/fourcolor/v24\_integer\_flow\_gate.py} and
\texttt{tools/fourcolor/v24\_vertex\_split\_gate.py}, with corresponding
JSON receipts under \texttt{results/fourcolor/}.  These are generic
executable controls, not Lean-instantiated examples or zero-Count graphs.

The remaining global steps are to connect the linkage to the full family
of intermediate contours, make the small cuts compatible, and assemble
the recursive decomposition with complete seam accounting.  The cubic
factor three cannot be transferred unchanged to contracted hubs.  No
global supply theorem, fixed-mesh exclusion or explicit generating base
is claimed by this checkpoint; the fixed ledger remains $57/100$.

\subsection{Exact meshes, forced nesting, and a mesh-free dual-distance bound}
\label{sec:contour-mesh-node-chain}

Write $H$ for the actual face-dual graph, $R_d$ for the far-rooted
component beyond distance $d$, and $C_d$ for its complete mixed primal
frontier.  These are the regions and cycles constructed in
Section~\ref{sec:integral-contour-linkage}, not supplied local patches.

\begin{theorem}[all intermediate contour crossings]
\Sverified{Mettapedia.GraphTheory.FourColor.FaceDistanceFrontiers.walk_meets_mixed}
\Sverified{Mettapedia.GraphTheory.FourColor.SphericalContourSeparators.adjacent_share_face}
\Sverified{Mettapedia.GraphTheory.FourColor.SphericalContourSeparators.walk_meets_intermediate}
Every ambient walk from $C_i$ to $C_j$ meets $C_d$ whenever
$i\le d\le j$.  The walk is not required to stay inside the intervening
band.  Generically, a walk leaving the vertices incident to a face region
meets the mixed frontier, provided adjacent vertices share an incident face.
\end{theorem}

An original edge supplies the common face at its two endpoints.  The
incident-face clique property puts $C_i$ outside the vertices touching
$R_d$ for $i<d$, whereas $C_j$ touches $R_d$ for $d\le j$.  Applying the
generic crossing lemma to the whole walk derives all row--column
intersections, including excursions outside the selected band.

\begin{theorem}[the exact mesh-or-bond alternative]
\Sverified{Mettapedia.GraphTheory.FourColor.RotationWalkMeshPath.exists_vert_iff}
\Sverified{Mettapedia.GraphTheory.FourColor.SphericalContourSeparators.exists_injectiveMesh_of_contour_linkage}
\Sverified{Mettapedia.GraphTheory.FourColor.SphericalContourSeparators.exists_injectiveMesh_or_ordered_bond}
\Sverified{Mettapedia.GraphTheory.FourColor.SphericalContourSeparators.exists_ordered_bond_of_no_injectiveMesh}
Let $a\ge2$ and $\ell+a\le d_H(r,z)$.  Either the rotation
multigraph has an injective $a$-by-$b$ mesh in the exact predicate of
\texttt{RawDecompositionOfNoInjectiveMesh}, or a complete bond retains
$C_\ell$ and excludes $C_{\ell+a-1}$.  Both induced shores are connected;
the entire crossing-edge and exposed-port carriers have size at most
$3(b-1)$; touched-face uniqueness and inverse complementary first-return
orders hold on that bond.
\end{theorem}

Rows are the $a$ actual consecutive contour cycles; columns are the
constructed vertex-disjoint linkage.  The walk adapter chooses original
rotation-multigraph edges without changing the vertex carrier or support.
Distinct depths make rows vertex-disjoint and the linkage makes columns
vertex-disjoint.  The preceding crossing theorem supplies every
intersection.  The conclusion uses the stack's mesh predicate, not an
unstated grid-minor or ordered-crossings surrogate.  No row, column,
intersection, separator or small cut is an input to the alternative.

\begin{theorem}[localization forces strict nesting]
\Sverified{Mettapedia.GraphTheory.FourColor.SphericalContourSeparators.contour_bond_avoids_touches}
\Sverified{Mettapedia.GraphTheory.FourColor.SphericalContourSeparators.contour_bond_contains_outside}
\Sverified{Mettapedia.GraphTheory.FourColor.SphericalContourSeparators.contour_bonds_strict}
A connected bond retaining $C_i$ and excluding $C_j$, $i<j$, retains
every vertex incident to a face outside $R_i$ and excludes every vertex
incident to a face inside $R_j$.  Consequently any two such bonds in
windows $i<j\le p<q$ have strictly nested retained vertex sets.
\end{theorem}

This removes the need to uncross independently chosen bonds in separated
windows.  A path within either connected shore cannot cross the opposite
terminal frontier.  The later first frontier witnesses strictness.
Connectedness matters; the claim does not extend to arbitrary vertex sets
having only the required terminal inclusions.

\begin{theorem}[constructed nested bonds and exact-support nodes]
\Sverified{Mettapedia.GraphTheory.FourColor.SphericalContourSeparators.exists_nested_contour_bonds}
\Sverified{Mettapedia.GraphTheory.FourColor.SphericalContourSeparators.crossingEquivBoundary}
\Sverified{Mettapedia.GraphTheory.FourColor.SphericalContourSeparators.OrderedContourBond.toConnectedNode}
\Sverified{Mettapedia.GraphTheory.FourColor.SphericalContourSeparators.exists_nested_contour_nodes}
Suppose the exact injective $a$-by-$b$ mesh is absent and
$na\le d_H(r,z)$.  The windows $[ta,ta+a-1]$, $0\le t<n$,
construct $n$ complete ordered bonds of width at most $w=3(b-1)$,
strictly nested as vertex shores and full incident-edge shores.
In a graph-backed least counterexample these produce $n$ strictly nested
connected exact-support nodes, each with literal middle width at most $w$.
\end{theorem}

The node conversion reuses the existing connected incident-edge adapter.
The two terminal contour cycles guarantee nonempty majority sides;
equality of the original vertex side with its majority side is unnecessary.
Crossing darts and actual exposed ports are equivalent, so the
middle-set bound is on the quantity consumed by the splice theorem.
The conversion does not silently identify the bond's vertex boundary
order with that of a possibly different majority side: the existing
literal-node constructor derives the latter from the connected edge shores.

Define the existing raw varying-width state bound by
\[
 N(w)=(6w+1)\sum_{j=0}^{w}j!\,2^{3^j}.
\]
\begin{theorem}[explicit dual-distance bound in a mesh-free least counterexample]
\Sverified{Mettapedia.GraphTheory.FourColor.SphericalContourSeparators.le_contourStateBound_of_nested_nodes}
\Sverified{Mettapedia.GraphTheory.FourColor.SphericalContourSeparators.dual_distance_lt_of_no_injectiveMesh}
Under the preceding exact no-mesh premise and $a\ge2$, every pair of
quotient faces in a graph-backed vertex-minimal Tait counterexample satisfies
\[
 d_H(r,z)<a\bigl(N(3(b-1))+1\bigr).
\]
\end{theorem}

The supplied-chain assumption is absent: if the displayed bound failed,
the construction would give $N(w)+1$ nested nodes.  Their card-phased
states are pairwise distinct by the already proved physical splice and
minimality theorem, contradicting the exact state count.  This is a raw
finite-state bound, not a new claim of efficient enumeration.

The crossing gate replays all $4950$ depth triples in the verified
$GP(k,0)$ archive, $1\le k\le8$, deleting the entire intermediate
frontier in the full ambient graph.  An abstract two-vertex control
demonstrates failure without common-face incidence.  The nesting gate
constructs three deterministic grown-and-filled bond samples in every depth
window (repetitions allowed): $2592$ samples and $239{,}598$ separated-window comparisons.  It
checks full-shore connectedness, localization, strict vertex and
incident-edge inclusion, and the actual middle-width bound.  Adding a
deep isolated vertex gives a negative connected-shore control.  These are
generic geometric controls, not zero-Count examples or colouring
certificates.  Replays are
\texttt{v24\_contour\_separation\_gate.py} and
\texttt{v24\_contour\_bond\_nesting\_gate.py} under
\texttt{tools/fourcolor/}, with corresponding JSON receipts under
\texttt{results/fourcolor/}.

\paragraph{Remaining global obligations.}
Fixed-mesh exclusion is still an explicit open route premise.  The new
bound is on face-dual distance, not vertex count: turning bounded face
size and this distance bound into an explicit size bound, and handling
the long-face case, remain geometric work.  A complete raw recursive
branch decomposition and the actual generating-base audit are not proved
here.  This checkpoint closes intermediate crossings, separated-window
compatibility and the exact-node conversion, but not all of D1b or D1c.
The fixed obligation ledger remains $57/100$.

\subsection{Bounded-face size transfer and an explicit vertex threshold}
\label{sec:bounded-face-size-transfer}

The distance bound of Section~\ref{sec:contour-mesh-node-chain} now has
an actual vertex-count consequence when all quotient-face boundaries have
at most $B$ edges.  This is the bounded-face component of the geometric
transfer; the long-face component remains separate.

\begin{theorem}[degree and Euler on the actual face carrier]
\Sverified{Mettapedia.GraphTheory.FourColor.SphericalContourSeparators.orbitFaceDualGraph_degree_le_boundary}
\Sverified{Mettapedia.GraphTheory.FourColor.SphericalContourSeparators.card_vertices_add_four_eq_two_mul_faces}
The degree of a quotient face in the full face-dual graph is at most the
cardinality of its edge boundary.  A cubic rotation system satisfying the
quotient-face spherical Euler equation satisfies
\[
 |V|+4=2|F|.
\]
Here $F$ is the quotient of darts by face orbits, not a deduplicated
collection of face-edge sets.
\end{theorem}

The degree proof transports the existing at-most-two-face incidence bound
through the full ambient-face subtype.  The Euler calculation combines
$3|V|=2|E|$ with $|V|-|E|+|F|=2$.  Thus both counts refer to the same
literal rotation system as the preceding contour construction.

\begin{theorem}[vertex largeness supplies a long dual axis]
\Sverified{Mettapedia.GraphTheory.FourColor.SphericalContourSeparators.exists_dual_distance_ge_of_large_vertex_card}
\Sverified{Mettapedia.GraphTheory.FourColor.SphericalContourSeparators.vertex_card_le_of_dual_distances_lt}
For a bridgeless spherical cubic map with every face boundary of size at
most $B$, if
\[
 2(B+1)^r-4<|V|,
\]
then two actual quotient faces have dual distance at least $r$.
Consequently, if all dual distances are strictly below $r$, then
$|V|\le 2(B+1)^r-4$.  Subtraction in the formal threshold is natural-number
subtraction.
\end{theorem}

This reuses the existing bounded-degree closed-neighborhood estimate:
a radius-$r$ reachable set has at most $(B+1)^r$ vertices.  Dual
connectedness, that estimate and the literal Euler count supply the axis;
no path, contour or decomposition is a hypothesis of this transfer.

\begin{theorem}[bounded-face reduction from vertex largeness]
\Sverified{Mettapedia.GraphTheory.FourColor.SphericalContourSeparators.exists_nested_nodes_of_bounded_faces_large}
\Sverified{Mettapedia.GraphTheory.FourColor.SphericalContourSeparators.vertex_card_le_of_noMesh_and_bounded_faces}
Let $a\ge2$ and consider a graph-backed vertex-minimal Tait
counterexample with face sizes at most $B$ and no exact injective
$a$-by-$b$ mesh.  Set $w=3(b-1)$ and retain the explicit $N(w)$ from
Section~\ref{sec:contour-mesh-node-chain}.  Vertex count greater than
$2(B+1)^{na}-4$ constructs $n$ strictly nested connected exact-support
nodes of width at most $w$.  In particular,
\[
 |V|\le 2(B+1)^{a(N(w)+1)}-4.
\]
\end{theorem}

Taking $n=N(w)+1$ contradicts the existing physical-splice state bound.
The threshold depends only on the displayed natural-number parameters,
not on the unknown least-counterexample size.  The proof is
universe-polymorphic.  It does not discharge the exact no-mesh premise.

The preregistered gate revalidates the $GP(k,0)$ archive for $1\le k\le8$,
calibrated by $|V|=20k^2$ and $|F|=10k^2+2$.  Every root and reachable
radius is checked using the complete dual graph and cubic Euler counts.
Four thickness-two annulus controls have circumferences $8,16,32,64$:
their dual diameters are all $4$, while their vertex counts are
$48,96,192,384$ and their largest face sizes grow with circumference.
There are $38{,}432$ reach/count checks in total.  These are finite
generic geometric controls; they are neither zero-Count instances nor
a formal unbounded-family refutation.  The replay is
\texttt{tools/fourcolor/v24\_bounded\_face\_size\_gate.py}, with its
JSON receipt under \texttt{results/fourcolor/}.

\paragraph{Remaining global obligations.}
The bounded-face component now supplies actual nodes from vertex count.
Long faces cannot simply be handled by dropping the face-size hypothesis
from the counting step.  Their complete-interface construction, fixed-mesh
exclusion, a full raw branch-decomposition supply and the generating-base
audit remain open.  Since the fixed D1c obligation includes both geometric
transfers, the ledger remains $57/100$.

\subsection{Ranked dual forests and the long-face tree--cotree gate}
\label{sec:ranked-dual-forest}

A possible way to handle long faces is to choose a breadth-first tree in
the full face dual and retain the complementary original edges in the
primal graph.  The intended small cuts are the fundamental cuts of that
primal tree, rather than cuts forced to keep an entire long face on one
side.  The following generic lemmas establish the parity and connectivity
parts of this construction.

\begin{theorem}[ranked forest parity]
\Sverified{Mettapedia.GraphTheory.FourColor.RankedForestParity.exists_singleton_incidence}
\Sverified{Mettapedia.GraphTheory.FourColor.RankedForestParity.eq_empty_of_even_incidence}
Orient a finite edge set from child to parent, with distinct children and
a natural-number rank strictly decreasing along every edge.  Every
nonempty selected edge set has a vertex incident to exactly one selected
edge.  Consequently an edge set with even incidence at every vertex is
empty.
\end{theorem}

Choose a selected edge whose child has maximum rank.  No other selected
edge can enter that child, since its own child would have larger rank;
injectivity of the child assignment excludes a second outgoing edge.
This argument uses neither Euler counting nor planarity.

\begin{theorem}[the primal complement of a ranked dual forest]
\Sverified{Mettapedia.GraphTheory.FourColor.RankedDualForestComplement.RankedDualForest.crossing_eq_empty}
\Sverified{Mettapedia.GraphTheory.FourColor.RankedDualForestComplement.RankedDualForest.complement_connected}
\Sverified{Mettapedia.GraphTheory.FourColor.RankedDualForestComplement.RankedDualForest.complement_isTree}
In a connected graph-backed rotation system with two-sided quotient
faces, let a selected original-edge set carry a ranked forest on the
incident faces.  Every full primal cut contained in that set is empty.
Deleting the selected edges preserves primal connectedness.  If the
quotient-face Euler equation is spherical and the selected set has
$|F|-1$ edges, its primal complement is a spanning tree.
\end{theorem}

The forest incidence field is exact: the quotient faces incident to each
selected original edge are its child and parent.  Existing literal facial
transition parity makes every full primal cut even in this incidence.
The preceding lemma therefore rules out a nonempty cut inside the forest.
Applying this to a component of the edge-deleted graph proves
connectedness.  Euler then gives $|V|-1$ complementary edges, and the
existing finite connected-graph cardinality theorem gives acyclicity.
The spanning and metric properties of an actual dual breadth-first tree
are not asserted by these relative lemmas.

The preregistered tree--cotree gate checks $24$ spherical maps: $GP(k,0)$
for $1\le k\le8$ and annuli of thickness $1$ through $4$ and circumference
$8,16,32,64$.  At up to three deterministic dual roots per map, it
constructs the dual breadth-first tree, its primal complement, and every
fundamental primal cut: $18{,}885$ full cuts.  Each cut equals the entire
dual fundamental cycle and has size at most $2h+1$, where $h$ is that
dual tree's height.  Both induced shores are connected; the complete dual
boundary is cyclic; the two literal facial first-return permutations are
inverse; and the incident-edge middle set has size at most the cut size.
Along longest primal tree paths, $2147$ consecutive comparisons verify
strict vertex and incident-edge nesting.  The gate also records which
cuts have cycles on both shores.  No colouring is enumerated.

On the sampled annuli the maximum cut widths are $6,8,10,12$, respectively,
independent of the four tested circumferences.  The auxiliary parity gate
exhausts all $1100$ labelled simple graphs on at most five vertices;
among the $340$ forests, every one with a nonempty edge set has an odd-degree vertex.
These results are executable controls, not a formal general metric bound
or examples in the zero-Count target class.  The replay and receipt are
\texttt{tools/fourcolor/v24\_tree\_cotree\_gate.py} and its corresponding
JSON under \texttt{results/fourcolor/}.

\paragraph{Next construction.}
The following subsection constructs the ranked forest and its $|F|-1$
edge count from actual dual distances.  The full fundamental-cut width
bound and extraction of enough cyclic nested shores remain to be proved.
Until those steps are complete, this is a gated long-face construction,
not a global supply theorem.
The fixed ledger remains $57/100$.

\subsection{Constructing the dual BFS tree on original edges}
\label{sec:dual-bfs-tree}

\begin{theorem}[actual distance-ranked dual tree]
\Sverified{Mettapedia.GraphTheory.FourColor.SphericalDualBFSTree.exists_adj_dist_pred}
\Sverified{Mettapedia.GraphTheory.FourColor.SphericalDualBFSTree.edge_incidence_iff_pair}
\Sverified{Mettapedia.GraphTheory.FourColor.SphericalDualBFSTree.of_connected}
For any root quotient face in a connected full facial dual, select one
original edge from every nonroot face to a predecessor at distance one
less.  These edges form a ranked dual forest with exactly $|F|-1$ edges.
Every nonroot face is the child of a selected edge; its rank is its
actual full-dual distance to the root, and each parent has rank exactly
one less.
\end{theorem}

The predecessor comes from the first edge of a shortest dual walk.
Its adjacency supplies an original shared edge.  Two distinct incident
faces exhaust that edge's quotient-face incidence.  If two distinct
children chose the same edge, each would be the other's predecessor,
contradicting strict distance decrease.  This proves injectivity of the
selection and hence the $|F|-1$ count.  No tree, contour, bounded-face
condition or decomposition is supplied as an input to the construction.

\begin{theorem}[selected paths and the primal spanning tree]
\Sverified{Mettapedia.GraphTheory.FourColor.SphericalDualBFSTree.DualBFSTree.faceGraph_le}
\Sverified{Mettapedia.GraphTheory.FourColor.SphericalDualBFSTree.DualBFSTree.exists_root_walk}
\Sverified{Mettapedia.GraphTheory.FourColor.SphericalDualBFSTree.DualBFSTree.faceGraph_connected}
\Sverified{Mettapedia.GraphTheory.FourColor.SphericalDualBFSTree.DualBFSTree.faceGraph_dist_le}
\Sverified{Mettapedia.GraphTheory.FourColor.SphericalDualBFSTree.DualBFSTree.dist_root_eq}
\Sverified{Mettapedia.GraphTheory.FourColor.SphericalDualBFSTree.DualBFSTree.exists_root_path}
\Sverified{Mettapedia.GraphTheory.FourColor.SphericalDualBFSTree.DualBFSTree.faceGraph_isTree}
\Sverified{Mettapedia.GraphTheory.FourColor.SphericalDualBFSTree.of_connected_complement_isTree}
The selected face graph is a spanning tree of the full dual and preserves exact distances to
the root.  It supplies simple root paths of those lengths, and the
selected-graph distance between $f$ and $g$ is at most
$d(f,r)+d(g,r)$.  For a graph-backed bridgeless spherical cubic map with
two-sided quotient faces, the primal complement of the constructed
dual BFS edges is a spanning tree.
\end{theorem}

Induction on rank constructs the root walks.  Every selected edge changes
rank by one, so no shorter root walk exists.  The selected dual graph is
itself a tree by connectedness and the constructed edge count.  The primal conclusion
instantiates the ranked-forest complement theorem with the constructed
edge set and count, using the actual map's connectedness and spherical
Euler equation.  Two-sidedness is an explicit hypothesis already present
in the vertex-minimal target class.  The construction itself only needs
connected facial geometry.

\paragraph{Exact remaining metric bridge.}
The estimate $d(f,r)+d(g,r)\leq 2h$ concerns selected dual paths.
The following subsection identifies each \emph{complete} primal
fundamental cut by original-edge parity and proves its width is at most
$2h+1$.  Long nested cyclic shores, the global supply
consumer, fixed-mesh exclusion and the generating base are not discharged
by this construction.  The fixed ledger remains $57/100$.

\subsection{Complete cotree cuts without a face-size bound}
\label{sec:complete-cotree-cuts}

Let $T_D$ be the selected original edges of the constructed dual BFS
tree, rooted at the quotient face $r$, and let $T_P=G\setminus T_D$ be
its primal spanning-tree complement.  Incidence parity below is computed
on original edges.  Different edges joining the same two quotient faces
are never identified.

\begin{theorem}[parity identifies the entire fundamental cut]
\Sverified{Mettapedia.GraphTheory.FourColor.FiniteEdgeParity.boundary_symmDiff}
\Sverified{Mettapedia.GraphTheory.FourColor.SphericalCotreeCutWidth.exists_root_support}
\Sverified{Mettapedia.GraphTheory.FourColor.SphericalCotreeCutWidth.eq_of_boundary_eq_of_sdiff_eq}
\Sverified{Mettapedia.GraphTheory.FourColor.SphericalCotreeCutWidth.crossing_eq_fundamental_support}
\Sverified{Mettapedia.GraphTheory.FourColor.SphericalCotreeCutWidth.crossing_card_le_of_singleton_sdiff}
Suppose a full primal cut $C$ meets $T_P$ in just the original edge $e$,
whose incident quotient faces are $f,g$.  With $P_f,P_g\subseteq T_D$
the constructed root-path supports,
\[
 C=P_f\mathbin{\triangle}P_g\mathbin{\triangle}\{e\},
 \qquad |C|\leq d(f,r)+d(g,r)+1\leq 2h+1
\]
whenever every face has root distance at most $h$.
\end{theorem}

The finite incidence boundary takes values in $\mathbb Z/2\mathbb Z$;
symmetric difference adds these boundaries.  Rank induction constructs
$P_f$ with at most $d(f,r)$ original edges and boundary supported exactly
at $f,r$ (with cancellation when they coincide).  The two root boundaries
and the boundary of $e$ cancel.  Full primal cuts also have zero boundary,
by literal facial-transition parity.  Their symmetric difference is an
even edge set supported entirely in the ranked forest $T_D$, so the
previous forest-parity lemma makes it empty.  This proves equality of
the whole cut, not merely containment of a local carrier, and needs no
unique-shared-edge hypothesis.

\begin{theorem}[constructing both complete connected shores]
\Sverified{Mettapedia.GraphTheory.FourColor.SphericalCotreeShore.reachable_or_reachable_after_delete}
\Sverified{Mettapedia.GraphTheory.FourColor.SphericalCotreeShore.exists_bounded_fundamental_shore}
\Sverified{Mettapedia.GraphTheory.FourColor.SphericalCotreeShore.exists_bounded_cotree_shore}
For each edge $uv$ of the actual primal cotree $T_P$, deleting $uv$
constructs a side containing $u$ and avoiding $v$.  Both this side and its
complement induce connected subgraphs of the original graph $G$.  Its
complete original cut meets $T_P$ in exactly $uv$ and has at most $2h+1$
edges.
\end{theorem}

The side is the reachable region of $u$ in $T_P-uv$.  Tree acyclicity
separates $u$ and $v$, while connectedness of $T_P$ puts every vertex in
one of their two reachable regions after deletion.  These regions give
connected original shores.  Every retained tree edge stays within a
region, so the full cut has precisely one non-$T_D$ edge.  The preceding
parity theorem now applies to the constructed cut.

\begin{theorem}[mesh-free target-class width, with no face-size premise]
\Sverified{Mettapedia.GraphTheory.FourColor.SphericalCotreeShore.exists_bounded_cotree_shore_of_no_injectiveMesh}
In a graph-backed vertex-minimal counterexample with no vertex-injective
$a\times b$ mesh, $a\geq2$, every edge of the constructed primal cotree
supplies such a pair of connected shores of width at most
\[
 W(a,b)=2\bigl(N(3(b-1))+1\bigr)a+1,
 \qquad
 N(w)=(6w+1)\sum_{j=0}^{w}j!\,2^{3^j}.
\]
\end{theorem}

Here $N$ is the already proved normalized contour-state bound.  Its
mesh-free dual-distance consequence supplies the root-radius premise of
the generic cotree construction.  The width $W$ is an explicit
natural-number expression independent of face sizes and vertex count.
The mesh-exclusion premise is retained visibly; this theorem does not
establish fixed-mesh exclusion.

The original preregistered spherical tree--cotree gate was replayed
successfully.  An additional original-edge parity regression exhausts
all labelled trees on one through five vertices, every root and every
possible extra endpoint pair, including parallel copies of tree edges.
The calibrated tree counts are $1,1,3,16,125$.  It checks $6663$
fundamental supports, including $2712$ parallel-edge cases, and $206368$
subset-boundary comparisons.  No Tait or ring colouring is enumerated.
The script is \texttt{tools/fourcolor/v24\_fundamental\_parity\_gate.py},
with its replay receipt under \texttt{results/fourcolor/}.

\paragraph{Consumer.}
The complete cuts now exist with a uniform target-class width.  The next
section constructs their nested cyclic subchain from large vertex count
and feeds it to the exact-support state bound.  Fixed-mesh exclusion and
the generating-base audit are not supplied by the width theorem.

\subsection{Cotree paths give a mesh-free size bound and the raw supplier}
\label{sec:cotree-path-size}

\begin{lemma}[strict nested fundamental shores]
\Sverified{Mettapedia.GraphTheory.FourColor.TreePathShore.getVert_mem_iff}
\Sverified{Mettapedia.GraphTheory.FourColor.TreePathShore.shore_ssubset}
Let $p=(v_0,\ldots,v_L)$ be a simple path in a tree $H$.
For $0\le i<L$, let $S_i$ be the component of $v_0$ after deleting
$v_iv_{i+1}$.  Then $v_j\in S_i$ exactly when $j\le i$.
For $i<j<L$, $S_i\subsetneq S_j$.  Moreover
\[
                    |S_i|\ge i+1,
        \qquad |V(H)\setminus S_i|\ge L-i.
\]
\Sverified{Mettapedia.GraphTheory.FourColor.TreePathShore.prefix_card_le}
\Sverified{Mettapedia.GraphTheory.FourColor.TreePathShore.suffix_card_le}
\end{lemma}

The proof transfers path prefixes and suffixes to the edge-deleted graph.
Path simplicity prevents either segment from using the removed edge;
tree acyclicity separates its endpoints.  A later removed edge lies
outside the earlier start component, which gives inclusion.  The path
vertex between the cuts witnesses strictness.  The size estimates inject
the whole prefix and suffix, not only their endpoints.

\begin{theorem}[large cardinality supplies nested exact-support nodes]
\Sverified{Mettapedia.GraphTheory.FourColor.SphericalCotreePathChain.exists_nested_nodes_of_large_card}
Let $G$ belong to the graph-backed vertex-minimal target class.  Suppose
all full-dual distances to an actual root face are at most $h$, and put
$w=2h+1$.  If
\[
                         |V(G)|>4^{n+2w},
\]
then $G$ has $n$ strictly nested connected-shore nodes at width at most
$w$, with their literal exact-support states.
\end{theorem}

The primal cotree is a spanning tree of maximum degree at most three.
The existing bounded-degree estimate therefore constructs a simple path
of length at least $n+2w$.  Choose its cuts at positions $w+i$,
$0\le i<n$.  Both sides contain more than $w$ vertices.  The previous
complete-cut theorem bounds their entire original boundary by $w$ and
gives connected original shores.
\Sverified{Mettapedia.GraphTheory.FourColor.SphericalCotreePathChain.path_shore_bounds}
In a cubic graph a connected acyclic shore has at most its boundary size
minus two vertices, so both selected shores contain cycles.
\Sverified{Mettapedia.GraphTheory.FourColor.SphericalCotreePathChain.hasCycle_of_boundary_lt_card}
The existing cyclic-cut adapter supplies every connected-shore node
field.  Connected complements preserve strict incident-edge nesting;
the literal node constructor supplies its existing planar boundary-order
and exact-support data.  No chain, path, cyclic shore or ordering receipt
is assumed in the displayed large-cardinality theorem.

\begin{theorem}[mesh-free size without bounded faces]
\Sverified{Mettapedia.GraphTheory.FourColor.SphericalCotreePathChain.vertex_card_le_of_no_injectiveMesh}
For $a\ge2$, define $W(a,b)$ as in Section~\ref{sec:complete-cotree-cuts}
and
\[
             U(a,b)=4^{N(W(a,b))+1+2W(a,b)}.
\]
Every graph-backed vertex-minimal Tait counterexample containing no
vertex-injective permissive $a\times b$ mesh has at most $U(a,b)$ vertices.
\end{theorem}

Its actual full-dual radius is supplied by the previously proved
mesh-free distance theorem.  Were the vertex bound violated, the
large-cardinality theorem would supply $N(W)+1$ strictly nested nodes,
contradicting their checked state bound $N(W)$.  The root face is the
orbit of the distinguished outer dart.  No bounded-face, supplied
decomposition or unknown-counterexample-size parameter remains.
Equivalently, any larger instance in that exact target class supplies a
vertex-injective permissive $a\times b$ mesh.
\Sverified{Mettapedia.GraphTheory.FourColor.SphericalCotreePathChain.exists_injectiveMesh_of_large_card}
This positive structural alternative assumes neither a chain nor mesh
absence.  Its mesh is permissive; the ordered upgrade, at a larger explicit
threshold, is proved in Section~\ref{sec:ordered-mesh-supply}.

\begin{proposition}[the permissive-mesh raw premise is discharged]
\Sverified{Mettapedia.GraphTheory.FourColor.SphericalMeshFreeRawDecomposition.rawDecompositionOfNoInjectiveMesh}
For $a\ge2$, the existing proposition
\texttt{RawDecompositionOfNoInjectiveMesh} holds with parameters
$(a,b,U(a,b))$.
\end{proposition}

This adapter is deliberately coarse.  A cubic graph-backed instance has
at least three distinct edges.  List all its edges without repetition,
use two as the root and one leaf, and place the nonempty remainder in a
binary comb.  The complete rooted leaf encoding is checked, and every
middle set is a subset of $V(G)$, so its width is at most $U$.
\Sverified{Mettapedia.GraphTheory.FourColor.SphericalMeshFreeRawDecomposition.exists_rooted_decomposition}
\Sverified{Mettapedia.GraphTheory.FourColor.SphericalMeshFreeRawDecomposition.widthAtMost_of_card_le}
Consequently the existing permissive-mesh headline specialization now
needs only fixed-mesh exclusion and the route-native base through
$B(U,U)$; its decomposition premise is no longer supplied externally.
\Sverified{Mettapedia.GraphTheory.FourColor.SphericalMeshFreeRawDecomposition.combinatorialFourColorStatement_of_fixedMeshExclusion_and_base}
This is not the general planar grid theorem, nor a claim that this very
large base is computationally feasible.

The sharper consumer bypasses that coarse compatibility adapter: apply
the Count-derived size bound $U$ directly to the existing
minimal-counterexample selection theorem.  Fixed permissive-mesh exclusion
and \texttt{TaitBaseVerifiedAt}$(U(a,b))$ already imply the same
combinatorial spherical headline.
\Sverified{Mettapedia.GraphTheory.FourColor.SphericalMeshFreeRawDecomposition.combinatorialFourColorStatement_of_fixedMeshExclusion_and_direct_base}
Thus this specialization does not require the additional expansion to
$B(U,U)$.  Both the fixed exclusion and the actual base verification remain
visible unproved hypotheses; the smaller required bound is still enormous.

\paragraph{Tests and remaining boundary.}
Before formalization, the tree-path gate exhausted all labelled trees
through six vertices and all ordered endpoint pairs, calibrated by
$1,1,3,16,125,1296$ trees.  It checked $50069$ paths, $87858$ full tree
cuts and $46266$ strict comparisons, including the prefix and suffix
cardinality bounds.  The separate spherical full-cut/order gate was
replayed.  A further comb control checked $1144$ cuts on cycles of sizes
$3$--$32$ and complete graphs of sizes $3$--$8$.  These are generic
controls, not claimed zero-Count examples; no colouring or configuration
catalogue is enumerated.

The wall-facing \emph{ordered}-mesh exclusion hypothesis is weaker than
permissive-mesh exclusion.  The next section supplies its decomposition
premise without silently strengthening that exclusion hypothesis.
The boundary/frozen-data generality of
Proposition~\ref{prop:target-forced-cut-chain}, fixed-mesh exclusion,
optimized physical closure and the actual generating base are not
discharged here.  Thus no complete fixed ledger row is newly closed:
the score remains $57/100$, despite the proved mathematical state change.

\subsection{First-hit ordering closes the ordered-mesh supplier}
\label{sec:ordered-mesh-supply}

\begin{lemma}[simultaneous monotone thinning]
\Sverified{Mettapedia.GraphTheory.FourColor.SimultaneousMonotoneThinning.exists_orderEmbedding}
Let $f_1,\ldots,f_a$ be injective natural-number coordinates on an ordered set of $q$
indices.  If $q>b^{2^a}$, there is an increasing selection of $b+1$
indices on which each $f_i$ is strictly increasing or strictly decreasing.
The direction may differ between coordinates.
\end{lemma}

The proof iterates the finite Erd\H{o}s--Szekeres theorem already in
mathlib's \texttt{Archive.Wiedijk100Theorems.AscendingDescendingSequences}.
At each step the increasing embedding of the surviving indices preserves
all earlier coordinate orders.  The induction squares the threshold;
it does not assume that the original row orders agree.

\begin{lemma}[ordered first encounters with full contours]
\Sverified{Mettapedia.GraphTheory.FourColor.SphericalContourSeparators.exists_strictMono_frontier_hits}
A full-ambient walk from distance contour $\ell$ to contour
$\ell+a-1$ has strictly increasing first-hit positions on the $a$
consecutive contours.  Repeated encounters and excursions outside the
displayed band are permitted.
\end{lemma}

Apply the previously checked intermediate-contour crossing theorem to
the prefix ending at the first encounter with a later contour.  It must
already meet each earlier one.  Distinct complete distance frontiers
share no vertices, so the first-hit positions are strictly ordered.
A simple contour cycle becomes a simple row path by removing its first
edge; all its vertices remain, including the old start at the path's end.
The walk-to-multigraph adapter preserves path-vertex injectivity.
\Sverified{Mettapedia.GraphTheory.FourColor.RotationWalkMeshPath.cycle_tail_support}
\Sverified{Mettapedia.GraphTheory.FourColor.RotationWalkMeshPath.vert_injective_of_isPath}

\begin{theorem}[ordered mesh or complete bounded bond]
\Sverified{Mettapedia.GraphTheory.FourColor.SphericalContourSeparators.exists_orderedMesh_of_contour_linkage}
\Sverified{Mettapedia.GraphTheory.FourColor.SphericalContourSeparators.exists_orderedMesh_or_contour_bond}
Put $q(a,b)=b^{2^a}+1$ and $c(a,b)=3(q(a,b)-1)$.
In a bridgeless spherical cubic instance, a window of $a\ge2$
consecutive full distance contours supplies either an ordered injective
$a\times b$ mesh or a complete contour bond of width at most $c(a,b)$.
The bond has both original shores connected, bounded exposed-port count,
unique boundary-face incidences and inverse opposite first-return orders.
\end{theorem}

The existing ambient linkage--separator alternative is used at size $q$.
For the linkage branch, take the first-hit vertices on each column and
their unique positions on each opened simple row.  Distinct columns give
injective row-position coordinates.  The thinning lemma selects common
columns; reverse each decreasing row and retain $b$ of the selected
$b+1$ columns.  Both branch-position orders and every field of the
existing \texttt{OrderedInjectiveMesh} are then proved.
\Sverified{Mettapedia.GraphTheory.FourColor.OrderedContourExtraction.exists_orderedMesh_of_paths}
No unique-intersection claim is made: the carrier permits additional
row--column intersections or shared path segments, as required by cubic
wall representations.  Each individual path is simple and each family
is vertex-disjoint.  For the separator branch, the prior full-bond
constructor retains all lateral attachments and its complete order data.

\begin{theorem}[ordered-mesh-free size and raw decomposition]
\Sverified{Mettapedia.GraphTheory.FourColor.SphericalOrderedMeshReduction.vertex_card_le_of_no_orderedMesh}
\Sverified{Mettapedia.GraphTheory.FourColor.SphericalOrderedMeshReduction.rawDecompositionOfNoOrderedInjectiveMesh}
Write $N$ for the contour-state bound, and define
\[
 h_o(a,b)=(N(c(a,b))+1)a,\qquad w_o(a,b)=2h_o(a,b)+1,
 \qquad U_o(a,b)=4^{N(w_o(a,b))+1+2w_o(a,b)}.
\]
For $a\ge2$, every graph-backed vertex-minimal Tait counterexample with
no ordered injective $a\times b$ mesh has at most $U_o(a,b)$ vertices.
The pre-existing \texttt{RawDecompositionOfNoOrderedInjectiveMesh}
holds at parameters $(a,b,U_o(a,b))$.
\end{theorem}

Ordered-mesh absence forces a complete bond in every depth window.
The geometry-only nesting argument is shared with the permissive branch.
\Sverified{Mettapedia.GraphTheory.FourColor.SphericalContourSeparators.ordered_contour_bonds_nested}
The terminal cycles supply cyclic shores and the literal Count nodes;
their finite state bound gives full-dual diameter less than $h_o$.
\Sverified{Mettapedia.GraphTheory.FourColor.SphericalOrderedMeshReduction.exists_nested_nodes_of_no_orderedMesh}
\Sverified{Mettapedia.GraphTheory.FourColor.SphericalOrderedMeshReduction.dual_distance_lt_of_no_orderedMesh}
The complete cotree-path construction then gives the displayed vertex
bound, with no bounded-face or supplied-decomposition hypothesis.
Equivalently, every larger instance in that target class supplies an
actual ordered mesh.
\Sverified{Mettapedia.GraphTheory.FourColor.SphericalOrderedMeshReduction.exists_orderedMesh_of_large_card}
The checked complete edge comb supplies the compatibility raw decomposition.

Consequently fixed ordered-mesh exclusion and the route-native base
through $B(U_o,U_o)$ suffice for the existing headline, without a separate
decomposition premise.
\Sverified{Mettapedia.GraphTheory.FourColor.SphericalOrderedMeshReduction.combinatorialFourColorStatement_of_fixedOrderedMeshExclusion_and_base}
The direct consumer needs only \texttt{TaitBaseVerifiedAt}$(U_o(a,b))$.
\Sverified{Mettapedia.GraphTheory.FourColor.SphericalOrderedMeshReduction.combinatorialFourColorStatement_of_fixedOrderedMeshExclusion_and_direct_base}
Neither fixed exclusion nor that actual base audit is established here.

\paragraph{Gate and scope.}
Before formalization the ordered-contour gate replayed all eight original
GP$(k,0)$ ambient receipts, checked increasing first hits and extracted
simple ordered submeshes.  Independent order controls exhausted $34979$
permutation families in one through three coordinates and $34$ seeded
cases beyond the strict threshold.  Of $183$ checked depth traces,
$179$ revisit a depth.  Negative controls reject arbitrary branch choices
and weakening the strict threshold.  The existing full-cut/order and
nesting receipts are unchanged.  These controls do not assert zero Count
and enumerate no configuration colouring.

This closes the exact ordered supplier, not the general planar grid
theorem and not fixed ordered-mesh exclusion.  The full boundary/frozen
target of Proposition~\ref{prop:target-forced-cut-chain}, optimized
physical closure and the actual generating base remain open.  The fixed
ledger therefore remains $57/100$; the completed implication is recorded
without assigning partial credit to a larger unresolved row.

\subsection{Marked material in the complete cotree chain}
\label{sec:marked-cotree-material}

The designated-boundary clause of
Proposition~\ref{prop:target-forced-cut-chain} requires control of the
actual sides used by Count.  Avoidance by preliminary cotree vertex shores
does not suffice: conversion to the majority side can add vertices.
The following construction instead marks those majority sides directly.

\begin{lemma}[a constant marked window]
\Sverified{Mettapedia.GraphTheory.FourColor.FiniteMarkedChain.exists_constant_window}
A monotone sequence of $(t+1)n+1$ vertex sides, with $t$ designated
vertices, contains $n+1$ consecutive sides on which membership of every
designated vertex is constant.
\end{lemma}

Each designated vertex can enter only once.  The proof samples the
$t+2$ endpoints of $t+1$ disjoint length-$n$ intervals.  Two samples
have equal marked-intersection cardinality; monotonicity gives equality
through an intervening interval.  Thus the loss is linear in $t$, not
the number $2^t$ of all subsets.  No colouring states are enumerated.

The cotree supplier itself now uses only bridgeless spherical cubic
geometry and two-sided face orbits, before imposing non-colourability
or vertex minimality.  Bounded full-dual radius $h$ and sufficiently
large vertex count construct the cotree, its long path and the connected
edge-shore nodes; no path or chain is an input.
\Sverified{Mettapedia.GraphTheory.FourColor.SphericalCotreePathChain.exists_nested_nodes_of_large_card_of_spherical}
The old target-class API remains a specialization.

\begin{theorem}[marked-free strict material from spherical geometry]
\Sverified{Mettapedia.GraphTheory.FourColor.SphericalMarkedCotreeChain.exists_marked_nodes_of_large_card}
Put $w=2h+1$, $r=(t+1)n$ and
\[
 U_{\mathrm{mark}}(h,t,n)=4^{(6w+1)r+1+2w}.
\]
Let a bridgeless spherical cubic map with two-sided face orbits have
all full-dual distances from a chosen face at most $h$, and more than
$U_{\mathrm{mark}}(h,t,n)$ vertices.  For any $t$ marked vertices it
supplies $n+1$ strictly nested connected edge-shore nodes, each of
middle width at most $w$.  Marked membership in their actual majority
sides is constant.  Between every two selected nodes there is an
unmarked vertex lying strictly between their majority sides, with all
three incident edges in the edge-shore difference.
\end{theorem}

Select every $(6w+1)$st node before applying the marked-window lemma.
Strict edge-shore inclusion forces cardinality growth at least the
index gap.  The cubic middle-set bound then places a full vertex star
inside each selected slab, which forces strict majority-side growth.
\Sverified{Mettapedia.GraphTheory.FourColor.SphericalMarkedCotreeChain.exists_star_of_index_gap}
Only indices change: the original complete cotree cuts and their
boundary orders are not altered.  In the least-counterexample class,
the existing node completion supplies the literal Count nodes.  This
does not itself prove the fixed-boundary target statement.

\paragraph{Frozen edges: persistence is not avoidance.}
If both endpoints of every frozen edge are marked, their crossing set
is identical at all selected majority-side cuts.
\Sverified{Mettapedia.GraphTheory.FourColor.SphericalMarkedCotreeChain.frozen_crossing_invariant}
That set need not be empty.  A checked octagonal-prism control takes
the four prefix strips of $2,3,4,5$ rungs.  Both shores remain connected
and cyclic, each complete cut has width $4$, and the frozen edge joining
positions $0$ and $7$ crosses all four cuts although neither endpoint
lies in an intervening slab.  The control also checks the spherical
rotation, inverse opposite boundary returns and equality with the
majority sides.  It is a reproducible finite geometric test, not a
Lean-certified spherical counterexample or a zero-Count instance.

The gate exhausts $351144$ marked-entry schedules and $20$ sharpness
controls, and replays complete cuts on $25$ spherical maps, including
nonempty marked sets in the strict-star tests.  Its receipt is
\texttt{results/fourcolor/v24\_marked\_cotree\_gate.json}.
It shows why a frozen-data splice must either exclude persistent
crossings or preserve their identities, colours and boundary order;
none of these follows from vertex avoidance alone.  Designated holes,
the remaining frozen-data transport, common seam typing and the
source-facing global supply remain open.  No whole ledger row is
closed by this partial construction, so the fixed score stays $57/100$.

\subsection{Physical Count across nested cuts with persistent wires}
\label{sec:nested-side-wires}

The persistent crossings in Section~\ref{sec:marked-cotree-material}
expose a representation restriction, not a failure of Count composition.
The vertex-attached slab interface requires every inner boundary edge
to enter a vertex of the slab.  In particular, its enlarged side contains
the opposite endpoint of each inner boundary dart.
\Sverified{Mettapedia.GraphTheory.FourColor.TubeSlab.NestedSideWire.attachedSlab_enters}
An edge crossing both nested cuts does not have such an endpoint in the
intervening material.  The following construction works directly with
the ambient vertex sides and leaves the vertex-attached API unchanged.

Let $A\subseteq B$ be vertex sides of a finite rotation system, and put
$S=B\setminus A$.  A colouring of $B$ restricts to colourings of $A$
and $S$ agreeing on every edge joining the two pieces.  Conversely these
restrictions glue uniquely, using the original ambient darts.  Properness,
nonzero colours and internal-edge equality are proved in both directions.
\Sverified{Mettapedia.GraphTheory.FourColor.TubeSlab.NestedSideWire.coloringEquiv}

For input and output words $x$ and $y$, a layer colouring must satisfy
three boundary conditions: it reads $x$ at edges entering the material
from $A$, it reads $y$ at edges leaving the material toward the complement
of $B$, and every persistent crossing copies the same \emph{nonzero}
colour from $x$ to $y$.  A persistent crossing has no vertex in $S$,
so nonzeroness must be checked explicitly rather than inferred from a
layer vertex.  There is no new edge identity or freely chosen permutation
at such a wire.

\begin{theorem}[exact nested-side Count composition]
\Sverified{Mettapedia.GraphTheory.FourColor.TubeSlab.NestedSideWire.wordColoringEquiv}
\Sverified{Mettapedia.GraphTheory.FourColor.TubeSlab.NestedSideWire.sideCount_eq_sum}
Write $C_A(x)$ for the number of proper colourings of $A$ with word $x$,
and $M_{A,B}(x,y)$ for the number of proper colourings of $S$ satisfying
the material and nonzero-wire boundary conditions above.  Then
\[
 C_B(y)=\sum_x C_A(x)M_{A,B}(x,y).
\]
The intermediate word is uniquely determined by the inner colouring,
so this is a bijection of finite colouring fibres, not merely an
equivalence of nonemptiness assertions.
\end{theorem}

The matrix's support is exactly the physical layer relation, and the
side count's support is exactly the existing side support.
\Sverified{Mettapedia.GraphTheory.FourColor.TubeSlab.NestedSideWire.slabCount_ne_zero_iff}
\Sverified{Mettapedia.GraphTheory.FourColor.TubeSlab.NestedSideWire.sideCount_ne_zero_iff}
For $A=B$, the material is empty and the transfer is the identity on
nonzero Tait words, with zero entries for any word containing zero.
\Sverified{Mettapedia.GraphTheory.FourColor.TubeSlab.NestedSideWire.slabCount_refl}
A mismatched wire gives count zero; a zero input colour also gives
count zero, whether the port enters material or traverses a wire.
\Sverified{Mettapedia.GraphTheory.FourColor.TubeSlab.NestedSideWire.slabCount_eq_zero_of_wire_mismatch}
\Sverified{Mettapedia.GraphTheory.FourColor.TubeSlab.NestedSideWire.slabCount_eq_zero_of_zero_input}

\paragraph{Frozen data and the existing consumer.}
Prescriptions are placed on original ambient darts.  If none of their
bases lies in $S$, restriction and gluing induce a bijection preserving
all prescribed colours.
\Sverified{Mettapedia.GraphTheory.FourColor.TubeSlab.NestedSideWire.fixedColoringEquiv}
Marking the endpoints of frozen edges and using the constant majority-side
window supplies precisely this avoidance of prescribed dart bases.
\Sverified{Mettapedia.GraphTheory.FourColor.SphericalMarkedCotreeChain.frozen_restriction_iff}
The support-composition theorem transports to the stack's existing
\texttt{innerSupport} under explicit input and output coordinates.
\Sverified{Mettapedia.GraphTheory.FourColor.TubeSlab.NestedSideWire.mem_innerSupport_iff}
The coordinate of a through wire is determined by its ambient dart,
and its colour is copied to the corresponding output coordinate.
\Sverified{Mettapedia.GraphTheory.FourColor.TubeSlab.NestedSideWire.wireIndex_dart}
\Sverified{Mettapedia.GraphTheory.FourColor.TubeSlab.NestedSideWire.coordinateRel_wire}

The diagnostic receipt
\texttt{results/fourcolor/v24\_nested\_side\_wire\_gate.json}
checks all ten nested pairs, including equal pairs, of the four prism
cuts.  It tests $41020$ dart assignments with $60$ known proper controls,
and rejects $168$ mismatched-wire and $168$ equal-but-zero-wire mutations.
It also rejects $36$ internal-seam mutations that leave both pieces
individually properly coloured but make their gluing improper.
This is a test of generic composition, not a configuration support table.
The Lean results quantify over arbitrary finite rotation systems and
nested vertex sides, including disconnected or empty material.

\paragraph{Remaining geometric obligation.}
These theorems establish exact colouring semantics and frozen-dart
preservation for decomposition and reassembly.  They do not establish
that deleting a repeated block preserves designated holes, cyclic seam
order, spherical admissibility or the target minimality class.  Coordinate
transport is not a proof that an arbitrary coordinate permutation is a
planar splice.  The source-facing global supply and those replacement
obligations remain open; no whole fixed ledger row is discharged, and
the score remains $57/100$.

\subsection{Exact cyclic order of persistent crossing wires}
\label{sec:persistent-wire-order}

The order compatibility needed by the preceding wire semantics is now
proved for actual nested spherical vertex sides.  Let $A\subseteq B$,
with $A$, $B$ and the complement of $B$ inducing connected graphs in
a connected spherical cubic map with cyclic vertex rotations and two-sided
face orbits.  Let $W$
be the original outward darts based in $A$ whose opposite endpoints lie
outside $B$.  Restrict each complete capped boundary successor to $W$.
The resulting permutations are equal under literal ambient dart identity.
No new port permutation, cyclic shift or reversal is chosen.
\Sverified{Mettapedia.GraphTheory.FourColor.SphericalNestedBoundaryOrder.commonWire_boundaryOrder}
The result is uniform in the boundary sizes and the number of surviving
wires; the material between $A$ and $B$ need not be connected, and the
complement of $A$ is not an additional hypothesis.

\paragraph{Euler equality supplies order preservation.}
First return to a marked subset is independent of an intermediate
restriction.  A transposition cannot reorder marks that lie on one cycle
both before and after surgery.
\Sverified{Mettapedia.GraphTheory.FourColor.PermutationFirstReturnRestriction.nextHit_nested}
\Sverified{Mettapedia.GraphTheory.FourColor.PermutationFirstReturnRestriction.nextHit_swap_mul_of_one_cycle}
For an edge presentation, genus-zero Euler equality is hereditary under
edge deletion.  Adding an edge within an existing component must split
a face; a face merge there would create genus that later additions cannot
remove.  Consequently marked darts of one original component which lie
on one final face retain their first-return order throughout the extension.
Intermediate one-face hypotheses are derived, not supplied.
\Sverified{Mettapedia.GraphTheory.FourColor.PlanarMapBoundaryRestriction.plane_tail_and_split}
\Sverified{Mettapedia.GraphTheory.FourColor.PlanarMapBoundaryRestriction.nextHit_face_extension}

The actual internal-edge involutions of $A$ and $B$ are presented in
that order by restricting the ambient flip.  The complementary parts have
exactly the original edge count, and ambient sphericity supplies the Euler
equality used by the deletion theorem.
\Sverified{Mettapedia.GraphTheory.FourColor.InvolutionRestrictionEdgeList.length_restrict_add_compl}
\Sverified{Mettapedia.GraphTheory.FourColor.SphericalNestedBoundaryOrder.exists_capped_edge_extension}
Existing planar-bond closure gives one capped face for the outer boundary;
the retained-dart inclusion identifies its first return with the original
boundary API.
\Sverified{Mettapedia.GraphTheory.FourColor.SphericalNestedBoundaryOrder.retained_boundary_sameCycle}
\Sverified{Mettapedia.GraphTheory.FourColor.SphericalNestedBoundaryOrder.boundaryOrderOn_eq}
For the supplied cotree nodes, shore connectivity and cubicity discharge
the required majority-side connectivity.  Thus the wire-order result
applies to the same nodes as the frozen-colour restriction theorem.
\Sverified{Mettapedia.GraphTheory.FourColor.SphericalMarkedCotreeChain.node_wire_boundaryOrder}

\paragraph{Diagnostic and exact remaining boundary.}
The reproducible geometry gate checks all cyclic connected/coconnected
sides of small prisms and all admissible middle-band choices between two
fixed caps in a $24$-vertex annulus.  It rejects $99725$ deliberate
reversals of interfaces with at least three common wires.  Independent
permutation tests check $12566$ surgeries, $2127$ nested restrictions and
$101473$ planar extensions; omitting planarity gives $7428$ failing
controls.  The receipt is
\texttt{results/fourcolor/v24\_persistent\_wire\_order\_gate.json}.
These are geometric tests, not colouring catalogues or target
counterexamples.  The formal theorem proves the universal order statement.

\paragraph{Equal order is not a full matching fixing the wires.}
A stronger proposed implication fails: even if the two complete boundaries
have the same size, their cyclic bijection need not fix every persistent
wire.  In the hexagonal prism, take vertex sides
$A=\{0,1,2,6,7\}$ and $B=\{0,1,2,5,6,7,8\}$ in the receipt's
numbering.  Both shores and their complements are connected and cyclic;
both cuts have width $5$.  Relative to one common wire, the other is at
cyclic position $3$ on the inner boundary but position $2$ on the outer.
A full cyclic matching fixing the first wire is unique and cannot fix
the second.  The normalized five-position obstruction is kernel-checked.
\Sverified{Mettapedia.GraphTheory.FourColor.PersistentWireMatchingObstruction.no_full_matching}
The finite geometry gate finds $24$ failures among $156$ eligible nested
pairs in that prism, and none among the $10$ pentagonal-prism controls.
Its receipt is
\texttt{results/fourcolor/v24\_persistent\_wire\_matching\_gate.json}.
The prism's geometry is checked by the executable gate, not asserted to
be a Lean-certified target counterexample.

Thus the order theorem alone does not construct a full seam fixing frozen
wires: their positions, or equivalently the intervening boundary gaps,
must also be matched.  The following refinement supplies those positions
and constructs the full named seam.  Designated-hole and target-specific
prescription transport, and the source-facing global supply, remain open.
The fixed ledger is unchanged at $57/100$.

\subsection{Marked positions supply the full physical seam}
\label{sec:marked-seam-positions}

The existing normalized state records a hub permutation and a set of
realizable Tait boundary words, but no names for persistent edges.
For a fixed family $M$ of $t$ designated ambient darts, augment the state
at width $j$ by a function
\[
 M\longrightarrow \{\bot\}\sqcup\operatorname{Port}_j.
\]
It records the complete-boundary coordinate of each named dart, or its
absence.  This is not merely its relative rank among the surviving wires.
The exact fixed-width cardinality is
$j!\,2^{3^j}(j+1)^t$.  Including all widths $j\leq k$ gives the bound
\[
 Q(k,t)=\left(\sum_{j=0}^{k} j!\,2^{3^j}\right)(k+1)^t.
\]
\Sverified{Mettapedia.GraphTheory.FourColor.MarkedSeamState.card_marked}
\Sverified{Mettapedia.GraphTheory.FourColor.MarkedSeamState.card_bounded_le}
The cost depends on the number of named darts, not the number of all
ambient edges.  Duplicate labels are allowed and no colour table is evaluated.

Equality of these records gives the same literal boundary width, the old
exact state equality, and identity of each named dart under the coordinate
change.  Membership on the replacement boundary is derived as well.
\Sverified{Mettapedia.GraphTheory.FourColor.MarkedSeamState.boundedOfShore_eq_elim}
\Sverified{Mettapedia.GraphTheory.FourColor.MarkedSeamState.matching_fixes_mark}
For the existing physical replacement matching, both orientations retain
the ambient edge partner; the matching constructor is not replaced by an
abstract existence hypothesis.
\Sverified{Mettapedia.GraphTheory.FourColor.MarkedSeamState.replacementMatching_mark}
\Sverified{Mettapedia.GraphTheory.FourColor.MarkedSeamState.replacementMatching_symm_mark}

\paragraph{Geometric supply and its actual consumer.}
The connected-shore node constructor requires only the ambient connected,
bridgeless spherical cubic data, not minimality or non-colourability.
\Sverified{Mettapedia.GraphTheory.FourColor.GoertzelV24ConnectedShoreLiteralNode.ConnectedShoreNode.toLiteralOfSpherical}
Let every dual face be at distance at most $h$ from a named root face,
let $k=2h+1$, and let $s$ vertices be marked.  If
\[
 |V|>4^{(6k+1)(s+1)Q(k,t)+1+2k},
\]
the actual marked cotree supplier gives $Q(k,t)+1$ literal nodes.
Finite-state repetition produces two strictly nested shores with equal
marked states, a strict cubic star between them, and no marked vertex
in the removed material.  Neither a chain nor a seam is supplied as a
premise of this theorem.
\Sverified{Mettapedia.GraphTheory.FourColor.MarkedCotreeSeam.exists_seamPair_of_large_card}
If named darts are based at marked vertices and the family contains both
orientations of each named edge, all those darts remain present and their
edge involution is exactly the original one.
\Sverified{Mettapedia.GraphTheory.FourColor.MarkedCotreeSeam.SeamPair.mark_survives}
\Sverified{Mettapedia.GraphTheory.FourColor.MarkedCotreeSeam.SeamPair.candidate_alpha_mark}
With zero ambient Count, the existing physical replacement theorem supplies
a strictly smaller connected, bridgeless spherical cubic zero-Count map,
using precisely this matching and preserving the named edge pairs.
\Sverified{Mettapedia.GraphTheory.FourColor.MarkedCotreeSeam.exists_marked_replacement_of_large_card}

\paragraph{Diagnostic and remaining boundary.}
The position gate checks $4096$ combinations of actual nested prism cuts,
marked subsets, absent labels and relative cyclic roots.  All $24$ previous
full-matching obstructions are rejected when every common wire is marked;
empty and singleton mark controls also pass.  The receipt is
\texttt{results/fourcolor/v24\_marked\_seam\_position\_gate.json}.
This is a generic geometric test, not configuration enumeration.

The theorem assumes bounded \emph{complete dual radius}; it does not
derive that condition from the source's wall-free hypotheses.  Preserving
named edge pairs also does not by itself prove preservation of designated
face orbits or the joint support under arbitrary prescribed frozen colours.
The face-orbit transport is established in Section~\ref{sec:marked-face-reduction};
joint frozen-colour transport is established in Section~\ref{sec:frozen-joint-reduction}.
The global D1b/D1c implications remain open. The next construction removes the separately supplied radius bound
under ordered-mesh exclusion for the closed-map Count system. No whole
fixed ledger row is newly discharged: $57/100$.

\subsection{Ordered mesh or an actual smaller Count instance}
\label{sec:meshfree-full-seam}

The complete contour and cotree constructions can be combined before
minimality is invoked. This matters because deriving a radius bound from
minimality and then invoking a radius-dependent supplier does not exhibit
a reduction on arbitrary large zero-Count instances.

Let $a\geq2$, fix an ordered-mesh width $b$, and let $s$ vertices and $t$
ambient darts be designated. Using $Q$ from
Section~\ref{sec:marked-seam-positions}, define
\[
 c=3b^{2^a},\qquad
 N=(6c+1)(s+1)Q(c,t)+1,\qquad H=aN,\qquad k=2H+1,
\]
and the explicit size threshold
\[
 V_{\mathrm{seam}}(a,b,s,t)
   =4^{(6k+1)(s+1)Q(k,t)+1+2k}.
\]
Every connected bridgeless spherical cubic graph with two-sided face
orbits and more than this many vertices supplies either an actual
$a$-by-$b$ ordered injective mesh, or a full marked seam of width at most
$\max(c,k)$. The latter includes two connected and coconnected literal
shores, their common exact state and boundary width, and a strict unmarked
cubic star between them. No radius, chain, or matching is an input.
\Sverified{Mettapedia.GraphTheory.FourColor.SphericalMeshFreeSeam.exists_orderedMesh_or_seamPair}

\paragraph{Construction, not a minimality contradiction.}
Choose a root face. If some face has distance greater than $H$, use the
mesh-or-complete-bond alternative in $N$ disjoint windows of $a$ depth
levels. In the absence of the ordered mesh, the resulting full bonds
give strictly nested connected-shore nodes of width $c$.
The contour-node supplier requires only the ambient spherical class;
its former minimal-counterexample interface is retained as a specialization.
\Sverified{Mettapedia.GraphTheory.FourColor.SphericalOrderedMeshReduction.exists_nested_nodes_of_no_orderedMesh_of_spherical}
Spacing by $6c+1$ gives strict cubic material, and a window avoiding all
marked transitions leaves $Q(c,t)+1$ nodes. The same full marked-state
pigeonhole used for cotree nodes supplies the seam.
\Sverified{Mettapedia.GraphTheory.FourColor.SphericalMarkedCotreeChain.exists_marked_nodes_of_nested_nodes}
\Sverified{Mettapedia.GraphTheory.FourColor.MarkedCotreeSeam.exists_seamPair_of_marked_nodes}

Otherwise every face has distance at most $H$ from the root. The existing
marked cotree theorem then applies directly at width $k$ using the displayed
size threshold. Thus mesh exclusion alone selects one of these two
constructions; it need not imply that the graph has small dual radius.
\Sverified{Mettapedia.GraphTheory.FourColor.SphericalMeshFreeSeam.exists_seamPair_of_no_orderedMesh}

If ambient Count is zero, the produced seam feeds the existing physical
replacement theorem. For designated darts based at marked vertices and
closed under edge reversal, this returns a strictly smaller connected,
bridgeless spherical cubic zero-Count map, with those original edge
partners still present.
\Sverified{Mettapedia.GraphTheory.FourColor.SphericalMeshFreeSeam.exists_marked_replacement_of_no_orderedMesh}
Equivalently, every sufficiently large zero-Count map has an ordered mesh
or an actual smaller zero-Count map. The mesh is an output alternative,
not an assumed external witness or a configuration catalogue.
\Sverified{Mettapedia.GraphTheory.FourColor.SphericalMeshFreeSeam.exists_orderedMesh_or_marked_replacement}

\paragraph{Verification and scope.}
The complete-contour nesting and marked-cotree geometry receipts were
replayed without modification. The branch gate additionally checks $3500$
radius-boundary cases and $1750$ spaced windows, including equality and
both adjacent sides of the radius threshold. Its receipt is
\texttt{results/fourcolor/v24\_meshfree\_seam\_branch\_gate.json}.
These are generic geometric and arithmetic controls, not zero-Count
examples or colouring enumerations; the universal claim is the Lean theorem.

The closed-map direct Count consumer has no designated-hole or frozen-colour
input. This result therefore does not assert the corresponding transport
for the more general boundary-conditioned species in
Proposition~\ref{prop:target-forced-cut-chain}. Nor does it exclude the
ordered mesh, identify its absence with every informal wall-free premise,
or verify the generating base. The former smaller target-class bounds are
not replaced by this coarser constructive bound. Fixed-wall exclusion and
the full D1b/D1c obligations remain open; the ledger stays at $57/100$.

\subsection{Designated faces survive the supplied physical reduction}
\label{sec:marked-face-reduction}

Preserving a named edge partner is not, by itself, a theorem about faces:
one must also preserve vertex rotation and show that an entire facial orbit
survives, without merging into another orbit. The actual sewn map now has
these properties on designated faces. Its ambient dart map is injective,
because the retained inner shore and the exterior of the outer shore are
disjoint. It intertwines vertex rotation on every surviving dart, independently
of which seam matching is used.
\Sverified{Mettapedia.GraphTheory.FourColor.MarkedCotreeSeam.SeamPair.ambientDart_injective}
\Sverified{Mettapedia.GraphTheory.FourColor.MarkedCotreeSeam.SeamPair.ambientDart_rho}

Let a designated facial orbit be covered by the named darts, with the named
family closed under edge reversal. The existing edge-partner theorem then
intertwines the facial successor $\phi=\rho\alpha$ on that orbit. Induction
intertwines every natural power of $\phi$; finiteness and injectivity give
both directions of the same-face relation. Consequently the image of the
entire candidate face is exactly the original face, with the same length
and rooted cyclic sequence. This does not require the whole marked family
to be closed under vertex rotation, a substantially stronger condition.
\Sverified{Mettapedia.GraphTheory.FourColor.MarkedCotreeSeam.SeamPair.sameCycle_iff}
\Sverified{Mettapedia.GraphTheory.FourColor.MarkedCotreeSeam.SeamPair.designated_face_survives}

\paragraph{Marks are constructed, not assumed.}
Given a finite set of ambient face roots, take the union of their complete
facial dart sets, add all reversed darts, and mark the base vertices of
the resulting set. These operations discharge every endpoint and
edge-reversal premise. All vertex visits of the designated faces avoid
the removed material; persistent seam crossings are allowed.
\Sverified{Mettapedia.GraphTheory.FourColor.SphericalMarkedFaceReduction.face_avoids_material}
Applying the previous mesh-or-seam supplier with these computed marks gives
an ordered mesh or an actual strictly smaller connected, bridgeless,
spherical cubic zero-Count map preserving every designated face. No radius,
chain, matching, endpoint marking or facial transport is supplied as a
hypothesis. The bound is exactly $V_{\mathrm{seam}}(a,b,s,t)$ above, for
the cardinalities $s,t$ of these computed vertex and dart sets.
\Sverified{Mettapedia.GraphTheory.FourColor.SphericalMarkedFaceReduction.exists_orderedMesh_or_replacement}

\paragraph{Gate and remaining scope.}
On $35$ nested full-cut pairs in prisms, the edge-fixing matching preserves
the marked face crossing the seam; $105$ deliberately shifted cyclic
matchings change it. Both shores and complements are connected and cyclic,
the cuts have width four, and the removed material contains a strict
cubic star. These are colourable geometric controls, not instances of the
minimal-counterexample class. The receipt is
\texttt{results/fourcolor/v24\_marked\_face\_splice\_gate.json}.

The threshold depends on the total designated perimeter, not merely the
number of holes. In particular this does not bound an arbitrarily long
designated face by treating its perimeter as a fixed parameter. Ordinary
closed-map zero Count is preserved; joint supports under prescribed frozen
colours are handled by the refinement below. Fixed-mesh exclusion, the complete
source-facing D1b/D1c claims and the generating base remain open.
The fixed ledger stays at $57/100$.

\subsection{A supplied reduction preserving joint frozen-colour support}
\label{sec:frozen-joint-reduction}

Preserving the marginal set of seam words is insufficient for a prescribed
colour problem: the colour of an observed dart may be correlated with that
word. For a family $M$ of named darts, each piece now records the joint set
of pairs $(x,f)$, where $x$ is its seam word and $f:M\to\mathrm{Color}$ is
a prescription realized by the same proper colouring. A name absent from
a piece imposes no constraint on that piece. Gluing requires a common $x$
and the same $f$ on both sides. The existing colouring bijection proves this
identity while retaining the witnesses and their observations.
\Sverified{Mettapedia.GraphTheory.FourColor.FrozenCountSupport.closedAccepts_iff}

The semantic-to-physical bridge reads these colours on the actual computed
edge orbits of the sewn rotation system. Exact dart transport strengthens
the old colourability-only reassembly theorem by a pointwise colour equation.
For complementary shores, the resulting prescription is exactly the one
on the original ambient dart names. Reindexing a boundary changes word
coordinates but does not change the observed input darts.
\Sverified{Mettapedia.GraphTheory.FourColor.FrozenPhysicalCount.physicalAccepts_iff}
\Sverified{Mettapedia.GraphTheory.FourColor.FrozenVertexSideCount.physicalAccepts_reassembly_iff}
\Sverified{Mettapedia.GraphTheory.FourColor.FrozenJointReindex.joint_eq_of_coordinate_eq}

\paragraph{The equality is supplied, not assumed.}
At width $j$, append the joint set to the full marked seam state. With
$t=|M|$, its additional carrier has cardinality $2^{4^{j+t}}$; zero colours
are included in this coarse carrier and rejected by proper colourings.
Thus a bound for the varying-width refined state is
\[
 Q_f(k,t)=Q(k,t)\,2^{4^{k+t}}.
\]
\Sverified{Mettapedia.GraphTheory.FourColor.FrozenSeamState.card_bounded_le}
Use $Q_f$ in place of $Q$ throughout the contour/cotree construction of
Section~\ref{sec:meshfree-full-seam}: in $N$, in $H=aN$, and in the final
cotree threshold. The larger geometric chain supplies two equal refined
states. The old components determine the same orientation-compatible,
ambient-name-fixing matching; the new component supplies joint-support
equality in precisely its coordinates. Every named dart actually survives,
so the prescription cannot be preserved vacuously by removing its names.
\Sverified{Mettapedia.GraphTheory.FourColor.SphericalFrozenReduction.exists_joint_seam_of_no_orderedMesh}
\Sverified{Mettapedia.GraphTheory.FourColor.MarkedCotreeSeam.SeamPair.accepts_iff}

The resulting alternative is an actual ordered mesh or a strictly smaller
connected, bridgeless spherical cubic map accepting exactly the same
prescriptions, for every $f$ simultaneously. Neither ordinary zero Count
nor minimality is a construction premise. Structural closure and strict
decrease are proved independently of the acceptance predicate. In
particular, rejection of a fixed prescription transfers even if the ambient
map is ordinarily colourable.
\Sverified{Mettapedia.GraphTheory.FourColor.SphericalFrozenReduction.exists_orderedMesh_or_frozen_replacement}

\paragraph{Faces and prescriptions on one candidate.}
Take all darts of the designated faces together with the supplied frozen
darts, close under edge reversal, and mark their base vertices. Applying
the refined supplier to these computed names yields one candidate retaining
every designated face's complete rooted cyclic orbit and each named edge
partner, while preserving all prescriptions on the named family.
\Sverified{Mettapedia.GraphTheory.FourColor.SphericalFrozenFaceReduction.exists_orderedMesh_or_replacement}

\paragraph{Diagnostic and remaining scope.}
The preregistered control glues two three-port cubic vertices to the planar
theta multigraph. Its six proper words give $144$ prescription controls;
all $48$ equal-joint-support controls pass. Forgetting the observation
distinguishes acceptance in $36$ cases, despite identical marginal seam
support. This is a generic gluing control, not a configuration certificate
or a target counterexample. Receipt:
\texttt{results/fourcolor/v24\_frozen\_joint\_support\_gate.json}.

The universal results are the Lean theorems, not extrapolations from that
control. They preserve acceptance, not equality of numerical colouring
multiplicities after deletion. The bound depends on the total designated
perimeter and the number of frozen names. It supplies neither a uniform
long-face bound nor fixed-mesh exclusion, nor the complete saturated-chain
claims of Proposition~\ref{prop:target-forced-cut-chain}. Those D1b/D1c
obligations and the generating base remain open. The ledger stays at
$57/100$.

\subsection{Ordered mesh corners do not automatically bound cells}
\label{sec:ordered-mesh-cell-obstruction}

The ordered-mesh alternative is not yet a decomposition into disk cells.
One attempted construction takes adjacent rows and columns, cuts their
four branch-to-branch subpaths, and concatenates them into a rectangular
rim. The existing ordered-mesh definition deliberately permits shared
row--column segments. Thus strict branch order and simple paths do not
make this rim simple, even when each family is vertex-disjoint.

\paragraph{An exact obstruction.}
In the dodecahedral control, take the trimmed rows
$(10,14,12)$ and $(17,19)$, and columns $(10,11,17)$ and $(12,14,19)$.
The designated branch matrix is
\[
 \begin{pmatrix}10&12\\17&19\end{pmatrix}.
\]
These paths inhabit the actual \texttt{OrderedInjectiveMesh} carrier;
within-family vertex-disjointness and ambient cubic degree are checked.
\Sverified{Mettapedia.GraphTheory.FourColor.OrderedMeshCellBoundaryCounterexample.ordered_cell_obstruction}
\Sverified{Mettapedia.GraphTheory.FourColor.OrderedMeshCellBoundaryCounterexample.families_vertex_disjoint}
\Sverified{Mettapedia.GraphTheory.FourColor.OrderedMeshCellBoundaryCounterexample.ambient_cubic}
The four sides concatenate to
$(10,14,12,14,19,17,11,10)$. Corner $12$ has only neighbour $14$ in
their entire edge union. No simple cycle in that union contains $12$,
regardless of where the cycle is rooted. In particular, loop erasure
cannot retain all four designated corners.
\Sverified{Mettapedia.GraphTheory.FourColor.OrderedMeshCellBoundaryCounterexample.corner_not_on_cycle}
Erasing the backtrack does leave the simple cycle
$(10,14,19,17,11,10)$, but it loses corner $12$.
\Sverified{Mettapedia.GraphTheory.FourColor.OrderedMeshCellBoundaryCounterexample.erased_is_cycle}

\paragraph{Complete-interface diagnostic.}
The union contains every ambient edge induced on its vertex set. The
full ambient edge boundary has size six, but corner $12$ has just one
internal incident edge, violating the inside-majority condition needed
for this vertex set to be a good side.
\Sverified{Mettapedia.GraphTheory.FourColor.OrderedMeshCellBoundaryCounterexample.cellEdges_induced}
\Sverified{Mettapedia.GraphTheory.FourColor.OrderedMeshCellBoundaryCounterexample.complete_boundary_card}
\Sverified{Mettapedia.GraphTheory.FourColor.OrderedMeshCellBoundaryCounterexample.corner_internal_edge_count}
The lab separately checks that both induced sides are connected and
contain cycles. Thus connectivity alone does not repair this candidate.

The preregistered gate replays the eight full spherical $GP(k,0)$ controls,
$1\le k\le8$, with their previously selected ordered meshes. Of their
$316$ adjacent cells, $252$ have non-simple candidate rims and a pendant
designated corner. A pentagonal facial-cycle control passes. The receipt
\texttt{results/fourcolor/v24\_ordered\_mesh\_cell\_gate.json} includes
the complete ambient graph and face data for the first obstruction.
Sphericity of that receipt is checked by the lab; the Lean statements above
check the finite ordered-mesh and cycle obstruction, not an additional
rotation-system theorem.

\paragraph{What fails, and what remains.}
This refutes the corner-preserving four-subpath cell construction, not
ordered-mesh extraction or fixed-mesh exclusion in the counterexample
class. A repair must allow branch relocation or another boundary
construction, then prove global consistency, full lateral-interface
coverage and the required saturated nested sides. Independent loop
erasure does not supply those conclusions. No non-colourability premise
or colouring enumeration was used. D1b/D1c and fixed-mesh exclusion remain
open; the fixed ledger stays at $57/100$.

\subsection{Face-region hulls repair the local mesh-cell construction}
\label{sec:face-region-hulls}

The obstruction above motivates a different construction, allowing the
designated corners to disappear. Retain the full contour cycles and
linkage paths as an ambient edge barrier $B$. Delete dual adjacencies
across those edges and take an actual connected component $R$ of the
remaining face graph. This construction is global across neighbouring
cells, rather than an independent loop erasure of each four-path rim.
Its connectivity in the full facial dual is proved; every original edge
between $R$ and a face outside $R$ belongs to $B$.
\Sverified{Mettapedia.GraphTheory.FourColor.FaceRegionHull.cell_dual_connected}
\Sverified{Mettapedia.GraphTheory.FourColor.FaceRegionHull.cell_crossing_in_barrier}

\paragraph{From faces to a complete vertex bond.}
Work in a connected cubic map with two-sided facial edges. Let $S(R)$
contain every vertex incident with a face of $R$. Facial cycles
and the shared endpoints of adjacent faces construct connectedness of
$G[S(R)]$ from connectedness of the face region. Its vertex complement
need not be connected, even when the complementary face region is.
Fix $b\notin S(R)$ and set
\[
 H_b(R)=V\setminus\operatorname{Comp}_{G[V\setminus S(R)]}(b).
\]
Thus every complementary pocket except the component containing $b$ is
filled. Both induced sides of $H_b(R)$ are connected. Every edge of its
complete ambient cut already crosses $S(R)$; filling introduces no new
cut edge. The operation is monotone and idempotent as a vertex-set hull.
\Sverified{Mettapedia.GraphTheory.FourColor.FaceRegionHull.closed_connected}
\Sverified{Mettapedia.GraphTheory.FourColor.FaceRegionHull.hull_connected_sides}
\Sverified{Mettapedia.GraphTheory.FourColor.FaceRegionHull.hull_boundary_subset}
\Sverified{Mettapedia.GraphTheory.FourColor.FaceRegionHull.hull_mono}
\Sverified{Mettapedia.GraphTheory.FourColor.FaceRegionHull.filled_idem}

An outside face disjoint from $S(R)$ and containing $b$ stays wholly in
the retained exterior. Together with an inside face this supplies cycles
on both sides. With sphericity and cyclic vertex rotations, the existing planar-bond theorem
therefore computes unique touched-face occurrences and inverse
complementary boundary-successor orders for the full constructed cut.
The incident-edge adapter constructs an existing exact Count node whose
middle-set width is at most $|\delta S(R)|$. Shore connectivity and the
middle-set bound are outputs, not extra inputs to that constructor.
\Sverified{Mettapedia.GraphTheory.FourColor.FaceRegionHull.hull_hasCycles}
\Sverified{Mettapedia.GraphTheory.FourColor.FaceRegionHull.hull_boundary_order}
\Sverified{Mettapedia.GraphTheory.FourColor.FaceRegionHull.toConnectedNode}
\Sverified{Mettapedia.GraphTheory.FourColor.FaceRegionHull.cellNode}

\paragraph{Strictness and width are still separate tasks.}
For $S\subseteq T$, the filled side grows strictly if and only if $T$
contains a vertex in the old retained exterior. Adding faces or vertices
inside an already filled pocket does not suffice.
\Sverified{Mettapedia.GraphTheory.FourColor.FaceRegionHull.filled_strict_iff}
The diagnostic \texttt{v24\_mesh\_region\_gate.json} replays the eight
existing spherical controls. All $384$ barrier cells have simple full
rims. Of $368$ cumulative band prefixes, $88$ have disconnected vertex
complements before filling, including $38$ proper prefixes whose face
complements are connected. Filling restores two connected cyclic sides
and complete boundary order in all $368$ cases. A whole outside face
containing the chosen anchor survives each time, and strict growth is
checked in this sample. These finite measurements are not universal
theorems about prefix growth.

The largest filled prefix cut at frequency eight still has $80$ edges.
This is a measurement of this construction, not a lower bound on optimal
width. The repair supplies local regions and their valid Count interfaces;
it does not yet force a sufficiently long chain of uniformly bounded
interfaces in the mesh alternative. Fixed-mesh exclusion and the global
D1b/D1c obligations remain open. The fixed ledger remains $57/100$.

\subsection{Reordering filled cells has an intermediate-width cost}
\label{sec:mesh-region-order-obstruction}

The hull construction repairs connectivity, but it does not make the
largest intermediate cut equal to the largest endpoint or individual-cell
cut. Write $A_i$ for the vertex set of cell $i$, fix the same exterior
anchor $b$, and put
\[
 w(s)=\left|\delta\left(V\setminus
   \operatorname{Comp}_{G[V\setminus\bigcup_{i\in s}A_i]}(b)\right)\right|.
\]
The attempted bound was that some order of the cells keeps every filled
prefix at width at most
$\max\{\max_i w(\{i\}),w(\text{all cells})\}$.

\paragraph{An exact minimax obstruction.}
The numbered $GP(3,0)$ control has a band with eight selected face cells.
Its ambient graph has $180$ vertices and $270$ distinct edges, with cubic
degree checked in Lean.
\Sverified{Mettapedia.GraphTheory.FourColor.MeshRegionOrderObstruction.edge_pairs_nodup}
\Sverified{Mettapedia.GraphTheory.FourColor.MeshRegionOrderObstruction.ambient_cubic}
For this band,
\[
 \forall i,\quad w(\{i\})\le15,\qquad
 w(\{0,\ldots,7\})=15,\qquad
 \forall s\ (|s|=6),\quad w(s)\ge16.
\]
The six-cell statement covers all $\binom86=28$ choices, including choices
that are not face-connected. Every permutation therefore has a six-cell
prefix of width greater than fifteen.
\Sverified{Mettapedia.GraphTheory.FourColor.MeshRegionOrderObstruction.local_and_final_width}
\Sverified{Mettapedia.GraphTheory.FourColor.MeshRegionOrderObstruction.six_cells_width}
\Sverified{Mettapedia.GraphTheory.FourColor.MeshRegionOrderObstruction.every_order_exceeds_fifteen}
The explicit local order $(1,2,3,4,5,7,0,6)$ keeps every nonempty prefix
at width at most sixteen. Thus the optimal peak over these cell orders is
exactly sixteen; neither half of this claim relies on trusting an optimizer.
\Sverified{Mettapedia.GraphTheory.FourColor.MeshRegionOrderObstruction.optimal_order_width_at_most_sixteen}

\paragraph{Checking the complete physical hull.}
The encoded raw side is proved to be precisely the union of the selected
cell vertex sets. For each recorded side, the retained exterior contains
$b$, avoids the raw side, is closed under every allowed ambient edge, and
has a rank-decreasing parent edge from each other vertex. A generic
rank-induction theorem identifies it with the entire exterior component.
Its complement is therefore the actual hull, and direct endpoint tests
are proved to give its complete ambient edge boundary.
\Sverified{Mettapedia.GraphTheory.FourColor.MeshRegionOrderObstruction.mem_rawMask_iff}
\Sverified{Mettapedia.GraphTheory.FourColor.RankedRegionComponent.flood_eq_of_rank}
\Sverified{Mettapedia.GraphTheory.FourColor.MeshRegionOrderObstruction.component_exact}
\Sverified{Mettapedia.GraphTheory.FourColor.MeshRegionOrderObstruction.boundary_exact}
The finite checks use flat byte pages and individually checked component
witnesses. Sphericity, face-cell extraction and the geometric admissibility
of the displayed order are checked in the ambient lab receipt; they are
not additional Lean rotation-system claims here.

\paragraph{Scope of the diagnostic.}
The gate replays all eight archived spherical controls and searches the
complete cell-subset lattices of their $52$ bands. Allowed states have
connected face regions, complete connected cyclic vertex shores,
inside-majority and full boundary orders; steps must grow strictly.
The proposed bound fails in $44$ bands. The recurrence agrees with full
permutation searches in $14$ controls, including all $40320$ orders of
the eight-cell specimen. At frequency eight the largest prefix width
across the bands falls from $80$ to $48$ after optimizing these orders;
the largest candidate bound there is $42$.

This refutes the proposed exact bound for this region-selection method.
It does not refute a larger quantitative bound, refinement or regrouping
of cells, arbitrary branch decompositions, or the target-counterexample
implication of Proposition~\ref{prop:target-forced-cut-chain}. No
non-colourability premise or colouring enumeration is involved. The
remaining construction must control complete intermediate interfaces,
not just local and final cuts. D1b/D1c and fixed-mesh exclusion remain
open, and the fixed ledger stays at $57/100$.

\subsection{Marked face returns rather than whole-face perimeter}
\label{sec:marked-face-return-rewire}

The designated-face supplier preserves complete facial cycles, and its
threshold consequently depends on their total perimeter. For a transfer
which needs only named edges to stay cofacial, preserving every intervening
dart is stronger than necessary. The following compositional geometry
provides a finite record of the required observation. The refined geometric
supplier in Subsection~\ref{sec:supplied-named-cofacial-reduction} below
uses this record in its marked-state choice.

\paragraph{Compression commutes with supported rewiring.}
Let $p$ be a permutation of a finite dart carrier, let $H$ be a set of
recorded points, and let $r_H(p)$ be the next-return permutation on $H$.
For a permutation $\tau$ of $H$, write $\bar\tau$ for its extension by
the identity off $H$. With rightmost factors acting first,
\[
 r_H(p\bar\tau)=r_H(p)\tau.
\]
Indeed, starting at $x\in H$, the rewired traversal follows the old arc
from $\tau(x)$ until its first return. There are no intervening recorded
points at which another rewire can occur. Both the exact return time and
the return permutation identity are checked generically.
\Sverified{Mettapedia.GraphTheory.FourColor.MarkedFaceReturnRewire.returnTime_rewire}
\Sverified{Mettapedia.GraphTheory.FourColor.MarkedFaceReturnRewire.nextHitPerm_rewire}
The cycles of this finite product compute precisely the new cofaciality
of the recorded darts. Equal records in common coordinates remain
interchangeable under every common rewire, even for different ambient
carrier sizes; no equality of unmarked arc lengths is needed.
\Sverified{Mettapedia.GraphTheory.FourColor.MarkedFaceReturnRewire.sameCycle_rewire_iff}
\Sverified{Mettapedia.GraphTheory.FourColor.MarkedFaceReturnRewire.replacement_congr}

\paragraph{Actual open-tangle sewing.}
For an open tangle, temporarily fix each boundary dart under the edge
involution and compose with its vertex rotation to obtain an open face
permutation. On the disjoint union of two tangles, actual sewing changes
only boundary partners. After the existing dart reassociation, its face
permutation is exactly the disjoint open permutation followed by the
concrete boundary swap. Applying the return law therefore computes face
membership in the existing \texttt{composeRotationSystem}, not merely in
an abstract permutation claimed to model it.
\Sverified{Mettapedia.GraphTheory.FourColor.OpenTangleMarkedFace.compose_phi}
\Sverified{Mettapedia.GraphTheory.FourColor.OpenTangleMarkedFace.compose_sameCycle_iff}
Take $H$ to contain every port and the requested dart observations. For
boundary sizes $k_1,k_2$ and $t$ named darts, $|H|\le k_1+k_2+t$; the
return-permutation component has at most $(k_1+k_2+t)!$ possibilities.
This is not the cardinality of the entire refined Count state: support,
presence and naming coordinates still have to be retained.
\Sverified{Mettapedia.GraphTheory.FourColor.OpenTangleMarkedFace.card_observed_le}
\Sverified{Mettapedia.GraphTheory.FourColor.OpenTangleMarkedFace.card_return_records_le}

\paragraph{Why the cofacial partition alone is insufficient.}
The cycles $(0\,1\,2\,3)$ and $(0\,2\,1\,3)$ have identical cofacial
partitions. After right multiplication by $(0\,1)$, the first still puts
$0$ and $2$ together, while the second separates them. Both claims are
kernel-checked. Cyclic return order, not just existing face equivalence,
is necessary for compositional rewiring.
\Sverified{Mettapedia.GraphTheory.FourColor.MarkedFaceReturnRewire.same_partition_before}
\Sverified{Mettapedia.GraphTheory.FourColor.MarkedFaceReturnRewire.different_partition_after}
The independent diagnostic checks $232503$ small supported rewirings,
$288$ unmarked-arc subdivisions, and $1824$ physical prism sewing cases.
No colouring or configuration catalogue is enumerated. The generic
permutation theorem, unlike these diagnostics, has no carrier-size bound.

\paragraph{What the local law alone leaves.}
A supplied Count seam needs compatible named-face records on its two
actual shores, with their naming and port coordinates. Equality of Count
support alone does not assert this refinement. The next construction
discharges this compatibility through an enlarged finite state and the
existing geometric suppliers. Neither construction identifies preservation
of cofaciality with preservation of an entire hole word or facial cycle.

\subsection{Supplied reduction preserving named cofaciality}
\label{sec:supplied-named-cofacial-reduction}

There is now a supplied version of the preceding face semantics. Given a
sufficiently large bridgeless spherical cubic map with two-sided orbit
faces and finitely many named darts, the construction returns either an
actual ordered injective mesh, or a strictly smaller structural map with
the same named-colour acceptance and the same cofaciality relation among
the named darts. The size threshold depends on the mesh parameters and
the number of names, not on the lengths of their faces.

\paragraph{A record that keeps absence and aliases.}
Let $\ell:L\to\operatorname{Option}(D)$ locate names on one side and
let $H$ be its observed image. Embed $H$ in $\operatorname{Fin}(|L|)$,
extend its first-return permutation by fixed points on unused coordinates,
and retain the optional coordinate of every name. The record carrier has
exactly
\[
 |L|!\,(|L|+1)^{|L|}
\]
elements. Equality of two records constructs a bijection of their actual
observed dart carriers: equality of the name maps identifies the occupied
coordinate ranges, and equality of the padded permutations then preserves
first returns. Absent names and multiple names for the same dart are not
silently discarded.
\Sverified{Mettapedia.GraphTheory.FourColor.NamedFaceRecord.card_record}
\Sverified{Mettapedia.GraphTheory.FourColor.NamedFaceRecord.matching_named_iff}
\Sverified{Mettapedia.GraphTheory.FourColor.NamedFaceRecord.matching_permCongr}
\Sverified{Mettapedia.GraphTheory.FourColor.NamedFaceRecord.replacement_congr}

For a width-$k$ shore and $t$ persistent names, use $L$ equal to the sum of
the standard port carrier and the name carrier. All ports are present;
dart names may be absent or may alias each other or a port. The unchanged
exterior's return permutation combines with that of the inner side by
disjoint union. The induced observed-dart bijection commutes with the
actual seam swap because it fixes the standard port labels. Hence equal
inner records preserve and reflect named cofaciality after sewing to the
same exterior in \texttt{composeRotationSystem}.
\Sverified{Mettapedia.GraphTheory.FourColor.FaceReturnSum.nextHitPerm_sum}
\Sverified{Mettapedia.GraphTheory.FourColor.NamedOpenFace.totalMatching_swap_permCongr}
\Sverified{Mettapedia.GraphTheory.FourColor.NamedOpenFace.physicalCofacial_iff}
Complementary vertex-side sewing reconstructs the original face
permutation; the observations read the literal ambient names. The resulting
equivalence on an actual marked seam therefore concerns the original and
replacement maps, rather than two auxiliary open permutations.
\Sverified{Mettapedia.GraphTheory.FourColor.VertexSideCofacial.reassemblyEquiv_phi}
\Sverified{Mettapedia.GraphTheory.FourColor.VertexSideCofacial.physicalCofacial_reassembly_iff}
\Sverified{Mettapedia.GraphTheory.FourColor.MarkedCotreeSeam.SeamPair.cofacial_iff}

\paragraph{The geometric supplier chooses the equality.}
Retain the complete structural and frozen joint-colour state, and append
the named return record. If $Q(k,t)$ is the preceding bounded frozen-state
bound, a bound for the refined state is
\[
 \widehat Q(k,t)=Q(k,t)\,(k+t)!\,(k+t+1)^{k+t}.
\]
The bounded-width dependent sum is counted explicitly, rather than using
the state count at one exact width as its cardinality.
\Sverified{Mettapedia.GraphTheory.FourColor.CofacialSeamState.card_bounded_le}
Both the far-face contour branch and the bounded-radius cotree branch
now request enough literal nodes for $\widehat Q+1$. Pigeonhole selects
equal refined states; their projections supply the original structural
and joint-support equalities as well as the new face equality. Full
boundary order, connected shores and complements, fixed marked vertices,
and a strict unmarked star come from those geometric constructions.
\Sverified{Mettapedia.GraphTheory.FourColor.SphericalCofacialReduction.exists_joint_face_seam_of_nodes}
\Sverified{Mettapedia.GraphTheory.FourColor.SphericalCofacialReduction.exists_joint_face_seam_of_no_orderedMesh}
The final theorem supplies the smaller map, its structural certificate,
strict vertex-count decrease and lifted named darts. For every prescribed
colour assignment it preserves acceptance, and for every pair of names
it preserves and reflects cofaciality. No chain, radius bound, repeated
state, face-compatibility receipt, or ordinary non-colourability assumption
is an input to this final alternative.
\Sverified{Mettapedia.GraphTheory.FourColor.SphericalCofacialReduction.exists_orderedMesh_or_cofacial_replacement}

\paragraph{Scope and remaining crux.}
The independent gate checks $3446$ small named records, $8474$ supported
replacement comparisons and $288$ unmarked-arc subdivisions. Negative
controls show why absence and alias coordinates cannot be dropped. These
are permutation diagnostics, not colouring enumerations or claimed
counterexamples in the target class.

The ordered-mesh branch is still explicit. This construction does not
exclude all fixed meshes, supply the entire source-facing D1b/D1c
implication, or verify the generating base. It removes designated-face
perimeter from this named-cofaciality reduction's threshold, not from any
consumer which genuinely needs complete hole words, lengths or facial
cycles. The stronger whole-face construction remains available for those
consumers. The fixed obligation ledger remains $57/100$.

\subsection{Separated mesh deletions: exact gluing and the unrooted boundary limit}
\label{sec:separated-deletion-boundary}

The ordered-mesh alternative still needs a target-facing reduction. Its
common-core route admits a useful simplification, but the simplification
also exposes why the existing boundary observable is not a supply theorem.
No wall exclusion or new ledger credit is asserted here.

\paragraph{Two patches, with full local coverage.}
For a finite vertex set $S$, write $N[S]$ for its closed neighborhood in
the full ambient graph. Let $S$ and $T$ be two adjacent deleted vertex
pairs. If
\[
                         N[S]\cap N[T]=\varnothing,
\]
the complements $G-S$ and $G-T$ jointly contain every ambient edge and
every incident pair of ambient edges. Indeed, any endpoint of an edge
removed by $S$ lies in $N[S]$. If neither deletion retained an edge, its
endpoint would lie in both neighborhoods. If neither retained an incident
pair, their shared vertex would lie in both neighborhoods, using one
removed edge from each deletion.
\Sverified{Mettapedia.GraphTheory.FourColor.SeparatedPairDeletion.retained_or_retained}
\Sverified{Mettapedia.GraphTheory.FourColor.SeparatedPairDeletion.adjacent_retained_or_retained}

Two proper nonzero deletion colourings which agree literally on
$G-(S\cup T)$ therefore glue through the existing atlas theorem to an
ambient Tait colouring. In a graph-backed vertex-minimal counterexample,
no choices of the two deletion colourings can have equal common-core
restrictions. In particular, a proper Kempe-repaired representative cannot
agree either: the strict-repair alternative of
Corollary~\ref{cor:common-core-disagreement-residue} is impossible at this
pair, leaving its branching/boundary disjunction.
\Sverified{Mettapedia.GraphTheory.FourColor.SeparatedPairDeletion.not_commonCore_eq_of_minimal}
\Sverified{Mettapedia.GraphTheory.FourColor.SeparatedPairDeletion.branching_or_taitReachesSecondPair_of_separated}

\paragraph{Ambient separation is constructed, not presumed.}
In maximum degree $d$,
\[
                  |N[S]|\leq(d+1)|S|.
\]
Thus for a cubic graph and a deleted pair $S$, the twice-closed
neighborhood $N[N[S]]$ has at most $32$ vertices. The first and second
endpoint maps of the designated steps on an ordered mesh row are
separately injective. At most $64$ such steps meet this forbidden set.
If the row has $n\geq65$ branch intervals, every designated source step
has a different target step with both endpoints outside $N[N[S]]$.
Symmetry of adjacency gives disjoint closed neighborhoods. All ambient
shortcuts and lateral attachments are included in this argument; row
distance alone is not used as a separation certificate.
\Sverified{Mettapedia.GraphTheory.FourColor.SeparatedPairDeletion.card_close_le}
\Sverified{Mettapedia.GraphTheory.FourColor.SeparatedPairDeletion.disjoint_close_of_avoid_twice}
\Sverified{Mettapedia.GraphTheory.FourColor.SeparatedOrderedMeshSites.exists_separated_rowSite}
\Sverified{Mettapedia.GraphTheory.FourColor.SeparatedOrderedMeshSites.exists_rowSite_with_no_commonCore_agreement}
\Sverified{Mettapedia.GraphTheory.FourColor.SeparatedOrderedMeshSites.exists_rowSite_forcing_branching_or_boundary}

\paragraph{The boundary horn is already true locally.}
Fix any proper nonzero colouring of $G-S$. Choose a port $z$ of the
second deleted pair $T$, with exposed edge $tz$ for $t\in T$.
The neighborhoods are disjoint, so $z\notin N[S]$ and every edge at $z$
survives the first deletion. Cubicity at $z$ supplies a neighbor
$w\notin T$. Consequently $zw$ is an edge of the common core, while
$tz$ is an adjacent exposed edge of $G-S$.

Select their two colours. They are distinct by properness and nonzero by
the Tait condition. The common-core two-colour component containing $zw$
touches $tz$, exactly witnessing
\texttt{FirstComponentReachesSecondPair}. Taking the current colouring
to be the original base and the preceding Kempe sequence to be empty
then witnesses \texttt{FirstTaitOrbitReachesSecondPair}.
This construction requires neither a target colouring nor planarity,
minimality, a mesh, or any path from the first deletion.
\Sverified{Mettapedia.GraphTheory.FourColor.SeparatedDeletionBoundary.atBase_of_adjacent_edges}
\Sverified{Mettapedia.GraphTheory.FourColor.SeparatedDeletionBoundary.atBase_of_separated}
\Sverified{Mettapedia.GraphTheory.FourColor.SeparatedDeletionBoundary.orbit_of_separated}
\Sverified{Mettapedia.GraphTheory.FourColor.SeparatedDeletionBoundary.exists_rowSite_with_atBase_boundary}

The actual nine-site horn predicate is therefore true for every assignment
whenever its selected family contains a separated pair. Its universal
no-horn premise is then false, not a surviving difficult synchronization
case. This does not assert separation for every arbitrary nine-site
selection, nor refute the earlier conditional theorems.
\Sverified{Mettapedia.GraphTheory.FourColor.SeparatedDeletionBoundary.nineSite_horn_of_separated_pair}
\Sverified{Mettapedia.GraphTheory.FourColor.SeparatedDeletionBoundary.not_nineSite_no_horn_of_separated_pair}

\paragraph{Verification and remaining implication.}
The structural gate checks $12900$ labelled subcubic graphs on four to
six vertices, $186270$ deletion-pair comparisons, and $780$ designated
sources in full ambient cubic row controls. Its calibrations are all
$64$ labelled four-vertex graphs and the $70$ labelled cubic six-vertex
graphs. Negative controls distinguish disjoint deletions, coverage of
individual edges, and coverage of incident pairs; an off-row shortcut
can invalidate separation. Every long-row control also supplies the
local two-edge boundary witness. No colourings or configuration
certificates are enumerated.

The next high-width implication must retain more information than this
unrooted encounter: for example, a component selected by an actual
discrepancy-repair sequence, with the anchoring or separating data needed
by a physical reduction. The present theorem does not supply those data
or justify assuming them. It removes a common-core repair branch at
supplied separated sites and identifies the exact weakness of its
remaining boundary observable. Fixed ordered-mesh exclusion, the full
D1b/D1c supply, and the generating base remain open; the ledger is still
$57/100$.

\subsection{Discrepancy-rooted component fusion at joint orbit minima}
\label{sec:discrepancy-rooted-fusion}

The unrooted boundary limitation in
Section~\ref{sec:separated-deletion-boundary} is addressed by an exact
single-switch criterion and an actual disagreement-minimizing construction.
The result supplies two-sided rooted witnesses in one explicit alternative;
it does not yet supply bounded interfaces or eliminate the other alternative.

\paragraph{The exact single-switch obstruction.}
For a graph embedding $f:H\hookrightarrow G$, let $f_*$ be its map on
connected components. A small component $K$ \emph{fuses} if some
$K'\ne K$ has $f_*(K')=f_*(K)$. On the selected two-colour line graphs,
for distinct colours $a,b$,
\[
 \text{the switch on $K$ is the restriction of one ambient $(a,b)$ switch}
 \quad\Longleftrightarrow\quad
 \neg\exists K'\ne K:\ f_*(K')=f_*(K).
\]
The forward implication follows because the changed-edge set is exactly
the switched component. The reverse uses the ambient component $f_*(K)$.
Boundary contact is allowed in this reverse implication: an excursion
outside the image need not merge two small components.
\Sverified{Mettapedia.GraphTheory.FourColor.KempeEmbeddingFusion.exists_lifted_switch_iff_not_fuses}

Fusion gives a walk between different small components, and every such
walk leaves the image. This is stronger than an arbitrary outside neighbor.
\Sverified{Mettapedia.GraphTheory.FourColor.KempeEmbeddingFusion.fuses_iff_reachable}
\Sverified{Mettapedia.GraphTheory.FourColor.KempeEmbeddingFusion.walk_leaves_range}
\Sverified{Mettapedia.GraphTheory.FourColor.KempeEmbeddingFusion.exists_fusion_walk}

A prescribed finite common-core repair can be lifted step by step until
a genuinely fusing step, retaining an exactly lifted prefix and a suffix
to the target. This does not establish productive progress by itself:
an existential suffix may undo the switch. The stronger argument below
uses literal disagreement, not a suffix as a progress certificate.
\Sverified{Mettapedia.GraphTheory.FourColor.KempeTargetFusion.targetFusion_or_lifted}
\Sverified{Mettapedia.GraphTheory.FourColor.KempeTargetFusion.TargetFusion.exists_failed_step}

\paragraph{A wrong edge forces a correct-edge blocker.}
Fix a proper nonzero target colouring $t$ on the common core $C$ and
a deletion colouring $c_0$. Minimize
\[
             d(c,t)=|\{e\in E(C):c(e)\ne t(e)\}|
\]
over the valid-pair Kempe orbit of $c_0$. This minimum is constructed
by well-ordering the natural-valued measure; it is not supplied as a
hypothesis. For any wrong edge $e$, switch its ambient component using
the pair $(c(e),t(e))$. The switch fixes $e$. Minimality forces some
previously correct common-core edge $f$ to become wrong, and $f$ lies
in that same ambient component.
\Sverified{Mettapedia.GraphTheory.FourColor.KempeTargetFusion.exists_minimum}
\Sverified{Mettapedia.GraphTheory.FourColor.KempeTargetFusion.blocking_edge_of_minimum}

If $c|_C$ and $t$ are locally related by one transposition at each vertex,
the existing seed-switch theorem says that switching the small component
of $e$ fixes all its edges and creates no errors. Hence a correct-edge
blocker cannot belong to this small component: it belongs to a different
component fused with it upstairs. Every wrong edge therefore roots a
wrong-to-correct fusion witness. No target-reachability premise is needed.
\Sverified{Mettapedia.GraphTheory.FourColor.KempeTargetFusion.seed_fuses_of_minimum_of_local}
\Sverified{Mettapedia.GraphTheory.FourColor.KempeTargetFusion.seedBlockingFusion_of_minimum_of_local}

\paragraph{Both deletions, with the same root.}
For separated deleted pairs $S,T$, minimize common-core disagreement
over the product of the two deletion-colouring orbits. Holding either
representative fixed gives the requisite minimum on the other side.
The two-patch gluing theorem rules out zero disagreement, so an actual
wrong edge exists. At this constructed pair, either some common-core
vertex has a non-transposition discrepancy, or \emph{every} wrong edge
roots blocking fusions in both deletion graphs. The ordered-row
construction supplies the separated target for every source when the
row has at least $65$ branch intervals, and the conclusion holds for
every pair of initial proper nonzero deletion colourings.
\Sverified{Mettapedia.GraphTheory.FourColor.KempeTargetFusion.exists_joint_minimum}
\Sverified{Mettapedia.GraphTheory.FourColor.SeparatedDeletionFusion.jointMinimumObstruction_of_separated}
\Sverified{Mettapedia.GraphTheory.FourColor.SeparatedDeletionFusion.exists_rowSite_with_joint_minimum_obstruction}

Instantiating the existing minimality-supplied base colourings leaves
no colouring, minimum, discrepancy, or reachability premise to supply.
\Sverified{Mettapedia.GraphTheory.FourColor.SeparatedDeletionFusion.exists_selected_joint_minimum_obstruction}

In the fusion alternative, the walk in $G-S$ encounters an edge at $T$;
the walk in $G-T$ encounters an edge at $S$. They share the chosen
common-edge root and use the same unordered pair of colours, but belong
to two different deletion colourings. They are not claimed to be a
single ambient Tait component, disjoint strands, or the boundary of a noose.
\Sverified{Mettapedia.GraphTheory.FourColor.SeparatedDeletionFusion.physical_walk_of_seedBlockingFusion}
\Sverified{Mettapedia.GraphTheory.FourColor.SeparatedDeletionFusion.physical_second_walk_of_seedBlockingFusion}

\paragraph{Gate and remaining geometric obligation.}
The generic structural gate checks all $1099$ labelled simple graphs
on one through five vertices, giving $47785$ component/subset cases.
Path controls distinguish harmless boundary contact from actual fusion.
Independent finite-set toggle checks cover $9996$ one-sided balance cases
and $4190$ joint-minimum cases. These are component and disagreement-set
diagnostics, not configuration or ring-colouring enumerations, and are
not asserted to be examples in the least-counterexample class.

The surviving obligations are precise: resolve the local
non-transposition alternative, and turn the paired rooted fusion walks
into a physically valid reduction with controlled complete interfaces.
The present argument does not prove either implication, a long nested
chain, fixed ordered-mesh exclusion, or the generating base. It does
replace arbitrary boundary encounters by discrepancy-driven witnesses
on an actually supplied pair. The fixed ledger remains $57/100$.

\subsection{Removing the choice of Kempe orbits from the minimum}
\label{sec:absolute-deletion-minimum}

The minimum in Section~\ref{sec:discrepancy-rooted-fusion} was taken
over two specified Kempe orbits. A calibrated reconfiguration experiment
shows why positivity of that restricted minimum is not itself a geometric
contradiction, even with planarity and no local branching. The construction
below removes this dependence on the initial orbit choices.

\paragraph{Restriction and simultaneous compatibility.}
One ambient component can restrict to several smaller components.
Switching it need not restrict to one switch, but does restrict to a
finite valid-pair Kempe sequence. The proof first restricts the local
transposition relation and then uses local-swap generation. Thus an ambient
Kempe orbit lies in one deletion orbit under restriction; no converse
lifting assertion is made.
\Sverified{Mettapedia.GraphTheory.FourColor.KempeRestriction.locallySwapRelated_switch}
\Sverified{Mettapedia.GraphTheory.FourColor.KempeRestriction.locallySwapRelated_pullback}
\Sverified{Mettapedia.GraphTheory.FourColor.KempeRestriction.reachable_pullback_of_step}
\Sverified{Mettapedia.GraphTheory.FourColor.KempeRestriction.reachable_pullback}

For separated deletions $S,T$, two specified deletion orbits have
representatives agreeing on the common core if and only if there is
\emph{one} ambient Tait colouring whose restrictions belong to both
orbits. This follows by constructing the glued colouring with its two
prescribed restrictions. Individual extendability of each orbit does
not replace simultaneous extendability.
\Sverified{Mettapedia.GraphTheory.FourColor.DeletionOrbitCompatibility.glued_restrict}
\Sverified{Mettapedia.GraphTheory.FourColor.DeletionOrbitCompatibility.exists_extension_of_agreement}
\Sverified{Mettapedia.GraphTheory.FourColor.DeletionOrbitCompatibility.compatible_iff_simultaneouslyExtendable}
\Sverified{Mettapedia.GraphTheory.FourColor.DeletionOrbitCompatibility.simultaneous_membership_of_ambient_reachable}

Allowing arbitrary orbit choices makes existence of compatible choices
exactly ambient Tait colourability. This is an equivalence describing the
remaining task, not a new assumed supplier of compatible orbits.
\Sverified{Mettapedia.GraphTheory.FourColor.DeletionOrbitCompatibility.exists_compatible_iff_taitColorable}

\paragraph{A colourable negative control.}
The structural lab uses two independent complete enumerators, one ordered
by edges and one by vertices. They recover $6$, $12$, and $60$ Tait
colourings for $K_4$, $K_{3,3}$, and the dodecahedron, respectively.
The last graph's ten six-element Kempe classes agree with the description
in Khovanov--Robert, \emph{Conical $SL(3)$ foams}, Section~3.2,
\href{https://doi.org/10.4171/JCA/61}{doi:10.4171/JCA/61}.
Its spherical cubic rotation is checked independently of the colouring
enumeration. No least-counterexample assumption is asserted for it.

Each of its $30$ adjacent-pair deletions has $108$ colourings in two
Kempe orbits of sizes $36$ and $72$. The $135$ separated deletion pairs
give $540$ orbit-pair comparisons. Agreement is possible in $525$ cases;
the other $15$ have minimum disagreement $7$. All $360$ minimum pairs
in these positive cases are locally transposition-related. For one
minimum in each exceptional case, the lab directly checks a correct-edge
fusion blocker and an outside-image path on both sides of every wrong
edge ($210$ rooted checks). Both orbits in the recorded incompatible
pair have nonempty, disjoint sets of ambient Kempe classes extending
them. Consequently the observed rooted fusions and positive orbit minimum
do not alone contradict planarity or ambient colourability.
These are reproducible computational diagnostics, not Lean-certified
graph specimens, ring-word tables, or reducibility certificates.

\paragraph{The absolute minimum is constructed.}
Instead minimize
\[
  d(c,d)=|\{e\in E(G-(S\cup T)):c(e)\ne d(e)\}|
\]
over \emph{all} proper nonzero colourings of $G-S$ and $G-T$.
The input colourings witness nonemptiness only. Well-ordering of the
natural-valued measure supplies a pair whose disagreement cannot be
decreased by changing either colouring, or both, into any other Kempe
orbit. This minimum supplies both one-sided orbit minima needed by the
existing rooted-blocker theorem.
\Sverified{Mettapedia.GraphTheory.FourColor.DeletionAbsoluteMinimum.exists_absolute_minimum}
\Sverified{Mettapedia.GraphTheory.FourColor.DeletionAbsoluteMinimum.AbsoluteMinimum.orbit_minima}

For a colourable ambient graph every such absolute minimum is zero.
For separated deletions the converse also holds, by gluing; on all $135$
separated dodecahedral controls the unrestricted minimum is indeed zero.
\Sverified{Mettapedia.GraphTheory.FourColor.DeletionAbsoluteMinimum.absolute_minimum_eq_of_tait}
\Sverified{Mettapedia.GraphTheory.FourColor.DeletionAbsoluteMinimum.absolute_minimum_eq_iff_taitColorable}

In a least counterexample, zero is impossible. At the same constructed
absolute minimizers, either there is a local non-transposition discrepancy,
or every wrong edge roots blocking fusions in both deletion graphs.
The final ordered-row corollary supplies the separated target and uses
minimality to supply the nonempty colouring spaces; its inputs contain
no colouring, chosen Kempe orbit, minimum, or reachability premise.
\Sverified{Mettapedia.GraphTheory.FourColor.DeletionAbsoluteMinimum.absoluteObstruction_of_separated}
\Sverified{Mettapedia.GraphTheory.FourColor.DeletionAbsoluteMinimum.exists_rowSite_with_absolute_obstruction}

This strengthens the target-facing minimum and excludes the arbitrary-orbit
artefact exposed by the lab. It does not eliminate branching or construct
a bounded complete interface from the paired walks. Fixed ordered-mesh
exclusion, full D1b/D1c supply, and the generating base remain open.
The fixed obligation ledger stays $57/100$.

\subsection{Whole-component exchange at the absolute minimum}
\label{sec:whole-component-exchange}

Let $D$ be the common-core edges where the two deletion colourings differ,
and connect two edges of $D$ when they share a vertex. This is the graph
of \emph{all} disagreement, not one selected two-colour Kempe graph.
Its components may branch. A whole component can be copied from the
donor colouring to the recipient whenever no incident recipient edge
leaves the common embedding. At the component frontier, every unselected
common edge already has equal colours. Hence the copied assignment is
proper, remains nonzero, and strictly decreases disagreement.
\Sverified{Mettapedia.GraphTheory.FourColor.DisagreementComponentExchange.copyColoring_tait}
\Sverified{Mettapedia.GraphTheory.FourColor.DisagreementComponentExchange.copyColoring_decreases}

The constructed absolute minimum permits this change even if it leaves
the original Kempe orbit. Applying exchange in each direction proves
that \emph{every} nonempty disagreement component has a boundary contact
in both deletion embeddings. Both contacts lie in the same component;
there is no exceptional non-transposition case.
\Sverified{Mettapedia.GraphTheory.FourColor.DisagreementComponentExchange.boundary_of_minimum}
\Sverified{Mettapedia.GraphTheory.FourColor.DisagreementComponentExchange.absolute_component_touches_both}

An adjacent edge outside the first common-core embedding leads to a
vertex of the second deleted pair. Its common endpoint is therefore one
of that pair's four ports; the reverse direction is symmetric. Thus every
component joins the two sets of ports. Distinct components cannot use the
same port, so there are at most four components. The ordered-row supplier
constructs this situation, including the colourings, absolute minimum,
and nonempty disagreement, from the least-counterexample and ordered-mesh
inputs already used in Section~\ref{sec:absolute-deletion-minimum}.
\Sverified{Mettapedia.GraphTheory.FourColor.DeletionDisagreementSpanning.allComponentsSpan_of_absolute}
\Sverified{Mettapedia.GraphTheory.FourColor.DeletionDisagreementSpanning.card_components_le_four}
\Sverified{Mettapedia.GraphTheory.FourColor.DeletionDisagreementSpanning.exists_rowSite_with_spanning_minimum}

\paragraph{Controls and rejected shortcuts.}
A complete, independently cross-checked enumeration for one separated
dodecahedral site pair tests all $108^2=11664$ deletion-colouring pairs.
It exercises $528$ boundary-avoiding component copies, all valid and
strictly improving, and finds $26688$ invalid copies when the boundary
condition is omitted. Counts are over attempted oriented component copies,
not distinct counterexamples. A separate three-edge-star control proves
in the kernel that the unrestricted copying rule is false.
\Sverified{Mettapedia.GraphTheory.FourColor.ComponentExchangeControls.unchecked_copy_fails}

Closed-map face potentials cannot simply be reused on a deleted graph:
the sum of two distinct nonzero Klein colours at a degree-two port is
nonzero. A relative boundary construction would be needed before invoking
the closed spherical cubic face-potential theorem.
\Sverified{Mettapedia.GraphTheory.FourColor.ComponentExchangeControls.degree_two_flux_ne_zero}

The next attempted interpolation also has an exact local limitation.
Normalize the three incident colours of one proper cubic star to the
identity. If a second and a proposed mixed star are permutations $p,q$
of these colours, and $q(i)\in\{i,p(i)\}$ on every edge, then $q$ is
either the identity or $p$. Merely choosing between the two existing
colours edge by edge cannot create a third proper star, including in the
three-cycle case. The lab exhausts all $36$ ordered pairs of star
permutations and their eight binary selections; the same normalized
statement is kernel-decided.
\Sverified{Mettapedia.GraphTheory.FourColor.ComponentExchangeControls.cubic_binary_mixing_rigid}

\paragraph{What remains geometric.}
Four components need not mean four short paths, small components, or
bounded complete interfaces. The present proofs supply no such bound.
Likewise mesh isoperimetry bounds branch vertices on a small side, not
every ambient vertex or every attached corridor. Neither it nor a finite
control census refutes all size-sensitive reductions. A distance argument
would need ambient separation of the sites (not just distant mesh indices),
and port-to-site attachment edges must be included in its constants.
The next missing implication is a route-native reduction or controlled
interface construction for these anchored disagreement components; a
catalogue of replaceable patches is not supplied or used. This is progress
inside the open D1b/D1c attack, not discharge of those obligations or of
fixed ordered-mesh exclusion. The ledger remains $57/100$.
The exchange statements concern ordinary Tait colourings; preservation of
the source's additional marked-pair equality constraints in a refined
count is not asserted.

\subsection{Adversarial audit of the bounded-interface pumping step}
\label{sec:bounded-interface-pumping-audit}

The proposed bounded-interface continuation survives the algebraic attack
\emph{when the repeated state is complete}. It does not follow from merely
revisiting an interface type. A single interface state and a single Count
letter $[2]$ give the distinct cumulative rows $[2^n]$. There is no finite
lossless code for these exact rows, although their supports are identical.
Thus finitely many interface objects and finitely many generators do not
imply finite exact Count image. This control refutes that inference, not
the finiteness of a Booleanized interface semantics.
\Sverified{Mettapedia.GraphTheory.FourColor.BoundedInterfacePumpingAudit.doublingRow_step}
\Sverified{Mettapedia.GraphTheory.FourColor.BoundedInterfacePumpingAudit.no_finite_exact_count_code}
\Sverified{Mettapedia.GraphTheory.FourColor.BoundedInterfacePumpingAudit.doublingRow_support}

Nor may a support be replaced by its cardinality. A two-state swap maps
$\{0\}$ to $\{1\}$; deleting this one-letter endomorphism preserves the
interface and the support cardinality but changes acceptance against
$\{0\}$. More generally, two abstract supports are indistinguishable by
\emph{all} terminal-set intersection tests exactly when they are equal.
Any code preserving all those tests on all subsets of a state set $Q$
therefore needs at least $2^{|Q|}$ codes. This is not a lower bound on the
actually realizable planar-tangle supports: restricting supports or
contexts can permit further compression.
\Sverified{Mettapedia.GraphTheory.FourColor.BoundedInterfacePumpingAudit.cardinality_repeat_does_not_license_deletion}
\Sverified{Mettapedia.GraphTheory.FourColor.BoundedInterfacePumpingAudit.support_eq_iff_terminal_tests}
\Sverified{Mettapedia.GraphTheory.FourColor.BoundedInterfacePumpingAudit.complete_code_card_lower_bound}

Topology is a separate field of the receipt. A three-cycle and its inverse
allow the same nonempty proper-triple support, but identity matching of two
three-port seams is orientation reversing for opposite rotations and not
for equal rotations. This is a kernel-checked seam-algebra example, not a
constructed least-counterexample or a claim that every non-reversing glue
is nonspherical.
\Sverified{Mettapedia.GraphTheory.FourColor.BoundedInterfacePumpingAudit.three_port_support_ignores_orientation}
\Sverified{Mettapedia.GraphTheory.FourColor.BoundedInterfacePumpingAudit.seam_order_matters}

\paragraph{Surviving implication and the remaining target.}
The existing support-deletion theorem already preserves every terminal
test after a repeated complete prefix support. The physical serial-tangle
composition theorem identifies this relational semantics with colouring,
and the literal-shore chain theorem separately supplies the physical
minimality contradiction from its certified shore hypotheses. These
theorems were re-audited, not re-proved or newly credited.
\Sverified{Mettapedia.GraphTheory.FourColor.GoertzelV24CountSupportSynchronization.runSupport_delete_block}
\Sverified{Mettapedia.GraphTheory.FourColor.GoertzelV24RawNooseCountPumping.supportRel_series}
\Sverified{Mettapedia.GraphTheory.FourColor.GoertzelV24MajorityShoreStateDescent.length_le_of_literalShoreChain}

An independent two-state control exhausts all $262144$ combinations of
four Boolean transfer letters and an initial support, with $1048576$
terminal acceptance comparisons after deleting a repeated-support block;
there are no failures. Of these runs, $93255$ have nonempty final support,
so the control is not merely testing rejection of empty supports. The
control concerns transfer algebra, not a census of planar configurations.

The open implication is still geometric: a sufficiently large anchored
disagreement region must supply sufficiently many properly nested,
bounded, complete interfaces, or some other route-native reduction.
Neither unbounded exact counts nor the failure of coarse summaries
refutes that implication. No classical configuration machinery is used;
the obligation ledger remains $57/100$.

\subsection{Agreement saturation: a construction and its exact obstruction}
\label{sec:agreement-closed-cuts}

One attempted supplier enlarges a seed shore along every edge on which
the two colourings agree. Let $H$ be this agreeing-edge subgraph, with
the full ambient vertex set. The enlargement is constructed as the set
of vertices reachable in $H$ from the seeds. It is the least $H$-closed
set containing the seeds; the operation is monotone and idempotent.
\Sverified{Mettapedia.GraphTheory.FourColor.AgreementClosedCuts.mem_hull_iff}
\Sverified{Mettapedia.GraphTheory.FourColor.AgreementClosedCuts.hull_least}
\Sverified{Mettapedia.GraphTheory.FourColor.AgreementClosedCuts.hull_mono}
\Sverified{Mettapedia.GraphTheory.FourColor.AgreementClosedCuts.hull_idempotent}

This gives an exact, constructive criterion: there is a shore containing
$A$, avoiding $B$, and having only disagreeing ambient crossing edges
if and only if no $H$-path joins a vertex of $A$ to a vertex of $B$.
Necessity follows by propagating shore membership along a path;
sufficiency uses the constructed reachable set. Closure is explicitly
equivalent to the assertion about \emph{all ambient} crossing edges.
No width estimate or long chain is part of this theorem.
\Sverified{Mettapedia.GraphTheory.FourColor.AgreementClosedCuts.closed_iff_boundary_disagrees}
\Sverified{Mettapedia.GraphTheory.FourColor.AgreementClosedCuts.exists_disagreeing_separator_iff}

\paragraph{A proper-colouring negative control.}
On the cube, alternate two colours around both squares and use the third
colour on the spokes. Swap the first two colours on the upper square
only. Both edge colourings are proper and nonzero, and disagreement is
exactly the upper square. The agreeing subgraph nevertheless connects all
eight vertices through the lower square and the spokes. Every nonempty
seed therefore saturates to the whole cube, and no $H$-closed shore
separates upper vertices $0$ and $2$. These facts are kernel-checked.
\Sverified{Mettapedia.GraphTheory.FourColor.AgreementHullPrismControl.proper_nonzero}
\Sverified{Mettapedia.GraphTheory.FourColor.AgreementHullPrismControl.disagrees_exactly_on_top}
\Sverified{Mettapedia.GraphTheory.FourColor.AgreementHullPrismControl.agreeing_preconnected}
\Sverified{Mettapedia.GraphTheory.FourColor.AgreementHullPrismControl.every_nonempty_hull_is_univ}
\Sverified{Mettapedia.GraphTheory.FourColor.AgreementHullPrismControl.no_closed_separation}

The accompanying lab checks the same fixed colour pair on $19$ even
prisms of circumferences $4$ through $40$, over $399$ proper top prefixes.
A prefix of $t$ upper vertices has two disagreeing crossing edges but
$t+2$ full ambient crossing edges. Adding the corresponding lower prefix,
instead of saturating along every agreeing edge, gives a strictly nested
chain of connected, co-connected shores with four full crossing edges.
Thus the lab contains a surviving selective construction on the same
family that defeats indiscriminate saturation.

\paragraph{Scope of the failed attack.}
These controls are ordinary colourable graphs, not positive absolute
deletion minima, least counterexamples, or instances of the ordered-mesh
supplier. They refute only the proposed generic enlargement rule. They
do not refute the existence of another controlled-interface construction;
indeed the paired-prefix control exhibits one on prisms. The open step
is to construct an appropriate \emph{selective} sweep or another
route-native reduction from the full absolute-minimum and ordered-mesh
hypotheses, bounding the agreeing crossing edges as well as the
disagreeing ones. No hypothesis asserting that sweep has been substituted
for its construction. No obligation row is discharged; the ledger remains
$57/100$.

\subsection{A disk-local repair can be overconstrained by its exterior}
\label{sec:disk-exchange-boundary-control}

The next attempted construction recolours both deletion graphs inside a
disk containing the two sites and the entire disagreement, while fixing
every edge outside its interior. This is stronger than merely assuming
sphericity. It can fail even when arbitrary joint recolouring is allowed,
with no restriction to Kempe orbits.

The control is the spherical dodecahedron with the explicit edge and star
presentation in \texttt{DiskExchangeBoundaryControl.lean}. Its deleted
adjacent pairs are $\{0,1\}$ and $\{17,19\}$, whose closed neighbourhoods
are disjoint. The chosen disk consists of ten faces; its complement is
two adjacent pentagons. The lab checks the spherical rotation, both
connected facial sides, the complete simple eight-edge rim, and the disk
Euler count. All edges incident with either deletion site are strictly
inside the disk. These geometric checks are computational; the Lean
theorems below certify the colouring obstruction on the stated edge
presentation, not a separate geometric realization theorem.

There are nine frozen edges: the eight rim edges and the shared edge of
the two exterior pentagons. In edge-index order they are
\[
  (10,11,12,13,14,17,23,25,27),
\]
with colour codes $(2,0,1,0,0,2,1,1,2)$. Here $0,1,2$ denote the three
nonzero Tait colours, not the zero Klein element. Their dictionary and
the incident-star presentation are kernel checked.
\Sverified{Mettapedia.GraphTheory.FourColor.DiskExchangeBoundaryControl.color_encoding}
\Sverified{Mettapedia.GraphTheory.FourColor.DiskExchangeBoundaryControl.presentation_valid}

Two independent complete enumerators find $60$ full colourings and $108$
colourings of each deletion. With this exterior fixed, those populations
shrink to $0$, $2$, and $1$ respectively. Both remaining deletion pairs
disagree on exactly six common-core edges, all inside the disk. Thus even
the best boundary-fixed exchange cannot improve them. The two-colouring
existence witnesses, impossibility of a full extension, and the gluing
implication are proved in Lean. A least-disagreement argument then proves
existence of a positive boundary-fixed minimum and refutes the universal
strict-improvement rule, without trusting the enumeration for minimality.
\Sverified{Mettapedia.GraphTheory.FourColor.DiskExchangeBoundaryControl.deletion_witnesses}
\Sverified{Mettapedia.GraphTheory.FourColor.DiskExchangeBoundaryControl.no_full_extension}
\Sverified{Mettapedia.GraphTheory.FourColor.DiskExchangeBoundaryControl.glue_proper}
\Sverified{Mettapedia.GraphTheory.FourColor.DiskExchangeBoundaryControl.no_fixed_agreement}
\Sverified{Mettapedia.GraphTheory.FourColor.DiskExchangeBoundaryControl.exists_frozen_positive_minimum}
\Sverified{Mettapedia.GraphTheory.FourColor.DiskExchangeBoundaryControl.no_universal_fixed_exchange}

The graph itself remains colourable. An explicit full colouring changes
exactly one of the nine frozen edges, edge $11=(5,9)$; Lean verifies this
witness. The lab also verifies that adding either exterior face to the
disk leaves two full extensions of the new exterior. There is no claim
that one-face enlargement suffices in general.
\Sverified{Mettapedia.GraphTheory.FourColor.DiskExchangeBoundaryControl.escape_changes_one}
The receipts and replay are in
\texttt{results/fourcolor/v24\_disk\_exchange\_gate.json} and
\texttt{tools/fourcolor/v24\_disk\_exchange\_gate.py}.

\paragraph{What changes in the proof plan.}
This refutes our proposed fixed-exterior localization rule, not an
assertion that the source requires that rule. It is not a positive
boundary-unrestricted absolute minimum or a least counterexample.
Likewise a general theorem making every positive absolute minimum
strictly improve would force agreement and hence ambient colourability;
that is not automatically an easier intermediate target than the source's
large-instance descent theorem.

The source-compatible target remains an actual smaller admissible piece
whose complete boundary behaviour transfers zero Count. The existing
support-based physical replacement theorem already gives that semantic
transfer; it does not require repairing an arbitrarily selected exterior
colouring. It still requires construction of the replacement, a valid
physical seam and strict size decrease. No such construction is assumed
or claimed here, and no configuration catalogue is introduced. This is a
refutation checkpoint on the global reduction attack, with the fixed
obligation ledger unchanged at $57/100$.

\subsection{An erased colour mark is not a surviving edge name}
\label{sec:marked-square-observable}

Addendum~XXVI, flag L4, requires transport when a ladder reduction touches
one of the two named edges.  The ordinary square count identity does not
provide this by assigning the erased name to a surviving seam edge.
This is a different question from the retired arbitrary-target
Kempe-orbit transport claim: the obstruction here is already an ordinary
refined count at a fixed boundary word.

Let the four outer colours of a facial square be
$w=(r,r,r,r)$, and mark the internal edge $x_0$ between corners $0$ and $1$.
Use the colour of outer port $2$ as the reference; that port is not incident
with $x_0$.  The two square extensions are
\[
             (b,p,b,p),\qquad(p,b,p,b).
\]
In both planar smoothings, both surviving seam edges have colour $r$.
Thus the target $x_0=w_2$ has count zero upstairs and count one on
\emph{each} smoothing under \emph{either} literal seam-edge image.
The unmarked counts still obey $2=1+1$.
\Sverified{Mettapedia.GraphTheory.FourColor.MarkedSquareObservable.explicit_extensions}
\Sverified{Mettapedia.GraphTheory.FourColor.MarkedSquareObservable.literal_image_zero_failure}
\Sverified{Mettapedia.GraphTheory.FourColor.MarkedSquareObservable.literal_count_failure}

Even choosing a boundary-dependent extension for each labelled smoothing
cannot be done equivariantly without extra information: the colour
transposition exchanging $b,p$ fixes $w$ and each labelled smoothing
colouring, but exchanges the two square extensions.  The kernel proves
that no selector into the square fibre is fixed by this transposition.
This does not forbid a non-equivariant choice, or a richer transported
observable; it rules out treating the unmarked cardinal identity as an
equivariant fibre identification with an internal colour observation.
\Sverified{Mettapedia.GraphTheory.FourColor.MarkedSquareObservable.no_equivariant_selector}

There is an exact finite local repair. Retain a virtual seed $t=x_0$ and
reconstruct all erased colours in the Klein group by
\[
 R(w,t)=\bigl(t,\ w_1+t,\ w_2+w_1+t,\ w_3+w_2+w_1+t\bigr).
\]
For a nonzero boundary word, validity is precisely
\[
       w_0+w_1+w_2+w_3=0,
       \qquad R(w,t)_i\ne0\quad(0\le i<4).
\]
Every square extension is uniquely $R(w,x_0)$.  The reconstruction
commutes with additive colour automorphisms when the seed transforms too.
For any predicate $P$ inspecting any or all internal colours, there is
an actual equivalence
\[
 \{x:\operatorname{SquareExt}(w,x)\land P(x)\}
 \simeq
 \{t:\operatorname{SquareExt}(w,R(w,t))\land P(R(w,t))\}.
\]
Consequently the two counts, and their zero predicates, agree.  The seed
has at most three nonzero values; it is not a physical edge of a smoothing.
\Sverified{Mettapedia.GraphTheory.FourColor.MarkedSquareObservable.reconstruct_map}
\Sverified{Mettapedia.GraphTheory.FourColor.MarkedSquareObservable.extension_reconstruct}
\Sverified{Mettapedia.GraphTheory.FourColor.MarkedSquareObservable.seedOK_iff}
\Sverified{Mettapedia.GraphTheory.FourColor.MarkedSquareObservable.targetEquiv}
\Sverified{Mettapedia.GraphTheory.FourColor.MarkedSquareObservable.target_count}
\Sverified{Mettapedia.GraphTheory.FourColor.MarkedSquareObservable.target_empty_iff}

The constructor gate independently compares all $81$ nonzero boundary
words, their $18$ total extensions, and $972$ tagged-edge equality tests.
Its receipt is \texttt{results/fourcolor/v24\_marked\_square\_gate.json}.
This is the source's generic square operator, not a configuration
catalogue or a reducibility check.  The new equivalence corrects its
observable semantics; it does not identify the seed fibre with colourings
of a smaller admissible physical map.  Closing that realization and
closure obligation, or avoiding the erased marks in the geometric
construction, is still necessary for descent.  In particular this does
not discharge the ordered-mesh branch, D1c, or the finite generating base.
The ledger remains $57/100$.

\paragraph{A tested but insufficient alternative.}
The preliminary attempt to explain every cofacial forced-equal edge pair
by a two-edge cut fails on the pentagonal prism: its $30$ labelled Tait
colourings force five such pairs equal, yet deleting any indicated pair
leaves the graph connected.  It has square faces, so it is not a
normal-form counterexample.  A separate exhaustive generator gate
\texttt{-a -d -m5 -c5} found no forced-equal cofacial pair on its $303$
maps through $38$ vertices.  The receipt
\texttt{results/fourcolor/v24\_cofacial\_rigidity\_gate.json}
records hashes, populations and complete dual-enumerator checks, not a
kernel proof of that stronger universal claim.  More importantly, a
reduction need not preserve this girth-five normal form: the two planar
smoothings at edge $01$ of the checked dodecahedron give $18$-vertex
spherical maps, one with two square faces and the other with two triangular
faces.  The positive
finite gate cannot be used to bypass the erased-mark transport problem.

\subsection{A smaller physical piece for one equality observation}
\label{sec:marked-square-boundary-reduction}

The complete erased-colour profile is stronger than a particular equality
test. This distinction changes the local realization problem of
Section~\ref{sec:marked-square-observable}. For the observation $x_0=w_2$,
replace the square by a digon attached to ports $0,1$, together with a wire
joining ports $2,3$. Label the two parallel edges $c_0,c_1$, the two
terminal edges $c_2,c_3$, and the separate wire $c_4$. The mark moves to
$c_0$ and the reference edge is $c_4$; these edges are nonincident.
There are two internal cubic vertices instead of four, with the same four
labelled boundary terminals. The Lean fibre is identified with nonzero
proper edge colourings of these explicit endpoint pairs, rather than
declared to be a physical state by definition alone.
\Sverified{Mettapedia.GraphTheory.FourColor.MarkedSquareBoundaryReduction.extension_iff_edgeProper}
\Sverified{Mettapedia.GraphTheory.FourColor.MarkedSquareBoundaryReduction.marked_reference_nonadjacent}
\Sverified{Mettapedia.GraphTheory.FourColor.MarkedSquareBoundaryReduction.fewer_internal_vertices}

There is an exact equivalence at every boundary word:
\[
 \{x:\operatorname{SquareExt}(w,x),\ x_0=w_2\}
 \simeq
 \{c:\operatorname{DigonWireExt}(w,c),\ c_0=c_4\}.
\]
The forward map is
$(x_0,x_1,w_0,w_1,w_2)$; its inverse is
$(c_0,c_1,w_0,c_1)$. Both validity and inverse laws are kernel checked.
The equivalence retains the entire exterior state when applied fibrewise
over its boundary word; it is not a claim that an arbitrarily frozen
exterior colouring can be changed.
\Sverified{Mettapedia.GraphTheory.FourColor.MarkedSquareBoundaryReduction.equalityEquiv}
\Sverified{Mettapedia.GraphTheory.FourColor.MarkedSquareBoundaryReduction.equality_count}
\Sverified{Mettapedia.GraphTheory.FourColor.MarkedSquareBoundaryReduction.physicalEqualityEquiv}
\Sverified{Mettapedia.GraphTheory.FourColor.MarkedSquareBoundaryReduction.exteriorEquiv}

Suppressing the remaining digon exposes the change of predicate. For
nonzero $w$ the count is exactly
\[
 N_{x_0=w_2}(w)=
 \mathbf 1\{w_0=w_1,\ w_2=w_3,\ w_0\ne w_2\}.
\]
The equality of the marked digon edge and reference wire has become
\emph{inequality} of the two surviving wire colours. In particular the
all-red boundary is an ordinary colourable smoothing with an empty
inequality target. This is not closure of the original equality-target
class under digon suppression; any enlarged state class needs its own
proved closure and generating floor.
\Sverified{Mettapedia.GraphTheory.FourColor.MarkedSquareBoundaryReduction.suppressed_count}
\Sverified{Mettapedia.GraphTheory.FourColor.MarkedSquareBoundaryReduction.zero_iff}
\Sverified{Mettapedia.GraphTheory.FourColor.MarkedSquareBoundaryReduction.all_red_suppressed_target_empty}

The constructor gate also tests the stronger realization claim. A cubic
four-terminal endpoint model with fewer than four internal vertices has
zero or two internal vertices, since $3v+4$ is even. All three matchings
of four endpoint slots and all $945$ matchings of ten slots were checked,
with every choice of marked edge. Loops, parallel edges, disconnected
pieces and nonplanar pairings were allowed, so planarity was not used to
exclude a prospective realization. No nonempty marked-colour support is
contained in the square's support. Thus no nonnegative sum of these
smaller pieces gives the complete marked profile: a positive summand
would necessarily introduce a forbidden state. This is a complete
computational gate for that endpoint model, not a new Lean theorem or a
refutation of representations with extra constraints. Vertex-free circle
components are outside the enumerated endpoint model.

The receipt is
\texttt{results/fourcolor/v24\_square\_physical\_observable\_gate.json}.
It also checks the explicit digon-and-wire drawing: adding the four
boundary-cycle edges gives a connected cubic rotation with six vertices,
nine edges and five face cycles; the boundary cycle is a face in the
specified terminal order. This embedding check is computational. The
Lean results certify incidence semantics and fibre transport, not a
general ambient spherical-splice theorem for this replacement.

The result is therefore an equality-specific physical repair and a warning
against demanding full-profile transport unnecessarily. It does not
handle an arbitrary remote reference edge, close L4 as a whole, settle the
ordered-mesh alternative, or verify the enlarged generating base. No
ordinary non-colourability claim follows from the empty refined fibre:
a colourable smaller map may have no colouring satisfying its carried
equality or inequality observation. No
obligation weight changes: the ledger remains $57/100$.

\subsection{A protected mesh carrier is constructed, not assumed}
\label{sec:protected-mesh-carrier}

The ordered-mesh branch can be separated from finite marked-edge
bookkeeping without changing the colouring predicate. Let $S$ be a finite
set of ambient vertices. For an edge set $F$, at most $|F|$ rows and at
most $|F|$ columns of an ordered injective mesh meet $F$: each blocked
path selects an edge of $F$, and within-family edge-disjointness makes
the selection injective. Discard these paths and take order-preserving
embeddings of the remaining index sets. The output retains the entire
original paths, including their tails, and the strict order of the
designated crossings. Indices need not be consecutive.
\Sverified{Mettapedia.GraphTheory.FourColor.ProtectedOrderedMesh.card_blocked_le}
\Sverified{Mettapedia.GraphTheory.FourColor.ProtectedOrderedMesh.exists_restriction_avoiding_edges}

Take $F$ to be the set of all ambient edges incident to $S$. In a cubic
rotation system $|F|\leq 3|S|$. Every vertex of a nontrivial path is
incident to one of its edges; target dimensions $p,q\geq2$ guarantee
that every retained path is nontrivial. Consequently an ordered mesh
with at least $p+3|S|$ rows and $q+3|S|$ columns contains a $p\times q$
ordered submesh whose full row and column vertex sets avoid $S$.
\Sverified{Mettapedia.GraphTheory.FourColor.ProtectedOrderedMesh.path_vertex_incident}
\Sverified{Mettapedia.GraphTheory.FourColor.ProtectedOrderedMesh.star_card_le_three}
\Sverified{Mettapedia.GraphTheory.FourColor.ProtectedOrderedMesh.exists_unmarked_submesh}

The protected set can include an actual ambient neighbourhood, including
lateral attachments outside the mesh. If $B_0(S)=S$ and
$B_{r+1}(S)=N[B_r(S)]$, the degree bound gives
$|B_r(S)|\leq |S|4^r$. Thus the dimension loss
\[
 d=3|S|4^r
\]
is sufficient to construct a $p\times q$ mesh avoiding $B_r(S)$ from a
$(p+d)\times(q+d)$ mesh. No minimality, colouring choice, frozen
exterior, supplied clean carrier, or planarity assumption is required
for this extraction.
\Sverified{Mettapedia.GraphTheory.FourColor.ProtectedOrderedMesh.card_buffer_le}
\Sverified{Mettapedia.GraphTheory.FourColor.ProtectedOrderedMesh.exists_buffer_avoiding_submesh}

Substituting the inflated dimensions into the existing spherical
mesh-or-replacement construction gives an actual protected mesh or the
existing strictly smaller spherical, bridgeless cubic zero-Count map
preserving the designated edges. The input is the same ambient class,
two-sided orbit faces, ordinary non-colourability, an alpha-closed list
of marked darts with endpoints in $S$, and the explicit size threshold
at the inflated dimensions. No clean-mesh supply hypothesis is added.
The replacement branch retains ordinary non-colourability; it is not
the empty equality-constrained fibre of the preceding subsection.
\Sverified{Mettapedia.GraphTheory.FourColor.ProtectedSphericalMeshSeam.exists_protectedMesh_or_marked_replacement}

\paragraph{The precise geometric limitation.}
A protected path carrier is not a protected disk interior. A direct
integer-coordinate control has twenty-two vertices and is cubic and
bridgeless: subdivide two opposite edges of a split-grid cell and join the
subdivision vertices through a diamond, then fill the exterior degree
deficits with an eight-cycle. The full mesh paths avoid the marked diamond
vertex, but the cell encloses it. The drawing,
path incidence, disjointness, strict interior test, edge-deletion
connectedness, and an explicit proper Tait colouring are checked in
\texttt{results/fourcolor/v24\_protected\_mesh\_gate.json}. This is a
computational cubic planar control, not a minimal-counterexample
specimen or a Lean topological theorem. Positive buffer gates cover
radii zero, one and two; an off-mesh attachment regression checks that
the buffer is ambient rather than mesh-relative.

The extraction therefore removes finite mark avoidance from the open
mesh branch, but supplies neither clean disk interiors nor complete
bounded-width interfaces nor a strict reduction on that branch.
These are separate geometric and semantic obligations. No fixed
obligation has been discharged by this refinement: the ledger remains
$57/100$.

The stronger supplier in Section~\ref{sec:protected-linked-band} repairs
the region-avoidance limitation by retaining the original contours before
the mesh projection. It does not infer a disk from arbitrary mesh paths.

\subsection{A size-independent full-vertex obstruction to generic pumping}
\label{sec:split-grid-chain-obstruction}

The mesh isoperimetric theorem controls changes in the designated branch
set, not arbitrary vertices attached to the mesh. An explicit dense family
closes that distinction without assuming strict branch growth. At each
position $(i,j)$ of an $r\times r$ grid put vertices $A_{ij},B_{ij}$ and
edge $A_{ij}B_{ij}$. Add edges $B_{ij}A_{i,j+1}$ and
$B_{ij}A_{i+1,j}$ whenever their indices exist. There are $2r^2$ vertices.
The full row and column paths alternate $A,B$ vertices and give an actual
ordered injective mesh with branches $A_{ij}$. Its vertices have degree at
most three. Colouring internal, horizontal and vertical edges with three
distinct colours is proper on the entire family.
\Sverified{Mettapedia.GraphTheory.FourColor.SplitGridMesh.ordered}
\Sverified{Mettapedia.GraphTheory.FourColor.SplitGridMesh.vertex_card}
\Sverified{Mettapedia.GraphTheory.FourColor.SplitGridMesh.edgeKind_proper}
\Sverified{Mettapedia.GraphTheory.FourColor.SplitGridMesh.degree_le_three}

Let $\delta(S)$ be the complete cut in this graph. The matching edges
give the estimate
\[
 |S|\leq 2\,|\{(i,j):A_{ij}\in S\}|+|\delta(S)|.
\]
Indeed a partner $B_{ij}\in S$ with $A_{ij}\notin S$ consumes its own
distinct crossing matching edge. If $|\delta(S)|\leq k<r$, the existing
mesh theorem says one side has at most $k^2$ branch vertices. Applying
the estimate to that side yields
\[
 \min(|S|,|S^c|)\leq 2k^2+k.
\]
\Sverified{Mettapedia.GraphTheory.FourColor.SplitGridMesh.extra_card_le_cut}
\Sverified{Mettapedia.GraphTheory.FourColor.SplitGridMesh.card_le_twice_selected_add_cut}
\Sverified{Mettapedia.GraphTheory.FourColor.SplitGridMesh.small_side_of_cut}

Consequently every strictly nested chain $S_0\subsetneq\cdots\subsetneq
S_m$ with all complete cuts of size at most $k$ satisfies
\[
 m\leq 4k^2+2k+1.
\]
This is independent of $r$ and holds for every width: when $k\geq r$,
the total vertex count already bounds the chain by $2r^2\leq2k^2$.
No condition on branch-set growth, connectedness, majority saturation,
cell order or one-vertex increments is needed. Imposing additional
admissibility requirements only restricts the allowed chains.
\Sverified{Mettapedia.GraphTheory.FourColor.SplitGridMesh.full_chain_length_le}

For $k\geq1$ the polynomial bound is strictly less than $2^{3^k}$.
Thus neither increasing $r$ nor choosing another positive width reaches
the generic pigeonhole threshold for all subsets of $3^k$ boundary words
on this family. A theorem explicitly constructs an example above every
prescribed vertex bound. At width five there are at most $112$ sets,
irrespective of wall size; this is much less than the previously audited
$3{,}809{,}280$ state-carrier bound for nested good sides. That numerical
comparison is not an identification of the open-wall carrier with the
minimal-counterexample class.
\Sverified{Mettapedia.GraphTheory.FourColor.SplitGridMesh.polynomial_lt_support_carrier}
\Sverified{Mettapedia.GraphTheory.FourColor.SplitGridMesh.no_full_support_pigeonhole_chain}
\Sverified{Mettapedia.GraphTheory.FourColor.SplitGridMesh.width_five_chain_sets}
\Sverified{Mettapedia.GraphTheory.FourColor.SplitGridMesh.arbitrarily_large_no_full_support_chain}

The pre-formalization gate checks every vertex subset at side lengths
one, two and three, with every distinct width threshold. A Boolean-lattice
dynamic program computes the longest valid nested chain while allowing
jumps over invalid intermediate sets. Its recurrence is independently
compared with all orders for every validity assignment on a three-element
ground set. Boundary counts are compared with direct edge scans. The
receipt is \texttt{results/fourcolor/v24\_split\_grid\_chain\_gate.json};
the all-size Lean bound does not trust the optimizer.

\paragraph{Scope of the obstruction.}
These are colourable subcubic open walls, not closed cubic spherical
minimal counterexamples. The cut is complete in the explicitly defined
graph. Restoring external terminal incidences can only increase the cost
of a fixed side; arbitrary nested growth in an attached exterior is not
bounded by this theorem. No rotation-system sphericity seal is claimed
for this family. Most importantly, failure to reach a worst-case
pigeonhole threshold does not mean that actual supports cannot repeat
earlier. Stronger semantic compression, early stabilization,
non-pumping reductions, or a theorem genuinely using non-colourability
remain possible. What is excluded is rescuing this generic full-support
counting argument merely by adjusting size and width. The fixed ledger
remains $57/100$.

\subsection{Physical contexts do not by themselves give a small exact code}
\label{sec:physical-context-bits}

The unrestricted support lower bound uses singleton suffixes which need
not be physically realizable. A different construction shows that an
exponential lower bound survives for a particular family of literal
physical contexts, without frozen colours or abstract singleton tests.
It does not establish such a lower bound for the admissible
minimal-counterexample contexts.

Use four ports in cyclic order. Piece $A$ is two digons, pairing ports
$01\mid23$; piece $B$ pairs $03\mid12$. Both pieces have four cubic
vertices and four internal edges. Piece $T$ is a triangle with a stem
to a fourth vertex: ports $0,1$ attach to two triangle vertices and
ports $2,3$ to the stem's far endpoint. The remaining triangle vertex
meets the stem. Its four interior edges also join distinct vertices.
Explicit dart involutions and cyclic vertex rotations construct all
three as \texttt{OpenTangleData}.

An $A$-extension has $w_0=w_1$, since each digon forces its two port
colours equal. A $T$-extension has $w_0\ne w_1$, by the triangle's
proper-colouring constraints. Their supports cannot meet. Conversely
$B$ and $T$ have a common extension, for example over boundary word
$(r,b,b,r)$. Both $A$ and $B$ can be closed colourably by an $A$ piece.
These claims are proved on the actual dart-incidence predicates.
\Sverified{Mettapedia.GraphTheory.FourColor.PhysicalContextBits.a_rejected}
\Sverified{Mettapedia.GraphTheory.FourColor.PhysicalContextBits.b_accepted}
\Sverified{Mettapedia.GraphTheory.FourColor.PhysicalContextBits.neutral_exists}

For $s\in\{0,1\}^n$, let $P_s$ be the disjoint union of $n$ pieces,
using $A$ or $B$ according to each bit. Let context $C_i$ use $T$ at
block $i$ and $A$ everywhere else. Reflect the exterior rotations before
gluing the identically labelled ports. The resulting literal sewn
rotation system is cubic, and
\[
  \operatorname{TaitColorable}(C_i\circ P_s)
  \quad\Longleftrightarrow\quad s_i=1.
\]
The proof transports the block colourings through the existing physical
Count/gluing equivalence; it is not a new uninterpreted acceptance
predicate. Every entire boundary word is retained. There are $4n$ ports,
$4n$ vertices on each side, and no vertex-free wire components.
\Sverified{Mettapedia.GraphTheory.FourColor.PhysicalContextBits.block_realizes_iff}
\Sverified{Mettapedia.GraphTheory.FourColor.PhysicalContextBits.physical_test}
\Sverified{Mettapedia.GraphTheory.FourColor.PhysicalContextBits.sewn_cubic}
\Sverified{Mettapedia.GraphTheory.FourColor.PhysicalContextBits.port_card}

If a deterministic code answers all these physical tests exactly, equality
of codes implies equality of every bit. Thus it needs at least $2^n$
states at width $4n$. This is smaller than the full powerset bound but
still exponential. For $n\ge16$ it already exceeds
$4(4n)^2+2(4n)+2$, the full-set chain bound from the split-grid family.
The two families are different: this comparison limits a \emph{uniform}
exact code for all these pieces, not a quotient restricted to the states
actually attained on one split-grid sweep.
The distinction is about the number of semantic states, not their storage:
the family is represented by just $n$ bits. Symbolic representations may
remain efficient without reducing the number relevant to pigeonholing.
\Sverified{Mettapedia.GraphTheory.FourColor.PhysicalContextBits.exact_code_injective}
\Sverified{Mettapedia.GraphTheory.FourColor.PhysicalContextBits.two_pow_le_card_code}
\Sverified{Mettapedia.GraphTheory.FourColor.PhysicalContextBits.exponential_exceeds_quadratic}

\paragraph{The admissibility boundary is essential.}
The rejecting $A$--$T$ closure has a bridge (the stem). Other neutral
blocks may be separate components. Accordingly these tests do not belong
to a family of connected bridgeless minimal-counterexample splices.
Planarity alone does not imply this stronger admissibility, and the
lower bound must not be transferred to that class. A target-restricted
quotient, an exterior-dependent replacement, or a one-sided support
simulation remains possible. Each requires its own constructive
soundness and physical class-preservation proof; exact coding of every
planar physical context is a stronger requirement.

The pre-formalization gate checks the three ordered-disc embeddings by
adding a terminal cycle and verifying the explicit dart rotation and
Euler count, then checks the four elementary closures componentwise.
Their colouring counts are respectively $144,48,0,24$ for $AA,BA,AT,BT$;
the direct closed-edge search agrees with independent composition of
local extensions. Only $AT$ has a bridge. A crossing port-order control
fails the disc check. Two-block closures are also searched directly, and
the generated bit tests are checked through six bits. These geometric
embedding receipts are computational checks, not a new general Lean
planarity theorem. The generic physical colourability and code bounds
above are kernel theorems. Receipt:
\texttt{results/fourcolor/v24\_planar\_context\_code\_gate.json}.
This narrows the semantic-compression alternative but discharges no
global reduction or generating-base obligation; the ledger remains
$57/100$.

\subsection{An adaptive conjunction need not be a physical replacement}
\label{sec:physical-support-intersection}

There is a further realization issue in a proposed target-aware repair.
Suppose every word of the actual exterior is rejected by at least one of
several smaller candidates. Then the exterior misses the intersection of
their supports. A tempting construction is to realize that intersection by
one new tangle, sharing the same boundary variables between the candidates.
The following obstruction rules out this \emph{generic} compilation step.
It is not a claim that the source manuscript requires such a step.

\begin{theorem}[no literal colour fanout]
Let $T$ be a finite cubic port tangle with at least one proper Tait
colouring. If $p,q,r$ are three distinct boundary ports, then they are not
equal in every proper colouring of $T$. More generally, every equivalence
class of ports forced equal in all proper colourings has cardinality at
most two.
\Sverified{Mettapedia.GraphTheory.FourColor.PhysicalColorNoFanout.not_three_forced_equal}
\Sverified{Mettapedia.GraphTheory.FourColor.PhysicalColorNoFanout.forcedEqual_class_card_le_two}
\end{theorem}

\begin{proof}
If some proper colouring already separates the three ports, there is
nothing to prove. Otherwise let their common colour be $a$ and choose a
different nonzero colour $b$. The physical $(a,b)$ Kempe component through
$p$ has exactly two boundary endpoints. It cannot contain both $q$ and
$r$. Switch that whole component. Properness is preserved, $p$ changes
colour, and at least one of $q,r$ retains colour $a$. This is an actual
colouring construction on the common dart carrier, not a separately chosen
boundary matching.
\Sverified{Mettapedia.GraphTheory.FourColor.PhysicalColorNoFanout.exists_split_coloring}
The cardinality assertion follows by choosing three ports from any larger
forced-equality class.
\end{proof}

The interior can be arbitrarily large. No planarity or connectedness is
needed. Additional \emph{free} boundary ports do not help: the second
endpoint of the switched component may be one of them. Nonemptiness is
essential, since an uncolourable tangle satisfies every universal equality
vacuously. The assertion is about individual boundary-edge colours, not
encoded signals, internal-edge observations, or a family restricted by
frozen exterior colours.

\begin{theorem}[nonempty physical supports are not closed under conjunction]
Let $A$ and $B$ be the two literal four-port digon pairings of
Section~\ref{sec:physical-context-bits}. Their support intersection is
nonempty, but no finite cubic four-port tangle $R$ with nonempty support
satisfies
\[
       \operatorname{Supp}(R)\subseteq
       \operatorname{Supp}(A)\cap\operatorname{Supp}(B).
\]
In particular, no such tangle realizes the intersection exactly, regardless
of its number of internal vertices.
\Sverified{Mettapedia.GraphTheory.FourColor.PhysicalSupportIntersection.intersection_nonempty}
\Sverified{Mettapedia.GraphTheory.FourColor.PhysicalSupportIntersection.no_nonempty_support_below_intersection}
\Sverified{Mettapedia.GraphTheory.FourColor.PhysicalSupportIntersection.no_exact_intersection}
\end{theorem}

\begin{proof}
Both pairings accept the all-red word, with their parallel internal edges
coloured blue and purple. An $A$-extension forces $w_0=w_1$; a
$B$-extension forces $w_1=w_2$. Hence every common word has three equal
designated ports. The no-fanout theorem rules out every nonempty physical
support contained in their intersection.
\Sverified{Mettapedia.GraphTheory.FourColor.PhysicalSupportIntersection.mem_support_iff}
\Sverified{Mettapedia.GraphTheory.FourColor.PhysicalSupportIntersection.intersection_forces_three_equal}
\end{proof}

This obstruction concerns the class allowing disconnected pieces, as does
the source's disjoint-union operation: $A$ and $B$ themselves are
disconnected. No analogous connected-source claim is made. The conclusion
is stronger than failure of one copying gadget but weaker than failure of
all adaptive reduction: a replacement for one actual exterior need not
realize this intersection at all.

The constructor gate uses explicit colourings, not a boundary-word census.
The two-port digon is the positive equality control; even cycles with
$4,6,8,20$ ports and two disjoint digons exhibit the two-ended switch.
Tests vary the root and colour names, check the inverse switch and free
auxiliary ports, and reject a non-cubic copying vertex. Receipt:
\texttt{results/fourcolor/v24\_copy\_gate.json}.

The proof-plan correction is to keep physical realizability alongside
every proposed operation on supports. Disjoint union tensors independent
interfaces; gluing consumes matched ports. Neither supplies a diagonal
copying map or arbitrary intersection. The protected-mesh branch still
requires an actual smaller admissible zero-Count map. This failed
construction discharges no fixed-ledger obligation; the total stays
$57/100$.

\subsection{Why exposing a closed seam does not give free-port rigidity}
\label{sec:internal-edge-opening}

An attempted use of Section~\ref{sec:physical-support-intersection} in
the physical mesh reduction is to open two seam edges of a smaller map,
transfer their forced colour relation to the new free ports, and invoke
no-fanout. The middle implication is false. Internal-edge observations
and free-port observations have different quantifiers.

The literal triangular prism has vertices $0,\ldots,5$, triangle edges
$01,12,20$ and $34,45,53$, and joining edges $03,14,25$.
Every proper three-edge-colouring assigns $01$ and $34$ the same colour.
The proof uses only uniqueness of the third colour at a cubic vertex;
the graph has an explicit proper colouring.
\Sverified{Mettapedia.GraphTheory.FourColor.InternalEdgeOpening.closed_forces_equal}
\Sverified{Mettapedia.GraphTheory.FourColor.InternalEdgeOpening.closedWitness_proper}

Cut $01$ and $34$ into four terminal edges, ordered by their incident
vertices $0,1,3,4$. Retain all seven other edges. The free opening has a
proper colouring with terminal word $(0,1,1,0)$, in the three-colour
palette $\{0,1,2\}$ (these are colour names, not a zero flow value).
In particular the two chosen half-edges at $0$ and $3$ no longer agree.
\Sverified{Mettapedia.GraphTheory.FourColor.InternalEdgeOpening.escape_proper}
\Sverified{Mettapedia.GraphTheory.FourColor.InternalEdgeOpening.escape_separates_signals}
\Sverified{Mettapedia.GraphTheory.FourColor.InternalEdgeOpening.free_opening_does_not_preserve_rigidity}

There is an exact positive statement. Cutting and rejoining give a
bijection between closed proper colourings and proper colourings of the
opening \emph{subject to} $w_0=w_1$ and $w_2=w_3$. Within this conditioned
family the original forced equality still holds. The displayed escape
violates reclosure and is not the cut image of any closed colouring.
\Sverified{Mettapedia.GraphTheory.FourColor.InternalEdgeOpening.colouringEquiv}
\Sverified{Mettapedia.GraphTheory.FourColor.InternalEdgeOpening.reclosable_forces_equal}
\Sverified{Mettapedia.GraphTheory.FourColor.InternalEdgeOpening.escape_not_in_cut_image}

The constructor gate checks the triangular and pentagonal prisms using
two independent whole-graph colouring enumerators. Their numbers of
labelled colourings are respectively $6$ and $30$; opening the two
selected edges gives $30$ and $102$, of which precisely $6$ and $30$
reclose. Opening just one edge is a positive control: it gives no extra
colourings. The gate also verifies explicit spherical rotations,
bridge-freeness, cofaciality of the chosen edges, connectedness after the
two-edge deletion, and properness of a free-path switch. Switching the
whole corresponding closed component preserves the forced equality.
Receipt: \texttt{results/fourcolor/v24\_internal\_edge\_opening\_gate.json}.

The Lean witness verifies literal incidence and the cut/glue bijection;
the drawing and the two numerical counts are computational gates, not
new kernel planarity or counting theorems. These colourable controls have
short faces, so they do not refute an implication genuinely using the
protected mesh or minimal-counterexample hypotheses. The existing
adjacent-pair surgery and its spherical class theorems are reused, not
reproved. This attempted passage from internal rigidity to free-port
rigidity supplies no strict ordinary zero-Count reduction. Such a
reduction still needs to respect the actual reclosure equations; the
ledger remains $57/100$.

\subsection{Retaining complete contours supplies a protected linked band}
\label{sec:protected-linked-band}

The geometric supplier already constructs more than its ordered-mesh
output records: complete nested distance contours and vertex-disjoint
ambient paths between them. Projection to an abstract mesh forgets the
facial regions. Instead of reconstructing simple cell rims from that
lossy output, retain the constructed regions and linkage.

Let $R_d$ be the connected component of the far face among dual faces
at distance greater than $d$ from the root, and let $F(v)$ be the set
of actual face sectors incident at $v$. Define the \emph{closed vertex
band}
\[
 B_{\ell,h}=\{v:F(v)\cap R_\ell\ne\varnothing,
                      \ F(v)\setminus R_h\ne\varnothing\}.
\]
It includes both complete boundary contours. If $h<p$, then
$B_{\ell,h}$ and $B_{p,q}$ are vertex-disjoint. A vertex in both would
be mixed at depth $h$ yet touch the deeper region $R_p$, contradicting
the cubic incident-face distance lemma. Thus $m+1$ separated bands
contain a wholly unmarked one for any set of $m$ marked vertices.
This protects the material between the contours, not just their paths.
\Sverified{Mettapedia.GraphTheory.FourColor.ProtectedContourBand.bands_disjoint}
\Sverified{Mettapedia.GraphTheory.FourColor.ProtectedContourBand.exists_unmarked_band}

Every ambient path between the extreme contours can be shortened to a
path inside any chosen intervening closed band, using only vertices of
the original path. Minimize length among candidate paths on that support.
An outer excursion gives a later outer-contour encounter; an inner
excursion gives an earlier inner-contour encounter. Either produces a
shorter candidate. Applying this independently to the disjoint columns
retains their number and their vertex-disjointness.
\Sverified{Mettapedia.GraphTheory.FourColor.ProtectedContourBand.exists_confined_path}
\Sverified{Mettapedia.GraphTheory.FourColor.ProtectedContourBand.exists_confined_linkage}

Every original edge leaving the band has its inside endpoint on one of
the two actual contours. A confined column joins these connected
contours; together they form a connected seed containing the whole
inside boundary. An ambient walk from any band vertex first reaches
this seed without leaving the band, proving that the entire induced
band, including lateral material, is connected.
\Sverified{Mettapedia.GraphTheory.FourColor.ProtectedContourBand.boundary_based_on_contours}
\Sverified{Mettapedia.GraphTheory.FourColor.ProtectedContourBand.band_connected}

\paragraph{The strengthened large-map alternative.}
Re-run the existing dual-radius, contour-linkage and cotree construction,
returning the linkage source whenever it occurs, instead of first
testing whether an arbitrary ordered mesh exists. If no window supplies
that linkage, the complete ordered bonds in all windows give the same
repeated full marked seam as before.
\Sverified{Mettapedia.GraphTheory.FourColor.SphericalProtectedLinkedBand.exists_linkedContours_or_seamPair}

For $a\geq2$, put $A=(|S|+1)a$. At the existing size threshold
\texttt{sizeBound}$(A,b,|S|,|M|)$, a spherical bridgeless cubic
ordinary zero-Count map with two-sided orbit faces supplies either:
(i) a wholly $S$-avoiding connected closed band with $a$ complete
consecutive contours and \texttt{orderedLinkageSize}$(A,b)$ confined
disjoint columns; or (ii) the existing strictly smaller spherical
bridgeless cubic ordinary zero-Count map preserving the alpha-closed
marked dart pairs. The band is constructed from the source hypotheses;
it is not a supplied-region premise.
\Sverified{Mettapedia.GraphTheory.FourColor.SphericalProtectedLinkedBand.exists_linkedBand}
\Sverified{Mettapedia.GraphTheory.FourColor.SphericalProtectedLinkedBand.exists_linkedBand_or_marked_replacement}
\Sverified{Mettapedia.GraphTheory.FourColor.SphericalProtectedLinkedBand.LinkedBand.complete_contour}
\Sverified{Mettapedia.GraphTheory.FourColor.SphericalProtectedLinkedBand.LinkedBand.connected}

The gate independently replays the eight existing spherical
$GP(k,0)$ controls, $1\leq k\leq8$, and checks $6330$ trimmed column
paths, complete ambient cut edges, and whole-band mark avoidance.
Boundary-only avoidance and windows sharing an endpoint contour are
negative controls. Receipt:
\texttt{results/fourcolor/v24\_protected\_contour\_band\_gate.json}.

This changes the input available for the remaining geometric surgery,
not its semantic conclusion. No fixed contour-width bound, periodicity,
support equality, or deletion of the linked band is proved. In particular
a two-boundary band is not automatically a connected-complement Count
shore; its complement can have two or more components. The next reduction must
use these actual annular regions and full interfaces without dropping
reclosure constraints. No fixed-ledger obligation is fully discharged:
the total remains $57/100$.

\subsection{Buffered absorption constructs connected caps and full ordered interfaces}
\label{sec:buffered-contour-shell}

There is a second geometric distinction: two complete nested contours
do not imply that the complement of their closed vertex band has exactly
two connected components. The gate records a spherical cubic graph on
$36$ vertices with only pentagonal and hexagonal faces. Its dual remains
connected after every deletion of at most four vertices, checked directly
in $6196$ cases. The band between depths $1$ and $2$ from root face $13$
towards far face $2$ has three complementary components of sizes $1,2,5$,
despite $8$ vertex-disjoint ambient paths between its complete contours
of lengths $13$ and $15$. Thus the raw band cannot simply be read as
a serial piece between two connected caps. This is a computational
geometric counterexample, not a kernel-certified literal map or a
counterexample to an implication using ordinary zero Count.

\paragraph{The repaired construction.}
Take a protected linked band $B_{\ell,h}$ with at least four consecutive
contours and at least one column. Its middle is
$S=B_{\ell+1,h-1}$. Choose vertices $x$ and $y$ on the buffered extreme
contours, and let $C,D$ be their respective connected components in
$G-S$. Set
\[
                    H=V(G)\setminus(C\cup D).
\]
Every ambient $x$--$y$ path crosses an intermediate complete contour,
so $C$ and $D$ are disjoint. The confined column connects the entire
middle band. Every other component of $G-S$ attaches to $S$, so absorbing
those components makes $H$ connected. Each cap and its complement are
connected as well.
\Sverified{Mettapedia.GraphTheory.FourColor.BufferedContourShell.floods_disjoint}
\Sverified{Mettapedia.GraphTheory.FourColor.BufferedContourShell.shell_connected}

Absorption must not swallow protected material. Vertices touching
$R_\ell^c$ form a connected set containing the outer contour: every
inside endpoint of an edge leaving that set is mixed at depth $\ell$.
The analogous statement holds for vertices touching $R_h$. The cubic
incident-face distance lemma separates both sets from $S$. They therefore
lie entirely in $C$ and $D$, respectively. Every vertex outside the
buffer belongs to one of those touched sets. Consequently
\[
                 B_{\ell+1,h-1}\subseteq H\subseteq B_{\ell,h}.
\]
This proves buffer containment rather than adding it as an assumption.
No original edge joins the two caps; every shell-side endpoint of a
shell cut edge lies in $S$. Marked vertices avoid $H$, and all original
neighbours of each marked vertex lie in a single cap.
\Sverified{Mettapedia.GraphTheory.FourColor.BufferedContourShell.touches_disjoint_compl}
\Sverified{Mettapedia.GraphTheory.FourColor.BufferedContourShell.exists_shell}
\Sverified{Mettapedia.GraphTheory.FourColor.BufferedContourShell.ShellOf.unmarked}
\Sverified{Mettapedia.GraphTheory.FourColor.BufferedContourShell.ShellOf.marked_star}

\paragraph{The actual interfaces, not a selected-path boundary.}
Apply the existing planar-bond theorem to the shores $C$ and $D^c$.
It supplies unique retained boundary darts on crossing faces and the
opposite complementary first-return orders. The resulting bonds lie
between depths $(\ell,\ell+1)$ and $(h-1,h)$; their vertex shores and
full incident-edge shores are strictly nested. Their widths are the
actual sizes $|\delta C|$ and $|\delta D^c|$, with no uniform bound
inferred from the chosen number of columns. The existing cycle-based
adapter then supplies two genuine connected Count nodes. This does
not assert that majority vertex sides equal the cap vertex sets.
\Sverified{Mettapedia.GraphTheory.FourColor.BufferedContourShell.ShellOf.exists_ordered_bonds}
\Sverified{Mettapedia.GraphTheory.FourColor.BufferedContourShell.ShellOf.exists_connected_nodes}

The global alternative of Section~\ref{sec:protected-linked-band}, with
$a\geq4$, now constructs this buffered shell in its linked-band branch;
the other branch remains the existing strict ordinary zero-Count
replacement preserving the marked edges. There is no new shell-supply
or cap-connectedness hypothesis.
\Sverified{Mettapedia.GraphTheory.FourColor.BufferedContourShell.exists_bufferedShell_or_marked_replacement}

The gate checks $1388$ buffered windows: $672$ on the eight existing
$GP(k,0)$ controls and $716$ on a $144$-vertex triangular refinement of
the irregular specimen, over all $74$ root faces. The latter has twelve
pentagons and $62$ hexagons. Twelve windows require absorption, with
$24$ absorbed-vertex occurrences in total; all stay in their buffers.
Complete ambient cut edges, both cap complements, and the largest
allowed marked set $V(G)\setminus B_{\ell,h}$ are checked. Receipt:
\texttt{results/fourcolor/v24\_buffered\_shell\_gate.json}.

This repairs the geometric carrier of the intended composition, not
its shrinking law. Neither equal boundary support nor a smaller
reclosable replacement is proved. A reduction must still preserve
ordinary zero Count, the spherical cubic bridgeless class, and marked
edges on these full interfaces. The fixed ledger remains $57/100$.

\subsection{A transverse shell sweep must account for both caps}
\label{sec:unsplit-cap-sweep}

The buffered construction does not justify treating a narrow shell as
a bounded-interface corridor while keeping its caps intact. There is
an all-orders obstruction to that additional restriction, even for the
following elementary closed family. For $n\geq3$, take four cyclic rows
$O_i,D_i,E_i,I_i$, indexed modulo $n$, with the six edge families
\[
 O_iO_{i+1},\quad O_iD_i,\quad D_iE_i,\quad
 E_iD_{i+1},\quad E_iI_i,\quad I_iI_{i+1}.
\]
The middle rows form a single zigzag ring; the outer and inner rows are
the two caps. The width measured here is the number of \emph{all original
edges} crossing a vertex cut, not the cut restricted to the ring and
not a branch-decomposition middle-vertex measure.

\paragraph{The obstruction.}
Let $S$ contain exactly $n$ of the $2n$ shell vertices. Suppose each cap
is wholly in $S$ or wholly outside it, allowing all four choices.
Then $|\delta S|\geq n$. If both caps are outside, charge the $n$
included shell vertices to their distinct cap spokes. If both are
inside, charge the $n$ excluded vertices instead. If the caps are on
opposite sides, each of the $n$ edge-disjoint paths
$O_i,D_i,E_i,I_i$ contains a crossing edge. These are proved injections
into the actual full cut.
\Sverified{Mettapedia.GraphTheory.FourColor.UnsplitCapSweep.outside_caps}
\Sverified{Mettapedia.GraphTheory.FourColor.UnsplitCapSweep.inside_caps}
\Sverified{Mettapedia.GraphTheory.FourColor.UnsplitCapSweep.opposite_caps}
\Sverified{Mettapedia.GraphTheory.FourColor.UnsplitCapSweep.balanced_cut}

Consequently the middle prefix of \emph{every} shell-vertex ordering
has width at least $n$, even if whole caps change their ownership at
every stage. Explicit prefix cuts realize these hypotheses, giving an
unboundedness theorem quantified over every ordering and every
stagewise cap assignment, not a conditional assertion about a supplied
sweep.
\Sverified{Mettapedia.GraphTheory.FourColor.UnsplitCapSweep.sweep_width_lower_bound}
\Sverified{Mettapedia.GraphTheory.FourColor.UnsplitCapSweep.unbounded_unsplit_cap_sweeps}

\paragraph{A positive escape control on the same family.}
Include all four vertices in columns $0\leq i<t$, for $0<t<n$.
Now both caps participate in the cut. Only edge families $0,3,5$
can cross, each at columns $t-1$ and $n-1$, so the full cut has at most
six edges, independently of $n$.
\Sverified{Mettapedia.GraphTheory.FourColor.UnsplitCapSweep.slice_width_le_six}
This is why the preceding lower bound is not a claim of unbounded
optimal branchwidth. It rules out the unsplit-cap restriction, not
coordinated cuts through the caps.

The gate exhausts all half-shell subsets and four cap assignments for
$3\leq n\leq8$: $70272$ full cuts, with exact minima
$4,6,6,8,8,10$. It checks explicit oriented faces, spherical Euler
counts, simplicity, cubicity, bridgelessness, connected caps and
connected cap complements. The ring-only consecutive sweep has width
two, while the whole-graph column slices have width six. Receipt:
\texttt{results/fourcolor/v24\_unsplit\_cap\_sweep\_gate.json}.
The parametric Lean theorems concern the displayed endpoint graph and
its full cuts; they do not additionally formalize its rotation system.
The finite geometry checks are not an all-orders kernel proof of
sphericity. Nor is this a family of least counterexamples or a
bounded-face family: its two cap faces have length $n$.

The resulting design correction is precise. Preserving designated
marks does not require preserving \emph{all} unmarked cap material as
an indivisible geometric object. A prospective general shortening may
cut through that material, but must then prove its complete-interface
bound, reclosure conditions, colouring lift and strict decrease. The
six-edge slices above do not supply those assertions for arbitrary
caps, nor do they discharge a source-facing reduction obligation.
The fixed ledger remains $57/100$.

\subsection{Optimal complete path sweeps have connected shores and derived boundary orders}
\label{sec:canonical-terminal-bonds}

Permitting caps to split must give an actual family of complete
interfaces, not a new assumption that suitable cap partitions exist.
The canonical terminal cuts already constructed earlier minimize the
number of original crossing edges, with side cardinality breaking ties.
They nest as source terminals grow and sink terminals shrink. The
remaining connectivity condition can now be proved.

\paragraph{Every minimum cut is a bond.}
Let $A,B$ be disjoint nonempty connected vertex sets in a finite connected
graph. Let $S$ minimize its full edge boundary among sides containing
$A$ and avoiding $B$. Flood from $A$ using only vertices of $S$, giving
$R\subseteq S$. Every edge leaving $R$ also leaves $S$, since an edge
to another vertex of $S$ would extend the flood. Thus
$\delta R\subseteq\delta S$. Minimality gives equality of these full
boundaries. In a connected graph, equal edge boundaries determine a
vertex side up to complementation: equality of membership propagates
along a walk once fixed at a shared vertex. A vertex of $A$ fixes it,
so $R=S$ and $S$ is connected. Applying the same argument with $A,B$
exchanged proves $S^c$ connected. No tie-break is needed for this part.
\Sverified{Mettapedia.GraphTheory.FourColor.CanonicalTerminalBonds.eq_of_boundary_eq}
\Sverified{Mettapedia.GraphTheory.FourColor.CanonicalTerminalBonds.minimum_bond}
\Sverified{Mettapedia.GraphTheory.FourColor.CanonicalTerminalBonds.canonical_bond}

The endpoint multigraph retains edge multiplicity; simple-graph
adjacency is used only for walks through the proved rotation adapter.
For a spherical cubic map with two-sided orbit faces, the existing
planar-bond theorem now supplies the unique retained dart on each
crossing face and opposite complementary first-return orders on the
\emph{constructed} canonical cut. Advancing terminals also gives strict
nesting on the full original incident-edge shores.
\Sverified{Mettapedia.GraphTheory.FourColor.SphericalCanonicalTerminalBonds.full_boundary_orders}
\Sverified{Mettapedia.GraphTheory.FourColor.SphericalCanonicalTerminalBonds.full_shores_strict}

\paragraph{No terminal-connectivity supplier is needed for a path sweep.}
Given an actual simple path $p$ and $0\leq t<\operatorname{length}(p)$,
take its first $t+1$ vertices as $A_t$ and the remaining vertices as
$B_t$. Both sets are connected, disjoint and nonempty. Canonical finite
minimization then constructs $S_t$, allowing \emph{all} off-path
vertices to move, including cap material. The construction returns
the minimum-width property, connected shores, complete boundary orders,
and strict nesting of the full incident-edge shores for every $i<j$.
A complete simple contour cycle supplies the path by removing its first
edge, using the existing \texttt{Walk.IsCycle.isPath\_tail} theorem.
\Sverified{Mettapedia.GraphTheory.FourColor.SphericalPathBondSweep.side_spec}
\Sverified{Mettapedia.GraphTheory.FourColor.SphericalPathBondSweep.side_strict}

The gate checks $2574$ connected-terminal pairs on $K_4$, a six-cycle
and the cube, examining all $3656$ optimal sides and cross-checking the
canonical choice against the flow solver. Dropping terminal or ambient
connectedness gives explicit negative controls. It also constructs and
certifies $148$ strictly nested full cuts on eight spherical controls.
The ring-and-caps controls have maximal optimum widths $4,6,6,6$ at
circumferences $3,5,8,12$. On the selected middle contours of
$GP(k,0)$ for $1\leq k\leq4$, the corresponding widths are
$6,11,16,21$. Equal-sized edge-disjoint terminal path packings certify
optimality, and the complete dual cycles certify the boundary orders.
Receipt: \texttt{results/fourcolor/v24\_canonical\_bond\_gate.json}.
The receipt retains the maximum-width cut and packing for each control;
the replay checks all intermediate cuts.

This constructs optimally split-cap geometry but does not remove its
cost. In particular, the width-$21$ packing prevents a narrower cut for
that \emph{specified terminal split}; it is not a lower bound over
every choice of contour, terminal placement, or branch decomposition.
Nor are arbitrary path shores asserted to be the good cyclic sides
required by every Count-node adapter. No uniform width, sufficient
state repetition, marked-edge preservation through a splice, or strict
ordinary zero-Count replacement follows from this construction alone.
Those target-facing conditions remain to be proved, and the fixed
ledger remains $57/100$.

\subsection{Buffered path cuts reach Count descent, or force a wide terminal separation}
\label{sec:path-count-sweep}

The connection from the preceding constructed bonds to the Count-node
consumer can be made without assuming cyclicity of each side. Endpoint
cuts must first be excluded: even an optimal connected bond of width
three on the cube can isolate a vertex. Connectedness and optimality
alone therefore do not supply the cycles required by the node adapter.

For a nonempty connected vertex side $S$ in a cubic graph, an acyclic
induced side is a tree, and degree counting gives
$|\delta S|=|S|+2$. The existing cubic-side counting theorem consequently
implies that $|\delta S|\leq |S|$ forces a cycle in $S$. The same applies
to the complement. The cut used here is the set of \emph{all} original
crossing edges; its equality with the boundary minimized by the
rotation-to-multigraph adapter is proved.
\Sverified{Mettapedia.GraphTheory.FourColor.SphericalPathCountSweep.cut_eq_boundaryEdges}
\Sverified{Mettapedia.GraphTheory.FourColor.SphericalPathCountSweep.cycle_of_card}

Let $p$ be an actual simple path of length $L$. Its prefix at index $t$
contains $t+1$ distinct vertices, and its suffix contains $L-t$.
For the canonical full cut $S_t$, suppose $|\delta S_t|\leq w$,
$w\leq t+1$, and $w\leq L-t$. Feasibility puts the prefix in $S_t$
and the suffix in $S_t^c$. Their cardinalities force cycles on both
sides. The existing cyclic-cut adapter then constructs a connected
Count node with boundary and middle bounds at most $w$, deriving its
majority-region nonemptiness and edge-shore connectedness. Its shore
is exactly the original incident-edge shore of $S_t$; no new majority
identity for the original vertex partition is assumed.
\Sverified{Mettapedia.GraphTheory.FourColor.SphericalPathCountSweep.realization}
\Sverified{Mettapedia.GraphTheory.FourColor.SphericalPathCountSweep.node}
\Sverified{Mettapedia.GraphTheory.FourColor.SphericalPathCountSweep.node_shore}

Write the existing exact phased-state bound as
\[
 Q(w)=(6w+1)\sum_{j=0}^{w} j!\,2^{3^j}.
\]
In a graph-backed vertex-minimal Tait counterexample, if every canonical
cut of $p$ has full width at most $w$, then
\[
 L\leq Q(w)+2w.
\]
Indeed, discard an endpoint buffer and use the $L-2w$ interior cuts.
The preceding construction supplies their actual Count nodes, while
canonical nesting supplies strict incident-edge-shore nesting. The
already proved literal-shore chain theorem bounds their number by
$Q(w)$. This does not replace the physical splice with an abstract
repeated-interface assumption: the existing exact phased-state
descent is the consumer.
\Sverified{Mettapedia.GraphTheory.FourColor.SphericalPathCountSweep.buffered_length_le}
\Sverified{Mettapedia.GraphTheory.FourColor.SphericalPathCountSweep.length_le}

Consequently, if $L>Q(w)+2w$, there is an index $t$ such that
\emph{every} full cut containing the path prefix and avoiding its suffix
has more than $w$ crossing edges. Canonical optimality upgrades the
conclusion from the chosen side to all feasible cap assignments.
\Sverified{Mettapedia.GraphTheory.FourColor.SphericalPathCountSweep.exists_wide_separation}

The geometric gate checks all $166$ nonempty proper connected sides of
the cube: $100$ tree sides obey the exact tree equation and $66$ sides
contain cycles. On the eight preceding spherical path controls,
$18$ cuts meet the conservative terminal-size criterion. All such cuts
have cycles on both sides. All $47$ actually cyclic cuts pass the full
incident-edge-shore connectivity, majority nonemptiness, middle-width,
and strict-nesting checks. The width-three endpoint cut is retained as
a negative control. No colourings or configuration certificates are
enumerated. Receipt:
\texttt{results/fourcolor/v24\_path\_count\_gate.json}.

This completes the narrow-cut implication for the constructed path
sweep. A complete source contour supplies such a path by deleting its
first edge. It does not prove that a sufficiently large source map
supplies a long \emph{narrow} path sweep, nor that a forced wide
separation gives a marked colour-lifting reduction. That latter
geometric-to-semantic implication remains open. The formula is the
existing exact-state carrier, not a new small-constant claim; no
fixed-ledger obligation is discharged, and the ledger remains $57/100$.

\subsection{A common-cycle lift is stronger than colourability}
\label{sec:smoothing-cycle-obstruction}

A concrete attack on the remaining wide case deletes an adjacent pair
and smooths each of its two degree-two vertices. Its inverse subdivides
the two new edges and joins the subdivision vertices. If a colouring
of the smaller graph puts both new edges on one bichromatic cycle,
the inverse insertion can be coloured by alternating along the expanded
cycle and assigning its third colour to the inserted edge. This is a
sufficient lift, not an automatic consequence of colourability.

The generic assertion that two nonincident cofacial edges can always
be put on one bichromatic cycle is false. The closed control in
\texttt{results/fourcolor/v24\_smoothing\_cycle\_gate.json} has $18$
vertices, $27$ edges and $11$ faces. Its designated edges are
$\{3,9\}$ and $\{4,10\}$, indices $8$ and $9$ in the recorded edge
list. They lie on the common facial cycle
$(0,3,9,15,16,10,4,1)$. The full spherical rotation and
three-connectivity are checked by the constructor replay.

The graph has $36$ labelled Tait colourings, all in one Kempe orbit.
None connects the designated edges in any two-colour subgraph.
All nine assignments to the two designated colours nevertheless occur,
four times each: colour values alone do not detect the missing
connectivity. Independent edge-order and vertex-order colouring
searches agree. A separate perfect-matching enumeration finds $23$
matchings and $13$ even complementary $2$-factors: eight Hamiltonian
cycles, four factors of cycle lengths $(6,12)$, and one of lengths
$(8,10)$. None contains a cycle through both designated edges.

The Lean proof does not assume completeness of these searches.
After a colour permutation fixes the three colours at vertex $0$,
explicit third-colour deductions classify the six remaining colourings.
For each colouring and each omitted colour, a finite separating edge
set is checked to be closed under actual edge adjacency. Induction on
a connecting walk proves that the two designated edges cannot be
joined. Normalization then removes the initial choice of colour names.
\Sverified{Mettapedia.GraphTheory.FourColor.SmoothingCycleObstruction.normalized_cases}
\Sverified{Mettapedia.GraphTheory.FourColor.SmoothingCycleObstruction.samples_not_joined}
\Sverified{Mettapedia.GraphTheory.FourColor.SmoothingCycleObstruction.exists_normalization}
\Sverified{Mettapedia.GraphTheory.FourColor.SmoothingCycleObstruction.no_coloring_joins}
\Sverified{Mettapedia.GraphTheory.FourColor.SmoothingCycleObstruction.colorable_but_never_joined}
These kernel statements certify literal colouring and reachability;
the embedding and inverse surgery are separate constructor checks.

Inverse insertion is spherical and three-connected, has $20$ vertices,
and admits $72$ labelled Tait colourings. Thus failure to find a common
cycle in \emph{any} colouring of the smaller graph does not imply
non-colourability of the enlarged graph. In particular, changing Kempe
orbit cannot repair this particular sufficient lifting rule.
The same cofacial-cycle assertion is stated as Conjecture~2 in
E.~Kendall, \emph{Simple Numerical Algorithm for Generating Hamiltonian
Cycles and Edge Labels on Planar Cubic Maps} (2021), \S2.3,
\url{https://arxiv.org/abs/2104.05896}; the control rejects that assertion
as stated, not the Four-Color Theorem.

The domain restriction is essential: the smaller graph has triangles
and the enlarged graph has squares. Neither is a least-counterexample
normal-form specimen. Nor is either a zero-Count example. The result
does not refute a lifting theorem that uses the actual source of the
smoothing inside a protected, normal-form ordered mesh. A future such
theorem must retain those hypotheses or use a different lifting
construction; it cannot silently appeal to the false generic cofacial
cycle assertion. No target-facing wide-case implication is discharged,
and the fixed ledger remains $57/100$.

\subsection{Actual smoothing provenance does not select a Kempe orbit}
\label{sec:smoothing-provenance-gate}

The next constructor gate starts with a spherical cubic graph of girth
five whose dual remains connected after every deletion of at most four
vertices. It smooths an actual adjacent pair, preserving the old cyclic
orders. This tests provenance from the stronger geometric population,
not an arbitrary cofacial pair in an unrelated smaller graph. The
original graph is colourable, not a least counterexample; no large
protected ordered mesh is asserted by this control.

The fixed original has $30$ vertices and $96$ labelled Tait colourings.
Deleting edge $\{0,1\}$ and smoothing its endpoints gives a spherical
bridgeless cubic graph with $28$ vertices, girth four, and $192$ labelled
Tait colourings. The two new edges have indices $0,5$ in the recorded
edge list. Its four Kempe orbits have the following exact counts:
\[
\begin{array}{c|rrrr}
 \text{orbit size} & 6 & 36 & 6 & 144\\
 \text{colourings with a common marked cycle} & 6 & 0 & 6 & 72
\end{array}
\]
Two independent complete colouring enumerators agree, and exhaustive
component switching closes these four orbits. Consequently, the claim
that \emph{every} smaller colouring can be repaired by Kempe moves is
false even for this actual smoothing of a normal-form geometric
control. The claim that \emph{some} smaller colouring has the cycle is
not refuted. Nor may normal form be transferred silently to the
smaller map: its squares are created by the surgery.

\paragraph{An orbit-independent construction was attempted.}
Let $M$ be a colour-class perfect matching of a trapped colouring.
Search for a supported matching $N$ such that $M\mathbin{\triangle}N$
is one alternating cycle, the residual factor $H-N$ has only even
cycles, and one residual cycle contains both marked edges. Alternate
two colours around the residual cycles, assign the third colour to
$N$, and apply the literal inverse-insertion lift. This is not an
exchange of two fixed colour classes: the selected alternating carrier
can contain all three old colours, and its resulting colouring can
lie outside the initial Kempe orbit.

The gate constructs such an exchange and an actual lifted colouring
for every one of the $36$ trapped colourings of the fixed control.
The background algebra is already available, rather than newly
assumed: supported matching exchange changes the residual graph by
the selected symmetric difference, and a bipartite residual factor
supplies a Tait colouring.
\Sverified{Mettapedia.GraphTheory.FourColor.GoertzelV24ResidualExchange.exchange_supportedBy_and_residualGraph_eq_symmDiff}
\Sverified{Mettapedia.GraphTheory.FourColor.GoertzelV24ResidualTwoFactor.taitColorable_of_residualGraph_isBipartite}

The attempted existence proof fails at the simultaneous parity and
joining requirement. The smaller control has $102$ perfect matchings:
\[
\begin{array}{c|rr}
 & \text{marks on different cycles or absent} & \text{marks on one cycle}\\
 \text{all residual cycles even} & 45 & 12\\
 \text{two odd residual cycles} & 30 & 15
\end{array}
\]
The receipt includes a single alternating-cycle exchange which joins
the marks but creates two odd residual cycles. Thus graph support,
one connected exchange carrier, and joining do not imply that the
new factor can be alternately coloured. The theorem of A.~Diwan,
\emph{Chords of 2-factors in planar cubic bridgeless graphs} (2021),
\url{https://arxiv.org/abs/2110.12584}, supplies cofacial joining in an
ordinary $2$-factor, not the missing all-even condition. Minimizing
oddness over joining factors does not prove that its minimum is zero.
Indeed the previous $18$-vertex control shows that zero need not be
attained in the broader geometric class.

More explicitly, cutting the old even residual cycles along the
alternating carrier leaves arcs. A new residual cycle has parity equal
to the sum of its surviving arc lengths and its newly inserted
matching edges. Evenness of the old cycles constrains the total of
the old arcs, not each new grouping. In the recorded failed exchange,
the new grouping gives cycle lengths $23$ and $5$. Thus the attempted
proof by ``one even alternating carrier preserves evenness'' fails
at an explicit parity implication, not at a missing exchange API.

\paragraph{Finite reach and remaining obligation.}
The optional gate over the $10$ generated dual $5$-connected controls
of cubic orders $20$ through $30$ checks all $411$ adjacent-pair
smoothings. Some lifting colouring exists in each tested smoothing,
but there are $26$ trapped Kempe orbits containing $936$ colourings.
The matching construction lifts all $936$; the chosen exchange
carriers have at most $24$ edges in this finite run. This is neither a
universal bound nor a proof of the exchange-existence statement.
No frozen exterior or named-colour preservation beyond the explicitly
unchanged part of the lifted cycle is claimed.

The executable attempt and its finite receipt are
\texttt{tools/fourcolor/v24\_smoothing\_provenance\_gate.py} and
\texttt{results/fourcolor/v24\_smoothing\_provenance\_gate.json}.
These new conclusions are constructor and exhaustive finite checks,
not new Lean seals. The source-facing wide-case implication still
requires constructing an admissible reduction with a physical lift
from its exact protected geometry. Replacing that requirement by an
assumed all-even joining factor would only rename it. This is the
second sustained wide-case attack without that implication; the
fixed ledger remains $57/100$, and the proof plan is due for review.

\subsection{Branching is not an escape from bounded-interface generation}
\label{sec:branching-interface-obstruction}

A different proof-plan experiment replaces a corridor by an arbitrary
binary composition tree. The finite-state tree-pumping theorem is already
available; changing the tree shape does not supply a bounded-interface
presentation. The following obstruction concerns the presentation itself,
before any colour-state count or pumping threshold is chosen.

Let every leaf carry at most $a$ ambient vertices, and let every parent
carry the union of its children's carriers. At each subtree require the
\emph{complete ambient edge cut} to have size at most $k$. No connectedness,
strict nesting, disjointness of atoms, or choice of balanced tree is assumed.
For the split-grid family of Section~\ref{sec:split-grid-chain-obstruction},
put $B=2k^2+k$. If $k<n+1$, the existing small-cut theorem says each such
carrier or its complement has at most $B$ vertices. Any binary union-tree
presentation of the whole graph therefore satisfies
\[
  2(n+1)^2 \leq 3\max\{a,2k^2+k\}.
\]
\Sverified{Mettapedia.GraphTheory.FourColor.BranchingInterfaceObstruction.splitGrid_size_bound}

Here is the complete combinatorial reason. Put $t=\max\{a,B\}$. A union
tree with more than $t$ vertices at its root has a subtree $S$ with
$t<|S|\leq2t$: descend into a child larger than $t$ until neither child
is larger, then take their parent. Union subadditivity suffices even
when the atoms overlap. If the whole graph has more than $3t$ vertices,
both $S$ and its complement are larger than $B$, contradicting the
small-cut theorem.
\Sverified{Mettapedia.GraphTheory.FourColor.BranchingInterfaceObstruction.exists_intermediate}
\Sverified{Mettapedia.GraphTheory.FourColor.BranchingInterfaceObstruction.card_le_three_max}

Consequently, for every fixed atom bound $a$, interface width $k$, and
size threshold $N$, an explicit larger split grid admits no such
presentation, whatever the binary branching shape.
\Sverified{Mettapedia.GraphTheory.FourColor.BranchingInterfaceObstruction.no_uniform_bounded_presentation}
The calibrated executable gate checks all binary shapes and leaf
orders through five singleton atoms, overlapping-atom controls, and
all $30{,}192$ intermediate cuts of the $18$-vertex control.
The files are \texttt{tools/fourcolor/v24\_branching\_interface\_gate.py}
and \texttt{results/fourcolor/v24\_branching\_interface\_gate.json}.

\paragraph{Scope and route consequence.}
These are open, properly three-edge-colourable split grids, not least
counterexamples. This is a full vertex-carrier cut theorem, not a
claimed identification with standard edge-middle-set branchwidth or
with every possible graph grammar. It excludes bounded-atom,
bounded-interface generation \emph{when the intermediate pieces are
retained as ambient vertex subsets}. It does not exclude growing
interfaces, semantic rewrites whose intermediate pieces are not such
subsets, or a counterexample-specific structural theorem.

Thus changing a chain to a tree is not by itself the missing mathematical
step. A revised compositional proof must justify its interface bounds
on the actual target class, or explicitly use a different kind of descent.
The source's finite-interface semantics remains distinct from a theorem
that the relevant class has a bounded-interface generating presentation.
No global reduction or finite base is proved here; the fixed ledger
remains $57/100$.

\subsection{Count reconnection: a graph-changing move need not preserve zero}
\label{sec:four-port-count-rewire}

The next experiment changes the graph itself, rather than changing its
decomposition tree. At four cyclic ports, let $I$ be two cubic vertices
joined by an edge, with port pairs $01$ and $23$. Let $H$ be the alternative
connection, with pairs $03$ and $12$. These are not the two wire smoothings
$A=01|23$ and $B=03|12$: the latter contain no cubic vertices.

For a boundary word $w$ let $i(w),h(w),a(w),b(w)$ be the numbers of
extensions into those pieces. Including words containing the zero colour,
the exact identity is
\[
 i(w)+a(w)=h(w)+b(w),
 \qquad q(w)=a(w)+b(w),
\]
where $q(w)$ is the existing square-extension count. Multiplying by any
nonnegative exterior multiplicity $W(w)$ and summing preserves both
identities. In particular, if the contextual $I$ count is zero, the $H$
count is zero exactly when the two contextual smoothing counts agree.
\Sverified{Mettapedia.GraphTheory.FourColor.FourPortCountRewire.weighted_balance}
\Sverified{Mettapedia.GraphTheory.FourColor.FourPortCountRewire.flip_zero_iff_smoothing_balance}
The square can be expanded from zero Count only when \emph{both} smoothing
counts vanish; zero for one branch does not suffice.
\Sverified{Mettapedia.GraphTheory.FourColor.FourPortCountRewire.square_zero_iff_both_smoothings_zero}
This reuses the source's SQ1 count law, not its retired weak-completability
claim. A bijection identifies the two-vertex extension count with the
actual port-tangle colouring fibre, preserving multiplicities.
\Sverified{Mettapedia.GraphTheory.FourColor.FourPortCountRewire.card_cap_col}

\paragraph{An arbitrary-exterior obstruction, not just a numerical test.}
Let $T$ be a finite cubic exterior with four unrestricted ports, and suppose
it has a proper Tait colouring. If the actual serial composite $T\circ I$
is uncolourable, then $T\circ H$ is colourable.
\Sverified{Mettapedia.GraphTheory.FourColor.FourPortCountRewire.flip_colorable_of_original_not_colorable}
The proof needs neither planarity nor connectedness. Cut parity gives
$w_0+w_1+w_2+w_3=0$. On nonzero words satisfying this equation, $I$ extends
exactly when $w_0\ne w_1$, and $H$ extends exactly when $w_0\ne w_3$.
If both closures failed, every exterior colouring would have
$w_0=w_1=w_3$. The already proved physical no-fanout theorem excludes
that: a bichromatic component switch separates at least two of these
three distinct ports.

The nonempty-exterior hypothesis is essential. It is not supplied here
from a hypothetical minimal counterexample. Frozen extra boundaries
would also change the hypothesis. The current formal exterior carriers
are in universe zero, with no numerical bound on their finite sizes.

\paragraph{Nonvacuity, control, and remaining task.}
A kernel-checked control reuses two disjoint digons with one port at each
vertex. Its $I$ closure is uncolourable and its $H$ closure is colourable.
\Sverified{Mettapedia.GraphTheory.FourColor.FourPortCountRewireControl.nonvacuous_rewire_obstruction}
A separate executable control uses two copies of $K_4$ with one edge
removed. Two independent complete colouring enumerators yield
\[
 (\#I,\#H,\#A,\#B,\#Q)=(0,24,36,12,48).
\]
The executable also checks simplicity, cubicity, planarity and bridges;
its $I$ closure has a bridge and its $A$ closure is disconnected. Neither
control is claimed to be a least-counterexample specimen. The files are
\texttt{tools/fourcolor/v24\_count\_rewire\_gate.py} and
\texttt{results/fourcolor/v24\_count\_rewire\_gate.json}.

Thus a chain of elementary two-vertex flips cannot simply be declared
zero-preserving. This does not exclude a composite graph rewrite,
growing interfaces, or a separately justified invariant allowing
colourable intermediate graphs. Such a construction must still recover
an admissible, strictly smaller zero-Count graph; the exact additive
identities alone do not construct it. No global descent or base
obligation is discharged here. The fixed ledger remains $57/100$.
